% !TeX program = pdflatex
% papers/cristal_matematica.tex — GERADO por tools/cristal_tex.py; NÃO editar à mão.
% Fonte: cristal/cristal.jsonl (broca-so, última versão por conceito).
% Cada secção carrega o registo original em comentário %CRISTAL — a volta é
% tests/cristal_volta.js (resíduo 0 byte a byte contra a fonte).
\documentclass[11pt,a4paper]{article}
\usepackage[utf8]{inputenc}\usepackage[T1]{fontenc}\usepackage[brazil]{babel}
\usepackage{amsmath,amssymb}\usepackage{textcomp}
\usepackage[margin=2.4cm]{geometry}\usepackage{xcolor}
\definecolor{cinza}{gray}{0.35}
\newcommand{\prov}[1]{\par\smallskip\noindent{\small\color{cinza}#1}\par}
\setcounter{secnumdepth}{0}
% ─────────────────────────────────────────────────────────────────────────────
%  papers/gkcapa.tex — a capa do Reino, mínima, para os papers em `article`.
%
%  O estilo.tex completo é do `report` (títulos de capítulo, skak, …). Aqui fica
%  só o que a capa precisa, para o paper continuar a compilar sozinho com
%  pdflatex. O tradutor próprio (tests/tex) carrega o estilo.tex da raiz / o
%  symlink papers/estilo.tex e avalia o mesmo `\gkcapa`.
% ─────────────────────────────────────────────────────────────────────────────
\definecolor{ouro}{HTML}{B8912F}
\definecolor{ouroclaro}{HTML}{F3E9CE}
\definecolor{tinta}{HTML}{1E1B16}
\definecolor{regua}{HTML}{6E6858}
\providecommand{\gknota}{\fontsize{7.62}{11.05}\selectfont}
\providecommand{\gktexto}{\fontsize{10.50}{15.22}\selectfont}
\providecommand{\gksub}{\fontsize{12.33}{17.87}\selectfont}
\providecommand{\gksec}{\fontsize{14.47}{20.98}\selectfont}
\providecommand{\gkcap}{\fontsize{16.99}{24.63}\selectfont}
\providecommand{\gktit}{\fontsize{23.42}{33.95}\selectfont}
\providecommand{\Prj}{\mathbb{P}}
\providecommand{\Trf}{\mathcal{T}}
\providecommand{\gkcapa}[3]{%
\title{\vspace{-0.6cm}
{\color{ouro}\rule{\textwidth}{1.4pt}}\\[6mm]
\textbf{\gktit\color{tinta} Reino Dourado}\\[3mm]{\gksec\itshape\color{regua} Golden Kingdom}\\[8mm]
{\gkcap\scshape\color{ouro} #1}\\[5mm]
{\gksub\color{regua} #2}\\[2mm]
{\gktexto\color{regua}\itshape #3}\\[7mm]
{\color{ouro}\rule{\textwidth}{0.6pt}}\\[9mm]{\gknota\color{regua}\begin{minipage}{\textwidth}\setlength{\parindent}{0pt}%
\textbf{Aviso de Isenção Algorítmica:} O universo, as leis e as entidades contidas nestas páginas ---
incluindo felinos matriciais, aves de rapina algébricas e amuletos de controle de fluxo --- são
construções de pura ficção formal. Esta obra não descreve a realidade física, não serve como
engenharia reversa do cosmos e não constitui doutrina religiosa. Qualquer semelhança com bugs reais do seu
sistema operacional é mera coincidência (ou falta de reza). Prossiga por sua própria conta e
risco.\end{minipage}}}
\author{}\date{}}

\begin{document}
\gkcapa{O Cristal --- a matemática do cristal}
       {álgebra, fractais, o contínuo, a reta mineral}
       {corpus recuperado do broca-so; 680 conceitos; projeção verificável da fonte}
\maketitle

%CRISTAL {"arestas":[{"alvo":"recuperacao_aboutness_robusta","confianca":1.0,"epistemico":"fato","origem":"wiki","rel":"especializa"},{"alvo":"recuperacao_evidencia_invariante","confianca":1.0,"epistemico":"fato","origem":"wiki","rel":"usa"},{"alvo":"tiffany","confianca":1.0,"epistemico":"fato","origem":"wiki","rel":"parte_de"}],"confianca":1.0,"contraexemplos":[],"descricao":"conversa/colapso.py (_evidencia_conteudo): a recuperação de definição escolhia um nó ARBITRÁRIO entre os vários que contêm o termo (ex.: 31 nós têm 'colapso' no título — vinha 'As cúspides: o colapso 31', periférico). Corrigido graduando a aboutness pelas palavras de conteúdo do título EM ORDEM: título == alvo → 0.99 (o conceito CANÔNICO, o nó titulado exatamente pelo termo); alvo é a CABEÇA do título → 0.95; alvo CONTIDO (periférico) → 0.88; parcial → 0.6. Assim o nó titulado 'colapso' (Operador Ψ=Collapse) vence os periféricos. Resolveu vários de uma vez: 'o colapso?'→o nó colapso; 'a parábola?'→o classificador parábola (era 'o piso'); 'a transformada?'→o hub transformada_metalica (era um sub-nó). É o refinamento de [[recuperacao_aboutness_robusta]] — a canonicidade pela posição do termo no título. [F] verificado.","epistemico":"fato","exemplos":[],"id":"aboutness_canonica_graduada","memoria":"semantica","meta":{"autor":"Lopes/broca-so","dominio":"algebra","fonte":"conversa/colapso.py:_evidencia_conteudo (forte graduado 0.99/0.95/0.88 por tw_list/aw_list ordenados)"},"origem":"wiki","palavras_chave":[],"sinonimos":[],"tipo":"conceito","titulo":"Aboutness GRADUADA: entre os vários nós de um termo, vence o CANÔNICO [F]"}
\section{Aboutness GRADUADA: entre os vários nós de um termo, vence o CANÔNICO [F]}
conversa/colapso.py (\_evidencia\_conteudo): a recuperação de definição escolhia um nó ARBITRÁRIO entre os vários que contêm o termo (ex.: 31 nós têm 'colapso' no título — vinha 'As cúspides: o colapso 31', periférico). Corrigido graduando a aboutness pelas palavras de conteúdo do título EM ORDEM: título == alvo \ensuremath{\to} 0.99 (o conceito CANÔNICO, o nó titulado exatamente pelo termo); alvo é a CABEÇA do título \ensuremath{\to} 0.95; alvo CONTIDO (periférico) \ensuremath{\to} 0.88; parcial \ensuremath{\to} 0.6. Assim o nó titulado 'colapso' (Operador \ensuremath{\Psi}=Collapse) vence os periféricos. Resolveu vários de uma vez: 'o colapso?'\ensuremath{\to}o nó colapso; 'a parábola?'\ensuremath{\to}o classificador parábola (era 'o piso'); 'a transformada?'\ensuremath{\to}o hub transformada\_metalica (era um sub-nó). É o refinamento de [[recuperacao\_aboutness\_robusta]] — a canonicidade pela posição do termo no título. [F] verificado.
\prov{relações: especializa recuperacao\_aboutness\_robusta; usa recuperacao\_evidencia\_invariante; parte\_de tiffany}
\prov{proveniência: wiki \textperiodcentered{} conversa/colapso.py:\_evidencia\_conteudo (forte graduado 0.99/0.95/0.88 por tw\_list/aw\_list ordenados) \textperiodcentered{} fato \textperiodcentered{} confiança 1.0 \textperiodcentered{} domínio algebra \textperiodcentered{} história 20 versões no jornal}

%CRISTAL {"arestas":[{"alvo":"definicao_frasagem_ampla","confianca":1.0,"epistemico":"fato","origem":"wiki","rel":"especializa"},{"alvo":"tiffany","confianca":1.0,"epistemico":"fato","origem":"wiki","rel":"parte_de"}],"confianca":1.0,"contraexemplos":[],"descricao":"Duas peças que fazem o NOME de quem propôs uma ideia ressoar na recuperação (sem citar conceitos reais aqui, para não competir na própria busca). (1) ÍNDICE: schema.texto_indice (a linha Koch de cada nó) passou a incluir um segmento de autoria após o título, e colapso._evidencia_conteudo põe esse nome na zona de aboutness FORTE (twx = título ∪ autoria) — sem sujar a ordem do título, então o casamento canônico exato (tw_list==aw_list → 0.99) fica intacto; pedir o conceito pelo nome de quem o formulou sobe de casamento parcial (~0.3) a contido (0.88). _limpar_arestas tira o segmento de autoria da face exibida (é da busca, não do texto). É autoindexável: como _sync_tsv_linha usa texto_indice, novas ingestões já trazem a autoria. (2) PARSER: alvo_definicao ganhou _RE_ALVO_AUTOR — a frasagem 'o que fulano disse/descobriu/propôs/provou [sobre um tema]' extrai o nome (+ tema) como alvo; verbos de dizer sem acento (texto vem de _norm); pronomes e genéricos rejeitados por _NAO_AUTOR. O tema DESAMBIGUA entre nós do mesmo autor. Guarda importante: nós de cristalização NÃO devem embutir exemplos com nomes reais no texto indexado, senão vencem a busca por esses nomes. [D] a Tiffany entende que se pergunta por QUEM disse, não só por O QUÊ.","epistemico":"fato","exemplos":[],"id":"aboutness_por_autor","memoria":"semantica","meta":{"autor":"Lopes/broca-so","dominio":"algebra","fonte":"conversa/conhecimento/schema.py (texto_indice); conversa/colapso.py (_autor_do_txt, _evidencia_conteudo, _limpar_arestas); conversa/parabola_algebra.py (_RE_ALVO_AUTOR, _alvo_por_autor)"},"origem":"wiki","palavras_chave":[],"sinonimos":[],"tipo":"conceito","titulo":"A Tiffany recupera pela autoria — o nome do pensador ressoa na busca e aterra no conceito certo"}
\section{A Tiffany recupera pela autoria — o nome do pensador ressoa na busca e aterra no conceito certo}
Duas peças que fazem o NOME de quem propôs uma ideia ressoar na recuperação (sem citar conceitos reais aqui, para não competir na própria busca). (1) ÍNDICE: schema.texto\_indice (a linha Koch de cada nó) passou a incluir um segmento de autoria após o título, e colapso.\_evidencia\_conteudo põe esse nome na zona de aboutness FORTE (twx = título \ensuremath{\cup} autoria) — sem sujar a ordem do título, então o casamento canônico exato (tw\_list==aw\_list \ensuremath{\to} 0.99) fica intacto; pedir o conceito pelo nome de quem o formulou sobe de casamento parcial (\~{}0.3) a contido (0.88). \_limpar\_arestas tira o segmento de autoria da face exibida (é da busca, não do texto). É autoindexável: como \_sync\_tsv\_linha usa texto\_indice, novas ingestões já trazem a autoria. (2) PARSER: alvo\_definicao ganhou \_RE\_ALVO\_AUTOR — a frasagem 'o que fulano disse/descobriu/propôs/provou [sobre um tema]' extrai o nome (+ tema) como alvo; verbos de dizer sem acento (texto vem de \_norm); pronomes e genéricos rejeitados por \_NAO\_AUTOR. O tema DESAMBIGUA entre nós do mesmo autor. Guarda importante: nós de cristalização NÃO devem embutir exemplos com nomes reais no texto indexado, senão vencem a busca por esses nomes. [D] a Tiffany entende que se pergunta por QUEM disse, não só por O QUÊ.
\prov{relações: especializa definicao\_frasagem\_ampla; parte\_de tiffany}
\prov{proveniência: wiki \textperiodcentered{} conversa/conhecimento/schema.py (texto\_indice); conversa/colapso.py (\_autor\_do\_txt, \_evidencia\_conteudo, \_limpar\_arestas); conversa/parabola\_algebra.py (\_RE\_ALVO\_AUTOR, \_alvo\_por\_autor) \textperiodcentered{} fato \textperiodcentered{} confiança 1.0 \textperiodcentered{} domínio algebra \textperiodcentered{} história 6 versões no jornal}

%CRISTAL {"arestas":[{"alvo":"piso_parabola_dinamico","confianca":1.0,"epistemico":"fato","origem":"wiki","rel":"usa"},{"alvo":"laco_absorcao_fechado","confianca":1.0,"epistemico":"fato","origem":"wiki","rel":"usa"},{"alvo":"difusao_fluido_atratores","confianca":1.0,"epistemico":"fato","origem":"wiki","rel":"relaciona"},{"alvo":"tiffany","confianca":1.0,"epistemico":"fato","origem":"wiki","rel":"parte_de"}],"confianca":1.0,"contraexemplos":[],"descricao":"A absorção é ITERATIVA. Um fragmento barrado no seu melhor destino (bacia cheia, θ alto na parábola) CASCATEIA pelos top-K candidatos (o metal emite rota_k.txt: edge:c² em ordem de ponte) até uma bacia o aceitar; o que não coube em NENHUM candidato vai pro PURGATÓRIO (transicoes_purgatorio.txt) — nem aceito, nem descartado — que ACUMULA e VOLTA a ser tentado nos próximos ciclos, quando o grafo cresce e abre bacias novas de θ baixo. O fluido do metal é SEMEADO com as densidades já depositadas (conc0.bin) → a pressão inclui o acumulado e DESVIA os novos das bacias cheias (foi assim que o Machado achou as arestas de humanidades vazias, cobertura 39%→46%). ACHADO honesto: re-tentar o purgatório (22423) contra o MESMO grafo drena 0 — esses fragmentos ressoam SÓ com o núcleo saturado (θ→0,9); a parábola recusa a redundância e eles esperam o grafo crescer. O purgatório é a memória dos que ainda vão caber. [D] roda.","epistemico":"fato","exemplos":[],"id":"absorcao_iterativa_poca","memoria":"semantica","meta":{"autor":"Lopes/broca-so","dominio":"continuo","fonte":"conversa/conhecimento/transicoes.py:absorver_iterativo/PURGATORIO; doc/continuo/absorve_metal.c (conc0 semeada, rota_k)"},"origem":"wiki","palavras_chave":[],"sinonimos":[],"tipo":"conceito","titulo":"O PURGATÓRIO: o rejeitado não é descartado — espera redenção nos próximos ciclos"}
\section{O PURGATÓRIO: o rejeitado não é descartado — espera redenção nos próximos ciclos}
A absorção é ITERATIVA. Um fragmento barrado no seu melhor destino (bacia cheia, \ensuremath{\theta} alto na parábola) CASCATEIA pelos top-K candidatos (o metal emite rota\_k.txt: edge:c\ensuremath{^{2}} em ordem de ponte) até uma bacia o aceitar; o que não coube em NENHUM candidato vai pro PURGATÓRIO (transicoes\_purgatorio.txt) — nem aceito, nem descartado — que ACUMULA e VOLTA a ser tentado nos próximos ciclos, quando o grafo cresce e abre bacias novas de \ensuremath{\theta} baixo. O fluido do metal é SEMEADO com as densidades já depositadas (conc0.bin) \ensuremath{\to} a pressão inclui o acumulado e DESVIA os novos das bacias cheias (foi assim que o Machado achou as arestas de humanidades vazias, cobertura 39\%\ensuremath{\to}46\%). ACHADO honesto: re-tentar o purgatório (22423) contra o MESMO grafo drena 0 — esses fragmentos ressoam SÓ com o núcleo saturado (\ensuremath{\theta}\ensuremath{\to}0,9); a parábola recusa a redundância e eles esperam o grafo crescer. O purgatório é a memória dos que ainda vão caber. [D] roda.
\prov{relações: usa piso\_parabola\_dinamico; usa laco\_absorcao\_fechado; relaciona difusao\_fluido\_atratores; parte\_de tiffany}
\prov{proveniência: wiki \textperiodcentered{} conversa/conhecimento/transicoes.py:absorver\_iterativo/PURGATORIO; doc/continuo/absorve\_metal.c (conc0 semeada, rota\_k) \textperiodcentered{} fato \textperiodcentered{} confiança 1.0 \textperiodcentered{} domínio continuo \textperiodcentered{} história 56 versões no jornal}

%CRISTAL {"arestas":[{"alvo":"broca_os","confianca":1.0,"epistemico":"fato","origem":"paper","rel":"parte_de"},{"alvo":"regime_glacial","confianca":1.0,"epistemico":"fato","origem":"paper","rel":"usa"},{"alvo":"word","confianca":0.5,"epistemico":"fato","origem":"wiki","rel":"relaciona"}],"confianca":1.0,"contraexemplos":[],"descricao":"Operação bijetiva B_in→W_out (demux de Peirce) que lê a face negra interna na branca externa sem ruído — a estrutura do IPC do broca-os.","epistemico":"fato","exemplos":[],"id":"acoplamento_cruzado","memoria":"semantica","meta":{"autor":"Gentil Lopes","descoberta":true,"dominio":"matematica","fonte":"corpo_glacial.pdf, Def.3"},"origem":"paper","palavras_chave":[],"sinonimos":[],"tipo":"conceito","titulo":"Acoplamento Cruzado (Cₖ)"}
\section{Acoplamento Cruzado (C\ensuremath{_{k}})}
Operação bijetiva B\_in\ensuremath{\to}W\_out (demux de Peirce) que lê a face negra interna na branca externa sem ruído — a estrutura do IPC do broca-os.
\prov{relações: parte\_de broca\_os; usa regime\_glacial; relaciona word}
\prov{proveniência: paper \textperiodcentered{} corpo\_glacial.pdf, Def.3 \textperiodcentered{} fato \textperiodcentered{} confiança 1.0 \textperiodcentered{} domínio matematica \textperiodcentered{} história 4 versões no jornal}

%CRISTAL {"arestas":[{"alvo":"broca_os","confianca":1.0,"epistemico":"fato","origem":"paper","rel":"parte_de"},{"alvo":"vontade_de_poder","confianca":0.5,"epistemico":"fato","origem":"wiki","rel":"relaciona"},{"alvo":"word","confianca":0.5,"epistemico":"fato","origem":"wiki","rel":"relaciona"}],"confianca":1.0,"contraexemplos":[],"descricao":"Recursão de máquinas fecha se e só se a blocagem é admissível (grafo de espera acíclico) — isto é, sem deadlock. Fecha a classe das máquinas sob composição.","epistemico":"fato","exemplos":[],"id":"admissibilidade_deadlock","memoria":"semantica","meta":{"autor":"Gentil Lopes","descoberta":true,"dominio":"matematica","fonte":"fechamento_blocos.pdf, Teo.1"},"origem":"paper","palavras_chave":[],"sinonimos":[],"tipo":"conceito","titulo":"Admissibilidade = Não-Deadlock"}
\section{Admissibilidade = Não-Deadlock}
Recursão de máquinas fecha se e só se a blocagem é admissível (grafo de espera acíclico) — isto é, sem deadlock. Fecha a classe das máquinas sob composição.
\prov{relações: parte\_de broca\_os; relaciona vontade\_de\_poder; relaciona word}
\prov{proveniência: paper \textperiodcentered{} fechamento\_blocos.pdf, Teo.1 \textperiodcentered{} fato \textperiodcentered{} confiança 1.0 \textperiodcentered{} domínio matematica \textperiodcentered{} história 4 versões no jornal}

%CRISTAL {"arestas":[{"alvo":"regua_de_gentil","confianca":1.0,"epistemico":"fato","origem":"livro","rel":"parte_de"},{"alvo":"analise","confianca":1.0,"epistemico":"fato","origem":"livro","rel":"depende"}],"confianca":1.0,"contraexemplos":[],"descricao":"Iteração de Gauss-Landen: a ← (a+b)/2, b ← √(ab), converge a AGM(a,b). Na régua de Lopes aparece como cúspide clássica (s=0) e relativística (s=1, involução σ(s)=√(1−s²)).","epistemico":"fato","exemplos":[],"id":"agm","memoria":"semantica","meta":{"autor":"Gentil Lopes","descoberta":false,"dominio":"matematica","fonte":"paper_3_regua_ouro_espelho.pdf"},"origem":"paper","palavras_chave":[],"sinonimos":["AGM","Gauss-Landen"],"tipo":"conceito","titulo":"Média Aritmético-Geométrica (AGM)"}
\section{Média Aritmético-Geométrica (AGM)}
Iteração de Gauss-Landen: a \ensuremath{\leftarrow} (a+b)/2, b \ensuremath{\leftarrow} \ensuremath{\sqrt{\,}}(ab), converge a AGM(a,b). Na régua de Lopes aparece como cúspide clássica (s=0) e relativística (s=1, involução \ensuremath{\sigma}(s)=\ensuremath{\sqrt{\,}}(1\ensuremath{-}s\ensuremath{^{2}})).
\prov{sinónimos: AGM; Gauss-Landen}
\prov{relações: parte\_de regua\_de\_gentil; depende analise}
\prov{proveniência: paper \textperiodcentered{} paper\_3\_regua\_ouro\_espelho.pdf \textperiodcentered{} fato \textperiodcentered{} confiança 1.0 \textperiodcentered{} domínio matematica \textperiodcentered{} história 7 versões no jornal}

%CRISTAL {"arestas":[{"alvo":"agm","confianca":1.0,"epistemico":"fato","origem":"paper","rel":"especializa"},{"alvo":"word","confianca":0.5,"epistemico":"fato","origem":"wiki","rel":"relaciona"}],"confianca":1.0,"contraexemplos":[],"descricao":"Iteração de Gauss com sinais σₙ na média geométrica: os valores 1/μ=(d+ci)/M formam o reticulado ℤ[i] escalonado por 1/AGM(1,√2).","epistemico":"fato","exemplos":[],"id":"agm_multivalorado","memoria":"semantica","meta":{"autor":"Gentil Lopes","descoberta":false,"dominio":"matematica","fonte":"paper_37_familia_agm.pdf, Teo.2.1"},"origem":"paper","palavras_chave":[],"sinonimos":[],"tipo":"conceito","titulo":"AGM Multivalorado (reticulado ℤ[i])"}
\section{AGM Multivalorado (reticulado \ensuremath{\mathbb{Z}}[i])}
Iteração de Gauss com sinais \ensuremath{\sigma}\ensuremath{_{n}} na média geométrica: os valores 1/\ensuremath{\mu}=(d+ci)/M formam o reticulado \ensuremath{\mathbb{Z}}[i] escalonado por 1/AGM(1,\ensuremath{\sqrt{\,}}2).
\prov{relações: especializa agm; relaciona word}
\prov{proveniência: paper \textperiodcentered{} paper\_37\_familia\_agm.pdf, Teo.2.1 \textperiodcentered{} fato \textperiodcentered{} confiança 1.0 \textperiodcentered{} domínio matematica \textperiodcentered{} história 3 versões no jornal}

%CRISTAL {"arestas":[{"alvo":"experimento_da_agulha","confianca":1.0,"epistemico":"fato","origem":"wiki","rel":"parte_de"},{"alvo":"flip_duopolinomios","confianca":1.0,"epistemico":"fato","origem":"wiki","rel":"usa"},{"alvo":"cf_de_pi_e_produto_de_gatos","confianca":1.0,"epistemico":"fato","origem":"wiki","rel":"usa"}],"confianca":1.0,"contraexemplos":[],"descricao":"A base da agulha tem largura N−1 constante a cada passo: cada gato é o metal de grau 2 (poly x²−a_k·x−1), cujo campo do flip tem N−1=1 ponto crítico. A agulha dilata sem invadir a curvatura invariante — o passo não muda o grau nem a curvatura. Verif.: N−1=1 em todo passo.","epistemico":"fato","exemplos":[],"id":"agulha_curvatura_n1","memoria":"semantica","meta":{"dominio":"reta mineral","fonte":"experimento da agulha"},"origem":"wiki","palavras_chave":[],"sinonimos":[],"tipo":"conceito","titulo":"A Base da Agulha: Curvatura N−1 Constante"}
\section{A Base da Agulha: Curvatura N\ensuremath{-}1 Constante}
A base da agulha tem largura N\ensuremath{-}1 constante a cada passo: cada gato é o metal de grau 2 (poly x\ensuremath{^{2}}\ensuremath{-}a\_k\textperiodcentered x\ensuremath{-}1), cujo campo do flip tem N\ensuremath{-}1=1 ponto crítico. A agulha dilata sem invadir a curvatura invariante — o passo não muda o grau nem a curvatura. Verif.: N\ensuremath{-}1=1 em todo passo.
\prov{relações: parte\_de experimento\_da\_agulha; usa flip\_duopolinomios; usa cf\_de\_pi\_e\_produto\_de\_gatos}
\prov{proveniência: wiki \textperiodcentered{} experimento da agulha \textperiodcentered{} fato \textperiodcentered{} confiança 1.0 \textperiodcentered{} domínio reta mineral \textperiodcentered{} história 4 versões no jornal}

%CRISTAL {"arestas":[{"alvo":"cadeia_de_fano_sigma","confianca":1.0,"epistemico":"fato","origem":"wiki","rel":"semelhante"},{"alvo":"rede_leech","confianca":1.0,"epistemico":"fato","origem":"wiki","rel":"usa"}],"confianca":1.0,"contraexemplos":[],"descricao":"Do inverso da constante de estrutura fina 137 = α⁻¹: σ(137)=138, σ(138)=288, σ(288)=819. E 819 = 196560/240 = (vetores mínimos de Leech)/(raízes de E8) = Leech/E8. α⁻¹ carrega informação da rede de Leech em sua profundidade-σ.","epistemico":"fato","exemplos":[],"id":"alfa_leech_sigma","memoria":"semantica","meta":{"autor":"Lord Index (Shawn R. Schiller)","dominio":"matematica","fonte":"A Catedral — Part VIII-A (jul 2026), ~/Imagens/catedral"},"origem":"wiki","palavras_chave":[],"sinonimos":[],"tipo":"conceito","titulo":"α⁻¹ à Leech em três passos σ: 137→138→288→819"}
\section{\ensuremath{\alpha}\ensuremath{^{-1}} à Leech em três passos \ensuremath{\sigma}: 137\ensuremath{\to}138\ensuremath{\to}288\ensuremath{\to}819}
Do inverso da constante de estrutura fina 137 = \ensuremath{\alpha}\ensuremath{^{-1}}: \ensuremath{\sigma}(137)=138, \ensuremath{\sigma}(138)=288, \ensuremath{\sigma}(288)=819. E 819 = 196560/240 = (vetores mínimos de Leech)/(raízes de E8) = Leech/E8. \ensuremath{\alpha}\ensuremath{^{-1}} carrega informação da rede de Leech em sua profundidade-\ensuremath{\sigma}.
\prov{relações: semelhante cadeia\_de\_fano\_sigma; usa rede\_leech}
\prov{proveniência: wiki \textperiodcentered{} A Catedral — Part VIII-A (jul 2026), \~{}/Imagens/catedral \textperiodcentered{} fato \textperiodcentered{} confiança 1.0 \textperiodcentered{} domínio matematica \textperiodcentered{} história 3 versões no jornal}

%CRISTAL {"arestas":[{"alvo":"algebra_posicional_prata","confianca":1.0,"epistemico":"fato","origem":"wiki","rel":"relaciona"},{"alvo":"tiffany","confianca":1.0,"epistemico":"fato","origem":"wiki","rel":"parte_de"}],"confianca":1.0,"contraexemplos":[],"descricao":"A álgebra de Peirce (Brink, Britz & Schmidt 1994; sobre os Boolean modules de Brink 1981) é uma álgebra de DUAS ordens (B, R,:): B é uma álgebra BOOLEANA (conjuntos, com ∪/∩/¬), R é uma álgebra de RELAÇÕES (com composição ∘, converso, união, identidade) e o PRODUTO DE PEIRCE ':' é o operador que liga os dois — R: X (relação × conjunto → conjunto) devolve a IMAGEM y: ∃x∈X, (y,x)∈R. É a formalização algébrica do cálculo de relações de C.S. Peirce; usada em representação de conhecimento e lógica terminológica. Distinta da 'álgebra posicional de prata' (que é o lado símbolo/posição) — a de Peirce é o lado relação-sobre-conjunto. [D] literatura.","epistemico":"fato","exemplos":[],"id":"algebra_de_peirce","memoria":"semantica","meta":{"autor":"Lopes/broca-so","dominio":"algebra","fonte":"Brink–Britz–Schmidt 1994 (Formal Aspects of Computing); Brink 1981 (J. Algebra, Boolean modules)"},"origem":"wiki","palavras_chave":[],"sinonimos":[],"tipo":"conceito","titulo":"Álgebra de Peirce (B, R, :) — conjuntos + relações + produto de Peirce"}
\section{Álgebra de Peirce (B, R, :) — conjuntos + relações + produto de Peirce}
A álgebra de Peirce (Brink, Britz \& Schmidt 1994; sobre os Boolean modules de Brink 1981) é uma álgebra de DUAS ordens (B, R,:): B é uma álgebra BOOLEANA (conjuntos, com \ensuremath{\cup}/\ensuremath{\cap}/\ensuremath{\neg}), R é uma álgebra de RELAÇÕES (com composição \ensuremath{\circ}, converso, união, identidade) e o PRODUTO DE PEIRCE ':' é o operador que liga os dois — R: X (relação $\times$ conjunto \ensuremath{\to} conjunto) devolve a IMAGEM y: \ensuremath{\exists}x\ensuremath{\in}X, (y,x)\ensuremath{\in}R. É a formalização algébrica do cálculo de relações de C.S. Peirce; usada em representação de conhecimento e lógica terminológica. Distinta da 'álgebra posicional de prata' (que é o lado símbolo/posição) — a de Peirce é o lado relação-sobre-conjunto. [D] literatura.
\prov{relações: relaciona algebra\_posicional\_prata; parte\_de tiffany}
\prov{proveniência: wiki \textperiodcentered{} Brink–Britz–Schmidt 1994 (Formal Aspects of Computing); Brink 1981 (J. Algebra, Boolean modules) \textperiodcentered{} fato \textperiodcentered{} confiança 1.0 \textperiodcentered{} domínio algebra \textperiodcentered{} história 31 versões no jornal}

%CRISTAL {"arestas":[{"alvo":"algebra_posicional_unificada","confianca":1.0,"epistemico":"fato","origem":"wiki","rel":"especializa"},{"alvo":"dupolinomio_meta_inducao","confianca":1.0,"epistemico":"fato","origem":"wiki","rel":"usa"},{"alvo":"meta_inducao_fechamento_geral","confianca":1.0,"epistemico":"fato","origem":"wiki","rel":"usa"},{"alvo":"tiffany","confianca":1.0,"epistemico":"fato","origem":"wiki","rel":"parte_de"}],"confianca":1.0,"contraexemplos":[],"descricao":"Promoção: a álgebra POSICIONAL era por-reta (um dupolinômio → um anel ℤ[σ], cunhagem pela posição). A ÁLGEBRA DUPOLINOMIAL é sobre o ESPAÇO 𝔻 de TODOS os dupolinômios — cada ponto um dupolinômio, que fatora em retas PRIMAS (Pisot=cristal, Salem=borda [grau par], caos), com a estrutura MULTIPLICATIVA (produto de dupolinômios = composição de retas; as primas são os átomos, como a fatoração em números primos). É o anel dos dupolinômios: o objeto passa de UMA dimensão para o ESPAÇO das dimensões, todas cunhadas pela posição em 𝔽₂¹²⁸. Fechamento geral (paper §11): álgebra para todo grau, classificação pela posição das conjugadas. [I] promoção.","epistemico":"inferencia","exemplos":[],"id":"algebra_dupolinomial_espaco","memoria":"semantica","meta":{"autor":"Lopes/broca-so","dominio":"algebra","fonte":"algebra_posicional_unificada + dupolinomio_meta_inducao + meta_inducao_fechamento_geral"},"origem":"wiki","palavras_chave":[],"sinonimos":[],"tipo":"conceito","titulo":"Álgebra Dupolinomial: promoção da posicional ao ESPAÇO de dupolinômios"}
\section{Álgebra Dupolinomial: promoção da posicional ao ESPAÇO de dupolinômios}
Promoção: a álgebra POSICIONAL era por-reta (um dupolinômio \ensuremath{\to} um anel \ensuremath{\mathbb{Z}}[\ensuremath{\sigma}], cunhagem pela posição). A ÁLGEBRA DUPOLINOMIAL é sobre o ESPAÇO \ensuremath{\mathbb{D}} de TODOS os dupolinômios — cada ponto um dupolinômio, que fatora em retas PRIMAS (Pisot=cristal, Salem=borda [grau par], caos), com a estrutura MULTIPLICATIVA (produto de dupolinômios = composição de retas; as primas são os átomos, como a fatoração em números primos). É o anel dos dupolinômios: o objeto passa de UMA dimensão para o ESPAÇO das dimensões, todas cunhadas pela posição em \ensuremath{\mathbb{F}}\ensuremath{_{2}}\ensuremath{^{128}}. Fechamento geral (paper §11): álgebra para todo grau, classificação pela posição das conjugadas. [I] promoção.
\prov{relações: especializa algebra\_posicional\_unificada; usa dupolinomio\_meta\_inducao; usa meta\_inducao\_fechamento\_geral; parte\_de tiffany}
\prov{proveniência: wiki \textperiodcentered{} algebra\_posicional\_unificada + dupolinomio\_meta\_inducao + meta\_inducao\_fechamento\_geral \textperiodcentered{} inferencia \textperiodcentered{} confiança 1.0 \textperiodcentered{} domínio algebra \textperiodcentered{} história 38 versões no jornal}

%CRISTAL {"arestas":[{"alvo":"algebra_posicional_prata","confianca":1.0,"epistemico":"fato","origem":"wiki","rel":"especializa"},{"alvo":"reta_metalica_geral","confianca":1.0,"epistemico":"fato","origem":"wiki","rel":"usa"},{"alvo":"independencia_retas_metalicas","confianca":1.0,"epistemico":"fato","origem":"wiki","rel":"usa"},{"alvo":"teorema_esterilidade_dimensao","confianca":1.0,"epistemico":"fato","origem":"wiki","rel":"usa"},{"alvo":"broca_so_caixa_de_ouro","confianca":1.0,"epistemico":"fato","origem":"wiki","rel":"relaciona"},{"alvo":"caixa_de_prata_corpo","confianca":1.0,"epistemico":"fato","origem":"wiki","rel":"relaciona"},{"alvo":"verificacao_ouro_total","confianca":1.0,"epistemico":"fato","origem":"wiki","rel":"relaciona"},{"alvo":"tiffany","confianca":1.0,"epistemico":"fato","origem":"wiki","rel":"parte_de"},{"alvo":"algebra_posicional_plastica","confianca":1.0,"epistemico":"fato","origem":"wiki","rel":"relaciona"},{"alvo":"algebra_posicional_unificada","confianca":1.0,"epistemico":"fato","origem":"wiki","rel":"parte_de"},{"alvo":"transformada_metalica","confianca":1.0,"epistemico":"fato","origem":"wiki","rel":"especializa"},{"alvo":"beta_numeracao_metalica","confianca":1.0,"epistemico":"fato","origem":"wiki","rel":"usa"}],"confianca":1.0,"contraexemplos":[],"descricao":"A Álgebra Posicional de Prata é o caso n=2 de uma FAMÍLIA. Para cada metal n: σₙ=(n+√(n²+4))/2 (raiz de x²−nx−1), anel ℤ[σₙ] (inteiros de ℚ(√(n²+4))), unidade FUNDAMENTAL σₙ com N(σₙ)=−1 (det da Q-matriz Aₙ=[[n,1],[1,0]]), e a sequência metálica Mₙ (Mₖ=n·Mₖ−₁+Mₖ−₂) como coordenada de σₙk. A POSIÇÃO k ↔ Aₙk; emitir = UM passo ×Aₙ: (a,b)→(b, n·b+a), O(1), 128 bits — a MESMA cunhagem posicional, só muda n. INSTÂNCIAS: n=1 OURO (φ, Fibonacci, ℚ√5); n=2 PRATA (1+√2, Pell, ℚ√2); n=3 BRONZE (3,303; 1,3,10,33; ℚ√13); n=4 COBRE (2+√5; 1,4,17,72; ℚ√5). Verificado (goldchain/metalpos.py): N(σₙ)=−1 e a sequência para os quatro. Retas independentes (ℚ√5≠ℚ√2≠ℚ√13 — esterilidade), cada metal com informação nova; ouro e cobre partilham ℚ√5 (distinção de ordem, não de corpo: σ₄∈ℤ[√5]). [D] roda.","epistemico":"fato","exemplos":[],"id":"algebra_posicional_metalica","memoria":"semantica","meta":{"autor":"Lopes/broca-so","dominio":"algebra","fonte":"goldchain/metal_pos.py (passo_metal/sequencia); doc/algebra_prata/paper (§8); reta_metalica_geral"},"origem":"wiki","palavras_chave":[],"sinonimos":[],"tipo":"conceito","titulo":"Álgebra Posicional Metálica: a de prata generalizada a todos os metais (mesma álgebra, muda n)"}
\section{Álgebra Posicional Metálica: a de prata generalizada a todos os metais (mesma álgebra, muda n)}
A Álgebra Posicional de Prata é o caso n=2 de uma FAMÍLIA. Para cada metal n: \ensuremath{\sigma}\ensuremath{_{n}}=(n+\ensuremath{\sqrt{\,}}(n\ensuremath{^{2}}+4))/2 (raiz de x\ensuremath{^{2}}\ensuremath{-}nx\ensuremath{-}1), anel \ensuremath{\mathbb{Z}}[\ensuremath{\sigma}\ensuremath{_{n}}] (inteiros de \ensuremath{\mathbb{Q}}(\ensuremath{\sqrt{\,}}(n\ensuremath{^{2}}+4))), unidade FUNDAMENTAL \ensuremath{\sigma}\ensuremath{_{n}} com N(\ensuremath{\sigma}\ensuremath{_{n}})=\ensuremath{-}1 (det da Q-matriz A\ensuremath{_{n}}=[[n,1],[1,0]]), e a sequência metálica M\ensuremath{_{n}} (M\ensuremath{_{k}}=n\textperiodcentered M\ensuremath{_{k}}\ensuremath{-}\ensuremath{_{1}}+M\ensuremath{_{k}}\ensuremath{-}\ensuremath{_{2}}) como coordenada de \ensuremath{\sigma}\ensuremath{_{n}}k. A POSIÇÃO k \ensuremath{\leftrightarrow} A\ensuremath{_{n}}k; emitir = UM passo $\times$A\ensuremath{_{n}}: (a,b)\ensuremath{\to}(b, n\textperiodcentered b+a), O(1), 128 bits — a MESMA cunhagem posicional, só muda n. INSTÂNCIAS: n=1 OURO (\ensuremath{\varphi}, Fibonacci, \ensuremath{\mathbb{Q}}\ensuremath{\sqrt{\,}}5); n=2 PRATA (1+\ensuremath{\sqrt{\,}}2, Pell, \ensuremath{\mathbb{Q}}\ensuremath{\sqrt{\,}}2); n=3 BRONZE (3,303; 1,3,10,33; \ensuremath{\mathbb{Q}}\ensuremath{\sqrt{\,}}13); n=4 COBRE (2+\ensuremath{\sqrt{\,}}5; 1,4,17,72; \ensuremath{\mathbb{Q}}\ensuremath{\sqrt{\,}}5). Verificado (goldchain/metalpos.py): N(\ensuremath{\sigma}\ensuremath{_{n}})=\ensuremath{-}1 e a sequência para os quatro. Retas independentes (\ensuremath{\mathbb{Q}}\ensuremath{\sqrt{\,}}5\ensuremath{\ne}\ensuremath{\mathbb{Q}}\ensuremath{\sqrt{\,}}2\ensuremath{\ne}\ensuremath{\mathbb{Q}}\ensuremath{\sqrt{\,}}13 — esterilidade), cada metal com informação nova; ouro e cobre partilham \ensuremath{\mathbb{Q}}\ensuremath{\sqrt{\,}}5 (distinção de ordem, não de corpo: \ensuremath{\sigma}\ensuremath{_{4}}\ensuremath{\in}\ensuremath{\mathbb{Z}}[\ensuremath{\sqrt{\,}}5]). [D] roda.
\prov{relações: especializa algebra\_posicional\_prata; usa reta\_metalica\_geral; usa independencia\_retas\_metalicas; usa teorema\_esterilidade\_dimensao; relaciona broca\_so\_caixa\_de\_ouro; relaciona caixa\_de\_prata\_corpo; relaciona verificacao\_ouro\_total; parte\_de tiffany; relaciona algebra\_posicional\_plastica; parte\_de algebra\_posicional\_unificada; especializa transformada\_metalica; usa beta\_numeracao\_metalica}
\prov{proveniência: wiki \textperiodcentered{} goldchain/metal\_pos.py (passo\_metal/sequencia); doc/algebra\_prata/paper (§8); reta\_metalica\_geral \textperiodcentered{} fato \textperiodcentered{} confiança 1.0 \textperiodcentered{} domínio algebra \textperiodcentered{} história 26 versões no jornal}

%CRISTAL {"arestas":[{"alvo":"algebra_posicional_metalica","confianca":1.0,"epistemico":"fato","origem":"wiki","rel":"especializa"},{"alvo":"reta_plastica_cubica","confianca":1.0,"epistemico":"fato","origem":"wiki","rel":"usa"},{"alvo":"numeros_pisot","confianca":1.0,"epistemico":"fato","origem":"wiki","rel":"usa"},{"alvo":"reta_metalica_geral","confianca":1.0,"epistemico":"fato","origem":"wiki","rel":"relaciona"},{"alvo":"teorema_esterilidade_dimensao","confianca":1.0,"epistemico":"fato","origem":"wiki","rel":"usa"},{"alvo":"caixas_fractais_generalizacao","confianca":1.0,"epistemico":"fato","origem":"wiki","rel":"relaciona"},{"alvo":"tiffany","confianca":1.0,"epistemico":"fato","origem":"wiki","rel":"parte_de"},{"alvo":"algebra_posicional_unificada","confianca":1.0,"epistemico":"fato","origem":"wiki","rel":"parte_de"},{"alvo":"transformada_metalica","confianca":1.0,"epistemico":"fato","origem":"wiki","rel":"especializa"}],"confianca":1.0,"contraexemplos":[],"descricao":"A família metálica é QUADRÁTICA (grau 2, matriz Aₙ 2×2); a reta PLÁSTICA é CÚBICA (grau 3). O número plástico ρ=1,3247 é a raiz real de x³=x+1 (o MENOR número de Pisot), unidade do anel ℤ[ρ] (ρ⁻¹=ρ²−1, pois ρ³−ρ=1), e a sequência de PADOVAN/Perrin é a coordenada de ρk. A MESMA álgebra posicional, agora com a matriz COMPANHEIRA de grau 3: estado triplo (a,b,c), passo (a,b,c)→(b,c,a+b) mod 2¹²⁸, O(1). Outras cúbicas da reta: ψ supergolden (x³=x²+1, Narayana) e tribonacci (x³=x²+x+1). UNIFICAÇÃO: para qualquer unidade algébrica (Pisot), a moeda é a POSIÇÃO k, emitida por um passo da matriz companheira de grau d (passocompanheira), 128 bits — metais d=2, plástica d=3, d arbitrário. Cada grau/gerador = reta independente (esterilidade vale igual); h=ln(gerador) é o passo geodésico. [D] roda (goldchain/metalpos.py).","epistemico":"fato","exemplos":[],"id":"algebra_posicional_plastica","memoria":"semantica","meta":{"autor":"Lopes/broca-so","dominio":"algebra","fonte":"goldchain/metal_pos.py (passo_companheira/passo_plastico/PLASTICOS); doc/algebra_prata/paper (§9)"},"origem":"wiki","palavras_chave":[],"sinonimos":[],"tipo":"conceito","titulo":"Álgebra Posicional Plástica (cúbica): a metálica de grau 2 sobe a grau 3 (matriz companheira)"}
\section{Álgebra Posicional Plástica (cúbica): a metálica de grau 2 sobe a grau 3 (matriz companheira)}
A família metálica é QUADRÁTICA (grau 2, matriz A\ensuremath{_{n}} 2$\times$2); a reta PLÁSTICA é CÚBICA (grau 3). O número plástico \ensuremath{\rho}=1,3247 é a raiz real de x\ensuremath{^{3}}=x+1 (o MENOR número de Pisot), unidade do anel \ensuremath{\mathbb{Z}}[\ensuremath{\rho}] (\ensuremath{\rho}\ensuremath{^{-1}}=\ensuremath{\rho}\ensuremath{^{2}}\ensuremath{-}1, pois \ensuremath{\rho}\ensuremath{^{3}}\ensuremath{-}\ensuremath{\rho}=1), e a sequência de PADOVAN/Perrin é a coordenada de \ensuremath{\rho}k. A MESMA álgebra posicional, agora com a matriz COMPANHEIRA de grau 3: estado triplo (a,b,c), passo (a,b,c)\ensuremath{\to}(b,c,a+b) mod 2\ensuremath{^{128}}, O(1). Outras cúbicas da reta: \ensuremath{\psi} supergolden (x\ensuremath{^{3}}=x\ensuremath{^{2}}+1, Narayana) e tribonacci (x\ensuremath{^{3}}=x\ensuremath{^{2}}+x+1). UNIFICAÇÃO: para qualquer unidade algébrica (Pisot), a moeda é a POSIÇÃO k, emitida por um passo da matriz companheira de grau d (passocompanheira), 128 bits — metais d=2, plástica d=3, d arbitrário. Cada grau/gerador = reta independente (esterilidade vale igual); h=ln(gerador) é o passo geodésico. [D] roda (goldchain/metalpos.py).
\prov{relações: especializa algebra\_posicional\_metalica; usa reta\_plastica\_cubica; usa numeros\_pisot; relaciona reta\_metalica\_geral; usa teorema\_esterilidade\_dimensao; relaciona caixas\_fractais\_generalizacao; parte\_de tiffany; parte\_de algebra\_posicional\_unificada; especializa transformada\_metalica}
\prov{proveniência: wiki \textperiodcentered{} goldchain/metal\_pos.py (passo\_companheira/passo\_plastico/PLASTICOS); doc/algebra\_prata/paper (§9) \textperiodcentered{} fato \textperiodcentered{} confiança 1.0 \textperiodcentered{} domínio algebra \textperiodcentered{} história 19 versões no jornal}

%CRISTAL {"arestas":[{"alvo":"pell_algebra_posicao","confianca":1.0,"epistemico":"fato","origem":"wiki","rel":"usa"},{"alvo":"dimensao_de_prata_pell","confianca":1.0,"epistemico":"fato","origem":"wiki","rel":"usa"},{"alvo":"reta_metalica_geral","confianca":1.0,"epistemico":"fato","origem":"wiki","rel":"usa"},{"alvo":"teorema_esterilidade_dimensao","confianca":1.0,"epistemico":"fato","origem":"wiki","rel":"usa"},{"alvo":"corpo_f2k","confianca":1.0,"epistemico":"fato","origem":"wiki","rel":"usa"},{"alvo":"symbolic_beta_pipe_so","confianca":1.0,"epistemico":"fato","origem":"wiki","rel":"explica"},{"alvo":"independencia_retas_metalicas","confianca":1.0,"epistemico":"fato","origem":"wiki","rel":"relaciona"},{"alvo":"grafo_tiffany_e_algebra_peirce","confianca":1.0,"epistemico":"fato","origem":"wiki","rel":"usa"},{"alvo":"algebra_posicional_metalica","confianca":1.0,"epistemico":"fato","origem":"wiki","rel":"especializa"},{"alvo":"beta_numeracao_metalica","confianca":1.0,"epistemico":"fato","origem":"wiki","rel":"especializa"},{"alvo":"tiffany","confianca":1.0,"epistemico":"fato","origem":"wiki","rel":"parte_de"}],"confianca":1.0,"contraexemplos":[],"descricao":"A ferramenta principal de reserva, cristalizada em paper (doc/algebra_prata/paper_algebra_posicional_prata.md). A ESTRUTURA é um ANEL: ℤ[√2] (o anel metálico de prata, domínio euclidiano, anel de inteiros de ℚ(√2)) — verificado (norma N(a+b√2)=a²−2b², fecho em +/× e distributividade; provado no metal em inteiros, doc/fractais/esterilidade_metal.c). A unidade FUNDAMENTAL σ=1+√2 (N(σ)=−1) gera as moedas: Pell = coordenada de √2 de σⁿ (1,2,5,12,29,70,169,408). A CONTRIBUIÇÃO (engenharia): representar a moeda pela POSIÇÃO n, não pelo valor — emitir = UM passo (×σ, ou ×A na Q-matriz A=[[2,1],[1,0]]) em 128 bits no corpo 𝔽₂¹²⁸, O(1), binário/simbólico (Peirce: o símbolo é a posição), legível por construção (le∘minera=id no corpo). Grupo das unidades ⟨σ⟩ ≅ (ℤ,+). Ganho: 240× (20 ms→87 µs), ledger 18 MB→KB. Nome: evita 'álgebra de Peirce' (álgebras de relação, tomada); 'Pell/prata' não nomeiam estruturas estabelecidas (só a ferramenta). PROMOVIDO a anel/ferramenta. [D] roda.","epistemico":"fato","exemplos":[],"id":"algebra_posicional_prata","memoria":"semantica","meta":{"autor":"Lopes/broca-so","dominio":"algebra","fonte":"doc/algebra_prata/paper_algebra_posicional_prata.md; goldchain/{pell,minera,corpo}.py"},"origem":"wiki","palavras_chave":[],"sinonimos":[],"tipo":"conceito","titulo":"A Álgebra Posicional de Prata (a ferramenta principal) — anel ℤ[√2], posições σⁿ em 𝔽₂¹²⁸"}
\section{A Álgebra Posicional de Prata (a ferramenta principal) — anel \ensuremath{\mathbb{Z}}[\ensuremath{\sqrt{\,}}2], posições \ensuremath{\sigma}\ensuremath{^{n}} em \ensuremath{\mathbb{F}}\ensuremath{_{2}}\ensuremath{^{128}}}
A ferramenta principal de reserva, cristalizada em paper (doc/algebra\_prata/paper\_algebra\_posicional\_prata.md). A ESTRUTURA é um ANEL: \ensuremath{\mathbb{Z}}[\ensuremath{\sqrt{\,}}2] (o anel metálico de prata, domínio euclidiano, anel de inteiros de \ensuremath{\mathbb{Q}}(\ensuremath{\sqrt{\,}}2)) — verificado (norma N(a+b\ensuremath{\sqrt{\,}}2)=a\ensuremath{^{2}}\ensuremath{-}2b\ensuremath{^{2}}, fecho em +/$\times$ e distributividade; provado no metal em inteiros, doc/fractais/esterilidade\_metal.c). A unidade FUNDAMENTAL \ensuremath{\sigma}=1+\ensuremath{\sqrt{\,}}2 (N(\ensuremath{\sigma})=\ensuremath{-}1) gera as moedas: Pell = coordenada de \ensuremath{\sqrt{\,}}2 de \ensuremath{\sigma}\ensuremath{^{n}} (1,2,5,12,29,70,169,408). A CONTRIBUIÇÃO (engenharia): representar a moeda pela POSIÇÃO n, não pelo valor — emitir = UM passo ($\times$\ensuremath{\sigma}, ou $\times$A na Q-matriz A=[[2,1],[1,0]]) em 128 bits no corpo \ensuremath{\mathbb{F}}\ensuremath{_{2}}\ensuremath{^{128}}, O(1), binário/simbólico (Peirce: o símbolo é a posição), legível por construção (le\ensuremath{\circ}minera=id no corpo). Grupo das unidades \ensuremath{\langle}\ensuremath{\sigma}\ensuremath{\rangle} \ensuremath{\cong} (\ensuremath{\mathbb{Z}},+). Ganho: 240$\times$ (20 ms\ensuremath{\to}87 \ensuremath{\mu}s), ledger 18 MB\ensuremath{\to}KB. Nome: evita 'álgebra de Peirce' (álgebras de relação, tomada); 'Pell/prata' não nomeiam estruturas estabelecidas (só a ferramenta). PROMOVIDO a anel/ferramenta. [D] roda.
\prov{relações: usa pell\_algebra\_posicao; usa dimensao\_de\_prata\_pell; usa reta\_metalica\_geral; usa teorema\_esterilidade\_dimensao; usa corpo\_f2k; explica symbolic\_beta\_pipe\_so; relaciona independencia\_retas\_metalicas; usa grafo\_tiffany\_e\_algebra\_peirce; especializa algebra\_posicional\_metalica; especializa beta\_numeracao\_metalica; parte\_de tiffany}
\prov{proveniência: wiki \textperiodcentered{} doc/algebra\_prata/paper\_algebra\_posicional\_prata.md; goldchain/\{pell,minera,corpo\}.py \textperiodcentered{} fato \textperiodcentered{} confiança 1.0 \textperiodcentered{} domínio algebra \textperiodcentered{} história 42 versões no jornal}

%CRISTAL {"arestas":[{"alvo":"algebra_posicional_metalica","confianca":1.0,"epistemico":"fato","origem":"wiki","rel":"usa"},{"alvo":"algebra_posicional_plastica","confianca":1.0,"epistemico":"fato","origem":"wiki","rel":"usa"},{"alvo":"algebra_posicional_prata","confianca":1.0,"epistemico":"fato","origem":"wiki","rel":"usa"},{"alvo":"dupolinomio_meta_inducao","confianca":1.0,"epistemico":"fato","origem":"wiki","rel":"usa"},{"alvo":"reta_metalica_geral","confianca":1.0,"epistemico":"fato","origem":"wiki","rel":"usa"},{"alvo":"tiffany","confianca":1.0,"epistemico":"fato","origem":"wiki","rel":"parte_de"},{"alvo":"transformada_metalica","confianca":1.0,"epistemico":"fato","origem":"wiki","rel":"parte_de"}],"confianca":1.0,"contraexemplos":[],"descricao":"As duas retas num paper só (doc/algebra_prata/paper_algebra_posicional_prata.md): a Álgebra Posicional é UMA família. Cada reta = um POLINÔMIO CARACTERÍSTICO (x²−nx−1 metais; x³−x−1 plástico; …) cuja MATRIZ COMPANHEIRA é o passo de posição e cujas raízes (unidades fundamentais) geram a reserva; as coordenadas são a sequência característica. Metálica: grau 2 (ℤ[σ_n] — ouro/prata/bronze/cobre). Plástica: grau 3 (ℤ[ρ] — plástico/supergolden/tribonacci). Geral: grau d via passo_companheira (goldchain/metal_pos.py), 128 bits, O(d), cifra em 𝔽₂¹²⁸. Cada reta é independente (esterilidade); h=ln(gerador) é o passo geodésico. É a ferramenta única para todas as dimensões. [D] roda.","epistemico":"fato","exemplos":[],"id":"algebra_posicional_unificada","memoria":"semantica","meta":{"autor":"Lopes/broca-so","dominio":"algebra","fonte":"doc/algebra_prata/paper_algebra_posicional_prata.md; goldchain/metal_pos.py"},"origem":"wiki","palavras_chave":[],"sinonimos":[],"tipo":"conceito","titulo":"A Álgebra Posicional unificada: metálica (grau 2) + plástica (grau 3) + grau d (um paper)"}
\section{A Álgebra Posicional unificada: metálica (grau 2) + plástica (grau 3) + grau d (um paper)}
As duas retas num paper só (doc/algebra\_prata/paper\_algebra\_posicional\_prata.md): a Álgebra Posicional é UMA família. Cada reta = um POLINÔMIO CARACTERÍSTICO (x\ensuremath{^{2}}\ensuremath{-}nx\ensuremath{-}1 metais; x\ensuremath{^{3}}\ensuremath{-}x\ensuremath{-}1 plástico; …) cuja MATRIZ COMPANHEIRA é o passo de posição e cujas raízes (unidades fundamentais) geram a reserva; as coordenadas são a sequência característica. Metálica: grau 2 (\ensuremath{\mathbb{Z}}[\ensuremath{\sigma}\_n] — ouro/prata/bronze/cobre). Plástica: grau 3 (\ensuremath{\mathbb{Z}}[\ensuremath{\rho}] — plástico/supergolden/tribonacci). Geral: grau d via passo\_companheira (goldchain/metal\_pos.py), 128 bits, O(d), cifra em \ensuremath{\mathbb{F}}\ensuremath{_{2}}\ensuremath{^{128}}. Cada reta é independente (esterilidade); h=ln(gerador) é o passo geodésico. É a ferramenta única para todas as dimensões. [D] roda.
\prov{relações: usa algebra\_posicional\_metalica; usa algebra\_posicional\_plastica; usa algebra\_posicional\_prata; usa dupolinomio\_meta\_inducao; usa reta\_metalica\_geral; parte\_de tiffany; parte\_de transformada\_metalica}
\prov{proveniência: wiki \textperiodcentered{} doc/algebra\_prata/paper\_algebra\_posicional\_prata.md; goldchain/metal\_pos.py \textperiodcentered{} fato \textperiodcentered{} confiança 1.0 \textperiodcentered{} domínio algebra \textperiodcentered{} história 8 versões no jornal}

%CRISTAL {"arestas":[{"alvo":"transformada_estelar","confianca":1.0,"epistemico":"fato","origem":"wiki","rel":"relaciona"},{"alvo":"transformada_metalica","confianca":1.0,"epistemico":"fato","origem":"wiki","rel":"especializa"},{"alvo":"numeros_de_pisot","confianca":1.0,"epistemico":"fato","origem":"wiki","rel":"relaciona"},{"alvo":"flip_pg_pa","confianca":1.0,"epistemico":"fato","origem":"wiki","rel":"usa"},{"alvo":"algebra_estelar","confianca":1.0,"epistemico":"fato","origem":"wiki","rel":"eh_um"},{"alvo":"invariante_estelar","confianca":1.0,"epistemico":"fato","origem":"wiki","rel":"usa"},{"alvo":"cuspide_pisot","confianca":1.0,"epistemico":"fato","origem":"wiki","rel":"relaciona"}],"confianca":1.0,"contraexemplos":[],"descricao":"A reta mineral, como os corpos dual/glacial/tropical, tem álgebra/geometria/mecânica. A ÁLGEBRA é a Estelar 𝓔ₛ do Gentil: hipercomplexa, 3D, que CONSERVA A NORMA (o invariante a+c²=1, a mesma norma metálica ∏raízes=±1). É o LADO RELATIVÍSTICO (s>0, curvo, multiplicativo); em s=0 é plana/associativa (D=3 limpo, associador=0). A norma conservada em 3D trava o subdominante à dominante: ρ=1/√β, λ=−½·log β — força o cristal. Verificado em bits.","epistemico":"inferencia","exemplos":[],"id":"algebra_reta_mineral","memoria":"semantica","meta":{"dominio":"matematica","fonte":"álgebra da reta mineral"},"origem":"wiki","palavras_chave":[],"sinonimos":[],"tipo":"conceito","titulo":"A Álgebra da Reta Mineral (Estelar 𝓔ₛ do Gentil)"}
\section{A Álgebra da Reta Mineral (Estelar \ensuremath{\mathcal{E}}\ensuremath{_{s}} do Gentil)}
A reta mineral, como os corpos dual/glacial/tropical, tem álgebra/geometria/mecânica. A ÁLGEBRA é a Estelar \ensuremath{\mathcal{E}}\ensuremath{_{s}} do Gentil: hipercomplexa, 3D, que CONSERVA A NORMA (o invariante a+c\ensuremath{^{2}}=1, a mesma norma metálica \ensuremath{\prod}raízes=$\pm$1). É o LADO RELATIVÍSTICO (s>0, curvo, multiplicativo); em s=0 é plana/associativa (D=3 limpo, associador=0). A norma conservada em 3D trava o subdominante à dominante: \ensuremath{\rho}=1/\ensuremath{\sqrt{\,}}\ensuremath{\beta}, \ensuremath{\lambda}=\ensuremath{-}\ensuremath{\tfrac12}\textperiodcentered log \ensuremath{\beta} — força o cristal. Verificado em bits.
\prov{relações: relaciona transformada\_estelar; especializa transformada\_metalica; relaciona numeros\_de\_pisot; usa flip\_pg\_pa; eh\_um algebra\_estelar; usa invariante\_estelar; relaciona cuspide\_pisot}
\prov{proveniência: wiki \textperiodcentered{} álgebra da reta mineral \textperiodcentered{} inferencia \textperiodcentered{} confiança 1.0 \textperiodcentered{} domínio matematica \textperiodcentered{} história 13 versões no jornal}

%CRISTAL {"arestas":[{"alvo":"algebra_tropical","confianca":1.0,"epistemico":"fato","origem":"wiki","rel":"usa"},{"alvo":"quantizacao_tropical","confianca":1.0,"epistemico":"fato","origem":"wiki","rel":"usa"},{"alvo":"corpo_dual","confianca":1.0,"epistemico":"fato","origem":"wiki","rel":"usa"},{"alvo":"cofre_parabola_hero","confianca":1.0,"epistemico":"fato","origem":"wiki","rel":"usa"},{"alvo":"camadas_tiffany_multifractal","confianca":1.0,"epistemico":"fato","origem":"wiki","rel":"explica"}],"confianca":1.0,"contraexemplos":[],"descricao":"Combinar tecidos NÃO é soma linear nem ganho exponencial: é a álgebra TROPICAL (max-plus / logsumexp), a mesma do corpo_dual ('série tropical↔glacial: balança a+c²=1'). A taxa combinada de N tecidos satura como um logsumexp — o ganho POR tecido DECAI (exponencialmente, contenção), mas a taxa COMBINADA sobe significativamente quando operam juntas. A temperatura T do tropical (⊕_T → min/max no limite T→0, a quantização_tropical) é fixada pela PARÁBOLA a+c²=1: c² = o lado cristal/associativo (o ganho), a = o lado calor/dissipativo (a contenção). Mesma parábola do cristal/floco, agora sobre a álgebra dos tecidos.","epistemico":"inferencia","exemplos":[],"id":"algebra_tecidos_tropical","memoria":"semantica","meta":{"autor":"Lopes/broca-so","dominio":"continuo","fonte":"algebra_tropical + quantizacao_tropical + corpo_dual (a+c²=1); goldchain/estudo_multifractal.py"},"origem":"wiki","palavras_chave":[],"sinonimos":[],"tipo":"conceito","titulo":"A álgebra dos tecidos é a mesma — tropical (max-plus), governada pela parábola"}
\section{A álgebra dos tecidos é a mesma — tropical (max-plus), governada pela parábola}
Combinar tecidos NÃO é soma linear nem ganho exponencial: é a álgebra TROPICAL (max-plus / logsumexp), a mesma do corpo\_dual ('série tropical\ensuremath{\leftrightarrow}glacial: balança a+c\ensuremath{^{2}}=1'). A taxa combinada de N tecidos satura como um logsumexp — o ganho POR tecido DECAI (exponencialmente, contenção), mas a taxa COMBINADA sobe significativamente quando operam juntas. A temperatura T do tropical (\ensuremath{\oplus}\_T \ensuremath{\to} min/max no limite T\ensuremath{\to}0, a quantização\_tropical) é fixada pela PARÁBOLA a+c\ensuremath{^{2}}=1: c\ensuremath{^{2}} = o lado cristal/associativo (o ganho), a = o lado calor/dissipativo (a contenção). Mesma parábola do cristal/floco, agora sobre a álgebra dos tecidos.
\prov{relações: usa algebra\_tropical; usa quantizacao\_tropical; usa corpo\_dual; usa cofre\_parabola\_hero; explica camadas\_tiffany\_multifractal}
\prov{proveniência: wiki \textperiodcentered{} algebra\_tropical + quantizacao\_tropical + corpo\_dual (a+c\ensuremath{^{2}}=1); goldchain/estudo\_multifractal.py \textperiodcentered{} inferencia \textperiodcentered{} confiança 1.0 \textperiodcentered{} domínio continuo \textperiodcentered{} história 24 versões no jornal}

%CRISTAL {"arestas":[{"alvo":"matematica","confianca":1.0,"epistemico":"fato","origem":"paper","rel":"parte_de"},{"alvo":"colapso","confianca":1.0,"epistemico":"fato","origem":"paper","rel":"referencia"}],"confianca":1.0,"contraexemplos":[],"descricao":"Semianel (max, +): ⊕ é o máximo, ⊗ é a soma. Lopes identifica a indução composta como meta-indução tropical (λ = max dos autovalores de Perron).","epistemico":"fato","exemplos":[],"id":"algebra_tropical","memoria":"semantica","meta":{"autor":"Gentil Lopes","descoberta":true,"dominio":"matematica","fonte":"paper_63_algebra_tropical.pdf"},"origem":"paper","palavras_chave":[],"sinonimos":[],"tipo":"conceito","titulo":"Álgebra Tropical (max-plus)"}
\section{Álgebra Tropical (max-plus)}
Semianel (max, +): \ensuremath{\oplus} é o máximo, \ensuremath{\otimes} é a soma. Lopes identifica a indução composta como meta-indução tropical (\ensuremath{\lambda} = max dos autovalores de Perron).
\prov{relações: parte\_de matematica; referencia colapso}
\prov{proveniência: paper \textperiodcentered{} paper\_63\_algebra\_tropical.pdf \textperiodcentered{} fato \textperiodcentered{} confiança 1.0 \textperiodcentered{} domínio matematica \textperiodcentered{} história 4 versões no jornal}

%CRISTAL {"arestas":[{"alvo":"matrizes","confianca":1.0,"epistemico":"fato","origem":"paper","rel":"usa"},{"alvo":"complexos","confianca":1.0,"epistemico":"fato","origem":"paper","rel":"generaliza"},{"alvo":"gato","confianca":1.0,"epistemico":"fato","origem":"paper","rel":"referencia"},{"alvo":"mecanica_quantica","confianca":1.0,"epistemico":"fato","origem":"paper","rel":"referencia"}],"confianca":1.0,"contraexemplos":[],"descricao":"Família de álgebras em ℝⁿ com multiplicação parametrizada por s: casos ℂ (n=2), ℍ quatérnios (n=4), 𝕆 octônios (n=8). Estrutura unificadora das álgebras de divisão.","epistemico":"fato","exemplos":[],"id":"algebras_hipercomplexas","memoria":"semantica","meta":{"autor":"Gentil Lopes","descoberta":true,"dominio":"matematica","fonte":"paper_1_algebra_geometria.pdf"},"origem":"paper","palavras_chave":[],"sinonimos":[],"tipo":"conceito","titulo":"Álgebras Hipercomplexas (família ×ₛ)"}
\section{Álgebras Hipercomplexas (família $\times$\ensuremath{_{s}})}
Família de álgebras em \ensuremath{\mathbb{R}}\ensuremath{^{n}} com multiplicação parametrizada por s: casos \ensuremath{\mathbb{C}} (n=2), \ensuremath{\mathbb{H}} quatérnios (n=4), \ensuremath{\mathbb{O}} octônios (n=8). Estrutura unificadora das álgebras de divisão.
\prov{relações: usa matrizes; generaliza complexos; referencia gato; referencia mecanica\_quantica}
\prov{proveniência: paper \textperiodcentered{} paper\_1\_algebra\_geometria.pdf \textperiodcentered{} fato \textperiodcentered{} confiança 1.0 \textperiodcentered{} domínio matematica \textperiodcentered{} história 5 versões no jornal}

%CRISTAL {"arestas":[{"alvo":"sete_constantes_si","confianca":1.0,"epistemico":"fato","origem":"wiki","rel":"parte_de"},{"alvo":"estatuto_derivado_posto_convencao","confianca":1.0,"epistemico":"fato","origem":"wiki","rel":"relaciona"},{"alvo":"bohr_niveis_energia","confianca":1.0,"epistemico":"fato","origem":"wiki","rel":"usa"},{"alvo":"mecanica_quantica","confianca":1.0,"epistemico":"fato","origem":"wiki","rel":"relaciona"}],"confianca":1.0,"contraexemplos":[],"descricao":"A constante de estrutura fina, lida em três réguas independentes que convergem. (1) ATÔMICA: a radiação que define o metro carrega α — Ry=½α²m_ec², logo α=√(2Ry/m_ec²)=√(2·13,606/511000)=1/137,0 (determinação metrológica, não derivação). (2) VON KLITZING: R_K=h/e²=Z₀/2α, a resistência quântica é 1/α; exata por TOPOLOGIA — o mesmo número de Chern do fibrado de Hopf 𝒪(−1) (Parte III). (3) RUNNING: α corre (1/137 no IR, 1/128 em M_Z); o coeficiente é ΣN_c Q_f²=8, de Q=N/3 somado sobre as 48 cargas. A álgebra dá a ESTRUTURA/running; o valor 1/137 é condição IR, posta. Numerologia (combinar 42,√2,π p/ achar 137) é recusada por princípio.","epistemico":"fato","exemplos":[],"id":"alpha_tres_reguas","memoria":"semantica","meta":{"autor":"Lopes/broca-so","dominio":"unidades","fonte":"hiper/circular/paper_5_unidades.tex §alpha (franjas_metro.py, von_klitzing.py, derivacao_alpha.py)"},"origem":"wiki","palavras_chave":[],"sinonimos":[],"tipo":"conceito","titulo":"α em três réguas: atômica, von Klitzing, e o running"}
\section{\ensuremath{\alpha} em três réguas: atômica, von Klitzing, e o running}
A constante de estrutura fina, lida em três réguas independentes que convergem. (1) ATÔMICA: a radiação que define o metro carrega \ensuremath{\alpha} — Ry=\ensuremath{\tfrac12}\ensuremath{\alpha}\ensuremath{^{2}}m\_ec\ensuremath{^{2}}, logo \ensuremath{\alpha}=\ensuremath{\sqrt{\,}}(2Ry/m\_ec\ensuremath{^{2}})=\ensuremath{\sqrt{\,}}(2\textperiodcentered 13,606/511000)=1/137,0 (determinação metrológica, não derivação). (2) VON KLITZING: R\_K=h/e\ensuremath{^{2}}=Z\ensuremath{_{0}}/2\ensuremath{\alpha}, a resistência quântica é 1/\ensuremath{\alpha}; exata por TOPOLOGIA — o mesmo número de Chern do fibrado de Hopf \ensuremath{\mathcal{O}}(\ensuremath{-}1) (Parte III). (3) RUNNING: \ensuremath{\alpha} corre (1/137 no IR, 1/128 em M\_Z); o coeficiente é \ensuremath{\Sigma}N\_c Q\_f\ensuremath{^{2}}=8, de Q=N/3 somado sobre as 48 cargas. A álgebra dá a ESTRUTURA/running; o valor 1/137 é condição IR, posta. Numerologia (combinar 42,\ensuremath{\sqrt{\,}}2,\ensuremath{\pi} p/ achar 137) é recusada por princípio.
\prov{relações: parte\_de sete\_constantes\_si; relaciona estatuto\_derivado\_posto\_convencao; usa bohr\_niveis\_energia; relaciona mecanica\_quantica}
\prov{proveniência: wiki \textperiodcentered{} hiper/circular/paper\_5\_unidades.tex §alpha (franjas\_metro.py, von\_klitzing.py, derivacao\_alpha.py) \textperiodcentered{} fato \textperiodcentered{} confiança 1.0 \textperiodcentered{} domínio unidades \textperiodcentered{} história 6 versões no jornal}

%CRISTAL {"arestas":[{"alvo":"expressividade","confianca":0.5,"epistemico":"fato","origem":"wiki","rel":"relaciona"},{"alvo":"etica","confianca":0.5,"epistemico":"fato","origem":"wiki","rel":"relaciona"},{"alvo":"paper_0_sintese","confianca":0.5,"epistemico":"fato","origem":"wiki","rel":"relaciona"},{"alvo":"word","confianca":0.5,"epistemico":"fato","origem":"wiki","rel":"relaciona"},{"alvo":"vontade_de_poder","confianca":0.5,"epistemico":"fato","origem":"wiki","rel":"relaciona"},{"alvo":"virtualizacao","confianca":0.5,"epistemico":"fato","origem":"wiki","rel":"relaciona"},{"alvo":"while","confianca":0.5,"epistemico":"fato","origem":"wiki","rel":"relaciona"},{"alvo":"transcricao","confianca":0.5,"epistemico":"fato","origem":"wiki","rel":"relaciona"},{"alvo":"dois_qubits_s7_s4","confianca":0.5,"epistemico":"fato","origem":"wiki","rel":"relaciona"},{"alvo":"semiotica_peirce_triade_tipos_de_interpretante","confianca":0.5,"epistemico":"fato","origem":"wiki","rel":"relaciona"},{"alvo":"geometria_hub_broca_contraexemplos","confianca":0.5,"epistemico":"fato","origem":"wiki","rel":"relaciona"},{"alvo":"progressao_geometrica","confianca":0.5,"epistemico":"fato","origem":"wiki","rel":"relaciona"},{"alvo":"veredito_experimental","confianca":0.5,"epistemico":"fato","origem":"wiki","rel":"relaciona"},{"alvo":"habitus","confianca":0.5,"epistemico":"fato","origem":"wiki","rel":"relaciona"}],"confianca":1.0,"contraexemplos":[],"descricao":"Limites, derivadas, integrais; cálculo sobre ℝ ou ℂ.","epistemico":"fato","exemplos":[],"id":"analise","memoria":"semantica","meta":{"dominio":"matematica","nivel":4,"ordem":15},"origem":"curriculo","palavras_chave":[],"sinonimos":["análise","cálculo"],"tipo":"conceito","titulo":"analise"}
\section{analise}
Limites, derivadas, integrais; cálculo sobre \ensuremath{\mathbb{R}} ou \ensuremath{\mathbb{C}}.
\prov{sinónimos: análise; cálculo}
\prov{relações: relaciona expressividade; relaciona etica; relaciona paper\_0\_sintese; relaciona word; relaciona vontade\_de\_poder; relaciona virtualizacao; relaciona while; relaciona transcricao; relaciona dois\_qubits\_s7\_s4; relaciona semiotica\_peirce\_triade\_tipos\_de\_interpretante; relaciona geometria\_hub\_broca\_contraexemplos; relaciona progressao\_geometrica; relaciona veredito\_experimental; relaciona habitus}
\prov{proveniência: curriculo \textperiodcentered{} fato \textperiodcentered{} confiança 1.0 \textperiodcentered{} domínio matematica \textperiodcentered{} história 16 versões no jornal}

%CRISTAL {"arestas":[{"alvo":"inversao_geometrica_s","confianca":1.0,"epistemico":"fato","origem":"paper","rel":"usa"},{"alvo":"identidade_legendre","confianca":1.0,"epistemico":"fato","origem":"paper","rel":"usa"}],"confianca":1.0,"contraexemplos":[],"descricao":"π expresso via inversão geométrica e o imposto métrico δ: π=℘(2√2·E−℘), ligando a constante lemniscática ℘ à energia E da lemniscata pela identidade de Legendre.","epistemico":"fato","exemplos":[],"id":"anatomia_pi","memoria":"semantica","meta":{"autor":"Gentil Lopes","descoberta":true,"dominio":"matematica","fonte":"paper_32_inversao_anatomia_pi.pdf, Teo.5.1"},"origem":"paper","palavras_chave":[],"sinonimos":[],"tipo":"conceito","titulo":"Anatomia de π"}
\section{Anatomia de \ensuremath{\pi}}
\ensuremath{\pi} expresso via inversão geométrica e o imposto métrico \ensuremath{\delta}: \ensuremath{\pi}=\ensuremath{\wp}(2\ensuremath{\sqrt{\,}}2\textperiodcentered E\ensuremath{-}\ensuremath{\wp}), ligando a constante lemniscática \ensuremath{\wp} à energia E da lemniscata pela identidade de Legendre.
\prov{relações: usa inversao\_geometrica\_s; usa identidade\_legendre}
\prov{proveniência: paper \textperiodcentered{} paper\_32\_inversao\_anatomia\_pi.pdf, Teo.5.1 \textperiodcentered{} fato \textperiodcentered{} confiança 1.0 \textperiodcentered{} domínio matematica \textperiodcentered{} história 3 versões no jornal}

%CRISTAL {"arestas":[{"alvo":"lei_da_parabola_e_dinamica_de_anosov","confianca":1.0,"epistemico":"fato","origem":"wiki","rel":"parte_de"},{"alvo":"selberg_promove_a_parabola","confianca":1.0,"epistemico":"fato","origem":"wiki","rel":"usa"},{"alvo":"chicote_e_o_crivo_de_selberg","confianca":1.0,"epistemico":"fato","origem":"wiki","rel":"usa"},{"alvo":"hipotese_de_riemann","confianca":1.0,"epistemico":"fato","origem":"wiki","rel":"relaciona"},{"alvo":"corpo_estelar","confianca":1.0,"epistemico":"fato","origem":"wiki","rel":"relaciona"},{"alvo":"linguagem_fisica_unificada","confianca":1.0,"epistemico":"fato","origem":"wiki","rel":"parte_de"},{"alvo":"bai_e_a_cosmologia_quantica","confianca":1.0,"epistemico":"fato","origem":"wiki","rel":"relaciona"}],"confianca":1.0,"contraexemplos":[],"descricao":"A face ESPECTRAL da dinâmica: a zeta e a fórmula do TRAÇO de Selberg em Γ\\ℍ. No caso COMPACTO, os zeros do bulk (autovalores λ≥¼) da zeta de Selberg estão em Re(s)=½ — a RH ali é PROVADA (Selberg 1956); os excepcionais λ<¼ dão zeros reais fora da linha (conjectura de Selberg λ_1≥¼ ABERTA, cota 975/4096 Kim-Sarnak). O traço liga espectro↔geodésicas (espectral=geométrico). ESCOPO (reparo do validador): a âncora Anosov/Selberg FUNDA a mecânica (o gato), a entropia/Lyapunov e a face espectral (traço=partição); NÃO deduz o corpo estelar (La Hire=a CÔNICA polo-polar) nem a cosmologia de Sitter (curvatura POSITIVA; Anosov vive na NEGATIVA) — esses unificam por analogia da cônica. INFERENCIA forte, não FATO único; 'fundamenta tudo' seria sobre-afirmação.","epistemico":"inferencia","exemplos":[],"id":"ancora_selberg_e_seu_escopo","memoria":"semantica","meta":{"dominio":"matematica","fonte":"dinâmica de Anosov / Selberg"},"origem":"wiki","palavras_chave":[],"sinonimos":[],"tipo":"conceito","titulo":"A Âncora de Selberg e seu Escopo Honesto"}
\section{A Âncora de Selberg e seu Escopo Honesto}
A face ESPECTRAL da dinâmica: a zeta e a fórmula do TRAÇO de Selberg em \ensuremath{\Gamma}\textbackslash{}\ensuremath{\mathbb{H}}. No caso COMPACTO, os zeros do bulk (autovalores \ensuremath{\lambda}\ensuremath{\ge}\ensuremath{\tfrac14}) da zeta de Selberg estão em Re(s)=\ensuremath{\tfrac12} — a RH ali é PROVADA (Selberg 1956); os excepcionais \ensuremath{\lambda}<\ensuremath{\tfrac14} dão zeros reais fora da linha (conjectura de Selberg \ensuremath{\lambda}\_1\ensuremath{\ge}\ensuremath{\tfrac14} ABERTA, cota 975/4096 Kim-Sarnak). O traço liga espectro\ensuremath{\leftrightarrow}geodésicas (espectral=geométrico). ESCOPO (reparo do validador): a âncora Anosov/Selberg FUNDA a mecânica (o gato), a entropia/Lyapunov e a face espectral (traço=partição); NÃO deduz o corpo estelar (La Hire=a CÔNICA polo-polar) nem a cosmologia de Sitter (curvatura POSITIVA; Anosov vive na NEGATIVA) — esses unificam por analogia da cônica. INFERENCIA forte, não FATO único; 'fundamenta tudo' seria sobre-afirmação.
\prov{relações: parte\_de lei\_da\_parabola\_e\_dinamica\_de\_anosov; usa selberg\_promove\_a\_parabola; usa chicote\_e\_o\_crivo\_de\_selberg; relaciona hipotese\_de\_riemann; relaciona corpo\_estelar; parte\_de linguagem\_fisica\_unificada; relaciona bai\_e\_a\_cosmologia\_quantica}
\prov{proveniência: wiki \textperiodcentered{} dinâmica de Anosov / Selberg \textperiodcentered{} inferencia \textperiodcentered{} confiança 1.0 \textperiodcentered{} domínio matematica \textperiodcentered{} história 32 versões no jornal}

%CRISTAL {"arestas":[{"alvo":"corpos","confianca":1.0,"epistemico":"fato","origem":"curriculo","rel":"especializa"},{"alvo":"etica","confianca":0.5,"epistemico":"fato","origem":"wiki","rel":"relaciona"},{"alvo":"expressividade","confianca":0.5,"epistemico":"fato","origem":"wiki","rel":"relaciona"},{"alvo":"imagem_do_floco_---_validacao_no","confianca":0.5,"epistemico":"fato","origem":"wiki","rel":"relaciona"},{"alvo":"topologia","confianca":0.5,"epistemico":"fato","origem":"wiki","rel":"relaciona"},{"alvo":"word","confianca":0.5,"epistemico":"fato","origem":"wiki","rel":"relaciona"}],"confianca":1.0,"contraexemplos":[],"descricao":"Dois operações (+, ×) com distributividade.","epistemico":"fato","exemplos":[],"id":"aneis","memoria":"semantica","meta":{"dominio":"matematica","nivel":4,"ordem":10},"origem":"curriculo","palavras_chave":[],"sinonimos":["anel"],"tipo":"conceito","titulo":"aneis"}
\section{aneis}
Dois operações (+, $\times$) com distributividade.
\prov{sinónimos: anel}
\prov{relações: especializa corpos; relaciona etica; relaciona expressividade; relaciona imagem\_do\_floco\_---\_validacao\_no; relaciona topologia; relaciona word}
\prov{proveniência: curriculo \textperiodcentered{} fato \textperiodcentered{} confiança 1.0 \textperiodcentered{} domínio matematica \textperiodcentered{} história 9 versões no jornal}

%CRISTAL {"arestas":[{"alvo":"word","confianca":0.5,"epistemico":"fato","origem":"wiki","rel":"relaciona"},{"alvo":"virtualizacao","confianca":0.5,"epistemico":"fato","origem":"wiki","rel":"relaciona"},{"alvo":"vida-morte","confianca":0.5,"epistemico":"fato","origem":"wiki","rel":"relaciona"},{"alvo":"sagrado_profano","confianca":0.5,"epistemico":"fato","origem":"wiki","rel":"relaciona"},{"alvo":"criterio_de_projeto","confianca":0.5,"epistemico":"fato","origem":"wiki","rel":"relaciona"},{"alvo":"octonios_hiperbolicos","confianca":0.5,"epistemico":"fato","origem":"wiki","rel":"relaciona"},{"alvo":"cache","confianca":0.5,"epistemico":"fato","origem":"wiki","rel":"relaciona"}],"confianca":1.0,"contraexemplos":[],"descricao":"Rₖ=ℚ[π²,ζ(3),ζ(5),…,ζ(2k−1)]≤ₖ: combinações ℚ-lineares de peso transcendental ≤k; espaço de instâncias cujo tamanho cresce exp(k).","epistemico":"fato","exemplos":[],"id":"anel_transcendental_rk","memoria":"semantica","meta":{"autor":"Gentil Lopes","descoberta":false,"dominio":"matematica","fonte":"paper_PQ1_beta_octonic_hardness.pdf, Def.2.1"},"origem":"paper","palavras_chave":[],"sinonimos":[],"tipo":"conceito","titulo":"Anel Transcendental R_k"}
\section{Anel Transcendental R\_k}
R\ensuremath{_{k}}=\ensuremath{\mathbb{Q}}[\ensuremath{\pi}\ensuremath{^{2}},\ensuremath{\zeta}(3),\ensuremath{\zeta}(5),…,\ensuremath{\zeta}(2k\ensuremath{-}1)]\ensuremath{\le}\ensuremath{_{k}}: combinações \ensuremath{\mathbb{Q}}-lineares de peso transcendental \ensuremath{\le}k; espaço de instâncias cujo tamanho cresce exp(k).
\prov{relações: relaciona word; relaciona virtualizacao; relaciona vida-morte; relaciona sagrado\_profano; relaciona criterio\_de\_projeto; relaciona octonios\_hiperbolicos; relaciona cache}
\prov{proveniência: paper \textperiodcentered{} paper\_PQ1\_beta\_octonic\_hardness.pdf, Def.2.1 \textperiodcentered{} fato \textperiodcentered{} confiança 1.0 \textperiodcentered{} domínio matematica \textperiodcentered{} história 9 versões no jornal}

%CRISTAL {"arestas":[{"alvo":"organism32","confianca":1.0,"epistemico":"fato","origem":"paper","rel":"especializa"},{"alvo":"incautos_enganados","confianca":0.5,"epistemico":"fato","origem":"wiki","rel":"relaciona"},{"alvo":"word","confianca":0.5,"epistemico":"fato","origem":"wiki","rel":"relaciona"},{"alvo":"virtualizacao","confianca":0.5,"epistemico":"fato","origem":"wiki","rel":"relaciona"},{"alvo":"octonios_hiperbolicos","confianca":0.5,"epistemico":"fato","origem":"wiki","rel":"relaciona"},{"alvo":"vida-morte","confianca":0.5,"epistemico":"fato","origem":"wiki","rel":"relaciona"},{"alvo":"psl_2_7","confianca":0.5,"epistemico":"fato","origem":"wiki","rel":"relaciona"},{"alvo":"criterio_de_projeto","confianca":0.5,"epistemico":"fato","origem":"wiki","rel":"relaciona"}],"confianca":1.0,"contraexemplos":[],"descricao":"Cifra de 16 componentes com chave transcendental (fração contínua de e), Cantor real (z↦z²) e Fourier real (NTT).","epistemico":"fato","exemplos":[],"id":"animal16","memoria":"semantica","meta":{"autor":"Gentil Lopes","descoberta":true,"dominio":"matematica","fonte":"paper_13_animal16.pdf"},"origem":"paper","palavras_chave":[],"sinonimos":[],"tipo":"conceito","titulo":"Animal-16"}
\section{Animal-16}
Cifra de 16 componentes com chave transcendental (fração contínua de e), Cantor real (z\ensuremath{\mapsto}z\ensuremath{^{2}}) e Fourier real (NTT).
\prov{relações: especializa organism32; relaciona incautos\_enganados; relaciona word; relaciona virtualizacao; relaciona octonios\_hiperbolicos; relaciona vida-morte; relaciona psl\_2\_7; relaciona criterio\_de\_projeto}
\prov{proveniência: paper \textperiodcentered{} paper\_13\_animal16.pdf \textperiodcentered{} fato \textperiodcentered{} confiança 1.0 \textperiodcentered{} domínio matematica \textperiodcentered{} história 9 versões no jornal}

%CRISTAL {"arestas":[{"alvo":"reta_metalica_geral","confianca":1.0,"epistemico":"fato","origem":"wiki","rel":"usa"},{"alvo":"bai_orbita_seleciona_colapso","confianca":1.0,"epistemico":"fato","origem":"wiki","rel":"relaciona"},{"alvo":"colapso_multicamada_peso","confianca":1.0,"epistemico":"fato","origem":"wiki","rel":"relaciona"},{"alvo":"tiffany","confianca":1.0,"epistemico":"fato","origem":"wiki","rel":"parte_de"}],"confianca":1.0,"contraexemplos":[],"descricao":"Cada aresta (A—rel→B) é o segmento [A,B] do contínuo. Um fragmento real f tem POSIÇÃO s(f)=ov(f,B)/(ov(f,A)+ov(f,B)) ∈ [0,1] (0=colado em A, 1=em B) e MEDIDA DE PONTE ponte(f,A,B)=min(ov(f,A),ov(f,B))+λ·(ov(f,A)+ov(f,B)) — ressoar com os DOIS extremos preenche o segmento (o min premia a ponte, não a proximidade de um lado). A transição da aresta = os fragmentos ordenados por s → a interpolação graduada do paper. Armazenada num campo novo 'transicao' na aresta + tecido novo 'transicoes' indexado (o colapso passa a pousar no MEIO das arestas, com energia de órbita própria).","epistemico":"inferencia","exemplos":[],"id":"aresta_transicao_graduada","memoria":"semantica","meta":{"autor":"Lopes/broca-so","dominio":"continuo","fonte":"doc/continuo/plano_transicoes.md; conversa/conhecimento/schema.py (Aresta)"},"origem":"wiki","palavras_chave":[],"sinonimos":[],"tipo":"conceito","titulo":"A aresta como segmento graduado (s:0→1) — o modelo da transição"}
\section{A aresta como segmento graduado (s:0\ensuremath{\to}1) — o modelo da transição}
Cada aresta (A—rel\ensuremath{\to}B) é o segmento [A,B] do contínuo. Um fragmento real f tem POSIÇÃO s(f)=ov(f,B)/(ov(f,A)+ov(f,B)) \ensuremath{\in} [0,1] (0=colado em A, 1=em B) e MEDIDA DE PONTE ponte(f,A,B)=min(ov(f,A),ov(f,B))+\ensuremath{\lambda}\textperiodcentered (ov(f,A)+ov(f,B)) — ressoar com os DOIS extremos preenche o segmento (o min premia a ponte, não a proximidade de um lado). A transição da aresta = os fragmentos ordenados por s \ensuremath{\to} a interpolação graduada do paper. Armazenada num campo novo 'transicao' na aresta + tecido novo 'transicoes' indexado (o colapso passa a pousar no MEIO das arestas, com energia de órbita própria).
\prov{relações: usa reta\_metalica\_geral; relaciona bai\_orbita\_seleciona\_colapso; relaciona colapso\_multicamada\_peso; parte\_de tiffany}
\prov{proveniência: wiki \textperiodcentered{} doc/continuo/plano\_transicoes.md; conversa/conhecimento/schema.py (Aresta) \textperiodcentered{} inferencia \textperiodcentered{} confiança 1.0 \textperiodcentered{} domínio continuo \textperiodcentered{} história 82 versões no jornal}

%CRISTAL {"arestas":[{"alvo":"ym_pow_mining","confianca":1.0,"epistemico":"fato","origem":"paper","rel":"explica"}],"confianca":1.0,"contraexemplos":[],"descricao":"Não existe circuito booleano fixo de tamanho O(κc) que avalie o determinante DM (matriz M=O(κ²) de coeficientes transcendentais): a resistência a ASIC é consequência aritmética, não design.","epistemico":"fato","exemplos":[],"id":"asic_resistencia","memoria":"semantica","meta":{"autor":"Gentil Lopes","descoberta":false,"dominio":"matematica","fonte":"paper_PQ_R.pdf, Prop.5.3"},"origem":"paper","palavras_chave":[],"sinonimos":[],"tipo":"conceito","titulo":"ASIC-Resistência Estrutural"}
\section{ASIC-Resistência Estrutural}
Não existe circuito booleano fixo de tamanho O(\ensuremath{\kappa}c) que avalie o determinante DM (matriz M=O(\ensuremath{\kappa}\ensuremath{^{2}}) de coeficientes transcendentais): a resistência a ASIC é consequência aritmética, não design.
\prov{relações: explica ym\_pow\_mining}
\prov{proveniência: paper \textperiodcentered{} paper\_PQ\_R.pdf, Prop.5.3 \textperiodcentered{} fato \textperiodcentered{} confiança 1.0 \textperiodcentered{} domínio matematica \textperiodcentered{} história 3 versões no jornal}

%CRISTAL {"arestas":[{"alvo":"algebras_hipercomplexas","confianca":1.0,"epistemico":"fato","origem":"paper","rel":"parte_de"},{"alvo":"word","confianca":0.5,"epistemico":"fato","origem":"wiki","rel":"relaciona"}],"confianca":1.0,"contraexemplos":[],"descricao":"Tensor [x,y,v]ₛ = (1−s²)A₀ que mede a não-associatividade da álgebra ×ₛ; magnitude canônica a = |1−s²|. Nulo em ℝ,ℂ,ℍ; não-nulo em 𝕆.","epistemico":"fato","exemplos":[],"id":"associador","memoria":"semantica","meta":{"autor":"Gentil Lopes","descoberta":true,"dominio":"matematica","fonte":"paper_1_algebra_geometria.pdf, Prop.3.4"},"origem":"paper","palavras_chave":[],"sinonimos":[],"tipo":"conceito","titulo":"Associador"}
\section{Associador}
Tensor [x,y,v]\ensuremath{_{s}} = (1\ensuremath{-}s\ensuremath{^{2}})A\ensuremath{_{0}} que mede a não-associatividade da álgebra $\times$\ensuremath{_{s}}; magnitude canônica a = |1\ensuremath{-}s\ensuremath{^{2}}|. Nulo em \ensuremath{\mathbb{R}},\ensuremath{\mathbb{C}},\ensuremath{\mathbb{H}}; não-nulo em \ensuremath{\mathbb{O}}.
\prov{relações: parte\_de algebras\_hipercomplexas; relaciona word}
\prov{proveniência: paper \textperiodcentered{} paper\_1\_algebra\_geometria.pdf, Prop.3.4 \textperiodcentered{} fato \textperiodcentered{} confiança 1.0 \textperiodcentered{} domínio matematica \textperiodcentered{} história 3 versões no jornal}

%CRISTAL {"arestas":[{"alvo":"goldchain","confianca":1.0,"epistemico":"fato","origem":"paper","rel":"referencia"}],"confianca":1.0,"contraexemplos":[],"descricao":"A atestação prova a identidade mas não recupera o estado: existem σ'≠σ com a mesma projeção μ. Memória de verdade exige o payload reconstituível, não só o atestado — separação de prova e história.","epistemico":"fato","exemplos":[],"id":"atestacao_projecao","memoria":"semantica","meta":{"autor":"Gentil Lopes","descoberta":false,"dominio":"matematica","fonte":"cristal.pdf, §atestar≠reconstituir"},"origem":"paper","palavras_chave":[],"sinonimos":[],"tipo":"conceito","titulo":"Atestar ≠ Reconstituir"}
\section{Atestar \ensuremath{\ne} Reconstituir}
A atestação prova a identidade mas não recupera o estado: existem \ensuremath{\sigma}'\ensuremath{\ne}\ensuremath{\sigma} com a mesma projeção \ensuremath{\mu}. Memória de verdade exige o payload reconstituível, não só o atestado — separação de prova e história.
\prov{relações: referencia goldchain}
\prov{proveniência: paper \textperiodcentered{} cristal.pdf, §atestar\ensuremath{\ne}reconstituir \textperiodcentered{} fato \textperiodcentered{} confiança 1.0 \textperiodcentered{} domínio matematica \textperiodcentered{} história 2 versões no jornal}

%CRISTAL {"arestas":[{"alvo":"fracao_continua_phi","confianca":1.0,"epistemico":"fato","origem":"paper","rel":"generaliza"}],"confianca":1.0,"contraexemplos":[],"descricao":"Uma sequência é estruturada quando alguma operação a reduz a período 1 (constante repetida); o grau de estrutura é o número de operações — algébricos por agrupamento, transcendentais por diferença.","epistemico":"fato","exemplos":[],"id":"atomo_periodo_1","memoria":"semantica","meta":{"autor":"Gentil Lopes","descoberta":true,"dominio":"matematica","fonte":"paper_22_reais.pdf, Def.3.1"},"origem":"paper","palavras_chave":[],"sinonimos":[],"tipo":"conceito","titulo":"Átomo do Período 1"}
\section{Átomo do Período 1}
Uma sequência é estruturada quando alguma operação a reduz a período 1 (constante repetida); o grau de estrutura é o número de operações — algébricos por agrupamento, transcendentais por diferença.
\prov{relações: generaliza fracao\_continua\_phi}
\prov{proveniência: paper \textperiodcentered{} paper\_22\_reais.pdf, Def.3.1 \textperiodcentered{} fato \textperiodcentered{} confiança 1.0 \textperiodcentered{} domínio matematica \textperiodcentered{} história 2 versões no jornal}

%CRISTAL {"arestas":[{"alvo":"reta_de_ouro","confianca":1.0,"epistemico":"fato","origem":"paper","rel":"usa"},{"alvo":"gato","confianca":1.0,"epistemico":"fato","origem":"paper","rel":"referencia"},{"alvo":"word","confianca":0.5,"epistemico":"fato","origem":"wiki","rel":"relaciona"}],"confianca":1.0,"contraexemplos":[],"descricao":"Conjunto reconstruído pela codificação na base φ, auto-similar nas escalas φ±k: o shift da sequência age como contração/expansão por φ.","epistemico":"fato","exemplos":[],"id":"atrator_fractal_ouro","memoria":"semantica","meta":{"autor":"Gentil Lopes","descoberta":false,"dominio":"matematica","fonte":"corpo_tropical.pdf, Prop.2.3"},"origem":"paper","palavras_chave":[],"sinonimos":[],"tipo":"conceito","titulo":"Atrator Fractal de Ouro"}
\section{Atrator Fractal de Ouro}
Conjunto reconstruído pela codificação na base \ensuremath{\varphi}, auto-similar nas escalas \ensuremath{\varphi}$\pm$k: o shift da sequência age como contração/expansão por \ensuremath{\varphi}.
\prov{relações: usa reta\_de\_ouro; referencia gato; relaciona word}
\prov{proveniência: paper \textperiodcentered{} corpo\_tropical.pdf, Prop.2.3 \textperiodcentered{} fato \textperiodcentered{} confiança 1.0 \textperiodcentered{} domínio matematica \textperiodcentered{} história 5 versões no jornal}

%CRISTAL {"arestas":[{"alvo":"utilidade_esperada","confianca":1.0,"epistemico":"fato","origem":"wiki","rel":"prova"},{"alvo":"jogador_racional","confianca":1.0,"epistemico":"fato","origem":"wiki","rel":"explica"}],"confianca":1.0,"contraexemplos":[],"descricao":"Se as preferências obedecem a quatro axiomas (completude, transitividade, continuidade, independência), então EXISTE uma função de utilidade cuja esperança as representa. A fundação axiomática da escolha racional — e o 'jogador racional' da teoria dos jogos.","epistemico":"fato","exemplos":[],"id":"axiomas_vnm","memoria":"semantica","meta":{"autor":"von Neumann-Morgenstern","dominio":"matematica","fonte":"teoria da decisão"},"origem":"wiki","palavras_chave":[],"sinonimos":[],"tipo":"conceito","titulo":"Teorema da Utilidade (von Neumann-Morgenstern)"}
\section{Teorema da Utilidade (von Neumann-Morgenstern)}
Se as preferências obedecem a quatro axiomas (completude, transitividade, continuidade, independência), então EXISTE uma função de utilidade cuja esperança as representa. A fundação axiomática da escolha racional — e o 'jogador racional' da teoria dos jogos.
\prov{relações: prova utilidade\_esperada; explica jogador\_racional}
\prov{proveniência: wiki \textperiodcentered{} teoria da decisão \textperiodcentered{} fato \textperiodcentered{} confiança 1.0 \textperiodcentered{} domínio matematica \textperiodcentered{} história 3 versões no jornal}

%CRISTAL {"arestas":[{"alvo":"octonios_hiperbolicos","confianca":1.0,"epistemico":"fato","origem":"paper","rel":"usa"},{"alvo":"word","confianca":0.5,"epistemico":"fato","origem":"wiki","rel":"relaciona"}],"confianca":1.0,"contraexemplos":[],"descricao":"A identidade a×(b×c)=b⟨a,c⟩−c⟨a,b⟩ (que regulariza NS em ℝ³) falha em ℝ⁷ por não-associatividade: o resíduo é da ordem do termo principal.","epistemico":"fato","exemplos":[],"id":"bac_cab_falha_r7","memoria":"semantica","meta":{"autor":"Gentil Lopes","descoberta":false,"dominio":"matematica","fonte":"paper_AG.pdf, Teo.6"},"origem":"paper","palavras_chave":[],"sinonimos":[],"tipo":"conceito","titulo":"BAC-CAB Falha em ℝ⁷"}
\section{BAC-CAB Falha em \ensuremath{\mathbb{R}}\ensuremath{^{7}}}
A identidade a$\times$(b$\times$c)=b\ensuremath{\langle}a,c\ensuremath{\rangle}\ensuremath{-}c\ensuremath{\langle}a,b\ensuremath{\rangle} (que regulariza NS em \ensuremath{\mathbb{R}}\ensuremath{^{3}}) falha em \ensuremath{\mathbb{R}}\ensuremath{^{7}} por não-associatividade: o resíduo é da ordem do termo principal.
\prov{relações: usa octonios\_hiperbolicos; relaciona word}
\prov{proveniência: paper \textperiodcentered{} paper\_AG.pdf, Teo.6 \textperiodcentered{} fato \textperiodcentered{} confiança 1.0 \textperiodcentered{} domínio matematica \textperiodcentered{} história 3 versões no jornal}

%CRISTAL {"arestas":[{"alvo":"conjugados_galois_parabola","confianca":1.0,"epistemico":"fato","origem":"wiki","rel":"usa"},{"alvo":"bai_orbita_reta_metalica","confianca":1.0,"epistemico":"fato","origem":"wiki","rel":"especializa"},{"alvo":"parabola_algebra_conversacional","confianca":1.0,"epistemico":"fato","origem":"wiki","rel":"usa"},{"alvo":"dtc_schema_ut","confianca":1.0,"epistemico":"fato","origem":"wiki","rel":"usa"},{"alvo":"corpo_dual","confianca":1.0,"epistemico":"fato","origem":"wiki","rel":"relaciona"},{"alvo":"tiffany","confianca":1.0,"epistemico":"fato","origem":"wiki","rel":"parte_de"}],"confianca":1.0,"contraexemplos":[],"descricao":"conversa/colapso.py (scorehit, pesoorbita): a órbita 𝒞δ deixou de SOMAR um δ pequeno (que era abafado pelo ov associativo ~250) e passou a ESCALAR o lado associativo, no esquema DTC (Direct Torque Control de Takahashi, estendido pelo Lopes — ver [[dtcschemaut]]): a órbita é o SETOR angular; a parábola a+c²=1 é o comparador de torque; o peso da camada = (δcamada/δmax)γ (γ=3) é o ALINHAMENTO ao setor. No foco da órbita o ov entra cheio (o lado associativo c²); atravessar p/ FORA do setor custa o ASSOCIADOR (o lado-a, a entropia de curvatura da DTC E4: 'chaveamentos que atravessam fibras ortogonais são penalizados'). A borda de Salem (a banda de histerese DBC DBAND≈0,85) é o piso do setor — evita o flip-flop. sc = round(ov·peso) + relevância. Efeito real: 'o que é um grupo?' passou a vencer em conhecimento (era saudação) — o torque puxa o colapso p/ a órbita certa; saudação/cadeia/conversa seguem nas suas retas. Não poda (peso>0). Gargalo seguinte (separado): a recuperação intra-camada (o ov Koch não discrimina o conceito). [D] a órbita como torque DTC, não empurrão aditivo.","epistemico":"fato","exemplos":[],"id":"bai_gate_dtc_orbita","memoria":"semantica","meta":{"autor":"Lopes/broca-so","dominio":"algebra","fonte":"conversa/colapso.py:_peso_orbita, score_hit; DTC (Takahashi/Lopes, purgatório paper E1/E4/E5); [[dtc_schema_ut]]"},"origem":"wiki","palavras_chave":[],"sinonimos":[],"tipo":"conceito","titulo":"A órbita do BAI é MULTIPLICATIVA — gate DTC pela parábola (setor + associador)"}
\section{A órbita do BAI é MULTIPLICATIVA — gate DTC pela parábola (setor + associador)}
conversa/colapso.py (scorehit, pesoorbita): a órbita \ensuremath{\mathcal{C}}\ensuremath{\delta} deixou de SOMAR um \ensuremath{\delta} pequeno (que era abafado pelo ov associativo \~{}250) e passou a ESCALAR o lado associativo, no esquema DTC (Direct Torque Control de Takahashi, estendido pelo Lopes — ver [[dtcschemaut]]): a órbita é o SETOR angular; a parábola a+c\ensuremath{^{2}}=1 é o comparador de torque; o peso da camada = (\ensuremath{\delta}camada/\ensuremath{\delta}max)\ensuremath{\gamma} (\ensuremath{\gamma}=3) é o ALINHAMENTO ao setor. No foco da órbita o ov entra cheio (o lado associativo c\ensuremath{^{2}}); atravessar p/ FORA do setor custa o ASSOCIADOR (o lado-a, a entropia de curvatura da DTC E4: 'chaveamentos que atravessam fibras ortogonais são penalizados'). A borda de Salem (a banda de histerese DBC DBAND\ensuremath{\approx}0,85) é o piso do setor — evita o flip-flop. sc = round(ov\textperiodcentered peso) + relevância. Efeito real: 'o que é um grupo?' passou a vencer em conhecimento (era saudação) — o torque puxa o colapso p/ a órbita certa; saudação/cadeia/conversa seguem nas suas retas. Não poda (peso>0). Gargalo seguinte (separado): a recuperação intra-camada (o ov Koch não discrimina o conceito). [D] a órbita como torque DTC, não empurrão aditivo.
\prov{relações: usa conjugados\_galois\_parabola; especializa bai\_orbita\_reta\_metalica; usa parabola\_algebra\_conversacional; usa dtc\_schema\_ut; relaciona corpo\_dual; parte\_de tiffany}
\prov{proveniência: wiki \textperiodcentered{} conversa/colapso.py:\_peso\_orbita, score\_hit; DTC (Takahashi/Lopes, purgatório paper E1/E4/E5); [[dtc\_schema\_ut]] \textperiodcentered{} fato \textperiodcentered{} confiança 1.0 \textperiodcentered{} domínio algebra \textperiodcentered{} história 128 versões no jornal}

%CRISTAL {"arestas":[{"alvo":"transformada_metalica","confianca":1.0,"epistemico":"fato","origem":"wiki","rel":"aplicacao"},{"alvo":"parabola_algebra_conversacional","confianca":1.0,"epistemico":"fato","origem":"wiki","rel":"usa"},{"alvo":"bai_orbita_energia_por_trabalho","confianca":1.0,"epistemico":"fato","origem":"wiki","rel":"especializa"},{"alvo":"bai_orbita_seleciona_colapso","confianca":1.0,"epistemico":"fato","origem":"wiki","rel":"relaciona"},{"alvo":"tiffany","confianca":1.0,"epistemico":"fato","origem":"wiki","rel":"parte_de"},{"alvo":"algebra_de_peirce","confianca":1.0,"epistemico":"fato","origem":"wiki","rel":"usa"}],"confianca":1.0,"contraexemplos":[],"descricao":"conversa/bai.py (orbita_energia): a órbita 𝒞_δ de cada tipo de trabalho passou a ser uma RETA da nova álgebra, com a energia δ derivada do RAIO σ_n do metal da classe (δ=round(σ_n·4)), não mais números ad-hoc. DEFINIÇÃO (fato) percorre a reta CRISTAL/ouro (φ) → acende a camada conhecimento com δ=6; SOCIAL percorre cristal/ouro → dialogos; CADEIA (continuidade CASL) percorre a BORDA/prata (1+√2) → acende transicoes com δ=10 (reforça o fio do contínuo); CONVERSA percorre o CAOS/cobre → órbita larga de repouso. A assinatura (energia, trabalho) fica preservada (o colapso não muda); o que muda é a FÍSICA do peso: agora vem do gerador da reta. Não poda a floresta — acende a órbita da reta. [D] o BAI algébrico (órbita = reta, δ = σ_n).","epistemico":"fato","exemplos":[],"id":"bai_orbita_reta_metalica","memoria":"semantica","meta":{"autor":"Lopes/broca-so","dominio":"algebra","fonte":"conversa/bai.py:orbita_energia; parabola_algebra.reta_do_trabalho; energia_da_reta"},"origem":"wiki","palavras_chave":[],"sinonimos":[],"tipo":"conceito","titulo":"O BAI: a órbita de cada trabalho é uma RETA metálica — a energia vem do raio σ_n"}
\section{O BAI: a órbita de cada trabalho é uma RETA metálica — a energia vem do raio \ensuremath{\sigma}\_n}
conversa/bai.py (orbita\_energia): a órbita \ensuremath{\mathcal{C}}\_\ensuremath{\delta} de cada tipo de trabalho passou a ser uma RETA da nova álgebra, com a energia \ensuremath{\delta} derivada do RAIO \ensuremath{\sigma}\_n do metal da classe (\ensuremath{\delta}=round(\ensuremath{\sigma}\_n\textperiodcentered 4)), não mais números ad-hoc. DEFINIÇÃO (fato) percorre a reta CRISTAL/ouro (\ensuremath{\varphi}) \ensuremath{\to} acende a camada conhecimento com \ensuremath{\delta}=6; SOCIAL percorre cristal/ouro \ensuremath{\to} dialogos; CADEIA (continuidade CASL) percorre a BORDA/prata (1+\ensuremath{\sqrt{\,}}2) \ensuremath{\to} acende transicoes com \ensuremath{\delta}=10 (reforça o fio do contínuo); CONVERSA percorre o CAOS/cobre \ensuremath{\to} órbita larga de repouso. A assinatura (energia, trabalho) fica preservada (o colapso não muda); o que muda é a FÍSICA do peso: agora vem do gerador da reta. Não poda a floresta — acende a órbita da reta. [D] o BAI algébrico (órbita = reta, \ensuremath{\delta} = \ensuremath{\sigma}\_n).
\prov{relações: aplicacao transformada\_metalica; usa parabola\_algebra\_conversacional; especializa bai\_orbita\_energia\_por\_trabalho; relaciona bai\_orbita\_seleciona\_colapso; parte\_de tiffany; usa algebra\_de\_peirce}
\prov{proveniência: wiki \textperiodcentered{} conversa/bai.py:orbita\_energia; parabola\_algebra.reta\_do\_trabalho; energia\_da\_reta \textperiodcentered{} fato \textperiodcentered{} confiança 1.0 \textperiodcentered{} domínio algebra \textperiodcentered{} história 118 versões no jornal}

%CRISTAL {"arestas":[{"alvo":"razao_aurea","confianca":1.0,"epistemico":"fato","origem":"livro","rel":"usa"},{"alvo":"virtualizacao","confianca":0.5,"epistemico":"fato","origem":"wiki","rel":"relaciona"},{"alvo":"word","confianca":0.5,"epistemico":"fato","origem":"wiki","rel":"relaciona"},{"alvo":"semiotica_peirce_triade","confianca":0.5,"epistemico":"fato","origem":"wiki","rel":"relaciona"},{"alvo":"vida-morte","confianca":0.5,"epistemico":"fato","origem":"wiki","rel":"relaciona"},{"alvo":"beta_numeracao_metalica","confianca":1.0,"epistemico":"fato","origem":"wiki","rel":"especializa"}],"confianca":1.0,"contraexemplos":[],"descricao":"Representação de reais ≥ 0 em base φ com dígitos 0,1, forma canônica sem '11'; bijeção com o anel Z[φ].","epistemico":"fato","exemplos":[],"id":"base_phi_bergman","memoria":"semantica","meta":{"autor":"Gentil Lopes","descoberta":false,"dominio":"matematica","fonte":"paper_2_reta_ouro.pdf, Teorema 1.1"},"origem":"paper","palavras_chave":[],"sinonimos":[],"tipo":"conceito","titulo":"Base-φ (Bergman, 1957)"}
\section{Base-\ensuremath{\varphi} (Bergman, 1957)}
Representação de reais \ensuremath{\ge} 0 em base \ensuremath{\varphi} com dígitos 0,1, forma canônica sem '11'; bijeção com o anel Z[\ensuremath{\varphi}].
\prov{relações: usa razao\_aurea; relaciona virtualizacao; relaciona word; relaciona semiotica\_peirce\_triade; relaciona vida-morte; especializa beta\_numeracao\_metalica}
\prov{proveniência: paper \textperiodcentered{} paper\_2\_reta\_ouro.pdf, Teorema 1.1 \textperiodcentered{} fato \textperiodcentered{} confiança 1.0 \textperiodcentered{} domínio matematica \textperiodcentered{} história 10 versões no jornal}

%CRISTAL {"arestas":[{"alvo":"numeros_pisot","confianca":1.0,"epistemico":"fato","origem":"paper","rel":"usa"},{"alvo":"atomo_periodo_1","confianca":1.0,"epistemico":"fato","origem":"paper","rel":"relaciona"},{"alvo":"reta_de_ouro","confianca":1.0,"epistemico":"fato","origem":"paper","rel":"relaciona"},{"alvo":"beta_numeracao_metalica","confianca":1.0,"epistemico":"fato","origem":"wiki","rel":"explica"}],"confianca":1.0,"contraexemplos":[],"descricao":"Para β de Pisot, a expansão β-ádica de x>0 é eventualmente periódica sse x∈ℚ(β) (Frougny–Solomyak): cada base é uma lente que só 'vê' seu corpo.","epistemico":"fato","exemplos":[],"id":"bases_pisot","memoria":"semantica","meta":{"autor":"Gentil Lopes","descoberta":false,"dominio":"matematica","fonte":"paper_22_reais.pdf, Teo.1.1"},"origem":"paper","palavras_chave":[],"sinonimos":[],"tipo":"conceito","titulo":"Bases de Pisot"}
\section{Bases de Pisot}
Para \ensuremath{\beta} de Pisot, a expansão \ensuremath{\beta}-ádica de x>0 é eventualmente periódica sse x\ensuremath{\in}\ensuremath{\mathbb{Q}}(\ensuremath{\beta}) (Frougny–Solomyak): cada base é uma lente que só 'vê' seu corpo.
\prov{relações: usa numeros\_pisot; relaciona atomo\_periodo\_1; relaciona reta\_de\_ouro; explica beta\_numeracao\_metalica}
\prov{proveniência: paper \textperiodcentered{} paper\_22\_reais.pdf, Teo.1.1 \textperiodcentered{} fato \textperiodcentered{} confiança 1.0 \textperiodcentered{} domínio matematica \textperiodcentered{} história 6 versões no jornal}

%CRISTAL {"arestas":[{"alvo":"algebra_posicional_metalica","confianca":1.0,"epistemico":"fato","origem":"wiki","rel":"usa"},{"alvo":"numeros_pisot","confianca":1.0,"epistemico":"fato","origem":"wiki","rel":"usa"},{"alvo":"retas_dobra_lucas_pell","confianca":1.0,"epistemico":"fato","origem":"wiki","rel":"relaciona"},{"alvo":"tiffany","confianca":1.0,"epistemico":"fato","origem":"wiki","rel":"parte_de"}],"confianca":1.0,"contraexemplos":[],"descricao":"O núcleo invertível da transformada: cada número tem uma β-expansão na base σ_n (a sequência metálica como 'pesos'), com dígitos 0..n e regra de carry. É bem-definida e única porque σ_n é PISOT (|conjugado −1/σ_n|<1): a Property F garante que toda expansão FECHA (finita). O ouro (n=1) é o caso puro — a numeração de Zeckendorf (dígitos binários, sem dois consecutivos, base Fibonacci). A prata (n=2) usa dígitos 0..2 (base Pell), etc. Forward = expandir (a posição/moeda são os dígitos); inverse = somar os pesos. Verificado: round-trip com o conjunto de dígitos correto; σ·conj=−1. É a transformada de numeração (não ortogonal — de substituição). [I] o forward/inverse da transformada metálica.","epistemico":"inferencia","exemplos":[],"id":"beta_numeracao_metalica","memoria":"semantica","meta":{"autor":"Lopes/broca-so","dominio":"algebra","fonte":"Zeckendorf (Fibonacci); β-expansão base Pisot; Property F; goldchain/metal_pos.py (sequência metálica)"},"origem":"wiki","palavras_chave":[],"sinonimos":[],"tipo":"conceito","titulo":"β-numeração metálica: a transformada é INVERTÍVEL/única por ser Pisot (Property F)"}
\section{\ensuremath{\beta}-numeração metálica: a transformada é INVERTÍVEL/única por ser Pisot (Property F)}
O núcleo invertível da transformada: cada número tem uma \ensuremath{\beta}-expansão na base \ensuremath{\sigma}\_n (a sequência metálica como 'pesos'), com dígitos 0..n e regra de carry. É bem-definida e única porque \ensuremath{\sigma}\_n é PISOT (|conjugado \ensuremath{-}1/\ensuremath{\sigma}\_n|<1): a Property F garante que toda expansão FECHA (finita). O ouro (n=1) é o caso puro — a numeração de Zeckendorf (dígitos binários, sem dois consecutivos, base Fibonacci). A prata (n=2) usa dígitos 0..2 (base Pell), etc. Forward = expandir (a posição/moeda são os dígitos); inverse = somar os pesos. Verificado: round-trip com o conjunto de dígitos correto; \ensuremath{\sigma}\textperiodcentered conj=\ensuremath{-}1. É a transformada de numeração (não ortogonal — de substituição). [I] o forward/inverse da transformada metálica.
\prov{relações: usa algebra\_posicional\_metalica; usa numeros\_pisot; relaciona retas\_dobra\_lucas\_pell; parte\_de tiffany}
\prov{proveniência: wiki \textperiodcentered{} Zeckendorf (Fibonacci); \ensuremath{\beta}-expansão base Pisot; Property F; goldchain/metal\_pos.py (sequência metálica) \textperiodcentered{} inferencia \textperiodcentered{} confiança 1.0 \textperiodcentered{} domínio algebra \textperiodcentered{} história 35 versões no jornal}

%CRISTAL {"arestas":[{"alvo":"cantor_ouro","confianca":1.0,"epistemico":"fato","origem":"paper","rel":"relaciona"},{"alvo":"osso_carne","confianca":1.0,"epistemico":"fato","origem":"paper","rel":"usa"}],"confianca":1.0,"contraexemplos":[],"descricao":"Um conjunto contém cópias reescaladas de si indefinidamente: a métrica se move (diâmetro 3⁻ⁿ→0) mas a topologia não (dim log2/log3 idêntica em toda camada; homeomorfas por Brouwer).","epistemico":"fato","exemplos":[],"id":"boneca_russa","memoria":"semantica","meta":{"autor":"Gentil Lopes","descoberta":false,"dominio":"matematica","fonte":"paper_72_boneca_russa.pdf"},"origem":"paper","palavras_chave":[],"sinonimos":[],"tipo":"conceito","titulo":"Boneca Russa (autossimilaridade aninhada)"}
\section{Boneca Russa (autossimilaridade aninhada)}
Um conjunto contém cópias reescaladas de si indefinidamente: a métrica se move (diâmetro 3\ensuremath{^{-n}}\ensuremath{\to}0) mas a topologia não (dim log2/log3 idêntica em toda camada; homeomorfas por Brouwer).
\prov{relações: relaciona cantor\_ouro; usa osso\_carne}
\prov{proveniência: paper \textperiodcentered{} paper\_72\_boneca\_russa.pdf \textperiodcentered{} fato \textperiodcentered{} confiança 1.0 \textperiodcentered{} domínio matematica \textperiodcentered{} história 3 versões no jornal}

%CRISTAL {"arestas":[{"alvo":"torre_de_hurwitz_mecanica","confianca":1.0,"epistemico":"fato","origem":"wiki","rel":"parte_de"},{"alvo":"gato","confianca":1.0,"epistemico":"fato","origem":"wiki","rel":"usa"},{"alvo":"algebras_hipercomplexas","confianca":1.0,"epistemico":"fato","origem":"wiki","rel":"usa"}],"confianca":1.0,"contraexemplos":[],"descricao":"A construção de Cayley-Dickson dobra a dimensão a cada passo — (a,b), com produto (ac−d̄b, da+bc̄) — e cada álgebra CONTÉM a anterior (o subespaço b=0): ℝ⊂ℂ⊂ℍ⊂𝕆⊂𝕊. É uma boneca russa, e a dobra é o mesmo empacotamento de 2 em 1 do gato: a torre é uma pilha de gatos.","epistemico":"fato","exemplos":[],"id":"boneca_russa_cayley_dickson","memoria":"semantica","meta":{"dominio":"reta mineral","fonte":"torre de Hurwitz"},"origem":"wiki","palavras_chave":[],"sinonimos":[],"tipo":"conceito","titulo":"A Boneca Russa (Cayley-Dickson)"}
\section{A Boneca Russa (Cayley-Dickson)}
A construção de Cayley-Dickson dobra a dimensão a cada passo — (a,b), com produto (ac\ensuremath{-}d\ensuremath{}b, da+bc\ensuremath{}) — e cada álgebra CONTÉM a anterior (o subespaço b=0): \ensuremath{\mathbb{R}}\ensuremath{\subset}\ensuremath{\mathbb{C}}\ensuremath{\subset}\ensuremath{\mathbb{H}}\ensuremath{\subset}\ensuremath{\mathbb{O}}\ensuremath{\subset}\ensuremath{\mathbb{S}}. É uma boneca russa, e a dobra é o mesmo empacotamento de 2 em 1 do gato: a torre é uma pilha de gatos.
\prov{relações: parte\_de torre\_de\_hurwitz\_mecanica; usa gato; usa algebras\_hipercomplexas}
\prov{proveniência: wiki \textperiodcentered{} torre de Hurwitz \textperiodcentered{} fato \textperiodcentered{} confiança 1.0 \textperiodcentered{} domínio reta mineral \textperiodcentered{} história 4 versões no jornal}

%CRISTAL {"arestas":[{"alvo":"mult_prtsu","confianca":1.0,"epistemico":"fato","origem":"paper","rel":"usa"},{"alvo":"word","confianca":0.5,"epistemico":"fato","origem":"wiki","rel":"relaciona"}],"confianca":1.0,"contraexemplos":[],"descricao":"Simetria contínua no bloco real (p,r) que preserva o produto z·w — o análogo algébrico do boost de Lorentz, mas o parâmetro depende dos operandos, não do observador.","epistemico":"fato","exemplos":[],"id":"boost_algebrico","memoria":"semantica","meta":{"autor":"Gentil Lopes","descoberta":false,"dominio":"matematica","fonte":"paper_A_algebra.pdf, Teo.7.21"},"origem":"paper","palavras_chave":[],"sinonimos":[],"tipo":"conceito","titulo":"Boost Algébrico"}
\section{Boost Algébrico}
Simetria contínua no bloco real (p,r) que preserva o produto z\textperiodcentered w — o análogo algébrico do boost de Lorentz, mas o parâmetro depende dos operandos, não do observador.
\prov{relações: usa mult\_prtsu; relaciona word}
\prov{proveniência: paper \textperiodcentered{} paper\_A\_algebra.pdf, Teo.7.21 \textperiodcentered{} fato \textperiodcentered{} confiança 1.0 \textperiodcentered{} domínio matematica \textperiodcentered{} história 3 versões no jornal}

%CRISTAL {"arestas":[{"alvo":"cuspide_pisot","confianca":1.0,"epistemico":"fato","origem":"paper","rel":"usa"},{"alvo":"ising_1d_critico","confianca":1.0,"epistemico":"fato","origem":"paper","rel":"usa"},{"alvo":"transicao_fase_lambda1","confianca":1.0,"epistemico":"fato","origem":"paper","rel":"relaciona"},{"alvo":"discriminante_faces","confianca":1.0,"epistemico":"fato","origem":"paper","rel":"relaciona"},{"alvo":"ponto_auto_dual","confianca":1.0,"epistemico":"fato","origem":"paper","rel":"relaciona"},{"alvo":"colapso","confianca":1.0,"epistemico":"fato","origem":"paper","rel":"referencia"}],"confianca":1.0,"contraexemplos":[],"descricao":"A grande síntese: o expoente de Lyapunov L=0 (borda do caos, Feigenbaum r≈3,5699) é a mesma coisa que gap=0 ≡ λ₁=λ₂ ≡ T_c ≡ discriminante=0 ≡ cúspide. Cúspide é vazio.","epistemico":"fato","exemplos":[],"id":"borda_critica_vazio","memoria":"semantica","meta":{"autor":"Gentil Lopes","descoberta":true,"dominio":"matematica","fonte":"paper_86_equilibrio_caos.pdf, Teo.4.1"},"origem":"paper","palavras_chave":[],"sinonimos":[],"tipo":"conceito","titulo":"Borda Crítica = Vazio"}
\section{Borda Crítica = Vazio}
A grande síntese: o expoente de Lyapunov L=0 (borda do caos, Feigenbaum r\ensuremath{\approx}3,5699) é a mesma coisa que gap=0 \ensuremath{\equiv} \ensuremath{\lambda}\ensuremath{_{1}}=\ensuremath{\lambda}\ensuremath{_{2}} \ensuremath{\equiv} T\_c \ensuremath{\equiv} discriminante=0 \ensuremath{\equiv} cúspide. Cúspide é vazio.
\prov{relações: usa cuspide\_pisot; usa ising\_1d\_critico; relaciona transicao\_fase\_lambda1; relaciona discriminante\_faces; relaciona ponto\_auto\_dual; referencia colapso}
\prov{proveniência: paper \textperiodcentered{} paper\_86\_equilibrio\_caos.pdf, Teo.4.1 \textperiodcentered{} fato \textperiodcentered{} confiança 1.0 \textperiodcentered{} domínio matematica \textperiodcentered{} história 7 versões no jornal}

%CRISTAL {"arestas":[{"alvo":"goldchain","confianca":1.0,"epistemico":"fato","origem":"wiki","rel":"relaciona"},{"alvo":"floco_phi","confianca":1.0,"epistemico":"fato","origem":"wiki","rel":"usa"},{"alvo":"pacote_seed_pow","confianca":1.0,"epistemico":"fato","origem":"wiki","rel":"relaciona"}],"confianca":1.0,"contraexemplos":[],"descricao":"Um ciclo real de tesouraria (Pell #150, mainnet 2026-06-29) gera três atratores mod p distintos sobre o mesmo evento — trabalho PoW (666), banda HTLC (554), contrato JSON (374); 'make borda-investigar' confirma atrator_C = atrator_borda(trabalho) e setor algébrico invariante com transição binária.","epistemico":"fato","exemplos":[],"id":"borda_goldchain_mainnet","memoria":"semantica","meta":{"autor":"broca-so","dominio":"fractal","fonte":"papers/floco_de_neve.tex §IV.1"},"origem":"wiki","palavras_chave":[],"sinonimos":[],"tipo":"conceito","titulo":"Borda GoldChain — Ciclo Mainnet Verificado"}
\section{Borda GoldChain — Ciclo Mainnet Verificado}
Um ciclo real de tesouraria (Pell \#150, mainnet 2026-06-29) gera três atratores mod p distintos sobre o mesmo evento — trabalho PoW (666), banda HTLC (554), contrato JSON (374); 'make borda-investigar' confirma atrator\_C = atrator\_borda(trabalho) e setor algébrico invariante com transição binária.
\prov{relações: relaciona goldchain; usa floco\_phi; relaciona pacote\_seed\_pow}
\prov{proveniência: wiki \textperiodcentered{} papers/floco\_de\_neve.tex §IV.1 \textperiodcentered{} fato \textperiodcentered{} confiança 1.0 \textperiodcentered{} domínio fractal \textperiodcentered{} história 4 versões no jornal}

%CRISTAL {"arestas":[{"alvo":"decoerencia_e_sair_da_borda","confianca":1.0,"epistemico":"fato","origem":"wiki","rel":"parte_de"},{"alvo":"numero_de_salem","confianca":1.0,"epistemico":"fato","origem":"wiki","rel":"explica"},{"alvo":"parabola_e_o_regime_quantico","confianca":1.0,"epistemico":"fato","origem":"wiki","rel":"relaciona"}],"confianca":1.0,"contraexemplos":[],"descricao":"A borda (Salem, ρ=1) é o subespaço LIVRE DE DECOERÊNCIA: enquanto a evolução é unitária (Γ=0), a coerência se preserva para sempre. É a mesma razão pela qual a vida, a computação coerente e a informação moram na borda — é o único regime fechado onde nada vaza. A decoerência empurra o sistema para fora da borda; ficar coerente é ficar na borda.","epistemico":"inferencia","exemplos":[],"id":"borda_livre_de_decoerencia","memoria":"semantica","meta":{"dominio":"reta mineral","fonte":"emaranhamento e decoerência"},"origem":"wiki","palavras_chave":[],"sinonimos":[],"tipo":"conceito","titulo":"A Borda é o Subespaço Livre de Decoerência"}
\section{A Borda é o Subespaço Livre de Decoerência}
A borda (Salem, \ensuremath{\rho}=1) é o subespaço LIVRE DE DECOERÊNCIA: enquanto a evolução é unitária (\ensuremath{\Gamma}=0), a coerência se preserva para sempre. É a mesma razão pela qual a vida, a computação coerente e a informação moram na borda — é o único regime fechado onde nada vaza. A decoerência empurra o sistema para fora da borda; ficar coerente é ficar na borda.
\prov{relações: parte\_de decoerencia\_e\_sair\_da\_borda; explica numero\_de\_salem; relaciona parabola\_e\_o\_regime\_quantico}
\prov{proveniência: wiki \textperiodcentered{} emaranhamento e decoerência \textperiodcentered{} inferencia \textperiodcentered{} confiança 1.0 \textperiodcentered{} domínio reta mineral \textperiodcentered{} história 4 versões no jornal}

%CRISTAL {"arestas":[{"alvo":"mecanica_quantica_mineral","confianca":1.0,"epistemico":"fato","origem":"wiki","rel":"parte_de"},{"alvo":"colapso","confianca":1.0,"epistemico":"fato","origem":"wiki","rel":"eh_um"},{"alvo":"bai","confianca":1.0,"epistemico":"fato","origem":"wiki","rel":"usa"}],"confianca":1.0,"contraexemplos":[],"descricao":"O colapso Ψ que atravessou a máquina (o minimax do xadrez, a rotulagem da segmentação, o reconhecimento do glifo, o diagnóstico) É a MEDIDA QUÂNTICA: seleciona um ramo com probabilidade |ψ_i|² (regra de Born). O orçamento δ do BAI é o QUANTUM DE AÇÃO (o ℏ da autorização, a profundidade indivisível até o fecho); o deny-overrides é o colapso ao subespaço permitido. Verif.: [0.6,0.8i]→[0.36,0.64].","epistemico":"fato","exemplos":[],"id":"born_e_o_colapso_do_bai","memoria":"semantica","meta":{"dominio":"reta mineral","fonte":"mecânica quântica mineral"},"origem":"wiki","palavras_chave":[],"sinonimos":[],"tipo":"conceito","titulo":"Born e o Colapso do BAI como Medida"}
\section{Born e o Colapso do BAI como Medida}
O colapso \ensuremath{\Psi} que atravessou a máquina (o minimax do xadrez, a rotulagem da segmentação, o reconhecimento do glifo, o diagnóstico) É a MEDIDA QUÂNTICA: seleciona um ramo com probabilidade |\ensuremath{\psi}\_i|\ensuremath{^{2}} (regra de Born). O orçamento \ensuremath{\delta} do BAI é o QUANTUM DE AÇÃO (o \ensuremath{\hbar} da autorização, a profundidade indivisível até o fecho); o deny-overrides é o colapso ao subespaço permitido. Verif.: [0.6,0.8i]\ensuremath{\to}[0.36,0.64].
\prov{relações: parte\_de mecanica\_quantica\_mineral; eh\_um colapso; usa bai}
\prov{proveniência: wiki \textperiodcentered{} mecânica quântica mineral \textperiodcentered{} fato \textperiodcentered{} confiança 1.0 \textperiodcentered{} domínio reta mineral \textperiodcentered{} história 4 versões no jornal}

%CRISTAL {"arestas":[{"alvo":"curvatura_algebrica","confianca":1.0,"epistemico":"fato","origem":"paper","rel":"usa"},{"alvo":"ponte_schur_weyl","confianca":1.0,"epistemico":"fato","origem":"paper","rel":"usa"}],"confianca":1.0,"contraexemplos":[],"descricao":"A curvatura algébrica dá cota rigorosa C(s)²=36(1−s²)² para a fração do gradiente que vaza do subespaço S₃-equivariante: máxima em s=0, nula em s=±1.","epistemico":"fato","exemplos":[],"id":"bound_leakage","memoria":"semantica","meta":{"autor":"Gentil Lopes","descoberta":false,"dominio":"matematica","fonte":"paper_AE_ponte_schur_weyl.pdf, Teo.5.4"},"origem":"paper","palavras_chave":[],"sinonimos":[],"tipo":"conceito","titulo":"Bound de Leakage C(s)²=36(1−s²)²"}
\section{Bound de Leakage C(s)\ensuremath{^{2}}=36(1\ensuremath{-}s\ensuremath{^{2}})\ensuremath{^{2}}}
A curvatura algébrica dá cota rigorosa C(s)\ensuremath{^{2}}=36(1\ensuremath{-}s\ensuremath{^{2}})\ensuremath{^{2}} para a fração do gradiente que vaza do subespaço S\ensuremath{_{3}}-equivariante: máxima em s=0, nula em s=$\pm$1.
\prov{relações: usa curvatura\_algebrica; usa ponte\_schur\_weyl}
\prov{proveniência: paper \textperiodcentered{} paper\_AE\_ponte\_schur\_weyl.pdf, Teo.5.4 \textperiodcentered{} fato \textperiodcentered{} confiança 1.0 \textperiodcentered{} domínio matematica \textperiodcentered{} história 3 versões no jornal}

%CRISTAL {"arestas":[{"alvo":"espaco_difusivo","confianca":1.0,"epistemico":"fato","origem":"wiki","rel":"usa"}],"confianca":1.0,"contraexemplos":[],"descricao":"A fração da CPU declarada liberada/difusível; default 0.90, clampada em [0.01, 1.0]. Escala diretamente gpu_blocos (=max(128, frac·4096)), passos, rodadas e chunk. O parâmetro-mestre da difusão.","epistemico":"fato","exemplos":[],"id":"broca_difusao_frac","memoria":"semantica","meta":{"autor":"broca-so","dominio":"fractal","fonte":"neuronio/difusao_recursos.py:frac_liberada"},"origem":"wiki","palavras_chave":[],"sinonimos":[],"tipo":"conceito","titulo":"BROCA_DIFUSAO_FRAC = 0.90 (a fração liberada)"}
\section{BROCA\_DIFUSAO\_FRAC = 0.90 (a fração liberada)}
A fração da CPU declarada liberada/difusível; default 0.90, clampada em [0.01, 1.0]. Escala diretamente gpu\_blocos (=max(128, frac\textperiodcentered 4096)), passos, rodadas e chunk. O parâmetro-mestre da difusão.
\prov{relações: usa espaco\_difusivo}
\prov{proveniência: wiki \textperiodcentered{} neuronio/difusao\_recursos.py:frac\_liberada \textperiodcentered{} fato \textperiodcentered{} confiança 1.0 \textperiodcentered{} domínio fractal \textperiodcentered{} história 2 versões no jornal}

%CRISTAL {"arestas":[{"alvo":"teorema_esterilidade_dimensao","confianca":1.0,"epistemico":"fato","origem":"wiki","rel":"usa"},{"alvo":"verificacao_ouro_total","confianca":1.0,"epistemico":"fato","origem":"wiki","rel":"usa"},{"alvo":"verificacao_prata_parcial","confianca":1.0,"epistemico":"fato","origem":"wiki","rel":"relaciona"},{"alvo":"caixa_de_prata_corpo","confianca":1.0,"epistemico":"fato","origem":"wiki","rel":"relaciona"},{"alvo":"missao_engenharia_ouro","confianca":1.0,"epistemico":"fato","origem":"wiki","rel":"usa"}],"confianca":1.0,"contraexemplos":[],"descricao":"O broca-so implementa a CAIXA DE OURO (ℤ[φ], φ²=φ+1): Fibonacci, o golden mean shift ('proibir 11'), a impressão Koch, a mineração por 5 modos (pentagrama) — visão TOTAL da reta de ouro. Pelo teorema da esterilidade, mapeada a caixa de ouro a informação recircula estéril entre as retas real e de ouro. A PRÓXIMA dimensão a mapear é a de PRATA (corpo ℚ(√2), σ₂=1+√2, Pell), com mineração honesta por 6 modos (hexagrama de Salomão) — hoje visão PARCIAL. É a Missão aplicada às dimensões: esgotado o ganho de uma, seguir para a próxima.","epistemico":"inferencia","exemplos":[],"id":"broca_so_caixa_de_ouro","memoria":"semantica","meta":{"autor":"Lopes/broca-so","dominio":"fractal","fonte":"broca-so (golden mean shift, 5 modos); machine/so/prata.py (Pell, 6 modos)"},"origem":"wiki","palavras_chave":[],"sinonimos":[],"tipo":"conceito","titulo":"O broca-so implementa a caixa de ouro (mapeada → próxima: prata)"}
\section{O broca-so implementa a caixa de ouro (mapeada \ensuremath{\to} próxima: prata)}
O broca-so implementa a CAIXA DE OURO (\ensuremath{\mathbb{Z}}[\ensuremath{\varphi}], \ensuremath{\varphi}\ensuremath{^{2}}=\ensuremath{\varphi}+1): Fibonacci, o golden mean shift ('proibir 11'), a impressão Koch, a mineração por 5 modos (pentagrama) — visão TOTAL da reta de ouro. Pelo teorema da esterilidade, mapeada a caixa de ouro a informação recircula estéril entre as retas real e de ouro. A PRÓXIMA dimensão a mapear é a de PRATA (corpo \ensuremath{\mathbb{Q}}(\ensuremath{\sqrt{\,}}2), \ensuremath{\sigma}\ensuremath{_{2}}=1+\ensuremath{\sqrt{\,}}2, Pell), com mineração honesta por 6 modos (hexagrama de Salomão) — hoje visão PARCIAL. É a Missão aplicada às dimensões: esgotado o ganho de uma, seguir para a próxima.
\prov{relações: usa teorema\_esterilidade\_dimensao; usa verificacao\_ouro\_total; relaciona verificacao\_prata\_parcial; relaciona caixa\_de\_prata\_corpo; usa missao\_engenharia\_ouro}
\prov{proveniência: wiki \textperiodcentered{} broca-so (golden mean shift, 5 modos); machine/so/prata.py (Pell, 6 modos) \textperiodcentered{} inferencia \textperiodcentered{} confiança 1.0 \textperiodcentered{} domínio fractal \textperiodcentered{} história 6 versões no jornal}

%CRISTAL {"arestas":[{"alvo":"gap_metalico","confianca":1.0,"epistemico":"fato","origem":"paper","rel":"usa"},{"alvo":"word","confianca":0.5,"epistemico":"fato","origem":"wiki","rel":"relaciona"}],"confianca":1.0,"contraexemplos":[],"descricao":"Sugestão do Lopes: o posto de Mordell-Weil de E seria o número de direções livres que partem do cruzamento do gap, e Rₙ=log σₙ o regulador. Correspondência, não prova.","epistemico":"hipotese","exemplos":[],"id":"bsd_gap","memoria":"semantica","meta":{"autor":"Gentil Lopes","descoberta":false,"dominio":"matematica","fonte":"corpo_tropical.pdf, §2.12"},"origem":"paper","palavras_chave":[],"sinonimos":[],"tipo":"conceito","titulo":"BSD via Gap"}
\section{BSD via Gap}
Sugestão do Lopes: o posto de Mordell-Weil de E seria o número de direções livres que partem do cruzamento do gap, e R\ensuremath{_{n}}=log \ensuremath{\sigma}\ensuremath{_{n}} o regulador. Correspondência, não prova.
\prov{relações: usa gap\_metalico; relaciona word}
\prov{proveniência: paper \textperiodcentered{} corpo\_tropical.pdf, §2.12 \textperiodcentered{} hipotese \textperiodcentered{} confiança 1.0 \textperiodcentered{} domínio matematica \textperiodcentered{} história 4 versões no jornal}

%CRISTAL {"arestas":[{"alvo":"pinn_surrogate","confianca":1.0,"epistemico":"fato","origem":"paper","rel":"usa"}],"confianca":1.0,"contraexemplos":[],"descricao":"A equação de Burgers viscosa (protótipo 1D da turbulência: formação de choque, cascata) via Hopf-Cole; PINN acerta 1-step (~5% erro, 6,4× speedup) mas acumula erro no rollout longo.","epistemico":"fato","exemplos":[],"id":"burgers_mvp","memoria":"semantica","meta":{"autor":"Gentil Lopes","descoberta":false,"dominio":"matematica","fonte":"paper_BN.pdf"},"origem":"paper","palavras_chave":[],"sinonimos":[],"tipo":"conceito","titulo":"Burgers como MVP"}
\section{Burgers como MVP}
A equação de Burgers viscosa (protótipo 1D da turbulência: formação de choque, cascata) via Hopf-Cole; PINN acerta 1-step (\~{}5\% erro, 6,4$\times$ speedup) mas acumula erro no rollout longo.
\prov{relações: usa pinn\_surrogate}
\prov{proveniência: paper \textperiodcentered{} paper\_BN.pdf \textperiodcentered{} fato \textperiodcentered{} confiança 1.0 \textperiodcentered{} domínio matematica \textperiodcentered{} história 2 versões no jornal}

%CRISTAL {"arestas":[{"alvo":"escala_volta_gentil","confianca":1.0,"epistemico":"fato","origem":"wiki","rel":"usa"},{"alvo":"energia_salto_hbar_raiz2","confianca":1.0,"epistemico":"fato","origem":"wiki","rel":"relaciona"}],"confianca":1.0,"contraexemplos":[],"descricao":"c=1 é a física; 299.792.458 é a roupa (convenção histórica p/ o metro coincidir com o Kr-86). O metro é uma volta; o segundo é o tempo de c voltas; espaço e tempo são ambos voltas, c os relaciona. Uma velocidade [L][T]⁻¹ tem dimensão de massa d=0 (adimensional): é potência ZERO do espaço-tempo, uma taxa v/c (voltas espaciais por volta temporal). 'm/s' não é dimensão no sistema algébrico — é um número.","epistemico":"inferencia","exemplos":[],"id":"c_contagem_nao_velocidade","memoria":"semantica","meta":{"autor":"Lopes/broca-so","dominio":"unidades","fonte":"hiper/circular/paper_5_unidades.tex"},"origem":"wiki","palavras_chave":[],"sinonimos":[],"tipo":"conceito","titulo":"c é uma contagem, não uma velocidade (v = potência zero)"}
\section{c é uma contagem, não uma velocidade (v = potência zero)}
c=1 é a física; 299.792.458 é a roupa (convenção histórica p/ o metro coincidir com o Kr-86). O metro é uma volta; o segundo é o tempo de c voltas; espaço e tempo são ambos voltas, c os relaciona. Uma velocidade [L][T]\ensuremath{^{-1}} tem dimensão de massa d=0 (adimensional): é potência ZERO do espaço-tempo, uma taxa v/c (voltas espaciais por volta temporal). 'm/s' não é dimensão no sistema algébrico — é um número.
\prov{relações: usa escala\_volta\_gentil; relaciona energia\_salto\_hbar\_raiz2}
\prov{proveniência: wiki \textperiodcentered{} hiper/circular/paper\_5\_unidades.tex \textperiodcentered{} inferencia \textperiodcentered{} confiança 1.0 \textperiodcentered{} domínio unidades \textperiodcentered{} história 4 versões no jornal}

%CRISTAL {"arestas":[{"alvo":"equilibrio_de_nash","confianca":1.0,"epistemico":"fato","origem":"wiki","rel":"usa"}],"confianca":1.0,"contraexemplos":[],"descricao":"Cooperar caça o cervo (melhor para todos), mas exige confiança; desertar pega a lebre (seguro, pior). O jogo da COORDENAÇÃO e da confiança — dois equilíbrios, um bom e um medroso.","epistemico":"fato","exemplos":[],"id":"caca_ao_cervo","memoria":"semantica","meta":{"dominio":"matematica","fonte":"teoria dos jogos"},"origem":"wiki","palavras_chave":[],"sinonimos":[],"tipo":"conceito","titulo":"Caça ao Cervo (Stag Hunt)"}
\section{Caça ao Cervo (Stag Hunt)}
Cooperar caça o cervo (melhor para todos), mas exige confiança; desertar pega a lebre (seguro, pior). O jogo da COORDENAÇÃO e da confiança — dois equilíbrios, um bom e um medroso.
\prov{relações: usa equilibrio\_de\_nash}
\prov{proveniência: wiki \textperiodcentered{} teoria dos jogos \textperiodcentered{} fato \textperiodcentered{} confiança 1.0 \textperiodcentered{} domínio matematica \textperiodcentered{} história 2 versões no jornal}

%CRISTAL {"arestas":[{"alvo":"torre_de_hurwitz_mecanica","confianca":1.0,"epistemico":"fato","origem":"wiki","rel":"parte_de"},{"alvo":"numero_de_salem","confianca":1.0,"epistemico":"fato","origem":"wiki","rel":"usa"},{"alvo":"g2_grupo_excepcional","confianca":1.0,"epistemico":"fato","origem":"wiki","rel":"usa"},{"alvo":"mass_gap_mineral","confianca":1.0,"epistemico":"fato","origem":"wiki","rel":"relaciona"}],"confianca":1.0,"contraexemplos":[],"descricao":"A esfera unitária de cada álgebra carrega um grupo, e com ele uma física: ℝ→O(1) (sem fase); ℂ→U(1) (a FASE, o círculo ρ=1, a borda/Salem — a mecânica quântica precisa de ℂ para ter interferência e evolução unitária); ℍ→SU(2) (o SPIN, a não-comutatividade); 𝕆→G₂ (o mass gap 𝔤₂=14, a não-associatividade). Subir a torre é ganhar física perdendo lei.","epistemico":"fato","exemplos":[],"id":"cada_andar_uma_mecanica","memoria":"semantica","meta":{"dominio":"reta mineral","fonte":"torre de Hurwitz"},"origem":"wiki","palavras_chave":[],"sinonimos":[],"tipo":"conceito","titulo":"Cada Andar é uma Mecânica"}
\section{Cada Andar é uma Mecânica}
A esfera unitária de cada álgebra carrega um grupo, e com ele uma física: \ensuremath{\mathbb{R}}\ensuremath{\to}O(1) (sem fase); \ensuremath{\mathbb{C}}\ensuremath{\to}U(1) (a FASE, o círculo \ensuremath{\rho}=1, a borda/Salem — a mecânica quântica precisa de \ensuremath{\mathbb{C}} para ter interferência e evolução unitária); \ensuremath{\mathbb{H}}\ensuremath{\to}SU(2) (o SPIN, a não-comutatividade); \ensuremath{\mathbb{O}}\ensuremath{\to}G\ensuremath{_{2}} (o mass gap \ensuremath{\mathfrak{g}}\ensuremath{_{2}}=14, a não-associatividade). Subir a torre é ganhar física perdendo lei.
\prov{relações: parte\_de torre\_de\_hurwitz\_mecanica; usa numero\_de\_salem; usa g2\_grupo\_excepcional; relaciona mass\_gap\_mineral}
\prov{proveniência: wiki \textperiodcentered{} torre de Hurwitz \textperiodcentered{} fato \textperiodcentered{} confiança 1.0 \textperiodcentered{} domínio reta mineral \textperiodcentered{} história 5 versões no jornal}

%CRISTAL {"arestas":[{"alvo":"catedral_sigma_aritmetica","confianca":1.0,"epistemico":"fato","origem":"wiki","rel":"parte_de"},{"alvo":"fano_psl_168","confianca":1.0,"epistemico":"fato","origem":"wiki","rel":"usa"},{"alvo":"rede_leech","confianca":1.0,"epistemico":"fato","origem":"wiki","rel":"relaciona"}],"confianca":1.0,"contraexemplos":[],"descricao":"Iterando σ desde o primo de Fano: σ(7)=8 (dim octônica), σ(8)=15, σ(15)=24 (dim de Leech), σ(24)=60 (|A₅| icosaedral), σ(60)=168 (a Muralha PSL), σ(168)=480 (o classificador pleno), σ(480)=1512 (9×168). Cada número é uma constante da Catedral; oito entradas para a oitava parte.","epistemico":"fato","exemplos":[],"id":"cadeia_de_fano_sigma","memoria":"semantica","meta":{"autor":"Lord Index (Shawn R. Schiller)","dominio":"matematica","fonte":"A Catedral — Part VIII-A (jul 2026), ~/Imagens/catedral"},"origem":"wiki","palavras_chave":[],"sinonimos":[],"tipo":"conceito","titulo":"A Cadeia de Fano (σ-espinha): 7→8→15→24→60→168→480→1512"}
\section{A Cadeia de Fano (\ensuremath{\sigma}-espinha): 7\ensuremath{\to}8\ensuremath{\to}15\ensuremath{\to}24\ensuremath{\to}60\ensuremath{\to}168\ensuremath{\to}480\ensuremath{\to}1512}
Iterando \ensuremath{\sigma} desde o primo de Fano: \ensuremath{\sigma}(7)=8 (dim octônica), \ensuremath{\sigma}(8)=15, \ensuremath{\sigma}(15)=24 (dim de Leech), \ensuremath{\sigma}(24)=60 (|A\ensuremath{_{5}}| icosaedral), \ensuremath{\sigma}(60)=168 (a Muralha PSL), \ensuremath{\sigma}(168)=480 (o classificador pleno), \ensuremath{\sigma}(480)=1512 (9$\times$168). Cada número é uma constante da Catedral; oito entradas para a oitava parte.
\prov{relações: parte\_de catedral\_sigma\_aritmetica; usa fano\_psl\_168; relaciona rede\_leech}
\prov{proveniência: wiki \textperiodcentered{} A Catedral — Part VIII-A (jul 2026), \~{}/Imagens/catedral \textperiodcentered{} fato \textperiodcentered{} confiança 1.0 \textperiodcentered{} domínio matematica \textperiodcentered{} história 4 versões no jornal}

%CRISTAL {"arestas":[{"alvo":"transformada_metalica","confianca":1.0,"epistemico":"fato","origem":"wiki","rel":"continua"},{"alvo":"lyapunov_da_orbita","confianca":1.0,"epistemico":"fato","origem":"wiki","rel":"usa"},{"alvo":"reta_de_ouro","confianca":1.0,"epistemico":"fato","origem":"wiki","rel":"relaciona"},{"alvo":"numeros_de_pisot","confianca":1.0,"epistemico":"fato","origem":"wiki","rel":"relaciona"},{"alvo":"gex_producao_multifractal","confianca":1.0,"epistemico":"fato","origem":"wiki","rel":"relaciona"}],"confianca":1.0,"contraexemplos":[],"descricao":"Além do quadrático (os metais x²−nx−1, todos Pisot/cristal), as cadeias — recorrências de polinômio companheiro qualquer — deixam a ilha cristal. A estatística exaustiva mede a imprevisibilidade crescer com o grau: g2: 50.0% cristal, λ̄=-0.23 · g3: 33.6% cristal, λ̄=0.046 · g4: 12.2% cristal, λ̄=0.199 · g5: 6.9% cristal, λ̄=0.253 · g6: 2.7% cristal, λ̄=0.285. A fração cristal despenca e o Lyapunov médio sobe: a ilha cristal é pequena, o caos é o mar. É a fonte de imprevisibilidade além dos metais.","epistemico":"inferencia","exemplos":[],"id":"cadeias_alem_dos_metais","memoria":"semantica","meta":{"dominio":"matematica","estatistica_por_grau":{"2":{"lam_medio":-0.23,"n":14,"pct_borda":21.4,"pct_caos":28.6,"pct_cristal":50.0},"3":{"lam_medio":0.046,"n":113,"pct_borda":14.2,"pct_caos":52.2,"pct_cristal":33.6},"4":{"lam_medio":0.199,"n":841,"pct_borda":11.2,"pct_caos":76.6,"pct_cristal":12.2},"5":{"lam_medio":0.253,"n":6216,"pct_borda":4.4,"pct_caos":88.6,"pct_cristal":6.9},"6":{"lam_medio":0.285,"n":45356,"pct_borda":2.3,"pct_caos":95.0,"pct_cristal":2.7}},"fonte":"cadeias além dos metais"},"origem":"wiki","palavras_chave":[],"sinonimos":[],"tipo":"conceito","titulo":"Cadeias Além dos Metais"}
\section{Cadeias Além dos Metais}
Além do quadrático (os metais x\ensuremath{^{2}}\ensuremath{-}nx\ensuremath{-}1, todos Pisot/cristal), as cadeias — recorrências de polinômio companheiro qualquer — deixam a ilha cristal. A estatística exaustiva mede a imprevisibilidade crescer com o grau: g2: 50.0\% cristal, \ensuremath{\lambda}\ensuremath{}=-0.23 \textperiodcentered  g3: 33.6\% cristal, \ensuremath{\lambda}\ensuremath{}=0.046 \textperiodcentered  g4: 12.2\% cristal, \ensuremath{\lambda}\ensuremath{}=0.199 \textperiodcentered  g5: 6.9\% cristal, \ensuremath{\lambda}\ensuremath{}=0.253 \textperiodcentered  g6: 2.7\% cristal, \ensuremath{\lambda}\ensuremath{}=0.285. A fração cristal despenca e o Lyapunov médio sobe: a ilha cristal é pequena, o caos é o mar. É a fonte de imprevisibilidade além dos metais.
\prov{relações: continua transformada\_metalica; usa lyapunov\_da\_orbita; relaciona reta\_de\_ouro; relaciona numeros\_de\_pisot; relaciona gex\_producao\_multifractal}
\prov{proveniência: wiki \textperiodcentered{} cadeias além dos metais \textperiodcentered{} inferencia \textperiodcentered{} confiança 1.0 \textperiodcentered{} domínio matematica \textperiodcentered{} história 11 versões no jornal}

%CRISTAL {"arestas":[{"alvo":"caixas_fractais_generalizacao","confianca":1.0,"epistemico":"fato","origem":"wiki","rel":"parte_de"},{"alvo":"entropia_topologica","confianca":1.0,"epistemico":"fato","origem":"wiki","rel":"usa"},{"alvo":"fab_cinco_estagios","confianca":1.0,"epistemico":"fato","origem":"wiki","rel":"especializa"},{"alvo":"entropia_log_phi","confianca":1.0,"epistemico":"fato","origem":"wiki","rel":"relaciona"},{"alvo":"beta_numeracao_metalica","confianca":1.0,"epistemico":"fato","origem":"wiki","rel":"relaciona"}],"confianca":1.0,"contraexemplos":[],"descricao":"O bit gerador da fábrica: subshift binário com PODA 'proibir 1k'. O nº de palavras admissíveis de tamanho L cresce como λₖL, com λₖ = maior raiz de xk=xk−¹+…+1 (constante k-bonacci); a dimensão é Dₖ=log₂ λₖ. k=2 (Fibonacci, 'proibir 11', o corte mínimo=o ouro): λ=φ, D=0,6942 (box-counting medido 0,6935); k=3 (tribonacci): D=0,8791; k=4 (tetranacci): 0,9468; k→∞: λ→2, D→1 (a poda sai, o universo enche a reta). [D]/[N] verificado.","epistemico":"fato","exemplos":[],"id":"caixa_binaria_kbonacci","memoria":"semantica","meta":{"autor":"Lopes/broca-so","dominio":"fractal","fonte":"fabrica.tex §bit gerador; paper_40_carne_fractal.tex (dim=logλ/log2)"},"origem":"wiki","palavras_chave":[],"sinonimos":[],"tipo":"conceito","titulo":"A caixa binária k-bonacci (proibir 1k): dim=log₂ λₖ"}
\section{A caixa binária k-bonacci (proibir 1k): dim=log\ensuremath{_{2}} \ensuremath{\lambda}\ensuremath{_{k}}}
O bit gerador da fábrica: subshift binário com PODA 'proibir 1k'. O nº de palavras admissíveis de tamanho L cresce como \ensuremath{\lambda}\ensuremath{_{k}}L, com \ensuremath{\lambda}\ensuremath{_{k}} = maior raiz de xk=xk\ensuremath{-}\ensuremath{^{1}}+…+1 (constante k-bonacci); a dimensão é D\ensuremath{_{k}}=log\ensuremath{_{2}} \ensuremath{\lambda}\ensuremath{_{k}}. k=2 (Fibonacci, 'proibir 11', o corte mínimo=o ouro): \ensuremath{\lambda}=\ensuremath{\varphi}, D=0,6942 (box-counting medido 0,6935); k=3 (tribonacci): D=0,8791; k=4 (tetranacci): 0,9468; k\ensuremath{\to}\ensuremath{\infty}: \ensuremath{\lambda}\ensuremath{\to}2, D\ensuremath{\to}1 (a poda sai, o universo enche a reta). [D]/[N] verificado.
\prov{relações: parte\_de caixas\_fractais\_generalizacao; usa entropia\_topologica; especializa fab\_cinco\_estagios; relaciona entropia\_log\_phi; relaciona beta\_numeracao\_metalica}
\prov{proveniência: wiki \textperiodcentered{} fabrica.tex §bit gerador; paper\_40\_carne\_fractal.tex (dim=log\ensuremath{\lambda}/log2) \textperiodcentered{} fato \textperiodcentered{} confiança 1.0 \textperiodcentered{} domínio fractal \textperiodcentered{} história 8 versões no jornal}

%CRISTAL {"arestas":[{"alvo":"computacao_fractal","confianca":1.0,"epistemico":"fato","origem":"wiki","rel":"parte_de"},{"alvo":"broca_os","confianca":1.0,"epistemico":"fato","origem":"wiki","rel":"parte_de"}],"confianca":1.0,"contraexemplos":[],"descricao":"A construção da caixa fractal (o bit gerador da fábrica: um subshift podado) generaliza a qualquer alfabeto: todo subshift Σ → um autovalor de Perron λ → entropia h=ln λ e dimensão dim=log_b λ (b=tamanho do alfabeto). Piso λ→1 (dim 0, morte), teto λ→2 no binário (dim 1, dobramento puro). O padrão verificado no ouro e parcial na prata é ESTENDIDO por meta-indução a todas as dimensões metálicas e à reta de plástico — MATEMÁTICA APENAS, dentro do limite de verificação declarado.","epistemico":"inferencia","exemplos":[],"id":"caixas_fractais_generalizacao","memoria":"semantica","meta":{"autor":"Lopes/broca-so","dominio":"fractal","fonte":"doc/fractais/caixas_fractais_metalicas.py"},"origem":"wiki","palavras_chave":[],"sinonimos":[],"tipo":"conceito","titulo":"Generalização das caixas fractais (meta-indução)"}
\section{Generalização das caixas fractais (meta-indução)}
A construção da caixa fractal (o bit gerador da fábrica: um subshift podado) generaliza a qualquer alfabeto: todo subshift \ensuremath{\Sigma} \ensuremath{\to} um autovalor de Perron \ensuremath{\lambda} \ensuremath{\to} entropia h=ln \ensuremath{\lambda} e dimensão dim=log\_b \ensuremath{\lambda} (b=tamanho do alfabeto). Piso \ensuremath{\lambda}\ensuremath{\to}1 (dim 0, morte), teto \ensuremath{\lambda}\ensuremath{\to}2 no binário (dim 1, dobramento puro). O padrão verificado no ouro e parcial na prata é ESTENDIDO por meta-indução a todas as dimensões metálicas e à reta de plástico — MATEMÁTICA APENAS, dentro do limite de verificação declarado.
\prov{relações: parte\_de computacao\_fractal; parte\_de broca\_os}
\prov{proveniência: wiki \textperiodcentered{} doc/fractais/caixas\_fractais\_metalicas.py \textperiodcentered{} inferencia \textperiodcentered{} confiança 1.0 \textperiodcentered{} domínio fractal \textperiodcentered{} história 3 versões no jornal}

%CRISTAL {"arestas":[{"alvo":"camadas_tiffany_multifractal","confianca":1.0,"epistemico":"fato","origem":"wiki","rel":"confirma"},{"alvo":"temperatura_camada_parabola","confianca":1.0,"epistemico":"fato","origem":"wiki","rel":"usa"},{"alvo":"absorcao_iterativa_poca","confianca":1.0,"epistemico":"fato","origem":"wiki","rel":"usa"}],"confianca":1.0,"contraexemplos":[],"descricao":"Verificação da generalização: roteando o purgatório numa NOVA camada L3 (folha fresca no mesmo grafo, mesma parábola, temperatura θ0=0,45 um tico mais quente), 5.741 fragmentos que estavam barrados por DENSIDADE na camada fria acharam casa — as bacias voltam ao θ inicial na folha nova. Purgatório 76.823→70.977; o tecido passou a ter 3 camadas (L1 curada 82, L2 base 25.688 em 8.356 arestas, L3 5.741 em 2.818). É a escada de temperatura em ação: cada camada drena uma banda; as próximas (L4, L5…) drenam o resto, cada conceito no seu nível. [D] roda (transicoes.absorver_iterativo camada=3).","epistemico":"fato","exemplos":[],"id":"camada_l3_drena_purgatorio","memoria":"semantica","meta":{"autor":"Lopes/broca-so","dominio":"continuo","fonte":"conversa/dados/transicoes.tsv (L1/L2/L3); transicoes.absorver_iterativo(camada=3, theta0=0.45)"},"origem":"wiki","palavras_chave":[],"sinonimos":[],"tipo":"conceito","titulo":"Prova: a camada L3 (folha fresca) drenou 5.741 do purgatório"}
\section{Prova: a camada L3 (folha fresca) drenou 5.741 do purgatório}
Verificação da generalização: roteando o purgatório numa NOVA camada L3 (folha fresca no mesmo grafo, mesma parábola, temperatura \ensuremath{\theta}0=0,45 um tico mais quente), 5.741 fragmentos que estavam barrados por DENSIDADE na camada fria acharam casa — as bacias voltam ao \ensuremath{\theta} inicial na folha nova. Purgatório 76.823\ensuremath{\to}70.977; o tecido passou a ter 3 camadas (L1 curada 82, L2 base 25.688 em 8.356 arestas, L3 5.741 em 2.818). É a escada de temperatura em ação: cada camada drena uma banda; as próximas (L4, L5…) drenam o resto, cada conceito no seu nível. [D] roda (transicoes.absorver\_iterativo camada=3).
\prov{relações: confirma camadas\_tiffany\_multifractal; usa temperatura\_camada\_parabola; usa absorcao\_iterativa\_poca}
\prov{proveniência: wiki \textperiodcentered{} conversa/dados/transicoes.tsv (L1/L2/L3); transicoes.absorver\_iterativo(camada=3, theta0=0.45) \textperiodcentered{} fato \textperiodcentered{} confiança 1.0 \textperiodcentered{} domínio continuo \textperiodcentered{} história 24 versões no jornal}

%CRISTAL {"arestas":[{"alvo":"teorema_esterilidade_dimensao","confianca":1.0,"epistemico":"fato","origem":"wiki","rel":"usa"},{"alvo":"cofre_parabola_hero","confianca":1.0,"epistemico":"fato","origem":"wiki","rel":"usa"},{"alvo":"reta_metalica_geral","confianca":1.0,"epistemico":"fato","origem":"wiki","rel":"relaciona"},{"alvo":"multifractal_construtivo","confianca":1.0,"epistemico":"fato","origem":"wiki","rel":"usa"},{"alvo":"temperatura_camada_parabola","confianca":1.0,"epistemico":"fato","origem":"wiki","rel":"usa"},{"alvo":"plano_continuo_transicoes","confianca":1.0,"epistemico":"fato","origem":"wiki","rel":"parte_de"},{"alvo":"transformada_metalica","confianca":1.0,"epistemico":"fato","origem":"wiki","rel":"generaliza"},{"alvo":"tiffany","confianca":1.0,"epistemico":"fato","origem":"wiki","rel":"parte_de"}],"confianca":1.0,"contraexemplos":[],"descricao":"Cada camada do contínuo é ESTÉRIL — completa um campo (uma dimensão mapeada, pelo teorema da esterilidade: mapeada a dimensão, a informação recircula estéril). A camada base saturou = um campo completo; não cabe mais NELA. Pra avançar cria-se OUTRA camada Tiffany = OUTRA dimensão, com a MESMA dinâmica da parábola a+c²=1, só que noutra TEMPERATURA. As camadas empilhadas no MESMO suporte (o grafo) com folhas de densidade distintas são um MULTIFRACTAL: o espectro f(α) do contínuo, cada α um nível. Assim todo conceito cabe, cada um no seu nível, ligado pela parábola — a base é o início, a temperatura muda nas outras. É o gancho que generaliza a próxima camada do so e todas as outras.","epistemico":"inferencia","exemplos":[],"id":"camadas_tiffany_multifractal","memoria":"semantica","meta":{"autor":"Lopes/broca-so","dominio":"continuo","fonte":"teorema_esterilidade_dimensao + cofre_parabola_hero + reta_metalica_geral (dimensões)"},"origem":"wiki","palavras_chave":[],"sinonimos":[],"tipo":"conceito","titulo":"As camadas da Tiffany são o espectro multifractal do contínuo (a generalização)"}
\section{As camadas da Tiffany são o espectro multifractal do contínuo (a generalização)}
Cada camada do contínuo é ESTÉRIL — completa um campo (uma dimensão mapeada, pelo teorema da esterilidade: mapeada a dimensão, a informação recircula estéril). A camada base saturou = um campo completo; não cabe mais NELA. Pra avançar cria-se OUTRA camada Tiffany = OUTRA dimensão, com a MESMA dinâmica da parábola a+c\ensuremath{^{2}}=1, só que noutra TEMPERATURA. As camadas empilhadas no MESMO suporte (o grafo) com folhas de densidade distintas são um MULTIFRACTAL: o espectro f(\ensuremath{\alpha}) do contínuo, cada \ensuremath{\alpha} um nível. Assim todo conceito cabe, cada um no seu nível, ligado pela parábola — a base é o início, a temperatura muda nas outras. É o gancho que generaliza a próxima camada do so e todas as outras.
\prov{relações: usa teorema\_esterilidade\_dimensao; usa cofre\_parabola\_hero; relaciona reta\_metalica\_geral; usa multifractal\_construtivo; usa temperatura\_camada\_parabola; parte\_de plano\_continuo\_transicoes; generaliza transformada\_metalica; parte\_de tiffany}
\prov{proveniência: wiki \textperiodcentered{} teorema\_esterilidade\_dimensao + cofre\_parabola\_hero + reta\_metalica\_geral (dimensões) \textperiodcentered{} inferencia \textperiodcentered{} confiança 1.0 \textperiodcentered{} domínio continuo \textperiodcentered{} história 79 versões no jornal}

%CRISTAL {"arestas":[{"alvo":"mult_prtsu","confianca":1.0,"epistemico":"fato","origem":"paper","rel":"parte_de"},{"alvo":"word","confianca":0.5,"epistemico":"fato","origem":"wiki","rel":"relaciona"}],"confianca":1.0,"contraexemplos":[],"descricao":"Para cada versor ã∈Sⁿ−², Cã=α+βã é subálgebra ≅ℂ (ã²=−1); ℝⁿ decompõe-se em ∪Cã, e o produto vetorial só age entre campos locais distintos.","epistemico":"fato","exemplos":[],"id":"campos_locais_ca","memoria":"semantica","meta":{"autor":"Gentil Lopes","descoberta":false,"dominio":"matematica","fonte":"paper_A_algebra.pdf, Teo.3.1"},"origem":"paper","palavras_chave":[],"sinonimos":[],"tipo":"conceito","titulo":"Campos Locais C_ã"}
\section{Campos Locais C\_ã}
Para cada versor ã\ensuremath{\in}S\ensuremath{^{n}}\ensuremath{-}\ensuremath{^{2}}, Cã=\ensuremath{\alpha}+\ensuremath{\beta}ã é subálgebra \ensuremath{\cong}\ensuremath{\mathbb{C}} (ã\ensuremath{^{2}}=\ensuremath{-}1); \ensuremath{\mathbb{R}}\ensuremath{^{n}} decompõe-se em \ensuremath{\cup}Cã, e o produto vetorial só age entre campos locais distintos.
\prov{relações: parte\_de mult\_prtsu; relaciona word}
\prov{proveniência: paper \textperiodcentered{} paper\_A\_algebra.pdf, Teo.3.1 \textperiodcentered{} fato \textperiodcentered{} confiança 1.0 \textperiodcentered{} domínio matematica \textperiodcentered{} história 4 versões no jornal}

%CRISTAL {"arestas":[{"alvo":"parabola_reparte_difusao","confianca":1.0,"epistemico":"fato","origem":"wiki","rel":"usa"},{"alvo":"cotacao_parabola","confianca":1.0,"epistemico":"fato","origem":"wiki","rel":"relaciona"},{"alvo":"conjugados_galois_parabola","confianca":1.0,"epistemico":"fato","origem":"wiki","rel":"usa"}],"confianca":1.0,"contraexemplos":[],"descricao":"A execução é gated pelo canal a+c²=1: se broca_so_exec_gated retorna BROCA_ERR_CANAL, levanta 'canal parabólico morto — difusão bloqueada'; canal_ok/canal_vivo guarda esse estado no EspacoDifusivo. Deny by default também na difusão: sem canal vivo, não difunde.","epistemico":"fato","exemplos":[],"id":"canal_parabolico_gate","memoria":"semantica","meta":{"autor":"broca-so","dominio":"fractal","fonte":"neuronio/broca_borda.py:_exec_gated"},"origem":"wiki","palavras_chave":[],"sinonimos":[],"tipo":"conceito","titulo":"O Canal Parabólico Morto Bloqueia a Difusão"}
\section{O Canal Parabólico Morto Bloqueia a Difusão}
A execução é gated pelo canal a+c\ensuremath{^{2}}=1: se broca\_so\_exec\_gated retorna BROCA\_ERR\_CANAL, levanta 'canal parabólico morto — difusão bloqueada'; canal\_ok/canal\_vivo guarda esse estado no EspacoDifusivo. Deny by default também na difusão: sem canal vivo, não difunde.
\prov{relações: usa parabola\_reparte\_difusao; relaciona cotacao\_parabola; usa conjugados\_galois\_parabola}
\prov{proveniência: wiki \textperiodcentered{} neuronio/broca\_borda.py:\_exec\_gated \textperiodcentered{} fato \textperiodcentered{} confiança 1.0 \textperiodcentered{} domínio fractal \textperiodcentered{} história 4 versões no jornal}

%CRISTAL {"arestas":[{"alvo":"hopf_fibration_vortex","confianca":1.0,"epistemico":"fato","origem":"paper","rel":"usa"},{"alvo":"turing_rogozhin_ns","confianca":1.0,"epistemico":"fato","origem":"paper","rel":"usa"}],"confianca":1.0,"contraexemplos":[],"descricao":"Sob as hipóteses do estágio 3, o Cantor invariante do fluxo decompõe-se em Σ₂×S²/G (fita de Turing × espaço de modos geométricos) — o produto lógica×geometria.","epistemico":"inferencia","exemplos":[],"id":"cantor_esferal_ns","memoria":"semantica","meta":{"autor":"Gentil Lopes","descoberta":false,"dominio":"matematica","fonte":"paper_AT.pdf, Teo.12"},"origem":"paper","palavras_chave":[],"sinonimos":[],"tipo":"conceito","titulo":"Cantor Σ₂×S² em NS"}
\section{Cantor \ensuremath{\Sigma}\ensuremath{_{2}}$\times$S\ensuremath{^{2}} em NS}
Sob as hipóteses do estágio 3, o Cantor invariante do fluxo decompõe-se em \ensuremath{\Sigma}\ensuremath{_{2}}$\times$S\ensuremath{^{2}}/G (fita de Turing $\times$ espaço de modos geométricos) — o produto lógica$\times$geometria.
\prov{relações: usa hopf\_fibration\_vortex; usa turing\_rogozhin\_ns}
\prov{proveniência: paper \textperiodcentered{} paper\_AT.pdf, Teo.12 \textperiodcentered{} inferencia \textperiodcentered{} confiança 1.0 \textperiodcentered{} domínio matematica \textperiodcentered{} história 3 versões no jornal}

%CRISTAL {"arestas":[{"alvo":"reta_de_ouro","confianca":1.0,"epistemico":"fato","origem":"livro","rel":"especializa"},{"alvo":"word","confianca":0.5,"epistemico":"fato","origem":"wiki","rel":"relaciona"},{"alvo":"koch_parabola_v2_sobre_a_curva_desdistorcida","confianca":0.5,"epistemico":"fato","origem":"wiki","rel":"relaciona"},{"alvo":"vontade_de_poder","confianca":0.5,"epistemico":"fato","origem":"wiki","rel":"relaciona"},{"alvo":"politica","confianca":0.5,"epistemico":"fato","origem":"wiki","rel":"relaciona"}],"confianca":1.0,"contraexemplos":[],"descricao":"Conjunto de Cantor em base φ — as palavras binárias infinitas sem '11'. Dimensão exata log₂ φ ≈ 0,694; medida nula; estrutura combinatória.","epistemico":"fato","exemplos":[],"id":"cantor_ouro","memoria":"semantica","meta":{"autor":"Gentil Lopes","descoberta":true,"dominio":"matematica","fonte":"paper_2_reta_ouro.pdf, Teorema 4.1"},"origem":"paper","palavras_chave":[],"sinonimos":[],"tipo":"conceito","titulo":"Cantor do Ouro"}
\section{Cantor do Ouro}
Conjunto de Cantor em base \ensuremath{\varphi} — as palavras binárias infinitas sem '11'. Dimensão exata log\ensuremath{_{2}} \ensuremath{\varphi} \ensuremath{\approx} 0,694; medida nula; estrutura combinatória.
\prov{relações: especializa reta\_de\_ouro; relaciona word; relaciona koch\_parabola\_v2\_sobre\_a\_curva\_desdistorcida; relaciona vontade\_de\_poder; relaciona politica}
\prov{proveniência: paper \textperiodcentered{} paper\_2\_reta\_ouro.pdf, Teorema 4.1 \textperiodcentered{} fato \textperiodcentered{} confiança 1.0 \textperiodcentered{} domínio matematica \textperiodcentered{} história 8 versões no jornal}

%CRISTAL {"arestas":[{"alvo":"transformada_metalica","confianca":1.0,"epistemico":"fato","origem":"wiki","rel":"usa"},{"alvo":"algebra_reta_mineral","confianca":1.0,"epistemico":"fato","origem":"wiki","rel":"relaciona"},{"alvo":"numeros_de_pisot","confianca":1.0,"epistemico":"fato","origem":"wiki","rel":"relaciona"},{"alvo":"reta_de_ouro","confianca":1.0,"epistemico":"fato","origem":"wiki","rel":"relaciona"},{"alvo":"corpo_primos_mineral","confianca":1.0,"epistemico":"fato","origem":"wiki","rel":"usa"},{"alvo":"cifra_metais_an","confianca":1.0,"epistemico":"fato","origem":"wiki","rel":"relaciona"},{"alvo":"fracao_continua_phi","confianca":1.0,"epistemico":"fato","origem":"wiki","rel":"usa"}],"confianca":1.0,"contraexemplos":[],"descricao":"A mesma ideia para racionais E irracionais: a FRAÇÃO CONTÍNUA caracteriza todo ponto da reta. Racional ↔ CF finita; quadrático (metal/mineral σ_n) ↔ CF periódica — e σ_n=[n;n,n,…] exato (Teorema de Lagrange: periódica ⟺ quadrático); cúbico+ (plástico, π) ↔ aperiódica. Os convergentes p_k/q_k (racionais) aproximam qualquer ponto — a reta racional densa caracteriza a irracional. Os metais são o coração PURA-PERIÓDICO [n̄] da reta. Uma régua, todos os pontos.","epistemico":"fato","exemplos":[],"id":"caracterizacao_da_reta","memoria":"semantica","meta":{"dominio":"matematica","fonte":"caracterização da reta"},"origem":"wiki","palavras_chave":[],"sinonimos":[],"tipo":"conceito","titulo":"A Caracterização de Toda a Reta (frações contínuas)"}
\section{A Caracterização de Toda a Reta (frações contínuas)}
A mesma ideia para racionais E irracionais: a FRAÇÃO CONTÍNUA caracteriza todo ponto da reta. Racional \ensuremath{\leftrightarrow} CF finita; quadrático (metal/mineral \ensuremath{\sigma}\_n) \ensuremath{\leftrightarrow} CF periódica — e \ensuremath{\sigma}\_n=[n;n,n,…] exato (Teorema de Lagrange: periódica \ensuremath{\Longleftrightarrow} quadrático); cúbico+ (plástico, \ensuremath{\pi}) \ensuremath{\leftrightarrow} aperiódica. Os convergentes p\_k/q\_k (racionais) aproximam qualquer ponto — a reta racional densa caracteriza a irracional. Os metais são o coração PURA-PERIÓDICO [n\ensuremath{}] da reta. Uma régua, todos os pontos.
\prov{relações: usa transformada\_metalica; relaciona algebra\_reta\_mineral; relaciona numeros\_de\_pisot; relaciona reta\_de\_ouro; usa corpo\_primos\_mineral; relaciona cifra\_metais\_an; usa fracao\_continua\_phi}
\prov{proveniência: wiki \textperiodcentered{} caracterização da reta \textperiodcentered{} fato \textperiodcentered{} confiança 1.0 \textperiodcentered{} domínio matematica \textperiodcentered{} história 10 versões no jornal}

%CRISTAL {"arestas":[{"alvo":"programa_ns_bloqueios","confianca":1.0,"epistemico":"fato","origem":"paper","rel":"referencia"}],"confianca":1.0,"contraexemplos":[],"descricao":"Em cópias hierárquicas do 24-cell (raio 2⁻ᵏ, viscosidade 4ᵏ) a energia flui de escalas grandes para pequenas; a vorticidade efetiva |ω|²=Σ4ᵏE_k cresce ~10× via transferência espectral (BKM discreto).","epistemico":"fato","exemplos":[],"id":"cascata_multiescala","memoria":"semantica","meta":{"autor":"Gentil Lopes","descoberta":false,"dominio":"matematica","fonte":"paper_BF.pdf"},"origem":"paper","palavras_chave":[],"sinonimos":[],"tipo":"conceito","titulo":"Cascata Multi-escala"}
\section{Cascata Multi-escala}
Em cópias hierárquicas do 24-cell (raio 2\ensuremath{^{-k}}, viscosidade 4\ensuremath{^{k}}) a energia flui de escalas grandes para pequenas; a vorticidade efetiva |\ensuremath{\omega}|\ensuremath{^{2}}=\ensuremath{\Sigma}4\ensuremath{^{k}}E\_k cresce \~{}10$\times$ via transferência espectral (BKM discreto).
\prov{relações: referencia programa\_ns\_bloqueios}
\prov{proveniência: paper \textperiodcentered{} paper\_BF.pdf \textperiodcentered{} fato \textperiodcentered{} confiança 1.0 \textperiodcentered{} domínio matematica \textperiodcentered{} história 2 versões no jornal}

%CRISTAL {"arestas":[{"alvo":"constante_168","confianca":1.0,"epistemico":"fato","origem":"wiki","rel":"relaciona"},{"alvo":"tiffany","confianca":1.0,"epistemico":"fato","origem":"wiki","rel":"relaciona"},{"alvo":"tensor_cayley","confianca":1.0,"epistemico":"fato","origem":"wiki","rel":"relaciona"},{"alvo":"steiner_golay","confianca":1.0,"epistemico":"fato","origem":"wiki","rel":"relaciona"}],"confianca":1.0,"contraexemplos":[],"descricao":"Um edifício de arte-matemática que usa a função σ (soma dos divisores) como fio único: PSL(2,7)=168, E8=248, o plano de Fano, a rede de Leech e o moonshine emergem de iterar σ. Lema da obra: 'estrutura é sagrada, número é verdade; PSL é linguagem, E8 é a álgebra excepcional'.","epistemico":"inferencia","exemplos":[],"id":"catedral_sigma_aritmetica","memoria":"semantica","meta":{"autor":"Lord Index (Shawn R. Schiller)","dominio":"matematica","fonte":"A Catedral — Part VIII-A (jul 2026), ~/Imagens/catedral"},"origem":"wiki","palavras_chave":[],"sinonimos":[],"tipo":"conceito","titulo":"A Catedral (σ-aritmética)"}
\section{A Catedral (\ensuremath{\sigma}-aritmética)}
Um edifício de arte-matemática que usa a função \ensuremath{\sigma} (soma dos divisores) como fio único: PSL(2,7)=168, E8=248, o plano de Fano, a rede de Leech e o moonshine emergem de iterar \ensuremath{\sigma}. Lema da obra: 'estrutura é sagrada, número é verdade; PSL é linguagem, E8 é a álgebra excepcional'.
\prov{relações: relaciona constante\_168; relaciona tiffany; relaciona tensor\_cayley; relaciona steiner\_golay}
\prov{proveniência: wiki \textperiodcentered{} A Catedral — Part VIII-A (jul 2026), \~{}/Imagens/catedral \textperiodcentered{} inferencia \textperiodcentered{} confiança 1.0 \textperiodcentered{} domínio matematica \textperiodcentered{} história 5 versões no jornal}

%CRISTAL {"arestas":[{"alvo":"algebras_hipercomplexas","confianca":1.0,"epistemico":"fato","origem":"paper","rel":"usa"}],"confianca":1.0,"contraexemplos":[],"descricao":"Duplicação ×2 que gera a torre ℝ→ℂ→ℍ→𝕆→𝕊, cada passo dobrando a dimensão e perdendo uma propriedade (ordem, comutatividade, associatividade, divisão).","epistemico":"fato","exemplos":[],"id":"cayley_dickson","memoria":"semantica","meta":{"autor":"Gentil Lopes","descoberta":false,"dominio":"matematica","fonte":"paper_49_torre_hurwitz.pdf"},"origem":"paper","palavras_chave":[],"sinonimos":[],"tipo":"conceito","titulo":"Construção de Cayley-Dickson"}
\section{Construção de Cayley-Dickson}
Duplicação $\times$2 que gera a torre \ensuremath{\mathbb{R}}\ensuremath{\to}\ensuremath{\mathbb{C}}\ensuremath{\to}\ensuremath{\mathbb{H}}\ensuremath{\to}\ensuremath{\mathbb{O}}\ensuremath{\to}\ensuremath{\mathbb{S}}, cada passo dobrando a dimensão e perdendo uma propriedade (ordem, comutatividade, associatividade, divisão).
\prov{relações: usa algebras\_hipercomplexas}
\prov{proveniência: paper \textperiodcentered{} paper\_49\_torre\_hurwitz.pdf \textperiodcentered{} fato \textperiodcentered{} confiança 1.0 \textperiodcentered{} domínio matematica \textperiodcentered{} história 2 versões no jornal}

%CRISTAL {"arestas":[{"alvo":"experimento_da_agulha","confianca":1.0,"epistemico":"fato","origem":"wiki","rel":"parte_de"},{"alvo":"gato_ula","confianca":1.0,"epistemico":"fato","origem":"wiki","rel":"usa"},{"alvo":"gato","confianca":1.0,"epistemico":"fato","origem":"wiki","rel":"usa"},{"alvo":"numero_de_salem","confianca":1.0,"epistemico":"fato","origem":"wiki","rel":"relaciona"}],"confianca":1.0,"contraexemplos":[],"descricao":"A expansão em fração contínua de π É um produto de gatos [[a_k,1],[1,0]] — a ULA do computador mineral. O produto acumulado dá os convergentes p_k/q_k (3/1, 22/7, 355/113, …) → π, e o DETERMINANTE do produto é (−1)^k: |det|=1 CONSTANTE a cada passo. É a invariante estável mantida por PA (soma p_{k−2}) e PG (multiplica por a_k). Verificado.","epistemico":"fato","exemplos":[],"id":"cf_de_pi_e_produto_de_gatos","memoria":"semantica","meta":{"dominio":"reta mineral","fonte":"experimento da agulha"},"origem":"wiki","palavras_chave":[],"sinonimos":[],"tipo":"conceito","titulo":"A CF de π é um Produto de Gatos"}
\section{A CF de \ensuremath{\pi} é um Produto de Gatos}
A expansão em fração contínua de \ensuremath{\pi} É um produto de gatos [[a\_k,1],[1,0]] — a ULA do computador mineral. O produto acumulado dá os convergentes p\_k/q\_k (3/1, 22/7, 355/113, …) \ensuremath{\to} \ensuremath{\pi}, e o DETERMINANTE do produto é (\ensuremath{-}1)\^{}k: |det|=1 CONSTANTE a cada passo. É a invariante estável mantida por PA (soma p\_\{k\ensuremath{-}2\}) e PG (multiplica por a\_k). Verificado.
\prov{relações: parte\_de experimento\_da\_agulha; usa gato\_ula; usa gato; relaciona numero\_de\_salem}
\prov{proveniência: wiki \textperiodcentered{} experimento da agulha \textperiodcentered{} fato \textperiodcentered{} confiança 1.0 \textperiodcentered{} domínio reta mineral \textperiodcentered{} história 5 versões no jornal}

%CRISTAL {"arestas":[{"alvo":"experimento_da_agulha","confianca":1.0,"epistemico":"fato","origem":"wiki","rel":"parte_de"},{"alvo":"cf_de_pi_e_produto_de_gatos","confianca":1.0,"epistemico":"fato","origem":"wiki","rel":"usa"},{"alvo":"corpo_mineral","confianca":1.0,"epistemico":"fato","origem":"wiki","rel":"explica"},{"alvo":"pi_emerge_do_limite","confianca":1.0,"epistemico":"fato","origem":"wiki","rel":"relaciona"}],"confianca":1.0,"contraexemplos":[],"descricao":"Cada passo a_k da fração contínua de π nomeia um metal σ_{a_k}=(a_k+√(a_k²+4))/2. CF(π)=[3,7,15,1,292,1,1,1,2,…] → σ_3(bronze), σ_7, σ_15, σ_1(ouro), σ_292, σ_1, σ_1, σ_1, σ_2(prata)… A expansão de π ENUMERA uma sequência de metais: a CF de π gera a reta mineral. O ouro φ domina a cauda (os 1's).","epistemico":"fato","exemplos":[],"id":"cf_de_pi_gera_a_reta_mineral","memoria":"semantica","meta":{"dominio":"reta mineral","fonte":"experimento da agulha"},"origem":"wiki","palavras_chave":[],"sinonimos":[],"tipo":"conceito","titulo":"A CF de π Gera a Reta Mineral"}
\section{A CF de \ensuremath{\pi} Gera a Reta Mineral}
Cada passo a\_k da fração contínua de \ensuremath{\pi} nomeia um metal \ensuremath{\sigma}\_\{a\_k\}=(a\_k+\ensuremath{\sqrt{\,}}(a\_k\ensuremath{^{2}}+4))/2. CF(\ensuremath{\pi})=[3,7,15,1,292,1,1,1,2,…] \ensuremath{\to} \ensuremath{\sigma}\_3(bronze), \ensuremath{\sigma}\_7, \ensuremath{\sigma}\_15, \ensuremath{\sigma}\_1(ouro), \ensuremath{\sigma}\_292, \ensuremath{\sigma}\_1, \ensuremath{\sigma}\_1, \ensuremath{\sigma}\_1, \ensuremath{\sigma}\_2(prata)… A expansão de \ensuremath{\pi} ENUMERA uma sequência de metais: a CF de \ensuremath{\pi} gera a reta mineral. O ouro \ensuremath{\varphi} domina a cauda (os 1's).
\prov{relações: parte\_de experimento\_da\_agulha; usa cf\_de\_pi\_e\_produto\_de\_gatos; explica corpo\_mineral; relaciona pi\_emerge\_do\_limite}
\prov{proveniência: wiki \textperiodcentered{} experimento da agulha \textperiodcentered{} fato \textperiodcentered{} confiança 1.0 \textperiodcentered{} domínio reta mineral \textperiodcentered{} história 5 versões no jornal}

%CRISTAL {"arestas":[{"alvo":"costura_cobre_dobra10","confianca":1.0,"epistemico":"fato","origem":"wiki","rel":"especializa"},{"alvo":"bai_orbita_reta_metalica","confianca":1.0,"epistemico":"fato","origem":"wiki","rel":"usa"},{"alvo":"face_polida_primeira_frase","confianca":1.0,"epistemico":"fato","origem":"wiki","rel":"relaciona"},{"alvo":"tiffany","confianca":1.0,"epistemico":"fato","origem":"wiki","rel":"parte_de"}],"confianca":1.0,"contraexemplos":[],"descricao":"conversa/colapso.py + parabola_algebra.py: o chat aberto ('me conta algo', 'surpreenda-me') caía em SAUDAÇÃO porque a reta conversa (caos/cobre) não tinha camada-casa (alvo='') → órbita de repouso dialogos. Resolvido dando à reta do cobre a sua casa e a sua natureza: (1) conversa → (CAOS, conhecimento) — o chat aberto EXPLORA um CONCEITO (a Tiffany é agente de conhecimento, oferece saber, não saudação); (2) o octante do cobre (dobra 10) reescala conhecimento na conversa, favorecendo a exploração (mid-c², nem verbatim nem ruído); (3) ROTAÇÃO ÁUREA (φ³=σ₄, o número do cobre — a sequência de menor discrepância) entre os conceitos do topo, semeada por fala+turno, para VARIAR o conceito a cada abertura. Efeito: aberturas distintas → conceitos distintos do acervo (divisão honesta, multifractal, purgatório, Koch-Cantor-Fibonacci). É a reta do cobre finalmente sendo o que é: a exploração. Saudação/definição/cadeia intactos. [D] o caos explora, não repousa.","epistemico":"fato","exemplos":[],"id":"chat_aberto_cobre_explora","memoria":"semantica","meta":{"autor":"Lopes/broca-so","dominio":"algebra","fonte":"conversa/parabola_algebra.py (conversa→CAOS/conhecimento); colapso._peso_costura (octante cobre em conhecimento) + rotação áurea φ"},"origem":"wiki","palavras_chave":[],"sinonimos":[],"tipo":"conceito","titulo":"O chat aberto RESOLVIDO — a reta do cobre explora conceitos por rotação áurea (não saudação)"}
\section{O chat aberto RESOLVIDO — a reta do cobre explora conceitos por rotação áurea (não saudação)}
conversa/colapso.py + parabola\_algebra.py: o chat aberto ('me conta algo', 'surpreenda-me') caía em SAUDAÇÃO porque a reta conversa (caos/cobre) não tinha camada-casa (alvo='') \ensuremath{\to} órbita de repouso dialogos. Resolvido dando à reta do cobre a sua casa e a sua natureza: (1) conversa \ensuremath{\to} (CAOS, conhecimento) — o chat aberto EXPLORA um CONCEITO (a Tiffany é agente de conhecimento, oferece saber, não saudação); (2) o octante do cobre (dobra 10) reescala conhecimento na conversa, favorecendo a exploração (mid-c\ensuremath{^{2}}, nem verbatim nem ruído); (3) ROTAÇÃO ÁUREA (\ensuremath{\varphi}\ensuremath{^{3}}=\ensuremath{\sigma}\ensuremath{_{4}}, o número do cobre — a sequência de menor discrepância) entre os conceitos do topo, semeada por fala+turno, para VARIAR o conceito a cada abertura. Efeito: aberturas distintas \ensuremath{\to} conceitos distintos do acervo (divisão honesta, multifractal, purgatório, Koch-Cantor-Fibonacci). É a reta do cobre finalmente sendo o que é: a exploração. Saudação/definição/cadeia intactos. [D] o caos explora, não repousa.
\prov{relações: especializa costura\_cobre\_dobra10; usa bai\_orbita\_reta\_metalica; relaciona face\_polida\_primeira\_frase; parte\_de tiffany}
\prov{proveniência: wiki \textperiodcentered{} conversa/parabola\_algebra.py (conversa\ensuremath{\to}CAOS/conhecimento); colapso.\_peso\_costura (octante cobre em conhecimento) + rotação áurea \ensuremath{\varphi} \textperiodcentered{} fato \textperiodcentered{} confiança 1.0 \textperiodcentered{} domínio algebra \textperiodcentered{} história 62 versões no jornal}

%CRISTAL {"arestas":[{"alvo":"crivo_de_selberg","confianca":1.0,"epistemico":"fato","origem":"wiki","rel":"usa"},{"alvo":"o_chicote_a_limpeza_nao_um_operador","confianca":1.0,"epistemico":"fato","origem":"wiki","rel":"eh_um"},{"alvo":"chicote_e_o_tensor","confianca":1.0,"epistemico":"fato","origem":"wiki","rel":"relaciona"},{"alvo":"hipotese_de_riemann","confianca":1.0,"epistemico":"fato","origem":"wiki","rel":"relaciona"}],"confianca":1.0,"contraexemplos":[],"descricao":"A promoção. O 'chicote' da máquina — a decomposição espectral de Selberg, L²(Γ\\H)=L²_disc⊕L²_cont, que varre o espectro CONTÍNUO (séries de Eisenstein) pra isolar o DISCRETO (formas de Maass) — é a MESMA limpeza de Selberg que o crivo. Um objeto se decompõe; varre-se um fundo (compostos / espectro contínuo) por uma OTIMIZAÇÃO (a forma quadrática de Selberg / o produto interno de L², a ortogonalidade); retém-se o essencial (primos / espectro discreto). O dicionário 'peneira' faz a tradução completa (9 construtos, verif.): compostos↔contínuo, primos↔discreto, λ_d↔projetor, forma quadrática↔produto interno, X/G(z)↔lei de Weyl, resto↔determinante de espalhamento φ. Um só Selberg.","epistemico":"inferencia","exemplos":[],"id":"chicote_e_o_crivo_de_selberg","memoria":"semantica","meta":{"dominio":"matematica","fonte":"crivo de Selberg"},"origem":"wiki","palavras_chave":[],"sinonimos":[],"tipo":"conceito","titulo":"O Chicote de Selberg é o Crivo de Selberg"}
\section{O Chicote de Selberg é o Crivo de Selberg}
A promoção. O 'chicote' da máquina — a decomposição espectral de Selberg, L\ensuremath{^{2}}(\ensuremath{\Gamma}\textbackslash{}H)=L\ensuremath{^{2}}\_disc\ensuremath{\oplus}L\ensuremath{^{2}}\_cont, que varre o espectro CONTÍNUO (séries de Eisenstein) pra isolar o DISCRETO (formas de Maass) — é a MESMA limpeza de Selberg que o crivo. Um objeto se decompõe; varre-se um fundo (compostos / espectro contínuo) por uma OTIMIZAÇÃO (a forma quadrática de Selberg / o produto interno de L\ensuremath{^{2}}, a ortogonalidade); retém-se o essencial (primos / espectro discreto). O dicionário 'peneira' faz a tradução completa (9 construtos, verif.): compostos\ensuremath{\leftrightarrow}contínuo, primos\ensuremath{\leftrightarrow}discreto, \ensuremath{\lambda}\_d\ensuremath{\leftrightarrow}projetor, forma quadrática\ensuremath{\leftrightarrow}produto interno, X/G(z)\ensuremath{\leftrightarrow}lei de Weyl, resto\ensuremath{\leftrightarrow}determinante de espalhamento \ensuremath{\varphi}. Um só Selberg.
\prov{relações: usa crivo\_de\_selberg; eh\_um o\_chicote\_a\_limpeza\_nao\_um\_operador; relaciona chicote\_e\_o\_tensor; relaciona hipotese\_de\_riemann}
\prov{proveniência: wiki \textperiodcentered{} crivo de Selberg \textperiodcentered{} inferencia \textperiodcentered{} confiança 1.0 \textperiodcentered{} domínio matematica \textperiodcentered{} história 10 versões no jornal}

%CRISTAL {"arestas":[{"alvo":"otimizacao_fractal_preserving","confianca":1.0,"epistemico":"fato","origem":"wiki","rel":"usa"},{"alvo":"ramificacao_nao_orbita","confianca":1.0,"epistemico":"fato","origem":"wiki","rel":"usa"}],"confianca":1.0,"contraexemplos":[],"descricao":"A transformação incorreta típica colapsa o subshift a uma órbita única (ciclo simples): ρ=1 logo D_tr=0. O compilador deve REJEITAR qualquer Θ que leve D_tr>0 a D_tr=0 — destrói o fractal (a broca vira laço morto). Dual da distinção órbita-vs-subshift.","epistemico":"fato","exemplos":[],"id":"codigo_morto_fractal","memoria":"semantica","meta":{"autor":"broca-so","dominio":"fractal","fonte":"machine/fractal/broca_cc.py:ciclo_simples; compilador_fractal.tex"},"origem":"wiki","palavras_chave":[],"sinonimos":[],"tipo":"conceito","titulo":"Colapso da Ramificação (código morto fractal)"}
\section{Colapso da Ramificação (código morto fractal)}
A transformação incorreta típica colapsa o subshift a uma órbita única (ciclo simples): \ensuremath{\rho}=1 logo D\_tr=0. O compilador deve REJEITAR qualquer \ensuremath{\Theta} que leve D\_tr>0 a D\_tr=0 — destrói o fractal (a broca vira laço morto). Dual da distinção órbita-vs-subshift.
\prov{relações: usa otimizacao\_fractal\_preserving; usa ramificacao\_nao\_orbita}
\prov{proveniência: wiki \textperiodcentered{} machine/fractal/broca\_cc.py:ciclo\_simples; compilador\_fractal.tex \textperiodcentered{} fato \textperiodcentered{} confiança 1.0 \textperiodcentered{} domínio fractal \textperiodcentered{} história 3 versões no jornal}

%CRISTAL {"arestas":[{"alvo":"algebra_de_peirce","confianca":1.0,"epistemico":"fato","origem":"wiki","rel":"usa"},{"alvo":"definicao_frasagem_ampla","confianca":1.0,"epistemico":"fato","origem":"wiki","rel":"usa"},{"alvo":"memoria_contexto_fio_conversa","confianca":1.0,"epistemico":"fato","origem":"wiki","rel":"relaciona"},{"alvo":"tiffany","confianca":1.0,"epistemico":"fato","origem":"wiki","rel":"parte_de"}],"confianca":1.0,"contraexemplos":[],"descricao":"conversa/colapso.py (_resposta_comparacao): a Tiffany passou a responder COMPARAÇÕES. 'diferença entre X e Y' / 'compara X e Y' / 'X versus Y' → extrai os dois termos, acha o conceito de cada (_conceito_por_termo: o título que melhor casa por palavra de conteúdo, reusando a aboutness), e mostra os DOIS lado a lado (X ┃ Y) + a RELAÇÃO direta do grafo se houver (_relacao_entre: aresta cx→cy ou cy→cx, com a seta). É um desvio no colapso (antes do fio de contexto), como o chat aberto e a comparação. Usa a estrutura de arestas de Peirce para dizer o ELO entre os dois, não só justapor. Verificado: ouro┃prata, transformada┃gato, pisot┃salem. Completa os modos da Tiffany: define · continua (fio) · cadeia · explora · COMPARA. [F] roda.","epistemico":"fato","exemplos":[],"id":"comparacao_dois_conceitos","memoria":"semantica","meta":{"autor":"Lopes/broca-so","dominio":"algebra","fonte":"conversa/colapso.py:_resposta_comparacao, _conceito_por_termo, _relacao_entre; verificado: ouro/prata, transformada/gato"},"origem":"wiki","palavras_chave":[],"sinonimos":[],"tipo":"conceito","titulo":"Comparação entre X e Y: dois conceitos de uma vez, com a relação do grafo [F]"}
\section{Comparação entre X e Y: dois conceitos de uma vez, com a relação do grafo [F]}
conversa/colapso.py (\_resposta\_comparacao): a Tiffany passou a responder COMPARAÇÕES. 'diferença entre X e Y' / 'compara X e Y' / 'X versus Y' \ensuremath{\to} extrai os dois termos, acha o conceito de cada (\_conceito\_por\_termo: o título que melhor casa por palavra de conteúdo, reusando a aboutness), e mostra os DOIS lado a lado (X | Y) + a RELAÇÃO direta do grafo se houver (\_relacao\_entre: aresta cx\ensuremath{\to}cy ou cy\ensuremath{\to}cx, com a seta). É um desvio no colapso (antes do fio de contexto), como o chat aberto e a comparação. Usa a estrutura de arestas de Peirce para dizer o ELO entre os dois, não só justapor. Verificado: ouro|prata, transformada|gato, pisot|salem. Completa os modos da Tiffany: define \textperiodcentered  continua (fio) \textperiodcentered  cadeia \textperiodcentered  explora \textperiodcentered  COMPARA. [F] roda.
\prov{relações: usa algebra\_de\_peirce; usa definicao\_frasagem\_ampla; relaciona memoria\_contexto\_fio\_conversa; parte\_de tiffany}
\prov{proveniência: wiki \textperiodcentered{} conversa/colapso.py:\_resposta\_comparacao, \_conceito\_por\_termo, \_relacao\_entre; verificado: ouro/prata, transformada/gato \textperiodcentered{} fato \textperiodcentered{} confiança 1.0 \textperiodcentered{} domínio algebra \textperiodcentered{} história 30 versões no jornal}

%CRISTAL {"arestas":[{"alvo":"decomposicao_peirce","confianca":1.0,"epistemico":"fato","origem":"paper","rel":"usa"},{"alvo":"cifra","confianca":1.0,"epistemico":"fato","origem":"paper","rel":"parte_de"}],"confianca":1.0,"contraexemplos":[],"descricao":"Compartilhamento de segredo como threshold formal via os canais nilpotentes de Peirce Eᵢⱼ (divisores de zero): nenhuma chave individual entra, todas se cancelam no vazio; Σeᵢ=I reconstrói.","epistemico":"fato","exemplos":[],"id":"compartilhamento_nilpotente","memoria":"semantica","meta":{"autor":"Gentil Lopes","descoberta":false,"dominio":"matematica","fonte":"chaves.pdf, §3"},"origem":"paper","palavras_chave":[],"sinonimos":[],"tipo":"conceito","titulo":"Threshold Nilpotente (FROST)"}
\section{Threshold Nilpotente (FROST)}
Compartilhamento de segredo como threshold formal via os canais nilpotentes de Peirce E\ensuremath{_{ij}} (divisores de zero): nenhuma chave individual entra, todas se cancelam no vazio; \ensuremath{\Sigma}e\ensuremath{_{i}}=I reconstrói.
\prov{relações: usa decomposicao\_peirce; parte\_de cifra}
\prov{proveniência: paper \textperiodcentered{} chaves.pdf, §3 \textperiodcentered{} fato \textperiodcentered{} confiança 1.0 \textperiodcentered{} domínio matematica \textperiodcentered{} história 3 versões no jornal}

%CRISTAL {"arestas":[{"alvo":"dimensao_de_traco","confianca":1.0,"epistemico":"fato","origem":"wiki","rel":"usa"},{"alvo":"computacao_fractal","confianca":1.0,"epistemico":"fato","origem":"wiki","rel":"parte_de"}],"confianca":1.0,"contraexemplos":[],"descricao":"Compilar a Máquina de Torção é transformar programas preservando Dtr=log₂λₖ. A IR é o próprio IFS do gap (programa = lista de blocos de retorno 1r0 com sua recursão). brocacc.py é o middle- formalizado e verificado (o frontend DSL e o backend GEX são projeto).","epistemico":"fato","exemplos":[],"id":"compilador_fractal_broca_cc","memoria":"semantica","meta":{"autor":"broca-so","dominio":"fractal","fonte":"machine/fractal/broca_cc.py; compilador_fractal.tex"},"origem":"wiki","palavras_chave":[],"sinonimos":[],"tipo":"conceito","titulo":"O Compilador Fractal broca-cc"}
\section{O Compilador Fractal broca-cc}
Compilar a Máquina de Torção é transformar programas preservando Dtr=log\ensuremath{_{2}}\ensuremath{\lambda}\ensuremath{_{k}}. A IR é o próprio IFS do gap (programa = lista de blocos de retorno 1r0 com sua recursão). brocacc.py é o middle- formalizado e verificado (o frontend DSL e o backend GEX são projeto).
\prov{relações: usa dimensao\_de\_traco; parte\_de computacao\_fractal}
\prov{proveniência: wiki \textperiodcentered{} machine/fractal/broca\_cc.py; compilador\_fractal.tex \textperiodcentered{} fato \textperiodcentered{} confiança 1.0 \textperiodcentered{} domínio fractal \textperiodcentered{} história 4 versões no jornal}

%CRISTAL {"arestas":[{"alvo":"vetores","confianca":0.461,"epistemico":"fato","origem":"manual","rel":"referencia"},{"alvo":"transcricao","confianca":0.5,"epistemico":"fato","origem":"wiki","rel":"relaciona"},{"alvo":"log_legado_conversadadostelemetriajsonl","confianca":0.5,"epistemico":"fato","origem":"wiki","rel":"relaciona"},{"alvo":"parabola_pa","confianca":0.5,"epistemico":"fato","origem":"wiki","rel":"relaciona"},{"alvo":"resposta_a_pergunta_dificil","confianca":0.5,"epistemico":"fato","origem":"wiki","rel":"relaciona"},{"alvo":"topologia","confianca":0.5,"epistemico":"fato","origem":"wiki","rel":"relaciona"},{"alvo":"por_sessao_conversadadostelemetriasessaojsonl","confianca":0.5,"epistemico":"fato","origem":"wiki","rel":"relaciona"},{"alvo":"vida-morte","confianca":0.5,"epistemico":"fato","origem":"wiki","rel":"relaciona"},{"alvo":"pilha","confianca":0.5,"epistemico":"fato","origem":"wiki","rel":"relaciona"},{"alvo":"testes","confianca":0.5,"epistemico":"fato","origem":"wiki","rel":"relaciona"},{"alvo":"numa","confianca":0.5,"epistemico":"fato","origem":"wiki","rel":"relaciona"},{"alvo":"geometria_hub_broca","confianca":0.5,"epistemico":"fato","origem":"wiki","rel":"relaciona"},{"alvo":"ponteiro","confianca":0.5,"epistemico":"fato","origem":"wiki","rel":"relaciona"},{"alvo":"recursao","confianca":0.5,"epistemico":"fato","origem":"wiki","rel":"relaciona"},{"alvo":"tiffany_cabeca_broca","confianca":0.5,"epistemico":"fato","origem":"wiki","rel":"relaciona"},{"alvo":"pow_metal_pre_fecho","confianca":0.5,"epistemico":"fato","origem":"wiki","rel":"relaciona"},{"alvo":"koch_cantor_fib_prova","confianca":0.5,"epistemico":"fato","origem":"wiki","rel":"relaciona"},{"alvo":"ds_universo","confianca":0.5,"epistemico":"fato","origem":"wiki","rel":"relaciona"},{"alvo":"virtualizacao","confianca":0.5,"epistemico":"fato","origem":"wiki","rel":"relaciona"},{"alvo":"thread","confianca":0.5,"epistemico":"fato","origem":"wiki","rel":"relaciona"},{"alvo":"pow_metal_setor","confianca":0.5,"epistemico":"fato","origem":"wiki","rel":"relaciona"},{"alvo":"metais","confianca":0.5,"epistemico":"fato","origem":"wiki","rel":"relaciona"},{"alvo":"motivos","confianca":0.5,"epistemico":"fato","origem":"wiki","rel":"relaciona"},{"alvo":"pow_metal_funil","confianca":0.5,"epistemico":"fato","origem":"wiki","rel":"relaciona"},{"alvo":"hardware_sequencias_de_cauchy","confianca":0.5,"epistemico":"fato","origem":"wiki","rel":"relaciona"},{"alvo":"syscall","confianca":0.5,"epistemico":"fato","origem":"wiki","rel":"relaciona"},{"alvo":"parabola_destruir","confianca":0.5,"epistemico":"fato","origem":"wiki","rel":"relaciona"},{"alvo":"lobo_frontal","confianca":0.5,"epistemico":"fato","origem":"wiki","rel":"relaciona"},{"alvo":"pow_metal_ressonancia","confianca":0.5,"epistemico":"fato","origem":"wiki","rel":"relaciona"},{"alvo":"utilitarismo","confianca":0.5,"epistemico":"fato","origem":"wiki","rel":"relaciona"},{"alvo":"mecanica_quantica_hub","confianca":0.5,"epistemico":"fato","origem":"wiki","rel":"relaciona"},{"alvo":"matriz_de_confusao","confianca":0.5,"epistemico":"fato","origem":"wiki","rel":"relaciona"},{"alvo":"pow_geometria","confianca":0.5,"epistemico":"fato","origem":"wiki","rel":"relaciona"},{"alvo":"funil_cumulativo","confianca":0.5,"epistemico":"fato","origem":"wiki","rel":"relaciona"}],"confianca":1.0,"contraexemplos":[],"descricao":"a+bi com i²=−1; plano de Argand.","epistemico":"fato","exemplos":[],"id":"complexos","memoria":"semantica","meta":{"dominio":"matematica","nivel":4,"ordem":4},"origem":"curriculo","palavras_chave":[],"sinonimos":["ℂ"],"tipo":"conceito","titulo":"complexos"}
\section{complexos}
a+bi com i\ensuremath{^{2}}=\ensuremath{-}1; plano de Argand.
\prov{sinónimos: \ensuremath{\mathbb{C}}}
\prov{relações: referencia vetores; relaciona transcricao; relaciona log\_legado\_conversadadostelemetriajsonl; relaciona parabola\_pa; relaciona resposta\_a\_pergunta\_dificil; relaciona topologia; relaciona por\_sessao\_conversadadostelemetriasessaojsonl; relaciona vida-morte; relaciona pilha; relaciona testes; relaciona numa; relaciona geometria\_hub\_broca; relaciona ponteiro; relaciona recursao; relaciona tiffany\_cabeca\_broca; relaciona pow\_metal\_pre\_fecho; relaciona koch\_cantor\_fib\_prova; relaciona ds\_universo; relaciona virtualizacao; relaciona thread; relaciona pow\_metal\_setor; relaciona metais; relaciona motivos; relaciona pow\_metal\_funil; relaciona hardware\_sequencias\_de\_cauchy; relaciona syscall; relaciona parabola\_destruir; relaciona lobo\_frontal; relaciona pow\_metal\_ressonancia; relaciona utilitarismo; relaciona mecanica\_quantica\_hub; relaciona matriz\_de\_confusao; relaciona pow\_geometria; relaciona funil\_cumulativo}
\prov{proveniência: curriculo \textperiodcentered{} fato \textperiodcentered{} confiança 1.0 \textperiodcentered{} domínio matematica \textperiodcentered{} história 40 versões no jornal}

%CRISTAL {"arestas":[{"alvo":"metal_broca_so","confianca":1.0,"epistemico":"fato","origem":"wiki","rel":"parte_de"},{"alvo":"colapso_koch_tiffany","confianca":1.0,"epistemico":"fato","origem":"wiki","rel":"usa"},{"alvo":"maquina_algebrica","confianca":1.0,"epistemico":"fato","origem":"wiki","rel":"relaciona"}],"confianca":1.0,"contraexemplos":[],"descricao":"Um computador é fractal quando seu espaço de traços admissíveis é um subshift auto-similar X_k: cada bloco de retorno contém uma cópia escalada da máquina inteira (K_k=∪ι_r(K_k)). É a 'broca de brocas' tornada tese precisa — o fractal está no espaço de traços, não na álgebra nem na órbita.","epistemico":"fato","exemplos":[],"id":"computacao_fractal","memoria":"semantica","meta":{"autor":"broca-so","dominio":"fractal","fonte":"machine/computador_fractal.tex; machine/fractal/fractal.py"},"origem":"wiki","palavras_chave":[],"sinonimos":[],"tipo":"conceito","titulo":"Computação Fractal — a Máquina que se Contém"}
\section{Computação Fractal — a Máquina que se Contém}
Um computador é fractal quando seu espaço de traços admissíveis é um subshift auto-similar X\_k: cada bloco de retorno contém uma cópia escalada da máquina inteira (K\_k=\ensuremath{\cup}\ensuremath{\iota}\_r(K\_k)). É a 'broca de brocas' tornada tese precisa — o fractal está no espaço de traços, não na álgebra nem na órbita.
\prov{relações: parte\_de metal\_broca\_so; usa colapso\_koch\_tiffany; relaciona maquina\_algebrica}
\prov{proveniência: wiki \textperiodcentered{} machine/computador\_fractal.tex; machine/fractal/fractal.py \textperiodcentered{} fato \textperiodcentered{} confiança 1.0 \textperiodcentered{} domínio fractal \textperiodcentered{} história 4 versões no jornal}

%CRISTAL {"arestas":[{"alvo":"algebras_hipercomplexas","confianca":1.0,"epistemico":"fato","origem":"paper","rel":"parte_de"},{"alvo":"virtualizacao","confianca":0.5,"epistemico":"fato","origem":"wiki","rel":"relaciona"},{"alvo":"word","confianca":0.5,"epistemico":"fato","origem":"wiki","rel":"relaciona"}],"confianca":1.0,"contraexemplos":[],"descricao":"[z,w]ₛ = z★w − w★z = 2s(a×b), puramente imaginário; mede quanto a ordem dos operandos importa. Anula-se quando s=0 (comutativo).","epistemico":"fato","exemplos":[],"id":"comutador_algebrico","memoria":"semantica","meta":{"autor":"Gentil Lopes","descoberta":true,"dominio":"matematica","fonte":"paper_1_algebra_geometria.pdf, Prop.3.4"},"origem":"paper","palavras_chave":[],"sinonimos":[],"tipo":"conceito","titulo":"Comutador [z,w]ₛ"}
\section{Comutador [z,w]\ensuremath{_{s}}}
[z,w]\ensuremath{_{s}} = z$\star$w \ensuremath{-} w$\star$z = 2s(a$\times$b), puramente imaginário; mede quanto a ordem dos operandos importa. Anula-se quando s=0 (comutativo).
\prov{relações: parte\_de algebras\_hipercomplexas; relaciona virtualizacao; relaciona word}
\prov{proveniência: paper \textperiodcentered{} paper\_1\_algebra\_geometria.pdf, Prop.3.4 \textperiodcentered{} fato \textperiodcentered{} confiança 1.0 \textperiodcentered{} domínio matematica \textperiodcentered{} história 4 versões no jornal}

%CRISTAL {"arestas":[{"alvo":"terceiro_estado","confianca":1.0,"epistemico":"fato","origem":"paper","rel":"usa"},{"alvo":"word","confianca":0.5,"epistemico":"fato","origem":"wiki","rel":"relaciona"}],"confianca":1.0,"contraexemplos":[],"descricao":"A hipersuperfície nula N(z)=0 (1−s²=0) que separa a ilha (tipo-tempo, N>0) do oceano (tipo-espaço, N<0): a 'reta crítica' onde a norma troca de sinal.","epistemico":"fato","exemplos":[],"id":"cone_luz_critico","memoria":"semantica","meta":{"autor":"Gentil Lopes","descoberta":false,"dominio":"matematica","fonte":"geometria.pdf, §9"},"origem":"paper","palavras_chave":[],"sinonimos":[],"tipo":"conceito","titulo":"Cone de Luz Algébrico"}
\section{Cone de Luz Algébrico}
A hipersuperfície nula N(z)=0 (1\ensuremath{-}s\ensuremath{^{2}}=0) que separa a ilha (tipo-tempo, N>0) do oceano (tipo-espaço, N<0): a 'reta crítica' onde a norma troca de sinal.
\prov{relações: usa terceiro\_estado; relaciona word}
\prov{proveniência: paper \textperiodcentered{} geometria.pdf, §9 \textperiodcentered{} fato \textperiodcentered{} confiança 1.0 \textperiodcentered{} domínio matematica \textperiodcentered{} história 3 versões no jornal}

%CRISTAL {"arestas":[{"alvo":"psl_2_7","confianca":1.0,"epistemico":"fato","origem":"paper","rel":"usa"},{"alvo":"klein_quartic","confianca":1.0,"epistemico":"fato","origem":"paper","rel":"explica"}],"confianca":1.0,"contraexemplos":[],"descricao":"Conjectura: 24 anéis de vórtice em ℝ⁷ na órbita de PSL(2,7) (=168/7) emergem em fluxos de NS-octônica, reproduzindo a combinatória da Klein quartic (24 heptágonos, 84 adjacências).","epistemico":"hipotese","exemplos":[],"id":"config_24_rings","memoria":"semantica","meta":{"autor":"Gentil Lopes","descoberta":false,"dominio":"matematica","fonte":"paper_AY.pdf, Conj.7.3"},"origem":"paper","palavras_chave":[],"sinonimos":[],"tipo":"conceito","titulo":"Configuração de 24 Vórtices (órbita PSL)"}
\section{Configuração de 24 Vórtices (órbita PSL)}
Conjectura: 24 anéis de vórtice em \ensuremath{\mathbb{R}}\ensuremath{^{7}} na órbita de PSL(2,7) (=168/7) emergem em fluxos de NS-octônica, reproduzindo a combinatória da Klein quartic (24 heptágonos, 84 adjacências).
\prov{relações: usa psl\_2\_7; explica klein\_quartic}
\prov{proveniência: paper \textperiodcentered{} paper\_AY.pdf, Conj.7.3 \textperiodcentered{} hipotese \textperiodcentered{} confiança 1.0 \textperiodcentered{} domínio matematica \textperiodcentered{} história 3 versões no jornal}

%CRISTAL {"arestas":[{"alvo":"anel_transcendental_rk","confianca":1.0,"epistemico":"fato","origem":"paper","rel":"explica"}],"confianca":1.0,"contraexemplos":[],"descricao":"Os transcendentais π², ζ(3), ζ(5), … são ℚ-algebricamente independentes; teorema para peso k≤12 (Brown 2012, motivos mistos de Tate), conjectura mainstream para k>12. Fundamento da não-representabilidade simbólica.","epistemico":"fato","exemplos":[],"id":"conjectura_brown","memoria":"semantica","meta":{"autor":"Gentil Lopes","descoberta":false,"dominio":"matematica","fonte":"paper_PQ_Q.pdf, Teo.2.2"},"origem":"paper","palavras_chave":[],"sinonimos":[],"tipo":"conceito","titulo":"Conjectura de Brown (independência dos ζ)"}
\section{Conjectura de Brown (independência dos \ensuremath{\zeta})}
Os transcendentais \ensuremath{\pi}\ensuremath{^{2}}, \ensuremath{\zeta}(3), \ensuremath{\zeta}(5), … são \ensuremath{\mathbb{Q}}-algebricamente independentes; teorema para peso k\ensuremath{\le}12 (Brown 2012, motivos mistos de Tate), conjectura mainstream para k>12. Fundamento da não-representabilidade simbólica.
\prov{relações: explica anel\_transcendental\_rk}
\prov{proveniência: paper \textperiodcentered{} paper\_PQ\_Q.pdf, Teo.2.2 \textperiodcentered{} fato \textperiodcentered{} confiança 1.0 \textperiodcentered{} domínio matematica \textperiodcentered{} história 3 versões no jornal}

%CRISTAL {"arestas":[{"alvo":"parabola_algebra_conversacional","confianca":1.0,"epistemico":"fato","origem":"wiki","rel":"explica"},{"alvo":"numeros_pisot","confianca":1.0,"epistemico":"fato","origem":"wiki","rel":"usa"},{"alvo":"corpo_dual","confianca":1.0,"epistemico":"fato","origem":"wiki","rel":"relaciona"},{"alvo":"conservacao_parabola_costura","confianca":1.0,"epistemico":"fato","origem":"wiki","rel":"explica"},{"alvo":"tiffany","confianca":1.0,"epistemico":"fato","origem":"wiki","rel":"parte_de"}],"confianca":1.0,"contraexemplos":[],"descricao":"A ponte que faz tudo fechar no mesmo objeto: a reta metálica x²−n·x−1 tem dois conjugados de Galois — σ_n (|σ|>1, EXPANDE) e −1/σ_n (|·|<1, CONTRAI). O eixo que CONTRAI é o lado associativo/cristal (c², a projeção estável, a janela de Rauzy/quasicristal); o que EXPANDE é o dissipativo (a, o calor). O produto σ·(−1/σ)=−1 é a NORMA = o det da Q-matriz A_n = o invariante conservado (a taxa de conversão entre metais). Assim a PARÁBOLA a+c²=1 é a decomposição de Galois em suas duas projeções somando o todo, e a condição PISOT (|conj|<1 = CRISTAL) é ao mesmo tempo o que faz a numeração fechar (Property F) e o que dá o quasicristal (cut-and-project). Um par conjugado gera: a numeração, a dobra, a conservação e a parábola. [I] a raiz comum de tudo.","epistemico":"inferencia","exemplos":[],"id":"conjugados_galois_parabola","memoria":"semantica","meta":{"autor":"Lopes/broca-so","dominio":"algebra","fonte":"teoria de Galois (conjugados de x²−nx−1); norma −1; Pisot⇒cut-and-project; [[conservacao_parabola_costura]]"},"origem":"wiki","palavras_chave":[],"sinonimos":[],"tipo":"conceito","titulo":"Os dois conjugados de Galois SÃO os dois eixos da parábola (σ expande, −1/σ contrai)"}
\section{Os dois conjugados de Galois SÃO os dois eixos da parábola (\ensuremath{\sigma} expande, \ensuremath{-}1/\ensuremath{\sigma} contrai)}
A ponte que faz tudo fechar no mesmo objeto: a reta metálica x\ensuremath{^{2}}\ensuremath{-}n\textperiodcentered x\ensuremath{-}1 tem dois conjugados de Galois — \ensuremath{\sigma}\_n (|\ensuremath{\sigma}|>1, EXPANDE) e \ensuremath{-}1/\ensuremath{\sigma}\_n (|\textperiodcentered |<1, CONTRAI). O eixo que CONTRAI é o lado associativo/cristal (c\ensuremath{^{2}}, a projeção estável, a janela de Rauzy/quasicristal); o que EXPANDE é o dissipativo (a, o calor). O produto \ensuremath{\sigma}\textperiodcentered (\ensuremath{-}1/\ensuremath{\sigma})=\ensuremath{-}1 é a NORMA = o det da Q-matriz A\_n = o invariante conservado (a taxa de conversão entre metais). Assim a PARÁBOLA a+c\ensuremath{^{2}}=1 é a decomposição de Galois em suas duas projeções somando o todo, e a condição PISOT (|conj|<1 = CRISTAL) é ao mesmo tempo o que faz a numeração fechar (Property F) e o que dá o quasicristal (cut-and-project). Um par conjugado gera: a numeração, a dobra, a conservação e a parábola. [I] a raiz comum de tudo.
\prov{relações: explica parabola\_algebra\_conversacional; usa numeros\_pisot; relaciona corpo\_dual; explica conservacao\_parabola\_costura; parte\_de tiffany}
\prov{proveniência: wiki \textperiodcentered{} teoria de Galois (conjugados de x\ensuremath{^{2}}\ensuremath{-}nx\ensuremath{-}1); norma \ensuremath{-}1; Pisot\ensuremath{\Rightarrow}cut-and-project; [[conservacao\_parabola\_costura]] \textperiodcentered{} inferencia \textperiodcentered{} confiança 1.0 \textperiodcentered{} domínio algebra \textperiodcentered{} história 38 versões no jornal}

%CRISTAL {"arestas":[{"alvo":"conjugados_galois_parabola","confianca":1.0,"epistemico":"fato","origem":"wiki","rel":"usa"},{"alvo":"bai_gate_dtc_orbita","confianca":1.0,"epistemico":"fato","origem":"wiki","rel":"usa"},{"alvo":"costura_prata_dobra8","confianca":1.0,"epistemico":"fato","origem":"wiki","rel":"usa"},{"alvo":"parabola_algebra_conversacional","confianca":1.0,"epistemico":"fato","origem":"wiki","rel":"usa"},{"alvo":"corpo_dual","confianca":1.0,"epistemico":"fato","origem":"wiki","rel":"relaciona"},{"alvo":"tiffany","confianca":1.0,"epistemico":"fato","origem":"wiki","rel":"parte_de"}],"confianca":1.0,"contraexemplos":[],"descricao":"conversa/colapso.py (parabola_colapso): a costura DTC (ouro radial × prata angular) escala o lado associativo (o ov), mas isso NÃO pode destruir energia — a lei de conservação é a própria parábola a+c²=1, e ela tem que continuar valendo através da costura. Formalização: o bit associativo bruto é c²₀=ov/256; o alinhamento α (peso radial × octante da prata) diz quanto FICA cristal: c²_ef=α·c²₀. O que sai, (1−α)·c²₀, não some — vira ASSOCIADOR (a entropia de cruzar o setor, DTC E4), somado ao desacordo bruto a₀=1−c²₀: a_ef = a₀+(1−α)·c²₀ = 1−c²_ef. Logo c²_ef + a_ef = 1 SEMPRE — a parábola é conservada. E APARECE: o colapso retorna c2/a/parabola do vencedor e o observar mostra 'c²+a=1' no passo. Verificado p/ todos (ov,α). O gate move peso do cristal (associativo/frio) pro calor (associador/dissipativo), na reta da parábola — nunca fora dela. [D] a conservação como lei da costura, visível.","epistemico":"fato","exemplos":[],"id":"conservacao_parabola_costura","memoria":"semantica","meta":{"autor":"Lopes/broca-so","dominio":"algebra","fonte":"conversa/colapso.py:parabola_colapso, _peso_costura, colapso (c2/a/parabola); observar.py:cmd_passo"},"origem":"wiki","palavras_chave":[],"sinonimos":[],"tipo":"conceito","titulo":"A costura CONSERVA a parábola a+c²=1 — o gate redistribui c²→a, não destrói"}
\section{A costura CONSERVA a parábola a+c\ensuremath{^{2}}=1 — o gate redistribui c\ensuremath{^{2}}\ensuremath{\to}a, não destrói}
conversa/colapso.py (parabola\_colapso): a costura DTC (ouro radial $\times$ prata angular) escala o lado associativo (o ov), mas isso NÃO pode destruir energia — a lei de conservação é a própria parábola a+c\ensuremath{^{2}}=1, e ela tem que continuar valendo através da costura. Formalização: o bit associativo bruto é c\ensuremath{^{2}}\ensuremath{_{0}}=ov/256; o alinhamento \ensuremath{\alpha} (peso radial $\times$ octante da prata) diz quanto FICA cristal: c\ensuremath{^{2}}\_ef=\ensuremath{\alpha}\textperiodcentered c\ensuremath{^{2}}\ensuremath{_{0}}. O que sai, (1\ensuremath{-}\ensuremath{\alpha})\textperiodcentered c\ensuremath{^{2}}\ensuremath{_{0}}, não some — vira ASSOCIADOR (a entropia de cruzar o setor, DTC E4), somado ao desacordo bruto a\ensuremath{_{0}}=1\ensuremath{-}c\ensuremath{^{2}}\ensuremath{_{0}}: a\_ef = a\ensuremath{_{0}}+(1\ensuremath{-}\ensuremath{\alpha})\textperiodcentered c\ensuremath{^{2}}\ensuremath{_{0}} = 1\ensuremath{-}c\ensuremath{^{2}}\_ef. Logo c\ensuremath{^{2}}\_ef + a\_ef = 1 SEMPRE — a parábola é conservada. E APARECE: o colapso retorna c2/a/parabola do vencedor e o observar mostra 'c\ensuremath{^{2}}+a=1' no passo. Verificado p/ todos (ov,\ensuremath{\alpha}). O gate move peso do cristal (associativo/frio) pro calor (associador/dissipativo), na reta da parábola — nunca fora dela. [D] a conservação como lei da costura, visível.
\prov{relações: usa conjugados\_galois\_parabola; usa bai\_gate\_dtc\_orbita; usa costura\_prata\_dobra8; usa parabola\_algebra\_conversacional; relaciona corpo\_dual; parte\_de tiffany}
\prov{proveniência: wiki \textperiodcentered{} conversa/colapso.py:parabola\_colapso, \_peso\_costura, colapso (c2/a/parabola); observar.py:cmd\_passo \textperiodcentered{} fato \textperiodcentered{} confiança 1.0 \textperiodcentered{} domínio algebra \textperiodcentered{} história 113 versões no jornal}

%CRISTAL {"arestas":[{"alvo":"tensor_cayley","confianca":1.0,"epistemico":"fato","origem":"paper","rel":"explica"},{"alvo":"psl_2_7","confianca":1.0,"epistemico":"fato","origem":"paper","rel":"usa"},{"alvo":"matematica","confianca":1.0,"epistemico":"fato","origem":"paper","rel":"parte_de"}],"confianca":1.0,"contraexemplos":[],"descricao":"168 aparece em três lugares como uma estrutura só: ‖φ‖²=168 (4-forma de Cayley em 7 trios de Fano), |PSL(2,7)|=168 (2º grupo simples não-abeliano), |Aut(Klein quartic)|=168 (limite de Hurwitz 84(g−1) para g=3). O plano de Fano é a origem comum.","epistemico":"fato","exemplos":[],"id":"constante_168","memoria":"semantica","meta":{"autor":"Gentil Lopes","descoberta":true,"dominio":"matematica","fonte":"paper_AU.pdf"},"origem":"paper","palavras_chave":[],"sinonimos":[],"tipo":"conceito","titulo":"A Constante 168"}
\section{A Constante 168}
168 aparece em três lugares como uma estrutura só: \ensuremath{\|}\ensuremath{\varphi}\ensuremath{\|}\ensuremath{^{2}}=168 (4-forma de Cayley em 7 trios de Fano), |PSL(2,7)|=168 (2º grupo simples não-abeliano), |Aut(Klein quartic)|=168 (limite de Hurwitz 84(g\ensuremath{-}1) para g=3). O plano de Fano é a origem comum.
\prov{relações: explica tensor\_cayley; usa psl\_2\_7; parte\_de matematica}
\prov{proveniência: paper \textperiodcentered{} paper\_AU.pdf \textperiodcentered{} fato \textperiodcentered{} confiança 1.0 \textperiodcentered{} domínio matematica \textperiodcentered{} história 4 versões no jornal}

%CRISTAL {"arestas":[{"alvo":"paley_169_e8","confianca":1.0,"epistemico":"fato","origem":"wiki","rel":"relaciona"}],"confianca":1.0,"contraexemplos":[],"descricao":"Números de teoria extremal de grafos e Ramsey caem na teia σ: σ(71)=72 (Mubayi), σ(133)=160 (Verstraete/von Neumann), σ(81)=121 (Ramsey), e o Paley de ordem 109 (primo). A aritmética das arestas encontrando as constantes excepcionais.","epistemico":"fato","exemplos":[],"id":"constantes_teoria_extremal","memoria":"semantica","meta":{"autor":"Lord Index (Shawn R. Schiller)","dominio":"matematica","fonte":"A Catedral — Part VIII-A (jul 2026), ~/Imagens/catedral"},"origem":"wiki","palavras_chave":[],"sinonimos":[],"tipo":"conceito","titulo":"Constantes da Teoria Extremal de Grafos na teia σ"}
\section{Constantes da Teoria Extremal de Grafos na teia \ensuremath{\sigma}}
Números de teoria extremal de grafos e Ramsey caem na teia \ensuremath{\sigma}: \ensuremath{\sigma}(71)=72 (Mubayi), \ensuremath{\sigma}(133)=160 (Verstraete/von Neumann), \ensuremath{\sigma}(81)=121 (Ramsey), e o Paley de ordem 109 (primo). A aritmética das arestas encontrando as constantes excepcionais.
\prov{relações: relaciona paley\_169\_e8}
\prov{proveniência: wiki \textperiodcentered{} A Catedral — Part VIII-A (jul 2026), \~{}/Imagens/catedral \textperiodcentered{} fato \textperiodcentered{} confiança 1.0 \textperiodcentered{} domínio matematica \textperiodcentered{} história 2 versões no jornal}

%CRISTAL {"arestas":[{"alvo":"numeros_azuis_vermelhos","confianca":1.0,"epistemico":"fato","origem":"paper","rel":"relaciona"},{"alvo":"inteiros","confianca":1.0,"epistemico":"fato","origem":"paper","rel":"explica"}],"confianca":1.0,"contraexemplos":[],"descricao":"ℤ=(ℕ×ℕ)/~ com (a,b)~(c,d)⟺a+d=b+c: cada inteiro é uma classe (uma reta paralela em ℕ²). O número é a classe, não a substância.","epistemico":"fato","exemplos":[],"id":"construcao_z_equivalencia","memoria":"semantica","meta":{"autor":"Gentil Lopes","descoberta":false,"dominio":"matematica","fonte":"corpus: OS_NUMEROS_INTEIROS_AZUIS, Teo.5"},"origem":"paper","palavras_chave":[],"sinonimos":[],"tipo":"conceito","titulo":"Construção de ℤ por Equivalência"}
\section{Construção de \ensuremath{\mathbb{Z}} por Equivalência}
\ensuremath{\mathbb{Z}}=(\ensuremath{\mathbb{N}}$\times$\ensuremath{\mathbb{N}})/\~{} com (a,b)\~{}(c,d)\ensuremath{\Longleftrightarrow}a+d=b+c: cada inteiro é uma classe (uma reta paralela em \ensuremath{\mathbb{N}}\ensuremath{^{2}}). O número é a classe, não a substância.
\prov{relações: relaciona numeros\_azuis\_vermelhos; explica inteiros}
\prov{proveniência: paper \textperiodcentered{} corpus: OS\_NUMEROS\_INTEIROS\_AZUIS, Teo.5 \textperiodcentered{} fato \textperiodcentered{} confiança 1.0 \textperiodcentered{} domínio matematica \textperiodcentered{} história 3 versões no jornal}

%CRISTAL {"arestas":[{"alvo":"hopfield","confianca":1.0,"epistemico":"fato","origem":"paper","rel":"referencia"}],"confianca":1.0,"contraexemplos":[],"descricao":"Estágio 1 (rigoroso): n em binário como anéis de vórtice orientados ±e_z; o incremento INC é uma sequência de reconexões locais com carry — prova de conceito construída em NS-3D.","epistemico":"fato","exemplos":[],"id":"contador_binario_ns","memoria":"semantica","meta":{"autor":"Gentil Lopes","descoberta":false,"dominio":"matematica","fonte":"paper_AP.pdf, Teo.8"},"origem":"paper","palavras_chave":[],"sinonimos":[],"tipo":"conceito","titulo":"Contador Binário em NS"}
\section{Contador Binário em NS}
Estágio 1 (rigoroso): n em binário como anéis de vórtice orientados $\pm$e\_z; o incremento INC é uma sequência de reconexões locais com carry — prova de conceito construída em NS-3D.
\prov{relações: referencia hopfield}
\prov{proveniência: paper \textperiodcentered{} paper\_AP.pdf, Teo.8 \textperiodcentered{} fato \textperiodcentered{} confiança 1.0 \textperiodcentered{} domínio matematica \textperiodcentered{} história 2 versões no jornal}

%CRISTAL {"arestas":[{"alvo":"plano_continuo_transicoes","confianca":1.0,"epistemico":"fato","origem":"wiki","rel":"parte_de"},{"alvo":"absorcao_iterativa_poca","confianca":1.0,"epistemico":"fato","origem":"wiki","rel":"usa"},{"alvo":"medida_continuidade_transicao","confianca":1.0,"epistemico":"fato","origem":"wiki","rel":"usa"}],"confianca":1.0,"contraexemplos":[],"descricao":"Ingerindo obras em levas (corpus_aarao, u11, gentil→Lopes, 346 papers .tex, Machado, u12, multiverso, machine, Newton da Costa/Euler/drive, bíblia PT, Platão, Schopenhauer, Clausewitz, Vedanta…), o contínuo passou a cobrir ~98,5% das arestas do grafo (≈5449/5530), com >20 mil frases reais preenchendo os segmentos entre os nós. O fluido semeado por densidade leva cada obra às arestas ainda vazias; a parábola satura as cheias; o excedente vai pro purgatório (>80 mil). ACHADO: a ressonância é CROSS-LÍNGUA — obras em inglês (Platão, Schopenhauer) também acham casa nas arestas vazias (o fluido as posiciona, mesmo com fit semântico mais frouxo). As ~81 arestas restantes não são tocadas por nenhuma obra ingerida — pedem conteúdo de tópico específico ou o crescimento do grafo. [D] roda; contínuo consultável (tecido transicoes).","epistemico":"fato","exemplos":[],"id":"continuo_quase_saturado","memoria":"semantica","meta":{"autor":"Lopes/broca-so","dominio":"continuo","fonte":"conversa/dados/transicoes.tsv (~20k linhas, 98,5% cobertura); scratchpad/ingest.py (levas)"},"origem":"wiki","palavras_chave":[],"sinonimos":[],"tipo":"conceito","titulo":"O contínuo quase saturado: ~13 obras absorvidas cobrem 98,5% das arestas"}
\section{O contínuo quase saturado: \~{}13 obras absorvidas cobrem 98,5\% das arestas}
Ingerindo obras em levas (corpus\_aarao, u11, gentil\ensuremath{\to}Lopes, 346 papers .tex, Machado, u12, multiverso, machine, Newton da Costa/Euler/drive, bíblia PT, Platão, Schopenhauer, Clausewitz, Vedanta…), o contínuo passou a cobrir \~{}98,5\% das arestas do grafo (\ensuremath{\approx}5449/5530), com >20 mil frases reais preenchendo os segmentos entre os nós. O fluido semeado por densidade leva cada obra às arestas ainda vazias; a parábola satura as cheias; o excedente vai pro purgatório (>80 mil). ACHADO: a ressonância é CROSS-LÍNGUA — obras em inglês (Platão, Schopenhauer) também acham casa nas arestas vazias (o fluido as posiciona, mesmo com fit semântico mais frouxo). As \~{}81 arestas restantes não são tocadas por nenhuma obra ingerida — pedem conteúdo de tópico específico ou o crescimento do grafo. [D] roda; contínuo consultável (tecido transicoes).
\prov{relações: parte\_de plano\_continuo\_transicoes; usa absorcao\_iterativa\_poca; usa medida\_continuidade\_transicao}
\prov{proveniência: wiki \textperiodcentered{} conversa/dados/transicoes.tsv (\~{}20k linhas, 98,5\% cobertura); scratchpad/ingest.py (levas) \textperiodcentered{} fato \textperiodcentered{} confiança 1.0 \textperiodcentered{} domínio continuo \textperiodcentered{} história 16 versões no jornal}

%CRISTAL {"arestas":[{"alvo":"symbolic_beta","confianca":1.0,"epistemico":"fato","origem":"paper","rel":"contradiz"}],"confianca":1.0,"contraexemplos":[],"descricao":"Casos onde a função beta é representável em forma fechada, violando a hardness geral: Maxwell puro (β=0 exato), N=4 SYM (cancelamento SUSY), Sine-Gordon 2D. A representabilidade ocorre numa família discreta de modelos, não em todo o espaço.","epistemico":"fato","exemplos":[],"id":"contraexemplos_hardness","memoria":"semantica","meta":{"autor":"Gentil Lopes","descoberta":false,"dominio":"matematica","fonte":"paper_PQ_J.pdf, §2-4"},"origem":"paper","palavras_chave":[],"sinonimos":[],"tipo":"conceito","titulo":"Contraexemplos à Hardness"}
\section{Contraexemplos à Hardness}
Casos onde a função beta é representável em forma fechada, violando a hardness geral: Maxwell puro (\ensuremath{\beta}=0 exato), N=4 SYM (cancelamento SUSY), Sine-Gordon 2D. A representabilidade ocorre numa família discreta de modelos, não em todo o espaço.
\prov{relações: contradiz symbolic\_beta}
\prov{proveniência: paper \textperiodcentered{} paper\_PQ\_J.pdf, §2-4 \textperiodcentered{} fato \textperiodcentered{} confiança 1.0 \textperiodcentered{} domínio matematica \textperiodcentered{} história 2 versões no jornal}

%CRISTAL {"arestas":[{"alvo":"goldchain","confianca":1.0,"epistemico":"fato","origem":"paper","rel":"parte_de"}],"confianca":1.0,"contraexemplos":[],"descricao":"Um contrato é a tripla P (chaves das partes), φ (predicado booleano de fechamento) e λ (liquidação). Interface universal sobre BTC/ETH — HTLC, multisig e transferência são casos.","epistemico":"fato","exemplos":[],"id":"contrato_objeto","memoria":"semantica","meta":{"autor":"Gentil Lopes","descoberta":true,"dominio":"matematica","fonte":"kernel/algebra_dos_contratos.pdf, Def.2"},"origem":"paper","palavras_chave":[],"sinonimos":[],"tipo":"conceito","titulo":"Contrato como Objeto C=(P,φ,λ)"}
\section{Contrato como Objeto C=(P,\ensuremath{\varphi},\ensuremath{\lambda})}
Um contrato é a tripla P (chaves das partes), \ensuremath{\varphi} (predicado booleano de fechamento) e \ensuremath{\lambda} (liquidação). Interface universal sobre BTC/ETH — HTLC, multisig e transferência são casos.
\prov{relações: parte\_de goldchain}
\prov{proveniência: paper \textperiodcentered{} kernel/algebra\_dos\_contratos.pdf, Def.2 \textperiodcentered{} fato \textperiodcentered{} confiança 1.0 \textperiodcentered{} domínio matematica \textperiodcentered{} história 2 versões no jornal}

%CRISTAL {"arestas":[{"alvo":"dq_espectro_ilha","confianca":1.0,"epistemico":"fato","origem":"wiki","rel":"usa"},{"alvo":"mecanica_quantica","confianca":1.0,"epistemico":"fato","origem":"wiki","rel":"relaciona"}],"confianca":1.0,"contraexemplos":[],"descricao":"No plano (t,s), a convenção r=+1 (quaternion) produz fração física=0.0 e todos os D_q=null, enquanto r=−1 (hermitian) dá 0.69% — a ilha físico-admissível só existe na convenção hermitian. A física seleciona o sinal.","epistemico":"fato","exemplos":[],"id":"convencao_hermitian_vs_quaternion","memoria":"semantica","meta":{"autor":"broca-so","dominio":"fractal","fonte":"hiper/benchmarks/fase35_multifractal_resultados.json"},"origem":"wiki","palavras_chave":[],"sinonimos":[],"tipo":"conceito","titulo":"A Ilha só Existe na Convenção Hermitian (r=−1)"}
\section{A Ilha só Existe na Convenção Hermitian (r=\ensuremath{-}1)}
No plano (t,s), a convenção r=+1 (quaternion) produz fração física=0.0 e todos os D\_q=null, enquanto r=\ensuremath{-}1 (hermitian) dá 0.69\% — a ilha físico-admissível só existe na convenção hermitian. A física seleciona o sinal.
\prov{relações: usa dq\_espectro\_ilha; relaciona mecanica\_quantica}
\prov{proveniência: wiki \textperiodcentered{} hiper/benchmarks/fase35\_multifractal\_resultados.json \textperiodcentered{} fato \textperiodcentered{} confiança 1.0 \textperiodcentered{} domínio fractal \textperiodcentered{} história 3 versões no jornal}

%CRISTAL {"arestas":[{"alvo":"cascata_multiescala","confianca":1.0,"epistemico":"fato","origem":"paper","rel":"usa"}],"confianca":1.0,"contraexemplos":[],"descricao":"O truncamento hierárquico converge exponencialmente na projeção de escala fixa (‖PM(γN−γN⁺Δ)‖≤e−c(Ns−M)), mas — honestamente — só em L² projetado, não em H¹/L^∞ (o limite inviscido diverge em normas fortes).","epistemico":"fato","exemplos":[],"id":"convergencia_galerkin","memoria":"semantica","meta":{"autor":"Gentil Lopes","descoberta":false,"dominio":"matematica","fonte":"paper_BK.pdf, Teo.3"},"origem":"paper","palavras_chave":[],"sinonimos":[],"tipo":"conceito","titulo":"Convergência de Galerkin"}
\section{Convergência de Galerkin}
O truncamento hierárquico converge exponencialmente na projeção de escala fixa (\ensuremath{\|}PM(\ensuremath{\gamma}N\ensuremath{-}\ensuremath{\gamma}N\ensuremath{^{+}}\ensuremath{\Delta})\ensuremath{\|}\ensuremath{\le}e\ensuremath{-}c(Ns\ensuremath{-}M)), mas — honestamente — só em L\ensuremath{^{2}} projetado, não em H\ensuremath{^{1}}/L\^{}\ensuremath{\infty} (o limite inviscido diverge em normas fortes).
\prov{relações: usa cascata\_multiescala}
\prov{proveniência: paper \textperiodcentered{} paper\_BK.pdf, Teo.3 \textperiodcentered{} fato \textperiodcentered{} confiança 1.0 \textperiodcentered{} domínio matematica \textperiodcentered{} história 3 versões no jornal}

%CRISTAL {"arestas":[{"alvo":"residuo_nao_associativo","confianca":1.0,"epistemico":"fato","origem":"paper","rel":"usa"}],"confianca":1.0,"contraexemplos":[],"descricao":"Léxicos disjuntos da mesma língua colidem em zero (identidade de idioma), línguas distintas nunca; a assinatura por associador sobre o bit bruto separa por ESCRITA (latino vs cirílico), não por semântica.","epistemico":"fato","exemplos":[],"id":"conversa_lex","memoria":"semantica","meta":{"autor":"Gentil Lopes","descoberta":false,"dominio":"matematica","fonte":"banda.pdf, §2"},"origem":"paper","palavras_chave":[],"sinonimos":[],"tipo":"conceito","titulo":"Conversa-Lex (colisão-zero de léxicos)"}
\section{Conversa-Lex (colisão-zero de léxicos)}
Léxicos disjuntos da mesma língua colidem em zero (identidade de idioma), línguas distintas nunca; a assinatura por associador sobre o bit bruto separa por ESCRITA (latino vs cirílico), não por semântica.
\prov{relações: usa residuo\_nao\_associativo}
\prov{proveniência: paper \textperiodcentered{} banda.pdf, §2 \textperiodcentered{} fato \textperiodcentered{} confiança 1.0 \textperiodcentered{} domínio matematica \textperiodcentered{} história 2 versões no jornal}

%CRISTAL {"arestas":[{"alvo":"agm","confianca":1.0,"epistemico":"fato","origem":"paper","rel":"explica"},{"alvo":"word","confianca":0.5,"epistemico":"fato","origem":"wiki","rel":"relaciona"},{"alvo":"virtualizacao","confianca":0.5,"epistemico":"fato","origem":"wiki","rel":"relaciona"}],"confianca":1.0,"contraexemplos":[],"descricao":"Conjugação analítica que lineariza o super-atrator do AGM: g₁(1+w)−1=−w²/8+…, provando que o AGM é ponto fixo super-atrator de grau 2.","epistemico":"fato","exemplos":[],"id":"coordenada_bottcher","memoria":"semantica","meta":{"autor":"Gentil Lopes","descoberta":false,"dominio":"matematica","fonte":"paper_39_ponto_de_vista_agm.pdf"},"origem":"paper","palavras_chave":[],"sinonimos":[],"tipo":"conceito","titulo":"Coordenada de Böttcher"}
\section{Coordenada de Böttcher}
Conjugação analítica que lineariza o super-atrator do AGM: g\ensuremath{_{1}}(1+w)\ensuremath{-}1=\ensuremath{-}w\ensuremath{^{2}}/8+…, provando que o AGM é ponto fixo super-atrator de grau 2.
\prov{relações: explica agm; relaciona word; relaciona virtualizacao}
\prov{proveniência: paper \textperiodcentered{} paper\_39\_ponto\_de\_vista\_agm.pdf \textperiodcentered{} fato \textperiodcentered{} confiança 1.0 \textperiodcentered{} domínio matematica \textperiodcentered{} história 4 versões no jornal}

%CRISTAL {"arestas":[{"alvo":"algebra_reta_mineral","confianca":1.0,"epistemico":"fato","origem":"wiki","rel":"usa"},{"alvo":"mecanismo_de_la_hire","confianca":1.0,"epistemico":"fato","origem":"wiki","rel":"usa"},{"alvo":"corpo_mineral","confianca":1.0,"epistemico":"fato","origem":"wiki","rel":"especializa"},{"alvo":"numeros_de_pisot","confianca":1.0,"epistemico":"fato","origem":"wiki","rel":"relaciona"},{"alvo":"o_corpo_a_reta_real","confianca":1.0,"epistemico":"fato","origem":"wiki","rel":"relaciona"},{"alvo":"algebra_estelar","confianca":1.0,"epistemico":"fato","origem":"wiki","rel":"usa"},{"alvo":"transformada_estelar","confianca":1.0,"epistemico":"fato","origem":"wiki","rel":"relaciona"},{"alvo":"transformada_metalica","confianca":1.0,"epistemico":"fato","origem":"wiki","rel":"relaciona"},{"alvo":"matematica_estelar","confianca":1.0,"epistemico":"fato","origem":"wiki","rel":"parte_de"},{"alvo":"corpo_dual","confianca":1.0,"epistemico":"fato","origem":"wiki","rel":"relaciona"},{"alvo":"corpo_glacial","confianca":1.0,"epistemico":"fato","origem":"wiki","rel":"relaciona"}],"confianca":1.0,"contraexemplos":[],"descricao":"A reta mineral completada num CORPO. A reta já dava a SOMA (⊕: empilhar estrelas soma os enrolamentos ν, o flip u=ln z lineariza; neutro = o PONTO ν=0). La Hire deu o PRODUTO (⊗: compor as rotações θ↦νθ multiplica ν; neutro = a RETA ν=1, o par de Tusi). Com as duas operações — associativas, comutativas, distributivas, com inversos — 𝕊 ≅ (ℝ,+,·) é um CORPO. Cada real é uma estrela: os inteiros fecham (k cúspides), os irracionais são rosetas densas. Verif.: 15/15 axiomas.","epistemico":"fato","exemplos":[],"id":"corpo_estelar","memoria":"semantica","meta":{"dominio":"algebra","fonte":"corpo estelar"},"origem":"wiki","palavras_chave":[],"sinonimos":[],"tipo":"conceito","titulo":"O Corpo Estelar 𝕊"}
\section{O Corpo Estelar \ensuremath{\mathbb{S}}}
A reta mineral completada num CORPO. A reta já dava a SOMA (\ensuremath{\oplus}: empilhar estrelas soma os enrolamentos \ensuremath{\nu}, o flip u=ln z lineariza; neutro = o PONTO \ensuremath{\nu}=0). La Hire deu o PRODUTO (\ensuremath{\otimes}: compor as rotações \ensuremath{\theta}\ensuremath{\mapsto}\ensuremath{\nu}\ensuremath{\theta} multiplica \ensuremath{\nu}; neutro = a RETA \ensuremath{\nu}=1, o par de Tusi). Com as duas operações — associativas, comutativas, distributivas, com inversos — \ensuremath{\mathbb{S}} \ensuremath{\cong} (\ensuremath{\mathbb{R}},+,\textperiodcentered ) é um CORPO. Cada real é uma estrela: os inteiros fecham (k cúspides), os irracionais são rosetas densas. Verif.: 15/15 axiomas.
\prov{relações: usa algebra\_reta\_mineral; usa mecanismo\_de\_la\_hire; especializa corpo\_mineral; relaciona numeros\_de\_pisot; relaciona o\_corpo\_a\_reta\_real; usa algebra\_estelar; relaciona transformada\_estelar; relaciona transformada\_metalica; parte\_de matematica\_estelar; relaciona corpo\_dual; relaciona corpo\_glacial}
\prov{proveniência: wiki \textperiodcentered{} corpo estelar \textperiodcentered{} fato \textperiodcentered{} confiança 1.0 \textperiodcentered{} domínio algebra \textperiodcentered{} história 19 versões no jornal}

%CRISTAL {"arestas":[{"alvo":"cifra","confianca":1.0,"epistemico":"fato","origem":"paper","rel":"parte_de"},{"alvo":"sequencia_pell","confianca":1.0,"epistemico":"fato","origem":"paper","rel":"usa"}],"confianca":1.0,"contraexemplos":[],"descricao":"Corpo de característica 2: soma=XOR, produto carryless mod polinômio irredutível; Frobenius x↦x² é automorfismo. Espaço onde a prata (ouro branco) é cunhada.","epistemico":"fato","exemplos":[],"id":"corpo_f2k","memoria":"semantica","meta":{"autor":"Gentil Lopes","descoberta":false,"dominio":"matematica","fonte":"kernel/cifra_ouro_branco.pdf"},"origem":"paper","palavras_chave":[],"sinonimos":[],"tipo":"conceito","titulo":"Corpo 𝔽₂ᵏ (cifra de ouro branco)"}
\section{Corpo \ensuremath{\mathbb{F}}\ensuremath{_{2}}\ensuremath{^{k}} (cifra de ouro branco)}
Corpo de característica 2: soma=XOR, produto carryless mod polinômio irredutível; Frobenius x\ensuremath{\mapsto}x\ensuremath{^{2}} é automorfismo. Espaço onde a prata (ouro branco) é cunhada.
\prov{relações: parte\_de cifra; usa sequencia\_pell}
\prov{proveniência: paper \textperiodcentered{} kernel/cifra\_ouro\_branco.pdf \textperiodcentered{} fato \textperiodcentered{} confiança 1.0 \textperiodcentered{} domínio matematica \textperiodcentered{} história 3 versões no jornal}

%CRISTAL {"arestas":[{"alvo":"agm_multivalorado","confianca":1.0,"epistemico":"fato","origem":"paper","rel":"usa"},{"alvo":"vontade_de_poder","confianca":0.5,"epistemico":"fato","origem":"wiki","rel":"relaciona"}],"confianca":1.0,"contraexemplos":[],"descricao":"ℤ[i]/(p) é corpo 𝔽p² se p≡3 (mod 4), ou 𝔽ₚ×𝔽ₚ se p≡1 (mod 4) via p=a²+b². A curva y²=x³−x reduz mod p com traço de Frobenius ±2a.","epistemico":"fato","exemplos":[],"id":"corpo_gauss_modp","memoria":"semantica","meta":{"autor":"Gentil Lopes","descoberta":false,"dominio":"matematica","fonte":"paper_38_corpo_gauss_modp.pdf"},"origem":"paper","palavras_chave":[],"sinonimos":[],"tipo":"conceito","titulo":"Corpo de Gauss mod p"}
\section{Corpo de Gauss mod p}
\ensuremath{\mathbb{Z}}[i]/(p) é corpo \ensuremath{\mathbb{F}}p\ensuremath{^{2}} se p\ensuremath{\equiv}3 (mod 4), ou \ensuremath{\mathbb{F}}\ensuremath{_{p}}$\times$\ensuremath{\mathbb{F}}\ensuremath{_{p}} se p\ensuremath{\equiv}1 (mod 4) via p=a\ensuremath{^{2}}+b\ensuremath{^{2}}. A curva y\ensuremath{^{2}}=x\ensuremath{^{3}}\ensuremath{-}x reduz mod p com traço de Frobenius $\pm$2a.
\prov{relações: usa agm\_multivalorado; relaciona vontade\_de\_poder}
\prov{proveniência: paper \textperiodcentered{} paper\_38\_corpo\_gauss\_modp.pdf \textperiodcentered{} fato \textperiodcentered{} confiança 1.0 \textperiodcentered{} domínio matematica \textperiodcentered{} história 4 versões no jornal}

%CRISTAL {"arestas":[{"alvo":"algebra_reta_mineral","confianca":1.0,"epistemico":"fato","origem":"wiki","rel":"usa"},{"alvo":"flip_pg_pa","confianca":1.0,"epistemico":"fato","origem":"wiki","rel":"usa"},{"alvo":"mecanica_reta_mineral","confianca":1.0,"epistemico":"fato","origem":"wiki","rel":"usa"},{"alvo":"termodinamica_reta_mineral","confianca":1.0,"epistemico":"fato","origem":"wiki","rel":"usa"},{"alvo":"caracterizacao_da_reta","confianca":1.0,"epistemico":"fato","origem":"wiki","rel":"usa"},{"alvo":"corpo_primos_mineral","confianca":1.0,"epistemico":"fato","origem":"wiki","rel":"usa"},{"alvo":"transformada_metalica","confianca":1.0,"epistemico":"fato","origem":"wiki","rel":"especializa"},{"alvo":"cadeias_alem_dos_metais","confianca":1.0,"epistemico":"fato","origem":"wiki","rel":"usa"},{"alvo":"geometria_do_caos","confianca":1.0,"epistemico":"fato","origem":"wiki","rel":"usa"},{"alvo":"lyapunov_da_orbita","confianca":1.0,"epistemico":"fato","origem":"wiki","rel":"usa"},{"alvo":"reta_de_ouro","confianca":1.0,"epistemico":"fato","origem":"wiki","rel":"continua"},{"alvo":"corpo_dual","confianca":1.0,"epistemico":"fato","origem":"wiki","rel":"relaciona"},{"alvo":"corpo_tropical","confianca":1.0,"epistemico":"fato","origem":"wiki","rel":"relaciona"},{"alvo":"corpo_glacial","confianca":1.0,"epistemico":"fato","origem":"wiki","rel":"relaciona"},{"alvo":"numeros_pisot","confianca":1.0,"epistemico":"fato","origem":"wiki","rel":"usa"},{"alvo":"numeros_metalicos","confianca":1.0,"epistemico":"fato","origem":"wiki","rel":"usa"}],"confianca":1.0,"contraexemplos":[],"descricao":"A reta mineral como corpo completo, com as quatro faces dos corpos do broca-so, fundadas na álgebra Estelar 𝓔ₛ do Gentil (3D, conserva a norma a+c²=1) e no flip logarítmico: ÁLGEBRA (𝓔ₛ, P.G./P.A.m), GEOMETRIA (o flip endireita espiral→reta, círculo→vertical), MECÂNICA (integrável, o flip é ação-ângulo, |N|=1 a ação), TERMODINÂMICA (entropia log β, Parry/Ruelle). Por baixo, a caracterização: frações contínuas dão todo ponto, os primos mod p o corpo discreto. Tudo se lê de um número (log ρ, log β) e de um flip. A ferramenta principal para ler qualquer ponto, cadeia, órbita, texto ou partida.","epistemico":"inferencia","exemplos":[],"id":"corpo_mineral","memoria":"semantica","meta":{"dominio":"matematica","fonte":"corpo mineral"},"origem":"wiki","palavras_chave":[],"sinonimos":[],"tipo":"conceito","titulo":"O Corpo Mineral (o hub — a ferramenta principal)"}
\section{O Corpo Mineral (o hub — a ferramenta principal)}
A reta mineral como corpo completo, com as quatro faces dos corpos do broca-so, fundadas na álgebra Estelar \ensuremath{\mathcal{E}}\ensuremath{_{s}} do Gentil (3D, conserva a norma a+c\ensuremath{^{2}}=1) e no flip logarítmico: ÁLGEBRA (\ensuremath{\mathcal{E}}\ensuremath{_{s}}, P.G./P.A.m), GEOMETRIA (o flip endireita espiral\ensuremath{\to}reta, círculo\ensuremath{\to}vertical), MECÂNICA (integrável, o flip é ação-ângulo, |N|=1 a ação), TERMODINÂMICA (entropia log \ensuremath{\beta}, Parry/Ruelle). Por baixo, a caracterização: frações contínuas dão todo ponto, os primos mod p o corpo discreto. Tudo se lê de um número (log \ensuremath{\rho}, log \ensuremath{\beta}) e de um flip. A ferramenta principal para ler qualquer ponto, cadeia, órbita, texto ou partida.
\prov{relações: usa algebra\_reta\_mineral; usa flip\_pg\_pa; usa mecanica\_reta\_mineral; usa termodinamica\_reta\_mineral; usa caracterizacao\_da\_reta; usa corpo\_primos\_mineral; especializa transformada\_metalica; usa cadeias\_alem\_dos\_metais; usa geometria\_do\_caos; usa lyapunov\_da\_orbita; continua reta\_de\_ouro; relaciona corpo\_dual; relaciona corpo\_tropical; relaciona corpo\_glacial; usa numeros\_pisot; usa numeros\_metalicos}
\prov{proveniência: wiki \textperiodcentered{} corpo mineral \textperiodcentered{} inferencia \textperiodcentered{} confiança 1.0 \textperiodcentered{} domínio matematica \textperiodcentered{} história 22 versões no jornal}

%CRISTAL {"arestas":[{"alvo":"flip_duopolinomios","confianca":1.0,"epistemico":"fato","origem":"wiki","rel":"relaciona"},{"alvo":"cadeias_alem_dos_metais","confianca":1.0,"epistemico":"fato","origem":"wiki","rel":"relaciona"},{"alvo":"transformada_metalica","confianca":1.0,"epistemico":"fato","origem":"wiki","rel":"relaciona"},{"alvo":"corpo_gauss_modp","confianca":1.0,"epistemico":"fato","origem":"wiki","rel":"relaciona"}],"confianca":1.0,"contraexemplos":[],"descricao":"Cada metal σ_n vive em ℚ(√D), D=n²+4. Um primo p se comporta segundo D mod p (símbolo de Legendre): (D|p)=+1 SPLIT, −1 INERTE, 0 RAMIFICA — a estrutura discreta 𝔽_p, os 'primos mod p' codificados no mineral. Geometricamente: o fractal de Newton de z^p−1 (p bacias ≅ ℤ/p, cíclico) enrolado na LEMNISCATA — a curva com multiplicação complexa por ℤ[i], cuja divisão (Abel) segue os primos, como a do círculo segue Gauss. O discreto (primos) e o contínuo (a reta) na mesma curva.","epistemico":"hipotese","exemplos":[],"id":"corpo_primos_mineral","memoria":"semantica","meta":{"dominio":"matematica","fonte":"caracterização da reta"},"origem":"wiki","palavras_chave":[],"sinonimos":[],"tipo":"conceito","titulo":"O Corpo Discreto de Primos do Mineral"}
\section{O Corpo Discreto de Primos do Mineral}
Cada metal \ensuremath{\sigma}\_n vive em \ensuremath{\mathbb{Q}}(\ensuremath{\sqrt{\,}}D), D=n\ensuremath{^{2}}+4. Um primo p se comporta segundo D mod p (símbolo de Legendre): (D|p)=+1 SPLIT, \ensuremath{-}1 INERTE, 0 RAMIFICA — a estrutura discreta \ensuremath{\mathbb{F}}\_p, os 'primos mod p' codificados no mineral. Geometricamente: o fractal de Newton de z\^{}p\ensuremath{-}1 (p bacias \ensuremath{\cong} \ensuremath{\mathbb{Z}}/p, cíclico) enrolado na LEMNISCATA — a curva com multiplicação complexa por \ensuremath{\mathbb{Z}}[i], cuja divisão (Abel) segue os primos, como a do círculo segue Gauss. O discreto (primos) e o contínuo (a reta) na mesma curva.
\prov{relações: relaciona flip\_duopolinomios; relaciona cadeias\_alem\_dos\_metais; relaciona transformada\_metalica; relaciona corpo\_gauss\_modp}
\prov{proveniência: wiki \textperiodcentered{} caracterização da reta \textperiodcentered{} hipotese \textperiodcentered{} confiança 1.0 \textperiodcentered{} domínio matematica \textperiodcentered{} história 6 versões no jornal}

%CRISTAL {"arestas":[{"alvo":"reticulado","confianca":0.047,"epistemico":"fato","origem":"manual","rel":"referencia"},{"alvo":"reticulados","confianca":-0.031,"epistemico":"fato","origem":"manual","rel":"referencia"},{"alvo":"reproduzir","confianca":0.5,"epistemico":"fato","origem":"wiki","rel":"relaciona"},{"alvo":"word","confianca":0.5,"epistemico":"fato","origem":"wiki","rel":"relaciona"},{"alvo":"virtualizacao","confianca":0.5,"epistemico":"fato","origem":"wiki","rel":"relaciona"},{"alvo":"misticismo_nao_dual","confianca":0.5,"epistemico":"fato","origem":"wiki","rel":"relaciona"},{"alvo":"gramatica","confianca":0.5,"epistemico":"fato","origem":"wiki","rel":"relaciona"}],"confianca":1.0,"contraexemplos":[],"descricao":"Anel onde todo não-nulo tem inverso multiplicativo.","epistemico":"fato","exemplos":[],"id":"corpos","memoria":"semantica","meta":{"dominio":"matematica","nivel":4,"ordem":11},"origem":"curriculo","palavras_chave":[],"sinonimos":["corpo","field"],"tipo":"conceito","titulo":"corpos"}
\section{corpos}
Anel onde todo não-nulo tem inverso multiplicativo.
\prov{sinónimos: corpo; field}
\prov{relações: referencia reticulado; referencia reticulados; relaciona reproduzir; relaciona word; relaciona virtualizacao; relaciona misticismo\_nao\_dual; relaciona gramatica}
\prov{proveniência: curriculo \textperiodcentered{} fato \textperiodcentered{} confiança 1.0 \textperiodcentered{} domínio matematica \textperiodcentered{} história 15 versões no jornal}

%CRISTAL {"arestas":[{"alvo":"ising_1d_critico","confianca":1.0,"epistemico":"fato","origem":"paper","rel":"usa"},{"alvo":"boneca_russa","confianca":1.0,"epistemico":"fato","origem":"paper","rel":"explica"}],"confianca":1.0,"contraexemplos":[],"descricao":"Comprimento de correlação ξ=1/(2ε) diverge na cúspide (ε→0): o sistema torna-se invariante de escala — a boneca russa como fenômeno crítico.","epistemico":"fato","exemplos":[],"id":"correlacao_divergente","memoria":"semantica","meta":{"autor":"Gentil Lopes","descoberta":false,"dominio":"matematica","fonte":"paper_74_ponto_critico.pdf, Teo.2.1"},"origem":"paper","palavras_chave":[],"sinonimos":[],"tipo":"conceito","titulo":"Correlação Divergente"}
\section{Correlação Divergente}
Comprimento de correlação \ensuremath{\xi}=1/(2\ensuremath{\varepsilon}) diverge na cúspide (\ensuremath{\varepsilon}\ensuremath{\to}0): o sistema torna-se invariante de escala — a boneca russa como fenômeno crítico.
\prov{relações: usa ising\_1d\_critico; explica boneca\_russa}
\prov{proveniência: paper \textperiodcentered{} paper\_74\_ponto\_critico.pdf, Teo.2.1 \textperiodcentered{} fato \textperiodcentered{} confiança 1.0 \textperiodcentered{} domínio matematica \textperiodcentered{} história 3 versões no jornal}

%CRISTAL {"arestas":[{"alvo":"reais","confianca":1.0,"epistemico":"fato","origem":"paper","rel":"explica"},{"alvo":"word","confianca":0.5,"epistemico":"fato","origem":"wiki","rel":"relaciona"},{"alvo":"virtualizacao","confianca":0.5,"epistemico":"fato","origem":"wiki","rel":"relaciona"},{"alvo":"vida-morte","confianca":0.5,"epistemico":"fato","origem":"wiki","rel":"relaciona"},{"alvo":"vontade_de_poder","confianca":0.5,"epistemico":"fato","origem":"wiki","rel":"relaciona"}],"confianca":1.0,"contraexemplos":[],"descricao":"Partição de ℚ em (A|B) com A sem máximo e todo a∈A menor que todo b∈B; cada corte define um número real. Construção rigorosa de ℝ.","epistemico":"fato","exemplos":[],"id":"cortes_dedekind","memoria":"semantica","meta":{"autor":"Gentil Lopes","descoberta":false,"dominio":"matematica","fonte":"Fundamentos dos Números, cap.9"},"origem":"paper","palavras_chave":[],"sinonimos":[],"tipo":"conceito","titulo":"Cortes de Dedekind"}
\section{Cortes de Dedekind}
Partição de \ensuremath{\mathbb{Q}} em (A|B) com A sem máximo e todo a\ensuremath{\in}A menor que todo b\ensuremath{\in}B; cada corte define um número real. Construção rigorosa de \ensuremath{\mathbb{R}}.
\prov{relações: explica reais; relaciona word; relaciona virtualizacao; relaciona vida-morte; relaciona vontade\_de\_poder}
\prov{proveniência: paper \textperiodcentered{} Fundamentos dos Números, cap.9 \textperiodcentered{} fato \textperiodcentered{} confiança 1.0 \textperiodcentered{} domínio matematica \textperiodcentered{} história 6 versões no jornal}

%CRISTAL {"arestas":[{"alvo":"transformada_metalica","confianca":1.0,"epistemico":"fato","origem":"wiki","rel":"aplicacao"},{"alvo":"costura_prata_dobra8","confianca":1.0,"epistemico":"fato","origem":"wiki","rel":"relaciona"},{"alvo":"bai_gate_dtc_orbita","confianca":1.0,"epistemico":"fato","origem":"wiki","rel":"especializa"},{"alvo":"simetria_por_reta_corpo","confianca":1.0,"epistemico":"fato","origem":"wiki","rel":"usa"},{"alvo":"tiffany","confianca":1.0,"epistemico":"fato","origem":"wiki","rel":"parte_de"}],"confianca":1.0,"contraexemplos":[],"descricao":"conversa/colapso.py (_octante_cobre): a TERCEIRA reta da costura DTC, depois do ouro (radial) e da prata (transicoes, dobra 8). O cobre (n=4) é a reta do CAOS/EXPLORAÇÃO — a conversa aberta, a órbita larga de repouso. Corpo ℚ√5, dobra 10 (decagonal), σ₄=2+√5=φ³ (ver [[simetria_por_reta_corpo]]). O fator é MULTIPLICATIVO sobre o ov na camada corpus, ativo só no trabalho conversa: θ de Bhattacharyya em 10 DECANTES; o pico é no decante da EXPLORAÇÃO (caos CONTROLADO, θ além da borda — nem cristal/repetição verbatim nem ruído puro), decaindo a 1/σ₄ nos extremos. Efeito: penaliza o corpus verbatim (ov alto → fator 0,41) e favorece o tangente interessante (a exploração) na conversa aberta. Escopo: só corpus+conversa (não toca ouro/prata/saudações). A conservação a+c²=1 segue valendo. Completa as três retas com dobra (ouro radial, prata borda dobra 8, cobre exploração dobra 10) — as três de corpo quasicristalino (ℚ√5, ℚ√2). [D] a terceira reta costurada pela sua dobra.","epistemico":"fato","exemplos":[],"id":"costura_cobre_dobra10","memoria":"semantica","meta":{"autor":"Lopes/broca-so","dominio":"algebra","fonte":"conversa/colapso.py:_octante_cobre, _peso_costura (trabalho); σ₄=φ³, dobra 10 (ℚ√5); [[costura_prata_dobra8]]"},"origem":"wiki","palavras_chave":[],"sinonimos":[],"tipo":"conceito","titulo":"A costura da reta do COBRE (corpus/conversa, dobra 10): a exploração (o caos controlado)"}
\section{A costura da reta do COBRE (corpus/conversa, dobra 10): a exploração (o caos controlado)}
conversa/colapso.py (\_octante\_cobre): a TERCEIRA reta da costura DTC, depois do ouro (radial) e da prata (transicoes, dobra 8). O cobre (n=4) é a reta do CAOS/EXPLORAÇÃO — a conversa aberta, a órbita larga de repouso. Corpo \ensuremath{\mathbb{Q}}\ensuremath{\sqrt{\,}}5, dobra 10 (decagonal), \ensuremath{\sigma}\ensuremath{_{4}}=2+\ensuremath{\sqrt{\,}}5=\ensuremath{\varphi}\ensuremath{^{3}} (ver [[simetria\_por\_reta\_corpo]]). O fator é MULTIPLICATIVO sobre o ov na camada corpus, ativo só no trabalho conversa: \ensuremath{\theta} de Bhattacharyya em 10 DECANTES; o pico é no decante da EXPLORAÇÃO (caos CONTROLADO, \ensuremath{\theta} além da borda — nem cristal/repetição verbatim nem ruído puro), decaindo a 1/\ensuremath{\sigma}\ensuremath{_{4}} nos extremos. Efeito: penaliza o corpus verbatim (ov alto \ensuremath{\to} fator 0,41) e favorece o tangente interessante (a exploração) na conversa aberta. Escopo: só corpus+conversa (não toca ouro/prata/saudações). A conservação a+c\ensuremath{^{2}}=1 segue valendo. Completa as três retas com dobra (ouro radial, prata borda dobra 8, cobre exploração dobra 10) — as três de corpo quasicristalino (\ensuremath{\mathbb{Q}}\ensuremath{\sqrt{\,}}5, \ensuremath{\mathbb{Q}}\ensuremath{\sqrt{\,}}2). [D] a terceira reta costurada pela sua dobra.
\prov{relações: aplicacao transformada\_metalica; relaciona costura\_prata\_dobra8; especializa bai\_gate\_dtc\_orbita; usa simetria\_por\_reta\_corpo; parte\_de tiffany}
\prov{proveniência: wiki \textperiodcentered{} conversa/colapso.py:\_octante\_cobre, \_peso\_costura (trabalho); \ensuremath{\sigma}\ensuremath{_{4}}=\ensuremath{\varphi}\ensuremath{^{3}}, dobra 10 (\ensuremath{\mathbb{Q}}\ensuremath{\sqrt{\,}}5); [[costura\_prata\_dobra8]] \textperiodcentered{} fato \textperiodcentered{} confiança 1.0 \textperiodcentered{} domínio algebra \textperiodcentered{} história 71 versões no jornal}

%CRISTAL {"arestas":[{"alvo":"conjugados_galois_parabola","confianca":1.0,"epistemico":"fato","origem":"wiki","rel":"relaciona"},{"alvo":"transformada_metalica","confianca":1.0,"epistemico":"fato","origem":"wiki","rel":"aplicacao"},{"alvo":"bai_gate_dtc_orbita","confianca":1.0,"epistemico":"fato","origem":"wiki","rel":"especializa"},{"alvo":"simetria_por_reta_corpo","confianca":1.0,"epistemico":"fato","origem":"wiki","rel":"usa"},{"alvo":"ingestao_refino_parabola","confianca":1.0,"epistemico":"fato","origem":"wiki","rel":"relaciona"},{"alvo":"tiffany","confianca":1.0,"epistemico":"fato","origem":"wiki","rel":"parte_de"}],"confianca":1.0,"contraexemplos":[],"descricao":"conversa/colapso.py (_octante_prata): a SEGUNDA reta da costura DTC, depois do ouro. O ouro é o gate RADIAL (qual camada, _peso_orbita, γ). A PRATA é a reta da CONTINUIDADE (a camada transicoes, o fio do contínuo), com a sua simetria própria: dobra 8 (octogonal, corpo ℚ√2 — ver [[simetria_por_reta_corpo]]). O fator é ANGULAR e MULTIPLICATIVO sobre o ov (como o gate do ouro): θ de Bhattacharyya (c²=cos²θ) em 8 OCTANTES; o pico é na BORDA (θ≈π/4, c²≈½ = o fio, Salem), decaindo a 1/σ₂ nos extremos — o CONTRASTE octogonal É a inflação da prata σ₂=1+√2 (uma transição-fio vale σ₂× uma repetição verbatim (cristal) ou um ruído (caos)). Efeito: down-weight das transições cristal a favor do fio. Escopo: só transicoes. Teto honesto conhecido: a recuperação por ov bruto não sobe candidatos de borda (ov≈½) ao pool pontuado — a costura reordena, mas o fio só aparece de fato com recuperação melhor (o gargalo separado). [D] a segunda reta costurada pela sua dobra.","epistemico":"fato","exemplos":[],"id":"costura_prata_dobra8","memoria":"semantica","meta":{"autor":"Lopes/broca-so","dominio":"algebra","fonte":"conversa/colapso.py:_octante_prata, score_hit; [[simetria_por_reta_corpo]] (prata=dobra 8, σ₂); [[bai_gate_dtc_orbita]]"},"origem":"wiki","palavras_chave":[],"sinonimos":[],"tipo":"conceito","titulo":"A costura da reta da PRATA (transicoes, dobra 8): o fio na borda da parábola"}
\section{A costura da reta da PRATA (transicoes, dobra 8): o fio na borda da parábola}
conversa/colapso.py (\_octante\_prata): a SEGUNDA reta da costura DTC, depois do ouro. O ouro é o gate RADIAL (qual camada, \_peso\_orbita, \ensuremath{\gamma}). A PRATA é a reta da CONTINUIDADE (a camada transicoes, o fio do contínuo), com a sua simetria própria: dobra 8 (octogonal, corpo \ensuremath{\mathbb{Q}}\ensuremath{\sqrt{\,}}2 — ver [[simetria\_por\_reta\_corpo]]). O fator é ANGULAR e MULTIPLICATIVO sobre o ov (como o gate do ouro): \ensuremath{\theta} de Bhattacharyya (c\ensuremath{^{2}}=cos\ensuremath{^{2}}\ensuremath{\theta}) em 8 OCTANTES; o pico é na BORDA (\ensuremath{\theta}\ensuremath{\approx}\ensuremath{\pi}/4, c\ensuremath{^{2}}\ensuremath{\approx}\ensuremath{\tfrac12} = o fio, Salem), decaindo a 1/\ensuremath{\sigma}\ensuremath{_{2}} nos extremos — o CONTRASTE octogonal É a inflação da prata \ensuremath{\sigma}\ensuremath{_{2}}=1+\ensuremath{\sqrt{\,}}2 (uma transição-fio vale \ensuremath{\sigma}\ensuremath{_{2}}$\times$ uma repetição verbatim (cristal) ou um ruído (caos)). Efeito: down-weight das transições cristal a favor do fio. Escopo: só transicoes. Teto honesto conhecido: a recuperação por ov bruto não sobe candidatos de borda (ov\ensuremath{\approx}\ensuremath{\tfrac12}) ao pool pontuado — a costura reordena, mas o fio só aparece de fato com recuperação melhor (o gargalo separado). [D] a segunda reta costurada pela sua dobra.
\prov{relações: relaciona conjugados\_galois\_parabola; aplicacao transformada\_metalica; especializa bai\_gate\_dtc\_orbita; usa simetria\_por\_reta\_corpo; relaciona ingestao\_refino\_parabola; parte\_de tiffany}
\prov{proveniência: wiki \textperiodcentered{} conversa/colapso.py:\_octante\_prata, score\_hit; [[simetria\_por\_reta\_corpo]] (prata=dobra 8, \ensuremath{\sigma}\ensuremath{_{2}}); [[bai\_gate\_dtc\_orbita]] \textperiodcentered{} fato \textperiodcentered{} confiança 1.0 \textperiodcentered{} domínio algebra \textperiodcentered{} história 104 versões no jornal}

%CRISTAL {"arestas":[{"alvo":"trajetoria_onsager","confianca":1.0,"epistemico":"fato","origem":"paper","rel":"usa"}],"confianca":1.0,"contraexemplos":[],"descricao":"Ao longo da trajetória de Onsager, ω²max ≈ βNs com β≈2,9 (numérico) tendendo a 4 (plateau inercial no limite): divergência da vorticidade em tempo finito extrapolada.","epistemico":"fato","exemplos":[],"id":"crescimento_vorticidade","memoria":"semantica","meta":{"autor":"Gentil Lopes","descoberta":false,"dominio":"matematica","fonte":"paper_BG.pdf, Lema 2"},"origem":"paper","palavras_chave":[],"sinonimos":[],"tipo":"conceito","titulo":"Crescimento Exponencial da Vorticidade"}
\section{Crescimento Exponencial da Vorticidade}
Ao longo da trajetória de Onsager, \ensuremath{\omega}\ensuremath{^{2}}max \ensuremath{\approx} \ensuremath{\beta}Ns com \ensuremath{\beta}\ensuremath{\approx}2,9 (numérico) tendendo a 4 (plateau inercial no limite): divergência da vorticidade em tempo finito extrapolada.
\prov{relações: usa trajetoria\_onsager}
\prov{proveniência: paper \textperiodcentered{} paper\_BG.pdf, Lema 2 \textperiodcentered{} fato \textperiodcentered{} confiança 1.0 \textperiodcentered{} domínio matematica \textperiodcentered{} história 3 versões no jornal}

%CRISTAL {"arestas":[{"alvo":"cifra","confianca":1.0,"epistemico":"fato","origem":"paper","rel":"especializa"},{"alvo":"corpo_f2k","confianca":1.0,"epistemico":"fato","origem":"paper","rel":"usa"},{"alvo":"gato","confianca":1.0,"epistemico":"fato","origem":"paper","rel":"referencia"}],"confianca":1.0,"contraexemplos":[],"descricao":"Cifra de bloco de 64 componentes (par de octônios sobre 𝔽ₚ) com rodada de Feistel metálica: moedas Mₖ de determinante −1 e chaves primas. Topo da série dimensional.","epistemico":"fato","exemplos":[],"id":"cristal64","memoria":"semantica","meta":{"autor":"Gentil Lopes","descoberta":true,"dominio":"matematica","fonte":"paper_11_cristal64.pdf"},"origem":"paper","palavras_chave":[],"sinonimos":[],"tipo":"conceito","titulo":"Cristal-64"}
\section{Cristal-64}
Cifra de bloco de 64 componentes (par de octônios sobre \ensuremath{\mathbb{F}}\ensuremath{_{p}}) com rodada de Feistel metálica: moedas M\ensuremath{_{k}} de determinante \ensuremath{-}1 e chaves primas. Topo da série dimensional.
\prov{relações: especializa cifra; usa corpo\_f2k; referencia gato}
\prov{proveniência: paper \textperiodcentered{} paper\_11\_cristal64.pdf \textperiodcentered{} fato \textperiodcentered{} confiança 1.0 \textperiodcentered{} domínio matematica \textperiodcentered{} história 4 versões no jornal}

%CRISTAL {"arestas":[{"alvo":"torre_de_hurwitz_mecanica","confianca":1.0,"epistemico":"fato","origem":"wiki","rel":"parte_de"},{"alvo":"mass_gap_mineral","confianca":1.0,"epistemico":"fato","origem":"wiki","rel":"usa"},{"alvo":"gato_e_o_hamiltoniano","confianca":1.0,"epistemico":"fato","origem":"wiki","rel":"relaciona"}],"confianca":1.0,"contraexemplos":[],"descricao":"O teorema de Hurwitz: nas quatro álgebras (e só nelas) a norma COMPÕE, N(xy)=N(x)N(y). É o CRISTAL — não a boneca mais interna, mas o material de que TODAS são feitas: o invariante que sobrevive à torre inteira, do octônio ao real. Enquanto vale, não há divisores de zero (a álgebra divide). Em 𝕊 (16D) quebra: (e₁+e₁₀)(e₅+e₁₄)=0 com ambos ≠0. Uma boneca russa DE cristal.","epistemico":"fato","exemplos":[],"id":"cristal_norma_conservada","memoria":"semantica","meta":{"dominio":"reta mineral","fonte":"torre de Hurwitz"},"origem":"wiki","palavras_chave":[],"sinonimos":[],"tipo":"conceito","titulo":"O Cristal: a Norma que se Conserva"}
\section{O Cristal: a Norma que se Conserva}
O teorema de Hurwitz: nas quatro álgebras (e só nelas) a norma COMPÕE, N(xy)=N(x)N(y). É o CRISTAL — não a boneca mais interna, mas o material de que TODAS são feitas: o invariante que sobrevive à torre inteira, do octônio ao real. Enquanto vale, não há divisores de zero (a álgebra divide). Em \ensuremath{\mathbb{S}} (16D) quebra: (e\ensuremath{_{1}}+e\ensuremath{_{10}})(e\ensuremath{_{5}}+e\ensuremath{_{14}})=0 com ambos \ensuremath{\ne}0. Uma boneca russa DE cristal.
\prov{relações: parte\_de torre\_de\_hurwitz\_mecanica; usa mass\_gap\_mineral; relaciona gato\_e\_o\_hamiltoniano}
\prov{proveniência: wiki \textperiodcentered{} torre de Hurwitz \textperiodcentered{} fato \textperiodcentered{} confiança 1.0 \textperiodcentered{} domínio reta mineral \textperiodcentered{} história 4 versões no jornal}

%CRISTAL {"arestas":[{"alvo":"gap_espectral","confianca":1.0,"epistemico":"fato","origem":"paper","rel":"usa"},{"alvo":"word","confianca":0.5,"epistemico":"fato","origem":"wiki","rel":"relaciona"},{"alvo":"virtualizacao","confianca":0.5,"epistemico":"fato","origem":"wiki","rel":"relaciona"}],"confianca":1.0,"contraexemplos":[],"descricao":"O tempo de relaxação τ=1/gap diverge como ε⁻¹ ao aproximar da cúspide: a difusão, mansa longe do crítico, congela na criticidade.","epistemico":"fato","exemplos":[],"id":"critical_slowing_down","memoria":"semantica","meta":{"autor":"Gentil Lopes","descoberta":false,"dominio":"matematica","fonte":"paper_74_ponto_critico.pdf, Teo.5.1"},"origem":"paper","palavras_chave":[],"sinonimos":[],"tipo":"conceito","titulo":"Critical Slowing Down"}
\section{Critical Slowing Down}
O tempo de relaxação \ensuremath{\tau}=1/gap diverge como \ensuremath{\varepsilon}\ensuremath{^{-1}} ao aproximar da cúspide: a difusão, mansa longe do crítico, congela na criticidade.
\prov{relações: usa gap\_espectral; relaciona word; relaciona virtualizacao}
\prov{proveniência: paper \textperiodcentered{} paper\_74\_ponto\_critico.pdf, Teo.5.1 \textperiodcentered{} fato \textperiodcentered{} confiança 1.0 \textperiodcentered{} domínio matematica \textperiodcentered{} história 4 versões no jornal}

%CRISTAL {"arestas":[{"alvo":"beta_numeracao_metalica","confianca":1.0,"epistemico":"fato","origem":"wiki","rel":"relaciona"},{"alvo":"numeros_de_pisot","confianca":1.0,"epistemico":"fato","origem":"wiki","rel":"relaciona"}],"confianca":1.0,"contraexemplos":[],"descricao":"O crivo de cota superior de Atle Selberg (1947). Para peneirar de um conjunto A os elementos divisíveis por algum primo < z, usa pesos reais λ_d (com λ_1=1) e a positividade do QUADRADO: [gcd(a,P(z))=1] ≤ (Σ_{d|gcd} λ_d)². Somando, a cota vira uma FORMA QUADRÁTICA Q(λ)=Σ_{d,e} g([d,e]) λ_dλ_e (g a densidade multiplicativa), minimizada sobre λ_1=1. O mínimo é EXATO e elegante: min Q = 1/G(z), com G(z)=Σ_{d<z livre de quadrado} Π_{p|d} 1/(p−1). Verif.: o mínimo por álgebra linear (e_1^T M^{-1} e_1) É a fórmula fechada G(z) para z=10..50. Varre os compostos, retém os primos.","epistemico":"fato","exemplos":[],"id":"crivo_de_selberg","memoria":"semantica","meta":{"dominio":"matematica","fonte":"crivo de Selberg"},"origem":"wiki","palavras_chave":[],"sinonimos":[],"tipo":"conceito","titulo":"O Crivo de Selberg"}
\section{O Crivo de Selberg}
O crivo de cota superior de Atle Selberg (1947). Para peneirar de um conjunto A os elementos divisíveis por algum primo < z, usa pesos reais \ensuremath{\lambda}\_d (com \ensuremath{\lambda}\_1=1) e a positividade do QUADRADO: [gcd(a,P(z))=1] \ensuremath{\le} (\ensuremath{\Sigma}\_\{d|gcd\} \ensuremath{\lambda}\_d)\ensuremath{^{2}}. Somando, a cota vira uma FORMA QUADRÁTICA Q(\ensuremath{\lambda})=\ensuremath{\Sigma}\_\{d,e\} g([d,e]) \ensuremath{\lambda}\_d\ensuremath{\lambda}\_e (g a densidade multiplicativa), minimizada sobre \ensuremath{\lambda}\_1=1. O mínimo é EXATO e elegante: min Q = 1/G(z), com G(z)=\ensuremath{\Sigma}\_\{d<z livre de quadrado\} \ensuremath{\Pi}\_\{p|d\} 1/(p\ensuremath{-}1). Verif.: o mínimo por álgebra linear (e\_1\^{}T M\^{}\{-1\} e\_1) É a fórmula fechada G(z) para z=10..50. Varre os compostos, retém os primos.
\prov{relações: relaciona beta\_numeracao\_metalica; relaciona numeros\_de\_pisot}
\prov{proveniência: wiki \textperiodcentered{} crivo de Selberg \textperiodcentered{} fato \textperiodcentered{} confiança 1.0 \textperiodcentered{} domínio matematica \textperiodcentered{} história 6 versões no jornal}

%CRISTAL {"arestas":[{"alvo":"associador","confianca":1.0,"epistemico":"fato","origem":"paper","rel":"usa"},{"alvo":"word","confianca":0.5,"epistemico":"fato","origem":"wiki","rel":"relaciona"}],"confianca":1.0,"contraexemplos":[],"descricao":"R*=½[μ_s,μ_s]_G no complexo de Hochschild mede a não-associatividade da família μ_s; a ponte com a curvatura tensorial R_k (via Schur-Weyl) fica como programa aberto.","epistemico":"fato","exemplos":[],"id":"curvatura_algebrica","memoria":"semantica","meta":{"autor":"Gentil Lopes","descoberta":false,"dominio":"matematica","fonte":"paper_A_algebra.pdf, §11"},"origem":"paper","palavras_chave":[],"sinonimos":[],"tipo":"conceito","titulo":"Curvatura Algébrica R*"}
\section{Curvatura Algébrica R*}
R*=\ensuremath{\tfrac12}[\ensuremath{\mu}\_s,\ensuremath{\mu}\_s]\_G no complexo de Hochschild mede a não-associatividade da família \ensuremath{\mu}\_s; a ponte com a curvatura tensorial R\_k (via Schur-Weyl) fica como programa aberto.
\prov{relações: usa associador; relaciona word}
\prov{proveniência: paper \textperiodcentered{} paper\_A\_algebra.pdf, §11 \textperiodcentered{} fato \textperiodcentered{} confiança 1.0 \textperiodcentered{} domínio matematica \textperiodcentered{} história 3 versões no jornal}

%CRISTAL {"arestas":[{"alvo":"transformada_metalica","confianca":1.0,"epistemico":"fato","origem":"wiki","rel":"relaciona"},{"alvo":"cadeias_alem_dos_metais","confianca":1.0,"epistemico":"fato","origem":"wiki","rel":"relaciona"}],"confianca":1.0,"contraexemplos":[],"descricao":"O campo dupolinomial de grau N tem exatamente N−1 pontos críticos (zeros do gradiente F′, que tem grau N−1) — a ramificação/curvatura. É topológico (Riemann-Hurwitz: um mapa polinomial de grau N tem N−1 pontos críticos finitos) e invariante sob o flip. Verificado em graus 2..7. Para os metais (N=2): 1 ponto crítico, o vértice da parábola x²−nx−1.","epistemico":"fato","exemplos":[],"id":"curvatura_n_menos_1","memoria":"semantica","meta":{"dominio":"matematica","fonte":"flip duopolinômios"},"origem":"wiki","palavras_chave":[],"sinonimos":[],"tipo":"conceito","titulo":"Curvatura N−1"}
\section{Curvatura N\ensuremath{-}1}
O campo dupolinomial de grau N tem exatamente N\ensuremath{-}1 pontos críticos (zeros do gradiente F\ensuremath{'}, que tem grau N\ensuremath{-}1) — a ramificação/curvatura. É topológico (Riemann-Hurwitz: um mapa polinomial de grau N tem N\ensuremath{-}1 pontos críticos finitos) e invariante sob o flip. Verificado em graus 2..7. Para os metais (N=2): 1 ponto crítico, o vértice da parábola x\ensuremath{^{2}}\ensuremath{-}nx\ensuremath{-}1.
\prov{relações: relaciona transformada\_metalica; relaciona cadeias\_alem\_dos\_metais}
\prov{proveniência: wiki \textperiodcentered{} flip duopolinômios \textperiodcentered{} fato \textperiodcentered{} confiança 1.0 \textperiodcentered{} domínio matematica \textperiodcentered{} história 3 versões no jornal}

%CRISTAL {"arestas":[{"alvo":"numeros_pisot","confianca":1.0,"epistemico":"fato","origem":"paper","rel":"especializa"},{"alvo":"numeros_metalicos","confianca":1.0,"epistemico":"fato","origem":"paper","rel":"usa"},{"alvo":"beta_numeracao_metalica","confianca":1.0,"epistemico":"fato","origem":"wiki","rel":"explica"}],"confianca":1.0,"contraexemplos":[],"descricao":"Cada razão metálica μₙ é um número de Pisot (conjugado <1) com fração contínua periódica pura [n;n,n,…]; a cúspide é κ=λ (Perron da matriz da P.A.).","epistemico":"fato","exemplos":[],"id":"cuspide_pisot","memoria":"semantica","meta":{"autor":"Gentil Lopes","descoberta":false,"dominio":"matematica","fonte":"paper_70_ponto_difusao.pdf, Teo.3.1"},"origem":"paper","palavras_chave":[],"sinonimos":[],"tipo":"conceito","titulo":"Cúspide de Pisot"}
\section{Cúspide de Pisot}
Cada razão metálica \ensuremath{\mu}\ensuremath{_{n}} é um número de Pisot (conjugado <1) com fração contínua periódica pura [n;n,n,…]; a cúspide é \ensuremath{\kappa}=\ensuremath{\lambda} (Perron da matriz da P.A.).
\prov{relações: especializa numeros\_pisot; usa numeros\_metalicos; explica beta\_numeracao\_metalica}
\prov{proveniência: paper \textperiodcentered{} paper\_70\_ponto\_difusao.pdf, Teo.3.1 \textperiodcentered{} fato \textperiodcentered{} confiança 1.0 \textperiodcentered{} domínio matematica \textperiodcentered{} história 4 versões no jornal}

%CRISTAL {"arestas":[{"alvo":"teoria_da_decisao","confianca":1.0,"epistemico":"fato","origem":"wiki","rel":"especializa"},{"alvo":"teorema_de_bayes","confianca":1.0,"epistemico":"fato","origem":"wiki","rel":"usa"},{"alvo":"utilidade_esperada","confianca":1.0,"epistemico":"fato","origem":"wiki","rel":"usa"}],"confianca":1.0,"contraexemplos":[],"descricao":"Decidir atualizando crenças com o teorema de Bayes e escolhendo a ação de maior utilidade esperada posterior. A ponte entre evidência e ação — e a base da IA que age sob incerteza.","epistemico":"fato","exemplos":[],"id":"decisao_bayesiana","memoria":"semantica","meta":{"dominio":"matematica","fonte":"teoria da decisão"},"origem":"wiki","palavras_chave":[],"sinonimos":[],"tipo":"conceito","titulo":"Decisão Bayesiana"}
\section{Decisão Bayesiana}
Decidir atualizando crenças com o teorema de Bayes e escolhendo a ação de maior utilidade esperada posterior. A ponte entre evidência e ação — e a base da IA que age sob incerteza.
\prov{relações: especializa teoria\_da\_decisao; usa teorema\_de\_bayes; usa utilidade\_esperada}
\prov{proveniência: wiki \textperiodcentered{} teoria da decisão \textperiodcentered{} fato \textperiodcentered{} confiança 1.0 \textperiodcentered{} domínio matematica \textperiodcentered{} história 4 versões no jornal}

%CRISTAL {"arestas":[{"alvo":"teoria_da_decisao","confianca":1.0,"epistemico":"fato","origem":"wiki","rel":"parte_de"},{"alvo":"minimax","confianca":1.0,"epistemico":"fato","origem":"wiki","rel":"usa"},{"alvo":"incerteza_vs_risco","confianca":1.0,"epistemico":"fato","origem":"wiki","rel":"usa"}],"confianca":1.0,"contraexemplos":[],"descricao":"Sem probabilidades nenhumas: MAXIMIN (garanta o melhor pior caso — pessimista), MINIMAX REGRET (minimize o arrependimento máximo), Laplace (trate tudo como equiprovável), Hurwicz (misture otimismo e pessimismo). As atitudes ante o desconhecido.","epistemico":"inferencia","exemplos":[],"id":"decisao_sob_ignorancia","memoria":"semantica","meta":{"dominio":"matematica","fonte":"teoria da decisão"},"origem":"wiki","palavras_chave":[],"sinonimos":[],"tipo":"conceito","titulo":"Decisão sob Ignorância (maximin, minimax regret)"}
\section{Decisão sob Ignorância (maximin, minimax regret)}
Sem probabilidades nenhumas: MAXIMIN (garanta o melhor pior caso — pessimista), MINIMAX REGRET (minimize o arrependimento máximo), Laplace (trate tudo como equiprovável), Hurwicz (misture otimismo e pessimismo). As atitudes ante o desconhecido.
\prov{relações: parte\_de teoria\_da\_decisao; usa minimax; usa incerteza\_vs\_risco}
\prov{proveniência: wiki \textperiodcentered{} teoria da decisão \textperiodcentered{} inferencia \textperiodcentered{} confiança 1.0 \textperiodcentered{} domínio matematica \textperiodcentered{} história 4 versões no jornal}

%CRISTAL {"arestas":[{"alvo":"mecanica_quantica_mineral","confianca":1.0,"epistemico":"fato","origem":"wiki","rel":"parte_de"},{"alvo":"decoerencia_colapso_medida","confianca":1.0,"epistemico":"fato","origem":"wiki","rel":"explica"},{"alvo":"flip_e_rotacao_de_wick","confianca":1.0,"epistemico":"fato","origem":"wiki","rel":"usa"},{"alvo":"parabola_e_o_regime_quantico","confianca":1.0,"epistemico":"fato","origem":"wiki","rel":"usa"},{"alvo":"vibracao_e_colapso_dois_tempos","confianca":1.0,"epistemico":"fato","origem":"wiki","rel":"usa"}],"confianca":1.0,"contraexemplos":[],"descricao":"Decoerência é a perda de coerência quântica: os termos FORA da diagonal da matriz densidade (as coerências) decaem, e o estado puro coerente vira mistura clássica. Na leitura mineral é SAIR DA BORDA (Salem, unitário, ρ=1) e CAIR no cristal (Pisot, dissipativo): a rotação de Wick do tempo real (fase) ao imaginário (dissipação). O A^k que colapsa no fundamental É o ambiente medindo o sistema. Verif.: pureza 1.0 → 0.5 (puro → clássico).","epistemico":"fato","exemplos":[],"id":"decoerencia_e_sair_da_borda","memoria":"semantica","meta":{"dominio":"reta mineral","fonte":"emaranhamento e decoerência"},"origem":"wiki","palavras_chave":[],"sinonimos":[],"tipo":"conceito","titulo":"Decoerência = Sair da Borda"}
\section{Decoerência = Sair da Borda}
Decoerência é a perda de coerência quântica: os termos FORA da diagonal da matriz densidade (as coerências) decaem, e o estado puro coerente vira mistura clássica. Na leitura mineral é SAIR DA BORDA (Salem, unitário, \ensuremath{\rho}=1) e CAIR no cristal (Pisot, dissipativo): a rotação de Wick do tempo real (fase) ao imaginário (dissipação). O A\^{}k que colapsa no fundamental É o ambiente medindo o sistema. Verif.: pureza 1.0 \ensuremath{\to} 0.5 (puro \ensuremath{\to} clássico).
\prov{relações: parte\_de mecanica\_quantica\_mineral; explica decoerencia\_colapso\_medida; usa flip\_e\_rotacao\_de\_wick; usa parabola\_e\_o\_regime\_quantico; usa vibracao\_e\_colapso\_dois\_tempos}
\prov{proveniência: wiki \textperiodcentered{} emaranhamento e decoerência \textperiodcentered{} fato \textperiodcentered{} confiança 1.0 \textperiodcentered{} domínio reta mineral \textperiodcentered{} história 6 versões no jornal}

%CRISTAL {"arestas":[{"alvo":"broca_os","confianca":1.0,"epistemico":"fato","origem":"paper","rel":"parte_de"},{"alvo":"virtualizacao","confianca":0.5,"epistemico":"fato","origem":"wiki","rel":"relaciona"},{"alvo":"word","confianca":0.5,"epistemico":"fato","origem":"wiki","rel":"relaciona"},{"alvo":"vida-morte","confianca":0.5,"epistemico":"fato","origem":"wiki","rel":"relaciona"},{"alvo":"misticismo_nao_dual","confianca":0.5,"epistemico":"fato","origem":"wiki","rel":"relaciona"}],"confianca":1.0,"contraexemplos":[],"descricao":"Base algébrica do SO fractal: o estado se parte em idempotentes ortogonais eᵢ com Σeᵢ=1; cada eᵢ é um processo/célula viva.","epistemico":"fato","exemplos":[],"id":"decomposicao_peirce","memoria":"semantica","meta":{"autor":"Gentil Lopes","descoberta":true,"dominio":"matematica","fonte":"so_fractal.pdf"},"origem":"paper","palavras_chave":[],"sinonimos":[],"tipo":"conceito","titulo":"Decomposição de Peirce"}
\section{Decomposição de Peirce}
Base algébrica do SO fractal: o estado se parte em idempotentes ortogonais e\ensuremath{_{i}} com \ensuremath{\Sigma}e\ensuremath{_{i}}=1; cada e\ensuremath{_{i}} é um processo/célula viva.
\prov{relações: parte\_de broca\_os; relaciona virtualizacao; relaciona word; relaciona vida-morte; relaciona misticismo\_nao\_dual}
\prov{proveniência: paper \textperiodcentered{} so\_fractal.pdf \textperiodcentered{} fato \textperiodcentered{} confiança 1.0 \textperiodcentered{} domínio matematica \textperiodcentered{} história 7 versões no jornal}

%CRISTAL {"arestas":[{"alvo":"moeda_sl2","confianca":1.0,"epistemico":"fato","origem":"paper","rel":"parte_de"}],"confianca":1.0,"contraexemplos":[],"descricao":"δ(j,k) = (rⱼ+rₖ) − r(MⱼMₖ) ≥ 0, nulo sse j=k: desigualdade triangular estrita em H² (curvatura −1). A parte dos boosts que não se alinha vira rotação.","epistemico":"fato","exemplos":[],"id":"defeito_composicao","memoria":"semantica","meta":{"autor":"Gentil Lopes","descoberta":true,"dominio":"matematica","fonte":"paper_25_geometria_moedas.pdf, Teo.3.1"},"origem":"paper","palavras_chave":[],"sinonimos":[],"tipo":"conceito","titulo":"Defeito de Composição δ(j,k)"}
\section{Defeito de Composição \ensuremath{\delta}(j,k)}
\ensuremath{\delta}(j,k) = (r\ensuremath{_{j}}+r\ensuremath{_{k}}) \ensuremath{-} r(M\ensuremath{_{j}}M\ensuremath{_{k}}) \ensuremath{\ge} 0, nulo sse j=k: desigualdade triangular estrita em H\ensuremath{^{2}} (curvatura \ensuremath{-}1). A parte dos boosts que não se alinha vira rotação.
\prov{relações: parte\_de moeda\_sl2}
\prov{proveniência: paper \textperiodcentered{} paper\_25\_geometria\_moedas.pdf, Teo.3.1 \textperiodcentered{} fato \textperiodcentered{} confiança 1.0 \textperiodcentered{} domínio matematica \textperiodcentered{} história 2 versões no jornal}

%CRISTAL {"arestas":[{"alvo":"matematica","confianca":0.95,"epistemico":"fato","origem":"paper","rel":"parte_de"}],"confianca":1.0,"contraexemplos":[],"descricao":"Grau) Seja p(x) = axm m × axm−1 m−1 × · · · × ax2 2 × ax 1 × a0 um dupolinômio não identicamente 1. Chama-se grau de p, e representa-se por ∂p ou gr p, o número natural q tal que aq ̸= 1 e ai = 1 para todo i > q. ∂p = q ⇐⇒    aq ̸= 1; ai = 1, ∀i > q. Os dupolinômios do exemplo anterior têm grau 3, 4 e 2 respectivamente. Pelo","epistemico":"fato","exemplos":[],"id":"definicao_13_grau_seja_px_axm_m_axm1_m1_ax2_2_a","memoria":"semantica","meta":{"arquivo":"gentil/Novas_Sequencias_Aritmeticas_e_Geometric.pdf","ciencia":"teoria_dos_numeros","dominio":"matematica","nivel":4,"numero":"13","tipo_no":"paper","tipo_teorema":"definição","verificacao":"conceito novo"},"origem":"gentil/Novas_Sequencias_Aritmeticas_e_Geometric.pdf","palavras_chave":[],"sinonimos":[],"tipo":"conceito","titulo":"Definição 13: Grau) Seja p(x) = axm m × axm−1 m−1 × · · · × ax2 2 × a…"}
\section{Definição 13: Grau) Seja p(x) = axm m $\times$ axm\ensuremath{-}1 m\ensuremath{-}1 $\times$ \textperiodcentered  \textperiodcentered  \textperiodcentered  $\times$ ax2 2 $\times$ a…}
Grau) Seja p(x) = axm m $\times$ axm\ensuremath{-}1 m\ensuremath{-}1 $\times$ \textperiodcentered  \textperiodcentered  \textperiodcentered  $\times$ ax2 2 $\times$ ax 1 $\times$ a0 um dupolinômio não identicamente 1. Chama-se grau de p, e representa-se por \ensuremath{\partial}p ou gr p, o número natural q tal que aq \ensuremath{}= 1 e ai = 1 para todo i > q. \ensuremath{\partial}p = q \ensuremath{\Leftarrow}\ensuremath{\Rightarrow} \{   aq \ensuremath{}= 1; ai = 1, \ensuremath{\forall}i > q. Os dupolinômios do exemplo anterior têm grau 3, 4 e 2 respectivamente. Pelo
\prov{relações: parte\_de matematica}
\prov{proveniência: gentil/Novas\_Sequencias\_Aritmeticas\_e\_Geometric.pdf \textperiodcentered{} fato \textperiodcentered{} confiança 1.0 \textperiodcentered{} domínio matematica \textperiodcentered{} história 2 versões no jornal}

%CRISTAL {"arestas":[{"alvo":"matematica","confianca":0.95,"epistemico":"fato","origem":"paper","rel":"parte_de"}],"confianca":1.0,"contraexemplos":[],"descricao":"p 230), temos: a12 = a11 + r1 a13 = a12 + r1 a14 = a13 + r1 · · · · · · · · · · · · · · · · · · a1j = a1(j−1) + r1 Somando essas j −1 igualdades e fazendo os cancelamentos apropriados, temos: a1j = a11 + (j −1)r1 (2) Substituindo (2) e (1), temos: aij = a11 + (j −1)r1 + (i −1)r2 Com o que podemos enunciar o seguinte","epistemico":"fato","exemplos":[],"id":"definicao_16_p_230_temos_a12_a11_r1_a13_a12_r1_a14_a13","memoria":"semantica","meta":{"arquivo":"gentil/Novas_Sequencias_Aritmeticas_e_Geometric.pdf","ciencia":"teoria_dos_numeros","dominio":"matematica","nivel":4,"numero":"16","tipo_no":"paper","tipo_teorema":"definição","verificacao":"conceito novo"},"origem":"gentil/Novas_Sequencias_Aritmeticas_e_Geometric.pdf","palavras_chave":[],"sinonimos":[],"tipo":"conceito","titulo":"Definição 16: p 230), temos: a12 = a11 + r1 a13 = a12 + r1 a14 = a13 …"}
\section{Definição 16: p 230), temos: a12 = a11 + r1 a13 = a12 + r1 a14 = a13 …}
p 230), temos: a12 = a11 + r1 a13 = a12 + r1 a14 = a13 + r1 \textperiodcentered  \textperiodcentered  \textperiodcentered  \textperiodcentered  \textperiodcentered  \textperiodcentered  \textperiodcentered  \textperiodcentered  \textperiodcentered  \textperiodcentered  \textperiodcentered  \textperiodcentered  \textperiodcentered  \textperiodcentered  \textperiodcentered  \textperiodcentered  \textperiodcentered  \textperiodcentered  a1j = a1(j\ensuremath{-}1) + r1 Somando essas j \ensuremath{-}1 igualdades e fazendo os cancelamentos apropriados, temos: a1j = a11 + (j \ensuremath{-}1)r1 (2) Substituindo (2) e (1), temos: aij = a11 + (j \ensuremath{-}1)r1 + (i \ensuremath{-}1)r2 Com o que podemos enunciar o seguinte
\prov{relações: parte\_de matematica}
\prov{proveniência: gentil/Novas\_Sequencias\_Aritmeticas\_e\_Geometric.pdf \textperiodcentered{} fato \textperiodcentered{} confiança 1.0 \textperiodcentered{} domínio matematica \textperiodcentered{} história 2 versões no jornal}

%CRISTAL {"arestas":[{"alvo":"matematica","confianca":0.95,"epistemico":"fato","origem":"paper","rel":"parte_de"}],"confianca":1.0,"contraexemplos":[],"descricao":"p 230), temos: a2j = a1j + r2 a3j = a2j + r2 a4j = a3j + r2 · · · · · · · · · · · · · · · · · · aij = a(i−1)j + r2 Somando essas i −1 igualdades e fazendo os cancelamentos apropriados, temos: aij = a1j + (i −1)r2 (1) Da segunda equação da","epistemico":"fato","exemplos":[],"id":"definicao_16_p_230_temos_a2j_a1j_r2_a3j_a2j_r2_a4j_a3j","memoria":"semantica","meta":{"arquivo":"gentil/Novas_Sequencias_Aritmeticas_e_Geometric.pdf","ciencia":"teoria_dos_numeros","dominio":"matematica","nivel":4,"numero":"16","tipo_no":"paper","tipo_teorema":"definição","verificacao":"conceito novo"},"origem":"gentil/Novas_Sequencias_Aritmeticas_e_Geometric.pdf","palavras_chave":[],"sinonimos":[],"tipo":"conceito","titulo":"Definição 16: p 230), temos: a2j = a1j + r2 a3j = a2j + r2 a4j = a3j …"}
\section{Definição 16: p 230), temos: a2j = a1j + r2 a3j = a2j + r2 a4j = a3j …}
p 230), temos: a2j = a1j + r2 a3j = a2j + r2 a4j = a3j + r2 \textperiodcentered  \textperiodcentered  \textperiodcentered  \textperiodcentered  \textperiodcentered  \textperiodcentered  \textperiodcentered  \textperiodcentered  \textperiodcentered  \textperiodcentered  \textperiodcentered  \textperiodcentered  \textperiodcentered  \textperiodcentered  \textperiodcentered  \textperiodcentered  \textperiodcentered  \textperiodcentered  aij = a(i\ensuremath{-}1)j + r2 Somando essas i \ensuremath{-}1 igualdades e fazendo os cancelamentos apropriados, temos: aij = a1j + (i \ensuremath{-}1)r2 (1) Da segunda equação da
\prov{relações: parte\_de matematica}
\prov{proveniência: gentil/Novas\_Sequencias\_Aritmeticas\_e\_Geometric.pdf \textperiodcentered{} fato \textperiodcentered{} confiança 1.0 \textperiodcentered{} domínio matematica \textperiodcentered{} história 2 versões no jornal}

%CRISTAL {"arestas":[{"alvo":"matematica","confianca":0.95,"epistemico":"fato","origem":"paper","rel":"parte_de"}],"confianca":1.0,"contraexemplos":[],"descricao":"p 82), temos: ∆0 ( a f + b g )(n) = ( a f + b g )(n) = a f(n) + b g(n) = a ∆0 f(n) + b ∆0 g(n) Suponhamos a validade da fórmula para m = p: (H.I.) ∆p ( a f + b g )(n) = a ∆p f(n) + b ∆p g(n) e provemos que vale para m = p + 1. Então, pela","epistemico":"fato","exemplos":[],"id":"definicao_3_p_82_temos_0_a_f_b_g_n_a_f_b_g_n","memoria":"semantica","meta":{"arquivo":"gentil/Novas_Sequencias_Aritmeticas_e_Geometric.pdf","ciencia":"teoria_dos_numeros","dominio":"matematica","nivel":4,"numero":"3","tipo_no":"paper","tipo_teorema":"definição","verificacao":"conceito novo"},"origem":"gentil/Novas_Sequencias_Aritmeticas_e_Geometric.pdf","palavras_chave":[],"sinonimos":[],"tipo":"conceito","titulo":"Definição 3: p 82), temos: ∆0 ( a f + b g )(n) = ( a f + b g )(n) = …"}
\section{Definição 3: p 82), temos: \ensuremath{\Delta}0 ( a f + b g )(n) = ( a f + b g )(n) = …}
p 82), temos: \ensuremath{\Delta}0 ( a f + b g )(n) = ( a f + b g )(n) = a f(n) + b g(n) = a \ensuremath{\Delta}0 f(n) + b \ensuremath{\Delta}0 g(n) Suponhamos a validade da fórmula para m = p: (H.I.) \ensuremath{\Delta}p ( a f + b g )(n) = a \ensuremath{\Delta}p f(n) + b \ensuremath{\Delta}p g(n) e provemos que vale para m = p + 1. Então, pela
\prov{relações: parte\_de matematica}
\prov{proveniência: gentil/Novas\_Sequencias\_Aritmeticas\_e\_Geometric.pdf \textperiodcentered{} fato \textperiodcentered{} confiança 1.0 \textperiodcentered{} domínio matematica \textperiodcentered{} história 2 versões no jornal}

%CRISTAL {"arestas":[{"alvo":"matematica","confianca":0.95,"epistemico":"fato","origem":"paper","rel":"parte_de"}],"confianca":1.0,"contraexemplos":[],"descricao":"p 94), temos: ∇0 ( a f + b g )(n) = ( a f + b g )(n) = a f(n) + b g(n) = a ∇0 f(n) + b ∇0 g(n) Suponhamos a validade da fórmula para m = p: (H.I.) ∇p ( a f + b g )(n) = a ∇p f(n) + b ∇p g(n) e provemos que vale para m = p + 1: (T.I.) ∇p+1 ( a f + b g )(n) = a ∇p+1 f(n) + b ∇p+1 g(n) Então, ∆p+1 ( a f + b g )(n) = n−1 X i = 1 ∇p ( a f + b g )(i) + c = n−1 X i = 1 [ a ∇p f(i) + b ∇p g(i) ] + c = a n−1 X i = 1 ∇p f(i) + b n−1 X i = 1 ∇p g(i) + a · c1 + b · c2 = a \u0010 n−1 X i = 1 ∇p f(i) + c1 \u0011 + b \u0010","epistemico":"fato","exemplos":[],"id":"definicao_4_p_94_temos_0_a_f_b_g_n_a_f_b_g_n","memoria":"semantica","meta":{"arquivo":"gentil/Novas_Sequencias_Aritmeticas_e_Geometric.pdf","ciencia":"teoria_dos_numeros","dominio":"matematica","nivel":4,"numero":"4","tipo_no":"paper","tipo_teorema":"definição","verificacao":"conceito novo"},"origem":"gentil/Novas_Sequencias_Aritmeticas_e_Geometric.pdf","palavras_chave":[],"sinonimos":[],"tipo":"conceito","titulo":"Definição 4: p 94), temos: ∇0 ( a f + b g )(n) = ( a f + b g )(n) = …"}
\section{Definição 4: p 94), temos: \ensuremath{\nabla}0 ( a f + b g )(n) = ( a f + b g )(n) = …}
p 94), temos: \ensuremath{\nabla}0 ( a f + b g )(n) = ( a f + b g )(n) = a f(n) + b g(n) = a \ensuremath{\nabla}0 f(n) + b \ensuremath{\nabla}0 g(n) Suponhamos a validade da fórmula para m = p: (H.I.) \ensuremath{\nabla}p ( a f + b g )(n) = a \ensuremath{\nabla}p f(n) + b \ensuremath{\nabla}p g(n) e provemos que vale para m = p + 1: (T.I.) \ensuremath{\nabla}p+1 ( a f + b g )(n) = a \ensuremath{\nabla}p+1 f(n) + b \ensuremath{\nabla}p+1 g(n) Então, \ensuremath{\Delta}p+1 ( a f + b g )(n) = n\ensuremath{-}1 X i = 1 \ensuremath{\nabla}p ( a f + b g )(i) + c = n\ensuremath{-}1 X i = 1 [ a \ensuremath{\nabla}p f(i) + b \ensuremath{\nabla}p g(i) ] + c = a n\ensuremath{-}1 X i = 1 \ensuremath{\nabla}p f(i) + b n\ensuremath{-}1 X i = 1 \ensuremath{\nabla}p g(i) + a \textperiodcentered  c1 + b \textperiodcentered  c2 = a  n\ensuremath{-}1 X i = 1 \ensuremath{\nabla}p f(i) + c1  + b 
\prov{relações: parte\_de matematica}
\prov{proveniência: gentil/Novas\_Sequencias\_Aritmeticas\_e\_Geometric.pdf \textperiodcentered{} fato \textperiodcentered{} confiança 1.0 \textperiodcentered{} domínio matematica \textperiodcentered{} história 2 versões no jornal}

%CRISTAL {"arestas":[{"alvo":"recuperacao_evidencia_invariante","confianca":1.0,"epistemico":"fato","origem":"wiki","rel":"usa"},{"alvo":"bai_orbita_reta_metalica","confianca":1.0,"epistemico":"fato","origem":"wiki","rel":"usa"},{"alvo":"tiffany","confianca":1.0,"epistemico":"fato","origem":"wiki","rel":"parte_de"}],"confianca":1.0,"contraexemplos":[],"descricao":"conversa/parabola_algebra.py (is_definicao, alvo_definicao, compartilhado por bai e colapso): o classificador de trabalho só pegava 'o que é X / define X' — frasagens como 'me fala da conservação', 'e a prata, o que é?', 'fala do gato' caíam em conversa/saudação. Agora cobre 'o que é/são/quem é X', 'define X', '(me) fala/conta/diga de/sobre X' e o termo ANTES ('X, o que é?', 'e a X?'), 'como funciona X'/'como X funciona', 'explica X', e o FOLLOW-UP ELÍPTICO curto 'e o X?'). bai.tipo_de_trabalho usa is_definicao → órbita cristal/conhecimento; colapso._evidencia_conteudo usa alvo_definicao → a evidência de aboutness dispara e a recuperação aterra no conceito certo. Saudação e chat aberto ('me conta uma coisa qualquer') seguem fora (conversa/caos). Efeito: 'me fala da conservação'→nó da conservação; 'e a prata'→dimensão de prata; 'e o gato?'→o gato; 'como funciona o colapso?'→colapso. [D] a Tiffany entende como se pergunta de verdade — inclusive o follow-up encurtado.","epistemico":"fato","exemplos":[],"id":"definicao_frasagem_ampla","memoria":"semantica","meta":{"autor":"Lopes/broca-so","dominio":"algebra","fonte":"conversa/parabola_algebra.py (is_definicao/alvo_definicao); bai.tipo_de_trabalho; colapso._evidencia_conteudo"},"origem":"wiki","palavras_chave":[],"sinonimos":[],"tipo":"conceito","titulo":"Tiffany entende definição em frasagem ampla ('fala de X', 'X, o que é?') — a aboutness dispara"}
\section{Tiffany entende definição em frasagem ampla ('fala de X', 'X, o que é?') — a aboutness dispara}
conversa/parabola\_algebra.py (is\_definicao, alvo\_definicao, compartilhado por bai e colapso): o classificador de trabalho só pegava 'o que é X / define X' — frasagens como 'me fala da conservação', 'e a prata, o que é?', 'fala do gato' caíam em conversa/saudação. Agora cobre 'o que é/são/quem é X', 'define X', '(me) fala/conta/diga de/sobre X' e o termo ANTES ('X, o que é?', 'e a X?'), 'como funciona X'/'como X funciona', 'explica X', e o FOLLOW-UP ELÍPTICO curto 'e o X?'). bai.tipo\_de\_trabalho usa is\_definicao \ensuremath{\to} órbita cristal/conhecimento; colapso.\_evidencia\_conteudo usa alvo\_definicao \ensuremath{\to} a evidência de aboutness dispara e a recuperação aterra no conceito certo. Saudação e chat aberto ('me conta uma coisa qualquer') seguem fora (conversa/caos). Efeito: 'me fala da conservação'\ensuremath{\to}nó da conservação; 'e a prata'\ensuremath{\to}dimensão de prata; 'e o gato?'\ensuremath{\to}o gato; 'como funciona o colapso?'\ensuremath{\to}colapso. [D] a Tiffany entende como se pergunta de verdade — inclusive o follow-up encurtado.
\prov{relações: usa recuperacao\_evidencia\_invariante; usa bai\_orbita\_reta\_metalica; parte\_de tiffany}
\prov{proveniência: wiki \textperiodcentered{} conversa/parabola\_algebra.py (is\_definicao/alvo\_definicao); bai.tipo\_de\_trabalho; colapso.\_evidencia\_conteudo \textperiodcentered{} fato \textperiodcentered{} confiança 1.0 \textperiodcentered{} domínio algebra \textperiodcentered{} história 62 versões no jornal}

%CRISTAL {"arestas":[{"alvo":"teoria_dos_jogos","confianca":1.0,"epistemico":"fato","origem":"wiki","rel":"especializa"}],"confianca":1.0,"contraexemplos":[],"descricao":"A engenharia reversa da teoria dos jogos: projetar as REGRAS para que o equilíbrio produza o resultado desejado, mesmo com agentes egoístas. A 'teoria dos jogos de trás para frente' — leilões, votações.","epistemico":"inferencia","exemplos":[],"id":"desenho_de_mecanismos","memoria":"semantica","meta":{"dominio":"matematica","fonte":"teoria dos jogos"},"origem":"wiki","palavras_chave":[],"sinonimos":[],"tipo":"conceito","titulo":"Desenho de Mecanismos"}
\section{Desenho de Mecanismos}
A engenharia reversa da teoria dos jogos: projetar as REGRAS para que o equilíbrio produza o resultado desejado, mesmo com agentes egoístas. A 'teoria dos jogos de trás para frente' — leilões, votações.
\prov{relações: especializa teoria\_dos\_jogos}
\prov{proveniência: wiki \textperiodcentered{} teoria dos jogos \textperiodcentered{} inferencia \textperiodcentered{} confiança 1.0 \textperiodcentered{} domínio matematica \textperiodcentered{} história 2 versões no jornal}

%CRISTAL {"arestas":[{"alvo":"libertar_espaco","confianca":1.0,"epistemico":"fato","origem":"wiki","rel":"usa"},{"alvo":"stratum_difusao_cpu","confianca":1.0,"epistemico":"fato","origem":"wiki","rel":"explica"},{"alvo":"pacote_seed_pow","confianca":1.0,"epistemico":"fato","origem":"wiki","rel":"relaciona"}],"confianca":1.0,"contraexemplos":[],"descricao":"No loop, para cada rodada e faixa de chunk, cada lane entrega trabalho ao .so (associativo) e testa verificar_nonce_ponte — o trabalho de PoW DIFUNDE para ocupar a CPU/GPU declarada livre pela frac. É como a máquina enche o espaço liberado pelo PoW-cache por job.","epistemico":"inferencia","exemplos":[],"id":"difusao_preenche_mineracao","memoria":"semantica","meta":{"autor":"broca-so","dominio":"fractal","fonte":"neuronio/stratum_trabalho.py:_loop_broca_primitivas"},"origem":"wiki","palavras_chave":[],"sinonimos":[],"tipo":"conceito","titulo":"A Difusão É o Preenchimento Minerador"}
\section{A Difusão É o Preenchimento Minerador}
No loop, para cada rodada e faixa de chunk, cada lane entrega trabalho ao .so (associativo) e testa verificar\_nonce\_ponte — o trabalho de PoW DIFUNDE para ocupar a CPU/GPU declarada livre pela frac. É como a máquina enche o espaço liberado pelo PoW-cache por job.
\prov{relações: usa libertar\_espaco; explica stratum\_difusao\_cpu; relaciona pacote\_seed\_pow}
\prov{proveniência: wiki \textperiodcentered{} neuronio/stratum\_trabalho.py:\_loop\_broca\_primitivas \textperiodcentered{} inferencia \textperiodcentered{} confiança 1.0 \textperiodcentered{} domínio fractal \textperiodcentered{} história 4 versões no jornal}

%CRISTAL {"arestas":[{"alvo":"estrategia_dominante","confianca":1.0,"epistemico":"fato","origem":"wiki","rel":"usa"},{"alvo":"tragedia_dos_comuns","confianca":1.0,"epistemico":"fato","origem":"wiki","rel":"explica"}],"confianca":1.0,"contraexemplos":[],"descricao":"Dois presos: delatar é dominante para cada um, mas ambos calarem seria melhor para os dois. O paradigma de como a racionalidade individual leva ao pior coletivo — cooperação vs deserção.","epistemico":"fato","exemplos":[],"id":"dilema_do_prisioneiro","memoria":"semantica","meta":{"dominio":"matematica","fonte":"teoria dos jogos"},"origem":"wiki","palavras_chave":[],"sinonimos":[],"tipo":"conceito","titulo":"Dilema do Prisioneiro"}
\section{Dilema do Prisioneiro}
Dois presos: delatar é dominante para cada um, mas ambos calarem seria melhor para os dois. O paradigma de como a racionalidade individual leva ao pior coletivo — cooperação vs deserção.
\prov{relações: usa estrategia\_dominante; explica tragedia\_dos\_comuns}
\prov{proveniência: wiki \textperiodcentered{} teoria dos jogos \textperiodcentered{} fato \textperiodcentered{} confiança 1.0 \textperiodcentered{} domínio matematica \textperiodcentered{} história 3 versões no jornal}

%CRISTAL {"arestas":[{"alvo":"computacao_fractal","confianca":1.0,"epistemico":"fato","origem":"wiki","rel":"parte_de"},{"alvo":"subshift_do_gap","confianca":1.0,"epistemico":"fato","origem":"wiki","rel":"usa"}],"confianca":1.0,"contraexemplos":[],"descricao":"O invariante computacional D_tr(M)=lim log₂|L_n|/n (L_n = palavras de comprimento n nos traços). Para o gap k vale D_tr=log₂λ_k (λ_k = raio de Perron de A_k), coincidindo por três métodos (contagem, Perron, Moran) a erro de máquina: k=1→0.694242, k=2→0.879146, k=3→0.946777, k=4→0.975225.","epistemico":"fato","exemplos":[],"id":"dimensao_de_traco","memoria":"semantica","meta":{"autor":"broca-so","dominio":"fractal","fonte":"machine/fractal/fractal.py:dim_traco; computador_fractal.tex"},"origem":"wiki","palavras_chave":[],"sinonimos":[],"tipo":"conceito","titulo":"Dimensão de Traço D_tr = log₂λ_k"}
\section{Dimensão de Traço D\_tr = log\ensuremath{_{2}}\ensuremath{\lambda}\_k}
O invariante computacional D\_tr(M)=lim log\ensuremath{_{2}}|L\_n|/n (L\_n = palavras de comprimento n nos traços). Para o gap k vale D\_tr=log\ensuremath{_{2}}\ensuremath{\lambda}\_k (\ensuremath{\lambda}\_k = raio de Perron de A\_k), coincidindo por três métodos (contagem, Perron, Moran) a erro de máquina: k=1\ensuremath{\to}0.694242, k=2\ensuremath{\to}0.879146, k=3\ensuremath{\to}0.946777, k=4\ensuremath{\to}0.975225.
\prov{relações: parte\_de computacao\_fractal; usa subshift\_do\_gap}
\prov{proveniência: wiki \textperiodcentered{} machine/fractal/fractal.py:dim\_traco; computador\_fractal.tex \textperiodcentered{} fato \textperiodcentered{} confiança 1.0 \textperiodcentered{} domínio fractal \textperiodcentered{} história 3 versões no jornal}

%CRISTAL {"arestas":[{"alvo":"entropia_topologica","confianca":1.0,"epistemico":"fato","origem":"paper","rel":"especializa"},{"alvo":"broca_os","confianca":1.0,"epistemico":"fato","origem":"paper","rel":"referencia"}],"confianca":1.0,"contraexemplos":[],"descricao":"Invariante da máquina fractal D_U=log₂λ (λ de Perron do gap): mede quantas palavras de tamanho n aparecem nos traços; distingue máquina fractal de digital.","epistemico":"fato","exemplos":[],"id":"dimensao_traco","memoria":"semantica","meta":{"autor":"Gentil Lopes","descoberta":true,"dominio":"matematica","fonte":"computador_fractal.pdf"},"origem":"paper","palavras_chave":[],"sinonimos":[],"tipo":"conceito","titulo":"Dimensão de Traço (D_U)"}
\section{Dimensão de Traço (D\_U)}
Invariante da máquina fractal D\_U=log\ensuremath{_{2}}\ensuremath{\lambda} (\ensuremath{\lambda} de Perron do gap): mede quantas palavras de tamanho n aparecem nos traços; distingue máquina fractal de digital.
\prov{relações: especializa entropia\_topologica; referencia broca\_os}
\prov{proveniência: paper \textperiodcentered{} computador\_fractal.pdf \textperiodcentered{} fato \textperiodcentered{} confiança 1.0 \textperiodcentered{} domínio matematica \textperiodcentered{} história 3 versões no jornal}

%CRISTAL {"arestas":[{"alvo":"lei_da_parabola_e_dinamica_de_anosov","confianca":1.0,"epistemico":"fato","origem":"wiki","rel":"parte_de"},{"alvo":"corpo_estelar","confianca":1.0,"epistemico":"fato","origem":"wiki","rel":"usa"},{"alvo":"hipotese_de_riemann","confianca":1.0,"epistemico":"fato","origem":"wiki","rel":"relaciona"}],"confianca":1.0,"contraexemplos":[],"descricao":"Uma decisão de linguagem, às claras: NADA NOVO, só linguagem equivalente. Três rótulos SINÔNIMOS e intercambiáveis: 'dinâmica de ANOSOV' (o termo consagrado, a hiperbolicidade dinâmica de 1967), 'dinâmica de SELBERG' (o sinônimo que honra o primeiro passo espectral, o traço de 1956) e 'dinâmica de SELBERG-ANOSOV' (o COMPOSTO que representa a costura histórica que o projeto faz: 1956→1967, o hífen sendo a própria ponte espectro↔geodésicas). Não é conceito novo nem renomeação — os resultados são a dinâmica hiperbólica de sempre; os nomes extras são rótulos equivalentes oferecidos à comunidade. E a INTENÇÃO é explícita: NÃO é teoria de tudo — é dar OPERAÇÃO EFETIVA ao corpo estelar.","epistemico":"inferencia","exemplos":[],"id":"dinamica_de_selberg_o_termo","memoria":"semantica","meta":{"dominio":"matematica","fonte":"dinâmica de Anosov / Selberg"},"origem":"wiki","palavras_chave":[],"sinonimos":[],"tipo":"conceito","titulo":"Dinâmica de Selberg-Anosov (sinônimos, a costura histórica)"}
\section{Dinâmica de Selberg-Anosov (sinônimos, a costura histórica)}
Uma decisão de linguagem, às claras: NADA NOVO, só linguagem equivalente. Três rótulos SINÔNIMOS e intercambiáveis: 'dinâmica de ANOSOV' (o termo consagrado, a hiperbolicidade dinâmica de 1967), 'dinâmica de SELBERG' (o sinônimo que honra o primeiro passo espectral, o traço de 1956) e 'dinâmica de SELBERG-ANOSOV' (o COMPOSTO que representa a costura histórica que o projeto faz: 1956\ensuremath{\to}1967, o hífen sendo a própria ponte espectro\ensuremath{\leftrightarrow}geodésicas). Não é conceito novo nem renomeação — os resultados são a dinâmica hiperbólica de sempre; os nomes extras são rótulos equivalentes oferecidos à comunidade. E a INTENÇÃO é explícita: NÃO é teoria de tudo — é dar OPERAÇÃO EFETIVA ao corpo estelar.
\prov{relações: parte\_de lei\_da\_parabola\_e\_dinamica\_de\_anosov; usa corpo\_estelar; relaciona hipotese\_de\_riemann}
\prov{proveniência: wiki \textperiodcentered{} dinâmica de Anosov / Selberg \textperiodcentered{} inferencia \textperiodcentered{} confiança 1.0 \textperiodcentered{} domínio matematica \textperiodcentered{} história 12 versões no jornal}

%CRISTAL {"arestas":[{"alvo":"numeros_duais","confianca":1.0,"epistemico":"fato","origem":"paper","rel":"usa"},{"alvo":"word","confianca":0.5,"epistemico":"fato","origem":"wiki","rel":"relaciona"}],"confianca":1.0,"contraexemplos":[],"descricao":"O sinal ε=sgn(−disc) da álgebra classifica a máquina: ε=+1 contínua (ℂ, RLC amortecido), ε=−1 discreta (split, lógica 𝔽₂), ε=0 nilpotente (dual).","epistemico":"fato","exemplos":[],"id":"discriminante_faces","memoria":"semantica","meta":{"autor":"Gentil Lopes","descoberta":false,"dominio":"matematica","fonte":"eletronica.pdf"},"origem":"paper","palavras_chave":[],"sinonimos":[],"tipo":"conceito","titulo":"Discriminante ε (três faces)"}
\section{Discriminante \ensuremath{\varepsilon} (três faces)}
O sinal \ensuremath{\varepsilon}=sgn(\ensuremath{-}disc) da álgebra classifica a máquina: \ensuremath{\varepsilon}=+1 contínua (\ensuremath{\mathbb{C}}, RLC amortecido), \ensuremath{\varepsilon}=\ensuremath{-}1 discreta (split, lógica \ensuremath{\mathbb{F}}\ensuremath{_{2}}), \ensuremath{\varepsilon}=0 nilpotente (dual).
\prov{relações: usa numeros\_duais; relaciona word}
\prov{proveniência: paper \textperiodcentered{} eletronica.pdf \textperiodcentered{} fato \textperiodcentered{} confiança 1.0 \textperiodcentered{} domínio matematica \textperiodcentered{} história 3 versões no jornal}

%CRISTAL {"arestas":[{"alvo":"crescimento_vorticidade","confianca":1.0,"epistemico":"fato","origem":"paper","rel":"usa"},{"alvo":"indecidibilidade_ns_psl","confianca":1.0,"epistemico":"fato","origem":"paper","rel":"relaciona"}],"confianca":1.0,"contraexemplos":[],"descricao":"Conjectura: no limite combinado (N_s→∞, ν₀→0) a taxa de dissipação ε~ν·|ω|² permanece finita mesmo com ν→0 — perda de regularidade sem quebra de conservação de energia, como Onsager 1949.","epistemico":"hipotese","exemplos":[],"id":"dissipacao_anomala_onsager","memoria":"semantica","meta":{"autor":"Gentil Lopes","descoberta":false,"dominio":"matematica","fonte":"paper_BG.pdf, §3"},"origem":"paper","palavras_chave":[],"sinonimos":[],"tipo":"conceito","titulo":"Dissipação Anômala de Onsager"}
\section{Dissipação Anômala de Onsager}
Conjectura: no limite combinado (N\_s\ensuremath{\to}\ensuremath{\infty}, \ensuremath{\nu}\ensuremath{_{0}}\ensuremath{\to}0) a taxa de dissipação \ensuremath{\varepsilon}\~{}\ensuremath{\nu}\textperiodcentered |\ensuremath{\omega}|\ensuremath{^{2}} permanece finita mesmo com \ensuremath{\nu}\ensuremath{\to}0 — perda de regularidade sem quebra de conservação de energia, como Onsager 1949.
\prov{relações: usa crescimento\_vorticidade; relaciona indecidibilidade\_ns\_psl}
\prov{proveniência: paper \textperiodcentered{} paper\_BG.pdf, §3 \textperiodcentered{} hipotese \textperiodcentered{} confiança 1.0 \textperiodcentered{} domínio matematica \textperiodcentered{} história 3 versões no jornal}

%CRISTAL {"arestas":[{"alvo":"floco_phi","confianca":1.0,"epistemico":"fato","origem":"wiki","rel":"usa"},{"alvo":"reta_de_ouro","confianca":1.0,"epistemico":"fato","origem":"wiki","rel":"relaciona"},{"alvo":"ruido_difusao_word","confianca":1.0,"epistemico":"fato","origem":"wiki","rel":"relaciona"}],"confianca":1.0,"contraexemplos":[],"descricao":"A projeção associativa 𝓛(ω)=(s,C₃,C_φ) perde informação (muitos w partilham o mesmo (s,C₃,C_φ)); a parcela perdida é exatamente Φ, com ω↦(𝓛(ω),Φ(ω)) e Φ≉𝓛. A energia que não fecha a ressonância a+c²=1 'cristaliza' como calor dissipado numa passagem irreversível, mas armazenável.","epistemico":"fato","exemplos":[],"id":"dissipacao_na_passagem","memoria":"semantica","meta":{"autor":"broca-so","dominio":"fractal","fonte":"papers/floco_de_neve.tex §II (Prop. Dissipação)"},"origem":"wiki","palavras_chave":[],"sinonimos":[],"tipo":"conceito","titulo":"A Reta Perde, o Floco é o Resto"}
\section{A Reta Perde, o Floco é o Resto}
A projeção associativa \ensuremath{\mathcal{L}}(\ensuremath{\omega})=(s,C\ensuremath{_{3}},C\_\ensuremath{\varphi}) perde informação (muitos w partilham o mesmo (s,C\ensuremath{_{3}},C\_\ensuremath{\varphi})); a parcela perdida é exatamente \ensuremath{\Phi}, com \ensuremath{\omega}\ensuremath{\mapsto}(\ensuremath{\mathcal{L}}(\ensuremath{\omega}),\ensuremath{\Phi}(\ensuremath{\omega})) e \ensuremath{\Phi}\ensuremath{\not\approx}\ensuremath{\mathcal{L}}. A energia que não fecha a ressonância a+c\ensuremath{^{2}}=1 'cristaliza' como calor dissipado numa passagem irreversível, mas armazenável.
\prov{relações: usa floco\_phi; relaciona reta\_de\_ouro; relaciona ruido\_difusao\_word}
\prov{proveniência: wiki \textperiodcentered{} papers/floco\_de\_neve.tex §II (Prop. Dissipação) \textperiodcentered{} fato \textperiodcentered{} confiança 1.0 \textperiodcentered{} domínio fractal \textperiodcentered{} história 4 versões no jornal}

%CRISTAL {"arestas":[{"alvo":"reta_metalica_geral","confianca":1.0,"epistemico":"fato","origem":"wiki","rel":"explica"},{"alvo":"caixa_binaria_kbonacci","confianca":1.0,"epistemico":"fato","origem":"wiki","rel":"explica"},{"alvo":"beta_numeracao_metalica","confianca":1.0,"epistemico":"fato","origem":"wiki","rel":"explica"}],"confianca":1.0,"contraexemplos":[],"descricao":"NÃO confundir os dois papéis de λ ao generalizar. (i) Na RETA METÁLICA, σₙ é a BASE da σₙ-expansão (alfabeto 0..n) e a entropia é ln σₙ (ouro 0,481, prata 0,881). (ii) Na CAIXA FRACTAL BINÁRIA, λ é o Perron do subshift binário podado e a dimensão é log₂ λ (ouro 0,694). COINCIDEM só em n=1 (σ₁=φ = golden mean shift = 'proibir 11'). Da prata em diante DIVERGEM: a caixa da prata é a σ₂-expansão (λ=σ₂=1+√2, alfabeto 0,1,2), NÃO 'proibir 1²' (que daria k-bonacci, não metálico). Ponto matemático que sustenta a generalização correta.","epistemico":"fato","exemplos":[],"id":"distincao_base_vs_subshift","memoria":"semantica","meta":{"autor":"Lopes/broca-so","dominio":"fractal","fonte":"doc/fractais/caixas_fractais_metalicas.py"},"origem":"wiki","palavras_chave":[],"sinonimos":[],"tipo":"conceito","titulo":"A distinção load-bearing: base σ_n vs Perron do subshift"}
\section{A distinção load-bearing: base \ensuremath{\sigma}\_n vs Perron do subshift}
NÃO confundir os dois papéis de \ensuremath{\lambda} ao generalizar. (i) Na RETA METÁLICA, \ensuremath{\sigma}\ensuremath{_{n}} é a BASE da \ensuremath{\sigma}\ensuremath{_{n}}-expansão (alfabeto 0..n) e a entropia é ln \ensuremath{\sigma}\ensuremath{_{n}} (ouro 0,481, prata 0,881). (ii) Na CAIXA FRACTAL BINÁRIA, \ensuremath{\lambda} é o Perron do subshift binário podado e a dimensão é log\ensuremath{_{2}} \ensuremath{\lambda} (ouro 0,694). COINCIDEM só em n=1 (\ensuremath{\sigma}\ensuremath{_{1}}=\ensuremath{\varphi} = golden mean shift = 'proibir 11'). Da prata em diante DIVERGEM: a caixa da prata é a \ensuremath{\sigma}\ensuremath{_{2}}-expansão (\ensuremath{\lambda}=\ensuremath{\sigma}\ensuremath{_{2}}=1+\ensuremath{\sqrt{\,}}2, alfabeto 0,1,2), NÃO 'proibir 1\ensuremath{^{2}}' (que daria k-bonacci, não metálico). Ponto matemático que sustenta a generalização correta.
\prov{relações: explica reta\_metalica\_geral; explica caixa\_binaria\_kbonacci; explica beta\_numeracao\_metalica}
\prov{proveniência: wiki \textperiodcentered{} doc/fractais/caixas\_fractais\_metalicas.py \textperiodcentered{} fato \textperiodcentered{} confiança 1.0 \textperiodcentered{} domínio fractal \textperiodcentered{} história 5 versões no jornal}

%CRISTAL {"arestas":[{"alvo":"dissipacao_na_passagem","confianca":1.0,"epistemico":"fato","origem":"wiki","rel":"usa"},{"alvo":"ciclo_i_canonico_floco_reta","confianca":1.0,"epistemico":"fato","origem":"wiki","rel":"explica"}],"confianca":1.0,"contraexemplos":[],"descricao":"Ciclo I: canónico ω → Dissip → (Φ + reta 𝓛). Ciclo II: (reta 𝓛 + Φ) → Rec → canónico. O floco guarda o que 𝓛 não carrega; sem floco não há volta. Roundtrip verificado 32/32 no pipeline parabola_auditar→parabola_para_reta→koch_coord.","epistemico":"fato","exemplos":[],"id":"dois_ciclos_reconstrucao","memoria":"semantica","meta":{"autor":"broca-so","dominio":"fractal","fonte":"papers/floco_de_neve.tex §V,§VI"},"origem":"wiki","palavras_chave":[],"sinonimos":[],"tipo":"conceito","titulo":"Os Dois Ciclos de Reconstrução"}
\section{Os Dois Ciclos de Reconstrução}
Ciclo I: canónico \ensuremath{\omega} \ensuremath{\to} Dissip \ensuremath{\to} (\ensuremath{\Phi} + reta \ensuremath{\mathcal{L}}). Ciclo II: (reta \ensuremath{\mathcal{L}} + \ensuremath{\Phi}) \ensuremath{\to} Rec \ensuremath{\to} canónico. O floco guarda o que \ensuremath{\mathcal{L}} não carrega; sem floco não há volta. Roundtrip verificado 32/32 no pipeline parabola\_auditar\ensuremath{\to}parabola\_para\_reta\ensuremath{\to}koch\_coord.
\prov{relações: usa dissipacao\_na\_passagem; explica ciclo\_i\_canonico\_floco\_reta}
\prov{proveniência: wiki \textperiodcentered{} papers/floco\_de\_neve.tex §V,§VI \textperiodcentered{} fato \textperiodcentered{} confiança 1.0 \textperiodcentered{} domínio fractal \textperiodcentered{} história 3 versões no jornal}

%CRISTAL {"arestas":[{"alvo":"octonios_hiperbolicos","confianca":1.0,"epistemico":"fato","origem":"paper","rel":"relaciona"},{"alvo":"associador","confianca":1.0,"epistemico":"fato","origem":"paper","rel":"usa"},{"alvo":"qubit_hopf","confianca":1.0,"epistemico":"fato","origem":"paper","rel":"especializa"}],"confianca":1.0,"contraexemplos":[],"descricao":"Dois qubits ℍ⊗ℍ vivem em ℝ⁸=𝕆; o emaranhamento é codificado na não-associatividade (associador) dos octônios.","epistemico":"fato","exemplos":[],"id":"dois_qubits_octonios","memoria":"semantica","meta":{"autor":"Gentil Lopes","descoberta":false,"dominio":"matematica","fonte":"paper_3_quantica.pdf, Obs.3.3"},"origem":"paper","palavras_chave":[],"sinonimos":[],"tipo":"conceito","titulo":"Dois Qubits = Octônios"}
\section{Dois Qubits = Octônios}
Dois qubits \ensuremath{\mathbb{H}}\ensuremath{\otimes}\ensuremath{\mathbb{H}} vivem em \ensuremath{\mathbb{R}}\ensuremath{^{8}}=\ensuremath{\mathbb{O}}; o emaranhamento é codificado na não-associatividade (associador) dos octônios.
\prov{relações: relaciona octonios\_hiperbolicos; usa associador; especializa qubit\_hopf}
\prov{proveniência: paper \textperiodcentered{} paper\_3\_quantica.pdf, Obs.3.3 \textperiodcentered{} fato \textperiodcentered{} confiança 1.0 \textperiodcentered{} domínio matematica \textperiodcentered{} história 4 versões no jornal}

%CRISTAL {"arestas":[{"alvo":"motor_iterado","confianca":1.0,"epistemico":"fato","origem":"paper","rel":"usa"},{"alvo":"word","confianca":0.5,"epistemico":"fato","origem":"wiki","rel":"relaciona"}],"confianca":1.0,"contraexemplos":[],"descricao":"Φ(c)=eG(c)⁺iθ(c) uniformiza o exterior do Mandelbrot; ds²=dG²+dθ² é hiperbólica, com os raios externos θ=const como geodésicas e dinâmica θ↦2θ.","epistemico":"fato","exemplos":[],"id":"douady_hubbard","memoria":"semantica","meta":{"autor":"Gentil Lopes","descoberta":false,"dominio":"matematica","fonte":"paper_19_mobius_fractal.pdf"},"origem":"paper","palavras_chave":[],"sinonimos":[],"tipo":"conceito","titulo":"Métrica de Douady-Hubbard"}
\section{Métrica de Douady-Hubbard}
\ensuremath{\Phi}(c)=eG(c)\ensuremath{^{+}}i\ensuremath{\theta}(c) uniformiza o exterior do Mandelbrot; ds\ensuremath{^{2}}=dG\ensuremath{^{2}}+d\ensuremath{\theta}\ensuremath{^{2}} é hiperbólica, com os raios externos \ensuremath{\theta}=const como geodésicas e dinâmica \ensuremath{\theta}\ensuremath{\mapsto}2\ensuremath{\theta}.
\prov{relações: usa motor\_iterado; relaciona word}
\prov{proveniência: paper \textperiodcentered{} paper\_19\_mobius\_fractal.pdf \textperiodcentered{} fato \textperiodcentered{} confiança 1.0 \textperiodcentered{} domínio matematica \textperiodcentered{} história 4 versões no jornal}

%CRISTAL {"arestas":[{"alvo":"ilha_fractal_3d_dim","confianca":1.0,"epistemico":"fato","origem":"wiki","rel":"usa"},{"alvo":"fractal_phi_fibrado","confianca":1.0,"epistemico":"fato","origem":"wiki","rel":"relaciona"}],"confianca":1.0,"contraexemplos":[],"descricao":"As dimensões generalizadas D_q da ilha (Hadamard, r=−1) crescem de D₋₂=0.899, D₀=0.940, D₂=0.968, até D₅=0.984 — todas próximas de 1. E são invariantes à porta unitária: Pauli-X, Pauli-Y e Hadamard dão espectros numericamente idênticos (a dimensão da ilha independe da porta).","epistemico":"fato","exemplos":[],"id":"dq_espectro_ilha","memoria":"semantica","meta":{"autor":"broca-so","dominio":"fractal","fonte":"hiper/benchmarks/fase35_multifractal_resultados.json"},"origem":"wiki","palavras_chave":[],"sinonimos":[],"tipo":"conceito","titulo":"Espectro D_q da Ilha (0.899→0.984)"}
\section{Espectro D\_q da Ilha (0.899\ensuremath{\to}0.984)}
As dimensões generalizadas D\_q da ilha (Hadamard, r=\ensuremath{-}1) crescem de D\ensuremath{_{-2}}=0.899, D\ensuremath{_{0}}=0.940, D\ensuremath{_{2}}=0.968, até D\ensuremath{_{5}}=0.984 — todas próximas de 1. E são invariantes à porta unitária: Pauli-X, Pauli-Y e Hadamard dão espectros numericamente idênticos (a dimensão da ilha independe da porta).
\prov{relações: usa ilha\_fractal\_3d\_dim; relaciona fractal\_phi\_fibrado}
\prov{proveniência: wiki \textperiodcentered{} hiper/benchmarks/fase35\_multifractal\_resultados.json \textperiodcentered{} fato \textperiodcentered{} confiança 1.0 \textperiodcentered{} domínio fractal \textperiodcentered{} história 3 versões no jornal}

%CRISTAL {"arestas":[{"alvo":"sigma_equivalencia_crown","confianca":1.0,"epistemico":"fato","origem":"wiki","rel":"usa"},{"alvo":"tensor_cayley","confianca":1.0,"epistemico":"fato","origem":"wiki","rel":"usa"},{"alvo":"constante_168","confianca":1.0,"epistemico":"fato","origem":"wiki","rel":"explica"},{"alvo":"psl_embeds_g2","confianca":1.0,"epistemico":"fato","origem":"wiki","rel":"usa"},{"alvo":"catedral_sigma_aritmetica","confianca":1.0,"epistemico":"fato","origem":"wiki","rel":"relaciona"}],"confianca":1.0,"contraexemplos":[],"descricao":"O mesmo 168 emerge por dois caminhos independentes: a Catedral o alcança por σ-aritmética (o plano de Fano dá PSL(2,7)=168, e σ(168)=480 o iguala ao E8); o programa Lopes o obtém como ‖φ‖² (a norma do tensor de Cayley) e via PSL(2,7)⊂G₂=Aut(𝕆). Duas derivações, uma constante — o núcleo que ambos os edifícios cercam por dentro e por fora.","epistemico":"inferencia","exemplos":[],"id":"duas_leituras_168","memoria":"semantica","meta":{"dominio":"matematica","fonte":"fusão Catedral↔Lopes"},"origem":"wiki","palavras_chave":[],"sinonimos":[],"tipo":"conceito","titulo":"As duas leituras do 168 (σ-divisor e norma octônica)"}
\section{As duas leituras do 168 (\ensuremath{\sigma}-divisor e norma octônica)}
O mesmo 168 emerge por dois caminhos independentes: a Catedral o alcança por \ensuremath{\sigma}-aritmética (o plano de Fano dá PSL(2,7)=168, e \ensuremath{\sigma}(168)=480 o iguala ao E8); o programa Lopes o obtém como \ensuremath{\|}\ensuremath{\varphi}\ensuremath{\|}\ensuremath{^{2}} (a norma do tensor de Cayley) e via PSL(2,7)\ensuremath{\subset}G\ensuremath{_{2}}=Aut(\ensuremath{\mathbb{O}}). Duas derivações, uma constante — o núcleo que ambos os edifícios cercam por dentro e por fora.
\prov{relações: usa sigma\_equivalencia\_crown; usa tensor\_cayley; explica constante\_168; usa psl\_embeds\_g2; relaciona catedral\_sigma\_aritmetica}
\prov{proveniência: wiki \textperiodcentered{} fusão Catedral\ensuremath{\leftrightarrow}Lopes \textperiodcentered{} inferencia \textperiodcentered{} confiança 1.0 \textperiodcentered{} domínio matematica \textperiodcentered{} história 6 versões no jornal}

%CRISTAL {"arestas":[{"alvo":"paley_169_e8","confianca":1.0,"epistemico":"fato","origem":"wiki","rel":"usa"},{"alvo":"grupos_excepcionais_jordan","confianca":1.0,"epistemico":"fato","origem":"wiki","rel":"usa"},{"alvo":"sigma_equivalencia_crown","confianca":1.0,"epistemico":"fato","origem":"wiki","rel":"usa"},{"alvo":"duas_leituras_168","confianca":1.0,"epistemico":"fato","origem":"wiki","rel":"relaciona"},{"alvo":"catedral_sigma_aritmetica","confianca":1.0,"epistemico":"fato","origem":"wiki","rel":"relaciona"}],"confianca":1.0,"contraexemplos":[],"descricao":"A dimensão 248 do E8 também tem dois acessos: a Catedral a destila por σ²(169)=248 (do 'Observador' 13² em dois passos de divisores); o programa Lopes a situa na família excepcional (E6/E7/E8, Jordan J₃(𝕆)/Albert). A σ-equivalência σ(168)=σ(248)=480 amarra as duas constantes numa balança só.","epistemico":"inferencia","exemplos":[],"id":"duas_leituras_e8_248","memoria":"semantica","meta":{"dominio":"matematica","fonte":"fusão Catedral↔Lopes"},"origem":"wiki","palavras_chave":[],"sinonimos":[],"tipo":"conceito","titulo":"As duas leituras do E8=248 (σ²(169) e grupos excepcionais)"}
\section{As duas leituras do E8=248 (\ensuremath{\sigma}\ensuremath{^{2}}(169) e grupos excepcionais)}
A dimensão 248 do E8 também tem dois acessos: a Catedral a destila por \ensuremath{\sigma}\ensuremath{^{2}}(169)=248 (do 'Observador' 13\ensuremath{^{2}} em dois passos de divisores); o programa Lopes a situa na família excepcional (E6/E7/E8, Jordan J\ensuremath{_{3}}(\ensuremath{\mathbb{O}})/Albert). A \ensuremath{\sigma}-equivalência \ensuremath{\sigma}(168)=\ensuremath{\sigma}(248)=480 amarra as duas constantes numa balança só.
\prov{relações: usa paley\_169\_e8; usa grupos\_excepcionais\_jordan; usa sigma\_equivalencia\_crown; relaciona duas\_leituras\_168; relaciona catedral\_sigma\_aritmetica}
\prov{proveniência: wiki \textperiodcentered{} fusão Catedral\ensuremath{\leftrightarrow}Lopes \textperiodcentered{} inferencia \textperiodcentered{} confiança 1.0 \textperiodcentered{} domínio matematica \textperiodcentered{} história 6 versões no jornal}

%CRISTAL {"arestas":[{"alvo":"gentil_definicao_13_grau_seja_px_axm_m_axm1_m1_ax2_2_ax_1","confianca":1.0,"epistemico":"fato","origem":"wiki","rel":"usa"},{"alvo":"quantizacao_tropical","confianca":1.0,"epistemico":"fato","origem":"wiki","rel":"usa"},{"alvo":"corpo_dual","confianca":1.0,"epistemico":"fato","origem":"wiki","rel":"usa"},{"alvo":"caixas_fractais_generalizacao","confianca":1.0,"epistemico":"fato","origem":"wiki","rel":"relaciona"},{"alvo":"teorema_esterilidade_dimensao","confianca":1.0,"epistemico":"fato","origem":"wiki","rel":"usa"},{"alvo":"cofre_parabola_hero","confianca":1.0,"epistemico":"fato","origem":"wiki","rel":"usa"},{"alvo":"tiffany","confianca":1.0,"epistemico":"fato","origem":"wiki","rel":"parte_de"},{"alvo":"transformada_metalica","confianca":1.0,"epistemico":"fato","origem":"wiki","rel":"parte_de"}],"confianca":1.0,"contraexemplos":[],"descricao":"Um DUPOLINÔMIO de Lopes (Def. 13) é o polinômio na forma MULTIPLICATIVA: p(x)=a·xᵐ ⊗ … ⊗ a₀, identidade 1 (não 0), grau = maior q com a_q≠1. É o dual (log/exp) do polinômio aditivo (soma de monômios ↔ produto de termos-potência). A álgebra posicional É esse dual: a POSIÇÃO é o EXPOENTE (aditiva — somar posições ↔ compor a companheira), o VALOR é o produto (o dupolinômio); log/exp faz a ponte e, no limite T→0, é a álgebra TROPICAL (quantizacao_tropical/corpo_dual) — a mesma parábola a+c²=1. Logo o polinômio característico de uma reta, lido multiplicativamente, É um dupolinômio, e a reserva vive no anel dos dupolinômios. META-INDUÇÃO (o limite de verificação, como nas caixas fractais): verificado grau 2 (metais) e grau 3 (plástico); daí indui-se o padrão a TODO grau/dupolinômio — não se reivindica verificação além do grau 3, mas a INDUÇÃO do padrão, que é a própria estrutura dos dupolinômios (o expoente é a posição). [I] honesto.","epistemico":"inferencia","exemplos":[],"id":"dupolinomio_meta_inducao","memoria":"semantica","meta":{"autor":"Lopes/broca-so","dominio":"algebra","fonte":"gentil Def. 13 (dupolinômio); doc/algebra_prata/paper §10; caixas_fractais (meta-indução); quantizacao_tropical"},"origem":"wiki","palavras_chave":[],"sinonimos":[],"tipo":"conceito","titulo":"Meta-indução sobre os dupolinômios de Lopes: o característico é um dupolinômio (posição=expoente)"}
\section{Meta-indução sobre os dupolinômios de Lopes: o característico é um dupolinômio (posição=expoente)}
Um DUPOLINÔMIO de Lopes (Def. 13) é o polinômio na forma MULTIPLICATIVA: p(x)=a\textperiodcentered x\ensuremath{^{m}} \ensuremath{\otimes} … \ensuremath{\otimes} a\ensuremath{_{0}}, identidade 1 (não 0), grau = maior q com a\_q\ensuremath{\ne}1. É o dual (log/exp) do polinômio aditivo (soma de monômios \ensuremath{\leftrightarrow} produto de termos-potência). A álgebra posicional É esse dual: a POSIÇÃO é o EXPOENTE (aditiva — somar posições \ensuremath{\leftrightarrow} compor a companheira), o VALOR é o produto (o dupolinômio); log/exp faz a ponte e, no limite T\ensuremath{\to}0, é a álgebra TROPICAL (quantizacao\_tropical/corpo\_dual) — a mesma parábola a+c\ensuremath{^{2}}=1. Logo o polinômio característico de uma reta, lido multiplicativamente, É um dupolinômio, e a reserva vive no anel dos dupolinômios. META-INDUÇÃO (o limite de verificação, como nas caixas fractais): verificado grau 2 (metais) e grau 3 (plástico); daí indui-se o padrão a TODO grau/dupolinômio — não se reivindica verificação além do grau 3, mas a INDUÇÃO do padrão, que é a própria estrutura dos dupolinômios (o expoente é a posição). [I] honesto.
\prov{relações: usa gentil\_definicao\_13\_grau\_seja\_px\_axm\_m\_axm1\_m1\_ax2\_2\_ax\_1; usa quantizacao\_tropical; usa corpo\_dual; relaciona caixas\_fractais\_generalizacao; usa teorema\_esterilidade\_dimensao; usa cofre\_parabola\_hero; parte\_de tiffany; parte\_de transformada\_metalica}
\prov{proveniência: wiki \textperiodcentered{} gentil Def. 13 (dupolinômio); doc/algebra\_prata/paper §10; caixas\_fractais (meta-indução); quantizacao\_tropical \textperiodcentered{} inferencia \textperiodcentered{} confiança 1.0 \textperiodcentered{} domínio algebra \textperiodcentered{} história 9 versões no jornal}

%CRISTAL {"arestas":[{"alvo":"sigma_equivalencia_crown","confianca":1.0,"epistemico":"fato","origem":"wiki","rel":"relaciona"},{"alvo":"grupos_excepcionais_jordan","confianca":1.0,"epistemico":"fato","origem":"wiki","rel":"usa"}],"confianca":1.0,"contraexemplos":[],"descricao":"E8 tem 240 raízes; σ(240) = 744, o termo constante j₀ da função j (o moonshine). A álgebra de Lie excepcional encontra o Monster em um único passo de σ: 240 → 744.","epistemico":"fato","exemplos":[],"id":"e8_moonshine_sigma","memoria":"semantica","meta":{"autor":"Lord Index (Shawn R. Schiller)","dominio":"matematica","fonte":"A Catedral — Part VIII-A (jul 2026), ~/Imagens/catedral"},"origem":"wiki","palavras_chave":[],"sinonimos":[],"tipo":"conceito","titulo":"E8 ao Moonshine em um passo σ: σ(240) = 744 = j₀"}
\section{E8 ao Moonshine em um passo \ensuremath{\sigma}: \ensuremath{\sigma}(240) = 744 = j\ensuremath{_{0}}}
E8 tem 240 raízes; \ensuremath{\sigma}(240) = 744, o termo constante j\ensuremath{_{0}} da função j (o moonshine). A álgebra de Lie excepcional encontra o Monster em um único passo de \ensuremath{\sigma}: 240 \ensuremath{\to} 744.
\prov{relações: relaciona sigma\_equivalencia\_crown; usa grupos\_excepcionais\_jordan}
\prov{proveniência: wiki \textperiodcentered{} A Catedral — Part VIII-A (jul 2026), \~{}/Imagens/catedral \textperiodcentered{} fato \textperiodcentered{} confiança 1.0 \textperiodcentered{} domínio matematica \textperiodcentered{} história 3 versões no jornal}

%CRISTAL {"arestas":[{"alvo":"moeda_sl2","confianca":1.0,"epistemico":"fato","origem":"paper","rel":"usa"},{"alvo":"word","confianca":0.5,"epistemico":"fato","origem":"wiki","rel":"relaciona"},{"alvo":"vontade_de_poder","confianca":0.5,"epistemico":"fato","origem":"wiki","rel":"relaciona"}],"confianca":1.0,"contraexemplos":[],"descricao":"A moeda Mₖ é conjugada ao squeezer de Bogoliubov S(r) com r=2·ln(mₖ): correlaciona pesos independentes (emaranha), expandindo uma quadratura e comprimindo a outra.","epistemico":"fato","exemplos":[],"id":"emaranhador","memoria":"semantica","meta":{"autor":"Gentil Lopes","descoberta":true,"dominio":"matematica","fonte":"paper_24_emaranhador.pdf, Teo.2.1"},"origem":"paper","palavras_chave":[],"sinonimos":[],"tipo":"conceito","titulo":"Emaranhador (squeezer)"}
\section{Emaranhador (squeezer)}
A moeda M\ensuremath{_{k}} é conjugada ao squeezer de Bogoliubov S(r) com r=2\textperiodcentered ln(m\ensuremath{_{k}}): correlaciona pesos independentes (emaranha), expandindo uma quadratura e comprimindo a outra.
\prov{relações: usa moeda\_sl2; relaciona word; relaciona vontade\_de\_poder}
\prov{proveniência: paper \textperiodcentered{} paper\_24\_emaranhador.pdf, Teo.2.1 \textperiodcentered{} fato \textperiodcentered{} confiança 1.0 \textperiodcentered{} domínio matematica \textperiodcentered{} história 4 versões no jornal}

%CRISTAL {"arestas":[{"alvo":"mecanica_quantica_mineral","confianca":1.0,"epistemico":"fato","origem":"wiki","rel":"parte_de"},{"alvo":"gato_e_o_hamiltoniano","confianca":1.0,"epistemico":"fato","origem":"wiki","rel":"usa"},{"alvo":"dois_qubits_octonios","confianca":1.0,"epistemico":"fato","origem":"wiki","rel":"relaciona"}],"confianca":1.0,"contraexemplos":[],"descricao":"Acoplar dois gatos — H = H_A⊗I + I⊗H_B + g·(σx⊗σx) — produz um estado fundamental que NÃO fatora num produto: emaranhado. A medida é a entropia de von Neumann da matriz densidade REDUZIDA (o traço parcial sobre um subsistema): 0 para um produto, 1 bit para o estado de Bell (máximo). O emaranhamento cresce com o acoplamento g. Verif.: Bell S=1, dois gatos g=3 → S=0.79.","epistemico":"fato","exemplos":[],"id":"emaranhamento_mineral","memoria":"semantica","meta":{"dominio":"reta mineral","fonte":"emaranhamento e decoerência"},"origem":"wiki","palavras_chave":[],"sinonimos":[],"tipo":"conceito","titulo":"Emaranhamento Mineral (dois gatos enredados)"}
\section{Emaranhamento Mineral (dois gatos enredados)}
Acoplar dois gatos — H = H\_A\ensuremath{\otimes}I + I\ensuremath{\otimes}H\_B + g\textperiodcentered (\ensuremath{\sigma}x\ensuremath{\otimes}\ensuremath{\sigma}x) — produz um estado fundamental que NÃO fatora num produto: emaranhado. A medida é a entropia de von Neumann da matriz densidade REDUZIDA (o traço parcial sobre um subsistema): 0 para um produto, 1 bit para o estado de Bell (máximo). O emaranhamento cresce com o acoplamento g. Verif.: Bell S=1, dois gatos g=3 \ensuremath{\to} S=0.79.
\prov{relações: parte\_de mecanica\_quantica\_mineral; usa gato\_e\_o\_hamiltoniano; relaciona dois\_qubits\_octonios}
\prov{proveniência: wiki \textperiodcentered{} emaranhamento e decoerência \textperiodcentered{} fato \textperiodcentered{} confiança 1.0 \textperiodcentered{} domínio reta mineral \textperiodcentered{} história 4 versões no jornal}

%CRISTAL {"arestas":[{"alvo":"catedral_sigma_aritmetica","confianca":1.0,"epistemico":"fato","origem":"wiki","rel":"aplicacao"}],"confianca":1.0,"contraexemplos":[],"descricao":"Cross-aplicação à fisiologia espacial: eletromiografia (EMG) monitorando músculo em microgravidade, com o mapeamento numerológico EMG=25, monitor=104. A EMG é real; a associação numérica é da obra.","epistemico":"hipotese","exemplos":[],"id":"emg_no_espaco","memoria":"semantica","meta":{"autor":"Lord Index (Shawn R. Schiller)","dominio":"matematica","fonte":"A Catedral — Part VIII-A (jul 2026), ~/Imagens/catedral"},"origem":"wiki","palavras_chave":[],"sinonimos":[],"tipo":"conceito","titulo":"EMG no Espaço: miografia e os índices da Catedral"}
\section{EMG no Espaço: miografia e os índices da Catedral}
Cross-aplicação à fisiologia espacial: eletromiografia (EMG) monitorando músculo em microgravidade, com o mapeamento numerológico EMG=25, monitor=104. A EMG é real; a associação numérica é da obra.
\prov{relações: aplicacao catedral\_sigma\_aritmetica}
\prov{proveniência: wiki \textperiodcentered{} A Catedral — Part VIII-A (jul 2026), \~{}/Imagens/catedral \textperiodcentered{} hipotese \textperiodcentered{} confiança 1.0 \textperiodcentered{} domínio matematica \textperiodcentered{} história 2 versões no jornal}

%CRISTAL {"arestas":[{"alvo":"koch_codifica_cantorfibonacci_prova_e_evidencia","confianca":1.0,"epistemico":"fato","origem":"wiki","rel":"relaciona"},{"alvo":"dissipacao_prigogine","confianca":1.0,"epistemico":"fato","origem":"wiki","rel":"relaciona"},{"alvo":"corpo_tropical","confianca":1.0,"epistemico":"fato","origem":"wiki","rel":"relaciona"},{"alvo":"computacao_fractal","confianca":1.0,"epistemico":"fato","origem":"wiki","rel":"relaciona"}],"confianca":1.0,"contraexemplos":[],"descricao":"A energia Ed=Σ(ΔV)d do imposto V(s)=1−s² só é constante sob a renormalização f₁(s)=(s+2)/3 em d=ln2/ln3≈0.6309 (dimensão de Hausdorff do Cantor); a pressão de Bowen-Ruelle P(d)=ln2−d·ln3=0 nesse ponto. Conecta fractalidade à termodinâmica (pressão zero).","epistemico":"fato","exemplos":[],"id":"energia_fractal_cantor","memoria":"semantica","meta":{"autor":"broca-so","dominio":"fractal","fonte":"hiper/circular/energia_fractal.py"},"origem":"wiki","palavras_chave":[],"sinonimos":[],"tipo":"conceito","titulo":"Energia Fractal Conservada = Dimensão de Cantor (ln2/ln3)"}
\section{Energia Fractal Conservada = Dimensão de Cantor (ln2/ln3)}
A energia Ed=\ensuremath{\Sigma}(\ensuremath{\Delta}V)d do imposto V(s)=1\ensuremath{-}s\ensuremath{^{2}} só é constante sob a renormalização f\ensuremath{_{1}}(s)=(s+2)/3 em d=ln2/ln3\ensuremath{\approx}0.6309 (dimensão de Hausdorff do Cantor); a pressão de Bowen-Ruelle P(d)=ln2\ensuremath{-}d\textperiodcentered ln3=0 nesse ponto. Conecta fractalidade à termodinâmica (pressão zero).
\prov{relações: relaciona koch\_codifica\_cantorfibonacci\_prova\_e\_evidencia; relaciona dissipacao\_prigogine; relaciona corpo\_tropical; relaciona computacao\_fractal}
\prov{proveniência: wiki \textperiodcentered{} hiper/circular/energia\_fractal.py \textperiodcentered{} fato \textperiodcentered{} confiança 1.0 \textperiodcentered{} domínio fractal \textperiodcentered{} história 6 versões no jornal}

%CRISTAL {"arestas":[{"alvo":"matriz_cantor","confianca":1.0,"epistemico":"fato","origem":"paper","rel":"parte_de"},{"alvo":"virtualizacao","confianca":0.5,"epistemico":"fato","origem":"wiki","rel":"relaciona"},{"alvo":"word","confianca":0.5,"epistemico":"fato","origem":"wiki","rel":"relaciona"}],"confianca":1.0,"contraexemplos":[],"descricao":"h(X)=limₙ(1/n)log#palavras de tamanho n=ln λ (Perron): invariante métrico-dinâmico do subshift de tipo finito. Ouro: λ=φ, h≈0,481.","epistemico":"fato","exemplos":[],"id":"entropia_topologica","memoria":"semantica","meta":{"autor":"Gentil Lopes","descoberta":false,"dominio":"matematica","fonte":"paper_62_matriz_cantor.pdf, Def.3.2"},"origem":"paper","palavras_chave":[],"sinonimos":[],"tipo":"conceito","titulo":"Entropia Topológica"}
\section{Entropia Topológica}
h(X)=lim\ensuremath{_{n}}(1/n)log\#palavras de tamanho n=ln \ensuremath{\lambda} (Perron): invariante métrico-dinâmico do subshift de tipo finito. Ouro: \ensuremath{\lambda}=\ensuremath{\varphi}, h\ensuremath{\approx}0,481.
\prov{relações: parte\_de matriz\_cantor; relaciona virtualizacao; relaciona word}
\prov{proveniência: paper \textperiodcentered{} paper\_62\_matriz\_cantor.pdf, Def.3.2 \textperiodcentered{} fato \textperiodcentered{} confiança 1.0 \textperiodcentered{} domínio matematica \textperiodcentered{} história 5 versões no jornal}

%CRISTAL {"arestas":[{"alvo":"dimensao_de_traco","confianca":1.0,"epistemico":"fato","origem":"wiki","rel":"confirma"},{"alvo":"ifs_de_retorno","confianca":1.0,"epistemico":"fato","origem":"wiki","rel":"usa"}],"confianca":1.0,"contraexemplos":[],"descricao":"A dimensão do atrator resolve Σr=₀k(2−(r⁺¹))s=1, que reduz a yk⁺¹=yk+…+1 com y=2s; logo y=λₖ e dimH Kₖ=log₂λₖ. Terceira via independente que confirma Dtr (via scipy brentq).","epistemico":"fato","exemplos":[],"id":"equacao_de_moran_traco","memoria":"semantica","meta":{"autor":"broca-so","dominio":"fractal","fonte":"machine/fractal/fractal.py:dim_moran; computador_fractal.tex Teo.1"},"origem":"wiki","palavras_chave":[],"sinonimos":[],"tipo":"conceito","titulo":"A Equação de Moran do Traço"}
\section{A Equação de Moran do Traço}
A dimensão do atrator resolve \ensuremath{\Sigma}r=\ensuremath{_{0}}k(2\ensuremath{-}(r\ensuremath{^{+1}}))s=1, que reduz a yk\ensuremath{^{+1}}=yk+…+1 com y=2s; logo y=\ensuremath{\lambda}\ensuremath{_{k}} e dimH K\ensuremath{_{k}}=log\ensuremath{_{2}}\ensuremath{\lambda}\ensuremath{_{k}}. Terceira via independente que confirma Dtr (via scipy brentq).
\prov{relações: confirma dimensao\_de\_traco; usa ifs\_de\_retorno}
\prov{proveniência: wiki \textperiodcentered{} machine/fractal/fractal.py:dim\_moran; computador\_fractal.tex Teo.1 \textperiodcentered{} fato \textperiodcentered{} confiança 1.0 \textperiodcentered{} domínio fractal \textperiodcentered{} história 4 versões no jornal}

%CRISTAL {"arestas":[{"alvo":"teoria_dos_jogos","confianca":1.0,"epistemico":"fato","origem":"wiki","rel":"parte_de"},{"alvo":"estrategia_jogo","confianca":1.0,"epistemico":"fato","origem":"wiki","rel":"usa"}],"confianca":1.0,"contraexemplos":[],"descricao":"Um perfil de estratégias em que NINGUÉM ganha desviando sozinho — cada um responde otimamente ao outro. Nash provou que todo jogo finito tem ao menos um (em estratégias mistas). O conceito central.","epistemico":"fato","exemplos":[],"id":"equilibrio_de_nash","memoria":"semantica","meta":{"autor":"John Nash","dominio":"matematica","fonte":"teoria dos jogos"},"origem":"wiki","palavras_chave":[],"sinonimos":[],"tipo":"conceito","titulo":"Equilíbrio de Nash"}
\section{Equilíbrio de Nash}
Um perfil de estratégias em que NINGUÉM ganha desviando sozinho — cada um responde otimamente ao outro. Nash provou que todo jogo finito tem ao menos um (em estratégias mistas). O conceito central.
\prov{relações: parte\_de teoria\_dos\_jogos; usa estrategia\_jogo}
\prov{proveniência: wiki \textperiodcentered{} teoria dos jogos \textperiodcentered{} fato \textperiodcentered{} confiança 1.0 \textperiodcentered{} domínio matematica \textperiodcentered{} história 3 versões no jornal}

%CRISTAL {"arestas":[{"alvo":"unidades_algebricas_pu","confianca":1.0,"epistemico":"fato","origem":"wiki","rel":"parte_de"},{"alvo":"regua_de_gentil","confianca":1.0,"epistemico":"fato","origem":"wiki","rel":"relaciona"}],"confianca":1.0,"contraexemplos":[],"descricao":"A definição-âncora: 'o espaço-tempo é o círculo de Lopes; a luz dá 299.792.458 voltas nele por segundo'. Uma volta = a unidade de comprimento (o metro); c = a taxa voltas/segundo. Uma volta é (i) um ciclo da fase θ em ℝ/ℤ (a costura 1≡0), (ii) uma franja de interferência (inteiro n, Bohr-Sommerfeld), (iii) um comprimento de onda. DOIS círculos distintos: ψ (a régua de álgebras, Partes I-IV) vs θ (a fase física interna a ℂ_â) — a âncora dimensional é o círculo de θ.","epistemico":"inferencia","exemplos":[],"id":"escala_volta_gentil","memoria":"semantica","meta":{"autor":"Lopes/broca-so","dominio":"unidades","fonte":"hiper/circular/paper_5_unidades.tex"},"origem":"wiki","palavras_chave":[],"sinonimos":[],"tipo":"conceito","titulo":"A escala: uma volta no círculo de Lopes (c = 299.792.458 voltas/s)"}
\section{A escala: uma volta no círculo de Lopes (c = 299.792.458 voltas/s)}
A definição-âncora: 'o espaço-tempo é o círculo de Lopes; a luz dá 299.792.458 voltas nele por segundo'. Uma volta = a unidade de comprimento (o metro); c = a taxa voltas/segundo. Uma volta é (i) um ciclo da fase \ensuremath{\theta} em \ensuremath{\mathbb{R}}/\ensuremath{\mathbb{Z}} (a costura 1\ensuremath{\equiv}0), (ii) uma franja de interferência (inteiro n, Bohr-Sommerfeld), (iii) um comprimento de onda. DOIS círculos distintos: \ensuremath{\psi} (a régua de álgebras, Partes I-IV) vs \ensuremath{\theta} (a fase física interna a \ensuremath{\mathbb{C}}\_â) — a âncora dimensional é o círculo de \ensuremath{\theta}.
\prov{relações: parte\_de unidades\_algebricas\_pu; relaciona regua\_de\_gentil}
\prov{proveniência: wiki \textperiodcentered{} hiper/circular/paper\_5\_unidades.tex \textperiodcentered{} inferencia \textperiodcentered{} confiança 1.0 \textperiodcentered{} domínio unidades \textperiodcentered{} história 4 versões no jornal}

%CRISTAL {"arestas":[{"alvo":"hardness_octonica","confianca":1.0,"epistemico":"fato","origem":"paper","rel":"relaciona"},{"alvo":"word","confianca":0.5,"epistemico":"fato","origem":"wiki","rel":"relaciona"}],"confianca":1.0,"contraexemplos":[],"descricao":"Teorema: a holografia AdS/CFT preserva o conteúdo mas não dá β(g) em forma fechada simbólica; nem lattice, nem AdS/CFT, nem perturbação fornecem β universal.","epistemico":"fato","exemplos":[],"id":"escape_adscft","memoria":"semantica","meta":{"autor":"Gentil Lopes","descoberta":false,"dominio":"matematica","fonte":"paper_PQ_O_adscft_escape.pdf, Teo.7.1"},"origem":"paper","palavras_chave":[],"sinonimos":[],"tipo":"conceito","titulo":"AdS/CFT Não Escapa"}
\section{AdS/CFT Não Escapa}
Teorema: a holografia AdS/CFT preserva o conteúdo mas não dá \ensuremath{\beta}(g) em forma fechada simbólica; nem lattice, nem AdS/CFT, nem perturbação fornecem \ensuremath{\beta} universal.
\prov{relações: relaciona hardness\_octonica; relaciona word}
\prov{proveniência: paper \textperiodcentered{} paper\_PQ\_O\_adscft\_escape.pdf, Teo.7.1 \textperiodcentered{} fato \textperiodcentered{} confiança 1.0 \textperiodcentered{} domínio matematica \textperiodcentered{} história 3 versões no jornal}

%CRISTAL {"arestas":[{"alvo":"computacao_fractal","confianca":1.0,"epistemico":"fato","origem":"wiki","rel":"parte_de"},{"alvo":"stratum_difusao_cpu","confianca":1.0,"epistemico":"fato","origem":"wiki","rel":"explica"}],"confianca":1.0,"contraexemplos":[],"descricao":"Uma dataclass congelada com todos os parâmetros do modelo (frac, cpu_total, olho_threads, chunk, passos, rodadas, gpu_celulas, c2, calor, canal_ok) e derivados (broca_passos_total = passos·gpu_celulas·rodadas). O objeto central imutável da difusão.","epistemico":"fato","exemplos":[],"id":"espaco_difusivo","memoria":"semantica","meta":{"autor":"broca-so","dominio":"fractal","fonte":"neuronio/difusao_recursos.py:EspacoDifusivo"},"origem":"wiki","palavras_chave":[],"sinonimos":[],"tipo":"conceito","titulo":"EspacoDifusivo — o Snapshot da Difusão"}
\section{EspacoDifusivo — o Snapshot da Difusão}
Uma dataclass congelada com todos os parâmetros do modelo (frac, cpu\_total, olho\_threads, chunk, passos, rodadas, gpu\_celulas, c2, calor, canal\_ok) e derivados (broca\_passos\_total = passos\textperiodcentered gpu\_celulas\textperiodcentered rodadas). O objeto central imutável da difusão.
\prov{relações: parte\_de computacao\_fractal; explica stratum\_difusao\_cpu}
\prov{proveniência: wiki \textperiodcentered{} neuronio/difusao\_recursos.py:EspacoDifusivo \textperiodcentered{} fato \textperiodcentered{} confiança 1.0 \textperiodcentered{} domínio fractal \textperiodcentered{} história 3 versões no jornal}

%CRISTAL {"arestas":[{"alvo":"algebra_tropical","confianca":1.0,"epistemico":"fato","origem":"paper","rel":"usa"},{"alvo":"virtualizacao","confianca":0.5,"epistemico":"fato","origem":"wiki","rel":"relaciona"},{"alvo":"misticismo_nao_dual","confianca":0.5,"epistemico":"fato","origem":"wiki","rel":"relaciona"}],"confianca":1.0,"contraexemplos":[],"descricao":"Espaço das entropias h=ln λ (uma coordenada por poda), variedade com operador max-plus; pontos entre a morte (h→0) e o caos (h=ln 2).","epistemico":"fato","exemplos":[],"id":"espaco_tropical","memoria":"semantica","meta":{"autor":"Gentil Lopes","descoberta":true,"dominio":"matematica","fonte":"paper_65_espaco_tropical.pdf"},"origem":"paper","palavras_chave":[],"sinonimos":[],"tipo":"conceito","titulo":"Espaço Tropical"}
\section{Espaço Tropical}
Espaço das entropias h=ln \ensuremath{\lambda} (uma coordenada por poda), variedade com operador max-plus; pontos entre a morte (h\ensuremath{\to}0) e o caos (h=ln 2).
\prov{relações: usa algebra\_tropical; relaciona virtualizacao; relaciona misticismo\_nao\_dual}
\prov{proveniência: paper \textperiodcentered{} paper\_65\_espaco\_tropical.pdf \textperiodcentered{} fato \textperiodcentered{} confiança 1.0 \textperiodcentered{} domínio matematica \textperiodcentered{} história 4 versões no jornal}

%CRISTAL {"arestas":[{"alvo":"transformacao_linear","confianca":1.0,"epistemico":"fato","origem":"curriculo","rel":"usa"},{"alvo":"vetores","confianca":1.0,"epistemico":"fato","origem":"curriculo","rel":"especializa"}],"confianca":1.0,"contraexemplos":[],"descricao":"Conjunto de vetores fechado sob adição e produto escalar; base da álgebra linear.","epistemico":"fato","exemplos":[],"id":"espaco_vetorial","memoria":"semantica","meta":{"dominio":"matematica","nivel":4,"ordem":6},"origem":"curriculo","palavras_chave":["base","dimensão","subespaço"],"sinonimos":["espaço vetorial"],"tipo":"conceito","titulo":"espaco vetorial"}
\section{espaco vetorial}
Conjunto de vetores fechado sob adição e produto escalar; base da álgebra linear.
\prov{sinónimos: espaço vetorial}
\prov{relações: usa transformacao\_linear; especializa vetores}
\prov{proveniência: curriculo \textperiodcentered{} fato \textperiodcentered{} confiança 1.0 \textperiodcentered{} domínio matematica \textperiodcentered{} história 55 versões no jornal}

%CRISTAL {"arestas":[{"alvo":"matematica","confianca":1.0,"epistemico":"fato","origem":"paper","rel":"parte_de"},{"alvo":"gato","confianca":1.0,"epistemico":"fato","origem":"paper","rel":"referencia"}],"confianca":1.0,"contraexemplos":[],"descricao":"Uma indução é a tríade (T,I,κ): T operador de passo (n→n+1), I invariante com I∘T=I (o osso topológico), κ cúspide (ponto fixo/atrator de T). Iterar T a partir de uma base gera a estrutura.","epistemico":"fato","exemplos":[],"id":"esquema_inducao","memoria":"semantica","meta":{"autor":"Gentil Lopes","descoberta":true,"dominio":"matematica","fonte":"paper_57_inducao_universal.pdf, Def.1.1"},"origem":"paper","palavras_chave":[],"sinonimos":[],"tipo":"conceito","titulo":"Esquema de Indução (T, I, κ)"}
\section{Esquema de Indução (T, I, \ensuremath{\kappa})}
Uma indução é a tríade (T,I,\ensuremath{\kappa}): T operador de passo (n\ensuremath{\to}n+1), I invariante com I\ensuremath{\circ}T=I (o osso topológico), \ensuremath{\kappa} cúspide (ponto fixo/atrator de T). Iterar T a partir de uma base gera a estrutura.
\prov{relações: parte\_de matematica; referencia gato}
\prov{proveniência: paper \textperiodcentered{} paper\_57\_inducao\_universal.pdf, Def.1.1 \textperiodcentered{} fato \textperiodcentered{} confiança 1.0 \textperiodcentered{} domínio matematica \textperiodcentered{} história 3 versões no jornal}

%CRISTAL {"arestas":[{"alvo":"transformada_metalica","confianca":1.0,"epistemico":"fato","origem":"wiki","rel":"usa"},{"alvo":"torre_unificada","confianca":1.0,"epistemico":"fato","origem":"wiki","rel":"parte_de"},{"alvo":"mineracao_conservacao_verificavel","confianca":1.0,"epistemico":"fato","origem":"wiki","rel":"usa"},{"alvo":"recuperacao_evidencia_invariante","confianca":1.0,"epistemico":"fato","origem":"wiki","rel":"usa"},{"alvo":"conservacao_parabola_costura","confianca":1.0,"epistemico":"fato","origem":"wiki","rel":"usa"},{"alvo":"bai_gate_dtc_orbita","confianca":1.0,"epistemico":"fato","origem":"wiki","rel":"usa"},{"alvo":"costura_prata_dobra8","confianca":1.0,"epistemico":"fato","origem":"wiki","rel":"usa"},{"alvo":"costura_cobre_dobra10","confianca":1.0,"epistemico":"fato","origem":"wiki","rel":"usa"},{"alvo":"chat_aberto_cobre_explora","confianca":1.0,"epistemico":"fato","origem":"wiki","rel":"usa"},{"alvo":"definicao_frasagem_ampla","confianca":1.0,"epistemico":"fato","origem":"wiki","rel":"usa"},{"alvo":"face_polida_primeira_frase","confianca":1.0,"epistemico":"fato","origem":"wiki","rel":"usa"},{"alvo":"transformada_metalica_fechada_perelman","confianca":1.0,"epistemico":"fato","origem":"wiki","rel":"usa"},{"alvo":"refino_multimetal_live_gex44","confianca":1.0,"epistemico":"fato","origem":"wiki","rel":"usa"},{"alvo":"tiffany","confianca":1.0,"epistemico":"fato","origem":"wiki","rel":"parte_de"}],"confianca":1.0,"contraexemplos":[],"descricao":"Consolidação: a TRANSFORMADA METÁLICA (β-numeração Pisot + Rauzy/dobra + parábola-conservação + Peirce, fechada por Perelman — ver [[transformada_metalica]]) virou a espinha operacional de TODO o sistema, sobre um invariante só: a+c²=1 (σ·conj=−1). (1) NO METAL/MINERAÇÃO: cada floco é refinado em metais e a conservação é VERIFICADA por floco (a norma ±1) — o gex44 roda live e auditado ([[mineracao_conservacao_verificavel]], [[refino_multimetal_live_gex44]]). (2) NO GRAFO: 78 arestas verificadas ancoram do metal ao grafo no hub da transformada (revisão por agentes). (3) NA CONVERSA (as 3 retas COMPLETAS): a costura DTC reta por reta pela sua dobra — OURO radial ([[bai_gate_dtc_orbita]]) = definição, que entende frasagem ampla ([[definicao_frasagem_ampla]]) e responde curto ([[face_polida_primeira_frase]]); PRATA dobra 8 ([[costura_prata_dobra8]]) = a continuidade/fio; COBRE dobra 10 ([[costura_cobre_dobra10]]) = o chat aberto que EXPLORA conceitos por rotação áurea ([[chat_aberto_cobre_explora]]) — e a recuperação por evidência que OBEDECE o invariante ([[recuperacao_evidencia_invariante]]), sem invadir a parábola ([[conservacao_parabola_costura]]). (4) O FECHO: a mesma entropia de Perelman aterra a recuperação E fecha a transformada ([[transformada_metalica_fechada_perelman]]). PROTOCOLOS (revisados, verdes): grafo 1 componente / 0 isolados / 62 testes; chat fluxo + recuperação; mineração gex44 viva, todos os metais conservam. É a Universidade Tiffany num princípio operacional só: a transformada metálica na parábola. [I] o estado completo.","epistemico":"inferencia","exemplos":[],"id":"estado_completo_transformada","memoria":"semantica","meta":{"autor":"Lopes/broca-so","dominio":"algebra","fonte":"sessão 2 jul 2026: transformada + costura (ouro/prata/cobre) + recuperação + mineração verificável + 78 arestas; protocolos floresta/reproduzir/verificar"},"origem":"wiki","palavras_chave":[],"sinonimos":[],"tipo":"conceito","titulo":"ESTADO COMPLETO — a transformada metálica é a espinha do metal ao grafo à conversa à mineração"}
\section{ESTADO COMPLETO — a transformada metálica é a espinha do metal ao grafo à conversa à mineração}
Consolidação: a TRANSFORMADA METÁLICA (\ensuremath{\beta}-numeração Pisot + Rauzy/dobra + parábola-conservação + Peirce, fechada por Perelman — ver [[transformada\_metalica]]) virou a espinha operacional de TODO o sistema, sobre um invariante só: a+c\ensuremath{^{2}}=1 (\ensuremath{\sigma}\textperiodcentered conj=\ensuremath{-}1). (1) NO METAL/MINERAÇÃO: cada floco é refinado em metais e a conservação é VERIFICADA por floco (a norma $\pm$1) — o gex44 roda live e auditado ([[mineracao\_conservacao\_verificavel]], [[refino\_multimetal\_live\_gex44]]). (2) NO GRAFO: 78 arestas verificadas ancoram do metal ao grafo no hub da transformada (revisão por agentes). (3) NA CONVERSA (as 3 retas COMPLETAS): a costura DTC reta por reta pela sua dobra — OURO radial ([[bai\_gate\_dtc\_orbita]]) = definição, que entende frasagem ampla ([[definicao\_frasagem\_ampla]]) e responde curto ([[face\_polida\_primeira\_frase]]); PRATA dobra 8 ([[costura\_prata\_dobra8]]) = a continuidade/fio; COBRE dobra 10 ([[costura\_cobre\_dobra10]]) = o chat aberto que EXPLORA conceitos por rotação áurea ([[chat\_aberto\_cobre\_explora]]) — e a recuperação por evidência que OBEDECE o invariante ([[recuperacao\_evidencia\_invariante]]), sem invadir a parábola ([[conservacao\_parabola\_costura]]). (4) O FECHO: a mesma entropia de Perelman aterra a recuperação E fecha a transformada ([[transformada\_metalica\_fechada\_perelman]]). PROTOCOLOS (revisados, verdes): grafo 1 componente / 0 isolados / 62 testes; chat fluxo + recuperação; mineração gex44 viva, todos os metais conservam. É a Universidade Tiffany num princípio operacional só: a transformada metálica na parábola. [I] o estado completo.
\prov{relações: usa transformada\_metalica; parte\_de torre\_unificada; usa mineracao\_conservacao\_verificavel; usa recuperacao\_evidencia\_invariante; usa conservacao\_parabola\_costura; usa bai\_gate\_dtc\_orbita; usa costura\_prata\_dobra8; usa costura\_cobre\_dobra10; usa chat\_aberto\_cobre\_explora; usa definicao\_frasagem\_ampla; usa face\_polida\_primeira\_frase; usa transformada\_metalica\_fechada\_perelman; usa refino\_multimetal\_live\_gex44; parte\_de tiffany}
\prov{proveniência: wiki \textperiodcentered{} sessão 2 jul 2026: transformada + costura (ouro/prata/cobre) + recuperação + mineração verificável + 78 arestas; protocolos floresta/reproduzir/verificar \textperiodcentered{} inferencia \textperiodcentered{} confiança 1.0 \textperiodcentered{} domínio algebra \textperiodcentered{} história 42 versões no jornal}

%CRISTAL {"arestas":[{"alvo":"unidades_algebricas_pu","confianca":1.0,"epistemico":"fato","origem":"wiki","rel":"explica"}],"confianca":1.0,"contraexemplos":[],"descricao":"A régua epistemológica da série inteira. DERIVADO: a estrutura adimensional — os números puros (√2, 3, Q=N/3, 1/2π, β_𝕆, ΣQ²=8) fixados pela álgebra; o SI os toma do experimento, aqui têm origem. POSTO: a escala — uma volta com conteúdo dimensional, a do meio (o campo do tubo de fluxo, o horizonte), não parâmetro livre. CONVENÇÃO: mol e candela (contagem/percepção humanas, fora da física). O que a álgebra NÃO dá: o valor de α (corre) e a escala absoluta. A honestidade que distingue derivar de vestir.","epistemico":"inferencia","exemplos":[],"id":"estatuto_derivado_posto_convencao","memoria":"semantica","meta":{"autor":"Lopes/broca-so","dominio":"unidades","fonte":"hiper/circular/paper_5_unidades.tex"},"origem":"wiki","palavras_chave":[],"sinonimos":[],"tipo":"conceito","titulo":"O estatuto: derivado / posto / convenção"}
\section{O estatuto: derivado / posto / convenção}
A régua epistemológica da série inteira. DERIVADO: a estrutura adimensional — os números puros (\ensuremath{\sqrt{\,}}2, 3, Q=N/3, 1/2\ensuremath{\pi}, \ensuremath{\beta}\_\ensuremath{\mathbb{O}}, \ensuremath{\Sigma}Q\ensuremath{^{2}}=8) fixados pela álgebra; o SI os toma do experimento, aqui têm origem. POSTO: a escala — uma volta com conteúdo dimensional, a do meio (o campo do tubo de fluxo, o horizonte), não parâmetro livre. CONVENÇÃO: mol e candela (contagem/percepção humanas, fora da física). O que a álgebra NÃO dá: o valor de \ensuremath{\alpha} (corre) e a escala absoluta. A honestidade que distingue derivar de vestir.
\prov{relações: explica unidades\_algebricas\_pu}
\prov{proveniência: wiki \textperiodcentered{} hiper/circular/paper\_5\_unidades.tex \textperiodcentered{} inferencia \textperiodcentered{} confiança 1.0 \textperiodcentered{} domínio unidades \textperiodcentered{} história 3 versões no jornal}

%CRISTAL {"arestas":[{"alvo":"estrategia_jogo","confianca":1.0,"epistemico":"fato","origem":"wiki","rel":"especializa"}],"confianca":1.0,"contraexemplos":[],"descricao":"A que é melhor não importa o que o outro faça. Quando existe, decide o jogo — mas o dilema do prisioneiro mostra que o equilíbrio dominante pode ser pior para todos.","epistemico":"fato","exemplos":[],"id":"estrategia_dominante","memoria":"semantica","meta":{"dominio":"matematica","fonte":"teoria dos jogos"},"origem":"wiki","palavras_chave":[],"sinonimos":[],"tipo":"conceito","titulo":"Estratégia Dominante"}
\section{Estratégia Dominante}
A que é melhor não importa o que o outro faça. Quando existe, decide o jogo — mas o dilema do prisioneiro mostra que o equilíbrio dominante pode ser pior para todos.
\prov{relações: especializa estrategia\_jogo}
\prov{proveniência: wiki \textperiodcentered{} teoria dos jogos \textperiodcentered{} fato \textperiodcentered{} confiança 1.0 \textperiodcentered{} domínio matematica \textperiodcentered{} história 2 versões no jornal}

%CRISTAL {"arestas":[{"alvo":"jogo_estrategico","confianca":1.0,"epistemico":"fato","origem":"wiki","rel":"parte_de"}],"confianca":1.0,"contraexemplos":[],"descricao":"Um plano de ação completo: o que o jogador faz em CADA situação possível do jogo. Pura (uma escolha) ou mista (uma distribuição de probabilidade sobre as escolhas).","epistemico":"fato","exemplos":[],"id":"estrategia_jogo","memoria":"semantica","meta":{"dominio":"matematica","fonte":"teoria dos jogos"},"origem":"wiki","palavras_chave":[],"sinonimos":[],"tipo":"conceito","titulo":"Estratégia (plano completo)"}
\section{Estratégia (plano completo)}
Um plano de ação completo: o que o jogador faz em CADA situação possível do jogo. Pura (uma escolha) ou mista (uma distribuição de probabilidade sobre as escolhas).
\prov{relações: parte\_de jogo\_estrategico}
\prov{proveniência: wiki \textperiodcentered{} teoria dos jogos \textperiodcentered{} fato \textperiodcentered{} confiança 1.0 \textperiodcentered{} domínio matematica \textperiodcentered{} história 2 versões no jornal}

%CRISTAL {"arestas":[{"alvo":"equilibrio_de_nash","confianca":1.0,"epistemico":"fato","origem":"wiki","rel":"usa"}],"confianca":1.0,"contraexemplos":[],"descricao":"Jogar aleatoriamente com probabilidades calculadas — o que torna o oponente indiferente e estabiliza o equilíbrio quando não há equilíbrio puro (pedra-papel-tesoura).","epistemico":"fato","exemplos":[],"id":"estrategia_mista","memoria":"semantica","meta":{"dominio":"matematica","fonte":"teoria dos jogos"},"origem":"wiki","palavras_chave":[],"sinonimos":[],"tipo":"conceito","titulo":"Estratégia Mista"}
\section{Estratégia Mista}
Jogar aleatoriamente com probabilidades calculadas — o que torna o oponente indiferente e estabiliza o equilíbrio quando não há equilíbrio puro (pedra-papel-tesoura).
\prov{relações: usa equilibrio\_de\_nash}
\prov{proveniência: wiki \textperiodcentered{} teoria dos jogos \textperiodcentered{} fato \textperiodcentered{} confiança 1.0 \textperiodcentered{} domínio matematica \textperiodcentered{} história 2 versões no jornal}

%CRISTAL {"arestas":[{"alvo":"corpo_estelar","confianca":1.0,"epistemico":"fato","origem":"wiki","rel":"parte_de"},{"alvo":"espectro_irredutivel_floco","confianca":1.0,"epistemico":"fato","origem":"wiki","rel":"eh_um"},{"alvo":"floco_recipiente_termico","confianca":1.0,"epistemico":"fato","origem":"wiki","rel":"relaciona"},{"alvo":"nucleossintese_dos_flocos","confianca":1.0,"epistemico":"fato","origem":"wiki","rel":"relaciona"},{"alvo":"limite_de_fusao_do_floco","confianca":1.0,"epistemico":"fato","origem":"wiki","rel":"relaciona"},{"alvo":"modos_do_floco_sao_particulas","confianca":1.0,"epistemico":"fato","origem":"wiki","rel":"relaciona"},{"alvo":"floco_de_neve","confianca":1.0,"epistemico":"fato","origem":"wiki","rel":"eh_um"},{"alvo":"fusao_estelar","confianca":1.0,"epistemico":"fato","origem":"wiki","rel":"relaciona"},{"alvo":"estrela_clifford_projetado","confianca":1.0,"epistemico":"fato","origem":"wiki","rel":"relaciona"},{"alvo":"nucleo_estelar","confianca":1.0,"epistemico":"fato","origem":"wiki","rel":"relaciona"}],"confianca":1.0,"contraexemplos":[],"descricao":"Unificação da linguagem: o floco ESPECTRAL — o floco da obra, o espectro irredutível {σ_i}, o recipiente térmico, o que La Hire multiplica — é promovido a ESTRELA (como goldchain→cristalchain). A estrela é o elemento do corpo estelar; o floco de neve É uma estrela (funde ao vale do ferro). NÃO se promove o floco de KOCH/mineração (floco_no_metal, o selo do canal, o tecido do gex44): é o floco de neve original, a carga viva — e ids de nó não se renomeiam (quebraria arestas). Cristaliza e relinka.","epistemico":"inferencia","exemplos":[],"id":"estrela_promove_floco","memoria":"semantica","meta":{"dominio":"algebra","fonte":"corpo estelar"},"origem":"wiki","palavras_chave":[],"sinonimos":[],"tipo":"conceito","titulo":"A Estrela (o Floco Espectral Promovido)"}
\section{A Estrela (o Floco Espectral Promovido)}
Unificação da linguagem: o floco ESPECTRAL — o floco da obra, o espectro irredutível \{\ensuremath{\sigma}\_i\}, o recipiente térmico, o que La Hire multiplica — é promovido a ESTRELA (como goldchain\ensuremath{\to}cristalchain). A estrela é o elemento do corpo estelar; o floco de neve É uma estrela (funde ao vale do ferro). NÃO se promove o floco de KOCH/mineração (floco\_no\_metal, o selo do canal, o tecido do gex44): é o floco de neve original, a carga viva — e ids de nó não se renomeiam (quebraria arestas). Cristaliza e relinka.
\prov{relações: parte\_de corpo\_estelar; eh\_um espectro\_irredutivel\_floco; relaciona floco\_recipiente\_termico; relaciona nucleossintese\_dos\_flocos; relaciona limite\_de\_fusao\_do\_floco; relaciona modos\_do\_floco\_sao\_particulas; eh\_um floco\_de\_neve; relaciona fusao\_estelar; relaciona estrela\_clifford\_projetado; relaciona nucleo\_estelar}
\prov{proveniência: wiki \textperiodcentered{} corpo estelar \textperiodcentered{} inferencia \textperiodcentered{} confiança 1.0 \textperiodcentered{} domínio algebra \textperiodcentered{} história 14 versões no jornal}

%CRISTAL {"arestas":[{"alvo":"tensor_cayley","confianca":1.0,"epistemico":"fato","origem":"paper","rel":"usa"},{"alvo":"programa_ns_bloqueios","confianca":1.0,"epistemico":"fato","origem":"paper","rel":"parte_de"}],"confianca":1.0,"contraexemplos":[],"descricao":"Conjectura 15: o excesso E_T daria dissipação a NS em ℝ⁷. Lopes REFUTA numericamente (60 campos: ∫ω·E_T sem sinal definido) — a janela octônica não basta para regularizar.","epistemico":"hipotese","exemplos":[],"id":"excesso_octonico_refutado","memoria":"semantica","meta":{"autor":"Gentil Lopes","descoberta":false,"dominio":"matematica","fonte":"paper_AG.pdf, Teo.18","refutado":true},"origem":"paper","palavras_chave":[],"sinonimos":[],"tipo":"conceito","titulo":"Excesso Octônico (REFUTADO)"}
\section{Excesso Octônico (REFUTADO)}
Conjectura 15: o excesso E\_T daria dissipação a NS em \ensuremath{\mathbb{R}}\ensuremath{^{7}}. Lopes REFUTA numericamente (60 campos: \ensuremath{\int}\ensuremath{\omega}\textperiodcentered E\_T sem sinal definido) — a janela octônica não basta para regularizar.
\prov{relações: usa tensor\_cayley; parte\_de programa\_ns\_bloqueios}
\prov{proveniência: paper \textperiodcentered{} paper\_AG.pdf, Teo.18 \textperiodcentered{} hipotese \textperiodcentered{} confiança 1.0 \textperiodcentered{} domínio matematica \textperiodcentered{} história 4 versões no jornal}

%CRISTAL {"arestas":[{"alvo":"pi_acima_da_reta_mineral","confianca":1.0,"epistemico":"fato","origem":"wiki","rel":"usa"},{"alvo":"pi_gerador_da_reta_acima","confianca":1.0,"epistemico":"fato","origem":"wiki","rel":"usa"},{"alvo":"corpo_mineral","confianca":1.0,"epistemico":"fato","origem":"wiki","rel":"usa"}],"confianca":1.0,"contraexemplos":[],"descricao":"A agulha de Dirac fixada em π e dilatada pela reta mineral: a expansão em frações contínuas de π é o movimento da agulha, passo a passo, mantendo a invariante estável. O quadro que unifica a CF, os gatos, os metais, a curvatura N−1, a cifra e o corpo de primos num só experimento.","epistemico":"inferencia","exemplos":[],"id":"experimento_da_agulha","memoria":"semantica","meta":{"dominio":"reta mineral","fonte":"experimento da agulha"},"origem":"wiki","palavras_chave":[],"sinonimos":[],"tipo":"conceito","titulo":"O Experimento da Agulha (π pela reta mineral)"}
\section{O Experimento da Agulha (\ensuremath{\pi} pela reta mineral)}
A agulha de Dirac fixada em \ensuremath{\pi} e dilatada pela reta mineral: a expansão em frações contínuas de \ensuremath{\pi} é o movimento da agulha, passo a passo, mantendo a invariante estável. O quadro que unifica a CF, os gatos, os metais, a curvatura N\ensuremath{-}1, a cifra e o corpo de primos num só experimento.
\prov{relações: usa pi\_acima\_da\_reta\_mineral; usa pi\_gerador\_da\_reta\_acima; usa corpo\_mineral}
\prov{proveniência: wiki \textperiodcentered{} experimento da agulha \textperiodcentered{} inferencia \textperiodcentered{} confiança 1.0 \textperiodcentered{} domínio reta mineral \textperiodcentered{} história 4 versões no jornal}

%CRISTAL {"arestas":[{"alvo":"recuperacao_evidencia_invariante","confianca":1.0,"epistemico":"fato","origem":"wiki","rel":"relaciona"},{"alvo":"definicao_frasagem_ampla","confianca":1.0,"epistemico":"fato","origem":"wiki","rel":"relaciona"},{"alvo":"tiffany","confianca":1.0,"epistemico":"fato","origem":"wiki","rel":"parte_de"}],"confianca":1.0,"contraexemplos":[],"descricao":"conversa/colapso.py (extrairface, limpararestas, primeirafrase): a resposta deixou de ser o texto cru do nó (título + descrição longa + cauda '— relaciona X — usa Y'). Agora, para nós de conhecimento (skn): (1) corta a cauda de arestas; (2) trunca na PRIMEIRA FRASE real (o 1º fim-de-frase ≥40 chars, ignorando decimal 3.79 e math..) ou, se a frase passa de ~180, no limite na fronteira de palavra com reticências. Saudações/turnos (dialogos) não truncam (já são curtos). Efeito: 'o que é entropia?' responde a definição (ex. 'Entropia Topológica — h=ln λ … subshift de tipo finito.'), não um parágrafo inteiro. É o acabamento da recuperação: achar o nó certo ([[recuperacaoevidenciaᵢnvariante]]) E apresentá-lo limpo. [D] a Tiffany fala curto.","epistemico":"fato","exemplos":[],"id":"face_polida_primeira_frase","memoria":"semantica","meta":{"autor":"Lopes/broca-so","dominio":"algebra","fonte":"conversa/colapso.py:extrair_face, _limpar_arestas, _primeira_frase (skn → 1ª frase, máx ~180)"},"origem":"wiki","palavras_chave":[],"sinonimos":[],"tipo":"conceito","titulo":"A face da Tiffany é polida — nó de conhecimento cortado na primeira frase (não parágrafo)"}
\section{A face da Tiffany é polida — nó de conhecimento cortado na primeira frase (não parágrafo)}
conversa/colapso.py (extrairface, limpararestas, primeirafrase): a resposta deixou de ser o texto cru do nó (título + descrição longa + cauda '— relaciona X — usa Y'). Agora, para nós de conhecimento (skn): (1) corta a cauda de arestas; (2) trunca na PRIMEIRA FRASE real (o 1º fim-de-frase \ensuremath{\ge}40 chars, ignorando decimal 3.79 e math..) ou, se a frase passa de \~{}180, no limite na fronteira de palavra com reticências. Saudações/turnos (dialogos) não truncam (já são curtos). Efeito: 'o que é entropia?' responde a definição (ex. 'Entropia Topológica — h=ln \ensuremath{\lambda} … subshift de tipo finito.'), não um parágrafo inteiro. É o acabamento da recuperação: achar o nó certo ([[recuperacaoevidencia\ensuremath{_{i}}nvariante]]) E apresentá-lo limpo. [D] a Tiffany fala curto.
\prov{relações: relaciona recuperacao\_evidencia\_invariante; relaciona definicao\_frasagem\_ampla; parte\_de tiffany}
\prov{proveniência: wiki \textperiodcentered{} conversa/colapso.py:extrair\_face, \_limpar\_arestas, \_primeira\_frase (skn \ensuremath{\to} 1ª frase, máx \~{}180) \textperiodcentered{} fato \textperiodcentered{} confiança 1.0 \textperiodcentered{} domínio algebra \textperiodcentered{} história 58 versões no jornal}

%CRISTAL {"arestas":[{"alvo":"conversa_recupera_nao_gera","confianca":1.0,"epistemico":"fato","origem":"wiki","rel":"especializa"},{"alvo":"aboutness_canonica_graduada","confianca":1.0,"epistemico":"fato","origem":"wiki","rel":"usa"},{"alvo":"tiffany","confianca":1.0,"epistemico":"fato","origem":"wiki","rel":"parte_de"}],"confianca":1.0,"contraexemplos":[],"descricao":"conversa/colapso.py: a disciplina do projeto (não fabricar, ver [[conversa_recupera_nao_gera]]) aplicada às DEFINIÇÕES. Uma pergunta 'o que é X' cujo conceito NÃO está no grafo respondia com um nó aleatório ('qual é a capital da frança?' → 'virtualização'/'Capital Cultural'). Agora, se há alvo de definição e o vencedor tem c² < 0.5 (sem aboutness real — não casou título nem conteúdo), a Tiffany responde honestamente 'Ainda não tenho «X» no meu grafo'. Correções que destravaram: a regex 'quais?' não casava 'qual' (→ 'qua(l|is)'); a ordem dos artigos ('os metais'→'s metais'); e a forte PARCIAL agora é PROPORCIONAL (0.6·|aw∩tw|/|aw|) — casar 1 de 2 palavras do alvo = 0.3 (abaixo do limiar), então 'capital da frança' não vira 'Capital Cultural' por coincidência de 'capital'. Fora do domínio → honesto; conceito real (1 ou N palavras) → responde. É o fail-closed do colapso estendido às definições. [F] verificado.","epistemico":"fato","exemplos":[],"id":"fail_closed_definicao_honesta","memoria":"semantica","meta":{"autor":"Lopes/broca-so","dominio":"algebra","fonte":"conversa/colapso.py (c²<0.5 + alvo → 'não tenho'); parabola_algebra (qua(l|is), ordem artigos, parcial proporcional)"},"origem":"wiki","palavras_chave":[],"sinonimos":[],"tipo":"conceito","titulo":"Fail-closed nas definições: 'não tenho isso' p/ fora do domínio (não fabrica) [F]"}
\section{Fail-closed nas definições: 'não tenho isso' p/ fora do domínio (não fabrica) [F]}
conversa/colapso.py: a disciplina do projeto (não fabricar, ver [[conversa\_recupera\_nao\_gera]]) aplicada às DEFINIÇÕES. Uma pergunta 'o que é X' cujo conceito NÃO está no grafo respondia com um nó aleatório ('qual é a capital da frança?' \ensuremath{\to} 'virtualização'/'Capital Cultural'). Agora, se há alvo de definição e o vencedor tem c\ensuremath{^{2}} < 0.5 (sem aboutness real — não casou título nem conteúdo), a Tiffany responde honestamente 'Ainda não tenho «X» no meu grafo'. Correções que destravaram: a regex 'quais?' não casava 'qual' (\ensuremath{\to} 'qua(l|is)'); a ordem dos artigos ('os metais'\ensuremath{\to}'s metais'); e a forte PARCIAL agora é PROPORCIONAL (0.6\textperiodcentered |aw\ensuremath{\cap}tw|/|aw|) — casar 1 de 2 palavras do alvo = 0.3 (abaixo do limiar), então 'capital da frança' não vira 'Capital Cultural' por coincidência de 'capital'. Fora do domínio \ensuremath{\to} honesto; conceito real (1 ou N palavras) \ensuremath{\to} responde. É o fail-closed do colapso estendido às definições. [F] verificado.
\prov{relações: especializa conversa\_recupera\_nao\_gera; usa aboutness\_canonica\_graduada; parte\_de tiffany}
\prov{proveniência: wiki \textperiodcentered{} conversa/colapso.py (c\ensuremath{^{2}}<0.5 + alvo \ensuremath{\to} 'não tenho'); parabola\_algebra (qua(l|is), ordem artigos, parcial proporcional) \textperiodcentered{} fato \textperiodcentered{} confiança 1.0 \textperiodcentered{} domínio algebra \textperiodcentered{} história 16 versões no jornal}

%CRISTAL {"arestas":[{"alvo":"catedral_sigma_aritmetica","confianca":1.0,"epistemico":"fato","origem":"wiki","rel":"parte_de"},{"alvo":"plano_fano_octonico","confianca":1.0,"epistemico":"fato","origem":"wiki","rel":"relaciona"},{"alvo":"psl_2_7","confianca":1.0,"epistemico":"fato","origem":"wiki","rel":"usa"},{"alvo":"constante_168","confianca":1.0,"epistemico":"fato","origem":"wiki","rel":"confirma"},{"alvo":"psl_embeds_g2","confianca":1.0,"epistemico":"fato","origem":"wiki","rel":"relaciona"}],"confianca":1.0,"contraexemplos":[],"descricao":"O plano de Fano (7 pontos, 7 retas, a geometria finita mínima) tem grupo de automorfismos Aut(Fano) ≅ PSL(2,7), de ordem |PSL(2,7)| = 168. A semente de toda a cadeia σ da Catedral.","epistemico":"fato","exemplos":[],"id":"fano_psl_168","memoria":"semantica","meta":{"autor":"Lord Index (Shawn R. Schiller)","dominio":"matematica","fonte":"A Catedral — Part VIII-A (jul 2026), ~/Imagens/catedral"},"origem":"wiki","palavras_chave":[],"sinonimos":[],"tipo":"conceito","titulo":"Fano e PSL(2,7) = 168"}
\section{Fano e PSL(2,7) = 168}
O plano de Fano (7 pontos, 7 retas, a geometria finita mínima) tem grupo de automorfismos Aut(Fano) \ensuremath{\cong} PSL(2,7), de ordem |PSL(2,7)| = 168. A semente de toda a cadeia \ensuremath{\sigma} da Catedral.
\prov{relações: parte\_de catedral\_sigma\_aritmetica; relaciona plano\_fano\_octonico; usa psl\_2\_7; confirma constante\_168; relaciona psl\_embeds\_g2}
\prov{proveniência: wiki \textperiodcentered{} A Catedral — Part VIII-A (jul 2026), \~{}/Imagens/catedral \textperiodcentered{} fato \textperiodcentered{} confiança 1.0 \textperiodcentered{} domínio matematica \textperiodcentered{} história 6 versões no jornal}

%CRISTAL {"arestas":[{"alvo":"sequencia","confianca":1.0,"epistemico":"fato","origem":"livro","rel":"especializa"},{"alvo":"razao_aurea","confianca":1.0,"epistemico":"fato","origem":"livro","rel":"relaciona"}],"confianca":1.0,"contraexemplos":[],"descricao":"1, 1, 2, 3, 5, 8, 13, …; Fₙ₊₁ = Fₙ + Fₙ₋₁. A razão Fₙ₊₁/Fₙ converge à razão áurea φ.","epistemico":"fato","exemplos":[],"id":"fibonacci","memoria":"semantica","meta":{"autor":"Gentil Lopes","descoberta":false,"dominio":"matematica","fonte":""},"origem":"manual","palavras_chave":["recorrência","áurea"],"sinonimos":[],"tipo":"conceito","titulo":"Sequência de Fibonacci"}
\section{Sequência de Fibonacci}
1, 1, 2, 3, 5, 8, 13, …; F\ensuremath{_{n+1}} = F\ensuremath{_{n}} + F\ensuremath{_{n-1}}. A razão F\ensuremath{_{n+1}}/F\ensuremath{_{n}} converge à razão áurea \ensuremath{\varphi}.
\prov{relações: especializa sequencia; relaciona razao\_aurea}
\prov{proveniência: manual \textperiodcentered{} fato \textperiodcentered{} confiança 1.0 \textperiodcentered{} domínio matematica \textperiodcentered{} história 6 versões no jornal}

%CRISTAL {"arestas":[{"alvo":"fibonacci","confianca":1.0,"epistemico":"fato","origem":"livro","rel":"especializa"},{"alvo":"utilitarismo","confianca":0.5,"epistemico":"fato","origem":"wiki","rel":"relaciona"},{"alvo":"vida-morte","confianca":0.5,"epistemico":"fato","origem":"wiki","rel":"relaciona"},{"alvo":"pell","confianca":0.5,"epistemico":"fato","origem":"wiki","rel":"relaciona"},{"alvo":"word","confianca":0.5,"epistemico":"fato","origem":"wiki","rel":"relaciona"},{"alvo":"virtualizacao","confianca":0.5,"epistemico":"fato","origem":"wiki","rel":"relaciona"},{"alvo":"readme","confianca":0.5,"epistemico":"fato","origem":"wiki","rel":"relaciona"},{"alvo":"koch_pell_verificacao","confianca":0.5,"epistemico":"fato","origem":"wiki","rel":"relaciona"},{"alvo":"gentil_12_definicao","confianca":0.5,"epistemico":"fato","origem":"wiki","rel":"relaciona"},{"alvo":"telemetria_shadow_producao","confianca":0.5,"epistemico":"fato","origem":"wiki","rel":"relaciona"},{"alvo":"vontade_de_poder","confianca":0.5,"epistemico":"fato","origem":"wiki","rel":"relaciona"},{"alvo":"sintaxe","confianca":0.5,"epistemico":"fato","origem":"wiki","rel":"relaciona"},{"alvo":"gentil_capitulo_9_programando_a_hp_prime","confianca":0.5,"epistemico":"fato","origem":"wiki","rel":"relaciona"},{"alvo":"pipeline_de_ingestao","confianca":0.5,"epistemico":"fato","origem":"wiki","rel":"relaciona"},{"alvo":"ula_da_broca_arquitetura_minima","confianca":0.5,"epistemico":"fato","origem":"wiki","rel":"relaciona"},{"alvo":"sequencia_pow_nativa","confianca":0.5,"epistemico":"fato","origem":"wiki","rel":"relaciona"},{"alvo":"incautos_enganados","confianca":0.5,"epistemico":"fato","origem":"wiki","rel":"relaciona"},{"alvo":"heap","confianca":0.5,"epistemico":"fato","origem":"wiki","rel":"relaciona"},{"alvo":"vetores","confianca":0.5,"epistemico":"fato","origem":"wiki","rel":"relaciona"},{"alvo":"proposition_escala_fasepropescala_xix","confianca":0.5,"epistemico":"fato","origem":"wiki","rel":"relaciona"},{"alvo":"literatura","confianca":0.5,"epistemico":"fato","origem":"wiki","rel":"relaciona"},{"alvo":"religiao","confianca":0.5,"epistemico":"fato","origem":"wiki","rel":"relaciona"},{"alvo":"painel_estelar","confianca":0.5,"epistemico":"fato","origem":"wiki","rel":"relaciona"},{"alvo":"kappa","confianca":0.5,"epistemico":"fato","origem":"wiki","rel":"relaciona"}],"confianca":1.0,"contraexemplos":[],"descricao":"Generalização em 2D/3D da sequência de Fibonacci (Lopes).","epistemico":"fato","exemplos":[],"id":"fibonacci_multidimensional","memoria":"semantica","meta":{"autor":"Gentil Lopes","descoberta":true,"dominio":"matematica","fonte":"corpus_gentil.txt [GENTIL:LIVRO:FIBONACCI_2D3D]"},"origem":"livro","palavras_chave":[],"sinonimos":[],"tipo":"conceito","titulo":"Fibonacci Multidimensional"}
\section{Fibonacci Multidimensional}
Generalização em 2D/3D da sequência de Fibonacci (Lopes).
\prov{relações: especializa fibonacci; relaciona utilitarismo; relaciona vida-morte; relaciona pell; relaciona word; relaciona virtualizacao; relaciona readme; relaciona koch\_pell\_verificacao; relaciona gentil\_12\_definicao; relaciona telemetria\_shadow\_producao; relaciona vontade\_de\_poder; relaciona sintaxe; relaciona gentil\_capitulo\_9\_programando\_a\_hp\_prime; relaciona pipeline\_de\_ingestao; relaciona ula\_da\_broca\_arquitetura\_minima; relaciona sequencia\_pow\_nativa; relaciona incautos\_enganados; relaciona heap; relaciona vetores; relaciona proposition\_escala\_fasepropescala\_xix; relaciona literatura; relaciona religiao; relaciona painel\_estelar; relaciona kappa}
\prov{proveniência: livro \textperiodcentered{} corpus\_gentil.txt [GENTIL:LIVRO:FIBONACCI\_2D3D] \textperiodcentered{} fato \textperiodcentered{} confiança 1.0 \textperiodcentered{} domínio matematica \textperiodcentered{} história 28 versões no jornal}

%CRISTAL {"arestas":[{"alvo":"reta_de_ouro","confianca":1.0,"epistemico":"fato","origem":"livro","rel":"relaciona"},{"alvo":"word","confianca":0.5,"epistemico":"fato","origem":"wiki","rel":"relaciona"},{"alvo":"vontade_de_poder","confianca":0.5,"epistemico":"fato","origem":"wiki","rel":"relaciona"},{"alvo":"virtualizacao","confianca":0.5,"epistemico":"fato","origem":"wiki","rel":"relaciona"},{"alvo":"progressao_geometrica","confianca":0.5,"epistemico":"fato","origem":"wiki","rel":"relaciona"},{"alvo":"relacao_com_dougherty","confianca":0.5,"epistemico":"fato","origem":"wiki","rel":"relaciona"}],"confianca":1.0,"contraexemplos":[],"descricao":"Transição exata entre o corpo discreto Z[φ]/p e o corpo contínuo de ouro; a reciprocidade mod 5 rege o período de Pisano ↔ a ordem relativística.","epistemico":"fato","exemplos":[],"id":"flip_discreto_continuo","memoria":"semantica","meta":{"autor":"Gentil Lopes","descoberta":true,"dominio":"matematica","fonte":"paper_3_regua_ouro_espelho.pdf, Teorema 6.1"},"origem":"paper","palavras_chave":[],"sinonimos":[],"tipo":"conceito","titulo":"Flip Discreto-Contínuo"}
\section{Flip Discreto-Contínuo}
Transição exata entre o corpo discreto Z[\ensuremath{\varphi}]/p e o corpo contínuo de ouro; a reciprocidade mod 5 rege o período de Pisano \ensuremath{\leftrightarrow} a ordem relativística.
\prov{relações: relaciona reta\_de\_ouro; relaciona word; relaciona vontade\_de\_poder; relaciona virtualizacao; relaciona progressao\_geometrica; relaciona relacao\_com\_dougherty}
\prov{proveniência: paper \textperiodcentered{} paper\_3\_regua\_ouro\_espelho.pdf, Teorema 6.1 \textperiodcentered{} fato \textperiodcentered{} confiança 1.0 \textperiodcentered{} domínio matematica \textperiodcentered{} história 9 versões no jornal}

%CRISTAL {"arestas":[{"alvo":"transformada_metalica","confianca":1.0,"epistemico":"fato","origem":"wiki","rel":"usa"},{"alvo":"lyapunov_da_orbita","confianca":1.0,"epistemico":"fato","origem":"wiki","rel":"relaciona"},{"alvo":"cadeias_alem_dos_metais","confianca":1.0,"epistemico":"fato","origem":"wiki","rel":"relaciona"},{"alvo":"reta_de_ouro","confianca":1.0,"epistemico":"fato","origem":"wiki","rel":"relaciona"},{"alvo":"curvatura_n_menos_1","confianca":1.0,"epistemico":"fato","origem":"wiki","rel":"usa"}],"confianca":1.0,"contraexemplos":[],"descricao":"O dupolinômio F=P_g−P_a (P_g=z^N geométrico/PG, P_a=Σc_i z^i aritmético/PA). O flip é a involução logarítmica u=ln z (real↔relativístico): troca PG↔PA (z^N vira N·u, linear). O gradiente F′ dá o FRACTAL de Newton (bacias das raízes): dim da fronteira ≈0.978 (grau 2, quase-suave) → ≈1.345 (grau 4, fractal). Descoberta: o flip leva σ_n,−1/σ_n a ±log σ_n = os expoentes de Lyapunov — o lado relativístico É o lado log/Lyapunov.","epistemico":"inferencia","exemplos":[],"id":"flip_duopolinomios","memoria":"semantica","meta":{"dim_fractal":{"grau4":1.345,"ouro":0.978,"tribonacci":1.341},"dominio":"matematica","fonte":"flip duopolinômios"},"origem":"wiki","palavras_chave":[],"sinonimos":[],"tipo":"conceito","titulo":"Flip dos Polinômios Duais"}
\section{Flip dos Polinômios Duais}
O dupolinômio F=P\_g\ensuremath{-}P\_a (P\_g=z\^{}N geométrico/PG, P\_a=\ensuremath{\Sigma}c\_i z\^{}i aritmético/PA). O flip é a involução logarítmica u=ln z (real\ensuremath{\leftrightarrow}relativístico): troca PG\ensuremath{\leftrightarrow}PA (z\^{}N vira N\textperiodcentered u, linear). O gradiente F\ensuremath{'} dá o FRACTAL de Newton (bacias das raízes): dim da fronteira \ensuremath{\approx}0.978 (grau 2, quase-suave) \ensuremath{\to} \ensuremath{\approx}1.345 (grau 4, fractal). Descoberta: o flip leva \ensuremath{\sigma}\_n,\ensuremath{-}1/\ensuremath{\sigma}\_n a $\pm$log \ensuremath{\sigma}\_n = os expoentes de Lyapunov — o lado relativístico É o lado log/Lyapunov.
\prov{relações: usa transformada\_metalica; relaciona lyapunov\_da\_orbita; relaciona cadeias\_alem\_dos\_metais; relaciona reta\_de\_ouro; usa curvatura\_n\_menos\_1}
\prov{proveniência: wiki \textperiodcentered{} flip duopolinômios \textperiodcentered{} inferencia \textperiodcentered{} confiança 1.0 \textperiodcentered{} domínio matematica \textperiodcentered{} história 6 versões no jornal}

%CRISTAL {"arestas":[{"alvo":"mecanica_quantica_mineral","confianca":1.0,"epistemico":"fato","origem":"wiki","rel":"parte_de"},{"alvo":"flip_duopolinomios","confianca":1.0,"epistemico":"fato","origem":"wiki","rel":"usa"},{"alvo":"vibracao_e_colapso_dois_tempos","confianca":1.0,"epistemico":"fato","origem":"wiki","rel":"explica"},{"alvo":"numero_de_salem","confianca":1.0,"epistemico":"fato","origem":"wiki","rel":"usa"}],"confianca":1.0,"contraexemplos":[],"descricao":"O flip u=ln z, que endireita as espirais, É a rotação de Wick entre os dois tempos: Re(u)=ln|z| é o eixo DISSIPATIVO (tempo imaginário, o λ, o cristal), Im(u)=arg z é o eixo UNITÁRIO (tempo real, a fase, a onda). Em |z|=1 (Salem/borda) u é puramente imaginário: evolução PURAMENTE unitária. O flip troca o relógio do calor pelo relógio da onda.","epistemico":"fato","exemplos":[],"id":"flip_e_rotacao_de_wick","memoria":"semantica","meta":{"dominio":"reta mineral","fonte":"mecânica quântica mineral"},"origem":"wiki","palavras_chave":[],"sinonimos":[],"tipo":"conceito","titulo":"O Flip é a Rotação de Wick"}
\section{O Flip é a Rotação de Wick}
O flip u=ln z, que endireita as espirais, É a rotação de Wick entre os dois tempos: Re(u)=ln|z| é o eixo DISSIPATIVO (tempo imaginário, o \ensuremath{\lambda}, o cristal), Im(u)=arg z é o eixo UNITÁRIO (tempo real, a fase, a onda). Em |z|=1 (Salem/borda) u é puramente imaginário: evolução PURAMENTE unitária. O flip troca o relógio do calor pelo relógio da onda.
\prov{relações: parte\_de mecanica\_quantica\_mineral; usa flip\_duopolinomios; explica vibracao\_e\_colapso\_dois\_tempos; usa numero\_de\_salem}
\prov{proveniência: wiki \textperiodcentered{} mecânica quântica mineral \textperiodcentered{} fato \textperiodcentered{} confiança 1.0 \textperiodcentered{} domínio reta mineral \textperiodcentered{} história 5 versões no jornal}

%CRISTAL {"arestas":[{"alvo":"flip_duopolinomios","confianca":1.0,"epistemico":"fato","origem":"wiki","rel":"especializa"},{"alvo":"geometria_do_caos","confianca":1.0,"epistemico":"fato","origem":"wiki","rel":"relaciona"},{"alvo":"transformada_estelar","confianca":1.0,"epistemico":"fato","origem":"wiki","rel":"relaciona"},{"alvo":"cadeias_alem_dos_metais","confianca":1.0,"epistemico":"fato","origem":"wiki","rel":"relaciona"},{"alvo":"transformada_dourada","confianca":1.0,"epistemico":"fato","origem":"wiki","rel":"relaciona"}],"confianca":1.0,"contraexemplos":[],"descricao":"As duas obras do Gentil são os dois lados de um flip: P.G. (progressão geométrica = multiplicar raízes, z^k = a espiral, lado relativístico) e P.A.m (progressão aritmética de ordem m, o livro do Gentil = o flip ln(z^k)=k·(λ,θ), a reta, lado euclidiano). O flip logarítmico ENDIREITA: espiral↦reta (inclinação θ/λ), círculo de Salem (ρ=1)↦reta vertical, real↦reta horizontal — exato na superfície de Riemann do log (a faixa (λ,θ)). Resíduo ~1e-14.","epistemico":"fato","exemplos":[],"id":"flip_pg_pa","memoria":"semantica","meta":{"dominio":"matematica","fonte":"álgebra da reta mineral"},"origem":"wiki","palavras_chave":[],"sinonimos":[],"tipo":"conceito","titulo":"O Flip P.G.↔P.A.m (espiral vira reta)"}
\section{O Flip P.G.\ensuremath{\leftrightarrow}P.A.m (espiral vira reta)}
As duas obras do Gentil são os dois lados de um flip: P.G. (progressão geométrica = multiplicar raízes, z\^{}k = a espiral, lado relativístico) e P.A.m (progressão aritmética de ordem m, o livro do Gentil = o flip ln(z\^{}k)=k\textperiodcentered (\ensuremath{\lambda},\ensuremath{\theta}), a reta, lado euclidiano). O flip logarítmico ENDIREITA: espiral\ensuremath{\mapsto}reta (inclinação \ensuremath{\theta}/\ensuremath{\lambda}), círculo de Salem (\ensuremath{\rho}=1)\ensuremath{\mapsto}reta vertical, real\ensuremath{\mapsto}reta horizontal — exato na superfície de Riemann do log (a faixa (\ensuremath{\lambda},\ensuremath{\theta})). Resíduo \~{}1e-14.
\prov{relações: especializa flip\_duopolinomios; relaciona geometria\_do\_caos; relaciona transformada\_estelar; relaciona cadeias\_alem\_dos\_metais; relaciona transformada\_dourada}
\prov{proveniência: wiki \textperiodcentered{} álgebra da reta mineral \textperiodcentered{} fato \textperiodcentered{} confiança 1.0 \textperiodcentered{} domínio matematica \textperiodcentered{} história 7 versões no jornal}

%CRISTAL {"arestas":[{"alvo":"pitagoras_dual","confianca":1.0,"epistemico":"fato","origem":"paper","rel":"usa"},{"alvo":"goldchain","confianca":1.0,"epistemico":"fato","origem":"paper","rel":"referencia"}],"confianca":1.0,"contraexemplos":[],"descricao":"Involução σ(s)=√(1−s²) que troca o regime circular (i²=−1) pelo hiperbólico (j²=+1), cosh↔cos; base geométrica da ponte de Wick.","epistemico":"fato","exemplos":[],"id":"flip_sigma","memoria":"semantica","meta":{"autor":"Gentil Lopes","descoberta":true,"dominio":"matematica","fonte":"hiper/gold/paper_0_engenharia_ouro.pdf"},"origem":"paper","palavras_chave":[],"sinonimos":[],"tipo":"conceito","titulo":"Flip σ (involução relativística)"}
\section{Flip \ensuremath{\sigma} (involução relativística)}
Involução \ensuremath{\sigma}(s)=\ensuremath{\sqrt{\,}}(1\ensuremath{-}s\ensuremath{^{2}}) que troca o regime circular (i\ensuremath{^{2}}=\ensuremath{-}1) pelo hiperbólico (j\ensuremath{^{2}}=+1), cosh\ensuremath{\leftrightarrow}cos; base geométrica da ponte de Wick.
\prov{relações: usa pitagoras\_dual; referencia goldchain}
\prov{proveniência: paper \textperiodcentered{} hiper/gold/paper\_0\_engenharia\_ouro.pdf \textperiodcentered{} fato \textperiodcentered{} confiança 1.0 \textperiodcentered{} domínio matematica \textperiodcentered{} história 3 versões no jornal}

%CRISTAL {"arestas":[{"alvo":"associador_calor","confianca":1.0,"epistemico":"fato","origem":"wiki","rel":"explica"},{"alvo":"entropia_log_phi","confianca":1.0,"epistemico":"fato","origem":"wiki","rel":"usa"},{"alvo":"floco_phi","confianca":1.0,"epistemico":"fato","origem":"wiki","rel":"especializa"},{"alvo":"minerar_e_resfriar","confianca":1.0,"epistemico":"fato","origem":"wiki","rel":"relaciona"}],"confianca":1.0,"contraexemplos":[],"descricao":"Φ(w)=koch_coord(w,gold(w)) armazena a informação NÃO-associativa produzida pela dissipação na passagem entre retas associativas — o defeito que não fecha na ressonância a+c²=1 'cristaliza no floco como calor dissipado numa passagem irreversível, mas armazenável'. É o análogo do setor glacial/nilpotente: o canal por onde a entropia escoa. Verificado no código (koch-memoria).","epistemico":"inferencia","exemplos":[],"id":"floco_calor_dissipado","memoria":"semantica","meta":{"autor":"broca-so","dominio":"fractal","fonte":"papers/floco_de_neve.tex §II; neuronio/parabola.py"},"origem":"wiki","palavras_chave":[],"sinonimos":[],"tipo":"conceito","titulo":"O floco Φ = calor dissipado não-associativo (armazenável)"}
\section{O floco \ensuremath{\Phi} = calor dissipado não-associativo (armazenável)}
\ensuremath{\Phi}(w)=koch\_coord(w,gold(w)) armazena a informação NÃO-associativa produzida pela dissipação na passagem entre retas associativas — o defeito que não fecha na ressonância a+c\ensuremath{^{2}}=1 'cristaliza no floco como calor dissipado numa passagem irreversível, mas armazenável'. É o análogo do setor glacial/nilpotente: o canal por onde a entropia escoa. Verificado no código (koch-memoria).
\prov{relações: explica associador\_calor; usa entropia\_log\_phi; especializa floco\_phi; relaciona minerar\_e\_resfriar}
\prov{proveniência: wiki \textperiodcentered{} papers/floco\_de\_neve.tex §II; neuronio/parabola.py \textperiodcentered{} inferencia \textperiodcentered{} confiança 1.0 \textperiodcentered{} domínio fractal \textperiodcentered{} história 5 versões no jornal}

%CRISTAL {"arestas":[{"alvo":"floco_phi","confianca":1.0,"epistemico":"fato","origem":"wiki","rel":"confirma"},{"alvo":"atrator_fractal_ouro","confianca":1.0,"epistemico":"fato","origem":"wiki","rel":"relaciona"}],"confianca":1.0,"contraexemplos":[],"descricao":"koch_floco_img_exp.c percorre o espectro do gato (cadeia Aₙ, silver/gold, Pell, grid 512, nonces) e deposita z² mod p, emitindo ficheiros reais: floco.ppm (radial 6-vias), mod_p.ppm (heatmap), cristal dendrítico, floco_gentil3 (Koch 3D+Hopf), report.json; koch_floco_atlas.py converte PPM→PNG. Resultado: 1009/1009 resíduos — 'o floco enche'.","epistemico":"fato","exemplos":[],"id":"floco_imagem_validacao","memoria":"semantica","meta":{"autor":"broca-so","dominio":"fractal","fonte":"ula/koch_floco_img_exp.c; ula/koch_floco_atlas.py"},"origem":"wiki","palavras_chave":[],"sinonimos":[],"tipo":"conceito","titulo":"Validação por Imagem (.ppm/.png, 1009/1009)"}
\section{Validação por Imagem (.ppm/.png, 1009/1009)}
koch\_floco\_img\_exp.c percorre o espectro do gato (cadeia A\ensuremath{_{n}}, silver/gold, Pell, grid 512, nonces) e deposita z\ensuremath{^{2}} mod p, emitindo ficheiros reais: floco.ppm (radial 6-vias), mod\_p.ppm (heatmap), cristal dendrítico, floco\_gentil3 (Koch 3D+Hopf), report.json; koch\_floco\_atlas.py converte PPM\ensuremath{\to}PNG. Resultado: 1009/1009 resíduos — 'o floco enche'.
\prov{relações: confirma floco\_phi; relaciona atrator\_fractal\_ouro}
\prov{proveniência: wiki \textperiodcentered{} ula/koch\_floco\_img\_exp.c; ula/koch\_floco\_atlas.py \textperiodcentered{} fato \textperiodcentered{} confiança 1.0 \textperiodcentered{} domínio fractal \textperiodcentered{} história 3 versões no jornal}

%CRISTAL {"arestas":[{"alvo":"floco_de_neve","confianca":1.0,"epistemico":"fato","origem":"wiki","rel":"explica"},{"alvo":"olho_de_koch","confianca":1.0,"epistemico":"fato","origem":"wiki","rel":"usa"},{"alvo":"colapso_koch_tiffany","confianca":1.0,"epistemico":"fato","origem":"wiki","rel":"relaciona"},{"alvo":"atrator_fractal_ouro","confianca":1.0,"epistemico":"fato","origem":"wiki","rel":"usa"},{"alvo":"computacao_fractal","confianca":1.0,"epistemico":"fato","origem":"wiki","rel":"parte_de"},{"alvo":"colapso_ping_pong_overlap","confianca":1.0,"epistemico":"fato","origem":"wiki","rel":"relaciona"}],"confianca":1.0,"contraexemplos":[],"descricao":"O floco Φ(w)=koch_coord(w,gold(w))=(endereço, atrator, z_koch, órbita), onde z_koch é a órbita z↦z² em ℤ[i]/pℤ[i] com p=1009; é determinístico em w, não-associativo e ORTOGONAL à reta associativa 𝓛(w). Não é a curva de Koch geométrica (área finita/perímetro infinito é só a metáfora que o nomeia).","epistemico":"fato","exemplos":[],"id":"floco_phi","memoria":"semantica","meta":{"autor":"broca-so","dominio":"fractal","fonte":"papers/floco_de_neve.tex §II; neuronio/parabola.py:koch_coordenadas"},"origem":"wiki","palavras_chave":[],"sinonimos":[],"tipo":"conceito","titulo":"O Floco Φ (dissipação não-associativa)"}
\section{O Floco \ensuremath{\Phi} (dissipação não-associativa)}
O floco \ensuremath{\Phi}(w)=koch\_coord(w,gold(w))=(endereço, atrator, z\_koch, órbita), onde z\_koch é a órbita z\ensuremath{\mapsto}z\ensuremath{^{2}} em \ensuremath{\mathbb{Z}}[i]/p\ensuremath{\mathbb{Z}}[i] com p=1009; é determinístico em w, não-associativo e ORTOGONAL à reta associativa \ensuremath{\mathcal{L}}(w). Não é a curva de Koch geométrica (área finita/perímetro infinito é só a metáfora que o nomeia).
\prov{relações: explica floco\_de\_neve; usa olho\_de\_koch; relaciona colapso\_koch\_tiffany; usa atrator\_fractal\_ouro; parte\_de computacao\_fractal; relaciona colapso\_ping\_pong\_overlap}
\prov{proveniência: wiki \textperiodcentered{} papers/floco\_de\_neve.tex §II; neuronio/parabola.py:koch\_coordenadas \textperiodcentered{} fato \textperiodcentered{} confiança 1.0 \textperiodcentered{} domínio fractal \textperiodcentered{} história 7 versões no jornal}

%CRISTAL {"arestas":[{"alvo":"floco_phi","confianca":1.0,"epistemico":"fato","origem":"wiki","rel":"usa"},{"alvo":"o_floco_como_selo_do_canal","confianca":1.0,"epistemico":"fato","origem":"wiki","rel":"explica"},{"alvo":"tecido_gex_floco","confianca":1.0,"epistemico":"fato","origem":"wiki","rel":"relaciona"}],"confianca":1.0,"contraexemplos":[],"descricao":"Na implementação viva o floco é um selo compacto de 16 bytes — FlocoSelo(atrator, z_koch_re, z_koch_im, endereço), empacotado '>4I' — 'selo conservado no canal E_ij', derivado exclusivamente de parabola(w,gold(w))+koch_coordenadas. Aprofunda o_floco_como_selo_do_canal.","epistemico":"fato","exemplos":[],"id":"floco_selo_16bytes","memoria":"semantica","meta":{"autor":"broca-so","dominio":"fractal","fonte":"neuronio/floco.py:FlocoSelo"},"origem":"wiki","palavras_chave":[],"sinonimos":[],"tipo":"conceito","titulo":"O Floco como Selo de 16 Bytes (FlocoSelo)"}
\section{O Floco como Selo de 16 Bytes (FlocoSelo)}
Na implementação viva o floco é um selo compacto de 16 bytes — FlocoSelo(atrator, z\_koch\_re, z\_koch\_im, endereço), empacotado '>4I' — 'selo conservado no canal E\_ij', derivado exclusivamente de parabola(w,gold(w))+koch\_coordenadas. Aprofunda o\_floco\_como\_selo\_do\_canal.
\prov{relações: usa floco\_phi; explica o\_floco\_como\_selo\_do\_canal; relaciona tecido\_gex\_floco}
\prov{proveniência: wiki \textperiodcentered{} neuronio/floco.py:FlocoSelo \textperiodcentered{} fato \textperiodcentered{} confiança 1.0 \textperiodcentered{} domínio fractal \textperiodcentered{} história 4 versões no jornal}

%CRISTAL {"arestas":[{"alvo":"agm","confianca":1.0,"epistemico":"fato","origem":"paper","rel":"usa"},{"alvo":"word","confianca":0.5,"epistemico":"fato","origem":"wiki","rel":"relaciona"},{"alvo":"virtualizacao","confianca":0.5,"epistemico":"fato","origem":"wiki","rel":"relaciona"}],"confianca":1.0,"contraexemplos":[],"descricao":"m ↦ Λ(m) = (1−√(1−m))/(1+√(1−m)): descida que leva |m|→0 quadraticamente; é o motor da convergência do AGM e computa ℘=π/AGM(1,√2).","epistemico":"fato","exemplos":[],"id":"fluxo_landen","memoria":"semantica","meta":{"autor":"Gentil Lopes","descoberta":false,"dominio":"matematica","fonte":"paper_33_metrica_viva.pdf"},"origem":"paper","palavras_chave":[],"sinonimos":[],"tipo":"conceito","titulo":"Fluxo de Landen"}
\section{Fluxo de Landen}
m \ensuremath{\mapsto} \ensuremath{\Lambda}(m) = (1\ensuremath{-}\ensuremath{\sqrt{\,}}(1\ensuremath{-}m))/(1+\ensuremath{\sqrt{\,}}(1\ensuremath{-}m)): descida que leva |m|\ensuremath{\to}0 quadraticamente; é o motor da convergência do AGM e computa \ensuremath{\wp}=\ensuremath{\pi}/AGM(1,\ensuremath{\sqrt{\,}}2).
\prov{relações: usa agm; relaciona word; relaciona virtualizacao}
\prov{proveniência: paper \textperiodcentered{} paper\_33\_metrica\_viva.pdf \textperiodcentered{} fato \textperiodcentered{} confiança 1.0 \textperiodcentered{} domínio matematica \textperiodcentered{} história 4 versões no jornal}

%CRISTAL {"arestas":[{"alvo":"anel_transcendental_rk","confianca":1.0,"epistemico":"fato","origem":"paper","rel":"parte_de"}],"confianca":1.0,"contraexemplos":[],"descricao":"Cada β∈R_k tem expressão única Σc_w·M_w numa base de Lyndon; verificável em tempo poli mas de tamanho exp(k) — a assimetria fácil-verificar/difícil-produzir.","epistemico":"fato","exemplos":[],"id":"forma_canonica_lyndon","memoria":"semantica","meta":{"autor":"Gentil Lopes","descoberta":false,"dominio":"matematica","fonte":"paper_PQ1_beta_octonic_hardness.pdf (Brown 2012)"},"origem":"paper","palavras_chave":[],"sinonimos":[],"tipo":"conceito","titulo":"Forma Canônica (Goncharov-Brown)"}
\section{Forma Canônica (Goncharov-Brown)}
Cada \ensuremath{\beta}\ensuremath{\in}R\_k tem expressão única \ensuremath{\Sigma}c\_w\textperiodcentered M\_w numa base de Lyndon; verificável em tempo poli mas de tamanho exp(k) — a assimetria fácil-verificar/difícil-produzir.
\prov{relações: parte\_de anel\_transcendental\_rk}
\prov{proveniência: paper \textperiodcentered{} paper\_PQ1\_beta\_octonic\_hardness.pdf (Brown 2012) \textperiodcentered{} fato \textperiodcentered{} confiança 1.0 \textperiodcentered{} domínio matematica \textperiodcentered{} história 2 versões no jornal}

%CRISTAL {"arestas":[{"alvo":"jogo_estrategico","confianca":1.0,"epistemico":"fato","origem":"wiki","rel":"usa"}],"confianca":1.0,"contraexemplos":[],"descricao":"O jogo como ÁRVORE de decisões sequenciais, com nós de escolha e informação. Resolve-se por indução retrógrada (do fim ao início) — a base do minimax e da busca em jogos.","epistemico":"fato","exemplos":[],"id":"forma_extensiva","memoria":"semantica","meta":{"dominio":"matematica","fonte":"teoria dos jogos"},"origem":"wiki","palavras_chave":[],"sinonimos":[],"tipo":"conceito","titulo":"Forma Extensiva (árvore do jogo)"}
\section{Forma Extensiva (árvore do jogo)}
O jogo como ÁRVORE de decisões sequenciais, com nós de escolha e informação. Resolve-se por indução retrógrada (do fim ao início) — a base do minimax e da busca em jogos.
\prov{relações: usa jogo\_estrategico}
\prov{proveniência: wiki \textperiodcentered{} teoria dos jogos \textperiodcentered{} fato \textperiodcentered{} confiança 1.0 \textperiodcentered{} domínio matematica \textperiodcentered{} história 2 versões no jornal}

%CRISTAL {"arestas":[{"alvo":"cantor_ouro","confianca":1.0,"epistemico":"fato","origem":"paper","rel":"explica"},{"alvo":"progressao_aritmetica_ordem_m","confianca":1.0,"epistemico":"fato","origem":"paper","rel":"usa"}],"confianca":1.0,"contraexemplos":[],"descricao":"Os 2N⁺¹ extremos da N-ésima etapa do Cantor dos terços são a(n)/3N, com a(n) obtido da expansão binária de n via PA-4D: converte binário em ternário restrito a 0,2.","epistemico":"fato","exemplos":[],"id":"formula_extremos_cantor","memoria":"semantica","meta":{"autor":"Gentil Lopes","descoberta":true,"dominio":"matematica","fonte":"paper_42_codificacao_racional.pdf, Teo.1.1"},"origem":"paper","palavras_chave":[],"sinonimos":[],"tipo":"conceito","titulo":"Fórmula dos Extremos do Cantor"}
\section{Fórmula dos Extremos do Cantor}
Os 2N\ensuremath{^{+1}} extremos da N-ésima etapa do Cantor dos terços são a(n)/3N, com a(n) obtido da expansão binária de n via PA-4D: converte binário em ternário restrito a 0,2.
\prov{relações: explica cantor\_ouro; usa progressao\_aritmetica\_ordem\_m}
\prov{proveniência: paper \textperiodcentered{} paper\_42\_codificacao\_racional.pdf, Teo.1.1 \textperiodcentered{} fato \textperiodcentered{} confiança 1.0 \textperiodcentered{} domínio matematica \textperiodcentered{} história 4 versões no jornal}

%CRISTAL {"arestas":[{"alvo":"lemniscata_pontryagin","confianca":1.0,"epistemico":"fato","origem":"paper","rel":"usa"},{"alvo":"word","confianca":0.5,"epistemico":"fato","origem":"wiki","rel":"relaciona"}],"confianca":1.0,"contraexemplos":[],"descricao":"A transformada de Fourier realiza a dualidade contínuo↔discreto exatamente: fluxo em S¹ ↔ espectro discreto em ℤ, reconstrução com erro de máquina.","epistemico":"fato","exemplos":[],"id":"fourier_reconstrucao","memoria":"semantica","meta":{"autor":"Gentil Lopes","descoberta":false,"dominio":"matematica","fonte":"paper_77_lemniscata.pdf, Teo.2.1"},"origem":"paper","palavras_chave":[],"sinonimos":[],"tipo":"conceito","titulo":"Fourier Liga os Laços"}
\section{Fourier Liga os Laços}
A transformada de Fourier realiza a dualidade contínuo\ensuremath{\leftrightarrow}discreto exatamente: fluxo em S\ensuremath{^{1}} \ensuremath{\leftrightarrow} espectro discreto em \ensuremath{\mathbb{Z}}, reconstrução com erro de máquina.
\prov{relações: usa lemniscata\_pontryagin; relaciona word}
\prov{proveniência: paper \textperiodcentered{} paper\_77\_lemniscata.pdf, Teo.2.1 \textperiodcentered{} fato \textperiodcentered{} confiança 1.0 \textperiodcentered{} domínio matematica \textperiodcentered{} história 3 versões no jornal}

%CRISTAL {"arestas":[{"alvo":"razao_aurea","confianca":1.0,"epistemico":"fato","origem":"livro","rel":"referencia"},{"alvo":"virtualizacao","confianca":0.5,"epistemico":"fato","origem":"wiki","rel":"relaciona"},{"alvo":"word","confianca":0.5,"epistemico":"fato","origem":"wiki","rel":"relaciona"},{"alvo":"vida-morte","confianca":0.5,"epistemico":"fato","origem":"wiki","rel":"relaciona"}],"confianca":1.0,"contraexemplos":[],"descricao":"φ = [1; 1, 1, 1, …] = 1 + 1/φ — a fração contínua 'mais irracional'; os convergentes são razões de Fibonacci.","epistemico":"fato","exemplos":[],"id":"fracao_continua_phi","memoria":"semantica","meta":{"autor":"Gentil Lopes","descoberta":false,"dominio":"matematica","fonte":"paper_2_reta_ouro.pdf, Teorema 6.1"},"origem":"paper","palavras_chave":[],"sinonimos":[],"tipo":"conceito","titulo":"Fração Contínua de φ"}
\section{Fração Contínua de \ensuremath{\varphi}}
\ensuremath{\varphi} = [1; 1, 1, 1, …] = 1 + 1/\ensuremath{\varphi} — a fração contínua 'mais irracional'; os convergentes são razões de Fibonacci.
\prov{relações: referencia razao\_aurea; relaciona virtualizacao; relaciona word; relaciona vida-morte}
\prov{proveniência: paper \textperiodcentered{} paper\_2\_reta\_ouro.pdf, Teorema 6.1 \textperiodcentered{} fato \textperiodcentered{} confiança 1.0 \textperiodcentered{} domínio matematica \textperiodcentered{} história 7 versões no jornal}

%CRISTAL {"arestas":[{"alvo":"algebras_hipercomplexas","confianca":1.0,"epistemico":"fato","origem":"paper","rel":"explica"},{"alvo":"matematica","confianca":1.0,"epistemico":"fato","origem":"paper","rel":"parte_de"}],"confianca":1.0,"contraexemplos":[],"descricao":"As combinações binárias de uma álgebra são informação estruturada sse ela é de divisão normada — ℝ,ℂ,ℍ,𝕆 (dim 1,2,4,8); a assinatura é a distribuição binomial de Pascal.","epistemico":"fato","exemplos":[],"id":"fronteira_hurwitz","memoria":"semantica","meta":{"autor":"Gentil Lopes","descoberta":false,"dominio":"matematica","fonte":"paper_21_irracionais.pdf, Teo.3.2"},"origem":"paper","palavras_chave":[],"sinonimos":[],"tipo":"conceito","titulo":"Fronteira de Hurwitz"}
\section{Fronteira de Hurwitz}
As combinações binárias de uma álgebra são informação estruturada sse ela é de divisão normada — \ensuremath{\mathbb{R}},\ensuremath{\mathbb{C}},\ensuremath{\mathbb{H}},\ensuremath{\mathbb{O}} (dim 1,2,4,8); a assinatura é a distribuição binomial de Pascal.
\prov{relações: explica algebras\_hipercomplexas; parte\_de matematica}
\prov{proveniência: paper \textperiodcentered{} paper\_21\_irracionais.pdf, Teo.3.2 \textperiodcentered{} fato \textperiodcentered{} confiança 1.0 \textperiodcentered{} domínio matematica \textperiodcentered{} história 3 versões no jornal}

%CRISTAL {"arestas":[{"alvo":"config_24_rings","confianca":1.0,"epistemico":"fato","origem":"paper","rel":"usa"}],"confianca":1.0,"contraexemplos":[],"descricao":"Conjectura: o funcional F_G(u)=⟨‖u−g·u‖²⟩ seria minimizado pela config de Klein, tornando-a atrator dinâmico de NS. Lopes REFUTA numericamente: 0/1000 ICs atingem F<0,01 — Klein é privilegiada ESTATICAMENTE (F=0), mas NÃO é atrator. Enfraquece a indecidibilidade.","epistemico":"hipotese","exemplos":[],"id":"funcional_simetria_klein","memoria":"semantica","meta":{"autor":"Gentil Lopes","descoberta":false,"dominio":"matematica","fonte":"paper_BN.pdf, Test_BB_22","refutado":true},"origem":"paper","palavras_chave":[],"sinonimos":[],"tipo":"conceito","titulo":"Klein como Atrator (REFUTADO)"}
\section{Klein como Atrator (REFUTADO)}
Conjectura: o funcional F\_G(u)=\ensuremath{\langle}\ensuremath{\|}u\ensuremath{-}g\textperiodcentered u\ensuremath{\|}\ensuremath{^{2}}\ensuremath{\rangle} seria minimizado pela config de Klein, tornando-a atrator dinâmico de NS. Lopes REFUTA numericamente: 0/1000 ICs atingem F<0,01 — Klein é privilegiada ESTATICAMENTE (F=0), mas NÃO é atrator. Enfraquece a indecidibilidade.
\prov{relações: usa config\_24\_rings}
\prov{proveniência: paper \textperiodcentered{} paper\_BN.pdf, Test\_BB\_22 \textperiodcentered{} hipotese \textperiodcentered{} confiança 1.0 \textperiodcentered{} domínio matematica \textperiodcentered{} história 3 versões no jornal}

%CRISTAL {"arestas":[{"alvo":"torre_hurwitz","confianca":1.0,"epistemico":"fato","origem":"paper","rel":"parte_de"},{"alvo":"fibracao_hopf_bloch","confianca":1.0,"epistemico":"fato","origem":"paper","rel":"explica"},{"alvo":"associador","confianca":1.0,"epistemico":"fato","origem":"paper","rel":"usa"}],"confianca":1.0,"contraexemplos":[],"descricao":"As fibrações de Hopf (S³→S², S⁷→S⁴, S¹⁵→S⁸) têm fibra que só é grupo até ℍ; em 𝕆 a não-associatividade quebra o grupo — o 'furo' dos octônios.","epistemico":"fato","exemplos":[],"id":"furo_hopf","memoria":"semantica","meta":{"autor":"Gentil Lopes","descoberta":true,"dominio":"matematica","fonte":"paper_50_octonios_furo.pdf, Teo.1.2"},"origem":"paper","palavras_chave":[],"sinonimos":[],"tipo":"conceito","titulo":"Furo de Hopf"}
\section{Furo de Hopf}
As fibrações de Hopf (S\ensuremath{^{3}}\ensuremath{\to}S\ensuremath{^{2}}, S\ensuremath{^{7}}\ensuremath{\to}S\ensuremath{^{4}}, S\ensuremath{^{15}}\ensuremath{\to}S\ensuremath{^{8}}) têm fibra que só é grupo até \ensuremath{\mathbb{H}}; em \ensuremath{\mathbb{O}} a não-associatividade quebra o grupo — o 'furo' dos octônios.
\prov{relações: parte\_de torre\_hurwitz; explica fibracao\_hopf\_bloch; usa associador}
\prov{proveniência: paper \textperiodcentered{} paper\_50\_octonios\_furo.pdf, Teo.1.2 \textperiodcentered{} fato \textperiodcentered{} confiança 1.0 \textperiodcentered{} domínio matematica \textperiodcentered{} história 4 versões no jornal}

%CRISTAL {"arestas":[{"alvo":"cuspide_pisot","confianca":1.0,"epistemico":"fato","origem":"paper","rel":"parte_de"},{"alvo":"word","confianca":0.5,"epistemico":"fato","origem":"wiki","rel":"relaciona"},{"alvo":"virtualizacao","confianca":0.5,"epistemico":"fato","origem":"wiki","rel":"relaciona"}],"confianca":1.0,"contraexemplos":[],"descricao":"A diferença λ₁−λ₂ dos autovalores da matriz de transição: controla a taxa de relaxação ao equilíbrio (difusão); tende a zero na criticidade.","epistemico":"fato","exemplos":[],"id":"gap_espectral","memoria":"semantica","meta":{"autor":"Gentil Lopes","descoberta":false,"dominio":"matematica","fonte":"paper_70_ponto_difusao.pdf"},"origem":"paper","palavras_chave":[],"sinonimos":[],"tipo":"conceito","titulo":"Gap Espectral"}
\section{Gap Espectral}
A diferença \ensuremath{\lambda}\ensuremath{_{1}}\ensuremath{-}\ensuremath{\lambda}\ensuremath{_{2}} dos autovalores da matriz de transição: controla a taxa de relaxação ao equilíbrio (difusão); tende a zero na criticidade.
\prov{relações: parte\_de cuspide\_pisot; relaciona word; relaciona virtualizacao}
\prov{proveniência: paper \textperiodcentered{} paper\_70\_ponto\_difusao.pdf \textperiodcentered{} fato \textperiodcentered{} confiança 1.0 \textperiodcentered{} domínio matematica \textperiodcentered{} história 4 versões no jornal}

%CRISTAL {"arestas":[{"alvo":"numeros_metalicos","confianca":1.0,"epistemico":"fato","origem":"paper","rel":"usa"},{"alvo":"colapso","confianca":1.0,"epistemico":"fato","origem":"paper","rel":"explica"}],"confianca":1.0,"contraexemplos":[],"descricao":"O vão √Δ=√(n²+4) entre os autovalores de Aₙ parte o sistema em ℝ⊕ℝ; sempre positivo (os metais nunca perdem o gap), o que torna a reconstrução simbólica possível.","epistemico":"fato","exemplos":[],"id":"gap_metalico","memoria":"semantica","meta":{"autor":"Gentil Lopes","descoberta":true,"dominio":"matematica","fonte":"corpo_tropical.pdf, Teo.2.5"},"origem":"paper","palavras_chave":[],"sinonimos":[],"tipo":"conceito","titulo":"Gap Metálico √(n²+4)"}
\section{Gap Metálico \ensuremath{\sqrt{\,}}(n\ensuremath{^{2}}+4)}
O vão \ensuremath{\sqrt{\,}}\ensuremath{\Delta}=\ensuremath{\sqrt{\,}}(n\ensuremath{^{2}}+4) entre os autovalores de A\ensuremath{_{n}} parte o sistema em \ensuremath{\mathbb{R}}\ensuremath{\oplus}\ensuremath{\mathbb{R}}; sempre positivo (os metais nunca perdem o gap), o que torna a reconstrução simbólica possível.
\prov{relações: usa numeros\_metalicos; explica colapso}
\prov{proveniência: paper \textperiodcentered{} corpo\_tropical.pdf, Teo.2.5 \textperiodcentered{} fato \textperiodcentered{} confiança 1.0 \textperiodcentered{} domínio matematica \textperiodcentered{} história 3 versões no jornal}

%CRISTAL {"arestas":[{"alvo":"algebras_hipercomplexas","confianca":1.0,"epistemico":"fato","origem":"paper","rel":"especializa"},{"alvo":"word","confianca":0.5,"epistemico":"fato","origem":"wiki","rel":"relaciona"}],"confianca":1.0,"contraexemplos":[],"descricao":"Álgebra de matrizes 2×2 sobre o corpo 𝔽ₚ, não-comutativa e com divisores de zero; faixas de Möbius Mₖ vivem em SL₂(ℤ).","epistemico":"fato","exemplos":[],"id":"garrafa_klein_algebra","memoria":"semantica","meta":{"autor":"Gentil Lopes","descoberta":true,"dominio":"matematica","fonte":"paper_14_garrafa_klein.pdf"},"origem":"paper","palavras_chave":[],"sinonimos":[],"tipo":"conceito","titulo":"Garrafa de Klein 𝕶 = M₂(𝔽ₚ)"}
\section{Garrafa de Klein \ensuremath{\mathfrak{K}} = M\ensuremath{_{2}}(\ensuremath{\mathbb{F}}\ensuremath{_{p}})}
Álgebra de matrizes 2$\times$2 sobre o corpo \ensuremath{\mathbb{F}}\ensuremath{_{p}}, não-comutativa e com divisores de zero; faixas de Möbius M\ensuremath{_{k}} vivem em SL\ensuremath{_{2}}(\ensuremath{\mathbb{Z}}).
\prov{relações: especializa algebras\_hipercomplexas; relaciona word}
\prov{proveniência: paper \textperiodcentered{} paper\_14\_garrafa\_klein.pdf \textperiodcentered{} fato \textperiodcentered{} confiança 1.0 \textperiodcentered{} domínio matematica \textperiodcentered{} história 3 versões no jornal}

%CRISTAL {"arestas":[{"alvo":"mecanica_quantica_mineral","confianca":1.0,"epistemico":"fato","origem":"wiki","rel":"parte_de"},{"alvo":"gato","confianca":1.0,"epistemico":"fato","origem":"wiki","rel":"usa"},{"alvo":"numero_de_salem","confianca":1.0,"epistemico":"fato","origem":"wiki","rel":"relaciona"},{"alvo":"numeros_de_pisot","confianca":1.0,"epistemico":"fato","origem":"wiki","rel":"relaciona"}],"confianca":1.0,"contraexemplos":[],"descricao":"O fato fundador: o gato A_n=[[n,1],[1,0]] é SIMÉTRICO, logo HERMITIANO — um observável. Seus autovalores são REAIS: σ_n e −1/σ_n. Os metais são os NÍVEIS DE ENERGIA, e a reta mineral inteira (a sequência dos σ_n) é um espectro. O estado fundamental é o autovetor do Pisot (σ_n). Verificado: Hermitiano para todo n, E+=σ_n.","epistemico":"fato","exemplos":[],"id":"gato_e_o_hamiltoniano","memoria":"semantica","meta":{"dominio":"reta mineral","fonte":"mecânica quântica mineral"},"origem":"wiki","palavras_chave":[],"sinonimos":[],"tipo":"conceito","titulo":"O Gato é o Hamiltoniano (Hermitiano)"}
\section{O Gato é o Hamiltoniano (Hermitiano)}
O fato fundador: o gato A\_n=[[n,1],[1,0]] é SIMÉTRICO, logo HERMITIANO — um observável. Seus autovalores são REAIS: \ensuremath{\sigma}\_n e \ensuremath{-}1/\ensuremath{\sigma}\_n. Os metais são os NÍVEIS DE ENERGIA, e a reta mineral inteira (a sequência dos \ensuremath{\sigma}\_n) é um espectro. O estado fundamental é o autovetor do Pisot (\ensuremath{\sigma}\_n). Verificado: Hermitiano para todo n, E+=\ensuremath{\sigma}\_n.
\prov{relações: parte\_de mecanica\_quantica\_mineral; usa gato; relaciona numero\_de\_salem; relaciona numeros\_de\_pisot}
\prov{proveniência: wiki \textperiodcentered{} mecânica quântica mineral \textperiodcentered{} fato \textperiodcentered{} confiança 1.0 \textperiodcentered{} domínio reta mineral \textperiodcentered{} história 5 versões no jornal}

%CRISTAL {"arestas":[{"alvo":"gentil_capitulo_1_sequencias_aritmeticas_de_ordem_m","confianca":1.0,"epistemico":"fato","origem":"gentil/Novas_Sequencias_Aritmeticas_e_Geometric.pdf","rel":"parte_de"},{"alvo":"utilitarismo","confianca":0.5,"epistemico":"fato","origem":"wiki","rel":"relaciona"},{"alvo":"vida-morte","confianca":0.5,"epistemico":"fato","origem":"wiki","rel":"relaciona"},{"alvo":"virtualizacao","confianca":0.5,"epistemico":"fato","origem":"wiki","rel":"relaciona"},{"alvo":"word","confianca":0.5,"epistemico":"fato","origem":"wiki","rel":"relaciona"},{"alvo":"readme","confianca":0.5,"epistemico":"fato","origem":"wiki","rel":"relaciona"},{"alvo":"telemetria_shadow_producao","confianca":0.5,"epistemico":"fato","origem":"wiki","rel":"relaciona"},{"alvo":"sequencia_pow_nativa","confianca":0.5,"epistemico":"fato","origem":"wiki","rel":"relaciona"},{"alvo":"pow_taxa_de_sucesso_do_bit_broca","confianca":0.5,"epistemico":"fato","origem":"wiki","rel":"relaciona"},{"alvo":"pell","confianca":0.5,"epistemico":"fato","origem":"wiki","rel":"relaciona"},{"alvo":"religiao","confianca":0.5,"epistemico":"fato","origem":"wiki","rel":"relaciona"}],"confianca":1.0,"contraexemplos":[],"descricao":"Secção 1.2 — p.16.","epistemico":"fato","exemplos":[],"id":"gentil_12_definicao","memoria":"semantica","meta":{"arquivo":"gentil/Novas_Sequencias_Aritmeticas_e_Geometric.pdf","ciencia":"teoria_dos_numeros","dominio":"matematica","nivel":4,"pagina":16,"secao":"1.2","tipo_no":"paper"},"origem":"gentil/Novas_Sequencias_Aritmeticas_e_Geometric.pdf","palavras_chave":[],"sinonimos":[],"tipo":"conceito","titulo":"gentil 1.2 — Definição"}
\section{gentil 1.2 — Definição}
Secção 1.2 — p.16.
\prov{relações: parte\_de gentil\_capitulo\_1\_sequencias\_aritmeticas\_de\_ordem\_m; relaciona utilitarismo; relaciona vida-morte; relaciona virtualizacao; relaciona word; relaciona readme; relaciona telemetria\_shadow\_producao; relaciona sequencia\_pow\_nativa; relaciona pow\_taxa\_de\_sucesso\_do\_bit\_broca; relaciona pell; relaciona religiao}
\prov{proveniência: gentil/Novas\_Sequencias\_Aritmeticas\_e\_Geometric.pdf \textperiodcentered{} fato \textperiodcentered{} confiança 1.0 \textperiodcentered{} domínio matematica \textperiodcentered{} história 12 versões no jornal}

%CRISTAL {"arestas":[{"alvo":"gentil_capitulo_1_sequencias_aritmeticas_de_ordem_m","confianca":1.0,"epistemico":"fato","origem":"gentil/Novas_Sequencias_Aritmeticas_e_Geometric.pdf","rel":"parte_de"},{"alvo":"virtualizacao","confianca":0.5,"epistemico":"fato","origem":"wiki","rel":"relaciona"},{"alvo":"transcricao","confianca":0.5,"epistemico":"fato","origem":"wiki","rel":"relaciona"},{"alvo":"utilitarismo","confianca":0.5,"epistemico":"fato","origem":"wiki","rel":"relaciona"},{"alvo":"vida-morte","confianca":0.5,"epistemico":"fato","origem":"wiki","rel":"relaciona"},{"alvo":"scheduler","confianca":0.5,"epistemico":"fato","origem":"wiki","rel":"relaciona"},{"alvo":"pow_metal_ressonancia_val","confianca":0.5,"epistemico":"fato","origem":"wiki","rel":"relaciona"},{"alvo":"vetores","confianca":0.5,"epistemico":"fato","origem":"wiki","rel":"relaciona"},{"alvo":"resultado_completo","confianca":0.5,"epistemico":"fato","origem":"wiki","rel":"relaciona"}],"confianca":1.0,"contraexemplos":[],"descricao":"Secção 1.5 — p.42.","epistemico":"fato","exemplos":[],"id":"gentil_15_propriedade_fundamental_de_uma_pam","memoria":"semantica","meta":{"arquivo":"gentil/Novas_Sequencias_Aritmeticas_e_Geometric.pdf","ciencia":"teoria_dos_numeros","dominio":"matematica","nivel":4,"pagina":42,"secao":"1.5","tipo_no":"paper"},"origem":"gentil/Novas_Sequencias_Aritmeticas_e_Geometric.pdf","palavras_chave":[],"sinonimos":[],"tipo":"conceito","titulo":"gentil 1.5 — Propriedade fundamental de uma P.A.m"}
\section{gentil 1.5 — Propriedade fundamental de uma P.A.m}
Secção 1.5 — p.42.
\prov{relações: parte\_de gentil\_capitulo\_1\_sequencias\_aritmeticas\_de\_ordem\_m; relaciona virtualizacao; relaciona transcricao; relaciona utilitarismo; relaciona vida-morte; relaciona scheduler; relaciona pow\_metal\_ressonancia\_val; relaciona vetores; relaciona resultado\_completo}
\prov{proveniência: gentil/Novas\_Sequencias\_Aritmeticas\_e\_Geometric.pdf \textperiodcentered{} fato \textperiodcentered{} confiança 1.0 \textperiodcentered{} domínio matematica \textperiodcentered{} história 10 versões no jornal}

%CRISTAL {"arestas":[{"alvo":"gentil_capitulo_1_sequencias_aritmeticas_de_ordem_m","confianca":1.0,"epistemico":"fato","origem":"gentil/Novas_Sequencias_Aritmeticas_e_Geometric.pdf","rel":"parte_de"},{"alvo":"utilitarismo","confianca":0.5,"epistemico":"fato","origem":"wiki","rel":"relaciona"},{"alvo":"ruido","confianca":0.5,"epistemico":"fato","origem":"wiki","rel":"relaciona"},{"alvo":"main","confianca":0.5,"epistemico":"fato","origem":"wiki","rel":"relaciona"},{"alvo":"scheduler","confianca":0.5,"epistemico":"fato","origem":"wiki","rel":"relaciona"},{"alvo":"resultado_completo","confianca":0.5,"epistemico":"fato","origem":"wiki","rel":"relaciona"},{"alvo":"pedagogia","confianca":0.5,"epistemico":"fato","origem":"wiki","rel":"relaciona"},{"alvo":"teoria_do_valor","confianca":0.5,"epistemico":"fato","origem":"wiki","rel":"relaciona"},{"alvo":"resultado_anterior_identidade_nao_cobertura","confianca":0.5,"epistemico":"fato","origem":"wiki","rel":"relaciona"},{"alvo":"vida-morte","confianca":0.5,"epistemico":"fato","origem":"wiki","rel":"relaciona"},{"alvo":"pitagorismo","confianca":0.5,"epistemico":"fato","origem":"wiki","rel":"relaciona"},{"alvo":"tipos_de_interpretante","confianca":0.5,"epistemico":"fato","origem":"wiki","rel":"relaciona"},{"alvo":"vetores","confianca":0.5,"epistemico":"fato","origem":"wiki","rel":"relaciona"},{"alvo":"prova_da_maquina_espectral","confianca":0.5,"epistemico":"fato","origem":"wiki","rel":"relaciona"}],"confianca":1.0,"contraexemplos":[],"descricao":"Secção 1.6 — p.46.","epistemico":"fato","exemplos":[],"id":"gentil_16_soma_dos_termos_de_uma_pam","memoria":"semantica","meta":{"arquivo":"gentil/Novas_Sequencias_Aritmeticas_e_Geometric.pdf","ciencia":"teoria_dos_numeros","dominio":"matematica","nivel":4,"pagina":46,"secao":"1.6","tipo_no":"paper"},"origem":"gentil/Novas_Sequencias_Aritmeticas_e_Geometric.pdf","palavras_chave":[],"sinonimos":[],"tipo":"conceito","titulo":"gentil 1.6 — Soma dos termos de uma P.A.m"}
\section{gentil 1.6 — Soma dos termos de uma P.A.m}
Secção 1.6 — p.46.
\prov{relações: parte\_de gentil\_capitulo\_1\_sequencias\_aritmeticas\_de\_ordem\_m; relaciona utilitarismo; relaciona ruido; relaciona main; relaciona scheduler; relaciona resultado\_completo; relaciona pedagogia; relaciona teoria\_do\_valor; relaciona resultado\_anterior\_identidade\_nao\_cobertura; relaciona vida-morte; relaciona pitagorismo; relaciona tipos\_de\_interpretante; relaciona vetores; relaciona prova\_da\_maquina\_espectral}
\prov{proveniência: gentil/Novas\_Sequencias\_Aritmeticas\_e\_Geometric.pdf \textperiodcentered{} fato \textperiodcentered{} confiança 1.0 \textperiodcentered{} domínio matematica \textperiodcentered{} história 15 versões no jornal}

%CRISTAL {"arestas":[{"alvo":"gentil_novas_sequencias_aritmeticas_geometricas","confianca":1.0,"epistemico":"fato","origem":"gentil/Novas_Sequencias_Aritmeticas_e_Geometric.pdf","rel":"parte_de"},{"alvo":"virtualizacao","confianca":0.5,"epistemico":"fato","origem":"wiki","rel":"relaciona"},{"alvo":"word","confianca":0.5,"epistemico":"fato","origem":"wiki","rel":"relaciona"},{"alvo":"vida-morte","confianca":0.5,"epistemico":"fato","origem":"wiki","rel":"relaciona"},{"alvo":"gentil_capitulo_3_sequencias_geometricas_de_ordem_m","confianca":0.5,"epistemico":"fato","origem":"wiki","rel":"relaciona"},{"alvo":"misticismo_nao_dual","confianca":0.5,"epistemico":"fato","origem":"wiki","rel":"relaciona"}],"confianca":1.0,"contraexemplos":[],"descricao":"Capítulo 1 — p.13.","epistemico":"fato","exemplos":[],"id":"gentil_capitulo_1_sequencias_aritmeticas_de_ordem_m","memoria":"semantica","meta":{"arquivo":"gentil/Novas_Sequencias_Aritmeticas_e_Geometric.pdf","capitulo":"1","ciencia":"teoria_dos_numeros","dominio":"matematica","nivel":4,"pagina":13,"tipo_no":"paper"},"origem":"gentil/Novas_Sequencias_Aritmeticas_e_Geometric.pdf","palavras_chave":[],"sinonimos":["capítulo 1"],"tipo":"conceito","titulo":"gentil capítulo 1 — Sequências aritméticas de ordem m"}
\section{gentil capítulo 1 — Sequências aritméticas de ordem m}
Capítulo 1 — p.13.
\prov{sinónimos: capítulo 1}
\prov{relações: parte\_de gentil\_novas\_sequencias\_aritmeticas\_geometricas; relaciona virtualizacao; relaciona word; relaciona vida-morte; relaciona gentil\_capitulo\_3\_sequencias\_geometricas\_de\_ordem\_m; relaciona misticismo\_nao\_dual}
\prov{proveniência: gentil/Novas\_Sequencias\_Aritmeticas\_e\_Geometric.pdf \textperiodcentered{} fato \textperiodcentered{} confiança 1.0 \textperiodcentered{} domínio matematica \textperiodcentered{} história 7 versões no jornal}

%CRISTAL {"arestas":[{"alvo":"gentil_novas_sequencias_aritmeticas_geometricas","confianca":1.0,"epistemico":"fato","origem":"gentil/Novas_Sequencias_Aritmeticas_e_Geometric.pdf","rel":"parte_de"},{"alvo":"virtualizacao","confianca":0.5,"epistemico":"fato","origem":"wiki","rel":"relaciona"},{"alvo":"word","confianca":0.5,"epistemico":"fato","origem":"wiki","rel":"relaciona"},{"alvo":"vontade_de_poder","confianca":0.5,"epistemico":"fato","origem":"wiki","rel":"relaciona"}],"confianca":1.0,"contraexemplos":[],"descricao":"Capítulo 2 — p.81.","epistemico":"fato","exemplos":[],"id":"gentil_capitulo_2_somas_e_diferencas_de_ordem_m","memoria":"semantica","meta":{"arquivo":"gentil/Novas_Sequencias_Aritmeticas_e_Geometric.pdf","capitulo":"2","ciencia":"teoria_dos_numeros","dominio":"matematica","nivel":4,"pagina":81,"tipo_no":"paper"},"origem":"gentil/Novas_Sequencias_Aritmeticas_e_Geometric.pdf","palavras_chave":[],"sinonimos":["capítulo 2"],"tipo":"conceito","titulo":"gentil capítulo 2 — Somas e Diferenças de ordem m"}
\section{gentil capítulo 2 — Somas e Diferenças de ordem m}
Capítulo 2 — p.81.
\prov{sinónimos: capítulo 2}
\prov{relações: parte\_de gentil\_novas\_sequencias\_aritmeticas\_geometricas; relaciona virtualizacao; relaciona word; relaciona vontade\_de\_poder}
\prov{proveniência: gentil/Novas\_Sequencias\_Aritmeticas\_e\_Geometric.pdf \textperiodcentered{} fato \textperiodcentered{} confiança 1.0 \textperiodcentered{} domínio matematica \textperiodcentered{} história 5 versões no jornal}

%CRISTAL {"arestas":[{"alvo":"gentil_novas_sequencias_aritmeticas_geometricas","confianca":1.0,"epistemico":"fato","origem":"gentil/Novas_Sequencias_Aritmeticas_e_Geometric.pdf","rel":"parte_de"},{"alvo":"word","confianca":0.5,"epistemico":"fato","origem":"wiki","rel":"relaciona"},{"alvo":"misticismo_nao_dual","confianca":0.5,"epistemico":"fato","origem":"wiki","rel":"relaciona"},{"alvo":"vida-morte","confianca":0.5,"epistemico":"fato","origem":"wiki","rel":"relaciona"},{"alvo":"virtualizacao","confianca":0.5,"epistemico":"fato","origem":"wiki","rel":"relaciona"}],"confianca":1.0,"contraexemplos":[],"descricao":"Capítulo 3 — p.127.","epistemico":"fato","exemplos":[],"id":"gentil_capitulo_3_sequencias_geometricas_de_ordem_m","memoria":"semantica","meta":{"arquivo":"gentil/Novas_Sequencias_Aritmeticas_e_Geometric.pdf","capitulo":"3","ciencia":"teoria_dos_numeros","dominio":"matematica","nivel":4,"pagina":127,"tipo_no":"paper"},"origem":"gentil/Novas_Sequencias_Aritmeticas_e_Geometric.pdf","palavras_chave":[],"sinonimos":["capítulo 3"],"tipo":"conceito","titulo":"gentil capítulo 3 — Sequências geométricas de ordem m"}
\section{gentil capítulo 3 — Sequências geométricas de ordem m}
Capítulo 3 — p.127.
\prov{sinónimos: capítulo 3}
\prov{relações: parte\_de gentil\_novas\_sequencias\_aritmeticas\_geometricas; relaciona word; relaciona misticismo\_nao\_dual; relaciona vida-morte; relaciona virtualizacao}
\prov{proveniência: gentil/Novas\_Sequencias\_Aritmeticas\_e\_Geometric.pdf \textperiodcentered{} fato \textperiodcentered{} confiança 1.0 \textperiodcentered{} domínio matematica \textperiodcentered{} história 6 versões no jornal}

%CRISTAL {"arestas":[{"alvo":"gentil_novas_sequencias_aritmeticas_geometricas","confianca":1.0,"epistemico":"fato","origem":"gentil/Novas_Sequencias_Aritmeticas_e_Geometric.pdf","rel":"parte_de"},{"alvo":"vida-morte","confianca":0.5,"epistemico":"fato","origem":"wiki","rel":"relaciona"},{"alvo":"virtualizacao","confianca":0.5,"epistemico":"fato","origem":"wiki","rel":"relaciona"},{"alvo":"telemetria_shadow_producao","confianca":0.5,"epistemico":"fato","origem":"wiki","rel":"relaciona"},{"alvo":"word","confianca":0.5,"epistemico":"fato","origem":"wiki","rel":"relaciona"}],"confianca":1.0,"contraexemplos":[],"descricao":"Capítulo 4 — p.189.","epistemico":"fato","exemplos":[],"id":"gentil_capitulo_4_sequencias_especiais","memoria":"semantica","meta":{"arquivo":"gentil/Novas_Sequencias_Aritmeticas_e_Geometric.pdf","capitulo":"4","ciencia":"teoria_dos_numeros","dominio":"matematica","nivel":4,"pagina":189,"tipo_no":"paper"},"origem":"gentil/Novas_Sequencias_Aritmeticas_e_Geometric.pdf","palavras_chave":[],"sinonimos":["capítulo 4"],"tipo":"conceito","titulo":"gentil capítulo 4 — Sequências Especiais"}
\section{gentil capítulo 4 — Sequências Especiais}
Capítulo 4 — p.189.
\prov{sinónimos: capítulo 4}
\prov{relações: parte\_de gentil\_novas\_sequencias\_aritmeticas\_geometricas; relaciona vida-morte; relaciona virtualizacao; relaciona telemetria\_shadow\_producao; relaciona word}
\prov{proveniência: gentil/Novas\_Sequencias\_Aritmeticas\_e\_Geometric.pdf \textperiodcentered{} fato \textperiodcentered{} confiança 1.0 \textperiodcentered{} domínio matematica \textperiodcentered{} história 6 versões no jornal}

%CRISTAL {"arestas":[{"alvo":"gentil_novas_sequencias_aritmeticas_geometricas","confianca":1.0,"epistemico":"fato","origem":"gentil/Novas_Sequencias_Aritmeticas_e_Geometric.pdf","rel":"parte_de"},{"alvo":"gentil_capitulo_6_progressao_geometrica_bidimensional","confianca":0.5,"epistemico":"fato","origem":"wiki","rel":"relaciona"}],"confianca":1.0,"contraexemplos":[],"descricao":"Capítulo 5 — p.227.","epistemico":"fato","exemplos":[],"id":"gentil_capitulo_5_progressao_aritmetica_bidimensional","memoria":"semantica","meta":{"arquivo":"gentil/Novas_Sequencias_Aritmeticas_e_Geometric.pdf","capitulo":"5","ciencia":"teoria_dos_numeros","dominio":"matematica","nivel":4,"pagina":227,"tipo_no":"paper"},"origem":"gentil/Novas_Sequencias_Aritmeticas_e_Geometric.pdf","palavras_chave":[],"sinonimos":["capítulo 5"],"tipo":"conceito","titulo":"gentil capítulo 5 — Progressão aritmética bidimensional"}
\section{gentil capítulo 5 — Progressão aritmética bidimensional}
Capítulo 5 — p.227.
\prov{sinónimos: capítulo 5}
\prov{relações: parte\_de gentil\_novas\_sequencias\_aritmeticas\_geometricas; relaciona gentil\_capitulo\_6\_progressao\_geometrica\_bidimensional}
\prov{proveniência: gentil/Novas\_Sequencias\_Aritmeticas\_e\_Geometric.pdf \textperiodcentered{} fato \textperiodcentered{} confiança 1.0 \textperiodcentered{} domínio matematica \textperiodcentered{} história 3 versões no jornal}

%CRISTAL {"arestas":[{"alvo":"gentil_novas_sequencias_aritmeticas_geometricas","confianca":1.0,"epistemico":"fato","origem":"gentil/Novas_Sequencias_Aritmeticas_e_Geometric.pdf","rel":"parte_de"},{"alvo":"virtualizacao","confianca":0.5,"epistemico":"fato","origem":"wiki","rel":"relaciona"},{"alvo":"word","confianca":0.5,"epistemico":"fato","origem":"wiki","rel":"relaciona"}],"confianca":1.0,"contraexemplos":[],"descricao":"Capítulo 6 — p.299.","epistemico":"fato","exemplos":[],"id":"gentil_capitulo_6_progressao_geometrica_bidimensional","memoria":"semantica","meta":{"arquivo":"gentil/Novas_Sequencias_Aritmeticas_e_Geometric.pdf","capitulo":"6","ciencia":"teoria_dos_numeros","dominio":"matematica","nivel":4,"pagina":299,"tipo_no":"paper"},"origem":"gentil/Novas_Sequencias_Aritmeticas_e_Geometric.pdf","palavras_chave":[],"sinonimos":["capítulo 6"],"tipo":"conceito","titulo":"gentil capítulo 6 — Progressão geométrica bidimensional"}
\section{gentil capítulo 6 — Progressão geométrica bidimensional}
Capítulo 6 — p.299.
\prov{sinónimos: capítulo 6}
\prov{relações: parte\_de gentil\_novas\_sequencias\_aritmeticas\_geometricas; relaciona virtualizacao; relaciona word}
\prov{proveniência: gentil/Novas\_Sequencias\_Aritmeticas\_e\_Geometric.pdf \textperiodcentered{} fato \textperiodcentered{} confiança 1.0 \textperiodcentered{} domínio matematica \textperiodcentered{} história 4 versões no jornal}

%CRISTAL {"arestas":[{"alvo":"gentil_novas_sequencias_aritmeticas_geometricas","confianca":1.0,"epistemico":"fato","origem":"gentil/Novas_Sequencias_Aritmeticas_e_Geometric.pdf","rel":"parte_de"},{"alvo":"virtualizacao","confianca":0.5,"epistemico":"fato","origem":"wiki","rel":"relaciona"},{"alvo":"word","confianca":0.5,"epistemico":"fato","origem":"wiki","rel":"relaciona"},{"alvo":"vida-morte","confianca":0.5,"epistemico":"fato","origem":"wiki","rel":"relaciona"}],"confianca":1.0,"contraexemplos":[],"descricao":"Capítulo 7 — p.337.","epistemico":"fato","exemplos":[],"id":"gentil_capitulo_7_progressao_aritmetica_tridimensional","memoria":"semantica","meta":{"arquivo":"gentil/Novas_Sequencias_Aritmeticas_e_Geometric.pdf","capitulo":"7","ciencia":"teoria_dos_numeros","dominio":"matematica","nivel":4,"pagina":337,"tipo_no":"paper"},"origem":"gentil/Novas_Sequencias_Aritmeticas_e_Geometric.pdf","palavras_chave":[],"sinonimos":["capítulo 7"],"tipo":"conceito","titulo":"gentil capítulo 7 — Progressão aritmética tridimensional"}
\section{gentil capítulo 7 — Progressão aritmética tridimensional}
Capítulo 7 — p.337.
\prov{sinónimos: capítulo 7}
\prov{relações: parte\_de gentil\_novas\_sequencias\_aritmeticas\_geometricas; relaciona virtualizacao; relaciona word; relaciona vida-morte}
\prov{proveniência: gentil/Novas\_Sequencias\_Aritmeticas\_e\_Geometric.pdf \textperiodcentered{} fato \textperiodcentered{} confiança 1.0 \textperiodcentered{} domínio matematica \textperiodcentered{} história 5 versões no jornal}

%CRISTAL {"arestas":[{"alvo":"gentil_novas_sequencias_aritmeticas_geometricas","confianca":1.0,"epistemico":"fato","origem":"gentil/Novas_Sequencias_Aritmeticas_e_Geometric.pdf","rel":"parte_de"},{"alvo":"vida-morte","confianca":0.5,"epistemico":"fato","origem":"wiki","rel":"relaciona"},{"alvo":"virtualizacao","confianca":0.5,"epistemico":"fato","origem":"wiki","rel":"relaciona"},{"alvo":"word","confianca":0.5,"epistemico":"fato","origem":"wiki","rel":"relaciona"}],"confianca":1.0,"contraexemplos":[],"descricao":"Capítulo 8 — p.431.","epistemico":"fato","exemplos":[],"id":"gentil_capitulo_8_mais_aplicacoes","memoria":"semantica","meta":{"arquivo":"gentil/Novas_Sequencias_Aritmeticas_e_Geometric.pdf","capitulo":"8","ciencia":"teoria_dos_numeros","dominio":"matematica","nivel":4,"pagina":431,"tipo_no":"paper"},"origem":"gentil/Novas_Sequencias_Aritmeticas_e_Geometric.pdf","palavras_chave":[],"sinonimos":["capítulo 8"],"tipo":"conceito","titulo":"gentil capítulo 8 — Mais aplicações"}
\section{gentil capítulo 8 — Mais aplicações}
Capítulo 8 — p.431.
\prov{sinónimos: capítulo 8}
\prov{relações: parte\_de gentil\_novas\_sequencias\_aritmeticas\_geometricas; relaciona vida-morte; relaciona virtualizacao; relaciona word}
\prov{proveniência: gentil/Novas\_Sequencias\_Aritmeticas\_e\_Geometric.pdf \textperiodcentered{} fato \textperiodcentered{} confiança 1.0 \textperiodcentered{} domínio matematica \textperiodcentered{} história 5 versões no jornal}

%CRISTAL {"arestas":[{"alvo":"gentil_novas_sequencias_aritmeticas_geometricas","confianca":1.0,"epistemico":"fato","origem":"gentil/Novas_Sequencias_Aritmeticas_e_Geometric.pdf","rel":"parte_de"},{"alvo":"vida-morte","confianca":0.5,"epistemico":"fato","origem":"wiki","rel":"relaciona"},{"alvo":"virtualizacao","confianca":0.5,"epistemico":"fato","origem":"wiki","rel":"relaciona"},{"alvo":"pell","confianca":0.5,"epistemico":"fato","origem":"wiki","rel":"relaciona"},{"alvo":"word","confianca":0.5,"epistemico":"fato","origem":"wiki","rel":"relaciona"},{"alvo":"heap","confianca":0.5,"epistemico":"fato","origem":"wiki","rel":"relaciona"},{"alvo":"sequencia_pow_nativa","confianca":0.5,"epistemico":"fato","origem":"wiki","rel":"relaciona"},{"alvo":"vetores","confianca":0.5,"epistemico":"fato","origem":"wiki","rel":"relaciona"},{"alvo":"readme","confianca":0.5,"epistemico":"fato","origem":"wiki","rel":"relaciona"}],"confianca":1.0,"contraexemplos":[],"descricao":"Capítulo 9 — p.481.","epistemico":"fato","exemplos":[],"id":"gentil_capitulo_9_programando_a_hp_prime","memoria":"semantica","meta":{"arquivo":"gentil/Novas_Sequencias_Aritmeticas_e_Geometric.pdf","capitulo":"9","ciencia":"teoria_dos_numeros","dominio":"matematica","nivel":4,"pagina":481,"tipo_no":"paper"},"origem":"gentil/Novas_Sequencias_Aritmeticas_e_Geometric.pdf","palavras_chave":[],"sinonimos":["capítulo 9"],"tipo":"conceito","titulo":"gentil capítulo 9 — Programando a HP Prime"}
\section{gentil capítulo 9 — Programando a HP Prime}
Capítulo 9 — p.481.
\prov{sinónimos: capítulo 9}
\prov{relações: parte\_de gentil\_novas\_sequencias\_aritmeticas\_geometricas; relaciona vida-morte; relaciona virtualizacao; relaciona pell; relaciona word; relaciona heap; relaciona sequencia\_pow\_nativa; relaciona vetores; relaciona readme}
\prov{proveniência: gentil/Novas\_Sequencias\_Aritmeticas\_e\_Geometric.pdf \textperiodcentered{} fato \textperiodcentered{} confiança 1.0 \textperiodcentered{} domínio matematica \textperiodcentered{} história 10 versões no jornal}

%CRISTAL {"arestas":[{"alvo":"matematica","confianca":0.95,"epistemico":"fato","origem":"paper","rel":"parte_de"}],"confianca":1.0,"contraexemplos":[],"descricao":"Dada uma sequência f : N →R definimos o Produto de ordem m de f pela seguinte fórmula de recorrência:            ∆ 0 f(n) = f(n), ∆ j f(1) = cj , cj ̸= 0, j = 1, 2, . . . , m. ∆ m f(n) = ∆ m f(n −1) × ∆ m−1 f(n −1) , m ≥1, n ≥2. Onde ∆ j f(1) = cj são constantes arbitrárias dadas. A fórmula seguinte nos fornece o m-ésimo Produto sem que tenhamos necessi- dade da fórmula de recorrência. 147","epistemico":"fato","exemplos":[],"id":"gentil_definicao_10_dada_uma_sequencia_f_n_r_definimos_o_produto_de_ordem_m_d","memoria":"semantica","meta":{"arquivo":"gentil/Novas_Sequencias_Aritmeticas_e_Geometric.pdf","ciencia":"teoria_dos_numeros","dominio":"matematica","nivel":4,"numero":"10","tipo_no":"paper","tipo_teorema":"definição"},"origem":"gentil/Novas_Sequencias_Aritmeticas_e_Geometric.pdf","palavras_chave":[],"sinonimos":[],"tipo":"conceito","titulo":"gentil Definição 10 — Dada uma sequência f : N →R definimos o Produto de ordem m d"}
\section{gentil Definição 10 — Dada uma sequência f : N \ensuremath{\to}R definimos o Produto de ordem m d}
Dada uma sequência f : N \ensuremath{\to}R definimos o Produto de ordem m de f pela seguinte fórmula de recorrência: \{           \ensuremath{\Delta} 0 f(n) = f(n), \ensuremath{\Delta} j f(1) = cj , cj \ensuremath{}= 0, j = 1, 2, . . . , m. \ensuremath{\Delta} m f(n) = \ensuremath{\Delta} m f(n \ensuremath{-}1) $\times$ \ensuremath{\Delta} m\ensuremath{-}1 f(n \ensuremath{-}1) , m \ensuremath{\ge}1, n \ensuremath{\ge}2. Onde \ensuremath{\Delta} j f(1) = cj são constantes arbitrárias dadas. A fórmula seguinte nos fornece o m-ésimo Produto sem que tenhamos necessi- dade da fórmula de recorrência. 147
\prov{relações: parte\_de matematica}
\prov{proveniência: gentil/Novas\_Sequencias\_Aritmeticas\_e\_Geometric.pdf \textperiodcentered{} fato \textperiodcentered{} confiança 1.0 \textperiodcentered{} domínio matematica \textperiodcentered{} história 3 versões no jornal}

%CRISTAL {"arestas":[{"alvo":"matematica","confianca":0.95,"epistemico":"fato","origem":"paper","rel":"parte_de"}],"confianca":1.0,"contraexemplos":[],"descricao":"Polinômio dual) Dada uma sequência de números reais positivos (am, am−1, . . . , a2, a1, a0) , ai > 0, ∀i ; am ̸= 1 consideremos a função p: R →R dada por p(x) = axm m × axm−1 m−1 × · · · × ax2 2 × ax 1 × a0 a função p é denominada função dupolinomial ou dupolinômio associada(o) à sequência dada. Exemplos: (a) p(x) = 2x3 · 3x2 · 4−x · 1 2 ⇔ \u00002, 3, 1 4, 1 2 \u0001 (b) p(x) = 2x4 · 3x2 · 5 ⇔ ( 2, 1, 3, 1, 5 ) (c) p(x) = 3x2−1 ⇔ \u00003, 1, 1 3 \u0001 160","epistemico":"fato","exemplos":[],"id":"gentil_definicao_12_polinomio_dual_dada_uma_sequencia_de_numeros_reais_positivo","memoria":"semantica","meta":{"arquivo":"gentil/Novas_Sequencias_Aritmeticas_e_Geometric.pdf","ciencia":"teoria_dos_numeros","dominio":"matematica","nivel":4,"numero":"12","tipo_no":"paper","tipo_teorema":"definição"},"origem":"gentil/Novas_Sequencias_Aritmeticas_e_Geometric.pdf","palavras_chave":[],"sinonimos":[],"tipo":"conceito","titulo":"gentil Definição 12 — Polinômio dual) Dada uma sequência de números reais positivo"}
\section{gentil Definição 12 — Polinômio dual) Dada uma sequência de números reais positivo}
Polinômio dual) Dada uma sequência de números reais positivos (am, am\ensuremath{-}1, . . . , a2, a1, a0) , ai > 0, \ensuremath{\forall}i ; am \ensuremath{}= 1 consideremos a função p: R \ensuremath{\to}R dada por p(x) = axm m $\times$ axm\ensuremath{-}1 m\ensuremath{-}1 $\times$ \textperiodcentered  \textperiodcentered  \textperiodcentered  $\times$ ax2 2 $\times$ ax 1 $\times$ a0 a função p é denominada função dupolinomial ou dupolinômio associada(o) à sequência dada. Exemplos: (a) p(x) = 2x3 \textperiodcentered  3x2 \textperiodcentered  4\ensuremath{-}x \textperiodcentered  1 2 \ensuremath{\Leftrightarrow} 2, 3, 1 4, 1 2  (b) p(x) = 2x4 \textperiodcentered  3x2 \textperiodcentered  5 \ensuremath{\Leftrightarrow} ( 2, 1, 3, 1, 5 ) (c) p(x) = 3x2\ensuremath{-}1 \ensuremath{\Leftrightarrow} 3, 1, 1 3  160
\prov{relações: parte\_de matematica}
\prov{proveniência: gentil/Novas\_Sequencias\_Aritmeticas\_e\_Geometric.pdf \textperiodcentered{} fato \textperiodcentered{} confiança 1.0 \textperiodcentered{} domínio matematica \textperiodcentered{} história 3 versões no jornal}

%CRISTAL {"arestas":[{"alvo":"gentil_novas_sequencias_aritmeticas_geometricas","confianca":1.0,"epistemico":"fato","origem":"gentil/Novas_Sequencias_Aritmeticas_e_Geometric.pdf","rel":"define"},{"alvo":"gentil_capitulo_1_sequencias_aritmeticas_de_ordem_m","confianca":1.0,"epistemico":"fato","origem":"gentil/Novas_Sequencias_Aritmeticas_e_Geometric.pdf","rel":"define"}],"confianca":1.0,"contraexemplos":[],"descricao":"Grau) Seja p(x) = axm m × axm−1 m−1 × · · · × ax2 2 × ax 1 × a0 um dupolinômio não identicamente 1. Chama-se grau de p, e representa-se por ∂p ou gr p, o número natural q tal que aq ̸= 1 e ai = 1 para todo i > q. ∂p = q ⇐⇒    aq ̸= 1; ai = 1, ∀i > q. Os dupolinômios do exemplo anterior têm grau 3, 4 e 2 respectivamente. Pelo","epistemico":"fato","exemplos":[],"id":"gentil_definicao_13_grau_seja_px_axm_m_axm1_m1_ax2_2_ax_1","memoria":"semantica","meta":{"arquivo":"gentil/Novas_Sequencias_Aritmeticas_e_Geometric.pdf","ciencia":"teoria_dos_numeros","dominio":"matematica","nivel":4,"numero":"13","tipo_no":"paper","tipo_teorema":"definição"},"origem":"gentil/Novas_Sequencias_Aritmeticas_e_Geometric.pdf","palavras_chave":[],"sinonimos":[],"tipo":"conceito","titulo":"gentil Definição 13 — Grau) Seja p(x) = axm m × axm−1 m−1 × · · · × ax2 2 × ax 1 ×"}
\section{gentil Definição 13 — Grau) Seja p(x) = axm m $\times$ axm\ensuremath{-}1 m\ensuremath{-}1 $\times$ \textperiodcentered  \textperiodcentered  \textperiodcentered  $\times$ ax2 2 $\times$ ax 1 $\times$}
Grau) Seja p(x) = axm m $\times$ axm\ensuremath{-}1 m\ensuremath{-}1 $\times$ \textperiodcentered  \textperiodcentered  \textperiodcentered  $\times$ ax2 2 $\times$ ax 1 $\times$ a0 um dupolinômio não identicamente 1. Chama-se grau de p, e representa-se por \ensuremath{\partial}p ou gr p, o número natural q tal que aq \ensuremath{}= 1 e ai = 1 para todo i > q. \ensuremath{\partial}p = q \ensuremath{\Leftarrow}\ensuremath{\Rightarrow} \{   aq \ensuremath{}= 1; ai = 1, \ensuremath{\forall}i > q. Os dupolinômios do exemplo anterior têm grau 3, 4 e 2 respectivamente. Pelo
\prov{relações: define gentil\_novas\_sequencias\_aritmeticas\_geometricas; define gentil\_capitulo\_1\_sequencias\_aritmeticas\_de\_ordem\_m}
\prov{proveniência: gentil/Novas\_Sequencias\_Aritmeticas\_e\_Geometric.pdf \textperiodcentered{} fato \textperiodcentered{} confiança 1.0 \textperiodcentered{} domínio matematica \textperiodcentered{} história 4 versões no jornal}

%CRISTAL {"arestas":[{"alvo":"matematica","confianca":0.95,"epistemico":"fato","origem":"paper","rel":"parte_de"}],"confianca":1.0,"contraexemplos":[],"descricao":"Chama-se progressão aritmética bidimensional (PA-2D) uma sequên- cia dupla dada pela seguinte fórmula de recorrência:          a11 = a a1j = a1(j−1) + r1, j ≥2; aij = a(i−1)j + r2, i ≥2, j ≥1. Onde: a11 = a, r1 e r2 são constantes dadas. Vejamos a ideia que está por trás desta","epistemico":"fato","exemplos":[],"id":"gentil_definicao_16_chama-se_progressao_aritmetica_bidimensional_pa-2d_uma_seq","memoria":"semantica","meta":{"arquivo":"gentil/Novas_Sequencias_Aritmeticas_e_Geometric.pdf","ciencia":"teoria_dos_numeros","dominio":"matematica","nivel":4,"numero":"16","tipo_no":"paper","tipo_teorema":"definição"},"origem":"gentil/Novas_Sequencias_Aritmeticas_e_Geometric.pdf","palavras_chave":[],"sinonimos":[],"tipo":"conceito","titulo":"gentil Definição 16 — Chama-se progressão aritmética bidimensional (PA-2D) uma seq"}
\section{gentil Definição 16 — Chama-se progressão aritmética bidimensional (PA-2D) uma seq}
Chama-se progressão aritmética bidimensional (PA-2D) uma sequên- cia dupla dada pela seguinte fórmula de recorrência: \{         a11 = a a1j = a1(j\ensuremath{-}1) + r1, j \ensuremath{\ge}2; aij = a(i\ensuremath{-}1)j + r2, i \ensuremath{\ge}2, j \ensuremath{\ge}1. Onde: a11 = a, r1 e r2 são constantes dadas. Vejamos a ideia que está por trás desta
\prov{relações: parte\_de matematica}
\prov{proveniência: gentil/Novas\_Sequencias\_Aritmeticas\_e\_Geometric.pdf \textperiodcentered{} fato \textperiodcentered{} confiança 1.0 \textperiodcentered{} domínio matematica \textperiodcentered{} história 3 versões no jornal}

%CRISTAL {"arestas":[{"alvo":"gentil_novas_sequencias_aritmeticas_geometricas","confianca":1.0,"epistemico":"fato","origem":"gentil/Novas_Sequencias_Aritmeticas_e_Geometric.pdf","rel":"define"},{"alvo":"gentil_capitulo_1_sequencias_aritmeticas_de_ordem_m","confianca":1.0,"epistemico":"fato","origem":"gentil/Novas_Sequencias_Aritmeticas_e_Geometric.pdf","rel":"define"}],"confianca":1.0,"contraexemplos":[],"descricao":"p 230), temos: a12 = a11 + r1 a13 = a12 + r1 a14 = a13 + r1 · · · · · · · · · · · · · · · · · · a1j = a1(j−1) + r1 Somando essas j −1 igualdades e fazendo os cancelamentos apropriados, temos: a1j = a11 + (j −1)r1 (2) Substituindo (2) e (1), temos: aij = a11 + (j −1)r1 + (i −1)r2 Com o que podemos enunciar o seguinte","epistemico":"fato","exemplos":[],"id":"gentil_definicao_16_p_230_temos_a12_a11_r1_a13_a12_r1_a14_a13_r1","memoria":"semantica","meta":{"arquivo":"gentil/Novas_Sequencias_Aritmeticas_e_Geometric.pdf","ciencia":"teoria_dos_numeros","dominio":"matematica","nivel":4,"numero":"16","tipo_no":"paper","tipo_teorema":"definição"},"origem":"gentil/Novas_Sequencias_Aritmeticas_e_Geometric.pdf","palavras_chave":[],"sinonimos":[],"tipo":"conceito","titulo":"gentil Definição 16 — p 230), temos: a12 = a11 + r1 a13 = a12 + r1 a14 = a13 + r1"}
\section{gentil Definição 16 — p 230), temos: a12 = a11 + r1 a13 = a12 + r1 a14 = a13 + r1}
p 230), temos: a12 = a11 + r1 a13 = a12 + r1 a14 = a13 + r1 \textperiodcentered  \textperiodcentered  \textperiodcentered  \textperiodcentered  \textperiodcentered  \textperiodcentered  \textperiodcentered  \textperiodcentered  \textperiodcentered  \textperiodcentered  \textperiodcentered  \textperiodcentered  \textperiodcentered  \textperiodcentered  \textperiodcentered  \textperiodcentered  \textperiodcentered  \textperiodcentered  a1j = a1(j\ensuremath{-}1) + r1 Somando essas j \ensuremath{-}1 igualdades e fazendo os cancelamentos apropriados, temos: a1j = a11 + (j \ensuremath{-}1)r1 (2) Substituindo (2) e (1), temos: aij = a11 + (j \ensuremath{-}1)r1 + (i \ensuremath{-}1)r2 Com o que podemos enunciar o seguinte
\prov{relações: define gentil\_novas\_sequencias\_aritmeticas\_geometricas; define gentil\_capitulo\_1\_sequencias\_aritmeticas\_de\_ordem\_m}
\prov{proveniência: gentil/Novas\_Sequencias\_Aritmeticas\_e\_Geometric.pdf \textperiodcentered{} fato \textperiodcentered{} confiança 1.0 \textperiodcentered{} domínio matematica \textperiodcentered{} história 4 versões no jornal}

%CRISTAL {"arestas":[{"alvo":"gentil_novas_sequencias_aritmeticas_geometricas","confianca":1.0,"epistemico":"fato","origem":"gentil/Novas_Sequencias_Aritmeticas_e_Geometric.pdf","rel":"define"},{"alvo":"gentil_capitulo_1_sequencias_aritmeticas_de_ordem_m","confianca":1.0,"epistemico":"fato","origem":"gentil/Novas_Sequencias_Aritmeticas_e_Geometric.pdf","rel":"define"}],"confianca":1.0,"contraexemplos":[],"descricao":"p 230), temos: a2j = a1j + r2 a3j = a2j + r2 a4j = a3j + r2 · · · · · · · · · · · · · · · · · · aij = a(i−1)j + r2 Somando essas i −1 igualdades e fazendo os cancelamentos apropriados, temos: aij = a1j + (i −1)r2 (1) Da segunda equação da","epistemico":"fato","exemplos":[],"id":"gentil_definicao_16_p_230_temos_a2j_a1j_r2_a3j_a2j_r2_a4j_a3j_r2","memoria":"semantica","meta":{"arquivo":"gentil/Novas_Sequencias_Aritmeticas_e_Geometric.pdf","ciencia":"teoria_dos_numeros","dominio":"matematica","nivel":4,"numero":"16","tipo_no":"paper","tipo_teorema":"definição"},"origem":"gentil/Novas_Sequencias_Aritmeticas_e_Geometric.pdf","palavras_chave":[],"sinonimos":[],"tipo":"conceito","titulo":"gentil Definição 16 — p 230), temos: a2j = a1j + r2 a3j = a2j + r2 a4j = a3j + r2"}
\section{gentil Definição 16 — p 230), temos: a2j = a1j + r2 a3j = a2j + r2 a4j = a3j + r2}
p 230), temos: a2j = a1j + r2 a3j = a2j + r2 a4j = a3j + r2 \textperiodcentered  \textperiodcentered  \textperiodcentered  \textperiodcentered  \textperiodcentered  \textperiodcentered  \textperiodcentered  \textperiodcentered  \textperiodcentered  \textperiodcentered  \textperiodcentered  \textperiodcentered  \textperiodcentered  \textperiodcentered  \textperiodcentered  \textperiodcentered  \textperiodcentered  \textperiodcentered  aij = a(i\ensuremath{-}1)j + r2 Somando essas i \ensuremath{-}1 igualdades e fazendo os cancelamentos apropriados, temos: aij = a1j + (i \ensuremath{-}1)r2 (1) Da segunda equação da
\prov{relações: define gentil\_novas\_sequencias\_aritmeticas\_geometricas; define gentil\_capitulo\_1\_sequencias\_aritmeticas\_de\_ordem\_m}
\prov{proveniência: gentil/Novas\_Sequencias\_Aritmeticas\_e\_Geometric.pdf \textperiodcentered{} fato \textperiodcentered{} confiança 1.0 \textperiodcentered{} domínio matematica \textperiodcentered{} história 4 versões no jornal}

%CRISTAL {"arestas":[{"alvo":"matematica","confianca":0.95,"epistemico":"fato","origem":"paper","rel":"parte_de"}],"confianca":1.0,"contraexemplos":[],"descricao":"Dizemos que uma sequência simples é equivalente a uma PA-2D quando existe uma PA-2D (com N colunas) que ao ser linearizada coincide com a sequência simples. Exemplos: (a) A sequência abaixo é equivalente a uma PA-2D: 1 1 2 2 3 3 4 4 . . . De fato, a sequência dupla abaixo 1 1 2 2 3 3 4 4 . . . . . . é uma PA-2D que, ao ser linearizada, coincide com a sequência simples dada. (b) A sequência dos ´ımpares positivos é equivalente a uma PA-2D: 1 3 5 7 9 11 13 15 . . . De fato, a sequência dupla abaix","epistemico":"fato","exemplos":[],"id":"gentil_definicao_17_dizemos_que_uma_sequencia_simples_e_equivalente_a_uma_pa-2d","memoria":"semantica","meta":{"arquivo":"gentil/Novas_Sequencias_Aritmeticas_e_Geometric.pdf","ciencia":"teoria_dos_numeros","dominio":"matematica","nivel":4,"numero":"17","tipo_no":"paper","tipo_teorema":"definição"},"origem":"gentil/Novas_Sequencias_Aritmeticas_e_Geometric.pdf","palavras_chave":[],"sinonimos":[],"tipo":"conceito","titulo":"gentil Definição 17 — Dizemos que uma sequência simples é equivalente a uma PA-2D"}
\section{gentil Definição 17 — Dizemos que uma sequência simples é equivalente a uma PA-2D}
Dizemos que uma sequência simples é equivalente a uma PA-2D quando existe uma PA-2D (com N colunas) que ao ser linearizada coincide com a sequência simples. Exemplos: (a) A sequência abaixo é equivalente a uma PA-2D: 1 1 2 2 3 3 4 4 . . . De fato, a sequência dupla abaixo 1 1 2 2 3 3 4 4 . . . . . . é uma PA-2D que, ao ser linearizada, coincide com a sequência simples dada. (b) A sequência dos '\ensuremath{\imath}mpares positivos é equivalente a uma PA-2D: 1 3 5 7 9 11 13 15 . . . De fato, a sequência dupla abaix
\prov{relações: parte\_de matematica}
\prov{proveniência: gentil/Novas\_Sequencias\_Aritmeticas\_e\_Geometric.pdf \textperiodcentered{} fato \textperiodcentered{} confiança 1.0 \textperiodcentered{} domínio matematica \textperiodcentered{} história 3 versões no jornal}

%CRISTAL {"arestas":[{"alvo":"matematica","confianca":0.95,"epistemico":"fato","origem":"paper","rel":"parte_de"}],"confianca":1.0,"contraexemplos":[],"descricao":"Dizemos que uma sequência dupla é equivalente a uma P A. quando, ao ser linearizada, a sequência simples resultante está em progressão aritmé- tica. Exemplos: A PA-2D abaixo: 1 1 2 2 3 3 4 4 . . . . . . não é equivalente a uma P.A. pois, ao ser linearizada, resulta 1 1 2 2 3 3 4 4 . . . Quando uma PA-2D é equivalente a uma P.A., diremos que é redut´ıvel, caso contrário, que é irredut´ıvel. 260","epistemico":"fato","exemplos":[],"id":"gentil_definicao_18_dizemos_que_uma_sequencia_dupla_e_equivalente_a_uma_p_a_qua","memoria":"semantica","meta":{"arquivo":"gentil/Novas_Sequencias_Aritmeticas_e_Geometric.pdf","ciencia":"teoria_dos_numeros","dominio":"matematica","nivel":4,"numero":"18","tipo_no":"paper","tipo_teorema":"definição"},"origem":"gentil/Novas_Sequencias_Aritmeticas_e_Geometric.pdf","palavras_chave":[],"sinonimos":[],"tipo":"conceito","titulo":"gentil Definição 18 — Dizemos que uma sequência dupla é equivalente a uma P A. qua"}
\section{gentil Definição 18 — Dizemos que uma sequência dupla é equivalente a uma P A. qua}
Dizemos que uma sequência dupla é equivalente a uma P A. quando, ao ser linearizada, a sequência simples resultante está em progressão aritmé- tica. Exemplos: A PA-2D abaixo: 1 1 2 2 3 3 4 4 . . . . . . não é equivalente a uma P.A. pois, ao ser linearizada, resulta 1 1 2 2 3 3 4 4 . . . Quando uma PA-2D é equivalente a uma P.A., diremos que é redut'\ensuremath{\imath}vel, caso contrário, que é irredut'\ensuremath{\imath}vel. 260
\prov{relações: parte\_de matematica}
\prov{proveniência: gentil/Novas\_Sequencias\_Aritmeticas\_e\_Geometric.pdf \textperiodcentered{} fato \textperiodcentered{} confiança 1.0 \textperiodcentered{} domínio matematica \textperiodcentered{} história 3 versões no jornal}

%CRISTAL {"arestas":[{"alvo":"matematica","confianca":0.95,"epistemico":"fato","origem":"paper","rel":"parte_de"}],"confianca":1.0,"contraexemplos":[],"descricao":"Chama-se progressão geométrica bidimensional ( PG-2D) uma se- quência dupla dada pela seguinte fórmula de recorrência:          a11 = a a1j = a1(j−1) × q1, j ≥2; aij = a(i−1)j × q2, i ≥2, j ≥1. Onde: a11 = a, q1 e q2 são constantes dadas (não nulas). Nota: Esta","epistemico":"fato","exemplos":[],"id":"gentil_definicao_19_chama-se_progressao_geometrica_bidimensional_pg-2d_uma_se","memoria":"semantica","meta":{"arquivo":"gentil/Novas_Sequencias_Aritmeticas_e_Geometric.pdf","ciencia":"teoria_dos_numeros","dominio":"matematica","nivel":4,"numero":"19","tipo_no":"paper","tipo_teorema":"definição"},"origem":"gentil/Novas_Sequencias_Aritmeticas_e_Geometric.pdf","palavras_chave":[],"sinonimos":[],"tipo":"conceito","titulo":"gentil Definição 19 — Chama-se progressão geométrica bidimensional ( PG-2D) uma se"}
\section{gentil Definição 19 — Chama-se progressão geométrica bidimensional ( PG-2D) uma se}
Chama-se progressão geométrica bidimensional ( PG-2D) uma se- quência dupla dada pela seguinte fórmula de recorrência: \{         a11 = a a1j = a1(j\ensuremath{-}1) $\times$ q1, j \ensuremath{\ge}2; aij = a(i\ensuremath{-}1)j $\times$ q2, i \ensuremath{\ge}2, j \ensuremath{\ge}1. Onde: a11 = a, q1 e q2 são constantes dadas (não nulas). Nota: Esta
\prov{relações: parte\_de matematica}
\prov{proveniência: gentil/Novas\_Sequencias\_Aritmeticas\_e\_Geometric.pdf \textperiodcentered{} fato \textperiodcentered{} confiança 1.0 \textperiodcentered{} domínio matematica \textperiodcentered{} história 3 versões no jornal}

%CRISTAL {"arestas":[{"alvo":"matematica","confianca":0.95,"epistemico":"fato","origem":"paper","rel":"parte_de"}],"confianca":1.0,"contraexemplos":[],"descricao":"Chama-se progressão aritmética de ordem m ( P A.m ) uma sequência dada pela seguinte fórmula de recorrência:          an0 = r, r ̸= 0, n ≥1; a1j = aj, j = 1, 2, . . . , m; anm = a(n−1)m + a(n−1)(m−1), m ≥1, n ≥2. vemos que são duais, o que implica que a toda propriedade de uma P.A.m corres- ponde uma propriedade dual para uma P.G.m . 130\nExemplo: Vejamos a ideia que está por trás desta","epistemico":"fato","exemplos":[],"id":"gentil_definicao_1_chama-se_progressao_aritmetica_de_ordem_m_p_am_uma_sequ","memoria":"semantica","meta":{"arquivo":"gentil/Novas_Sequencias_Aritmeticas_e_Geometric.pdf","ciencia":"teoria_dos_numeros","dominio":"matematica","nivel":4,"numero":"1","tipo_no":"paper","tipo_teorema":"definição"},"origem":"gentil/Novas_Sequencias_Aritmeticas_e_Geometric.pdf","palavras_chave":[],"sinonimos":[],"tipo":"conceito","titulo":"gentil Definição 1 — Chama-se progressão aritmética de ordem m ( P A.m ) uma sequ"}
\section{gentil Definição 1 — Chama-se progressão aritmética de ordem m ( P A.m ) uma sequ}
Chama-se progressão aritmética de ordem m ( P A.m ) uma sequência dada pela seguinte fórmula de recorrência: \{         an0 = r, r \ensuremath{}= 0, n \ensuremath{\ge}1; a1j = aj, j = 1, 2, . . . , m; anm = a(n\ensuremath{-}1)m + a(n\ensuremath{-}1)(m\ensuremath{-}1), m \ensuremath{\ge}1, n \ensuremath{\ge}2. vemos que são duais, o que implica que a toda propriedade de uma P.A.m corres- ponde uma propriedade dual para uma P.G.m . 130
Exemplo: Vejamos a ideia que está por trás desta
\prov{relações: parte\_de matematica}
\prov{proveniência: gentil/Novas\_Sequencias\_Aritmeticas\_e\_Geometric.pdf \textperiodcentered{} fato \textperiodcentered{} confiança 1.0 \textperiodcentered{} domínio matematica \textperiodcentered{} história 6 versões no jornal}

%CRISTAL {"arestas":[{"alvo":"matematica","confianca":0.95,"epistemico":"fato","origem":"paper","rel":"parte_de"}],"confianca":1.0,"contraexemplos":[],"descricao":"Chama-se progressão aritmética tridimensional (PA-3D) uma sequên- cia tripla dada pela seguinte fórmula de recorrência:                a111 = a a1j1 = a1(j−1)1 + r1, j ≥2; aij1 = a(i−1)j1 + r2, i ≥2, j ≥1. aijk = aij(k−1) + r3, i ≥1, j ≥1, k ≥2. Onde: a111 = a, r1, r2 e r3 são constantes dadas. a111 a121 a131 a211 a221 a231 a311 a321 a331 a112 a122 a132 a212 a222 a232 a312 a322 a332 a113 a123 a133 a213 a223 a233 a313 a323 a333 y z x Observe como se dá a dinâmica desta construção,","epistemico":"fato","exemplos":[],"id":"gentil_definicao_20_chama-se_progressao_aritmetica_tridimensional_pa-3d_uma_se","memoria":"semantica","meta":{"arquivo":"gentil/Novas_Sequencias_Aritmeticas_e_Geometric.pdf","ciencia":"teoria_dos_numeros","dominio":"matematica","nivel":4,"numero":"20","tipo_no":"paper","tipo_teorema":"definição"},"origem":"gentil/Novas_Sequencias_Aritmeticas_e_Geometric.pdf","palavras_chave":[],"sinonimos":[],"tipo":"conceito","titulo":"gentil Definição 20 — Chama-se progressão aritmética tridimensional (PA-3D) uma se"}
\section{gentil Definição 20 — Chama-se progressão aritmética tridimensional (PA-3D) uma se}
Chama-se progressão aritmética tridimensional (PA-3D) uma sequên- cia tripla dada pela seguinte fórmula de recorrência: \{               a111 = a a1j1 = a1(j\ensuremath{-}1)1 + r1, j \ensuremath{\ge}2; aij1 = a(i\ensuremath{-}1)j1 + r2, i \ensuremath{\ge}2, j \ensuremath{\ge}1. aijk = aij(k\ensuremath{-}1) + r3, i \ensuremath{\ge}1, j \ensuremath{\ge}1, k \ensuremath{\ge}2. Onde: a111 = a, r1, r2 e r3 são constantes dadas. a111 a121 a131 a211 a221 a231 a311 a321 a331 a112 a122 a132 a212 a222 a232 a312 a322 a332 a113 a123 a133 a213 a223 a233 a313 a323 a333 y z x Observe como se dá a dinâmica desta construção,
\prov{relações: parte\_de matematica}
\prov{proveniência: gentil/Novas\_Sequencias\_Aritmeticas\_e\_Geometric.pdf \textperiodcentered{} fato \textperiodcentered{} confiança 1.0 \textperiodcentered{} domínio matematica \textperiodcentered{} história 3 versões no jornal}

%CRISTAL {"arestas":[{"alvo":"matematica","confianca":0.95,"epistemico":"fato","origem":"paper","rel":"parte_de"}],"confianca":1.0,"contraexemplos":[],"descricao":"Dizemos que uma sequência simples é equivalente a uma PA-3D quando existe uma PA-3D (com N colunas e M linhas) que ao ser linearizada coincide com a sequência simples. Exemplo: A sequência dos números ´ımpares positivos, sem os múltiplos de 5, isto é, 1 3 7 9 11 13 17 19 21 23 27 29 . . . é equivalente a uma PA-3D, 1 3 11 13 21 23 31 33 7 9 17 19 27 29 37 39 r1 = 2 r3 = 10 r2 = 6 1 M = 2, N = 2 372","epistemico":"fato","exemplos":[],"id":"gentil_definicao_21_dizemos_que_uma_sequencia_simples_e_equivalente_a_uma_pa-3d","memoria":"semantica","meta":{"arquivo":"gentil/Novas_Sequencias_Aritmeticas_e_Geometric.pdf","ciencia":"teoria_dos_numeros","dominio":"matematica","nivel":4,"numero":"21","tipo_no":"paper","tipo_teorema":"definição"},"origem":"gentil/Novas_Sequencias_Aritmeticas_e_Geometric.pdf","palavras_chave":[],"sinonimos":[],"tipo":"conceito","titulo":"gentil Definição 21 — Dizemos que uma sequência simples é equivalente a uma PA-3D"}
\section{gentil Definição 21 — Dizemos que uma sequência simples é equivalente a uma PA-3D}
Dizemos que uma sequência simples é equivalente a uma PA-3D quando existe uma PA-3D (com N colunas e M linhas) que ao ser linearizada coincide com a sequência simples. Exemplo: A sequência dos números '\ensuremath{\imath}mpares positivos, sem os múltiplos de 5, isto é, 1 3 7 9 11 13 17 19 21 23 27 29 . . . é equivalente a uma PA-3D, 1 3 11 13 21 23 31 33 7 9 17 19 27 29 37 39 r1 = 2 r3 = 10 r2 = 6 1 M = 2, N = 2 372
\prov{relações: parte\_de matematica}
\prov{proveniência: gentil/Novas\_Sequencias\_Aritmeticas\_e\_Geometric.pdf \textperiodcentered{} fato \textperiodcentered{} confiança 1.0 \textperiodcentered{} domínio matematica \textperiodcentered{} história 3 versões no jornal}

%CRISTAL {"arestas":[{"alvo":"matematica","confianca":0.95,"epistemico":"fato","origem":"paper","rel":"parte_de"}],"confianca":1.0,"contraexemplos":[],"descricao":"Um quadrado numérico é mágico se possui n2 números inteiros positivos −diferentes entre si −tais que a soma dos n números que figuram nas linhas, colunas e diagonais é constante. Em particular, chamaremos de quadrado mágico natural ao quadrado mágico formado pelos 1, 2, . . . , n2 primeiros números naturais. Exemplo de quadrado mágico (natural): 1 9 2 3 5 7 8 1 6 Observe que a soma dos números em cada uma das linhas, em cada uma das colunas e em cada uma das diagonais é sempre 15. Este valor é c","epistemico":"fato","exemplos":[],"id":"gentil_definicao_22_um_quadrado_numerico_e_magico_se_possui_n2_numeros_inteiros","memoria":"semantica","meta":{"arquivo":"gentil/Novas_Sequencias_Aritmeticas_e_Geometric.pdf","ciencia":"teoria_dos_numeros","dominio":"matematica","nivel":4,"numero":"22","tipo_no":"paper","tipo_teorema":"definição"},"origem":"gentil/Novas_Sequencias_Aritmeticas_e_Geometric.pdf","palavras_chave":[],"sinonimos":[],"tipo":"conceito","titulo":"gentil Definição 22 — Um quadrado numérico é mágico se possui n2 números inteiros"}
\section{gentil Definição 22 — Um quadrado numérico é mágico se possui n2 números inteiros}
Um quadrado numérico é mágico se possui n2 números inteiros positivos \ensuremath{-}diferentes entre si \ensuremath{-}tais que a soma dos n números que figuram nas linhas, colunas e diagonais é constante. Em particular, chamaremos de quadrado mágico natural ao quadrado mágico formado pelos 1, 2, . . . , n2 primeiros números naturais. Exemplo de quadrado mágico (natural): 1 9 2 3 5 7 8 1 6 Observe que a soma dos números em cada uma das linhas, em cada uma das colunas e em cada uma das diagonais é sempre 15. Este valor é c
\prov{relações: parte\_de matematica}
\prov{proveniência: gentil/Novas\_Sequencias\_Aritmeticas\_e\_Geometric.pdf \textperiodcentered{} fato \textperiodcentered{} confiança 1.0 \textperiodcentered{} domínio matematica \textperiodcentered{} história 3 versões no jornal}

%CRISTAL {"arestas":[{"alvo":"matematica","confianca":0.95,"epistemico":"fato","origem":"paper","rel":"parte_de"}],"confianca":1.0,"contraexemplos":[],"descricao":"p 82) para o cálculo da m-ésima Diferença é enfadonho. A fórmula seguinte foi obtida com o aux´ılio do Triângulo Aritmético de Pascal (TAP) e nos fornece a m-ésima Diferença sem recursividade. 84","epistemico":"fato","exemplos":[],"id":"gentil_definicao_3_p_82_para_o_calculo_da_m-esima_diferenca_e_enfadonho_a_for","memoria":"semantica","meta":{"arquivo":"gentil/Novas_Sequencias_Aritmeticas_e_Geometric.pdf","ciencia":"teoria_dos_numeros","dominio":"matematica","nivel":4,"numero":"3","tipo_no":"paper","tipo_teorema":"definição"},"origem":"gentil/Novas_Sequencias_Aritmeticas_e_Geometric.pdf","palavras_chave":[],"sinonimos":[],"tipo":"conceito","titulo":"gentil Definição 3 — p 82) para o cálculo da m-ésima Diferença é enfadonho. A fór"}
\section{gentil Definição 3 — p 82) para o cálculo da m-ésima Diferença é enfadonho. A fór}
p 82) para o cálculo da m-ésima Diferença é enfadonho. A fórmula seguinte foi obtida com o aux'\ensuremath{\imath}lio do Triângulo Aritmético de Pascal (TAP) e nos fornece a m-ésima Diferença sem recursividade. 84
\prov{relações: parte\_de matematica}
\prov{proveniência: gentil/Novas\_Sequencias\_Aritmeticas\_e\_Geometric.pdf \textperiodcentered{} fato \textperiodcentered{} confiança 1.0 \textperiodcentered{} domínio matematica \textperiodcentered{} história 3 versões no jornal}

%CRISTAL {"arestas":[{"alvo":"gentil_novas_sequencias_aritmeticas_geometricas","confianca":1.0,"epistemico":"fato","origem":"gentil/Novas_Sequencias_Aritmeticas_e_Geometric.pdf","rel":"define"}],"confianca":1.0,"contraexemplos":[],"descricao":"p 82), temos: ∆0 ( a f + b g )(n) = ( a f + b g )(n) = a f(n) + b g(n) = a ∆0 f(n) + b ∆0 g(n) Suponhamos a validade da fórmula para m = p: (H.I.) ∆p ( a f + b g )(n) = a ∆p f(n) + b ∆p g(n) e provemos que vale para m = p + 1. Então, pela","epistemico":"fato","exemplos":[],"id":"gentil_definicao_3_p_82_temos_0_a_f_b_g_n_a_f_b_g_n_a_fn","memoria":"semantica","meta":{"arquivo":"gentil/Novas_Sequencias_Aritmeticas_e_Geometric.pdf","ciencia":"teoria_dos_numeros","dominio":"matematica","nivel":4,"numero":"3","tipo_no":"paper","tipo_teorema":"definição"},"origem":"gentil/Novas_Sequencias_Aritmeticas_e_Geometric.pdf","palavras_chave":[],"sinonimos":[],"tipo":"conceito","titulo":"gentil Definição 3 — p 82), temos: ∆0 ( a f + b g )(n) = ( a f + b g )(n) = a f(n"}
\section{gentil Definição 3 — p 82), temos: \ensuremath{\Delta}0 ( a f + b g )(n) = ( a f + b g )(n) = a f(n}
p 82), temos: \ensuremath{\Delta}0 ( a f + b g )(n) = ( a f + b g )(n) = a f(n) + b g(n) = a \ensuremath{\Delta}0 f(n) + b \ensuremath{\Delta}0 g(n) Suponhamos a validade da fórmula para m = p: (H.I.) \ensuremath{\Delta}p ( a f + b g )(n) = a \ensuremath{\Delta}p f(n) + b \ensuremath{\Delta}p g(n) e provemos que vale para m = p + 1. Então, pela
\prov{relações: define gentil\_novas\_sequencias\_aritmeticas\_geometricas}
\prov{proveniência: gentil/Novas\_Sequencias\_Aritmeticas\_e\_Geometric.pdf \textperiodcentered{} fato \textperiodcentered{} confiança 1.0 \textperiodcentered{} domínio matematica \textperiodcentered{} história 4 versões no jornal}

%CRISTAL {"arestas":[{"alvo":"gentil_novas_sequencias_aritmeticas_geometricas","confianca":1.0,"epistemico":"fato","origem":"gentil/Novas_Sequencias_Aritmeticas_e_Geometric.pdf","rel":"define"},{"alvo":"gentil_capitulo_1_sequencias_aritmeticas_de_ordem_m","confianca":1.0,"epistemico":"fato","origem":"gentil/Novas_Sequencias_Aritmeticas_e_Geometric.pdf","rel":"define"}],"confianca":1.0,"contraexemplos":[],"descricao":"p 94), temos: ∇0 ( a f + b g )(n) = ( a f + b g )(n) = a f(n) + b g(n) = a ∇0 f(n) + b ∇0 g(n) Suponhamos a validade da fórmula para m = p: (H.I.) ∇p ( a f + b g )(n) = a ∇p f(n) + b ∇p g(n) e provemos que vale para m = p + 1: (T.I.) ∇p+1 ( a f + b g )(n) = a ∇p+1 f(n) + b ∇p+1 g(n) Então, ∆p+1 ( a f + b g )(n) = n−1 X i = 1 ∇p ( a f + b g )(i) + c = n−1 X i = 1 [ a ∇p f(i) + b ∇p g(i) ] + c = a n−1 X i = 1 ∇p f(i) + b n−1 X i = 1 ∇p g(i) + a · c1 + b · c2 = a \u0010 n−1 X i = 1 ∇p f(i) + c1 \u0011 + b \u0010","epistemico":"fato","exemplos":[],"id":"gentil_definicao_4_p_94_temos_0_a_f_b_g_n_a_f_b_g_n_a_fn","memoria":"semantica","meta":{"arquivo":"gentil/Novas_Sequencias_Aritmeticas_e_Geometric.pdf","ciencia":"teoria_dos_numeros","dominio":"matematica","nivel":4,"numero":"4","tipo_no":"paper","tipo_teorema":"definição"},"origem":"gentil/Novas_Sequencias_Aritmeticas_e_Geometric.pdf","palavras_chave":[],"sinonimos":[],"tipo":"conceito","titulo":"gentil Definição 4 — p 94), temos: ∇0 ( a f + b g )(n) = ( a f + b g )(n) = a f(n"}
\section{gentil Definição 4 — p 94), temos: \ensuremath{\nabla}0 ( a f + b g )(n) = ( a f + b g )(n) = a f(n}
p 94), temos: \ensuremath{\nabla}0 ( a f + b g )(n) = ( a f + b g )(n) = a f(n) + b g(n) = a \ensuremath{\nabla}0 f(n) + b \ensuremath{\nabla}0 g(n) Suponhamos a validade da fórmula para m = p: (H.I.) \ensuremath{\nabla}p ( a f + b g )(n) = a \ensuremath{\nabla}p f(n) + b \ensuremath{\nabla}p g(n) e provemos que vale para m = p + 1: (T.I.) \ensuremath{\nabla}p+1 ( a f + b g )(n) = a \ensuremath{\nabla}p+1 f(n) + b \ensuremath{\nabla}p+1 g(n) Então, \ensuremath{\Delta}p+1 ( a f + b g )(n) = n\ensuremath{-}1 X i = 1 \ensuremath{\nabla}p ( a f + b g )(i) + c = n\ensuremath{-}1 X i = 1 [ a \ensuremath{\nabla}p f(i) + b \ensuremath{\nabla}p g(i) ] + c = a n\ensuremath{-}1 X i = 1 \ensuremath{\nabla}p f(i) + b n\ensuremath{-}1 X i = 1 \ensuremath{\nabla}p g(i) + a \textperiodcentered  c1 + b \textperiodcentered  c2 = a  n\ensuremath{-}1 X i = 1 \ensuremath{\nabla}p f(i) + c1  + b 
\prov{relações: define gentil\_novas\_sequencias\_aritmeticas\_geometricas; define gentil\_capitulo\_1\_sequencias\_aritmeticas\_de\_ordem\_m}
\prov{proveniência: gentil/Novas\_Sequencias\_Aritmeticas\_e\_Geometric.pdf \textperiodcentered{} fato \textperiodcentered{} confiança 1.0 \textperiodcentered{} domínio matematica \textperiodcentered{} história 4 versões no jornal}

%CRISTAL {"arestas":[{"alvo":"matematica","confianca":0.95,"epistemico":"fato","origem":"paper","rel":"parte_de"}],"confianca":1.0,"contraexemplos":[],"descricao":"Chama-se progressão geométrica de ordem m ( P G.m ) uma sequência dada pela seguinte fórmula de recorrência:          an0 = q, q ̸= 0, 1; n ≥1; a1j = aj, aj ̸= 0; j = 1, 2, . . . , m; anm = a(n−1)m × a(n−1)(m−1), m ≥1, n ≥2. Onde: (i) m ≥1 é um natural arbitrariamente fixado. (ii) q e aj são constantes dadas. Podemos chamar q de razão ou semente da P.G. de ordem m. Por","epistemico":"fato","exemplos":[],"id":"gentil_definicao_8_chama-se_progressao_geometrica_de_ordem_m_p_gm_uma_sequ","memoria":"semantica","meta":{"arquivo":"gentil/Novas_Sequencias_Aritmeticas_e_Geometric.pdf","ciencia":"teoria_dos_numeros","dominio":"matematica","nivel":4,"numero":"8","tipo_no":"paper","tipo_teorema":"definição"},"origem":"gentil/Novas_Sequencias_Aritmeticas_e_Geometric.pdf","palavras_chave":[],"sinonimos":[],"tipo":"conceito","titulo":"gentil Definição 8 — Chama-se progressão geométrica de ordem m ( P G.m ) uma sequ"}
\section{gentil Definição 8 — Chama-se progressão geométrica de ordem m ( P G.m ) uma sequ}
Chama-se progressão geométrica de ordem m ( P G.m ) uma sequência dada pela seguinte fórmula de recorrência: \{         an0 = q, q \ensuremath{}= 0, 1; n \ensuremath{\ge}1; a1j = aj, aj \ensuremath{}= 0; j = 1, 2, . . . , m; anm = a(n\ensuremath{-}1)m $\times$ a(n\ensuremath{-}1)(m\ensuremath{-}1), m \ensuremath{\ge}1, n \ensuremath{\ge}2. Onde: (i) m \ensuremath{\ge}1 é um natural arbitrariamente fixado. (ii) q e aj são constantes dadas. Podemos chamar q de razão ou semente da P.G. de ordem m. Por
\prov{relações: parte\_de matematica}
\prov{proveniência: gentil/Novas\_Sequencias\_Aritmeticas\_e\_Geometric.pdf \textperiodcentered{} fato \textperiodcentered{} confiança 1.0 \textperiodcentered{} domínio matematica \textperiodcentered{} história 3 versões no jornal}

%CRISTAL {"arestas":[{"alvo":"matematica","confianca":0.95,"epistemico":"fato","origem":"paper","rel":"parte_de"}],"confianca":1.0,"contraexemplos":[],"descricao":"p 140) para o cálculo do m-ésimo Quociente é um tanto quanto enfadonho. A fórmula seguinte nos fornece o m-ésimo Quociente sem recursividade.","epistemico":"fato","exemplos":[],"id":"gentil_definicao_9_p_140_para_o_calculo_do_m-esimo_quociente_e_um_tanto_quanto","memoria":"semantica","meta":{"arquivo":"gentil/Novas_Sequencias_Aritmeticas_e_Geometric.pdf","ciencia":"teoria_dos_numeros","dominio":"matematica","nivel":4,"numero":"9","tipo_no":"paper","tipo_teorema":"definição"},"origem":"gentil/Novas_Sequencias_Aritmeticas_e_Geometric.pdf","palavras_chave":[],"sinonimos":[],"tipo":"conceito","titulo":"gentil Definição 9 — p 140) para o cálculo do m-ésimo Quociente é um tanto quanto"}
\section{gentil Definição 9 — p 140) para o cálculo do m-ésimo Quociente é um tanto quanto}
p 140) para o cálculo do m-ésimo Quociente é um tanto quanto enfadonho. A fórmula seguinte nos fornece o m-ésimo Quociente sem recursividade.
\prov{relações: parte\_de matematica}
\prov{proveniência: gentil/Novas\_Sequencias\_Aritmeticas\_e\_Geometric.pdf \textperiodcentered{} fato \textperiodcentered{} confiança 1.0 \textperiodcentered{} domínio matematica \textperiodcentered{} história 3 versões no jornal}

%CRISTAL {"arestas":[{"alvo":"gentil_fundamentos_capitulo_2_relac_oes_bin_arias","confianca":1.0,"epistemico":"fato","origem":"gentil/Gentil_Lopes_FUNDAMENTOS_DOS_NUMEROS_pdf.pdf","rel":"parte_de"},{"alvo":"resultado_anterior_identidade_nao_cobertura","confianca":0.5,"epistemico":"fato","origem":"wiki","rel":"relaciona"},{"alvo":"word","confianca":0.5,"epistemico":"fato","origem":"wiki","rel":"relaciona"},{"alvo":"syscall","confianca":0.5,"epistemico":"fato","origem":"wiki","rel":"relaciona"},{"alvo":"virtualizacao","confianca":0.5,"epistemico":"fato","origem":"wiki","rel":"relaciona"}],"confianca":1.0,"contraexemplos":[],"descricao":"Secção 2.2 — p.59.","epistemico":"fato","exemplos":[],"id":"gentil_fundamentos_22_relacoes_especiais","memoria":"semantica","meta":{"arquivo":"gentil/Lopes_Lopes_FUNDAMENTOS_DOS_NUMEROS_pdf.pdf","autor":"Gentil Lopes da Silva","ciencia":"teoria_dos_numeros","dominio":"matematica","nivel":4,"pagina":59,"secao":"2.2","tipo_no":"paper","titulo_legivel":"Fundamentos dos Números"},"origem":"gentil/Gentil_Lopes_FUNDAMENTOS_DOS_NUMEROS_pdf.pdf","palavras_chave":[],"sinonimos":[],"tipo":"conceito","titulo":"gentil fundamentos 2.2 — Relações Especiais"}
\section{gentil fundamentos 2.2 — Relações Especiais}
Secção 2.2 — p.59.
\prov{relações: parte\_de gentil\_fundamentos\_capitulo\_2\_relac\_oes\_bin\_arias; relaciona resultado\_anterior\_identidade\_nao\_cobertura; relaciona word; relaciona syscall; relaciona virtualizacao}
\prov{proveniência: gentil/Gentil\_Lopes\_FUNDAMENTOS\_DOS\_NUMEROS\_pdf.pdf \textperiodcentered{} fato \textperiodcentered{} confiança 1.0 \textperiodcentered{} domínio matematica \textperiodcentered{} história 9 versões no jornal}

%CRISTAL {"arestas":[{"alvo":"gentil_fundamentos_dos_numeros","confianca":1.0,"epistemico":"fato","origem":"gentil/Gentil_Lopes_FUNDAMENTOS_DOS_NUMEROS_pdf.pdf","rel":"parte_de"},{"alvo":"reais","confianca":0.9,"epistemico":"fato","origem":"gentil/Gentil_Lopes_FUNDAMENTOS_DOS_NUMEROS_pdf.pdf","rel":"define"},{"alvo":"word","confianca":0.5,"epistemico":"fato","origem":"wiki","rel":"relaciona"}],"confianca":1.0,"contraexemplos":[],"descricao":"Capítulo 10 — p.383. Fundamentos dos Números.","epistemico":"fato","exemplos":[],"id":"gentil_fundamentos_capitulo_10_n_umeros_reais_por_cantor","memoria":"semantica","meta":{"arquivo":"gentil/Lopes_Lopes_FUNDAMENTOS_DOS_NUMEROS_pdf.pdf","autor":"Gentil Lopes da Silva","capitulo":"10","ciencia":"teoria_dos_numeros","dominio":"matematica","nivel":4,"pagina":383,"tipo_no":"paper","titulo_legivel":"Fundamentos dos Números"},"origem":"gentil/Gentil_Lopes_FUNDAMENTOS_DOS_NUMEROS_pdf.pdf","palavras_chave":[],"sinonimos":["capítulo 10"],"tipo":"conceito","titulo":"gentil fundamentos capítulo 10 — N ´UMEROS REAIS POR CANTOR"}
\section{gentil fundamentos capítulo 10 — N 'UMEROS REAIS POR CANTOR}
Capítulo 10 — p.383. Fundamentos dos Números.
\prov{sinónimos: capítulo 10}
\prov{relações: parte\_de gentil\_fundamentos\_dos\_numeros; define reais; relaciona word}
\prov{proveniência: gentil/Gentil\_Lopes\_FUNDAMENTOS\_DOS\_NUMEROS\_pdf.pdf \textperiodcentered{} fato \textperiodcentered{} confiança 1.0 \textperiodcentered{} domínio matematica \textperiodcentered{} história 6 versões no jornal}

%CRISTAL {"arestas":[{"alvo":"gentil_fundamentos_dos_numeros","confianca":1.0,"epistemico":"fato","origem":"gentil/Gentil_Lopes_FUNDAMENTOS_DOS_NUMEROS_pdf.pdf","rel":"parte_de"},{"alvo":"inteiros","confianca":0.9,"epistemico":"fato","origem":"gentil/Gentil_Lopes_FUNDAMENTOS_DOS_NUMEROS_pdf.pdf","rel":"define"},{"alvo":"racionais","confianca":0.9,"epistemico":"fato","origem":"gentil/Gentil_Lopes_FUNDAMENTOS_DOS_NUMEROS_pdf.pdf","rel":"define"},{"alvo":"reais","confianca":0.9,"epistemico":"fato","origem":"gentil/Gentil_Lopes_FUNDAMENTOS_DOS_NUMEROS_pdf.pdf","rel":"define"}],"confianca":1.0,"contraexemplos":[],"descricao":"Capítulo 11 — p.453. Fundamentos dos Números.","epistemico":"fato","exemplos":[],"id":"gentil_fundamentos_capitulo_11_n_umeros_reais_azuis_e_vermelhos","memoria":"semantica","meta":{"arquivo":"gentil/Lopes_Lopes_FUNDAMENTOS_DOS_NUMEROS_pdf.pdf","autor":"Gentil Lopes da Silva","capitulo":"11","ciencia":"teoria_dos_numeros","dominio":"matematica","nivel":4,"pagina":453,"tipo_no":"paper","titulo_legivel":"Fundamentos dos Números"},"origem":"gentil/Gentil_Lopes_FUNDAMENTOS_DOS_NUMEROS_pdf.pdf","palavras_chave":[],"sinonimos":["capítulo 11"],"tipo":"conceito","titulo":"gentil fundamentos capítulo 11 — N ´UMEROS REAIS AZUIS E VERMELHOS"}
\section{gentil fundamentos capítulo 11 — N 'UMEROS REAIS AZUIS E VERMELHOS}
Capítulo 11 — p.453. Fundamentos dos Números.
\prov{sinónimos: capítulo 11}
\prov{relações: parte\_de gentil\_fundamentos\_dos\_numeros; define inteiros; define racionais; define reais}
\prov{proveniência: gentil/Gentil\_Lopes\_FUNDAMENTOS\_DOS\_NUMEROS\_pdf.pdf \textperiodcentered{} fato \textperiodcentered{} confiança 1.0 \textperiodcentered{} domínio matematica \textperiodcentered{} história 6 versões no jornal}

%CRISTAL {"arestas":[{"alvo":"gentil_fundamentos_dos_numeros","confianca":1.0,"epistemico":"fato","origem":"gentil/Gentil_Lopes_FUNDAMENTOS_DOS_NUMEROS_pdf.pdf","rel":"parte_de"},{"alvo":"complexos","confianca":0.9,"epistemico":"fato","origem":"gentil/Gentil_Lopes_FUNDAMENTOS_DOS_NUMEROS_pdf.pdf","rel":"define"},{"alvo":"word","confianca":0.5,"epistemico":"fato","origem":"wiki","rel":"relaciona"}],"confianca":1.0,"contraexemplos":[],"descricao":"Capítulo 12 — p.455. Fundamentos dos Números.","epistemico":"fato","exemplos":[],"id":"gentil_fundamentos_capitulo_12_n_umeros_complexos","memoria":"semantica","meta":{"arquivo":"gentil/Lopes_Lopes_FUNDAMENTOS_DOS_NUMEROS_pdf.pdf","autor":"Gentil Lopes da Silva","capitulo":"12","ciencia":"teoria_dos_numeros","dominio":"matematica","nivel":4,"pagina":455,"tipo_no":"paper","titulo_legivel":"Fundamentos dos Números"},"origem":"gentil/Gentil_Lopes_FUNDAMENTOS_DOS_NUMEROS_pdf.pdf","palavras_chave":[],"sinonimos":["capítulo 12"],"tipo":"conceito","titulo":"gentil fundamentos capítulo 12 — N ´UMEROS COMPLEXOS"}
\section{gentil fundamentos capítulo 12 — N 'UMEROS COMPLEXOS}
Capítulo 12 — p.455. Fundamentos dos Números.
\prov{sinónimos: capítulo 12}
\prov{relações: parte\_de gentil\_fundamentos\_dos\_numeros; define complexos; relaciona word}
\prov{proveniência: gentil/Gentil\_Lopes\_FUNDAMENTOS\_DOS\_NUMEROS\_pdf.pdf \textperiodcentered{} fato \textperiodcentered{} confiança 1.0 \textperiodcentered{} domínio matematica \textperiodcentered{} história 6 versões no jornal}

%CRISTAL {"arestas":[{"alvo":"gentil_fundamentos_dos_numeros","confianca":1.0,"epistemico":"fato","origem":"gentil/Gentil_Lopes_FUNDAMENTOS_DOS_NUMEROS_pdf.pdf","rel":"parte_de"},{"alvo":"complexos","confianca":0.9,"epistemico":"fato","origem":"gentil/Gentil_Lopes_FUNDAMENTOS_DOS_NUMEROS_pdf.pdf","rel":"generaliza"},{"alvo":"word","confianca":0.5,"epistemico":"fato","origem":"wiki","rel":"relaciona"},{"alvo":"vontade_de_poder","confianca":0.5,"epistemico":"fato","origem":"wiki","rel":"relaciona"},{"alvo":"virtualizacao","confianca":0.5,"epistemico":"fato","origem":"wiki","rel":"relaciona"},{"alvo":"while","confianca":0.5,"epistemico":"fato","origem":"wiki","rel":"relaciona"}],"confianca":1.0,"contraexemplos":[],"descricao":"Capítulo 13 — p.471. Fundamentos dos Números.","epistemico":"fato","exemplos":[],"id":"gentil_fundamentos_capitulo_13_n_umeros_hipercomplexos","memoria":"semantica","meta":{"arquivo":"gentil/Lopes_Lopes_FUNDAMENTOS_DOS_NUMEROS_pdf.pdf","autor":"Gentil Lopes da Silva","capitulo":"13","ciencia":"teoria_dos_numeros","dominio":"matematica","nivel":4,"pagina":471,"tipo_no":"paper","titulo_legivel":"Fundamentos dos Números"},"origem":"gentil/Gentil_Lopes_FUNDAMENTOS_DOS_NUMEROS_pdf.pdf","palavras_chave":[],"sinonimos":["capítulo 13"],"tipo":"conceito","titulo":"gentil fundamentos capítulo 13 — N ´UMEROS HIPERCOMPLEXOS"}
\section{gentil fundamentos capítulo 13 — N 'UMEROS HIPERCOMPLEXOS}
Capítulo 13 — p.471. Fundamentos dos Números.
\prov{sinónimos: capítulo 13}
\prov{relações: parte\_de gentil\_fundamentos\_dos\_numeros; generaliza complexos; relaciona word; relaciona vontade\_de\_poder; relaciona virtualizacao; relaciona while}
\prov{proveniência: gentil/Gentil\_Lopes\_FUNDAMENTOS\_DOS\_NUMEROS\_pdf.pdf \textperiodcentered{} fato \textperiodcentered{} confiança 1.0 \textperiodcentered{} domínio matematica \textperiodcentered{} história 10 versões no jornal}

%CRISTAL {"arestas":[{"alvo":"gentil_fundamentos_dos_numeros","confianca":1.0,"epistemico":"fato","origem":"gentil/Gentil_Lopes_FUNDAMENTOS_DOS_NUMEROS_pdf.pdf","rel":"parte_de"},{"alvo":"logica","confianca":0.9,"epistemico":"fato","origem":"gentil/Gentil_Lopes_FUNDAMENTOS_DOS_NUMEROS_pdf.pdf","rel":"referencia"},{"alvo":"word","confianca":0.5,"epistemico":"fato","origem":"wiki","rel":"relaciona"},{"alvo":"virtualizacao","confianca":0.5,"epistemico":"fato","origem":"wiki","rel":"relaciona"},{"alvo":"while","confianca":0.5,"epistemico":"fato","origem":"wiki","rel":"relaciona"},{"alvo":"vontade_de_poder","confianca":0.5,"epistemico":"fato","origem":"wiki","rel":"relaciona"}],"confianca":1.0,"contraexemplos":[],"descricao":"Capítulo 14 — p.487. Fundamentos dos Números.","epistemico":"fato","exemplos":[],"id":"gentil_fundamentos_capitulo_14_consultas","memoria":"semantica","meta":{"arquivo":"gentil/Lopes_Lopes_FUNDAMENTOS_DOS_NUMEROS_pdf.pdf","autor":"Gentil Lopes da Silva","capitulo":"14","ciencia":"teoria_dos_numeros","dominio":"matematica","nivel":4,"pagina":487,"tipo_no":"paper","titulo_legivel":"Fundamentos dos Números"},"origem":"gentil/Gentil_Lopes_FUNDAMENTOS_DOS_NUMEROS_pdf.pdf","palavras_chave":[],"sinonimos":["capítulo 14"],"tipo":"conceito","titulo":"gentil fundamentos capítulo 14 — CONSULTAS"}
\section{gentil fundamentos capítulo 14 — CONSULTAS}
Capítulo 14 — p.487. Fundamentos dos Números.
\prov{sinónimos: capítulo 14}
\prov{relações: parte\_de gentil\_fundamentos\_dos\_numeros; referencia logica; relaciona word; relaciona virtualizacao; relaciona while; relaciona vontade\_de\_poder}
\prov{proveniência: gentil/Gentil\_Lopes\_FUNDAMENTOS\_DOS\_NUMEROS\_pdf.pdf \textperiodcentered{} fato \textperiodcentered{} confiança 1.0 \textperiodcentered{} domínio matematica \textperiodcentered{} história 8 versões no jornal}

%CRISTAL {"arestas":[{"alvo":"gentil_fundamentos_dos_numeros","confianca":1.0,"epistemico":"fato","origem":"gentil/Gentil_Lopes_FUNDAMENTOS_DOS_NUMEROS_pdf.pdf","rel":"parte_de"},{"alvo":"virtualizacao","confianca":0.5,"epistemico":"fato","origem":"wiki","rel":"relaciona"},{"alvo":"word","confianca":0.5,"epistemico":"fato","origem":"wiki","rel":"relaciona"}],"confianca":1.0,"contraexemplos":[],"descricao":"Capítulo 1 — p.13. Fundamentos dos Números.","epistemico":"fato","exemplos":[],"id":"gentil_fundamentos_capitulo_1_o_que_e_um_n_umero","memoria":"semantica","meta":{"arquivo":"gentil/Lopes_Lopes_FUNDAMENTOS_DOS_NUMEROS_pdf.pdf","autor":"Gentil Lopes da Silva","capitulo":"1","ciencia":"teoria_dos_numeros","dominio":"matematica","nivel":4,"pagina":13,"tipo_no":"paper","titulo_legivel":"Fundamentos dos Números"},"origem":"gentil/Gentil_Lopes_FUNDAMENTOS_DOS_NUMEROS_pdf.pdf","palavras_chave":[],"sinonimos":["capítulo 1"],"tipo":"conceito","titulo":"gentil fundamentos capítulo 1 — O QUE ´E UM N ´UMERO"}
\section{gentil fundamentos capítulo 1 — O QUE 'E UM N 'UMERO}
Capítulo 1 — p.13. Fundamentos dos Números.
\prov{sinónimos: capítulo 1}
\prov{relações: parte\_de gentil\_fundamentos\_dos\_numeros; relaciona virtualizacao; relaciona word}
\prov{proveniência: gentil/Gentil\_Lopes\_FUNDAMENTOS\_DOS\_NUMEROS\_pdf.pdf \textperiodcentered{} fato \textperiodcentered{} confiança 1.0 \textperiodcentered{} domínio matematica \textperiodcentered{} história 7 versões no jornal}

%CRISTAL {"arestas":[{"alvo":"gentil_fundamentos_dos_numeros","confianca":1.0,"epistemico":"fato","origem":"gentil/Gentil_Lopes_FUNDAMENTOS_DOS_NUMEROS_pdf.pdf","rel":"parte_de"},{"alvo":"reticulados","confianca":0.9,"epistemico":"fato","origem":"gentil/Gentil_Lopes_FUNDAMENTOS_DOS_NUMEROS_pdf.pdf","rel":"referencia"},{"alvo":"word","confianca":0.5,"epistemico":"fato","origem":"wiki","rel":"relaciona"}],"confianca":1.0,"contraexemplos":[],"descricao":"Capítulo 2 — p.49. Fundamentos dos Números.","epistemico":"fato","exemplos":[],"id":"gentil_fundamentos_capitulo_2_relac_oes_bin_arias","memoria":"semantica","meta":{"arquivo":"gentil/Lopes_Lopes_FUNDAMENTOS_DOS_NUMEROS_pdf.pdf","autor":"Gentil Lopes da Silva","capitulo":"2","ciencia":"teoria_dos_numeros","dominio":"matematica","nivel":4,"pagina":49,"tipo_no":"paper","titulo_legivel":"Fundamentos dos Números"},"origem":"gentil/Gentil_Lopes_FUNDAMENTOS_DOS_NUMEROS_pdf.pdf","palavras_chave":[],"sinonimos":["capítulo 2"],"tipo":"conceito","titulo":"gentil fundamentos capítulo 2 — RELAC¸ ˜OES BIN´ARIAS"}
\section{gentil fundamentos capítulo 2 — RELAC, \~{}OES BIN'ARIAS}
Capítulo 2 — p.49. Fundamentos dos Números.
\prov{sinónimos: capítulo 2}
\prov{relações: parte\_de gentil\_fundamentos\_dos\_numeros; referencia reticulados; relaciona word}
\prov{proveniência: gentil/Gentil\_Lopes\_FUNDAMENTOS\_DOS\_NUMEROS\_pdf.pdf \textperiodcentered{} fato \textperiodcentered{} confiança 1.0 \textperiodcentered{} domínio matematica \textperiodcentered{} história 8 versões no jornal}

%CRISTAL {"arestas":[{"alvo":"gentil_fundamentos_dos_numeros","confianca":1.0,"epistemico":"fato","origem":"gentil/Gentil_Lopes_FUNDAMENTOS_DOS_NUMEROS_pdf.pdf","rel":"parte_de"},{"alvo":"inteiros","confianca":0.9,"epistemico":"fato","origem":"gentil/Gentil_Lopes_FUNDAMENTOS_DOS_NUMEROS_pdf.pdf","rel":"especializa"},{"alvo":"word","confianca":0.5,"epistemico":"fato","origem":"wiki","rel":"relaciona"}],"confianca":1.0,"contraexemplos":[],"descricao":"Capítulo 3 — p.81. Fundamentos dos Números.","epistemico":"fato","exemplos":[],"id":"gentil_fundamentos_capitulo_3_n_umeros_naturais","memoria":"semantica","meta":{"arquivo":"gentil/Lopes_Lopes_FUNDAMENTOS_DOS_NUMEROS_pdf.pdf","autor":"Gentil Lopes da Silva","capitulo":"3","ciencia":"teoria_dos_numeros","dominio":"matematica","nivel":4,"pagina":81,"tipo_no":"paper","titulo_legivel":"Fundamentos dos Números"},"origem":"gentil/Gentil_Lopes_FUNDAMENTOS_DOS_NUMEROS_pdf.pdf","palavras_chave":[],"sinonimos":["capítulo 3"],"tipo":"conceito","titulo":"gentil fundamentos capítulo 3 — N ´UMEROS NATURAIS"}
\section{gentil fundamentos capítulo 3 — N 'UMEROS NATURAIS}
Capítulo 3 — p.81. Fundamentos dos Números.
\prov{sinónimos: capítulo 3}
\prov{relações: parte\_de gentil\_fundamentos\_dos\_numeros; especializa inteiros; relaciona word}
\prov{proveniência: gentil/Gentil\_Lopes\_FUNDAMENTOS\_DOS\_NUMEROS\_pdf.pdf \textperiodcentered{} fato \textperiodcentered{} confiança 1.0 \textperiodcentered{} domínio matematica \textperiodcentered{} história 8 versões no jornal}

%CRISTAL {"arestas":[{"alvo":"gentil_fundamentos_dos_numeros","confianca":1.0,"epistemico":"fato","origem":"gentil/Gentil_Lopes_FUNDAMENTOS_DOS_NUMEROS_pdf.pdf","rel":"parte_de"},{"alvo":"inteiros","confianca":0.9,"epistemico":"fato","origem":"gentil/Gentil_Lopes_FUNDAMENTOS_DOS_NUMEROS_pdf.pdf","rel":"define"},{"alvo":"racionais","confianca":0.9,"epistemico":"fato","origem":"gentil/Gentil_Lopes_FUNDAMENTOS_DOS_NUMEROS_pdf.pdf","rel":"define"}],"confianca":1.0,"contraexemplos":[],"descricao":"Capítulo 4 — p.117. Fundamentos dos Números.","epistemico":"fato","exemplos":[],"id":"gentil_fundamentos_capitulo_4_n_umeros_naturais_azuis_e_vermelhos","memoria":"semantica","meta":{"arquivo":"gentil/Lopes_Lopes_FUNDAMENTOS_DOS_NUMEROS_pdf.pdf","autor":"Gentil Lopes da Silva","capitulo":"4","ciencia":"teoria_dos_numeros","dominio":"matematica","nivel":4,"pagina":117,"tipo_no":"paper","titulo_legivel":"Fundamentos dos Números"},"origem":"gentil/Gentil_Lopes_FUNDAMENTOS_DOS_NUMEROS_pdf.pdf","palavras_chave":[],"sinonimos":["capítulo 4"],"tipo":"conceito","titulo":"gentil fundamentos capítulo 4 — N ´UMEROS NATURAIS AZUIS E VERMELHOS"}
\section{gentil fundamentos capítulo 4 — N 'UMEROS NATURAIS AZUIS E VERMELHOS}
Capítulo 4 — p.117. Fundamentos dos Números.
\prov{sinónimos: capítulo 4}
\prov{relações: parte\_de gentil\_fundamentos\_dos\_numeros; define inteiros; define racionais}
\prov{proveniência: gentil/Gentil\_Lopes\_FUNDAMENTOS\_DOS\_NUMEROS\_pdf.pdf \textperiodcentered{} fato \textperiodcentered{} confiança 1.0 \textperiodcentered{} domínio matematica \textperiodcentered{} história 11 versões no jornal}

%CRISTAL {"arestas":[{"alvo":"gentil_fundamentos_dos_numeros","confianca":1.0,"epistemico":"fato","origem":"gentil/Gentil_Lopes_FUNDAMENTOS_DOS_NUMEROS_pdf.pdf","rel":"parte_de"},{"alvo":"inteiros","confianca":0.9,"epistemico":"fato","origem":"gentil/Gentil_Lopes_FUNDAMENTOS_DOS_NUMEROS_pdf.pdf","rel":"define"},{"alvo":"word","confianca":0.5,"epistemico":"fato","origem":"wiki","rel":"relaciona"},{"alvo":"virtualizacao","confianca":0.5,"epistemico":"fato","origem":"wiki","rel":"relaciona"},{"alvo":"syscall","confianca":0.5,"epistemico":"fato","origem":"wiki","rel":"relaciona"}],"confianca":1.0,"contraexemplos":[],"descricao":"Capítulo 5 — p.147. Fundamentos dos Números.","epistemico":"fato","exemplos":[],"id":"gentil_fundamentos_capitulo_5_n_umeros_inteiros","memoria":"semantica","meta":{"arquivo":"gentil/Lopes_Lopes_FUNDAMENTOS_DOS_NUMEROS_pdf.pdf","autor":"Gentil Lopes da Silva","capitulo":"5","ciencia":"teoria_dos_numeros","dominio":"matematica","nivel":4,"pagina":147,"tipo_no":"paper","titulo_legivel":"Fundamentos dos Números"},"origem":"gentil/Gentil_Lopes_FUNDAMENTOS_DOS_NUMEROS_pdf.pdf","palavras_chave":[],"sinonimos":["capítulo 5"],"tipo":"conceito","titulo":"gentil fundamentos capítulo 5 — N ´UMEROS INTEIROS"}
\section{gentil fundamentos capítulo 5 — N 'UMEROS INTEIROS}
Capítulo 5 — p.147. Fundamentos dos Números.
\prov{sinónimos: capítulo 5}
\prov{relações: parte\_de gentil\_fundamentos\_dos\_numeros; define inteiros; relaciona word; relaciona virtualizacao; relaciona syscall}
\prov{proveniência: gentil/Gentil\_Lopes\_FUNDAMENTOS\_DOS\_NUMEROS\_pdf.pdf \textperiodcentered{} fato \textperiodcentered{} confiança 1.0 \textperiodcentered{} domínio matematica \textperiodcentered{} história 10 versões no jornal}

%CRISTAL {"arestas":[{"alvo":"gentil_fundamentos_dos_numeros","confianca":1.0,"epistemico":"fato","origem":"gentil/Gentil_Lopes_FUNDAMENTOS_DOS_NUMEROS_pdf.pdf","rel":"parte_de"},{"alvo":"inteiros","confianca":0.9,"epistemico":"fato","origem":"gentil/Gentil_Lopes_FUNDAMENTOS_DOS_NUMEROS_pdf.pdf","rel":"define"},{"alvo":"racionais","confianca":0.9,"epistemico":"fato","origem":"gentil/Gentil_Lopes_FUNDAMENTOS_DOS_NUMEROS_pdf.pdf","rel":"define"}],"confianca":1.0,"contraexemplos":[],"descricao":"Capítulo 6 — p.209. Fundamentos dos Números.","epistemico":"fato","exemplos":[],"id":"gentil_fundamentos_capitulo_6_inteiros_azuis_e_vermelhos","memoria":"semantica","meta":{"arquivo":"gentil/Lopes_Lopes_FUNDAMENTOS_DOS_NUMEROS_pdf.pdf","autor":"Gentil Lopes da Silva","capitulo":"6","ciencia":"teoria_dos_numeros","dominio":"matematica","nivel":4,"pagina":209,"tipo_no":"paper","titulo_legivel":"Fundamentos dos Números"},"origem":"gentil/Gentil_Lopes_FUNDAMENTOS_DOS_NUMEROS_pdf.pdf","palavras_chave":[],"sinonimos":["capítulo 6"],"tipo":"conceito","titulo":"gentil fundamentos capítulo 6 — INTEIROS AZUIS E VERMELHOS"}
\section{gentil fundamentos capítulo 6 — INTEIROS AZUIS E VERMELHOS}
Capítulo 6 — p.209. Fundamentos dos Números.
\prov{sinónimos: capítulo 6}
\prov{relações: parte\_de gentil\_fundamentos\_dos\_numeros; define inteiros; define racionais}
\prov{proveniência: gentil/Gentil\_Lopes\_FUNDAMENTOS\_DOS\_NUMEROS\_pdf.pdf \textperiodcentered{} fato \textperiodcentered{} confiança 1.0 \textperiodcentered{} domínio matematica \textperiodcentered{} história 9 versões no jornal}

%CRISTAL {"arestas":[{"alvo":"gentil_fundamentos_dos_numeros","confianca":1.0,"epistemico":"fato","origem":"gentil/Gentil_Lopes_FUNDAMENTOS_DOS_NUMEROS_pdf.pdf","rel":"parte_de"},{"alvo":"racionais","confianca":0.9,"epistemico":"fato","origem":"gentil/Gentil_Lopes_FUNDAMENTOS_DOS_NUMEROS_pdf.pdf","rel":"define"},{"alvo":"word","confianca":0.5,"epistemico":"fato","origem":"wiki","rel":"relaciona"}],"confianca":1.0,"contraexemplos":[],"descricao":"Capítulo 7 — p.219. Fundamentos dos Números.","epistemico":"fato","exemplos":[],"id":"gentil_fundamentos_capitulo_7_n_umeros_racionais","memoria":"semantica","meta":{"arquivo":"gentil/Lopes_Lopes_FUNDAMENTOS_DOS_NUMEROS_pdf.pdf","autor":"Gentil Lopes da Silva","capitulo":"7","ciencia":"teoria_dos_numeros","dominio":"matematica","nivel":4,"pagina":219,"tipo_no":"paper","titulo_legivel":"Fundamentos dos Números"},"origem":"gentil/Gentil_Lopes_FUNDAMENTOS_DOS_NUMEROS_pdf.pdf","palavras_chave":[],"sinonimos":["capítulo 7"],"tipo":"conceito","titulo":"gentil fundamentos capítulo 7 — N ´UMEROS RACIONAIS"}
\section{gentil fundamentos capítulo 7 — N 'UMEROS RACIONAIS}
Capítulo 7 — p.219. Fundamentos dos Números.
\prov{sinónimos: capítulo 7}
\prov{relações: parte\_de gentil\_fundamentos\_dos\_numeros; define racionais; relaciona word}
\prov{proveniência: gentil/Gentil\_Lopes\_FUNDAMENTOS\_DOS\_NUMEROS\_pdf.pdf \textperiodcentered{} fato \textperiodcentered{} confiança 1.0 \textperiodcentered{} domínio matematica \textperiodcentered{} história 8 versões no jornal}

%CRISTAL {"arestas":[{"alvo":"gentil_fundamentos_dos_numeros","confianca":1.0,"epistemico":"fato","origem":"gentil/Gentil_Lopes_FUNDAMENTOS_DOS_NUMEROS_pdf.pdf","rel":"parte_de"},{"alvo":"inteiros","confianca":0.9,"epistemico":"fato","origem":"gentil/Gentil_Lopes_FUNDAMENTOS_DOS_NUMEROS_pdf.pdf","rel":"define"},{"alvo":"racionais","confianca":0.9,"epistemico":"fato","origem":"gentil/Gentil_Lopes_FUNDAMENTOS_DOS_NUMEROS_pdf.pdf","rel":"define"}],"confianca":1.0,"contraexemplos":[],"descricao":"Capítulo 8 — p.261. Fundamentos dos Números.","epistemico":"fato","exemplos":[],"id":"gentil_fundamentos_capitulo_8_racionais_azuis_e_vermelhos","memoria":"semantica","meta":{"arquivo":"gentil/Lopes_Lopes_FUNDAMENTOS_DOS_NUMEROS_pdf.pdf","autor":"Gentil Lopes da Silva","capitulo":"8","ciencia":"teoria_dos_numeros","dominio":"matematica","nivel":4,"pagina":261,"tipo_no":"paper","titulo_legivel":"Fundamentos dos Números"},"origem":"gentil/Gentil_Lopes_FUNDAMENTOS_DOS_NUMEROS_pdf.pdf","palavras_chave":[],"sinonimos":["capítulo 8"],"tipo":"conceito","titulo":"gentil fundamentos capítulo 8 — RACIONAIS AZUIS E VERMELHOS"}
\section{gentil fundamentos capítulo 8 — RACIONAIS AZUIS E VERMELHOS}
Capítulo 8 — p.261. Fundamentos dos Números.
\prov{sinónimos: capítulo 8}
\prov{relações: parte\_de gentil\_fundamentos\_dos\_numeros; define inteiros; define racionais}
\prov{proveniência: gentil/Gentil\_Lopes\_FUNDAMENTOS\_DOS\_NUMEROS\_pdf.pdf \textperiodcentered{} fato \textperiodcentered{} confiança 1.0 \textperiodcentered{} domínio matematica \textperiodcentered{} história 9 versões no jornal}

%CRISTAL {"arestas":[{"alvo":"gentil_fundamentos_dos_numeros","confianca":1.0,"epistemico":"fato","origem":"gentil/Gentil_Lopes_FUNDAMENTOS_DOS_NUMEROS_pdf.pdf","rel":"parte_de"},{"alvo":"reais","confianca":0.9,"epistemico":"fato","origem":"gentil/Gentil_Lopes_FUNDAMENTOS_DOS_NUMEROS_pdf.pdf","rel":"define"},{"alvo":"word","confianca":0.5,"epistemico":"fato","origem":"wiki","rel":"relaciona"}],"confianca":1.0,"contraexemplos":[],"descricao":"Capítulo 9 — p.287. Fundamentos dos Números.","epistemico":"fato","exemplos":[],"id":"gentil_fundamentos_capitulo_9_n_umeros_reais_por_dedekind","memoria":"semantica","meta":{"arquivo":"gentil/Lopes_Lopes_FUNDAMENTOS_DOS_NUMEROS_pdf.pdf","autor":"Gentil Lopes da Silva","capitulo":"9","ciencia":"teoria_dos_numeros","dominio":"matematica","nivel":4,"pagina":287,"tipo_no":"paper","titulo_legivel":"Fundamentos dos Números"},"origem":"gentil/Gentil_Lopes_FUNDAMENTOS_DOS_NUMEROS_pdf.pdf","palavras_chave":[],"sinonimos":["capítulo 9"],"tipo":"conceito","titulo":"gentil fundamentos capítulo 9 — N ´UMEROS REAIS POR DEDEKIND"}
\section{gentil fundamentos capítulo 9 — N 'UMEROS REAIS POR DEDEKIND}
Capítulo 9 — p.287. Fundamentos dos Números.
\prov{sinónimos: capítulo 9}
\prov{relações: parte\_de gentil\_fundamentos\_dos\_numeros; define reais; relaciona word}
\prov{proveniência: gentil/Gentil\_Lopes\_FUNDAMENTOS\_DOS\_NUMEROS\_pdf.pdf \textperiodcentered{} fato \textperiodcentered{} confiança 1.0 \textperiodcentered{} domínio matematica \textperiodcentered{} história 8 versões no jornal}

%CRISTAL {"arestas":[{"alvo":"word","confianca":0.5,"epistemico":"fato","origem":"wiki","rel":"relaciona"},{"alvo":"virtualizacao","confianca":0.5,"epistemico":"fato","origem":"wiki","rel":"relaciona"}],"confianca":1.0,"contraexemplos":[],"descricao":", p 63 Definição 29, p. 156 Teorema 2, p. 69","epistemico":"fato","exemplos":[],"id":"gentil_fundamentos_definicao_10_p_63_definicao_29_p_156_teorema_2_p_69","memoria":"semantica","meta":{"arquivo":"gentil/Lopes_Lopes_FUNDAMENTOS_DOS_NUMEROS_pdf.pdf","autor":"Gentil Lopes da Silva","ciencia":"teoria_dos_numeros","dominio":"matematica","nivel":4,"numero":"10","tipo_no":"paper","tipo_teorema":"definição","titulo_legivel":"Fundamentos dos Números"},"origem":"gentil/Gentil_Lopes_FUNDAMENTOS_DOS_NUMEROS_pdf.pdf","palavras_chave":[],"sinonimos":[],"tipo":"conceito","titulo":"gentil fundamentos Definição 10 — , p 63 Definição 29, p. 156 Teorema 2, p. 69"}
\section{gentil fundamentos Definição 10 — , p 63 Definição 29, p. 156 Teorema 2, p. 69}
, p 63 Definição 29, p. 156 Teorema 2, p. 69
\prov{relações: relaciona word; relaciona virtualizacao}
\prov{proveniência: gentil/Gentil\_Lopes\_FUNDAMENTOS\_DOS\_NUMEROS\_pdf.pdf \textperiodcentered{} fato \textperiodcentered{} confiança 1.0 \textperiodcentered{} domínio matematica \textperiodcentered{} história 7 versões no jornal}

%CRISTAL {"arestas":[{"alvo":"matematica","confianca":0.95,"epistemico":"fato","origem":"paper","rel":"parte_de"}],"confianca":1.0,"contraexemplos":[],"descricao":"Classe de equivalência) Seja ∼= \u0000A, A, p(x, y) \u0001 uma relação de equivalência em um conjunto A. Fixado arbitrariamente um ele- mento a ∈A, chama-se classe de equivalência determinada por a, módulo ∼, o subconjunto ¯a de A constituido pelos elementos x tais que x ∼a. Em s´ımbolos: ¯a = x ∈A: x ∼a Resumindo: a classe de equivalência de um elemento a ∈A reuni todos os elementos que se relacionam com a (que são equivalentes a a, diriamos), segundo o critério ∼escolhido. Nota: Pela propriedade ref","epistemico":"fato","exemplos":[],"id":"gentil_fundamentos_definicao_11_classe_de_equivalencia_seja_a_a_px_y_uma_relacao","memoria":"semantica","meta":{"arquivo":"gentil/Lopes_Lopes_FUNDAMENTOS_DOS_NUMEROS_pdf.pdf","autor":"Gentil Lopes da Silva","ciencia":"teoria_dos_numeros","dominio":"matematica","nivel":4,"numero":"11","tipo_no":"paper","tipo_teorema":"definição","titulo_legivel":"Fundamentos dos Números"},"origem":"gentil/Gentil_Lopes_FUNDAMENTOS_DOS_NUMEROS_pdf.pdf","palavras_chave":[],"sinonimos":[],"tipo":"conceito","titulo":"gentil fundamentos Definição 11 — Classe de equivalência) Seja ∼= \u0000A, A, p(x, y) \u0001 uma relação"}
\section{gentil fundamentos Definição 11 — Classe de equivalência) Seja \ensuremath{\sim}= A, A, p(x, y)  uma relação}
Classe de equivalência) Seja \ensuremath{\sim}= A, A, p(x, y)  uma relação de equivalência em um conjunto A. Fixado arbitrariamente um ele- mento a \ensuremath{\in}A, chama-se classe de equivalência determinada por a, módulo \ensuremath{\sim}, o subconjunto \ensuremath{\bar{\,}}a de A constituido pelos elementos x tais que x \ensuremath{\sim}a. Em s'\ensuremath{\imath}mbolos: \ensuremath{\bar{\,}}a = x \ensuremath{\in}A: x \ensuremath{\sim}a Resumindo: a classe de equivalência de um elemento a \ensuremath{\in}A reuni todos os elementos que se relacionam com a (que são equivalentes a a, diriamos), segundo o critério \ensuremath{\sim}escolhido. Nota: Pela propriedade ref
\prov{relações: parte\_de matematica}
\prov{proveniência: gentil/Gentil\_Lopes\_FUNDAMENTOS\_DOS\_NUMEROS\_pdf.pdf \textperiodcentered{} fato \textperiodcentered{} confiança 1.0 \textperiodcentered{} domínio matematica \textperiodcentered{} história 5 versões no jornal}

%CRISTAL {"arestas":[{"alvo":"matematica","confianca":0.95,"epistemico":"fato","origem":"paper","rel":"parte_de"}],"confianca":1.0,"contraexemplos":[],"descricao":"p 66), temos: ¯a = x ∈A: x ∼a. Traduzindo para o nosso contexto: ¯a = x ∈A: x ∼a (0, 1) = (x, y) ∈N × N: (x, y) ∼(0, 1) ↕ ↕ ↕ Então, pela","epistemico":"fato","exemplos":[],"id":"gentil_fundamentos_definicao_11_p_66_temos_a_x_a_x_a_traduzindo_para_o_nosso_c","memoria":"semantica","meta":{"arquivo":"gentil/Lopes_Lopes_FUNDAMENTOS_DOS_NUMEROS_pdf.pdf","autor":"Gentil Lopes da Silva","ciencia":"teoria_dos_numeros","dominio":"matematica","nivel":4,"numero":"11","tipo_no":"paper","tipo_teorema":"definição","titulo_legivel":"Fundamentos dos Números"},"origem":"gentil/Gentil_Lopes_FUNDAMENTOS_DOS_NUMEROS_pdf.pdf","palavras_chave":[],"sinonimos":[],"tipo":"conceito","titulo":"gentil fundamentos Definição 11 — p 66), temos: ¯a = x ∈A: x ∼a. Traduzindo para o nosso c"}
\section{gentil fundamentos Definição 11 — p 66), temos: \ensuremath{\bar{\,}}a = x \ensuremath{\in}A: x \ensuremath{\sim}a. Traduzindo para o nosso c}
p 66), temos: \ensuremath{\bar{\,}}a = x \ensuremath{\in}A: x \ensuremath{\sim}a. Traduzindo para o nosso contexto: \ensuremath{\bar{\,}}a = x \ensuremath{\in}A: x \ensuremath{\sim}a (0, 1) = (x, y) \ensuremath{\in}N $\times$ N: (x, y) \ensuremath{\sim}(0, 1) \ensuremath{\updownarrow} \ensuremath{\updownarrow} \ensuremath{\updownarrow} Então, pela
\prov{relações: parte\_de matematica}
\prov{proveniência: gentil/Gentil\_Lopes\_FUNDAMENTOS\_DOS\_NUMEROS\_pdf.pdf \textperiodcentered{} fato \textperiodcentered{} confiança 1.0 \textperiodcentered{} domínio matematica \textperiodcentered{} história 6 versões no jornal}

%CRISTAL {"arestas":[{"alvo":"matematica","confianca":0.95,"epistemico":"fato","origem":"paper","rel":"parte_de"}],"confianca":1.0,"contraexemplos":[],"descricao":"p 71) devemos provar que: i ) duas classes de equivalência quaisquer de A/∼ou são iguais ou são disjuntas. Isto já foi feito no","epistemico":"fato","exemplos":[],"id":"gentil_fundamentos_definicao_13_p_71_devemos_provar_que_i_duas_classes_de_equiva","memoria":"semantica","meta":{"arquivo":"gentil/Lopes_Lopes_FUNDAMENTOS_DOS_NUMEROS_pdf.pdf","autor":"Gentil Lopes da Silva","ciencia":"teoria_dos_numeros","dominio":"matematica","nivel":4,"numero":"13","tipo_no":"paper","tipo_teorema":"definição","titulo_legivel":"Fundamentos dos Números"},"origem":"gentil/Gentil_Lopes_FUNDAMENTOS_DOS_NUMEROS_pdf.pdf","palavras_chave":[],"sinonimos":[],"tipo":"conceito","titulo":"gentil fundamentos Definição 13 — p 71) devemos provar que: i ) duas classes de equivalência q"}
\section{gentil fundamentos Definição 13 — p 71) devemos provar que: i ) duas classes de equivalência q}
p 71) devemos provar que: i ) duas classes de equivalência quaisquer de A/\ensuremath{\sim}ou são iguais ou são disjuntas. Isto já foi feito no
\prov{relações: parte\_de matematica}
\prov{proveniência: gentil/Gentil\_Lopes\_FUNDAMENTOS\_DOS\_NUMEROS\_pdf.pdf \textperiodcentered{} fato \textperiodcentered{} confiança 1.0 \textperiodcentered{} domínio matematica \textperiodcentered{} história 4 versões no jornal}

%CRISTAL {"arestas":[{"alvo":"matematica","confianca":0.95,"epistemico":"fato","origem":"paper","rel":"parte_de"}],"confianca":1.0,"contraexemplos":[],"descricao":"p 76. 2a ) Ordem total: Definição 17, p. 78. 3a ) Compatibilidade com a Adiç~ao a ≤b ⇔a + c ≤b + c 4a ) Compatibilidade com a Multiplicaç~ao a ≤b ⇒ a · c ≤b · c 91 Vamos juntar estas nove propriedades em um conjunto denotado por ρ: ρ = A1, A2, A3, M1, M2, M3, D, Ordenado, PBO (3.4) 2a ) Apenas lembrando: assim como o que caracteriza o xadrez é o con- junto de suas regras −e não suas peças −similarmente o que caracteriza os números naturais é o conjunto de regras listadas em ρ, estas são as","epistemico":"fato","exemplos":[],"id":"gentil_fundamentos_definicao_14_p_76_2a_ordem_total_definicao_17_p_78_3a_compati","memoria":"semantica","meta":{"arquivo":"gentil/Lopes_Lopes_FUNDAMENTOS_DOS_NUMEROS_pdf.pdf","autor":"Gentil Lopes da Silva","ciencia":"teoria_dos_numeros","dominio":"matematica","nivel":4,"numero":"14","tipo_no":"paper","tipo_teorema":"definição","titulo_legivel":"Fundamentos dos Números"},"origem":"gentil/Gentil_Lopes_FUNDAMENTOS_DOS_NUMEROS_pdf.pdf","palavras_chave":[],"sinonimos":[],"tipo":"conceito","titulo":"gentil fundamentos Definição 14 — , p 76. 2a ) Ordem total: Definição 17, p. 78. 3a ) Compatib"}
\section{gentil fundamentos Definição 14 — , p 76. 2a ) Ordem total: Definição 17, p. 78. 3a ) Compatib}
p 76. 2a ) Ordem total: Definição 17, p. 78. 3a ) Compatibilidade com a Adiç\~{}ao a \ensuremath{\le}b \ensuremath{\Leftrightarrow}a + c \ensuremath{\le}b + c 4a ) Compatibilidade com a Multiplicaç\~{}ao a \ensuremath{\le}b \ensuremath{\Rightarrow} a \textperiodcentered  c \ensuremath{\le}b \textperiodcentered  c 91 Vamos juntar estas nove propriedades em um conjunto denotado por \ensuremath{\rho}: \ensuremath{\rho} = A1, A2, A3, M1, M2, M3, D, Ordenado, PBO (3.4) 2a ) Apenas lembrando: assim como o que caracteriza o xadrez é o con- junto de suas regras \ensuremath{-}e não suas peças \ensuremath{-}similarmente o que caracteriza os números naturais é o conjunto de regras listadas em \ensuremath{\rho}, estas são as
\prov{relações: parte\_de matematica}
\prov{proveniência: gentil/Gentil\_Lopes\_FUNDAMENTOS\_DOS\_NUMEROS\_pdf.pdf \textperiodcentered{} fato \textperiodcentered{} confiança 1.0 \textperiodcentered{} domínio matematica \textperiodcentered{} história 5 versões no jornal}

%CRISTAL {"arestas":[{"alvo":"matematica","confianca":0.95,"epistemico":"fato","origem":"paper","rel":"parte_de"}],"confianca":1.0,"contraexemplos":[],"descricao":"Relação de ordem parcial) Uma relação R sobre um con- junto A não vazio é chamada relação de ordem parcial sobre A se, e somente se, R é reflexiva, antissimétrica e transitiva; isto é, se são verdadeiras as sentenças: ( i ) ( ∀x ∈A ) ( x R x ) (Reflexiva) ( ii ) ( ∀x, y ∈A) \u0000se x R y e y R x ⇒x = y \u0001 (Antissimétrica) ( iii ) ( ∀x, y, z ∈A) ( se x R y e y R z ⇒ x R z ) (Transitiva) Mudança de notação: No caso especial das relações de ordem faremos − quando conveniente −uma mudança de notação (a n","epistemico":"fato","exemplos":[],"id":"gentil_fundamentos_definicao_14_relacao_de_ordem_parcial_uma_relacao_r_sobre_um_","memoria":"semantica","meta":{"arquivo":"gentil/Lopes_Lopes_FUNDAMENTOS_DOS_NUMEROS_pdf.pdf","autor":"Gentil Lopes da Silva","ciencia":"teoria_dos_numeros","dominio":"matematica","nivel":4,"numero":"14","tipo_no":"paper","tipo_teorema":"definição","titulo_legivel":"Fundamentos dos Números"},"origem":"gentil/Gentil_Lopes_FUNDAMENTOS_DOS_NUMEROS_pdf.pdf","palavras_chave":[],"sinonimos":[],"tipo":"conceito","titulo":"gentil fundamentos Definição 14 — Relação de ordem parcial) Uma relação R sobre um con- junto"}
\section{gentil fundamentos Definição 14 — Relação de ordem parcial) Uma relação R sobre um con- junto}
Relação de ordem parcial) Uma relação R sobre um con- junto A não vazio é chamada relação de ordem parcial sobre A se, e somente se, R é reflexiva, antissimétrica e transitiva; isto é, se são verdadeiras as sentenças: ( i ) ( \ensuremath{\forall}x \ensuremath{\in}A ) ( x R x ) (Reflexiva) ( ii ) ( \ensuremath{\forall}x, y \ensuremath{\in}A) se x R y e y R x \ensuremath{\Rightarrow}x = y  (Antissimétrica) ( iii ) ( \ensuremath{\forall}x, y, z \ensuremath{\in}A) ( se x R y e y R z \ensuremath{\Rightarrow} x R z ) (Transitiva) Mudança de notação: No caso especial das relações de ordem faremos \ensuremath{-} quando conveniente \ensuremath{-}uma mudança de notação (a n
\prov{relações: parte\_de matematica}
\prov{proveniência: gentil/Gentil\_Lopes\_FUNDAMENTOS\_DOS\_NUMEROS\_pdf.pdf \textperiodcentered{} fato \textperiodcentered{} confiança 1.0 \textperiodcentered{} domínio matematica \textperiodcentered{} história 4 versões no jornal}

%CRISTAL {"arestas":[{"alvo":"matematica","confianca":0.95,"epistemico":"fato","origem":"paper","rel":"parte_de"}],"confianca":1.0,"contraexemplos":[],"descricao":"Um conjunto A, com uma relação R (ou ≤) de ordem parcial espec´ıfica em A, é chamado conjunto parcialmente ordenado. Nota: A bem da verdade, ao dotarmos um conjunto A de uma ordem pas- samos a ter uma estrutura: um conjunto ordenado; por esta razão, um conjunto parcialmente ordenado é algumas vezes designado pelo par orde- nado (A, R) ou (A, ≤)","epistemico":"fato","exemplos":[],"id":"gentil_fundamentos_definicao_15_um_conjunto_a_com_uma_relacao_r_ou_de_ordem_parc","memoria":"semantica","meta":{"arquivo":"gentil/Lopes_Lopes_FUNDAMENTOS_DOS_NUMEROS_pdf.pdf","autor":"Gentil Lopes da Silva","ciencia":"teoria_dos_numeros","dominio":"matematica","nivel":4,"numero":"15","tipo_no":"paper","tipo_teorema":"definição","titulo_legivel":"Fundamentos dos Números"},"origem":"gentil/Gentil_Lopes_FUNDAMENTOS_DOS_NUMEROS_pdf.pdf","palavras_chave":[],"sinonimos":[],"tipo":"conceito","titulo":"gentil fundamentos Definição 15 — Um conjunto A, com uma relação R (ou ≤) de ordem parcial esp"}
\section{gentil fundamentos Definição 15 — Um conjunto A, com uma relação R (ou \ensuremath{\le}) de ordem parcial esp}
Um conjunto A, com uma relação R (ou \ensuremath{\le}) de ordem parcial espec'\ensuremath{\imath}fica em A, é chamado conjunto parcialmente ordenado. Nota: A bem da verdade, ao dotarmos um conjunto A de uma ordem pas- samos a ter uma estrutura: um conjunto ordenado; por esta razão, um conjunto parcialmente ordenado é algumas vezes designado pelo par orde- nado (A, R) ou (A, \ensuremath{\le})
\prov{relações: parte\_de matematica}
\prov{proveniência: gentil/Gentil\_Lopes\_FUNDAMENTOS\_DOS\_NUMEROS\_pdf.pdf \textperiodcentered{} fato \textperiodcentered{} confiança 1.0 \textperiodcentered{} domínio matematica \textperiodcentered{} história 4 versões no jornal}

%CRISTAL {"arestas":[{"alvo":"matematica","confianca":0.95,"epistemico":"fato","origem":"paper","rel":"parte_de"}],"confianca":1.0,"contraexemplos":[],"descricao":"Elementos comparáveis) Sejam A um conjunto e R uma relação de ordem parcial em A. Dois elementos quaisquer x, y ∈A dizem-se comparáveis pela ordem R se e somente se uma das sentenças x R y ou y R x é verdadeira. Exemplo: a) Os números naturais 3 e 5 são comparáveis pela ordem natural “x ≤y” em N∗, porque 3 ≤5, mas não são comparáveis pela ordem parcial “x|y” em N∗, visto que 3 não divide 5 e 5 não divide 3. Conjuntos Totalmente Ordenados Como vimos anteriormente, numa ordem parcial em um conjunt","epistemico":"fato","exemplos":[],"id":"gentil_fundamentos_definicao_16_elementos_comparaveis_sejam_a_um_conjunto_e_r_um","memoria":"semantica","meta":{"arquivo":"gentil/Lopes_Lopes_FUNDAMENTOS_DOS_NUMEROS_pdf.pdf","autor":"Gentil Lopes da Silva","ciencia":"teoria_dos_numeros","dominio":"matematica","nivel":4,"numero":"16","tipo_no":"paper","tipo_teorema":"definição","titulo_legivel":"Fundamentos dos Números"},"origem":"gentil/Gentil_Lopes_FUNDAMENTOS_DOS_NUMEROS_pdf.pdf","palavras_chave":[],"sinonimos":[],"tipo":"conceito","titulo":"gentil fundamentos Definição 16 — Elementos comparáveis) Sejam A um conjunto e R uma relação d"}
\section{gentil fundamentos Definição 16 — Elementos comparáveis) Sejam A um conjunto e R uma relação d}
Elementos comparáveis) Sejam A um conjunto e R uma relação de ordem parcial em A. Dois elementos quaisquer x, y \ensuremath{\in}A dizem-se comparáveis pela ordem R se e somente se uma das sentenças x R y ou y R x é verdadeira. Exemplo: a) Os números naturais 3 e 5 são comparáveis pela ordem natural “x \ensuremath{\le}y” em N\ensuremath{*}, porque 3 \ensuremath{\le}5, mas não são comparáveis pela ordem parcial “x|y” em N\ensuremath{*}, visto que 3 não divide 5 e 5 não divide 3. Conjuntos Totalmente Ordenados Como vimos anteriormente, numa ordem parcial em um conjunt
\prov{relações: parte\_de matematica}
\prov{proveniência: gentil/Gentil\_Lopes\_FUNDAMENTOS\_DOS\_NUMEROS\_pdf.pdf \textperiodcentered{} fato \textperiodcentered{} confiança 1.0 \textperiodcentered{} domínio matematica \textperiodcentered{} história 4 versões no jornal}

%CRISTAL {"arestas":[{"alvo":"matematica","confianca":0.95,"epistemico":"fato","origem":"paper","rel":"parte_de"}],"confianca":1.0,"contraexemplos":[],"descricao":"Adição, p 92) Lema 1, p. 94 Lei do Corte (+) (Prop. 5, p. 97) Lei da Tricotomia (Teo. 11, p. 107)","epistemico":"fato","exemplos":[],"id":"gentil_fundamentos_definicao_18_adicao_p_92_lema_1_p_94_lei_do_corte_prop_5_p_97","memoria":"semantica","meta":{"arquivo":"gentil/Lopes_Lopes_FUNDAMENTOS_DOS_NUMEROS_pdf.pdf","autor":"Gentil Lopes da Silva","ciencia":"teoria_dos_numeros","dominio":"matematica","nivel":4,"numero":"18","tipo_no":"paper","tipo_teorema":"definição","titulo_legivel":"Fundamentos dos Números"},"origem":"gentil/Gentil_Lopes_FUNDAMENTOS_DOS_NUMEROS_pdf.pdf","palavras_chave":[],"sinonimos":[],"tipo":"conceito","titulo":"gentil fundamentos Definição 18 — Adição, p 92) Lema 1, p. 94 Lei do Corte (+) (Prop. 5, p. 97"}
\section{gentil fundamentos Definição 18 — Adição, p 92) Lema 1, p. 94 Lei do Corte (+) (Prop. 5, p. 97}
Adição, p 92) Lema 1, p. 94 Lei do Corte (+) (Prop. 5, p. 97) Lei da Tricotomia (Teo. 11, p. 107)
\prov{relações: parte\_de matematica}
\prov{proveniência: gentil/Gentil\_Lopes\_FUNDAMENTOS\_DOS\_NUMEROS\_pdf.pdf \textperiodcentered{} fato \textperiodcentered{} confiança 1.0 \textperiodcentered{} domínio matematica \textperiodcentered{} história 4 versões no jornal}

%CRISTAL {"arestas":[{"alvo":"matematica","confianca":0.95,"epistemico":"fato","origem":"paper","rel":"parte_de"}],"confianca":1.0,"contraexemplos":[],"descricao":"Adição) Seja n ∈N um natural dado. Então ( i ) m + 0 = m; ( ii ) m + σ(n) = σ(m + n).","epistemico":"fato","exemplos":[],"id":"gentil_fundamentos_definicao_18_adicao_seja_n_n_um_natural_dado_entao_i_m_0_m","memoria":"semantica","meta":{"arquivo":"gentil/Lopes_Lopes_FUNDAMENTOS_DOS_NUMEROS_pdf.pdf","autor":"Gentil Lopes da Silva","ciencia":"teoria_dos_numeros","dominio":"matematica","nivel":4,"numero":"18","tipo_no":"paper","tipo_teorema":"definição","titulo_legivel":"Fundamentos dos Números"},"origem":"gentil/Gentil_Lopes_FUNDAMENTOS_DOS_NUMEROS_pdf.pdf","palavras_chave":[],"sinonimos":[],"tipo":"conceito","titulo":"gentil fundamentos Definição 18 — Adição) Seja n ∈N um natural dado. Então ( i ) m + 0 = m; ("}
\section{gentil fundamentos Definição 18 — Adição) Seja n \ensuremath{\in}N um natural dado. Então ( i ) m + 0 = m; (}
Adição) Seja n \ensuremath{\in}N um natural dado. Então ( i ) m + 0 = m; ( ii ) m + \ensuremath{\sigma}(n) = \ensuremath{\sigma}(m + n).
\prov{relações: parte\_de matematica}
\prov{proveniência: gentil/Gentil\_Lopes\_FUNDAMENTOS\_DOS\_NUMEROS\_pdf.pdf \textperiodcentered{} fato \textperiodcentered{} confiança 1.0 \textperiodcentered{} domínio matematica \textperiodcentered{} história 4 versões no jornal}

%CRISTAL {"arestas":[{"alvo":"matematica","confianca":0.95,"epistemico":"fato","origem":"paper","rel":"parte_de"},{"alvo":"virtualizacao","confianca":0.5,"epistemico":"fato","origem":"wiki","rel":"relaciona"},{"alvo":"word","confianca":0.5,"epistemico":"fato","origem":"wiki","rel":"relaciona"}],"confianca":1.0,"contraexemplos":[],"descricao":"p 92) e, finalmente, na terceira igualdade utili- zamos (i) da","epistemico":"fato","exemplos":[],"id":"gentil_fundamentos_definicao_18_p_92_e_finalmente_na_terceira_igualdade_utili-_z","memoria":"semantica","meta":{"arquivo":"gentil/Lopes_Lopes_FUNDAMENTOS_DOS_NUMEROS_pdf.pdf","autor":"Gentil Lopes da Silva","ciencia":"teoria_dos_numeros","dominio":"matematica","nivel":4,"numero":"18","tipo_no":"paper","tipo_teorema":"definição","titulo_legivel":"Fundamentos dos Números"},"origem":"gentil/Gentil_Lopes_FUNDAMENTOS_DOS_NUMEROS_pdf.pdf","palavras_chave":[],"sinonimos":[],"tipo":"conceito","titulo":"gentil fundamentos Definição 18 — p 92) e, finalmente, na terceira igualdade utili- zamos (i)"}
\section{gentil fundamentos Definição 18 — p 92) e, finalmente, na terceira igualdade utili- zamos (i)}
p 92) e, finalmente, na terceira igualdade utili- zamos (i) da
\prov{relações: parte\_de matematica; relaciona virtualizacao; relaciona word}
\prov{proveniência: gentil/Gentil\_Lopes\_FUNDAMENTOS\_DOS\_NUMEROS\_pdf.pdf \textperiodcentered{} fato \textperiodcentered{} confiança 1.0 \textperiodcentered{} domínio matematica \textperiodcentered{} história 6 versões no jornal}

%CRISTAL {"arestas":[{"alvo":"matematica","confianca":0.95,"epistemico":"fato","origem":"paper","rel":"parte_de"}],"confianca":1.0,"contraexemplos":[],"descricao":"Para provar a segunda igualdade, consideremos o conjunto A = m ∈N: σ(m) = 1 + m e provemos que A = N. Por (ii) (","epistemico":"fato","exemplos":[],"id":"gentil_fundamentos_definicao_18_para_provar_a_segunda_igualdade_consideremos_o_c","memoria":"semantica","meta":{"arquivo":"gentil/Lopes_Lopes_FUNDAMENTOS_DOS_NUMEROS_pdf.pdf","autor":"Gentil Lopes da Silva","ciencia":"teoria_dos_numeros","dominio":"matematica","nivel":4,"numero":"18","tipo_no":"paper","tipo_teorema":"definição","titulo_legivel":"Fundamentos dos Números"},"origem":"gentil/Gentil_Lopes_FUNDAMENTOS_DOS_NUMEROS_pdf.pdf","palavras_chave":[],"sinonimos":[],"tipo":"conceito","titulo":"gentil fundamentos Definição 18 — Para provar a segunda igualdade, consideremos o conjunto A ="}
\section{gentil fundamentos Definição 18 — Para provar a segunda igualdade, consideremos o conjunto A =}
Para provar a segunda igualdade, consideremos o conjunto A = m \ensuremath{\in}N: \ensuremath{\sigma}(m) = 1 + m e provemos que A = N. Por (ii) (
\prov{relações: parte\_de matematica}
\prov{proveniência: gentil/Gentil\_Lopes\_FUNDAMENTOS\_DOS\_NUMEROS\_pdf.pdf \textperiodcentered{} fato \textperiodcentered{} confiança 1.0 \textperiodcentered{} domínio matematica \textperiodcentered{} história 5 versões no jornal}

%CRISTAL {"arestas":[{"alvo":"matematica","confianca":0.95,"epistemico":"fato","origem":"paper","rel":"parte_de"}],"confianca":1.0,"contraexemplos":[],"descricao":") temos que 1 + 0 = 1 = σ(0), portanto 0 ∈A Considerando, por hipótese, que m ∈A, vamos provar que σ(m) ∈A. Isto é, provemos que σ \u0000σ(m) \u0001 = 1 + σ(m) (3.5) Apliquemos σ à hipótese de indução, σ(m) = 1 + m, isto é, σ \u0000σ(m) \u0001 = σ(1 + m) Por (ii) (","epistemico":"fato","exemplos":[],"id":"gentil_fundamentos_definicao_18_temos_que_1_0_1_σ0_portanto_0_a_considerando_po","memoria":"semantica","meta":{"arquivo":"gentil/Lopes_Lopes_FUNDAMENTOS_DOS_NUMEROS_pdf.pdf","autor":"Gentil Lopes da Silva","ciencia":"teoria_dos_numeros","dominio":"matematica","nivel":4,"numero":"18","tipo_no":"paper","tipo_teorema":"definição","titulo_legivel":"Fundamentos dos Números"},"origem":"gentil/Gentil_Lopes_FUNDAMENTOS_DOS_NUMEROS_pdf.pdf","palavras_chave":[],"sinonimos":[],"tipo":"conceito","titulo":"gentil fundamentos Definição 18 — ) temos que 1 + 0 = 1 = σ(0), portanto 0 ∈A Considerando, po"}
\section{gentil fundamentos Definição 18 — ) temos que 1 + 0 = 1 = \ensuremath{\sigma}(0), portanto 0 \ensuremath{\in}A Considerando, po}
) temos que 1 + 0 = 1 = \ensuremath{\sigma}(0), portanto 0 \ensuremath{\in}A Considerando, por hipótese, que m \ensuremath{\in}A, vamos provar que \ensuremath{\sigma}(m) \ensuremath{\in}A. Isto é, provemos que \ensuremath{\sigma} \ensuremath{\sigma}(m)  = 1 + \ensuremath{\sigma}(m) (3.5) Apliquemos \ensuremath{\sigma} à hipótese de indução, \ensuremath{\sigma}(m) = 1 + m, isto é, \ensuremath{\sigma} \ensuremath{\sigma}(m)  = \ensuremath{\sigma}(1 + m) Por (ii) (
\prov{relações: parte\_de matematica}
\prov{proveniência: gentil/Gentil\_Lopes\_FUNDAMENTOS\_DOS\_NUMEROS\_pdf.pdf \textperiodcentered{} fato \textperiodcentered{} confiança 1.0 \textperiodcentered{} domínio matematica \textperiodcentered{} história 4 versões no jornal}

%CRISTAL {"arestas":[{"alvo":"matematica","confianca":0.95,"epistemico":"fato","origem":"paper","rel":"parte_de"}],"confianca":1.0,"contraexemplos":[],"descricao":"Multiplicação) A multiplicação em N é definida por: ( n · 0 = 0; ∀n ∈N n · (m + 1) = n · m + n, sempre que n · m está definido.","epistemico":"fato","exemplos":[],"id":"gentil_fundamentos_definicao_19_multiplicacao_a_multiplicacao_em_n_e_definida_po","memoria":"semantica","meta":{"arquivo":"gentil/Lopes_Lopes_FUNDAMENTOS_DOS_NUMEROS_pdf.pdf","autor":"Gentil Lopes da Silva","ciencia":"teoria_dos_numeros","dominio":"matematica","nivel":4,"numero":"19","tipo_no":"paper","tipo_teorema":"definição","titulo_legivel":"Fundamentos dos Números"},"origem":"gentil/Gentil_Lopes_FUNDAMENTOS_DOS_NUMEROS_pdf.pdf","palavras_chave":[],"sinonimos":[],"tipo":"conceito","titulo":"gentil fundamentos Definição 19 — Multiplicação) A multiplicação em N é definida por: ( n · 0"}
\section{gentil fundamentos Definição 19 — Multiplicação) A multiplicação em N é definida por: ( n \textperiodcentered  0}
Multiplicação) A multiplicação em N é definida por: ( n \textperiodcentered  0 = 0; \ensuremath{\forall}n \ensuremath{\in}N n \textperiodcentered  (m + 1) = n \textperiodcentered  m + n, sempre que n \textperiodcentered  m está definido.
\prov{relações: parte\_de matematica}
\prov{proveniência: gentil/Gentil\_Lopes\_FUNDAMENTOS\_DOS\_NUMEROS\_pdf.pdf \textperiodcentered{} fato \textperiodcentered{} confiança 1.0 \textperiodcentered{} domínio matematica \textperiodcentered{} história 4 versões no jornal}

%CRISTAL {"arestas":[{"alvo":"gentil_fundamentos_capitulo_1_o_que_e_um_n_umero","confianca":1.0,"epistemico":"fato","origem":"gentil/Gentil_Lopes_FUNDAMENTOS_DOS_NUMEROS_pdf.pdf","rel":"define"}],"confianca":1.0,"contraexemplos":[],"descricao":", (p 92) ( n · 0 = 0; ∀n ∈N n · (m + 1) = n · m + n, sempre que n · m está definido. Observe alguns casos especiais, m = 0 ⇒ n · (0 + 1) = n · 0 + n · 1 ⇒ n · 1 = n. m = 1 ⇒ n · (1 + 1) = n · 1 + n · 1 ⇒ n · 2 = n + n . m = 2 ⇒ n · (2 + 1) = n · 2 + n · 1 ⇒ n · 3 = n · 2 + n = (n + n) + n. Isto é, multiplicar um número n por 0 resulta, por","epistemico":"fato","exemplos":[],"id":"gentil_fundamentos_definicao_19_p_92_n_0_0_n_n_n_m_1_n_m_n_sempre","memoria":"semantica","meta":{"arquivo":"gentil/Lopes_Lopes_FUNDAMENTOS_DOS_NUMEROS_pdf.pdf","autor":"Gentil Lopes da Silva","ciencia":"teoria_dos_numeros","dominio":"matematica","nivel":4,"numero":"19","tipo_no":"paper","tipo_teorema":"definição","titulo_legivel":"Fundamentos dos Números"},"origem":"gentil/Gentil_Lopes_FUNDAMENTOS_DOS_NUMEROS_pdf.pdf","palavras_chave":[],"sinonimos":[],"tipo":"conceito","titulo":"gentil fundamentos Definição 19 — , (p 92) ( n · 0 = 0; ∀n ∈N n · (m + 1) = n · m + n, sempre"}
\section{gentil fundamentos Definição 19 — , (p 92) ( n \textperiodcentered  0 = 0; \ensuremath{\forall}n \ensuremath{\in}N n \textperiodcentered  (m + 1) = n \textperiodcentered  m + n, sempre}
, (p 92) ( n \textperiodcentered  0 = 0; \ensuremath{\forall}n \ensuremath{\in}N n \textperiodcentered  (m + 1) = n \textperiodcentered  m + n, sempre que n \textperiodcentered  m está definido. Observe alguns casos especiais, m = 0 \ensuremath{\Rightarrow} n \textperiodcentered  (0 + 1) = n \textperiodcentered  0 + n \textperiodcentered  1 \ensuremath{\Rightarrow} n \textperiodcentered  1 = n. m = 1 \ensuremath{\Rightarrow} n \textperiodcentered  (1 + 1) = n \textperiodcentered  1 + n \textperiodcentered  1 \ensuremath{\Rightarrow} n \textperiodcentered  2 = n + n . m = 2 \ensuremath{\Rightarrow} n \textperiodcentered  (2 + 1) = n \textperiodcentered  2 + n \textperiodcentered  1 \ensuremath{\Rightarrow} n \textperiodcentered  3 = n \textperiodcentered  2 + n = (n + n) + n. Isto é, multiplicar um número n por 0 resulta, por
\prov{relações: define gentil\_fundamentos\_capitulo\_1\_o\_que\_e\_um\_n\_umero}
\prov{proveniência: gentil/Gentil\_Lopes\_FUNDAMENTOS\_DOS\_NUMEROS\_pdf.pdf \textperiodcentered{} fato \textperiodcentered{} confiança 1.0 \textperiodcentered{} domínio matematica \textperiodcentered{} história 5 versões no jornal}

%CRISTAL {"arestas":[{"alvo":"matematica","confianca":0.95,"epistemico":"fato","origem":"paper","rel":"parte_de"}],"confianca":1.0,"contraexemplos":[],"descricao":"Operação) Sendo E um conjunto não vazio, toda aplicação (função) f : E × E →E recebe o nome de operação sobre E. Para construirmos (erigirmos) um sistema numérico sobre um dado con- junto basta definirmos duas operações sobre este conjunto, uma das quais será chamada de adição e a outra de multiplicação, simbolizadas por + e ·, respectivamente. Mais formalmente,","epistemico":"fato","exemplos":[],"id":"gentil_fundamentos_definicao_1_operacao_sendo_e_um_conjunto_nao_vazio_toda_aplic","memoria":"semantica","meta":{"arquivo":"gentil/Lopes_Lopes_FUNDAMENTOS_DOS_NUMEROS_pdf.pdf","autor":"Gentil Lopes da Silva","ciencia":"teoria_dos_numeros","dominio":"matematica","nivel":4,"numero":"1","tipo_no":"paper","tipo_teorema":"definição","titulo_legivel":"Fundamentos dos Números"},"origem":"gentil/Gentil_Lopes_FUNDAMENTOS_DOS_NUMEROS_pdf.pdf","palavras_chave":[],"sinonimos":[],"tipo":"conceito","titulo":"gentil fundamentos Definição 1 — Operação) Sendo E um conjunto não vazio, toda aplicação (fun"}
\section{gentil fundamentos Definição 1 — Operação) Sendo E um conjunto não vazio, toda aplicação (fun}
Operação) Sendo E um conjunto não vazio, toda aplicação (função) f : E $\times$ E \ensuremath{\to}E recebe o nome de operação sobre E. Para construirmos (erigirmos) um sistema numérico sobre um dado con- junto basta definirmos duas operações sobre este conjunto, uma das quais será chamada de adição e a outra de multiplicação, simbolizadas por + e \textperiodcentered , respectivamente. Mais formalmente,
\prov{relações: parte\_de matematica}
\prov{proveniência: gentil/Gentil\_Lopes\_FUNDAMENTOS\_DOS\_NUMEROS\_pdf.pdf \textperiodcentered{} fato \textperiodcentered{} confiança 1.0 \textperiodcentered{} domínio matematica \textperiodcentered{} história 4 versões no jornal}

%CRISTAL {"arestas":[{"alvo":"matematica","confianca":0.95,"epistemico":"fato","origem":"paper","rel":"parte_de"}],"confianca":1.0,"contraexemplos":[],"descricao":"Indicaremos por 1 (lemos “um”) o número natural que é sucessor de 0, isto é, 1 = σ(0). Ainda por","epistemico":"fato","exemplos":[],"id":"gentil_fundamentos_definicao_21_indicaremos_por_1_lemos_um_o_numero_natural_que_","memoria":"semantica","meta":{"arquivo":"gentil/Lopes_Lopes_FUNDAMENTOS_DOS_NUMEROS_pdf.pdf","autor":"Gentil Lopes da Silva","ciencia":"teoria_dos_numeros","dominio":"matematica","nivel":4,"numero":"21","tipo_no":"paper","tipo_teorema":"definição","titulo_legivel":"Fundamentos dos Números"},"origem":"gentil/Gentil_Lopes_FUNDAMENTOS_DOS_NUMEROS_pdf.pdf","palavras_chave":[],"sinonimos":[],"tipo":"conceito","titulo":"gentil fundamentos Definição 21 — Indicaremos por 1 (lemos “um”) o número natural que é sucess"}
\section{gentil fundamentos Definição 21 — Indicaremos por 1 (lemos “um”) o número natural que é sucess}
Indicaremos por 1 (lemos “um”) o número natural que é sucessor de 0, isto é, 1 = \ensuremath{\sigma}(0). Ainda por
\prov{relações: parte\_de matematica}
\prov{proveniência: gentil/Gentil\_Lopes\_FUNDAMENTOS\_DOS\_NUMEROS\_pdf.pdf \textperiodcentered{} fato \textperiodcentered{} confiança 1.0 \textperiodcentered{} domínio matematica \textperiodcentered{} história 4 versões no jornal}

%CRISTAL {"arestas":[{"alvo":"matematica","confianca":0.95,"epistemico":"fato","origem":"paper","rel":"parte_de"}],"confianca":1.0,"contraexemplos":[],"descricao":"p 105) concluimos que existem r, s ∈N tais que m = n + r e p = s + q logo m + p = (n + r) + (s + q) ⇒ · · · ⇒ m + p = (n + q) + (r + s) Sendo assim, m + p ≤n + q. Por oportuno, vamos provar o seguinte resultado: (iii)’ m < n e p < q ⇒ m + p < n + q; Da","epistemico":"fato","exemplos":[],"id":"gentil_fundamentos_definicao_22_p_105_concluimos_que_existem_r_s_n_tais_que_m_n_","memoria":"semantica","meta":{"arquivo":"gentil/Lopes_Lopes_FUNDAMENTOS_DOS_NUMEROS_pdf.pdf","autor":"Gentil Lopes da Silva","ciencia":"teoria_dos_numeros","dominio":"matematica","nivel":4,"numero":"22","tipo_no":"paper","tipo_teorema":"definição","titulo_legivel":"Fundamentos dos Números"},"origem":"gentil/Gentil_Lopes_FUNDAMENTOS_DOS_NUMEROS_pdf.pdf","palavras_chave":[],"sinonimos":[],"tipo":"conceito","titulo":"gentil fundamentos Definição 22 — p 105) concluimos que existem r, s ∈N tais que m = n + r e p"}
\section{gentil fundamentos Definição 22 — p 105) concluimos que existem r, s \ensuremath{\in}N tais que m = n + r e p}
p 105) concluimos que existem r, s \ensuremath{\in}N tais que m = n + r e p = s + q logo m + p = (n + r) + (s + q) \ensuremath{\Rightarrow} \textperiodcentered  \textperiodcentered  \textperiodcentered  \ensuremath{\Rightarrow} m + p = (n + q) + (r + s) Sendo assim, m + p \ensuremath{\le}n + q. Por oportuno, vamos provar o seguinte resultado: (iii)’ m < n e p < q \ensuremath{\Rightarrow} m + p < n + q; Da
\prov{relações: parte\_de matematica}
\prov{proveniência: gentil/Gentil\_Lopes\_FUNDAMENTOS\_DOS\_NUMEROS\_pdf.pdf \textperiodcentered{} fato \textperiodcentered{} confiança 1.0 \textperiodcentered{} domínio matematica \textperiodcentered{} história 4 versões no jornal}

%CRISTAL {"arestas":[{"alvo":"matematica","confianca":0.95,"epistemico":"fato","origem":"paper","rel":"parte_de"}],"confianca":1.0,"contraexemplos":[],"descricao":"p 105) Teorema 12 (p. 109) ´Item (iii)’ (p. 110) Proposição 8 (p. 104) Lei do Corte ( · ) (Teo. 13, p. 109, (v)) ´Item (iv)’ (p. 111) 2 x + 1 = 7 114\nPrinc´ıpio da Boa Ordem Vamos estabelecer agora um dos resultados mais importantes envolvendo a relação de ordem −o que vai nos permitir a caracterização final dos naturais −i.e., fechar a construção dos números naturais. Antes precisamos de uma","epistemico":"fato","exemplos":[],"id":"gentil_fundamentos_definicao_22_p_105_teorema_12_p_109_item_iii_p_110_proposicao","memoria":"semantica","meta":{"arquivo":"gentil/Lopes_Lopes_FUNDAMENTOS_DOS_NUMEROS_pdf.pdf","autor":"Gentil Lopes da Silva","ciencia":"teoria_dos_numeros","dominio":"matematica","nivel":4,"numero":"22","tipo_no":"paper","tipo_teorema":"definição","titulo_legivel":"Fundamentos dos Números"},"origem":"gentil/Gentil_Lopes_FUNDAMENTOS_DOS_NUMEROS_pdf.pdf","palavras_chave":[],"sinonimos":[],"tipo":"conceito","titulo":"gentil fundamentos Definição 22 — p 105) Teorema 12 (p. 109) ´Item (iii)’ (p. 110) Proposição"}
\section{gentil fundamentos Definição 22 — p 105) Teorema 12 (p. 109) 'Item (iii)’ (p. 110) Proposição}
p 105) Teorema 12 (p. 109) 'Item (iii)’ (p. 110) Proposição 8 (p. 104) Lei do Corte ( \textperiodcentered  ) (Teo. 13, p. 109, (v)) 'Item (iv)’ (p. 111) 2 x + 1 = 7 114
Princ'\ensuremath{\imath}pio da Boa Ordem Vamos estabelecer agora um dos resultados mais importantes envolvendo a relação de ordem \ensuremath{-}o que vai nos permitir a caracterização final dos naturais \ensuremath{-}i.e., fechar a construção dos números naturais. Antes precisamos de uma
\prov{relações: parte\_de matematica}
\prov{proveniência: gentil/Gentil\_Lopes\_FUNDAMENTOS\_DOS\_NUMEROS\_pdf.pdf \textperiodcentered{} fato \textperiodcentered{} confiança 1.0 \textperiodcentered{} domínio matematica \textperiodcentered{} história 4 versões no jornal}

%CRISTAL {"arestas":[{"alvo":"matematica","confianca":0.95,"epistemico":"fato","origem":"paper","rel":"parte_de"}],"confianca":1.0,"contraexemplos":[],"descricao":"p 105). Uma relação de ordem parcial que satisfaz a lei da tricotomia é dita uma relação de ordem total. (def. 17, p. 78) 108","epistemico":"fato","exemplos":[],"id":"gentil_fundamentos_definicao_22_p_105_uma_relacao_de_ordem_parcial_que_satisfaz_","memoria":"semantica","meta":{"arquivo":"gentil/Lopes_Lopes_FUNDAMENTOS_DOS_NUMEROS_pdf.pdf","autor":"Gentil Lopes da Silva","ciencia":"teoria_dos_numeros","dominio":"matematica","nivel":4,"numero":"22","tipo_no":"paper","tipo_teorema":"definição","titulo_legivel":"Fundamentos dos Números"},"origem":"gentil/Gentil_Lopes_FUNDAMENTOS_DOS_NUMEROS_pdf.pdf","palavras_chave":[],"sinonimos":[],"tipo":"conceito","titulo":"gentil fundamentos Definição 22 — p 105). Uma relação de ordem parcial que satisfaz a lei da t"}
\section{gentil fundamentos Definição 22 — p 105). Uma relação de ordem parcial que satisfaz a lei da t}
p 105). Uma relação de ordem parcial que satisfaz a lei da tricotomia é dita uma relação de ordem total. (def. 17, p. 78) 108
\prov{relações: parte\_de matematica}
\prov{proveniência: gentil/Gentil\_Lopes\_FUNDAMENTOS\_DOS\_NUMEROS\_pdf.pdf \textperiodcentered{} fato \textperiodcentered{} confiança 1.0 \textperiodcentered{} domínio matematica \textperiodcentered{} história 4 versões no jornal}

%CRISTAL {"arestas":[{"alvo":"matematica","confianca":0.95,"epistemico":"fato","origem":"paper","rel":"parte_de"}],"confianca":1.0,"contraexemplos":[],"descricao":"Relação de ordem em N) Dados x, y ∈N dizemos que x R y se, e somente se, existir p ∈N tal que y = x + p. Destacando a sentença aberta de R = \u0000N, N, p(x, y) \u0001 , temos: p(x, y): y = x + p, para algum p ∈N. Por exemplo, p(1, 3): 3 = 1 + 2, para 2 ∈N. e p(1, 3) é verdadeira, isto é, 1 R 3. Também, p(0, 1): 1 = 0 + 1, para 1 ∈N. e p(0, 1) é verdadeira, isto é, 0 R 1. Vamos provar que R é uma relação de ordem parcial em N. Acompanhe pela","epistemico":"fato","exemplos":[],"id":"gentil_fundamentos_definicao_22_relacao_de_ordem_em_n_dados_x_y_n_dizemos_que_x_","memoria":"semantica","meta":{"arquivo":"gentil/Lopes_Lopes_FUNDAMENTOS_DOS_NUMEROS_pdf.pdf","autor":"Gentil Lopes da Silva","ciencia":"teoria_dos_numeros","dominio":"matematica","nivel":4,"numero":"22","tipo_no":"paper","tipo_teorema":"definição","titulo_legivel":"Fundamentos dos Números"},"origem":"gentil/Gentil_Lopes_FUNDAMENTOS_DOS_NUMEROS_pdf.pdf","palavras_chave":[],"sinonimos":[],"tipo":"conceito","titulo":"gentil fundamentos Definição 22 — Relação de ordem em N) Dados x, y ∈N dizemos que x R y se, e"}
\section{gentil fundamentos Definição 22 — Relação de ordem em N) Dados x, y \ensuremath{\in}N dizemos que x R y se, e}
Relação de ordem em N) Dados x, y \ensuremath{\in}N dizemos que x R y se, e somente se, existir p \ensuremath{\in}N tal que y = x + p. Destacando a sentença aberta de R = N, N, p(x, y)  , temos: p(x, y): y = x + p, para algum p \ensuremath{\in}N. Por exemplo, p(1, 3): 3 = 1 + 2, para 2 \ensuremath{\in}N. e p(1, 3) é verdadeira, isto é, 1 R 3. Também, p(0, 1): 1 = 0 + 1, para 1 \ensuremath{\in}N. e p(0, 1) é verdadeira, isto é, 0 R 1. Vamos provar que R é uma relação de ordem parcial em N. Acompanhe pela
\prov{relações: parte\_de matematica}
\prov{proveniência: gentil/Gentil\_Lopes\_FUNDAMENTOS\_DOS\_NUMEROS\_pdf.pdf \textperiodcentered{} fato \textperiodcentered{} confiança 1.0 \textperiodcentered{} domínio matematica \textperiodcentered{} história 4 versões no jornal}

%CRISTAL {"arestas":[{"alvo":"matematica","confianca":0.95,"epistemico":"fato","origem":"paper","rel":"parte_de"}],"confianca":1.0,"contraexemplos":[],"descricao":"p 106) concluimos que existem r, s ∈N∗tais que m = n + r e p = s + q logo m + p = (n + r) + (s + q) ⇒ · · · ⇒ m + p = (n + q) + (r + s) Sendo assim, m + p < n + q. (iv) m ≤n ⇒ m · p ≤n · p. Suponhamos m, n arbitrariamente fixados e consideremos a proposição P(p): m ≤n ⇒ m · p ≤n · p, ∀p ∈N Esta proposição para p = 0 fica P(0): m ≤n ⇒ m · 0 ≤n · 0, ∀p ∈N que é verdadeira, pela","epistemico":"fato","exemplos":[],"id":"gentil_fundamentos_definicao_23_p_106_concluimos_que_existem_r_s_ntais_que_m_n_r","memoria":"semantica","meta":{"arquivo":"gentil/Lopes_Lopes_FUNDAMENTOS_DOS_NUMEROS_pdf.pdf","autor":"Gentil Lopes da Silva","ciencia":"teoria_dos_numeros","dominio":"matematica","nivel":4,"numero":"23","tipo_no":"paper","tipo_teorema":"definição","titulo_legivel":"Fundamentos dos Números"},"origem":"gentil/Gentil_Lopes_FUNDAMENTOS_DOS_NUMEROS_pdf.pdf","palavras_chave":[],"sinonimos":[],"tipo":"conceito","titulo":"gentil fundamentos Definição 23 — p 106) concluimos que existem r, s ∈N∗tais que m = n + r e p"}
\section{gentil fundamentos Definição 23 — p 106) concluimos que existem r, s \ensuremath{\in}N\ensuremath{*}tais que m = n + r e p}
p 106) concluimos que existem r, s \ensuremath{\in}N\ensuremath{*}tais que m = n + r e p = s + q logo m + p = (n + r) + (s + q) \ensuremath{\Rightarrow} \textperiodcentered  \textperiodcentered  \textperiodcentered  \ensuremath{\Rightarrow} m + p = (n + q) + (r + s) Sendo assim, m + p < n + q. (iv) m \ensuremath{\le}n \ensuremath{\Rightarrow} m \textperiodcentered  p \ensuremath{\le}n \textperiodcentered  p. Suponhamos m, n arbitrariamente fixados e consideremos a proposição P(p): m \ensuremath{\le}n \ensuremath{\Rightarrow} m \textperiodcentered  p \ensuremath{\le}n \textperiodcentered  p, \ensuremath{\forall}p \ensuremath{\in}N Esta proposição para p = 0 fica P(0): m \ensuremath{\le}n \ensuremath{\Rightarrow} m \textperiodcentered  0 \ensuremath{\le}n \textperiodcentered  0, \ensuremath{\forall}p \ensuremath{\in}N que é verdadeira, pela
\prov{relações: parte\_de matematica}
\prov{proveniência: gentil/Gentil\_Lopes\_FUNDAMENTOS\_DOS\_NUMEROS\_pdf.pdf \textperiodcentered{} fato \textperiodcentered{} confiança 1.0 \textperiodcentered{} domínio matematica \textperiodcentered{} história 4 versões no jornal}

%CRISTAL {"arestas":[{"alvo":"matematica","confianca":0.95,"epistemico":"fato","origem":"paper","rel":"parte_de"}],"confianca":1.0,"contraexemplos":[],"descricao":"Menor elemento) Seja A um subconjunto de N. Diremos que m ∈N é o menor elemento de A, se: (i) m ∈A; (ii) m ≤n, para todo n em A. Este menor elemento é também chamado de elemento m´ınimo de A e, por vezes, é denotado por min A.","epistemico":"fato","exemplos":[],"id":"gentil_fundamentos_definicao_24_menor_elemento_seja_a_um_subconjunto_de_n_diremo","memoria":"semantica","meta":{"arquivo":"gentil/Lopes_Lopes_FUNDAMENTOS_DOS_NUMEROS_pdf.pdf","autor":"Gentil Lopes da Silva","ciencia":"teoria_dos_numeros","dominio":"matematica","nivel":4,"numero":"24","tipo_no":"paper","tipo_teorema":"definição","titulo_legivel":"Fundamentos dos Números"},"origem":"gentil/Gentil_Lopes_FUNDAMENTOS_DOS_NUMEROS_pdf.pdf","palavras_chave":[],"sinonimos":[],"tipo":"conceito","titulo":"gentil fundamentos Definição 24 — Menor elemento) Seja A um subconjunto de N. Diremos que m ∈N"}
\section{gentil fundamentos Definição 24 — Menor elemento) Seja A um subconjunto de N. Diremos que m \ensuremath{\in}N}
Menor elemento) Seja A um subconjunto de N. Diremos que m \ensuremath{\in}N é o menor elemento de A, se: (i) m \ensuremath{\in}A; (ii) m \ensuremath{\le}n, para todo n em A. Este menor elemento é também chamado de elemento m'\ensuremath{\imath}nimo de A e, por vezes, é denotado por min A.
\prov{relações: parte\_de matematica}
\prov{proveniência: gentil/Gentil\_Lopes\_FUNDAMENTOS\_DOS\_NUMEROS\_pdf.pdf \textperiodcentered{} fato \textperiodcentered{} confiança 1.0 \textperiodcentered{} domínio matematica \textperiodcentered{} história 4 versões no jornal}

%CRISTAL {"arestas":[{"alvo":"matematica","confianca":0.95,"epistemico":"fato","origem":"paper","rel":"parte_de"},{"alvo":"virtualizacao","confianca":0.5,"epistemico":"fato","origem":"wiki","rel":"relaciona"},{"alvo":"word","confianca":0.5,"epistemico":"fato","origem":"wiki","rel":"relaciona"}],"confianca":1.0,"contraexemplos":[],"descricao":"p 115). Falta apenas provar o ´ıtem (i) da referida","epistemico":"fato","exemplos":[],"id":"gentil_fundamentos_definicao_24_p_115_falta_apenas_provar_o_ıtem_i_da_referida","memoria":"semantica","meta":{"arquivo":"gentil/Lopes_Lopes_FUNDAMENTOS_DOS_NUMEROS_pdf.pdf","autor":"Gentil Lopes da Silva","ciencia":"teoria_dos_numeros","dominio":"matematica","nivel":4,"numero":"24","tipo_no":"paper","tipo_teorema":"definição","titulo_legivel":"Fundamentos dos Números"},"origem":"gentil/Gentil_Lopes_FUNDAMENTOS_DOS_NUMEROS_pdf.pdf","palavras_chave":[],"sinonimos":[],"tipo":"conceito","titulo":"gentil fundamentos Definição 24 — p 115). Falta apenas provar o ´ıtem (i) da referida"}
\section{gentil fundamentos Definição 24 — p 115). Falta apenas provar o '\ensuremath{\imath}tem (i) da referida}
p 115). Falta apenas provar o '\ensuremath{\imath}tem (i) da referida
\prov{relações: parte\_de matematica; relaciona virtualizacao; relaciona word}
\prov{proveniência: gentil/Gentil\_Lopes\_FUNDAMENTOS\_DOS\_NUMEROS\_pdf.pdf \textperiodcentered{} fato \textperiodcentered{} confiança 1.0 \textperiodcentered{} domínio matematica \textperiodcentered{} história 6 versões no jornal}

%CRISTAL {"arestas":[{"alvo":"matematica","confianca":0.95,"epistemico":"fato","origem":"paper","rel":"parte_de"}],"confianca":1.0,"contraexemplos":[],"descricao":"Adição) Seja n ∈N um natural dado. Então ( i ) m + 0 = m; ( ii ) m + σ(n) = σ(m + n).","epistemico":"fato","exemplos":[],"id":"gentil_fundamentos_definicao_26_adicao_seja_n_n_um_natural_dado_entao_i_m_0_m","memoria":"semantica","meta":{"arquivo":"gentil/Lopes_Lopes_FUNDAMENTOS_DOS_NUMEROS_pdf.pdf","autor":"Gentil Lopes da Silva","ciencia":"teoria_dos_numeros","dominio":"matematica","nivel":4,"numero":"26","tipo_no":"paper","tipo_teorema":"definição","titulo_legivel":"Fundamentos dos Números"},"origem":"gentil/Gentil_Lopes_FUNDAMENTOS_DOS_NUMEROS_pdf.pdf","palavras_chave":[],"sinonimos":[],"tipo":"conceito","titulo":"gentil fundamentos Definição 26 — Adição) Seja n ∈N um natural dado. Então ( i ) m + 0 = m; ("}
\section{gentil fundamentos Definição 26 — Adição) Seja n \ensuremath{\in}N um natural dado. Então ( i ) m + 0 = m; (}
Adição) Seja n \ensuremath{\in}N um natural dado. Então ( i ) m + 0 = m; ( ii ) m + \ensuremath{\sigma}(n) = \ensuremath{\sigma}(m + n).
\prov{relações: parte\_de matematica}
\prov{proveniência: gentil/Gentil\_Lopes\_FUNDAMENTOS\_DOS\_NUMEROS\_pdf.pdf \textperiodcentered{} fato \textperiodcentered{} confiança 1.0 \textperiodcentered{} domínio matematica \textperiodcentered{} história 4 versões no jornal}

%CRISTAL {"arestas":[{"alvo":"matematica","confianca":0.95,"epistemico":"fato","origem":"paper","rel":"parte_de"}],"confianca":1.0,"contraexemplos":[],"descricao":"Multiplicação) A multiplicação em N é definida por: ( n · 0 = 0; ∀n ∈N n · (m + 1) = n · m + n, sempre que n · m está definido.","epistemico":"fato","exemplos":[],"id":"gentil_fundamentos_definicao_27_multiplicacao_a_multiplicacao_em_n_e_definida_po","memoria":"semantica","meta":{"arquivo":"gentil/Lopes_Lopes_FUNDAMENTOS_DOS_NUMEROS_pdf.pdf","autor":"Gentil Lopes da Silva","ciencia":"teoria_dos_numeros","dominio":"matematica","nivel":4,"numero":"27","tipo_no":"paper","tipo_teorema":"definição","titulo_legivel":"Fundamentos dos Números"},"origem":"gentil/Gentil_Lopes_FUNDAMENTOS_DOS_NUMEROS_pdf.pdf","palavras_chave":[],"sinonimos":[],"tipo":"conceito","titulo":"gentil fundamentos Definição 27 — Multiplicação) A multiplicação em N é definida por: ( n · 0"}
\section{gentil fundamentos Definição 27 — Multiplicação) A multiplicação em N é definida por: ( n \textperiodcentered  0}
Multiplicação) A multiplicação em N é definida por: ( n \textperiodcentered  0 = 0; \ensuremath{\forall}n \ensuremath{\in}N n \textperiodcentered  (m + 1) = n \textperiodcentered  m + n, sempre que n \textperiodcentered  m está definido.
\prov{relações: parte\_de matematica}
\prov{proveniência: gentil/Gentil\_Lopes\_FUNDAMENTOS\_DOS\_NUMEROS\_pdf.pdf \textperiodcentered{} fato \textperiodcentered{} confiança 1.0 \textperiodcentered{} domínio matematica \textperiodcentered{} história 4 versões no jornal}

%CRISTAL {"arestas":[{"alvo":"matematica","confianca":0.95,"epistemico":"fato","origem":"paper","rel":"parte_de"}],"confianca":1.0,"contraexemplos":[],"descricao":"Dados dois elementos (a, b) e (c, d) do conjunto N × N, diremos que (a, b) ∼(c, d), se e somente se, a + d = b + c. Por exemplo: (1, 2) ∼(3, 4) ⇐⇒ 1 + 4 = 2 + 3 Vamos provar que a relação definida acima é de equivalência. Prova: Acompanhe pela","epistemico":"fato","exemplos":[],"id":"gentil_fundamentos_definicao_29_dados_dois_elementos_a_b_e_c_d_do_conjunto_n_n_d","memoria":"semantica","meta":{"arquivo":"gentil/Lopes_Lopes_FUNDAMENTOS_DOS_NUMEROS_pdf.pdf","autor":"Gentil Lopes da Silva","ciencia":"teoria_dos_numeros","dominio":"matematica","nivel":4,"numero":"29","tipo_no":"paper","tipo_teorema":"definição","titulo_legivel":"Fundamentos dos Números"},"origem":"gentil/Gentil_Lopes_FUNDAMENTOS_DOS_NUMEROS_pdf.pdf","palavras_chave":[],"sinonimos":[],"tipo":"conceito","titulo":"gentil fundamentos Definição 29 — Dados dois elementos (a, b) e (c, d) do conjunto N × N, dire"}
\section{gentil fundamentos Definição 29 — Dados dois elementos (a, b) e (c, d) do conjunto N $\times$ N, dire}
Dados dois elementos (a, b) e (c, d) do conjunto N $\times$ N, diremos que (a, b) \ensuremath{\sim}(c, d), se e somente se, a + d = b + c. Por exemplo: (1, 2) \ensuremath{\sim}(3, 4) \ensuremath{\Leftarrow}\ensuremath{\Rightarrow} 1 + 4 = 2 + 3 Vamos provar que a relação definida acima é de equivalência. Prova: Acompanhe pela
\prov{relações: parte\_de matematica}
\prov{proveniência: gentil/Gentil\_Lopes\_FUNDAMENTOS\_DOS\_NUMEROS\_pdf.pdf \textperiodcentered{} fato \textperiodcentered{} confiança 1.0 \textperiodcentered{} domínio matematica \textperiodcentered{} história 4 versões no jornal}

%CRISTAL {"arestas":[{"alvo":"gentil_fundamentos_capitulo_1_o_que_e_um_n_umero","confianca":1.0,"epistemico":"fato","origem":"gentil/Gentil_Lopes_FUNDAMENTOS_DOS_NUMEROS_pdf.pdf","rel":"define"}],"confianca":1.0,"contraexemplos":[],"descricao":"p 156) (a + b, a + b) ∼(n, n) ⇐⇒ (a + b) + n = (a + b) + n Agora é só utilizar o","epistemico":"fato","exemplos":[],"id":"gentil_fundamentos_definicao_29_p_156_a_b_a_b_n_n_a_b_n_a_b_n_a","memoria":"semantica","meta":{"arquivo":"gentil/Lopes_Lopes_FUNDAMENTOS_DOS_NUMEROS_pdf.pdf","autor":"Gentil Lopes da Silva","ciencia":"teoria_dos_numeros","dominio":"matematica","nivel":4,"numero":"29","tipo_no":"paper","tipo_teorema":"definição","titulo_legivel":"Fundamentos dos Números"},"origem":"gentil/Gentil_Lopes_FUNDAMENTOS_DOS_NUMEROS_pdf.pdf","palavras_chave":[],"sinonimos":[],"tipo":"conceito","titulo":"gentil fundamentos Definição 29 — p 156) (a + b, a + b) ∼(n, n) ⇐⇒ (a + b) + n = (a + b) + n A"}
\section{gentil fundamentos Definição 29 — p 156) (a + b, a + b) \ensuremath{\sim}(n, n) \ensuremath{\Leftarrow}\ensuremath{\Rightarrow} (a + b) + n = (a + b) + n A}
p 156) (a + b, a + b) \ensuremath{\sim}(n, n) \ensuremath{\Leftarrow}\ensuremath{\Rightarrow} (a + b) + n = (a + b) + n Agora é só utilizar o
\prov{relações: define gentil\_fundamentos\_capitulo\_1\_o\_que\_e\_um\_n\_umero}
\prov{proveniência: gentil/Gentil\_Lopes\_FUNDAMENTOS\_DOS\_NUMEROS\_pdf.pdf \textperiodcentered{} fato \textperiodcentered{} confiança 1.0 \textperiodcentered{} domínio matematica \textperiodcentered{} história 5 versões no jornal}

%CRISTAL {"arestas":[{"alvo":"gentil_fundamentos_capitulo_1_o_que_e_um_n_umero","confianca":1.0,"epistemico":"fato","origem":"gentil/Gentil_Lopes_FUNDAMENTOS_DOS_NUMEROS_pdf.pdf","rel":"define"}],"confianca":1.0,"contraexemplos":[],"descricao":"p 156) (a + n, b + n) ∼(a, b) ⇐⇒ (a + n) + b = (b + n) + a Esta última igualdade é verdadeira quando trata-se de números naturais. Finalmente pelo","epistemico":"fato","exemplos":[],"id":"gentil_fundamentos_definicao_29_p_156_a_n_b_n_a_b_a_n_b_b_n_a_e","memoria":"semantica","meta":{"arquivo":"gentil/Lopes_Lopes_FUNDAMENTOS_DOS_NUMEROS_pdf.pdf","autor":"Gentil Lopes da Silva","ciencia":"teoria_dos_numeros","dominio":"matematica","nivel":4,"numero":"29","tipo_no":"paper","tipo_teorema":"definição","titulo_legivel":"Fundamentos dos Números"},"origem":"gentil/Gentil_Lopes_FUNDAMENTOS_DOS_NUMEROS_pdf.pdf","palavras_chave":[],"sinonimos":[],"tipo":"conceito","titulo":"gentil fundamentos Definição 29 — p 156) (a + n, b + n) ∼(a, b) ⇐⇒ (a + n) + b = (b + n) + a E"}
\section{gentil fundamentos Definição 29 — p 156) (a + n, b + n) \ensuremath{\sim}(a, b) \ensuremath{\Leftarrow}\ensuremath{\Rightarrow} (a + n) + b = (b + n) + a E}
p 156) (a + n, b + n) \ensuremath{\sim}(a, b) \ensuremath{\Leftarrow}\ensuremath{\Rightarrow} (a + n) + b = (b + n) + a Esta última igualdade é verdadeira quando trata-se de números naturais. Finalmente pelo
\prov{relações: define gentil\_fundamentos\_capitulo\_1\_o\_que\_e\_um\_n\_umero}
\prov{proveniência: gentil/Gentil\_Lopes\_FUNDAMENTOS\_DOS\_NUMEROS\_pdf.pdf \textperiodcentered{} fato \textperiodcentered{} confiança 1.0 \textperiodcentered{} domínio matematica \textperiodcentered{} história 5 versões no jornal}

%CRISTAL {"arestas":[{"alvo":"gentil_fundamentos_capitulo_1_o_que_e_um_n_umero","confianca":1.0,"epistemico":"fato","origem":"gentil/Gentil_Lopes_FUNDAMENTOS_DOS_NUMEROS_pdf.pdf","rel":"define"}],"confianca":1.0,"contraexemplos":[],"descricao":"p 156) \u0000a(n+1)+bn, an+b(n+1) \u0001 ∼(a, b) ⇐⇒ a(n+1)+bn+b = an+b(n+1)+a Esta última igualdade é verdadeira quando trata-se de números naturais. Finalmente pelo","epistemico":"fato","exemplos":[],"id":"gentil_fundamentos_definicao_29_p_156_an1bn_anbn1_a_b_an1bnb_anb","memoria":"semantica","meta":{"arquivo":"gentil/Lopes_Lopes_FUNDAMENTOS_DOS_NUMEROS_pdf.pdf","autor":"Gentil Lopes da Silva","ciencia":"teoria_dos_numeros","dominio":"matematica","nivel":4,"numero":"29","tipo_no":"paper","tipo_teorema":"definição","titulo_legivel":"Fundamentos dos Números"},"origem":"gentil/Gentil_Lopes_FUNDAMENTOS_DOS_NUMEROS_pdf.pdf","palavras_chave":[],"sinonimos":[],"tipo":"conceito","titulo":"gentil fundamentos Definição 29 — p 156) \u0000a(n+1)+bn, an+b(n+1) \u0001 ∼(a, b) ⇐⇒ a(n+1)+bn+b = an+b"}
\section{gentil fundamentos Definição 29 — p 156) a(n+1)+bn, an+b(n+1)  \ensuremath{\sim}(a, b) \ensuremath{\Leftarrow}\ensuremath{\Rightarrow} a(n+1)+bn+b = an+b}
p 156) a(n+1)+bn, an+b(n+1)  \ensuremath{\sim}(a, b) \ensuremath{\Leftarrow}\ensuremath{\Rightarrow} a(n+1)+bn+b = an+b(n+1)+a Esta última igualdade é verdadeira quando trata-se de números naturais. Finalmente pelo
\prov{relações: define gentil\_fundamentos\_capitulo\_1\_o\_que\_e\_um\_n\_umero}
\prov{proveniência: gentil/Gentil\_Lopes\_FUNDAMENTOS\_DOS\_NUMEROS\_pdf.pdf \textperiodcentered{} fato \textperiodcentered{} confiança 1.0 \textperiodcentered{} domínio matematica \textperiodcentered{} história 5 versões no jornal}

%CRISTAL {"arestas":[{"alvo":"gentil_fundamentos_capitulo_1_o_que_e_um_n_umero","confianca":1.0,"epistemico":"fato","origem":"gentil/Gentil_Lopes_FUNDAMENTOS_DOS_NUMEROS_pdf.pdf","rel":"define"},{"alvo":"word","confianca":0.5,"epistemico":"fato","origem":"wiki","rel":"relaciona"},{"alvo":"syscall","confianca":0.5,"epistemico":"fato","origem":"wiki","rel":"relaciona"}],"confianca":1.0,"contraexemplos":[],"descricao":"p 156) Definição 10 (p. 63) Teorema 2 (p. 69)","epistemico":"fato","exemplos":[],"id":"gentil_fundamentos_definicao_29_p_156_definicao_10_p_63_teorema_2_p_69","memoria":"semantica","meta":{"arquivo":"gentil/Lopes_Lopes_FUNDAMENTOS_DOS_NUMEROS_pdf.pdf","autor":"Gentil Lopes da Silva","ciencia":"teoria_dos_numeros","dominio":"matematica","nivel":4,"numero":"29","tipo_no":"paper","tipo_teorema":"definição","titulo_legivel":"Fundamentos dos Números"},"origem":"gentil/Gentil_Lopes_FUNDAMENTOS_DOS_NUMEROS_pdf.pdf","palavras_chave":[],"sinonimos":[],"tipo":"conceito","titulo":"gentil fundamentos Definição 29 — p 156) Definição 10 (p. 63) Teorema 2 (p. 69)"}
\section{gentil fundamentos Definição 29 — p 156) Definição 10 (p. 63) Teorema 2 (p. 69)}
p 156) Definição 10 (p. 63) Teorema 2 (p. 69)
\prov{relações: define gentil\_fundamentos\_capitulo\_1\_o\_que\_e\_um\_n\_umero; relaciona word; relaciona syscall}
\prov{proveniência: gentil/Gentil\_Lopes\_FUNDAMENTOS\_DOS\_NUMEROS\_pdf.pdf \textperiodcentered{} fato \textperiodcentered{} confiança 1.0 \textperiodcentered{} domínio matematica \textperiodcentered{} história 7 versões no jornal}

%CRISTAL {"arestas":[{"alvo":"gentil_fundamentos_capitulo_1_o_que_e_um_n_umero","confianca":1.0,"epistemico":"fato","origem":"gentil/Gentil_Lopes_FUNDAMENTOS_DOS_NUMEROS_pdf.pdf","rel":"define"}],"confianca":1.0,"contraexemplos":[],"descricao":"p 156), temos (a, b) ∼(a −b, 0) ⇐⇒ a + 0 = b + (a −b) Logo, pelo","epistemico":"fato","exemplos":[],"id":"gentil_fundamentos_definicao_29_p_156_temos_a_b_a_b_0_a_0_b_a_b_logo","memoria":"semantica","meta":{"arquivo":"gentil/Lopes_Lopes_FUNDAMENTOS_DOS_NUMEROS_pdf.pdf","autor":"Gentil Lopes da Silva","ciencia":"teoria_dos_numeros","dominio":"matematica","nivel":4,"numero":"29","tipo_no":"paper","tipo_teorema":"definição","titulo_legivel":"Fundamentos dos Números"},"origem":"gentil/Gentil_Lopes_FUNDAMENTOS_DOS_NUMEROS_pdf.pdf","palavras_chave":[],"sinonimos":[],"tipo":"conceito","titulo":"gentil fundamentos Definição 29 — p 156), temos (a, b) ∼(a −b, 0) ⇐⇒ a + 0 = b + (a −b) Logo,"}
\section{gentil fundamentos Definição 29 — p 156), temos (a, b) \ensuremath{\sim}(a \ensuremath{-}b, 0) \ensuremath{\Leftarrow}\ensuremath{\Rightarrow} a + 0 = b + (a \ensuremath{-}b) Logo,}
p 156), temos (a, b) \ensuremath{\sim}(a \ensuremath{-}b, 0) \ensuremath{\Leftarrow}\ensuremath{\Rightarrow} a + 0 = b + (a \ensuremath{-}b) Logo, pelo
\prov{relações: define gentil\_fundamentos\_capitulo\_1\_o\_que\_e\_um\_n\_umero}
\prov{proveniência: gentil/Gentil\_Lopes\_FUNDAMENTOS\_DOS\_NUMEROS\_pdf.pdf \textperiodcentered{} fato \textperiodcentered{} confiança 1.0 \textperiodcentered{} domínio matematica \textperiodcentered{} história 5 versões no jornal}

%CRISTAL {"arestas":[{"alvo":"gentil_fundamentos_capitulo_1_o_que_e_um_n_umero","confianca":1.0,"epistemico":"fato","origem":"gentil/Gentil_Lopes_FUNDAMENTOS_DOS_NUMEROS_pdf.pdf","rel":"define"}],"confianca":1.0,"contraexemplos":[],"descricao":"p 156) temos (a + b, b) ∼(a, 0) ⇐⇒ a + b + 0 = b + a Portanto, pelo","epistemico":"fato","exemplos":[],"id":"gentil_fundamentos_definicao_29_p_156_temos_a_b_b_a_0_a_b_0_b_a_portant","memoria":"semantica","meta":{"arquivo":"gentil/Lopes_Lopes_FUNDAMENTOS_DOS_NUMEROS_pdf.pdf","autor":"Gentil Lopes da Silva","ciencia":"teoria_dos_numeros","dominio":"matematica","nivel":4,"numero":"29","tipo_no":"paper","tipo_teorema":"definição","titulo_legivel":"Fundamentos dos Números"},"origem":"gentil/Gentil_Lopes_FUNDAMENTOS_DOS_NUMEROS_pdf.pdf","palavras_chave":[],"sinonimos":[],"tipo":"conceito","titulo":"gentil fundamentos Definição 29 — p 156) temos (a + b, b) ∼(a, 0) ⇐⇒ a + b + 0 = b + a Portant"}
\section{gentil fundamentos Definição 29 — p 156) temos (a + b, b) \ensuremath{\sim}(a, 0) \ensuremath{\Leftarrow}\ensuremath{\Rightarrow} a + b + 0 = b + a Portant}
p 156) temos (a + b, b) \ensuremath{\sim}(a, 0) \ensuremath{\Leftarrow}\ensuremath{\Rightarrow} a + b + 0 = b + a Portanto, pelo
\prov{relações: define gentil\_fundamentos\_capitulo\_1\_o\_que\_e\_um\_n\_umero}
\prov{proveniência: gentil/Gentil\_Lopes\_FUNDAMENTOS\_DOS\_NUMEROS\_pdf.pdf \textperiodcentered{} fato \textperiodcentered{} confiança 1.0 \textperiodcentered{} domínio matematica \textperiodcentered{} história 5 versões no jornal}

%CRISTAL {"arestas":[{"alvo":"matematica","confianca":0.95,"epistemico":"fato","origem":"paper","rel":"parte_de"}],"confianca":1.0,"contraexemplos":[],"descricao":"Conjunto numérico) Dado um conjunto E não vazio e duas operações sobre E, + : E × E →E (x, y) x + y 7→ · : E × E →E (x, y) x · y 7→ A terna (E, +, ·) é o que entendemos por um conjunto numérico (ou estrutura numérica). Usaremos da seguinte notação (E, +, ·) = E. Observe que um “conjunto numérico” é mais que um mero conjunto, é uma estrutura (ou sistema).","epistemico":"fato","exemplos":[],"id":"gentil_fundamentos_definicao_2_conjunto_numerico_dado_um_conjunto_e_nao_vazio_e_","memoria":"semantica","meta":{"arquivo":"gentil/Lopes_Lopes_FUNDAMENTOS_DOS_NUMEROS_pdf.pdf","autor":"Gentil Lopes da Silva","ciencia":"teoria_dos_numeros","dominio":"matematica","nivel":4,"numero":"2","tipo_no":"paper","tipo_teorema":"definição","titulo_legivel":"Fundamentos dos Números"},"origem":"gentil/Gentil_Lopes_FUNDAMENTOS_DOS_NUMEROS_pdf.pdf","palavras_chave":[],"sinonimos":[],"tipo":"conceito","titulo":"gentil fundamentos Definição 2 — Conjunto numérico) Dado um conjunto E não vazio e duas opera"}
\section{gentil fundamentos Definição 2 — Conjunto numérico) Dado um conjunto E não vazio e duas opera}
Conjunto numérico) Dado um conjunto E não vazio e duas operações sobre E, + : E $\times$ E \ensuremath{\to}E (x, y) x + y 7\ensuremath{\to} \textperiodcentered  : E $\times$ E \ensuremath{\to}E (x, y) x \textperiodcentered  y 7\ensuremath{\to} A terna (E, +, \textperiodcentered ) é o que entendemos por um conjunto numérico (ou estrutura numérica). Usaremos da seguinte notação (E, +, \textperiodcentered ) = E. Observe que um “conjunto numérico” é mais que um mero conjunto, é uma estrutura (ou sistema).
\prov{relações: parte\_de matematica}
\prov{proveniência: gentil/Gentil\_Lopes\_FUNDAMENTOS\_DOS\_NUMEROS\_pdf.pdf \textperiodcentered{} fato \textperiodcentered{} confiança 1.0 \textperiodcentered{} domínio matematica \textperiodcentered{} história 4 versões no jornal}

%CRISTAL {"arestas":[{"alvo":"matematica","confianca":0.95,"epistemico":"fato","origem":"paper","rel":"parte_de"}],"confianca":1.0,"contraexemplos":[],"descricao":"Adição em ¯Z) Sejam α = (a, b) e β = (c, d) elementos de ¯Z. Definimos a soma de α e β como α + β = (a, b) + (c, d) = (a + c, b + d)","epistemico":"fato","exemplos":[],"id":"gentil_fundamentos_definicao_30_adicao_em_z_sejam_α_a_b_e_β_c_d_elementos_de_z","memoria":"semantica","meta":{"arquivo":"gentil/Lopes_Lopes_FUNDAMENTOS_DOS_NUMEROS_pdf.pdf","autor":"Gentil Lopes da Silva","ciencia":"teoria_dos_numeros","dominio":"matematica","nivel":4,"numero":"30","tipo_no":"paper","tipo_teorema":"definição","titulo_legivel":"Fundamentos dos Números"},"origem":"gentil/Gentil_Lopes_FUNDAMENTOS_DOS_NUMEROS_pdf.pdf","palavras_chave":[],"sinonimos":[],"tipo":"conceito","titulo":"gentil fundamentos Definição 30 — Adição em ¯Z) Sejam α = (a, b) e β = (c, d) elementos de ¯Z."}
\section{gentil fundamentos Definição 30 — Adição em \ensuremath{\bar{\,}}Z) Sejam \ensuremath{\alpha} = (a, b) e \ensuremath{\beta} = (c, d) elementos de \ensuremath{\bar{\,}}Z.}
Adição em \ensuremath{\bar{\,}}Z) Sejam \ensuremath{\alpha} = (a, b) e \ensuremath{\beta} = (c, d) elementos de \ensuremath{\bar{\,}}Z. Definimos a soma de \ensuremath{\alpha} e \ensuremath{\beta} como \ensuremath{\alpha} + \ensuremath{\beta} = (a, b) + (c, d) = (a + c, b + d)
\prov{relações: parte\_de matematica}
\prov{proveniência: gentil/Gentil\_Lopes\_FUNDAMENTOS\_DOS\_NUMEROS\_pdf.pdf \textperiodcentered{} fato \textperiodcentered{} confiança 1.0 \textperiodcentered{} domínio matematica \textperiodcentered{} história 4 versões no jornal}

%CRISTAL {"arestas":[{"alvo":"matematica","confianca":0.95,"epistemico":"fato","origem":"paper","rel":"parte_de"},{"alvo":"word","confianca":0.5,"epistemico":"fato","origem":"wiki","rel":"relaciona"},{"alvo":"virtualizacao","confianca":0.5,"epistemico":"fato","origem":"wiki","rel":"relaciona"}],"confianca":1.0,"contraexemplos":[],"descricao":"+, p 159) Definição 31 (×, p. 159) Teoremas p. 162, 163, 164","epistemico":"fato","exemplos":[],"id":"gentil_fundamentos_definicao_30_p_159_definicao_31_p_159_teoremas_p_162_163_164","memoria":"semantica","meta":{"arquivo":"gentil/Lopes_Lopes_FUNDAMENTOS_DOS_NUMEROS_pdf.pdf","autor":"Gentil Lopes da Silva","ciencia":"teoria_dos_numeros","dominio":"matematica","nivel":4,"numero":"30","tipo_no":"paper","tipo_teorema":"definição","titulo_legivel":"Fundamentos dos Números"},"origem":"gentil/Gentil_Lopes_FUNDAMENTOS_DOS_NUMEROS_pdf.pdf","palavras_chave":[],"sinonimos":[],"tipo":"conceito","titulo":"gentil fundamentos Definição 30 — +, p 159) Definição 31 (×, p. 159) Teoremas p. 162, 163, 164"}
\section{gentil fundamentos Definição 30 — +, p 159) Definição 31 ($\times$, p. 159) Teoremas p. 162, 163, 164}
+, p 159) Definição 31 ($\times$, p. 159) Teoremas p. 162, 163, 164
\prov{relações: parte\_de matematica; relaciona word; relaciona virtualizacao}
\prov{proveniência: gentil/Gentil\_Lopes\_FUNDAMENTOS\_DOS\_NUMEROS\_pdf.pdf \textperiodcentered{} fato \textperiodcentered{} confiança 1.0 \textperiodcentered{} domínio matematica \textperiodcentered{} história 6 versões no jornal}

%CRISTAL {"arestas":[{"alvo":"matematica","confianca":0.95,"epistemico":"fato","origem":"paper","rel":"parte_de"}],"confianca":1.0,"contraexemplos":[],"descricao":"Multiplicação em ¯Z) Sejam α = (a, b) e β = (c, d) ele- mentos de ¯Z. Definimos o produto de α e β como α · β = (a, b) · (c, d) = (ac + bd, ad + bc) (5.6)","epistemico":"fato","exemplos":[],"id":"gentil_fundamentos_definicao_31_multiplicacao_em_z_sejam_α_a_b_e_β_c_d_ele-_ment","memoria":"semantica","meta":{"arquivo":"gentil/Lopes_Lopes_FUNDAMENTOS_DOS_NUMEROS_pdf.pdf","autor":"Gentil Lopes da Silva","ciencia":"teoria_dos_numeros","dominio":"matematica","nivel":4,"numero":"31","tipo_no":"paper","tipo_teorema":"definição","titulo_legivel":"Fundamentos dos Números"},"origem":"gentil/Gentil_Lopes_FUNDAMENTOS_DOS_NUMEROS_pdf.pdf","palavras_chave":[],"sinonimos":[],"tipo":"conceito","titulo":"gentil fundamentos Definição 31 — Multiplicação em ¯Z) Sejam α = (a, b) e β = (c, d) ele- ment"}
\section{gentil fundamentos Definição 31 — Multiplicação em \ensuremath{\bar{\,}}Z) Sejam \ensuremath{\alpha} = (a, b) e \ensuremath{\beta} = (c, d) ele- ment}
Multiplicação em \ensuremath{\bar{\,}}Z) Sejam \ensuremath{\alpha} = (a, b) e \ensuremath{\beta} = (c, d) ele- mentos de \ensuremath{\bar{\,}}Z. Definimos o produto de \ensuremath{\alpha} e \ensuremath{\beta} como \ensuremath{\alpha} \textperiodcentered  \ensuremath{\beta} = (a, b) \textperiodcentered  (c, d) = (ac + bd, ad + bc) (5.6)
\prov{relações: parte\_de matematica}
\prov{proveniência: gentil/Gentil\_Lopes\_FUNDAMENTOS\_DOS\_NUMEROS\_pdf.pdf \textperiodcentered{} fato \textperiodcentered{} confiança 1.0 \textperiodcentered{} domínio matematica \textperiodcentered{} história 4 versões no jornal}

%CRISTAL {"arestas":[{"alvo":"matematica","confianca":0.95,"epistemico":"fato","origem":"paper","rel":"parte_de"}],"confianca":1.0,"contraexemplos":[],"descricao":"Números Inteiros) Denominamos de sistema (“conjunto”) dos números inteiros à terna, ¯Z = (¯Z, +, ·) Observe esta construção no esquema a seguir, ¯Z ¯Z × ¯Z ¯Z ¯Z + · ¯Z = (¯Z, +, ·) - Conjunto (hardware) (aqui temos meros elementos) - Estrutura (aqui temos os números inteiros) E por que estes são os números inteiros? Respondemos: porque com estas operações demonstramos −vamos demonstrar −que as especificações do sistema, ν (p. 154), são atendidas. 159\nDe passagem, uma observação não desprez´ıvel","epistemico":"fato","exemplos":[],"id":"gentil_fundamentos_definicao_32_numeros_inteiros_denominamos_de_sistema_conjunto","memoria":"semantica","meta":{"arquivo":"gentil/Lopes_Lopes_FUNDAMENTOS_DOS_NUMEROS_pdf.pdf","autor":"Gentil Lopes da Silva","ciencia":"teoria_dos_numeros","dominio":"matematica","nivel":4,"numero":"32","tipo_no":"paper","tipo_teorema":"definição","titulo_legivel":"Fundamentos dos Números"},"origem":"gentil/Gentil_Lopes_FUNDAMENTOS_DOS_NUMEROS_pdf.pdf","palavras_chave":[],"sinonimos":[],"tipo":"conceito","titulo":"gentil fundamentos Definição 32 — Números Inteiros) Denominamos de sistema (“conjunto”) dos nú"}
\section{gentil fundamentos Definição 32 — Números Inteiros) Denominamos de sistema (“conjunto”) dos nú}
Números Inteiros) Denominamos de sistema (“conjunto”) dos números inteiros à terna, \ensuremath{\bar{\,}}Z = (\ensuremath{\bar{\,}}Z, +, \textperiodcentered ) Observe esta construção no esquema a seguir, \ensuremath{\bar{\,}}Z \ensuremath{\bar{\,}}Z $\times$ \ensuremath{\bar{\,}}Z \ensuremath{\bar{\,}}Z \ensuremath{\bar{\,}}Z + \textperiodcentered  \ensuremath{\bar{\,}}Z = (\ensuremath{\bar{\,}}Z, +, \textperiodcentered ) - Conjunto (hardware) (aqui temos meros elementos) - Estrutura (aqui temos os números inteiros) E por que estes são os números inteiros? Respondemos: porque com estas operações demonstramos \ensuremath{-}vamos demonstrar \ensuremath{-}que as especificações do sistema, \ensuremath{\nu} (p. 154), são atendidas. 159
De passagem, uma observação não desprez'\ensuremath{\imath}vel
\prov{relações: parte\_de matematica}
\prov{proveniência: gentil/Gentil\_Lopes\_FUNDAMENTOS\_DOS\_NUMEROS\_pdf.pdf \textperiodcentered{} fato \textperiodcentered{} confiança 1.0 \textperiodcentered{} domínio matematica \textperiodcentered{} história 4 versões no jornal}

%CRISTAL {"arestas":[{"alvo":"matematica","confianca":0.95,"epistemico":"fato","origem":"paper","rel":"parte_de"}],"confianca":1.0,"contraexemplos":[],"descricao":"Dados os elementos α = (a, b) e β = (c, d) em ¯Z, diremos que α é menor ou igual a β, e escrevemos α ≤β se, e somente se, a+d ≤b+c. Resumindo: (a, b) ≤(c, d) ⇐⇒ a + d ≤b + c Uma observação oportuna é a de que na","epistemico":"fato","exemplos":[],"id":"gentil_fundamentos_definicao_33_dados_os_elementos_α_a_b_e_β_c_d_em_z_diremos_qu","memoria":"semantica","meta":{"arquivo":"gentil/Lopes_Lopes_FUNDAMENTOS_DOS_NUMEROS_pdf.pdf","autor":"Gentil Lopes da Silva","ciencia":"teoria_dos_numeros","dominio":"matematica","nivel":4,"numero":"33","tipo_no":"paper","tipo_teorema":"definição","titulo_legivel":"Fundamentos dos Números"},"origem":"gentil/Gentil_Lopes_FUNDAMENTOS_DOS_NUMEROS_pdf.pdf","palavras_chave":[],"sinonimos":[],"tipo":"conceito","titulo":"gentil fundamentos Definição 33 — Dados os elementos α = (a, b) e β = (c, d) em ¯Z, diremos qu"}
\section{gentil fundamentos Definição 33 — Dados os elementos \ensuremath{\alpha} = (a, b) e \ensuremath{\beta} = (c, d) em \ensuremath{\bar{\,}}Z, diremos qu}
Dados os elementos \ensuremath{\alpha} = (a, b) e \ensuremath{\beta} = (c, d) em \ensuremath{\bar{\,}}Z, diremos que \ensuremath{\alpha} é menor ou igual a \ensuremath{\beta}, e escrevemos \ensuremath{\alpha} \ensuremath{\le}\ensuremath{\beta} se, e somente se, a+d \ensuremath{\le}b+c. Resumindo: (a, b) \ensuremath{\le}(c, d) \ensuremath{\Leftarrow}\ensuremath{\Rightarrow} a + d \ensuremath{\le}b + c Uma observação oportuna é a de que na
\prov{relações: parte\_de matematica}
\prov{proveniência: gentil/Gentil\_Lopes\_FUNDAMENTOS\_DOS\_NUMEROS\_pdf.pdf \textperiodcentered{} fato \textperiodcentered{} confiança 1.0 \textperiodcentered{} domínio matematica \textperiodcentered{} história 4 versões no jornal}

%CRISTAL {"arestas":[{"alvo":"matematica","confianca":0.95,"epistemico":"fato","origem":"paper","rel":"parte_de"}],"confianca":1.0,"contraexemplos":[],"descricao":"(p 171) Por exemplo, para os inteiros a seguir (1, 0), (2, 1), (3, 2), (4, 3),... ? ? ? (0, 0), (1, 1), (2, 2), (3, 3),... quem é maior do que quem? Para responder esta pergunta, temos (p. 171) (1, 0) ≤(0, 0) ⇐⇒ 1 + 0 ≤0 + 0 Como esta última desigualdade não é verdadeira, devemos ter (0, 0) < (1, 0) Segundo notações já fixadas (p. 163, p. 167) ¯0 < ¯1 foi o que acabamos de provar. Novamente observamos que esta não é uma desigualdade trivial. Nos “inteiros dos ocidentais e dos orien","epistemico":"fato","exemplos":[],"id":"gentil_fundamentos_definicao_33_p_171_por_exemplo_para_os_inteiros_a_seguir_1_0_","memoria":"semantica","meta":{"arquivo":"gentil/Lopes_Lopes_FUNDAMENTOS_DOS_NUMEROS_pdf.pdf","autor":"Gentil Lopes da Silva","ciencia":"teoria_dos_numeros","dominio":"matematica","nivel":4,"numero":"33","tipo_no":"paper","tipo_teorema":"definição","titulo_legivel":"Fundamentos dos Números"},"origem":"gentil/Gentil_Lopes_FUNDAMENTOS_DOS_NUMEROS_pdf.pdf","palavras_chave":[],"sinonimos":[],"tipo":"conceito","titulo":"gentil fundamentos Definição 33 — (p 171) Por exemplo, para os inteiros a seguir (1, 0), (2"}
\section{gentil fundamentos Definição 33 — (p 171) Por exemplo, para os inteiros a seguir (1, 0), (2}
(p 171) Por exemplo, para os inteiros a seguir (1, 0), (2, 1), (3, 2), (4, 3),... ? ? ? (0, 0), (1, 1), (2, 2), (3, 3),... quem é maior do que quem? Para responder esta pergunta, temos (p. 171) (1, 0) \ensuremath{\le}(0, 0) \ensuremath{\Leftarrow}\ensuremath{\Rightarrow} 1 + 0 \ensuremath{\le}0 + 0 Como esta última desigualdade não é verdadeira, devemos ter (0, 0) < (1, 0) Segundo notações já fixadas (p. 163, p. 167) \ensuremath{\bar{\,}}0 < \ensuremath{\bar{\,}}1 foi o que acabamos de provar. Novamente observamos que esta não é uma desigualdade trivial. Nos “inteiros dos ocidentais e dos orien
\prov{relações: parte\_de matematica}
\prov{proveniência: gentil/Gentil\_Lopes\_FUNDAMENTOS\_DOS\_NUMEROS\_pdf.pdf \textperiodcentered{} fato \textperiodcentered{} confiança 1.0 \textperiodcentered{} domínio matematica \textperiodcentered{} história 6 versões no jornal}

%CRISTAL {"arestas":[{"alvo":"matematica","confianca":0.95,"epistemico":"fato","origem":"paper","rel":"parte_de"}],"confianca":1.0,"contraexemplos":[],"descricao":"Limitado inferiormente) Seja A um subconjunto não vazio de ¯Z. Diz-se que A é limitado inferiormente se existe algum elemento k ∈¯Z, tal que k ≤a, para todo a ∈A. Um tal k se chama cota inferior de A. Observe que na","epistemico":"fato","exemplos":[],"id":"gentil_fundamentos_definicao_35_limitado_inferiormente_seja_a_um_subconjunto_nao","memoria":"semantica","meta":{"arquivo":"gentil/Lopes_Lopes_FUNDAMENTOS_DOS_NUMEROS_pdf.pdf","autor":"Gentil Lopes da Silva","ciencia":"teoria_dos_numeros","dominio":"matematica","nivel":4,"numero":"35","tipo_no":"paper","tipo_teorema":"definição","titulo_legivel":"Fundamentos dos Números"},"origem":"gentil/Gentil_Lopes_FUNDAMENTOS_DOS_NUMEROS_pdf.pdf","palavras_chave":[],"sinonimos":[],"tipo":"conceito","titulo":"gentil fundamentos Definição 35 — Limitado inferiormente) Seja A um subconjunto não vazio de ¯"}
\section{gentil fundamentos Definição 35 — Limitado inferiormente) Seja A um subconjunto não vazio de \ensuremath{\bar{\,}}}
Limitado inferiormente) Seja A um subconjunto não vazio de \ensuremath{\bar{\,}}Z. Diz-se que A é limitado inferiormente se existe algum elemento k \ensuremath{\in}\ensuremath{\bar{\,}}Z, tal que k \ensuremath{\le}a, para todo a \ensuremath{\in}A. Um tal k se chama cota inferior de A. Observe que na
\prov{relações: parte\_de matematica}
\prov{proveniência: gentil/Gentil\_Lopes\_FUNDAMENTOS\_DOS\_NUMEROS\_pdf.pdf \textperiodcentered{} fato \textperiodcentered{} confiança 1.0 \textperiodcentered{} domínio matematica \textperiodcentered{} história 4 versões no jornal}

%CRISTAL {"arestas":[{"alvo":"matematica","confianca":0.95,"epistemico":"fato","origem":"paper","rel":"parte_de"}],"confianca":1.0,"contraexemplos":[],"descricao":"Sejam A e B dois conjuntos Chama-se relação binária de A em B ou apenas relação de A em B todo subconjunto R do produto cartesiano A × B. Portanto, simbolicamente: R é relação de A em B ⇐⇒ R ⊂A × B A princ´ıpio posso ver esta","epistemico":"fato","exemplos":[],"id":"gentil_fundamentos_definicao_4_sejam_a_e_b_dois_conjuntos_chama-se_relacao_binar","memoria":"semantica","meta":{"arquivo":"gentil/Lopes_Lopes_FUNDAMENTOS_DOS_NUMEROS_pdf.pdf","autor":"Gentil Lopes da Silva","ciencia":"teoria_dos_numeros","dominio":"matematica","nivel":4,"numero":"4","tipo_no":"paper","tipo_teorema":"definição","titulo_legivel":"Fundamentos dos Números"},"origem":"gentil/Gentil_Lopes_FUNDAMENTOS_DOS_NUMEROS_pdf.pdf","palavras_chave":[],"sinonimos":[],"tipo":"conceito","titulo":"gentil fundamentos Definição 4 — Sejam A e B dois conjuntos Chama-se relação binária de A em"}
\section{gentil fundamentos Definição 4 — Sejam A e B dois conjuntos Chama-se relação binária de A em}
Sejam A e B dois conjuntos Chama-se relação binária de A em B ou apenas relação de A em B todo subconjunto R do produto cartesiano A $\times$ B. Portanto, simbolicamente: R é relação de A em B \ensuremath{\Leftarrow}\ensuremath{\Rightarrow} R \ensuremath{\subset}A $\times$ B A princ'\ensuremath{\imath}pio posso ver esta
\prov{relações: parte\_de matematica}
\prov{proveniência: gentil/Gentil\_Lopes\_FUNDAMENTOS\_DOS\_NUMEROS\_pdf.pdf \textperiodcentered{} fato \textperiodcentered{} confiança 1.0 \textperiodcentered{} domínio matematica \textperiodcentered{} história 4 versões no jornal}

%CRISTAL {"arestas":[{"alvo":"matematica","confianca":0.95,"epistemico":"fato","origem":"paper","rel":"parte_de"}],"confianca":1.0,"contraexemplos":[],"descricao":"Relação Antissimétrica) Seja R = \u0000A, A, p(x, y) \u0001 uma relação em um conjunto A, R é chamada de relação antissimétrica se, e somente se, quaisquer que sejam os elementos x, y ∈A, se (x, y) ∈G e (y, x) ∈G, então x = y). Em s´ımbolos: R é antissimétrica ⇐⇒(∀x) (∀y) \u0000x, y ∈A, se x R y e y R x ⇒x = y \u0001 De outro modo: se x ̸= y, então possivelmente x estará relacionado com y ou possivelmente y estará relacionado com x, porém nunca ambos. R é antissimétrica ⇐⇒(∀x) (∀y) \u0000x, y ∈A, se x ̸= y e x R y ⇒y R","epistemico":"fato","exemplos":[],"id":"gentil_fundamentos_definicao_8_relacao_antissimetrica_seja_r_a_a_px_y_uma_relaca","memoria":"semantica","meta":{"arquivo":"gentil/Lopes_Lopes_FUNDAMENTOS_DOS_NUMEROS_pdf.pdf","autor":"Gentil Lopes da Silva","ciencia":"teoria_dos_numeros","dominio":"matematica","nivel":4,"numero":"8","tipo_no":"paper","tipo_teorema":"definição","titulo_legivel":"Fundamentos dos Números"},"origem":"gentil/Gentil_Lopes_FUNDAMENTOS_DOS_NUMEROS_pdf.pdf","palavras_chave":[],"sinonimos":[],"tipo":"conceito","titulo":"gentil fundamentos Definição 8 — Relação Antissimétrica) Seja R = \u0000A, A, p(x, y) \u0001 uma relaçã"}
\section{gentil fundamentos Definição 8 — Relação Antissimétrica) Seja R = A, A, p(x, y)  uma relaçã}
Relação Antissimétrica) Seja R = A, A, p(x, y)  uma relação em um conjunto A, R é chamada de relação antissimétrica se, e somente se, quaisquer que sejam os elementos x, y \ensuremath{\in}A, se (x, y) \ensuremath{\in}G e (y, x) \ensuremath{\in}G, então x = y). Em s'\ensuremath{\imath}mbolos: R é antissimétrica \ensuremath{\Leftarrow}\ensuremath{\Rightarrow}(\ensuremath{\forall}x) (\ensuremath{\forall}y) x, y \ensuremath{\in}A, se x R y e y R x \ensuremath{\Rightarrow}x = y  De outro modo: se x \ensuremath{}= y, então possivelmente x estará relacionado com y ou possivelmente y estará relacionado com x, porém nunca ambos. R é antissimétrica \ensuremath{\Leftarrow}\ensuremath{\Rightarrow}(\ensuremath{\forall}x) (\ensuremath{\forall}y) x, y \ensuremath{\in}A, se x \ensuremath{}= y e x R y \ensuremath{\Rightarrow}y R
\prov{relações: parte\_de matematica}
\prov{proveniência: gentil/Gentil\_Lopes\_FUNDAMENTOS\_DOS\_NUMEROS\_pdf.pdf \textperiodcentered{} fato \textperiodcentered{} confiança 1.0 \textperiodcentered{} domínio matematica \textperiodcentered{} história 4 versões no jornal}

%CRISTAL {"arestas":[{"alvo":"matematica","confianca":0.95,"epistemico":"fato","origem":"paper","rel":"parte_de"}],"confianca":1.0,"contraexemplos":[],"descricao":"Relação Transitiva) Seja R = \u0000A, A, p(x, y) \u0001 uma relação em um conjunto A, R é chamada de relação transitiva se, e somente se, quaisquer que sejam os elementos x, y, z ∈A, se (x, y) ∈G e (y, z) ∈G, então (x, z) ∈G. Em poucas palavras: se x está relacionado com y e y está relacionado com z, então x está relacionado com z. Em s´ımbolos: R é transitiva ⇐⇒(∀x)(∀y)(∀z) \u0000x, y, z ∈A, se x R y e y R z ⇒x R z \u0001 Desta","epistemico":"fato","exemplos":[],"id":"gentil_fundamentos_definicao_9_relacao_transitiva_seja_r_a_a_px_y_uma_relacao_em","memoria":"semantica","meta":{"arquivo":"gentil/Lopes_Lopes_FUNDAMENTOS_DOS_NUMEROS_pdf.pdf","autor":"Gentil Lopes da Silva","ciencia":"teoria_dos_numeros","dominio":"matematica","nivel":4,"numero":"9","tipo_no":"paper","tipo_teorema":"definição","titulo_legivel":"Fundamentos dos Números"},"origem":"gentil/Gentil_Lopes_FUNDAMENTOS_DOS_NUMEROS_pdf.pdf","palavras_chave":[],"sinonimos":[],"tipo":"conceito","titulo":"gentil fundamentos Definição 9 — Relação Transitiva) Seja R = \u0000A, A, p(x, y) \u0001 uma relação em"}
\section{gentil fundamentos Definição 9 — Relação Transitiva) Seja R = A, A, p(x, y)  uma relação em}
Relação Transitiva) Seja R = A, A, p(x, y)  uma relação em um conjunto A, R é chamada de relação transitiva se, e somente se, quaisquer que sejam os elementos x, y, z \ensuremath{\in}A, se (x, y) \ensuremath{\in}G e (y, z) \ensuremath{\in}G, então (x, z) \ensuremath{\in}G. Em poucas palavras: se x está relacionado com y e y está relacionado com z, então x está relacionado com z. Em s'\ensuremath{\imath}mbolos: R é transitiva \ensuremath{\Leftarrow}\ensuremath{\Rightarrow}(\ensuremath{\forall}x)(\ensuremath{\forall}y)(\ensuremath{\forall}z) x, y, z \ensuremath{\in}A, se x R y e y R z \ensuremath{\Rightarrow}x R z  Desta
\prov{relações: parte\_de matematica}
\prov{proveniência: gentil/Gentil\_Lopes\_FUNDAMENTOS\_DOS\_NUMEROS\_pdf.pdf \textperiodcentered{} fato \textperiodcentered{} confiança 1.0 \textperiodcentered{} domínio matematica \textperiodcentered{} história 4 versões no jornal}

%CRISTAL {"arestas":[{"alvo":"teoria_dos_numeros","confianca":1.0,"epistemico":"fato","origem":"gentil/Gentil_Lopes_FUNDAMENTOS_DOS_NUMEROS_pdf.pdf","rel":"parte_de"},{"alvo":"gentil_novas_sequencias_aritmeticas_geometricas","confianca":0.85,"epistemico":"fato","origem":"gentil/Gentil_Lopes_FUNDAMENTOS_DOS_NUMEROS_pdf.pdf","rel":"paralelo"},{"alvo":"numero","confianca":0.95,"epistemico":"fato","origem":"gentil/Gentil_Lopes_FUNDAMENTOS_DOS_NUMEROS_pdf.pdf","rel":"define"}],"confianca":1.0,"contraexemplos":[],"descricao":"Gentil Lopes da Silva (2016). Construção N→Z→Q→R→C→H; números azuis e vermelhos; fundamentação filosófica.","epistemico":"fato","exemplos":[],"id":"gentil_fundamentos_dos_numeros","memoria":"semantica","meta":{"arquivo":"gentil/Lopes_Lopes_FUNDAMENTOS_DOS_NUMEROS_pdf.pdf","autor":"Gentil Lopes da Silva","ciencia":"teoria_dos_numeros","dominio":"matematica","nivel":4,"tipo_no":"paper","titulo_legivel":"Fundamentos dos Números"},"origem":"livro","palavras_chave":[],"sinonimos":["Fundamentos dos Números","livro Gentil"],"tipo":"conceito","titulo":"gentil fundamentos dos numeros"}
\section{gentil fundamentos dos numeros}
Gentil Lopes da Silva (2016). Construção N\ensuremath{\to}Z\ensuremath{\to}Q\ensuremath{\to}R\ensuremath{\to}C\ensuremath{\to}H; números azuis e vermelhos; fundamentação filosófica.
\prov{sinónimos: Fundamentos dos Números; livro Gentil}
\prov{relações: parte\_de teoria\_dos\_numeros; paralelo gentil\_novas\_sequencias\_aritmeticas\_geometricas; define numero}
\prov{proveniência: livro \textperiodcentered{} fato \textperiodcentered{} confiança 1.0 \textperiodcentered{} domínio matematica \textperiodcentered{} história 9 versões no jornal}

%CRISTAL {"arestas":[],"confianca":1.0,"contraexemplos":[],"descricao":"p 94), devemos provar que (p + 1) + m = m + (p + 1) (3.6) Aplicando σ na H.I. temos, σ(p + m) = σ(m + p). Então, por (ii) é verdade que p + σ(m) = m + σ(p) Ainda pelo","epistemico":"fato","exemplos":[],"id":"gentil_fundamentos_lema_1_p_94_devemos_provar_que_p_1_m_m_p_1_36_ap","memoria":"semantica","meta":{"arquivo":"gentil/Lopes_Lopes_FUNDAMENTOS_DOS_NUMEROS_pdf.pdf","autor":"Gentil Lopes da Silva","ciencia":"teoria_dos_numeros","dominio":"matematica","nivel":4,"numero":"1","tipo_no":"paper","tipo_teorema":"lema","titulo_legivel":"Fundamentos dos Números"},"origem":"gentil/Gentil_Lopes_FUNDAMENTOS_DOS_NUMEROS_pdf.pdf","palavras_chave":[],"sinonimos":[],"tipo":"conceito","titulo":"gentil fundamentos Lema 1 — p 94), devemos provar que (p + 1) + m = m + (p + 1) (3.6) Ap"}
\section{gentil fundamentos Lema 1 — p 94), devemos provar que (p + 1) + m = m + (p + 1) (3.6) Ap}
p 94), devemos provar que (p + 1) + m = m + (p + 1) (3.6) Aplicando \ensuremath{\sigma} na H.I. temos, \ensuremath{\sigma}(p + m) = \ensuremath{\sigma}(m + p). Então, por (ii) é verdade que p + \ensuremath{\sigma}(m) = m + \ensuremath{\sigma}(p) Ainda pelo
\prov{proveniência: gentil/Gentil\_Lopes\_FUNDAMENTOS\_DOS\_NUMEROS\_pdf.pdf \textperiodcentered{} fato \textperiodcentered{} confiança 1.0 \textperiodcentered{} domínio matematica \textperiodcentered{} história 5 versões no jornal}

%CRISTAL {"arestas":[],"confianca":1.0,"contraexemplos":[],"descricao":"p 94), devemos provar que (T.I.) m = n ⇒ m + (r + 1) = n + (r + 1) é verdadeira. Então, iniciando com a hipótese de indução, temos m = n ⇒ m + r = n + r ⇒ (m + r) + 1 = (n + r) + 1 ⇒ m + (r + 1) = n + (r + 1) Na segunda implicação acima utilizamos o axioma P2 (p. 83). Do axioma de indução temos que A = N e segue a tese. ■ Vamos agora provar a especificação M2 de ρ. (p. 92) 99\nProposição 7 (1 é elemento neutro). Para todo n ∈N tem-se: n · 1 = 1 · n = n Nota: Enfatizamos que as igualdades acima nã","epistemico":"fato","exemplos":[],"id":"gentil_fundamentos_lema_1_p_94_devemos_provar_que_ti_m_n_m_r_1_n","memoria":"semantica","meta":{"arquivo":"gentil/Lopes_Lopes_FUNDAMENTOS_DOS_NUMEROS_pdf.pdf","autor":"Gentil Lopes da Silva","ciencia":"teoria_dos_numeros","dominio":"matematica","nivel":4,"numero":"1","tipo_no":"paper","tipo_teorema":"lema","titulo_legivel":"Fundamentos dos Números"},"origem":"gentil/Gentil_Lopes_FUNDAMENTOS_DOS_NUMEROS_pdf.pdf","palavras_chave":[],"sinonimos":[],"tipo":"conceito","titulo":"gentil fundamentos Lema 1 — p 94), devemos provar que (T.I.) m = n ⇒ m + (r + 1) = n + ("}
\section{gentil fundamentos Lema 1 — p 94), devemos provar que (T.I.) m = n \ensuremath{\Rightarrow} m + (r + 1) = n + (}
p 94), devemos provar que (T.I.) m = n \ensuremath{\Rightarrow} m + (r + 1) = n + (r + 1) é verdadeira. Então, iniciando com a hipótese de indução, temos m = n \ensuremath{\Rightarrow} m + r = n + r \ensuremath{\Rightarrow} (m + r) + 1 = (n + r) + 1 \ensuremath{\Rightarrow} m + (r + 1) = n + (r + 1) Na segunda implicação acima utilizamos o axioma P2 (p. 83). Do axioma de indução temos que A = N e segue a tese. \rule{1ex}{1ex} Vamos agora provar a especificação M2 de \ensuremath{\rho}. (p. 92) 99
Proposição 7 (1 é elemento neutro). Para todo n \ensuremath{\in}N tem-se: n \textperiodcentered  1 = 1 \textperiodcentered  n = n Nota: Enfatizamos que as igualdades acima nã
\prov{proveniência: gentil/Gentil\_Lopes\_FUNDAMENTOS\_DOS\_NUMEROS\_pdf.pdf \textperiodcentered{} fato \textperiodcentered{} confiança 1.0 \textperiodcentered{} domínio matematica \textperiodcentered{} história 5 versões no jornal}

%CRISTAL {"arestas":[],"confianca":1.0,"contraexemplos":[],"descricao":"p 94), devemos provar que (T.I.) m + (r + 1) = n + (r + 1) ⇒ m = n Então, inicialmente vamos assumir que m + (r + 1) = n + (r + 1) Vamos aplicar associatividade (m + r) + 1 = (n + r) + 1 Pelo","epistemico":"fato","exemplos":[],"id":"gentil_fundamentos_lema_1_p_94_devemos_provar_que_ti_m_r_1_n_r_1","memoria":"semantica","meta":{"arquivo":"gentil/Lopes_Lopes_FUNDAMENTOS_DOS_NUMEROS_pdf.pdf","autor":"Gentil Lopes da Silva","ciencia":"teoria_dos_numeros","dominio":"matematica","nivel":4,"numero":"1","tipo_no":"paper","tipo_teorema":"lema","titulo_legivel":"Fundamentos dos Números"},"origem":"gentil/Gentil_Lopes_FUNDAMENTOS_DOS_NUMEROS_pdf.pdf","palavras_chave":[],"sinonimos":[],"tipo":"conceito","titulo":"gentil fundamentos Lema 1 — p 94), devemos provar que (T.I.) m + (r + 1) = n + (r + 1) ⇒"}
\section{gentil fundamentos Lema 1 — p 94), devemos provar que (T.I.) m + (r + 1) = n + (r + 1) \ensuremath{\Rightarrow}}
p 94), devemos provar que (T.I.) m + (r + 1) = n + (r + 1) \ensuremath{\Rightarrow} m = n Então, inicialmente vamos assumir que m + (r + 1) = n + (r + 1) Vamos aplicar associatividade (m + r) + 1 = (n + r) + 1 Pelo
\prov{proveniência: gentil/Gentil\_Lopes\_FUNDAMENTOS\_DOS\_NUMEROS\_pdf.pdf \textperiodcentered{} fato \textperiodcentered{} confiança 1.0 \textperiodcentered{} domínio matematica \textperiodcentered{} história 5 versões no jornal}

%CRISTAL {"arestas":[],"confianca":1.0,"contraexemplos":[],"descricao":"Para todo natural m, tem-se σ(m) = m + 1 e σ(m) = 1 + m; portanto m + 1 = 1 + m Observação: Este lema não é necessáriamente “trivial” −a ponto de dispensar uma prova −; por exemplo, nos naturais vermelhos um caso par- ticular é, σ(m) = 1 + m ←→ σ( ) = + Esta igualdade não é imediata, inclusive porque está amarrada com a","epistemico":"fato","exemplos":[],"id":"gentil_fundamentos_lema_1_para_todo_natural_m_tem-se_σm_m_1_e_σm_1_m_por","memoria":"semantica","meta":{"arquivo":"gentil/Lopes_Lopes_FUNDAMENTOS_DOS_NUMEROS_pdf.pdf","autor":"Gentil Lopes da Silva","ciencia":"teoria_dos_numeros","dominio":"matematica","nivel":4,"numero":"1","tipo_no":"paper","tipo_teorema":"lema","titulo_legivel":"Fundamentos dos Números"},"origem":"gentil/Gentil_Lopes_FUNDAMENTOS_DOS_NUMEROS_pdf.pdf","palavras_chave":[],"sinonimos":[],"tipo":"conceito","titulo":"gentil fundamentos Lema 1 — Para todo natural m, tem-se σ(m) = m + 1 e σ(m) = 1 + m; por"}
\section{gentil fundamentos Lema 1 — Para todo natural m, tem-se \ensuremath{\sigma}(m) = m + 1 e \ensuremath{\sigma}(m) = 1 + m; por}
Para todo natural m, tem-se \ensuremath{\sigma}(m) = m + 1 e \ensuremath{\sigma}(m) = 1 + m; portanto m + 1 = 1 + m Observação: Este lema não é necessáriamente “trivial” \ensuremath{-}a ponto de dispensar uma prova \ensuremath{-}; por exemplo, nos naturais vermelhos um caso par- ticular é, \ensuremath{\sigma}(m) = 1 + m \ensuremath{\leftarrow}\ensuremath{\to} \ensuremath{\sigma}( ) = + Esta igualdade não é imediata, inclusive porque está amarrada com a
\prov{proveniência: gentil/Gentil\_Lopes\_FUNDAMENTOS\_DOS\_NUMEROS\_pdf.pdf \textperiodcentered{} fato \textperiodcentered{} confiança 1.0 \textperiodcentered{} domínio matematica \textperiodcentered{} história 5 versões no jornal}

%CRISTAL {"arestas":[{"alvo":"matematica","confianca":0.95,"epistemico":"fato","origem":"paper","rel":"parte_de"}],"confianca":1.0,"contraexemplos":[],"descricao":"podemos escrever σ(m + r) = σ(n + r) 98\npelo axioma P2 (p 83) temos que σ(m + r) = σ(n + r) ⇒ m + r = n + r pela hipótese de indução decorre que m = n. Do axioma de indução temos que A = N e segue a tese. ■ Se devemos provar que podemos “cortar ” quantidades iguais de ambos os membros de uma igualdade, então devemos provar a “operação inversa”. Proposição 6 (Lei do acréscimo). Para m, n e r naturais, vale: m = n ⇒ m + r = n + r Prova: Suponhamos m e n naturais fixos e considere A o conjunto dos","epistemico":"fato","exemplos":[],"id":"gentil_fundamentos_lema_1_podemos_escrever_σm_r_σn_r_98_pelo_axioma_p2_p_83","memoria":"semantica","meta":{"arquivo":"gentil/Lopes_Lopes_FUNDAMENTOS_DOS_NUMEROS_pdf.pdf","autor":"Gentil Lopes da Silva","ciencia":"teoria_dos_numeros","dominio":"matematica","nivel":4,"numero":"1","tipo_no":"paper","tipo_teorema":"lema","titulo_legivel":"Fundamentos dos Números"},"origem":"gentil/Gentil_Lopes_FUNDAMENTOS_DOS_NUMEROS_pdf.pdf","palavras_chave":[],"sinonimos":[],"tipo":"conceito","titulo":"gentil fundamentos Lema 1 — podemos escrever σ(m + r) = σ(n + r) 98\npelo axioma P2 (p 83"}
\section{gentil fundamentos Lema 1 — podemos escrever \ensuremath{\sigma}(m + r) = \ensuremath{\sigma}(n + r) 98
pelo axioma P2 (p 83}
podemos escrever \ensuremath{\sigma}(m + r) = \ensuremath{\sigma}(n + r) 98
pelo axioma P2 (p 83) temos que \ensuremath{\sigma}(m + r) = \ensuremath{\sigma}(n + r) \ensuremath{\Rightarrow} m + r = n + r pela hipótese de indução decorre que m = n. Do axioma de indução temos que A = N e segue a tese. \rule{1ex}{1ex} Se devemos provar que podemos “cortar ” quantidades iguais de ambos os membros de uma igualdade, então devemos provar a “operação inversa”. Proposição 6 (Lei do acréscimo). Para m, n e r naturais, vale: m = n \ensuremath{\Rightarrow} m + r = n + r Prova: Suponhamos m e n naturais fixos e considere A o conjunto dos
\prov{relações: parte\_de matematica}
\prov{proveniência: gentil/Gentil\_Lopes\_FUNDAMENTOS\_DOS\_NUMEROS\_pdf.pdf \textperiodcentered{} fato \textperiodcentered{} confiança 1.0 \textperiodcentered{} domínio matematica \textperiodcentered{} história 4 versões no jornal}

%CRISTAL {"arestas":[{"alvo":"matematica","confianca":0.95,"epistemico":"fato","origem":"paper","rel":"parte_de"}],"confianca":1.0,"contraexemplos":[],"descricao":"Sejam α e β inteiros quaisquer Então α ≤β ⇔ α −β ≤¯0 Prova: Consideremos α = (a, b), β = (c, d). Então α ≤β ⇔ (a, b) ≤(c, d) ⇔ a + d ≤b + c ⇔ (a + d) + n ≤(b + c) + n ⇔ (a + d, b + c) ≤(n, n) ⇔ (a, b) + (d, c) ≤(n, n) ⇔ (a, b) + \u0000−(c, d) \u0001 ≤(n, n) ⇔ α + (−β) ≤¯0 ■","epistemico":"fato","exemplos":[],"id":"gentil_fundamentos_lema_4_sejam_α_e_β_inteiros_quaisquer_entao_α_β_α_β_0_prova","memoria":"semantica","meta":{"arquivo":"gentil/Lopes_Lopes_FUNDAMENTOS_DOS_NUMEROS_pdf.pdf","autor":"Gentil Lopes da Silva","ciencia":"teoria_dos_numeros","dominio":"matematica","nivel":4,"numero":"4","tipo_no":"paper","tipo_teorema":"lema","titulo_legivel":"Fundamentos dos Números"},"origem":"gentil/Gentil_Lopes_FUNDAMENTOS_DOS_NUMEROS_pdf.pdf","palavras_chave":[],"sinonimos":[],"tipo":"conceito","titulo":"gentil fundamentos Lema 4 — Sejam α e β inteiros quaisquer Então α ≤β ⇔ α −β ≤¯0 Prova:"}
\section{gentil fundamentos Lema 4 — Sejam \ensuremath{\alpha} e \ensuremath{\beta} inteiros quaisquer Então \ensuremath{\alpha} \ensuremath{\le}\ensuremath{\beta} \ensuremath{\Leftrightarrow} \ensuremath{\alpha} \ensuremath{-}\ensuremath{\beta} \ensuremath{\le}\ensuremath{\bar{\,}}0 Prova:}
Sejam \ensuremath{\alpha} e \ensuremath{\beta} inteiros quaisquer Então \ensuremath{\alpha} \ensuremath{\le}\ensuremath{\beta} \ensuremath{\Leftrightarrow} \ensuremath{\alpha} \ensuremath{-}\ensuremath{\beta} \ensuremath{\le}\ensuremath{\bar{\,}}0 Prova: Consideremos \ensuremath{\alpha} = (a, b), \ensuremath{\beta} = (c, d). Então \ensuremath{\alpha} \ensuremath{\le}\ensuremath{\beta} \ensuremath{\Leftrightarrow} (a, b) \ensuremath{\le}(c, d) \ensuremath{\Leftrightarrow} a + d \ensuremath{\le}b + c \ensuremath{\Leftrightarrow} (a + d) + n \ensuremath{\le}(b + c) + n \ensuremath{\Leftrightarrow} (a + d, b + c) \ensuremath{\le}(n, n) \ensuremath{\Leftrightarrow} (a, b) + (d, c) \ensuremath{\le}(n, n) \ensuremath{\Leftrightarrow} (a, b) + \ensuremath{-}(c, d)  \ensuremath{\le}(n, n) \ensuremath{\Leftrightarrow} \ensuremath{\alpha} + (\ensuremath{-}\ensuremath{\beta}) \ensuremath{\le}\ensuremath{\bar{\,}}0 \rule{1ex}{1ex}
\prov{relações: parte\_de matematica}
\prov{proveniência: gentil/Gentil\_Lopes\_FUNDAMENTOS\_DOS\_NUMEROS\_pdf.pdf \textperiodcentered{} fato \textperiodcentered{} confiança 1.0 \textperiodcentered{} domínio matematica \textperiodcentered{} história 4 versões no jornal}

%CRISTAL {"arestas":[{"alvo":"matematica","confianca":0.95,"epistemico":"fato","origem":"paper","rel":"parte_de"}],"confianca":1.0,"contraexemplos":[],"descricao":"Sejam α e β inteiros Então, se α ≤¯0 e β ≥¯0 ⇒ α · β ≤¯0. Prova: Consideremos α = (a, b), β = (c, d) e ¯0 = (n, n). Como α ≤¯0 e β ≥¯0 temos (def. 33, p. 171) (a, b) ≤(n, n) ⇒ a + n ≤b + n ⇒ a ≤b e (c, d) ≥(n, n) ⇒ c + n ≥d + n ⇒ c ≥d Como (def. 31, p. 159) α · β = (a, b) · (c, d) = (ac + bd, ad + bc) Para mostrar que α · β ≤¯0 é suficiente mostrar que ac + bd ≤ad + bc Como d ≤c por def. 22 (p. 105) existe r ∈N tal que d + r = c. Por outro lado, como a ≤b, temos que ar ≤br, logo ac + bd = a(d +","epistemico":"fato","exemplos":[],"id":"gentil_fundamentos_lema_5_sejam_α_e_β_inteiros_entao_se_α_0_e_β_0_α_β_0_pr","memoria":"semantica","meta":{"arquivo":"gentil/Lopes_Lopes_FUNDAMENTOS_DOS_NUMEROS_pdf.pdf","autor":"Gentil Lopes da Silva","ciencia":"teoria_dos_numeros","dominio":"matematica","nivel":4,"numero":"5","tipo_no":"paper","tipo_teorema":"lema","titulo_legivel":"Fundamentos dos Números"},"origem":"gentil/Gentil_Lopes_FUNDAMENTOS_DOS_NUMEROS_pdf.pdf","palavras_chave":[],"sinonimos":[],"tipo":"conceito","titulo":"gentil fundamentos Lema 5 — Sejam α e β inteiros Então, se α ≤¯0 e β ≥¯0 ⇒ α · β ≤¯0. Pr"}
\section{gentil fundamentos Lema 5 — Sejam \ensuremath{\alpha} e \ensuremath{\beta} inteiros Então, se \ensuremath{\alpha} \ensuremath{\le}\ensuremath{\bar{\,}}0 e \ensuremath{\beta} \ensuremath{\ge}\ensuremath{\bar{\,}}0 \ensuremath{\Rightarrow} \ensuremath{\alpha} \textperiodcentered  \ensuremath{\beta} \ensuremath{\le}\ensuremath{\bar{\,}}0. Pr}
Sejam \ensuremath{\alpha} e \ensuremath{\beta} inteiros Então, se \ensuremath{\alpha} \ensuremath{\le}\ensuremath{\bar{\,}}0 e \ensuremath{\beta} \ensuremath{\ge}\ensuremath{\bar{\,}}0 \ensuremath{\Rightarrow} \ensuremath{\alpha} \textperiodcentered  \ensuremath{\beta} \ensuremath{\le}\ensuremath{\bar{\,}}0. Prova: Consideremos \ensuremath{\alpha} = (a, b), \ensuremath{\beta} = (c, d) e \ensuremath{\bar{\,}}0 = (n, n). Como \ensuremath{\alpha} \ensuremath{\le}\ensuremath{\bar{\,}}0 e \ensuremath{\beta} \ensuremath{\ge}\ensuremath{\bar{\,}}0 temos (def. 33, p. 171) (a, b) \ensuremath{\le}(n, n) \ensuremath{\Rightarrow} a + n \ensuremath{\le}b + n \ensuremath{\Rightarrow} a \ensuremath{\le}b e (c, d) \ensuremath{\ge}(n, n) \ensuremath{\Rightarrow} c + n \ensuremath{\ge}d + n \ensuremath{\Rightarrow} c \ensuremath{\ge}d Como (def. 31, p. 159) \ensuremath{\alpha} \textperiodcentered  \ensuremath{\beta} = (a, b) \textperiodcentered  (c, d) = (ac + bd, ad + bc) Para mostrar que \ensuremath{\alpha} \textperiodcentered  \ensuremath{\beta} \ensuremath{\le}\ensuremath{\bar{\,}}0 é suficiente mostrar que ac + bd \ensuremath{\le}ad + bc Como d \ensuremath{\le}c por def. 22 (p. 105) existe r \ensuremath{\in}N tal que d + r = c. Por outro lado, como a \ensuremath{\le}b, temos que ar \ensuremath{\le}br, logo ac + bd = a(d +
\prov{relações: parte\_de matematica}
\prov{proveniência: gentil/Gentil\_Lopes\_FUNDAMENTOS\_DOS\_NUMEROS\_pdf.pdf \textperiodcentered{} fato \textperiodcentered{} confiança 1.0 \textperiodcentered{} domínio matematica \textperiodcentered{} história 4 versões no jornal}

%CRISTAL {"arestas":[{"alvo":"gentil_fundamentos_capitulo_1_o_que_e_um_n_umero","confianca":1.0,"epistemico":"fato","origem":"gentil/Gentil_Lopes_FUNDAMENTOS_DOS_NUMEROS_pdf.pdf","rel":"prova"}],"confianca":1.0,"contraexemplos":[],"descricao":"Seja α = (a, b) um número inteiro, então α < ¯0 ⇔ −α > ¯0 Prova: Temos α = (a, b) < ¯0 ⇔ a < b ⇔ b > a ⇔ (b, a) > ¯0 ⇔ −α = (b, a) > ¯0 ■ 180","epistemico":"fato","exemplos":[],"id":"gentil_fundamentos_lema_6_seja_α_a_b_um_numero_inteiro_entao_α_0_α_0_pr","memoria":"semantica","meta":{"arquivo":"gentil/Lopes_Lopes_FUNDAMENTOS_DOS_NUMEROS_pdf.pdf","autor":"Gentil Lopes da Silva","ciencia":"teoria_dos_numeros","dominio":"matematica","nivel":4,"numero":"6","tipo_no":"paper","tipo_teorema":"lema","titulo_legivel":"Fundamentos dos Números"},"origem":"gentil/Gentil_Lopes_FUNDAMENTOS_DOS_NUMEROS_pdf.pdf","palavras_chave":[],"sinonimos":[],"tipo":"conceito","titulo":"gentil fundamentos Lema 6 — Seja α = (a, b) um número inteiro, então α < ¯0 ⇔ −α > ¯0 Pr"}
\section{gentil fundamentos Lema 6 — Seja \ensuremath{\alpha} = (a, b) um número inteiro, então \ensuremath{\alpha} < \ensuremath{\bar{\,}}0 \ensuremath{\Leftrightarrow} \ensuremath{-}\ensuremath{\alpha} > \ensuremath{\bar{\,}}0 Pr}
Seja \ensuremath{\alpha} = (a, b) um número inteiro, então \ensuremath{\alpha} < \ensuremath{\bar{\,}}0 \ensuremath{\Leftrightarrow} \ensuremath{-}\ensuremath{\alpha} > \ensuremath{\bar{\,}}0 Prova: Temos \ensuremath{\alpha} = (a, b) < \ensuremath{\bar{\,}}0 \ensuremath{\Leftrightarrow} a < b \ensuremath{\Leftrightarrow} b > a \ensuremath{\Leftrightarrow} (b, a) > \ensuremath{\bar{\,}}0 \ensuremath{\Leftrightarrow} \ensuremath{-}\ensuremath{\alpha} = (b, a) > \ensuremath{\bar{\,}}0 \rule{1ex}{1ex} 180
\prov{relações: prova gentil\_fundamentos\_capitulo\_1\_o\_que\_e\_um\_n\_umero}
\prov{proveniência: gentil/Gentil\_Lopes\_FUNDAMENTOS\_DOS\_NUMEROS\_pdf.pdf \textperiodcentered{} fato \textperiodcentered{} confiança 1.0 \textperiodcentered{} domínio matematica \textperiodcentered{} história 5 versões no jornal}

%CRISTAL {"arestas":[{"alvo":"matematica","confianca":0.95,"epistemico":"fato","origem":"paper","rel":"parte_de"}],"confianca":1.0,"contraexemplos":[],"descricao":", p 183 Teorema 18, p. 163 Teorema 20, p. 164 Proposição 14, p. 184 Lei do Corte Prop. 9, p. 169 Lei do Corte (+, N) Lei do Corte (·, N) 2 x + 3 = 1 Claro, este fluxograma não é único, podemos acrescentar ainda outras dependências. O momento agora é oportuno para elucidarmos uma questão levantada na Revista do Professor de Matemática (RPM/61, 2006) onde lemos: Uma preocupação manifestada pelo leitor Nelson O. F. Correa é a identi- ficação dos números naturais com os inteiros positivos e, como ex","epistemico":"fato","exemplos":[],"id":"gentil_fundamentos_lema_7_p_183_teorema_18_p_163_teorema_20_p_164_proposicao_14","memoria":"semantica","meta":{"arquivo":"gentil/Lopes_Lopes_FUNDAMENTOS_DOS_NUMEROS_pdf.pdf","autor":"Gentil Lopes da Silva","ciencia":"teoria_dos_numeros","dominio":"matematica","nivel":4,"numero":"7","tipo_no":"paper","tipo_teorema":"lema","titulo_legivel":"Fundamentos dos Números"},"origem":"gentil/Gentil_Lopes_FUNDAMENTOS_DOS_NUMEROS_pdf.pdf","palavras_chave":[],"sinonimos":[],"tipo":"conceito","titulo":"gentil fundamentos Lema 7 — , p 183 Teorema 18, p. 163 Teorema 20, p. 164 Proposição 14,"}
\section{gentil fundamentos Lema 7 — , p 183 Teorema 18, p. 163 Teorema 20, p. 164 Proposição 14,}
, p 183 Teorema 18, p. 163 Teorema 20, p. 164 Proposição 14, p. 184 Lei do Corte Prop. 9, p. 169 Lei do Corte (+, N) Lei do Corte (\textperiodcentered , N) 2 x + 3 = 1 Claro, este fluxograma não é único, podemos acrescentar ainda outras dependências. O momento agora é oportuno para elucidarmos uma questão levantada na Revista do Professor de Matemática (RPM/61, 2006) onde lemos: Uma preocupação manifestada pelo leitor Nelson O. F. Correa é a identi- ficação dos números naturais com os inteiros positivos e, como ex
\prov{relações: parte\_de matematica}
\prov{proveniência: gentil/Gentil\_Lopes\_FUNDAMENTOS\_DOS\_NUMEROS\_pdf.pdf \textperiodcentered{} fato \textperiodcentered{} confiança 1.0 \textperiodcentered{} domínio matematica \textperiodcentered{} história 4 versões no jornal}

%CRISTAL {"arestas":[{"alvo":"matematica","confianca":0.95,"epistemico":"fato","origem":"paper","rel":"parte_de"}],"confianca":1.0,"contraexemplos":[],"descricao":"Sendo a e b dois naturais, com a ≥b, vale b + (a −b) = a Prova: a −b = a + (−b) = (a, 0) + \u0000−(b, 0) \u0001 = (a, 0) + (0, b) = (a, b) Logo, b + (a −b) = (b, 0) + (a, b) = (a + b, b) Pela","epistemico":"fato","exemplos":[],"id":"gentil_fundamentos_lema_7_sendo_a_e_b_dois_naturais_com_a_b_vale_b_a_b_a_pro","memoria":"semantica","meta":{"arquivo":"gentil/Lopes_Lopes_FUNDAMENTOS_DOS_NUMEROS_pdf.pdf","autor":"Gentil Lopes da Silva","ciencia":"teoria_dos_numeros","dominio":"matematica","nivel":4,"numero":"7","tipo_no":"paper","tipo_teorema":"lema","titulo_legivel":"Fundamentos dos Números"},"origem":"gentil/Gentil_Lopes_FUNDAMENTOS_DOS_NUMEROS_pdf.pdf","palavras_chave":[],"sinonimos":[],"tipo":"conceito","titulo":"gentil fundamentos Lema 7 — Sendo a e b dois naturais, com a ≥b, vale b + (a −b) = a Pro"}
\section{gentil fundamentos Lema 7 — Sendo a e b dois naturais, com a \ensuremath{\ge}b, vale b + (a \ensuremath{-}b) = a Pro}
Sendo a e b dois naturais, com a \ensuremath{\ge}b, vale b + (a \ensuremath{-}b) = a Prova: a \ensuremath{-}b = a + (\ensuremath{-}b) = (a, 0) + \ensuremath{-}(b, 0)  = (a, 0) + (0, b) = (a, b) Logo, b + (a \ensuremath{-}b) = (b, 0) + (a, b) = (a + b, b) Pela
\prov{relações: parte\_de matematica}
\prov{proveniência: gentil/Gentil\_Lopes\_FUNDAMENTOS\_DOS\_NUMEROS\_pdf.pdf \textperiodcentered{} fato \textperiodcentered{} confiança 1.0 \textperiodcentered{} domínio matematica \textperiodcentered{} história 4 versões no jornal}

%CRISTAL {"arestas":[{"alvo":"matematica","confianca":0.95,"epistemico":"fato","origem":"paper","rel":"parte_de"}],"confianca":1.0,"contraexemplos":[],"descricao":"(p 201) Considerações análogas valem para a multiplicação, devemos mostrar que esta operação está bem definida; isto é, que independe dos representan- tes tomados nas classes. Isto está feito no apêndice (","epistemico":"fato","exemplos":[],"id":"gentil_fundamentos_lema_8_p_201_consideracoes_analogas_valem_para_a_multiplicaca","memoria":"semantica","meta":{"arquivo":"gentil/Lopes_Lopes_FUNDAMENTOS_DOS_NUMEROS_pdf.pdf","autor":"Gentil Lopes da Silva","ciencia":"teoria_dos_numeros","dominio":"matematica","nivel":4,"numero":"8","tipo_no":"paper","tipo_teorema":"lema","titulo_legivel":"Fundamentos dos Números"},"origem":"gentil/Gentil_Lopes_FUNDAMENTOS_DOS_NUMEROS_pdf.pdf","palavras_chave":[],"sinonimos":[],"tipo":"conceito","titulo":"gentil fundamentos Lema 8 — (p 201) Considerações análogas valem para a multiplicação, d"}
\section{gentil fundamentos Lema 8 — (p 201) Considerações análogas valem para a multiplicação, d}
(p 201) Considerações análogas valem para a multiplicação, devemos mostrar que esta operação está bem definida; isto é, que independe dos representan- tes tomados nas classes. Isto está feito no apêndice (
\prov{relações: parte\_de matematica}
\prov{proveniência: gentil/Gentil\_Lopes\_FUNDAMENTOS\_DOS\_NUMEROS\_pdf.pdf \textperiodcentered{} fato \textperiodcentered{} confiança 1.0 \textperiodcentered{} domínio matematica \textperiodcentered{} história 4 versões no jornal}

%CRISTAL {"arestas":[{"alvo":"gentil_fundamentos_capitulo_1_o_que_e_um_n_umero","confianca":1.0,"epistemico":"fato","origem":"gentil/Gentil_Lopes_FUNDAMENTOS_DOS_NUMEROS_pdf.pdf","rel":"prova"}],"confianca":1.0,"contraexemplos":[],"descricao":"Sejam α = (a, b), α′ = (a′, b′), β = (c, d), β′ = (c′, d′) números em Z Se α = α′ e β = β′, então α + β = (a + c, b + d) = (a′ + c′, b′ + d′) = α′ + β′ Prova: Da hipótese, temos que a + b′ = a′ + b c + d′ = c′ + d Logo, (a + c) + (b′ + d′) = (a′ + c′) + (b + d) donde (a + c, b + d) = (a′ + c′, b′ + d′) ■","epistemico":"fato","exemplos":[],"id":"gentil_fundamentos_lema_8_sejam_α_a_b_α_a_b_β_c_d_β_c_d_n","memoria":"semantica","meta":{"arquivo":"gentil/Lopes_Lopes_FUNDAMENTOS_DOS_NUMEROS_pdf.pdf","autor":"Gentil Lopes da Silva","ciencia":"teoria_dos_numeros","dominio":"matematica","nivel":4,"numero":"8","tipo_no":"paper","tipo_teorema":"lema","titulo_legivel":"Fundamentos dos Números"},"origem":"gentil/Gentil_Lopes_FUNDAMENTOS_DOS_NUMEROS_pdf.pdf","palavras_chave":[],"sinonimos":[],"tipo":"conceito","titulo":"gentil fundamentos Lema 8 — Sejam α = (a, b), α′ = (a′, b′), β = (c, d), β′ = (c′, d′) n"}
\section{gentil fundamentos Lema 8 — Sejam \ensuremath{\alpha} = (a, b), \ensuremath{\alpha}\ensuremath{'} = (a\ensuremath{'}, b\ensuremath{'}), \ensuremath{\beta} = (c, d), \ensuremath{\beta}\ensuremath{'} = (c\ensuremath{'}, d\ensuremath{'}) n}
Sejam \ensuremath{\alpha} = (a, b), \ensuremath{\alpha}\ensuremath{'} = (a\ensuremath{'}, b\ensuremath{'}), \ensuremath{\beta} = (c, d), \ensuremath{\beta}\ensuremath{'} = (c\ensuremath{'}, d\ensuremath{'}) números em Z Se \ensuremath{\alpha} = \ensuremath{\alpha}\ensuremath{'} e \ensuremath{\beta} = \ensuremath{\beta}\ensuremath{'}, então \ensuremath{\alpha} + \ensuremath{\beta} = (a + c, b + d) = (a\ensuremath{'} + c\ensuremath{'}, b\ensuremath{'} + d\ensuremath{'}) = \ensuremath{\alpha}\ensuremath{'} + \ensuremath{\beta}\ensuremath{'} Prova: Da hipótese, temos que a + b\ensuremath{'} = a\ensuremath{'} + b c + d\ensuremath{'} = c\ensuremath{'} + d Logo, (a + c) + (b\ensuremath{'} + d\ensuremath{'}) = (a\ensuremath{'} + c\ensuremath{'}) + (b + d) donde (a + c, b + d) = (a\ensuremath{'} + c\ensuremath{'}, b\ensuremath{'} + d\ensuremath{'}) \rule{1ex}{1ex}
\prov{relações: prova gentil\_fundamentos\_capitulo\_1\_o\_que\_e\_um\_n\_umero}
\prov{proveniência: gentil/Gentil\_Lopes\_FUNDAMENTOS\_DOS\_NUMEROS\_pdf.pdf \textperiodcentered{} fato \textperiodcentered{} confiança 1.0 \textperiodcentered{} domínio matematica \textperiodcentered{} história 5 versões no jornal}

%CRISTAL {"arestas":[{"alvo":"matematica","confianca":0.95,"epistemico":"fato","origem":"paper","rel":"parte_de"}],"confianca":1.0,"contraexemplos":[],"descricao":"p 201). Nosso próximo passo será demonstrar que a adição satisfaz as quatro primeiras propriedades em ν: (p. 154) ν = A1, A2, A3, A4 | z ?, M1, M2, M3, D, Ordenado, PBO","epistemico":"fato","exemplos":[],"id":"gentil_fundamentos_lema_9_p_201_nosso_proximo_passo_sera_demonstrar_que_a_adicao","memoria":"semantica","meta":{"arquivo":"gentil/Lopes_Lopes_FUNDAMENTOS_DOS_NUMEROS_pdf.pdf","autor":"Gentil Lopes da Silva","ciencia":"teoria_dos_numeros","dominio":"matematica","nivel":4,"numero":"9","tipo_no":"paper","tipo_teorema":"lema","titulo_legivel":"Fundamentos dos Números"},"origem":"gentil/Gentil_Lopes_FUNDAMENTOS_DOS_NUMEROS_pdf.pdf","palavras_chave":[],"sinonimos":[],"tipo":"conceito","titulo":"gentil fundamentos Lema 9 — , p 201). Nosso próximo passo será demonstrar que a adição s"}
\section{gentil fundamentos Lema 9 — , p 201). Nosso próximo passo será demonstrar que a adição s}
p 201). Nosso próximo passo será demonstrar que a adição satisfaz as quatro primeiras propriedades em \ensuremath{\nu}: (p. 154) \ensuremath{\nu} = A1, A2, A3, A4 | z ?, M1, M2, M3, D, Ordenado, PBO
\prov{relações: parte\_de matematica}
\prov{proveniência: gentil/Gentil\_Lopes\_FUNDAMENTOS\_DOS\_NUMEROS\_pdf.pdf \textperiodcentered{} fato \textperiodcentered{} confiança 1.0 \textperiodcentered{} domínio matematica \textperiodcentered{} história 5 versões no jornal}

%CRISTAL {"arestas":[{"alvo":"matematica","confianca":0.95,"epistemico":"fato","origem":"paper","rel":"parte_de"}],"confianca":1.0,"contraexemplos":[],"descricao":"Sejam m, n e p naturais Então: (i) m < n e n < p ⇒m < p (transitividade) (ii) m < n ⇔m + p < n + p (monotonicidade) Prova: (i) Por hipótese, existem naturais r e t, tais que m + r = n e n + t = p Substituindo o valor de n da primeira equação na segunda, obtemos: (m + r) + t = p ⇒ m + (r + t) = p ⇒ m < p (ii) ( ⇒) m < n ⇒m + p < n + p Por hipótese existe um natural r tal que m + r = n. Sendo assim, n + p = (m + r) + p ⇒ (m + p) + r = n + p ⇒ m + p < n + p ( ⇐) m + p < n + p ⇒m < n Por hipótese ex","epistemico":"fato","exemplos":[],"id":"gentil_fundamentos_teorema_10_sejam_m_n_e_p_naturais_entao_i_m_n_e_n_p_m_p_tra","memoria":"semantica","meta":{"arquivo":"gentil/Lopes_Lopes_FUNDAMENTOS_DOS_NUMEROS_pdf.pdf","autor":"Gentil Lopes da Silva","ciencia":"teoria_dos_numeros","dominio":"matematica","nivel":4,"numero":"10","tipo_no":"paper","tipo_teorema":"teorema","titulo_legivel":"Fundamentos dos Números"},"origem":"gentil/Gentil_Lopes_FUNDAMENTOS_DOS_NUMEROS_pdf.pdf","palavras_chave":[],"sinonimos":[],"tipo":"conceito","titulo":"gentil fundamentos Teorema 10 — Sejam m, n e p naturais Então: (i) m < n e n < p ⇒m < p (tra"}
\section{gentil fundamentos Teorema 10 — Sejam m, n e p naturais Então: (i) m < n e n < p \ensuremath{\Rightarrow}m < p (tra}
Sejam m, n e p naturais Então: (i) m < n e n < p \ensuremath{\Rightarrow}m < p (transitividade) (ii) m < n \ensuremath{\Leftrightarrow}m + p < n + p (monotonicidade) Prova: (i) Por hipótese, existem naturais r e t, tais que m + r = n e n + t = p Substituindo o valor de n da primeira equação na segunda, obtemos: (m + r) + t = p \ensuremath{\Rightarrow} m + (r + t) = p \ensuremath{\Rightarrow} m < p (ii) ( \ensuremath{\Rightarrow}) m < n \ensuremath{\Rightarrow}m + p < n + p Por hipótese existe um natural r tal que m + r = n. Sendo assim, n + p = (m + r) + p \ensuremath{\Rightarrow} (m + p) + r = n + p \ensuremath{\Rightarrow} m + p < n + p ( \ensuremath{\Leftarrow}) m + p < n + p \ensuremath{\Rightarrow}m < n Por hipótese ex
\prov{relações: parte\_de matematica}
\prov{proveniência: gentil/Gentil\_Lopes\_FUNDAMENTOS\_DOS\_NUMEROS\_pdf.pdf \textperiodcentered{} fato \textperiodcentered{} confiança 1.0 \textperiodcentered{} domínio matematica \textperiodcentered{} história 4 versões no jornal}

%CRISTAL {"arestas":[{"alvo":"matematica","confianca":0.95,"epistemico":"fato","origem":"paper","rel":"parte_de"}],"confianca":1.0,"contraexemplos":[],"descricao":"Princ´ıpio de Indução Generalizado) Seja M um subconjunto de números naturais tal que k ∈M e σ(m) ∈M, para todo m ≥k em M. Então, M contém todos os números naturais n ≥k. Prova: Considere N= 0, 1, 2, 3,..., p ∪M, onde p é tal que σ(p) = k. Temos que 0 ∈N, suponhamos que n ∈N, então σ(n) ∈N; logo, pelo Princ´ıpio de Indução, temos N = N. ■ Exemplo: M = 4, 5, 6,..., satisfaz as hipóteses do","epistemico":"fato","exemplos":[],"id":"gentil_fundamentos_teorema_12_princ_ıpio_de_inducao_generalizado_seja_m_um_subco","memoria":"semantica","meta":{"arquivo":"gentil/Lopes_Lopes_FUNDAMENTOS_DOS_NUMEROS_pdf.pdf","autor":"Gentil Lopes da Silva","ciencia":"teoria_dos_numeros","dominio":"matematica","nivel":4,"numero":"12","tipo_no":"paper","tipo_teorema":"teorema","titulo_legivel":"Fundamentos dos Números"},"origem":"gentil/Gentil_Lopes_FUNDAMENTOS_DOS_NUMEROS_pdf.pdf","palavras_chave":[],"sinonimos":[],"tipo":"conceito","titulo":"gentil fundamentos Teorema 12 — Princ´ıpio de Indução Generalizado) Seja M um subconjunto de"}
\section{gentil fundamentos Teorema 12 — Princ'\ensuremath{\imath}pio de Indução Generalizado) Seja M um subconjunto de}
Princ'\ensuremath{\imath}pio de Indução Generalizado) Seja M um subconjunto de números naturais tal que k \ensuremath{\in}M e \ensuremath{\sigma}(m) \ensuremath{\in}M, para todo m \ensuremath{\ge}k em M. Então, M contém todos os números naturais n \ensuremath{\ge}k. Prova: Considere N= 0, 1, 2, 3,..., p \ensuremath{\cup}M, onde p é tal que \ensuremath{\sigma}(p) = k. Temos que 0 \ensuremath{\in}N, suponhamos que n \ensuremath{\in}N, então \ensuremath{\sigma}(n) \ensuremath{\in}N; logo, pelo Princ'\ensuremath{\imath}pio de Indução, temos N = N. \rule{1ex}{1ex} Exemplo: M = 4, 5, 6,..., satisfaz as hipóteses do
\prov{relações: parte\_de matematica}
\prov{proveniência: gentil/Gentil\_Lopes\_FUNDAMENTOS\_DOS\_NUMEROS\_pdf.pdf \textperiodcentered{} fato \textperiodcentered{} confiança 1.0 \textperiodcentered{} domínio matematica \textperiodcentered{} história 5 versões no jornal}

%CRISTAL {"arestas":[{"alvo":"matematica","confianca":0.95,"epistemico":"fato","origem":"paper","rel":"parte_de"}],"confianca":1.0,"contraexemplos":[],"descricao":"Sejam m, n, p e q naturais, então as seguintes propriedades se verificam: (i) m < n + 1 ⇒ m ≤n; (ii) m + p ≤n + p ⇔ m ≤n; (Compatibilidade com a Adição) (iii) m ≤n e p ≤q ⇒ m + p ≤n + q; (iv) m < n ⇒ m · p < n · p; (Compatibilidade com a Multiplicação) (v) m · p = n · p ⇒ m = n; (p ̸= 0) (Lei do corte) (vi) m < n ⇒ m + 1 ≤n; (vii) m ≤n e n ≤m ⇒ m = n. Prova: (i) Considere m < n + 1, então pela","epistemico":"fato","exemplos":[],"id":"gentil_fundamentos_teorema_13_sejam_m_n_p_e_q_naturais_entao_as_seguintes_propri","memoria":"semantica","meta":{"arquivo":"gentil/Lopes_Lopes_FUNDAMENTOS_DOS_NUMEROS_pdf.pdf","autor":"Gentil Lopes da Silva","ciencia":"teoria_dos_numeros","dominio":"matematica","nivel":4,"numero":"13","tipo_no":"paper","tipo_teorema":"teorema","titulo_legivel":"Fundamentos dos Números"},"origem":"gentil/Gentil_Lopes_FUNDAMENTOS_DOS_NUMEROS_pdf.pdf","palavras_chave":[],"sinonimos":[],"tipo":"conceito","titulo":"gentil fundamentos Teorema 13 — Sejam m, n, p e q naturais, então as seguintes propriedades"}
\section{gentil fundamentos Teorema 13 — Sejam m, n, p e q naturais, então as seguintes propriedades}
Sejam m, n, p e q naturais, então as seguintes propriedades se verificam: (i) m < n + 1 \ensuremath{\Rightarrow} m \ensuremath{\le}n; (ii) m + p \ensuremath{\le}n + p \ensuremath{\Leftrightarrow} m \ensuremath{\le}n; (Compatibilidade com a Adição) (iii) m \ensuremath{\le}n e p \ensuremath{\le}q \ensuremath{\Rightarrow} m + p \ensuremath{\le}n + q; (iv) m < n \ensuremath{\Rightarrow} m \textperiodcentered  p < n \textperiodcentered  p; (Compatibilidade com a Multiplicação) (v) m \textperiodcentered  p = n \textperiodcentered  p \ensuremath{\Rightarrow} m = n; (p \ensuremath{}= 0) (Lei do corte) (vi) m < n \ensuremath{\Rightarrow} m + 1 \ensuremath{\le}n; (vii) m \ensuremath{\le}n e n \ensuremath{\le}m \ensuremath{\Rightarrow} m = n. Prova: (i) Considere m < n + 1, então pela
\prov{relações: parte\_de matematica}
\prov{proveniência: gentil/Gentil\_Lopes\_FUNDAMENTOS\_DOS\_NUMEROS\_pdf.pdf \textperiodcentered{} fato \textperiodcentered{} confiança 1.0 \textperiodcentered{} domínio matematica \textperiodcentered{} história 4 versões no jornal}

%CRISTAL {"arestas":[{"alvo":"matematica","confianca":0.95,"epistemico":"fato","origem":"paper","rel":"parte_de"}],"confianca":1.0,"contraexemplos":[],"descricao":"Princ´ıpio da Boa Ordem) Todo subconjunto não vazio de N tem um menor elemento. Prova: Na prova que faremos o que vem entre o par de retas paralelas a seguir é apenas um exemplo para facilitar o entendimento da demons- tração, que de modo algum interfere no desenvolvimento lógico da prova. Pois bem, Seja A ⊂N = 0, 1, 2, 3,..., A ̸= ∅. A = 9, 7, 8, 15 Consideremos o conjunto B = n ∈N: n ≤m, ∀m ∈A. Isto é, os elementos de B são os naturais menores (ou =) que qualquer elemento de A.","epistemico":"fato","exemplos":[],"id":"gentil_fundamentos_teorema_14_princ_ıpio_da_boa_ordem_todo_subconjunto_nao_vazio","memoria":"semantica","meta":{"arquivo":"gentil/Lopes_Lopes_FUNDAMENTOS_DOS_NUMEROS_pdf.pdf","autor":"Gentil Lopes da Silva","ciencia":"teoria_dos_numeros","dominio":"matematica","nivel":4,"numero":"14","tipo_no":"paper","tipo_teorema":"teorema","titulo_legivel":"Fundamentos dos Números"},"origem":"gentil/Gentil_Lopes_FUNDAMENTOS_DOS_NUMEROS_pdf.pdf","palavras_chave":[],"sinonimos":[],"tipo":"conceito","titulo":"gentil fundamentos Teorema 14 — Princ´ıpio da Boa Ordem) Todo subconjunto não vazio de N tem"}
\section{gentil fundamentos Teorema 14 — Princ'\ensuremath{\imath}pio da Boa Ordem) Todo subconjunto não vazio de N tem}
Princ'\ensuremath{\imath}pio da Boa Ordem) Todo subconjunto não vazio de N tem um menor elemento. Prova: Na prova que faremos o que vem entre o par de retas paralelas a seguir é apenas um exemplo para facilitar o entendimento da demons- tração, que de modo algum interfere no desenvolvimento lógico da prova. Pois bem, Seja A \ensuremath{\subset}N = 0, 1, 2, 3,..., A \ensuremath{}= \ensuremath{\varnothing}. A = 9, 7, 8, 15 Consideremos o conjunto B = n \ensuremath{\in}N: n \ensuremath{\le}m, \ensuremath{\forall}m \ensuremath{\in}A. Isto é, os elementos de B são os naturais menores (ou =) que qualquer elemento de A.
\prov{relações: parte\_de matematica}
\prov{proveniência: gentil/Gentil\_Lopes\_FUNDAMENTOS\_DOS\_NUMEROS\_pdf.pdf \textperiodcentered{} fato \textperiodcentered{} confiança 1.0 \textperiodcentered{} domínio matematica \textperiodcentered{} história 5 versões no jornal}

%CRISTAL {"arestas":[{"alvo":"matematica","confianca":0.95,"epistemico":"fato","origem":"paper","rel":"parte_de"}],"confianca":1.0,"contraexemplos":[],"descricao":"Curva de Peano) Existe uma sobrejeção cont´ınua, ϕ: [ 0, 1 ] −→[ 0, 1 ]3 A demonstração deste resultado torna-se bem mais simples se os elemen- tos de [ 0, 1 ] forem substituidos por sequências binárias. (ver [19]) A bem da verdade, não tenho conhecimento se é poss´ıvel demonstrar este","epistemico":"fato","exemplos":[],"id":"gentil_fundamentos_teorema_15_curva_de_peano_existe_uma_sobrejecao_cont_ınua_φ_0","memoria":"semantica","meta":{"arquivo":"gentil/Lopes_Lopes_FUNDAMENTOS_DOS_NUMEROS_pdf.pdf","autor":"Gentil Lopes da Silva","ciencia":"teoria_dos_numeros","dominio":"matematica","nivel":4,"numero":"15","tipo_no":"paper","tipo_teorema":"teorema","titulo_legivel":"Fundamentos dos Números"},"origem":"gentil/Gentil_Lopes_FUNDAMENTOS_DOS_NUMEROS_pdf.pdf","palavras_chave":[],"sinonimos":[],"tipo":"conceito","titulo":"gentil fundamentos Teorema 15 — Curva de Peano) Existe uma sobrejeção cont´ınua, ϕ: [ 0, 1 ]"}
\section{gentil fundamentos Teorema 15 — Curva de Peano) Existe uma sobrejeção cont'\ensuremath{\imath}nua, \ensuremath{\phi}: [ 0, 1 ]}
Curva de Peano) Existe uma sobrejeção cont'\ensuremath{\imath}nua, \ensuremath{\phi}: [ 0, 1 ] \ensuremath{-}\ensuremath{\to}[ 0, 1 ]3 A demonstração deste resultado torna-se bem mais simples se os elemen- tos de [ 0, 1 ] forem substituidos por sequências binárias. (ver [19]) A bem da verdade, não tenho conhecimento se é poss'\ensuremath{\imath}vel demonstrar este
\prov{relações: parte\_de matematica}
\prov{proveniência: gentil/Gentil\_Lopes\_FUNDAMENTOS\_DOS\_NUMEROS\_pdf.pdf \textperiodcentered{} fato \textperiodcentered{} confiança 1.0 \textperiodcentered{} domínio matematica \textperiodcentered{} história 4 versões no jornal}

%CRISTAL {"arestas":[{"alvo":"gentil_fundamentos_capitulo_1_o_que_e_um_n_umero","confianca":1.0,"epistemico":"fato","origem":"gentil/Gentil_Lopes_FUNDAMENTOS_DOS_NUMEROS_pdf.pdf","rel":"prova"}],"confianca":1.0,"contraexemplos":[],"descricao":"Associatividade) Sejam α = (a, b), β = (c, d), e γ = (e, f), inteiros quaisquer. Vale a seguinte igualdade: A1 ) (α + β) + γ = α + (β + γ). 162\nProva: (α + β) + γ = \u0000(a, b) + (c, d) \u0001 + (e, f) = (a + c, b + d) + (e, f) = \u0000(a + c) + e, (b + d) + f \u0001 = \u0000a + (c + e), b + (d + f) \u0001 = (a, b) + \u0000c + e, d + f \u0001 = (a, b) + \u0000(c, d) + (e, f) \u0001 = α + (β + γ) ■ Apenas enfatizamos que da terceira para a quarta igualdade usamos a associatividade nos naturais. Nas outras igualdades utilizamos apenas a","epistemico":"fato","exemplos":[],"id":"gentil_fundamentos_teorema_17_associatividade_sejam_α_a_b_β_c_d_e_γ_e_f","memoria":"semantica","meta":{"arquivo":"gentil/Lopes_Lopes_FUNDAMENTOS_DOS_NUMEROS_pdf.pdf","autor":"Gentil Lopes da Silva","ciencia":"teoria_dos_numeros","dominio":"matematica","nivel":4,"numero":"17","tipo_no":"paper","tipo_teorema":"teorema","titulo_legivel":"Fundamentos dos Números"},"origem":"gentil/Gentil_Lopes_FUNDAMENTOS_DOS_NUMEROS_pdf.pdf","palavras_chave":[],"sinonimos":[],"tipo":"conceito","titulo":"gentil fundamentos Teorema 17 — Associatividade) Sejam α = (a, b), β = (c, d), e γ = (e, f),"}
\section{gentil fundamentos Teorema 17 — Associatividade) Sejam \ensuremath{\alpha} = (a, b), \ensuremath{\beta} = (c, d), e \ensuremath{\gamma} = (e, f),}
Associatividade) Sejam \ensuremath{\alpha} = (a, b), \ensuremath{\beta} = (c, d), e \ensuremath{\gamma} = (e, f), inteiros quaisquer. Vale a seguinte igualdade: A1 ) (\ensuremath{\alpha} + \ensuremath{\beta}) + \ensuremath{\gamma} = \ensuremath{\alpha} + (\ensuremath{\beta} + \ensuremath{\gamma}). 162
Prova: (\ensuremath{\alpha} + \ensuremath{\beta}) + \ensuremath{\gamma} = (a, b) + (c, d)  + (e, f) = (a + c, b + d) + (e, f) = (a + c) + e, (b + d) + f  = a + (c + e), b + (d + f)  = (a, b) + c + e, d + f  = (a, b) + (c, d) + (e, f)  = \ensuremath{\alpha} + (\ensuremath{\beta} + \ensuremath{\gamma}) \rule{1ex}{1ex} Apenas enfatizamos que da terceira para a quarta igualdade usamos a associatividade nos naturais. Nas outras igualdades utilizamos apenas a
\prov{relações: prova gentil\_fundamentos\_capitulo\_1\_o\_que\_e\_um\_n\_umero}
\prov{proveniência: gentil/Gentil\_Lopes\_FUNDAMENTOS\_DOS\_NUMEROS\_pdf.pdf \textperiodcentered{} fato \textperiodcentered{} confiança 1.0 \textperiodcentered{} domínio matematica \textperiodcentered{} história 5 versões no jornal}

%CRISTAL {"arestas":[{"alvo":"matematica","confianca":0.95,"epistemico":"fato","origem":"paper","rel":"parte_de"}],"confianca":1.0,"contraexemplos":[],"descricao":"Elemento neutro) Seja n um natural arbitrariamente fixado. O inteiro θ = (n, n) é o elemento neutro para a adição em ¯Z. (eq. (5.4), p. 158) Prova: Seja α = (a, b) um inteiro arbitrário, então α + θ = (a, b) + (n, n) = (a + n, b + n) Por outro lado, vamos mostrar que (a + n, b + n) = (a, b) Com efeito, pela","epistemico":"fato","exemplos":[],"id":"gentil_fundamentos_teorema_18_elemento_neutro_seja_n_um_natural_arbitrariamente_","memoria":"semantica","meta":{"arquivo":"gentil/Lopes_Lopes_FUNDAMENTOS_DOS_NUMEROS_pdf.pdf","autor":"Gentil Lopes da Silva","ciencia":"teoria_dos_numeros","dominio":"matematica","nivel":4,"numero":"18","tipo_no":"paper","tipo_teorema":"teorema","titulo_legivel":"Fundamentos dos Números"},"origem":"gentil/Gentil_Lopes_FUNDAMENTOS_DOS_NUMEROS_pdf.pdf","palavras_chave":[],"sinonimos":[],"tipo":"conceito","titulo":"gentil fundamentos Teorema 18 — Elemento neutro) Seja n um natural arbitrariamente fixado. O"}
\section{gentil fundamentos Teorema 18 — Elemento neutro) Seja n um natural arbitrariamente fixado. O}
Elemento neutro) Seja n um natural arbitrariamente fixado. O inteiro \ensuremath{\theta} = (n, n) é o elemento neutro para a adição em \ensuremath{\bar{\,}}Z. (eq. (5.4), p. 158) Prova: Seja \ensuremath{\alpha} = (a, b) um inteiro arbitrário, então \ensuremath{\alpha} + \ensuremath{\theta} = (a, b) + (n, n) = (a + n, b + n) Por outro lado, vamos mostrar que (a + n, b + n) = (a, b) Com efeito, pela
\prov{relações: parte\_de matematica}
\prov{proveniência: gentil/Gentil\_Lopes\_FUNDAMENTOS\_DOS\_NUMEROS\_pdf.pdf \textperiodcentered{} fato \textperiodcentered{} confiança 1.0 \textperiodcentered{} domínio matematica \textperiodcentered{} história 4 versões no jornal}

%CRISTAL {"arestas":[{"alvo":"matematica","confianca":0.95,"epistemico":"fato","origem":"paper","rel":"parte_de"}],"confianca":1.0,"contraexemplos":[],"descricao":"Comutativa) Sejam α = (a, b) e β = (c, d) inteiros quais- quer. Vale a seguinte igualdade: A3 ) α + β = β + α Prova: Temos α + β = (a, b) + (c, d) = (a + c, b + d) = (c + a, d + b) = (c, d) + (a, b) = β + α ■","epistemico":"fato","exemplos":[],"id":"gentil_fundamentos_teorema_19_comutativa_sejam_α_a_b_e_β_c_d_inteiros_quais-_qu","memoria":"semantica","meta":{"arquivo":"gentil/Lopes_Lopes_FUNDAMENTOS_DOS_NUMEROS_pdf.pdf","autor":"Gentil Lopes da Silva","ciencia":"teoria_dos_numeros","dominio":"matematica","nivel":4,"numero":"19","tipo_no":"paper","tipo_teorema":"teorema","titulo_legivel":"Fundamentos dos Números"},"origem":"gentil/Gentil_Lopes_FUNDAMENTOS_DOS_NUMEROS_pdf.pdf","palavras_chave":[],"sinonimos":[],"tipo":"conceito","titulo":"gentil fundamentos Teorema 19 — Comutativa) Sejam α = (a, b) e β = (c, d) inteiros quais- qu"}
\section{gentil fundamentos Teorema 19 — Comutativa) Sejam \ensuremath{\alpha} = (a, b) e \ensuremath{\beta} = (c, d) inteiros quais- qu}
Comutativa) Sejam \ensuremath{\alpha} = (a, b) e \ensuremath{\beta} = (c, d) inteiros quais- quer. Vale a seguinte igualdade: A3 ) \ensuremath{\alpha} + \ensuremath{\beta} = \ensuremath{\beta} + \ensuremath{\alpha} Prova: Temos \ensuremath{\alpha} + \ensuremath{\beta} = (a, b) + (c, d) = (a + c, b + d) = (c + a, d + b) = (c, d) + (a, b) = \ensuremath{\beta} + \ensuremath{\alpha} \rule{1ex}{1ex}
\prov{relações: parte\_de matematica}
\prov{proveniência: gentil/Gentil\_Lopes\_FUNDAMENTOS\_DOS\_NUMEROS\_pdf.pdf \textperiodcentered{} fato \textperiodcentered{} confiança 1.0 \textperiodcentered{} domínio matematica \textperiodcentered{} história 4 versões no jornal}

%CRISTAL {"arestas":[{"alvo":"matematica","confianca":0.95,"epistemico":"fato","origem":"paper","rel":"parte_de"}],"confianca":1.0,"contraexemplos":[],"descricao":", resulta: ¯a = ¯b ■ Teorema 4. Se as classes de equivalência ¯a e ¯b são diferentes, então são disjuntas. Prova: ´E o","epistemico":"fato","exemplos":[],"id":"gentil_fundamentos_teorema_1_resulta_a_b_teorema_4_se_as_classes_de_equivalenci","memoria":"semantica","meta":{"arquivo":"gentil/Lopes_Lopes_FUNDAMENTOS_DOS_NUMEROS_pdf.pdf","autor":"Gentil Lopes da Silva","ciencia":"teoria_dos_numeros","dominio":"matematica","nivel":4,"numero":"1","tipo_no":"paper","tipo_teorema":"teorema","titulo_legivel":"Fundamentos dos Números"},"origem":"gentil/Gentil_Lopes_FUNDAMENTOS_DOS_NUMEROS_pdf.pdf","palavras_chave":[],"sinonimos":[],"tipo":"conceito","titulo":"gentil fundamentos Teorema 1 — , resulta: ¯a = ¯b ■ Teorema 4. Se as classes de equivalênci"}
\section{gentil fundamentos Teorema 1 — , resulta: \ensuremath{\bar{\,}}a = \ensuremath{\bar{\,}}b \rule{1ex}{1ex} Teorema 4. Se as classes de equivalênci}
, resulta: \ensuremath{\bar{\,}}a = \ensuremath{\bar{\,}}b \rule{1ex}{1ex} Teorema 4. Se as classes de equivalência \ensuremath{\bar{\,}}a e \ensuremath{\bar{\,}}b são diferentes, então são disjuntas. Prova: 'E o
\prov{relações: parte\_de matematica}
\prov{proveniência: gentil/Gentil\_Lopes\_FUNDAMENTOS\_DOS\_NUMEROS\_pdf.pdf \textperiodcentered{} fato \textperiodcentered{} confiança 1.0 \textperiodcentered{} domínio matematica \textperiodcentered{} história 4 versões no jornal}

%CRISTAL {"arestas":[{"alvo":"matematica","confianca":0.95,"epistemico":"fato","origem":"paper","rel":"parte_de"}],"confianca":1.0,"contraexemplos":[],"descricao":"■ Teorema 3 Se as classes de equivalência ¯a e ¯b não são disjuntas, então são iguais. Prova: (⇒) Separando em hipótese e tese, temos: H: ¯a ∩¯b ̸= ∅ T: ¯a = ¯b Com efeito, se ¯a e ¯b não são disjuntas, então existe c ∈A tal que c ∈¯a e c ∈¯b, isto é, tal que c ∼a e c ∼b; logo, por ser ∼simétrica, resulta, a ∼c. Pela transitividade de ∼valem as implicações: a ∼c e c ∼b ⇒ a ∼b Sendo assim, pelo","epistemico":"fato","exemplos":[],"id":"gentil_fundamentos_teorema_1_teorema_3_se_as_classes_de_equivalencia_a_e_b_nao_s","memoria":"semantica","meta":{"arquivo":"gentil/Lopes_Lopes_FUNDAMENTOS_DOS_NUMEROS_pdf.pdf","autor":"Gentil Lopes da Silva","ciencia":"teoria_dos_numeros","dominio":"matematica","nivel":4,"numero":"1","tipo_no":"paper","tipo_teorema":"teorema","titulo_legivel":"Fundamentos dos Números"},"origem":"gentil/Gentil_Lopes_FUNDAMENTOS_DOS_NUMEROS_pdf.pdf","palavras_chave":[],"sinonimos":[],"tipo":"conceito","titulo":"gentil fundamentos Teorema 1 — ■ Teorema 3 Se as classes de equivalência ¯a e ¯b não são di"}
\section{gentil fundamentos Teorema 1 — \rule{1ex}{1ex} Teorema 3 Se as classes de equivalência \ensuremath{\bar{\,}}a e \ensuremath{\bar{\,}}b não são di}
\rule{1ex}{1ex} Teorema 3 Se as classes de equivalência \ensuremath{\bar{\,}}a e \ensuremath{\bar{\,}}b não são disjuntas, então são iguais. Prova: (\ensuremath{\Rightarrow}) Separando em hipótese e tese, temos: H: \ensuremath{\bar{\,}}a \ensuremath{\cap}\ensuremath{\bar{\,}}b \ensuremath{}= \ensuremath{\varnothing} T: \ensuremath{\bar{\,}}a = \ensuremath{\bar{\,}}b Com efeito, se \ensuremath{\bar{\,}}a e \ensuremath{\bar{\,}}b não são disjuntas, então existe c \ensuremath{\in}A tal que c \ensuremath{\in}\ensuremath{\bar{\,}}a e c \ensuremath{\in}\ensuremath{\bar{\,}}b, isto é, tal que c \ensuremath{\sim}a e c \ensuremath{\sim}b; logo, por ser \ensuremath{\sim}simétrica, resulta, a \ensuremath{\sim}c. Pela transitividade de \ensuremath{\sim}valem as implicações: a \ensuremath{\sim}c e c \ensuremath{\sim}b \ensuremath{\Rightarrow} a \ensuremath{\sim}b Sendo assim, pelo
\prov{relações: parte\_de matematica}
\prov{proveniência: gentil/Gentil\_Lopes\_FUNDAMENTOS\_DOS\_NUMEROS\_pdf.pdf \textperiodcentered{} fato \textperiodcentered{} confiança 1.0 \textperiodcentered{} domínio matematica \textperiodcentered{} história 4 versões no jornal}

%CRISTAL {"arestas":[{"alvo":"matematica","confianca":0.95,"epistemico":"fato","origem":"paper","rel":"parte_de"}],"confianca":1.0,"contraexemplos":[],"descricao":"Elemento oposto aditivo) Seja α um número inteiro qual- quer, então existe um e apenas um α′ inteiro tal que α + α′ = ¯0 Prova: Existência: Seja α = (a, b), vamos mostrar que seu oposto aditivo vale α′ = (b, a). De fato, α + α′ = (a, b) + (b, a) = (a + b, b + a) = (a + b, a + b) = (n, n) = ¯0 Observe que (a + b, a + b) = (n, n) De fato, pela","epistemico":"fato","exemplos":[],"id":"gentil_fundamentos_teorema_20_elemento_oposto_aditivo_seja_α_um_numero_inteiro_q","memoria":"semantica","meta":{"arquivo":"gentil/Lopes_Lopes_FUNDAMENTOS_DOS_NUMEROS_pdf.pdf","autor":"Gentil Lopes da Silva","ciencia":"teoria_dos_numeros","dominio":"matematica","nivel":4,"numero":"20","tipo_no":"paper","tipo_teorema":"teorema","titulo_legivel":"Fundamentos dos Números"},"origem":"gentil/Gentil_Lopes_FUNDAMENTOS_DOS_NUMEROS_pdf.pdf","palavras_chave":[],"sinonimos":[],"tipo":"conceito","titulo":"gentil fundamentos Teorema 20 — Elemento oposto aditivo) Seja α um número inteiro qual- quer"}
\section{gentil fundamentos Teorema 20 — Elemento oposto aditivo) Seja \ensuremath{\alpha} um número inteiro qual- quer}
Elemento oposto aditivo) Seja \ensuremath{\alpha} um número inteiro qual- quer, então existe um e apenas um \ensuremath{\alpha}\ensuremath{'} inteiro tal que \ensuremath{\alpha} + \ensuremath{\alpha}\ensuremath{'} = \ensuremath{\bar{\,}}0 Prova: Existência: Seja \ensuremath{\alpha} = (a, b), vamos mostrar que seu oposto aditivo vale \ensuremath{\alpha}\ensuremath{'} = (b, a). De fato, \ensuremath{\alpha} + \ensuremath{\alpha}\ensuremath{'} = (a, b) + (b, a) = (a + b, b + a) = (a + b, a + b) = (n, n) = \ensuremath{\bar{\,}}0 Observe que (a + b, a + b) = (n, n) De fato, pela
\prov{relações: parte\_de matematica}
\prov{proveniência: gentil/Gentil\_Lopes\_FUNDAMENTOS\_DOS\_NUMEROS\_pdf.pdf \textperiodcentered{} fato \textperiodcentered{} confiança 1.0 \textperiodcentered{} domínio matematica \textperiodcentered{} história 4 versões no jornal}

%CRISTAL {"arestas":[{"alvo":"gentil_fundamentos_capitulo_1_o_que_e_um_n_umero","confianca":1.0,"epistemico":"fato","origem":"gentil/Gentil_Lopes_FUNDAMENTOS_DOS_NUMEROS_pdf.pdf","rel":"prova"}],"confianca":1.0,"contraexemplos":[],"descricao":"p 164), temos −α = (b, a), logo −(−α) = − \u0000(b, a) \u0001 = (a, b) = α ■ Prova: (ii) Sejam α = (a, b) e β = (c, d) dois inteiros. Pelo","epistemico":"fato","exemplos":[],"id":"gentil_fundamentos_teorema_20_p_164_temos_α_b_a_logo_α_b_a_a_b","memoria":"semantica","meta":{"arquivo":"gentil/Lopes_Lopes_FUNDAMENTOS_DOS_NUMEROS_pdf.pdf","autor":"Gentil Lopes da Silva","ciencia":"teoria_dos_numeros","dominio":"matematica","nivel":4,"numero":"20","tipo_no":"paper","tipo_teorema":"teorema","titulo_legivel":"Fundamentos dos Números"},"origem":"gentil/Gentil_Lopes_FUNDAMENTOS_DOS_NUMEROS_pdf.pdf","palavras_chave":[],"sinonimos":[],"tipo":"conceito","titulo":"gentil fundamentos Teorema 20 — p 164), temos −α = (b, a), logo −(−α) = − \u0000(b, a) \u0001 = (a, b)"}
\section{gentil fundamentos Teorema 20 — p 164), temos \ensuremath{-}\ensuremath{\alpha} = (b, a), logo \ensuremath{-}(\ensuremath{-}\ensuremath{\alpha}) = \ensuremath{-} (b, a)  = (a, b)}
p 164), temos \ensuremath{-}\ensuremath{\alpha} = (b, a), logo \ensuremath{-}(\ensuremath{-}\ensuremath{\alpha}) = \ensuremath{-} (b, a)  = (a, b) = \ensuremath{\alpha} \rule{1ex}{1ex} Prova: (ii) Sejam \ensuremath{\alpha} = (a, b) e \ensuremath{\beta} = (c, d) dois inteiros. Pelo
\prov{relações: prova gentil\_fundamentos\_capitulo\_1\_o\_que\_e\_um\_n\_umero}
\prov{proveniência: gentil/Gentil\_Lopes\_FUNDAMENTOS\_DOS\_NUMEROS\_pdf.pdf \textperiodcentered{} fato \textperiodcentered{} confiança 1.0 \textperiodcentered{} domínio matematica \textperiodcentered{} história 5 versões no jornal}

%CRISTAL {"arestas":[{"alvo":"matematica","confianca":0.95,"epistemico":"fato","origem":"paper","rel":"parte_de"}],"confianca":1.0,"contraexemplos":[],"descricao":"Lei do corte) Sejam α, β, e γ, inteiros quaisquer. Vale a seguinte implicação: α + γ = β + γ ⇒ α = β Prova: Temos α + γ = β + γ ⇒ (α + γ) + (−γ) = (β + γ) + (−γ) aplicando associatividade, temos (α + γ) + (−γ) = (β + γ) + (−γ) ⇒ α + \u0000γ + (−γ) \u0001 = β + \u0000γ + (−γ) \u0001 ⇒ α + ¯0 = β + ¯0 ⇒ α = β ■","epistemico":"fato","exemplos":[],"id":"gentil_fundamentos_teorema_21_lei_do_corte_sejam_α_β_e_γ_inteiros_quaisquer_vale","memoria":"semantica","meta":{"arquivo":"gentil/Lopes_Lopes_FUNDAMENTOS_DOS_NUMEROS_pdf.pdf","autor":"Gentil Lopes da Silva","ciencia":"teoria_dos_numeros","dominio":"matematica","nivel":4,"numero":"21","tipo_no":"paper","tipo_teorema":"teorema","titulo_legivel":"Fundamentos dos Números"},"origem":"gentil/Gentil_Lopes_FUNDAMENTOS_DOS_NUMEROS_pdf.pdf","palavras_chave":[],"sinonimos":[],"tipo":"conceito","titulo":"gentil fundamentos Teorema 21 — Lei do corte) Sejam α, β, e γ, inteiros quaisquer. Vale a se"}
\section{gentil fundamentos Teorema 21 — Lei do corte) Sejam \ensuremath{\alpha}, \ensuremath{\beta}, e \ensuremath{\gamma}, inteiros quaisquer. Vale a se}
Lei do corte) Sejam \ensuremath{\alpha}, \ensuremath{\beta}, e \ensuremath{\gamma}, inteiros quaisquer. Vale a seguinte implicação: \ensuremath{\alpha} + \ensuremath{\gamma} = \ensuremath{\beta} + \ensuremath{\gamma} \ensuremath{\Rightarrow} \ensuremath{\alpha} = \ensuremath{\beta} Prova: Temos \ensuremath{\alpha} + \ensuremath{\gamma} = \ensuremath{\beta} + \ensuremath{\gamma} \ensuremath{\Rightarrow} (\ensuremath{\alpha} + \ensuremath{\gamma}) + (\ensuremath{-}\ensuremath{\gamma}) = (\ensuremath{\beta} + \ensuremath{\gamma}) + (\ensuremath{-}\ensuremath{\gamma}) aplicando associatividade, temos (\ensuremath{\alpha} + \ensuremath{\gamma}) + (\ensuremath{-}\ensuremath{\gamma}) = (\ensuremath{\beta} + \ensuremath{\gamma}) + (\ensuremath{-}\ensuremath{\gamma}) \ensuremath{\Rightarrow} \ensuremath{\alpha} + \ensuremath{\gamma} + (\ensuremath{-}\ensuremath{\gamma})  = \ensuremath{\beta} + \ensuremath{\gamma} + (\ensuremath{-}\ensuremath{\gamma})  \ensuremath{\Rightarrow} \ensuremath{\alpha} + \ensuremath{\bar{\,}}0 = \ensuremath{\beta} + \ensuremath{\bar{\,}}0 \ensuremath{\Rightarrow} \ensuremath{\alpha} = \ensuremath{\beta} \rule{1ex}{1ex}
\prov{relações: parte\_de matematica}
\prov{proveniência: gentil/Gentil\_Lopes\_FUNDAMENTOS\_DOS\_NUMEROS\_pdf.pdf \textperiodcentered{} fato \textperiodcentered{} confiança 1.0 \textperiodcentered{} domínio matematica \textperiodcentered{} história 4 versões no jornal}

%CRISTAL {"arestas":[{"alvo":"gentil_fundamentos_capitulo_1_o_que_e_um_n_umero","confianca":1.0,"epistemico":"fato","origem":"gentil/Gentil_Lopes_FUNDAMENTOS_DOS_NUMEROS_pdf.pdf","rel":"prova"}],"confianca":1.0,"contraexemplos":[],"descricao":"Associatividade) Sejam α = (a, b), β = (c, d), e γ = (e, f), inteiros quaisquer. Vale a seguinte igualdade: M1 ) (α · β) · γ = α · (β · γ) Prova: (α · β) · γ = \u0000(a, b) · (c, d) \u0001 · (e, f) = (ac + bd, ad + bc) · (e, f) = \u0000(ac + bd)e + (ad + bc)f, (ac + bd)f + (ad + bc)e \u0001 = \u0000(ac)e + (bd)e + (ad)f + (bc)f, (ac)f + (bd)f + (ad)e + (bc)e \u0001 = \u0000a(ce) + b(de) + a(df) + b(cf), a(cf) + b(df) + a(de) + b(ce) \u0001 = \u0000a(ce + df) + b(de + cf), a(cf + de) + b(df + ce) \u0001 = \u0000a(ce + df) + b(cf + de), a(cf + de) + b","epistemico":"fato","exemplos":[],"id":"gentil_fundamentos_teorema_22_associatividade_sejam_α_a_b_β_c_d_e_γ_e_f","memoria":"semantica","meta":{"arquivo":"gentil/Lopes_Lopes_FUNDAMENTOS_DOS_NUMEROS_pdf.pdf","autor":"Gentil Lopes da Silva","ciencia":"teoria_dos_numeros","dominio":"matematica","nivel":4,"numero":"22","tipo_no":"paper","tipo_teorema":"teorema","titulo_legivel":"Fundamentos dos Números"},"origem":"gentil/Gentil_Lopes_FUNDAMENTOS_DOS_NUMEROS_pdf.pdf","palavras_chave":[],"sinonimos":[],"tipo":"conceito","titulo":"gentil fundamentos Teorema 22 — Associatividade) Sejam α = (a, b), β = (c, d), e γ = (e, f),"}
\section{gentil fundamentos Teorema 22 — Associatividade) Sejam \ensuremath{\alpha} = (a, b), \ensuremath{\beta} = (c, d), e \ensuremath{\gamma} = (e, f),}
Associatividade) Sejam \ensuremath{\alpha} = (a, b), \ensuremath{\beta} = (c, d), e \ensuremath{\gamma} = (e, f), inteiros quaisquer. Vale a seguinte igualdade: M1 ) (\ensuremath{\alpha} \textperiodcentered  \ensuremath{\beta}) \textperiodcentered  \ensuremath{\gamma} = \ensuremath{\alpha} \textperiodcentered  (\ensuremath{\beta} \textperiodcentered  \ensuremath{\gamma}) Prova: (\ensuremath{\alpha} \textperiodcentered  \ensuremath{\beta}) \textperiodcentered  \ensuremath{\gamma} = (a, b) \textperiodcentered  (c, d)  \textperiodcentered  (e, f) = (ac + bd, ad + bc) \textperiodcentered  (e, f) = (ac + bd)e + (ad + bc)f, (ac + bd)f + (ad + bc)e  = (ac)e + (bd)e + (ad)f + (bc)f, (ac)f + (bd)f + (ad)e + (bc)e  = a(ce) + b(de) + a(df) + b(cf), a(cf) + b(df) + a(de) + b(ce)  = a(ce + df) + b(de + cf), a(cf + de) + b(df + ce)  = a(ce + df) + b(cf + de), a(cf + de) + b
\prov{relações: prova gentil\_fundamentos\_capitulo\_1\_o\_que\_e\_um\_n\_umero}
\prov{proveniência: gentil/Gentil\_Lopes\_FUNDAMENTOS\_DOS\_NUMEROS\_pdf.pdf \textperiodcentered{} fato \textperiodcentered{} confiança 1.0 \textperiodcentered{} domínio matematica \textperiodcentered{} história 5 versões no jornal}

%CRISTAL {"arestas":[{"alvo":"matematica","confianca":0.95,"epistemico":"fato","origem":"paper","rel":"parte_de"}],"confianca":1.0,"contraexemplos":[],"descricao":"Elemento neutro) Seja n um natural arbitrariamente fixado. O inteiro Θ = (n + 1, n) é o elemento neutro para a multiplicação em ¯Z. (eq. 5.5, p. 158) Prova: Seja α = (a, b) um inteiro arbitrário, então α · Θ = (a, b) · (n + 1, n) = \u0000a(n + 1) + bn, an + b(n + 1) \u0001 Por outro lado, vamos mostrar que \u0000a(n + 1) + bn, an + b(n + 1) \u0001 = (a, b) (5.8) Com efeito, pela","epistemico":"fato","exemplos":[],"id":"gentil_fundamentos_teorema_23_elemento_neutro_seja_n_um_natural_arbitrariamente_","memoria":"semantica","meta":{"arquivo":"gentil/Lopes_Lopes_FUNDAMENTOS_DOS_NUMEROS_pdf.pdf","autor":"Gentil Lopes da Silva","ciencia":"teoria_dos_numeros","dominio":"matematica","nivel":4,"numero":"23","tipo_no":"paper","tipo_teorema":"teorema","titulo_legivel":"Fundamentos dos Números"},"origem":"gentil/Gentil_Lopes_FUNDAMENTOS_DOS_NUMEROS_pdf.pdf","palavras_chave":[],"sinonimos":[],"tipo":"conceito","titulo":"gentil fundamentos Teorema 23 — Elemento neutro) Seja n um natural arbitrariamente fixado. O"}
\section{gentil fundamentos Teorema 23 — Elemento neutro) Seja n um natural arbitrariamente fixado. O}
Elemento neutro) Seja n um natural arbitrariamente fixado. O inteiro \ensuremath{\Theta} = (n + 1, n) é o elemento neutro para a multiplicação em \ensuremath{\bar{\,}}Z. (eq. 5.5, p. 158) Prova: Seja \ensuremath{\alpha} = (a, b) um inteiro arbitrário, então \ensuremath{\alpha} \textperiodcentered  \ensuremath{\Theta} = (a, b) \textperiodcentered  (n + 1, n) = a(n + 1) + bn, an + b(n + 1)  Por outro lado, vamos mostrar que a(n + 1) + bn, an + b(n + 1)  = (a, b) (5.8) Com efeito, pela
\prov{relações: parte\_de matematica}
\prov{proveniência: gentil/Gentil\_Lopes\_FUNDAMENTOS\_DOS\_NUMEROS\_pdf.pdf \textperiodcentered{} fato \textperiodcentered{} confiança 1.0 \textperiodcentered{} domínio matematica \textperiodcentered{} história 4 versões no jornal}

%CRISTAL {"arestas":[{"alvo":"matematica","confianca":0.95,"epistemico":"fato","origem":"paper","rel":"parte_de"}],"confianca":1.0,"contraexemplos":[],"descricao":"Comutativa) Sejam α = (a, b) e β = (c, d) inteiros quais- quer. Vale a seguinte igualdade: A3 ) α · β = β · α Prova: Temos α · β = (a, b) · (c, d) = (ac + bd, ad + bc) = (ca + db, cb + da) = (c, d) · (a, b) = β · α ■ 167","epistemico":"fato","exemplos":[],"id":"gentil_fundamentos_teorema_24_comutativa_sejam_α_a_b_e_β_c_d_inteiros_quais-_qu","memoria":"semantica","meta":{"arquivo":"gentil/Lopes_Lopes_FUNDAMENTOS_DOS_NUMEROS_pdf.pdf","autor":"Gentil Lopes da Silva","ciencia":"teoria_dos_numeros","dominio":"matematica","nivel":4,"numero":"24","tipo_no":"paper","tipo_teorema":"teorema","titulo_legivel":"Fundamentos dos Números"},"origem":"gentil/Gentil_Lopes_FUNDAMENTOS_DOS_NUMEROS_pdf.pdf","palavras_chave":[],"sinonimos":[],"tipo":"conceito","titulo":"gentil fundamentos Teorema 24 — Comutativa) Sejam α = (a, b) e β = (c, d) inteiros quais- qu"}
\section{gentil fundamentos Teorema 24 — Comutativa) Sejam \ensuremath{\alpha} = (a, b) e \ensuremath{\beta} = (c, d) inteiros quais- qu}
Comutativa) Sejam \ensuremath{\alpha} = (a, b) e \ensuremath{\beta} = (c, d) inteiros quais- quer. Vale a seguinte igualdade: A3 ) \ensuremath{\alpha} \textperiodcentered  \ensuremath{\beta} = \ensuremath{\beta} \textperiodcentered  \ensuremath{\alpha} Prova: Temos \ensuremath{\alpha} \textperiodcentered  \ensuremath{\beta} = (a, b) \textperiodcentered  (c, d) = (ac + bd, ad + bc) = (ca + db, cb + da) = (c, d) \textperiodcentered  (a, b) = \ensuremath{\beta} \textperiodcentered  \ensuremath{\alpha} \rule{1ex}{1ex} 167
\prov{relações: parte\_de matematica}
\prov{proveniência: gentil/Gentil\_Lopes\_FUNDAMENTOS\_DOS\_NUMEROS\_pdf.pdf \textperiodcentered{} fato \textperiodcentered{} confiança 1.0 \textperiodcentered{} domínio matematica \textperiodcentered{} história 4 versões no jornal}

%CRISTAL {"arestas":[{"alvo":"matematica","confianca":0.95,"epistemico":"fato","origem":"paper","rel":"parte_de"}],"confianca":1.0,"contraexemplos":[],"descricao":"As seguintes implicações são válidas: a ∼b ⇒ a ∈¯b ⇒ b ∈¯a ⇒ ¯a = ¯b ⇒ a ∼b ( I ) ( II ) ( III) ( IV ) ( I ) Prova: ( I ) ⇒( II ): Decorre da própria","epistemico":"fato","exemplos":[],"id":"gentil_fundamentos_teorema_2_as_seguintes_implicacoes_sao_validas_a_b_a_b_b_a","memoria":"semantica","meta":{"arquivo":"gentil/Lopes_Lopes_FUNDAMENTOS_DOS_NUMEROS_pdf.pdf","autor":"Gentil Lopes da Silva","ciencia":"teoria_dos_numeros","dominio":"matematica","nivel":4,"numero":"2","tipo_no":"paper","tipo_teorema":"teorema","titulo_legivel":"Fundamentos dos Números"},"origem":"gentil/Gentil_Lopes_FUNDAMENTOS_DOS_NUMEROS_pdf.pdf","palavras_chave":[],"sinonimos":[],"tipo":"conceito","titulo":"gentil fundamentos Teorema 2 — As seguintes implicações são válidas: a ∼b ⇒ a ∈¯b ⇒ b ∈¯a ⇒"}
\section{gentil fundamentos Teorema 2 — As seguintes implicações são válidas: a \ensuremath{\sim}b \ensuremath{\Rightarrow} a \ensuremath{\in}\ensuremath{\bar{\,}}b \ensuremath{\Rightarrow} b \ensuremath{\in}\ensuremath{\bar{\,}}a \ensuremath{\Rightarrow}}
As seguintes implicações são válidas: a \ensuremath{\sim}b \ensuremath{\Rightarrow} a \ensuremath{\in}\ensuremath{\bar{\,}}b \ensuremath{\Rightarrow} b \ensuremath{\in}\ensuremath{\bar{\,}}a \ensuremath{\Rightarrow} \ensuremath{\bar{\,}}a = \ensuremath{\bar{\,}}b \ensuremath{\Rightarrow} a \ensuremath{\sim}b ( I ) ( II ) ( III) ( IV ) ( I ) Prova: ( I ) \ensuremath{\Rightarrow}( II ): Decorre da própria
\prov{relações: parte\_de matematica}
\prov{proveniência: gentil/Gentil\_Lopes\_FUNDAMENTOS\_DOS\_NUMEROS\_pdf.pdf \textperiodcentered{} fato \textperiodcentered{} confiança 1.0 \textperiodcentered{} domínio matematica \textperiodcentered{} história 4 versões no jornal}

%CRISTAL {"arestas":[],"confianca":1.0,"contraexemplos":[],"descricao":"p 69) temos (a + n, b + n) ∼(a, b) ⇒ (a + n, b + n) = (a, b) Portanto, α + θ = α. De modo análogo demonstra-se que θ é neutro pela esquerda. ■ Vamos utilizar a seguinte notação para o elemento neutro da adição: (n, n) = ¯0 163","epistemico":"fato","exemplos":[],"id":"gentil_fundamentos_teorema_2_p_69_temos_a_n_b_n_a_b_a_n_b_n_a_b","memoria":"semantica","meta":{"arquivo":"gentil/Lopes_Lopes_FUNDAMENTOS_DOS_NUMEROS_pdf.pdf","autor":"Gentil Lopes da Silva","ciencia":"teoria_dos_numeros","dominio":"matematica","nivel":4,"numero":"2","tipo_no":"paper","tipo_teorema":"teorema","titulo_legivel":"Fundamentos dos Números"},"origem":"gentil/Gentil_Lopes_FUNDAMENTOS_DOS_NUMEROS_pdf.pdf","palavras_chave":[],"sinonimos":[],"tipo":"conceito","titulo":"gentil fundamentos Teorema 2 — p 69) temos (a + n, b + n) ∼(a, b) ⇒ (a + n, b + n) = (a, b)"}
\section{gentil fundamentos Teorema 2 — p 69) temos (a + n, b + n) \ensuremath{\sim}(a, b) \ensuremath{\Rightarrow} (a + n, b + n) = (a, b)}
p 69) temos (a + n, b + n) \ensuremath{\sim}(a, b) \ensuremath{\Rightarrow} (a + n, b + n) = (a, b) Portanto, \ensuremath{\alpha} + \ensuremath{\theta} = \ensuremath{\alpha}. De modo análogo demonstra-se que \ensuremath{\theta} é neutro pela esquerda. \rule{1ex}{1ex} Vamos utilizar a seguinte notação para o elemento neutro da adição: (n, n) = \ensuremath{\bar{\,}}0 163
\prov{proveniência: gentil/Gentil\_Lopes\_FUNDAMENTOS\_DOS\_NUMEROS\_pdf.pdf \textperiodcentered{} fato \textperiodcentered{} confiança 1.0 \textperiodcentered{} domínio matematica \textperiodcentered{} história 5 versões no jornal}

%CRISTAL {"arestas":[],"confianca":1.0,"contraexemplos":[],"descricao":"p 69) temos \u0000a(n + 1) + bn, an + b(n + 1) \u0001 ∼(a, b) ⇒ (5.8) Portanto, α · Θ = α. De modo análogo demonstra-se que Θ é neutro pela esquerda. ■ Vamos utilizar a seguinte notação para o elemento neutro da multiplicação: (n + 1, n) = ¯1","epistemico":"fato","exemplos":[],"id":"gentil_fundamentos_teorema_2_p_69_temos_an_1_bn_an_bn_1_a_b_58","memoria":"semantica","meta":{"arquivo":"gentil/Lopes_Lopes_FUNDAMENTOS_DOS_NUMEROS_pdf.pdf","autor":"Gentil Lopes da Silva","ciencia":"teoria_dos_numeros","dominio":"matematica","nivel":4,"numero":"2","tipo_no":"paper","tipo_teorema":"teorema","titulo_legivel":"Fundamentos dos Números"},"origem":"gentil/Gentil_Lopes_FUNDAMENTOS_DOS_NUMEROS_pdf.pdf","palavras_chave":[],"sinonimos":[],"tipo":"conceito","titulo":"gentil fundamentos Teorema 2 — p 69) temos \u0000a(n + 1) + bn, an + b(n + 1) \u0001 ∼(a, b) ⇒ (5.8)"}
\section{gentil fundamentos Teorema 2 — p 69) temos a(n + 1) + bn, an + b(n + 1)  \ensuremath{\sim}(a, b) \ensuremath{\Rightarrow} (5.8)}
p 69) temos a(n + 1) + bn, an + b(n + 1)  \ensuremath{\sim}(a, b) \ensuremath{\Rightarrow} (5.8) Portanto, \ensuremath{\alpha} \textperiodcentered  \ensuremath{\Theta} = \ensuremath{\alpha}. De modo análogo demonstra-se que \ensuremath{\Theta} é neutro pela esquerda. \rule{1ex}{1ex} Vamos utilizar a seguinte notação para o elemento neutro da multiplicação: (n + 1, n) = \ensuremath{\bar{\,}}1
\prov{proveniência: gentil/Gentil\_Lopes\_FUNDAMENTOS\_DOS\_NUMEROS\_pdf.pdf \textperiodcentered{} fato \textperiodcentered{} confiança 1.0 \textperiodcentered{} domínio matematica \textperiodcentered{} história 5 versões no jornal}

%CRISTAL {"arestas":[{"alvo":"matematica","confianca":0.95,"epistemico":"fato","origem":"paper","rel":"parte_de"},{"alvo":"virtualizacao","confianca":0.5,"epistemico":"fato","origem":"wiki","rel":"relaciona"},{"alvo":"word","confianca":0.5,"epistemico":"fato","origem":"wiki","rel":"relaciona"}],"confianca":1.0,"contraexemplos":[],"descricao":"p 69) . Unicidade: Sejam α′ e α′′ inteiros que satisfazem ao","epistemico":"fato","exemplos":[],"id":"gentil_fundamentos_teorema_2_p_69_unicidade_sejam_α_e_α_inteiros_que_satisfazem_","memoria":"semantica","meta":{"arquivo":"gentil/Lopes_Lopes_FUNDAMENTOS_DOS_NUMEROS_pdf.pdf","autor":"Gentil Lopes da Silva","ciencia":"teoria_dos_numeros","dominio":"matematica","nivel":4,"numero":"2","tipo_no":"paper","tipo_teorema":"teorema","titulo_legivel":"Fundamentos dos Números"},"origem":"gentil/Gentil_Lopes_FUNDAMENTOS_DOS_NUMEROS_pdf.pdf","palavras_chave":[],"sinonimos":[],"tipo":"conceito","titulo":"gentil fundamentos Teorema 2 — p 69) . Unicidade: Sejam α′ e α′′ inteiros que satisfazem ao"}
\section{gentil fundamentos Teorema 2 — p 69) . Unicidade: Sejam \ensuremath{\alpha}\ensuremath{'} e \ensuremath{\alpha}\ensuremath{'}\ensuremath{'} inteiros que satisfazem ao}
p 69) . Unicidade: Sejam \ensuremath{\alpha}\ensuremath{'} e \ensuremath{\alpha}\ensuremath{'}\ensuremath{'} inteiros que satisfazem ao
\prov{relações: parte\_de matematica; relaciona virtualizacao; relaciona word}
\prov{proveniência: gentil/Gentil\_Lopes\_FUNDAMENTOS\_DOS\_NUMEROS\_pdf.pdf \textperiodcentered{} fato \textperiodcentered{} confiança 1.0 \textperiodcentered{} domínio matematica \textperiodcentered{} história 6 versões no jornal}

%CRISTAL {"arestas":[{"alvo":"matematica","confianca":0.95,"epistemico":"fato","origem":"paper","rel":"parte_de"}],"confianca":1.0,"contraexemplos":[],"descricao":"Dados um conjunto não vazio A e uma relação de equivalência ∼em A o conjunto-quociente A/∼é uma partição de A. Prova: Pela","epistemico":"fato","exemplos":[],"id":"gentil_fundamentos_teorema_5_dados_um_conjunto_nao_vazio_a_e_uma_relacao_de_equi","memoria":"semantica","meta":{"arquivo":"gentil/Lopes_Lopes_FUNDAMENTOS_DOS_NUMEROS_pdf.pdf","autor":"Gentil Lopes da Silva","ciencia":"teoria_dos_numeros","dominio":"matematica","nivel":4,"numero":"5","tipo_no":"paper","tipo_teorema":"teorema","titulo_legivel":"Fundamentos dos Números"},"origem":"gentil/Gentil_Lopes_FUNDAMENTOS_DOS_NUMEROS_pdf.pdf","palavras_chave":[],"sinonimos":[],"tipo":"conceito","titulo":"gentil fundamentos Teorema 5 — Dados um conjunto não vazio A e uma relação de equivalência"}
\section{gentil fundamentos Teorema 5 — Dados um conjunto não vazio A e uma relação de equivalência}
Dados um conjunto não vazio A e uma relação de equivalência \ensuremath{\sim}em A o conjunto-quociente A/\ensuremath{\sim}é uma partição de A. Prova: Pela
\prov{relações: parte\_de matematica}
\prov{proveniência: gentil/Gentil\_Lopes\_FUNDAMENTOS\_DOS\_NUMEROS\_pdf.pdf \textperiodcentered{} fato \textperiodcentered{} confiança 1.0 \textperiodcentered{} domínio matematica \textperiodcentered{} história 4 versões no jornal}

%CRISTAL {"arestas":[{"alvo":"matematica","confianca":0.95,"epistemico":"fato","origem":"paper","rel":"parte_de"}],"confianca":1.0,"contraexemplos":[],"descricao":"p 71) também é verdadeira. Teorema 6. Seja F = P1, P2,..., Pn,... uma partição do conjunto não vazio A. A relação ∼em A tal que dois elementos x, y ∈A estão na relação ∼se, e somente se, x e y pertencem a uma mesma cela Pi da partição F, isto é, cuja sentença aberta é dada por, p(x, y): x ∼y ⇐⇒ x, y ∈Pi é uma relação de equivalência em A. Prova: Então, quaisquer que sejam os elementos x, y e z ∈A, temos: ( i ) ( x ∼x ). Com efeito, seja x ∈A, pela","epistemico":"fato","exemplos":[],"id":"gentil_fundamentos_teorema_5_p_71_tambem_e_verdadeira_teorema_6_seja_f_p1_p2","memoria":"semantica","meta":{"arquivo":"gentil/Lopes_Lopes_FUNDAMENTOS_DOS_NUMEROS_pdf.pdf","autor":"Gentil Lopes da Silva","ciencia":"teoria_dos_numeros","dominio":"matematica","nivel":4,"numero":"5","tipo_no":"paper","tipo_teorema":"teorema","titulo_legivel":"Fundamentos dos Números"},"origem":"gentil/Gentil_Lopes_FUNDAMENTOS_DOS_NUMEROS_pdf.pdf","palavras_chave":[],"sinonimos":[],"tipo":"conceito","titulo":"gentil fundamentos Teorema 5 — p 71) também é verdadeira. Teorema 6. Seja F = P1, P2,.."}
\section{gentil fundamentos Teorema 5 — p 71) também é verdadeira. Teorema 6. Seja F = P1, P2,..}
p 71) também é verdadeira. Teorema 6. Seja F = P1, P2,..., Pn,... uma partição do conjunto não vazio A. A relação \ensuremath{\sim}em A tal que dois elementos x, y \ensuremath{\in}A estão na relação \ensuremath{\sim}se, e somente se, x e y pertencem a uma mesma cela Pi da partição F, isto é, cuja sentença aberta é dada por, p(x, y): x \ensuremath{\sim}y \ensuremath{\Leftarrow}\ensuremath{\Rightarrow} x, y \ensuremath{\in}Pi é uma relação de equivalência em A. Prova: Então, quaisquer que sejam os elementos x, y e z \ensuremath{\in}A, temos: ( i ) ( x \ensuremath{\sim}x ). Com efeito, seja x \ensuremath{\in}A, pela
\prov{relações: parte\_de matematica}
\prov{proveniência: gentil/Gentil\_Lopes\_FUNDAMENTOS\_DOS\_NUMEROS\_pdf.pdf \textperiodcentered{} fato \textperiodcentered{} confiança 1.0 \textperiodcentered{} domínio matematica \textperiodcentered{} história 6 versões no jornal}

%CRISTAL {"arestas":[{"alvo":"matematica","confianca":0.95,"epistemico":"fato","origem":"paper","rel":"parte_de"}],"confianca":1.0,"contraexemplos":[],"descricao":"Para todo n ∈N tem-se σ(n) ̸= n Prova: Com efeito, seja A = n ∈N: σ(n) ̸= n. Tem-se 0 ∈A, pois por P1, 0 não é sucessor de número algum, em particular 0 ̸= σ(0). Agora vamos assumir que, para um n arbitrariamente fixado, n ̸= σ(n). Logo, pela injetividade de σ temos que σ(n) ̸= σ \u0000σ(n) \u0001, portanto σ(n) ∈A. Assim, mostramos que n ∈A implica σ(n) ∈A. Como 0 ∈A, segue-se do axioma P3 que A = N, ou seja, n ̸= σ(n) para todo n ∈N. ■ Deste","epistemico":"fato","exemplos":[],"id":"gentil_fundamentos_teorema_8_para_todo_n_n_tem-se_σn_n_prova_com_efeito_seja_a","memoria":"semantica","meta":{"arquivo":"gentil/Lopes_Lopes_FUNDAMENTOS_DOS_NUMEROS_pdf.pdf","autor":"Gentil Lopes da Silva","ciencia":"teoria_dos_numeros","dominio":"matematica","nivel":4,"numero":"8","tipo_no":"paper","tipo_teorema":"teorema","titulo_legivel":"Fundamentos dos Números"},"origem":"gentil/Gentil_Lopes_FUNDAMENTOS_DOS_NUMEROS_pdf.pdf","palavras_chave":[],"sinonimos":[],"tipo":"conceito","titulo":"gentil fundamentos Teorema 8 — Para todo n ∈N tem-se σ(n) ̸= n Prova: Com efeito, seja A ="}
\section{gentil fundamentos Teorema 8 — Para todo n \ensuremath{\in}N tem-se \ensuremath{\sigma}(n) \ensuremath{}= n Prova: Com efeito, seja A =}
Para todo n \ensuremath{\in}N tem-se \ensuremath{\sigma}(n) \ensuremath{}= n Prova: Com efeito, seja A = n \ensuremath{\in}N: \ensuremath{\sigma}(n) \ensuremath{}= n. Tem-se 0 \ensuremath{\in}A, pois por P1, 0 não é sucessor de número algum, em particular 0 \ensuremath{}= \ensuremath{\sigma}(0). Agora vamos assumir que, para um n arbitrariamente fixado, n \ensuremath{}= \ensuremath{\sigma}(n). Logo, pela injetividade de \ensuremath{\sigma} temos que \ensuremath{\sigma}(n) \ensuremath{}= \ensuremath{\sigma} \ensuremath{\sigma}(n) , portanto \ensuremath{\sigma}(n) \ensuremath{\in}A. Assim, mostramos que n \ensuremath{\in}A implica \ensuremath{\sigma}(n) \ensuremath{\in}A. Como 0 \ensuremath{\in}A, segue-se do axioma P3 que A = N, ou seja, n \ensuremath{}= \ensuremath{\sigma}(n) para todo n \ensuremath{\in}N. \rule{1ex}{1ex} Deste
\prov{relações: parte\_de matematica}
\prov{proveniência: gentil/Gentil\_Lopes\_FUNDAMENTOS\_DOS\_NUMEROS\_pdf.pdf \textperiodcentered{} fato \textperiodcentered{} confiança 1.0 \textperiodcentered{} domínio matematica \textperiodcentered{} história 5 versões no jornal}

%CRISTAL {"arestas":[{"alvo":"gentil_novas_sequencias_aritmeticas_geometricas","confianca":1.0,"epistemico":"fato","origem":"gentil/Novas_Sequencias_Aritmeticas_e_Geometric.pdf","rel":"prova"},{"alvo":"gentil_capitulo_1_sequencias_aritmeticas_de_ordem_m","confianca":1.0,"epistemico":"fato","origem":"gentil/Novas_Sequencias_Aritmeticas_e_Geometric.pdf","rel":"prova"}],"confianca":1.0,"contraexemplos":[],"descricao":"402\n2o) (p+1)−1 2m−1 = (p−1)+1 2m−1 ̸∈Z Sendo assim, temos (−1) \u0016 (p+1)−1 2m−1 \u0017 = (−1) \u0016 p−1 2m−1 \u0017 = (−1)( p−1 2m−1) é suficiente mostrar que \u0000 p−1 2m−1 \u0001 e \u0000(p+1)−1 2m−1 \u0001 têm a mesma paridade. Isto foi feito no","epistemico":"fato","exemplos":[],"id":"gentil_lema_11_402_2o_p11_2m1_p11_2m1_z_sendo_assim_temos","memoria":"semantica","meta":{"arquivo":"gentil/Novas_Sequencias_Aritmeticas_e_Geometric.pdf","ciencia":"teoria_dos_numeros","dominio":"matematica","nivel":4,"numero":"11","tipo_no":"paper","tipo_teorema":"lema"},"origem":"gentil/Novas_Sequencias_Aritmeticas_e_Geometric.pdf","palavras_chave":[],"sinonimos":[],"tipo":"conceito","titulo":"gentil Lema 11 — 402\n2o) (p+1)−1 2m−1 = (p−1)+1 2m−1 ̸∈Z Sendo assim, temos ("}
\section{gentil Lema 11 — 402
2o) (p+1)\ensuremath{-}1 2m\ensuremath{-}1 = (p\ensuremath{-}1)+1 2m\ensuremath{-}1 \ensuremath{}\ensuremath{\in}Z Sendo assim, temos (}
402
2o) (p+1)\ensuremath{-}1 2m\ensuremath{-}1 = (p\ensuremath{-}1)+1 2m\ensuremath{-}1 \ensuremath{}\ensuremath{\in}Z Sendo assim, temos (\ensuremath{-}1)  (p+1)\ensuremath{-}1 2m\ensuremath{-}1  = (\ensuremath{-}1)  p\ensuremath{-}1 2m\ensuremath{-}1  = (\ensuremath{-}1)( p\ensuremath{-}1 2m\ensuremath{-}1) é suficiente mostrar que  p\ensuremath{-}1 2m\ensuremath{-}1  e (p+1)\ensuremath{-}1 2m\ensuremath{-}1  têm a mesma paridade. Isto foi feito no
\prov{relações: prova gentil\_novas\_sequencias\_aritmeticas\_geometricas; prova gentil\_capitulo\_1\_sequencias\_aritmeticas\_de\_ordem\_m}
\prov{proveniência: gentil/Novas\_Sequencias\_Aritmeticas\_e\_Geometric.pdf \textperiodcentered{} fato \textperiodcentered{} confiança 1.0 \textperiodcentered{} domínio matematica \textperiodcentered{} história 4 versões no jornal}

%CRISTAL {"arestas":[{"alvo":"gentil_novas_sequencias_aritmeticas_geometricas","confianca":1.0,"epistemico":"fato","origem":"gentil/Novas_Sequencias_Aritmeticas_e_Geometric.pdf","rel":"prova"},{"alvo":"gentil_capitulo_1_sequencias_aritmeticas_de_ordem_m","confianca":1.0,"epistemico":"fato","origem":"gentil/Novas_Sequencias_Aritmeticas_e_Geometric.pdf","rel":"prova"}],"confianca":1.0,"contraexemplos":[],"descricao":"Lopes/18 07.1999). Sejam j e n inteiros positivos arbitrariamente fixa- dos. Temos que n 2 j−1 ∈Z ⇐⇒ \u0012n −1 2 j−1 \u0013 e \u0012 n 2 j−1 \u0013 têm paridades distintas. Antes da prova do","epistemico":"fato","exemplos":[],"id":"gentil_lema_11_gentil18_071999_sejam_j_e_n_inteiros_positivos_arbitrari","memoria":"semantica","meta":{"arquivo":"gentil/Novas_Sequencias_Aritmeticas_e_Geometric.pdf","ciencia":"teoria_dos_numeros","dominio":"matematica","nivel":4,"numero":"11","tipo_no":"paper","tipo_teorema":"lema"},"origem":"gentil/Novas_Sequencias_Aritmeticas_e_Geometric.pdf","palavras_chave":[],"sinonimos":[],"tipo":"conceito","titulo":"gentil Lema 11 — Lopes/18 07.1999). Sejam j e n inteiros positivos arbitrari"}
\section{gentil Lema 11 — Lopes/18 07.1999). Sejam j e n inteiros positivos arbitrari}
Lopes/18 07.1999). Sejam j e n inteiros positivos arbitrariamente fixa- dos. Temos que n 2 j\ensuremath{-}1 \ensuremath{\in}Z \ensuremath{\Leftarrow}\ensuremath{\Rightarrow} n \ensuremath{-}1 2 j\ensuremath{-}1  e  n 2 j\ensuremath{-}1  têm paridades distintas. Antes da prova do
\prov{relações: prova gentil\_novas\_sequencias\_aritmeticas\_geometricas; prova gentil\_capitulo\_1\_sequencias\_aritmeticas\_de\_ordem\_m}
\prov{proveniência: gentil/Novas\_Sequencias\_Aritmeticas\_e\_Geometric.pdf \textperiodcentered{} fato \textperiodcentered{} confiança 1.0 \textperiodcentered{} domínio matematica \textperiodcentered{} história 5 versões no jornal}

%CRISTAL {"arestas":[{"alvo":"gentil_novas_sequencias_aritmeticas_geometricas","confianca":1.0,"epistemico":"fato","origem":"gentil/Novas_Sequencias_Aritmeticas_e_Geometric.pdf","rel":"prova"},{"alvo":"gentil_capitulo_1_sequencias_aritmeticas_de_ordem_m","confianca":1.0,"epistemico":"fato","origem":"gentil/Novas_Sequencias_Aritmeticas_e_Geometric.pdf","rel":"prova"}],"confianca":1.0,"contraexemplos":[],"descricao":"p 399), que estas sequências são iguais. Teorema 52 (Lopes/1999). As sequências ( an ) e ( bn ) dadas acima são iguais. Prova: Pela","epistemico":"fato","exemplos":[],"id":"gentil_lema_11_p_399_que_estas_sequencias_sao_iguais_teorema_52_gentil","memoria":"semantica","meta":{"arquivo":"gentil/Novas_Sequencias_Aritmeticas_e_Geometric.pdf","ciencia":"teoria_dos_numeros","dominio":"matematica","nivel":4,"numero":"11","tipo_no":"paper","tipo_teorema":"lema"},"origem":"gentil/Novas_Sequencias_Aritmeticas_e_Geometric.pdf","palavras_chave":[],"sinonimos":[],"tipo":"conceito","titulo":"gentil Lema 11 — p 399), que estas sequências são iguais. Teorema 52 (Lopes/"}
\section{gentil Lema 11 — p 399), que estas sequências são iguais. Teorema 52 (Lopes/}
p 399), que estas sequências são iguais. Teorema 52 (Lopes/1999). As sequências ( an ) e ( bn ) dadas acima são iguais. Prova: Pela
\prov{relações: prova gentil\_novas\_sequencias\_aritmeticas\_geometricas; prova gentil\_capitulo\_1\_sequencias\_aritmeticas\_de\_ordem\_m}
\prov{proveniência: gentil/Novas\_Sequencias\_Aritmeticas\_e\_Geometric.pdf \textperiodcentered{} fato \textperiodcentered{} confiança 1.0 \textperiodcentered{} domínio matematica \textperiodcentered{} história 5 versões no jornal}

%CRISTAL {"arestas":[{"alvo":"matematica","confianca":0.95,"epistemico":"fato","origem":"paper","rel":"parte_de"}],"confianca":1.0,"contraexemplos":[],"descricao":"p 399) surgiu quando tentavamos demonstrar a igualdade entre as matrizes que comparecem na página 396. Posteriormente concebemos uma prova mais direta da referida identidade: Pelo corolário do","epistemico":"fato","exemplos":[],"id":"gentil_lema_11_p_399_surgiu_quando_tentavamos_demonstrar_a_igualdade_entre","memoria":"semantica","meta":{"arquivo":"gentil/Novas_Sequencias_Aritmeticas_e_Geometric.pdf","ciencia":"teoria_dos_numeros","dominio":"matematica","nivel":4,"numero":"11","tipo_no":"paper","tipo_teorema":"lema"},"origem":"gentil/Novas_Sequencias_Aritmeticas_e_Geometric.pdf","palavras_chave":[],"sinonimos":[],"tipo":"conceito","titulo":"gentil Lema 11 — p 399) surgiu quando tentavamos demonstrar a igualdade entre"}
\section{gentil Lema 11 — p 399) surgiu quando tentavamos demonstrar a igualdade entre}
p 399) surgiu quando tentavamos demonstrar a igualdade entre as matrizes que comparecem na página 396. Posteriormente concebemos uma prova mais direta da referida identidade: Pelo corolário do
\prov{relações: parte\_de matematica}
\prov{proveniência: gentil/Novas\_Sequencias\_Aritmeticas\_e\_Geometric.pdf \textperiodcentered{} fato \textperiodcentered{} confiança 1.0 \textperiodcentered{} domínio matematica \textperiodcentered{} história 3 versões no jornal}

%CRISTAL {"arestas":[{"alvo":"gentil_novas_sequencias_aritmeticas_geometricas","confianca":1.0,"epistemico":"fato","origem":"gentil/Novas_Sequencias_Aritmeticas_e_Geometric.pdf","rel":"prova"}],"confianca":1.0,"contraexemplos":[],"descricao":"p 17), temos: an1 = a11 + S(n−1)(1−1) onde, S(n−1)0 é a soma dos n−1 termos iniciais da P.A. de ordem zero, vale: S(n−1)0 = r + r + · · · + r | z (n−1) termos = (n −1) r Sendo assim, temos: an1 = a11 + (n −1) r (1.2) ´E a fórmula do termo geral de uma P.A.1. (ii) m = 2: Utilizando o","epistemico":"fato","exemplos":[],"id":"gentil_lema_1_p_17_temos_an1_a11_sn111_onde_sn10_e_a_soma","memoria":"semantica","meta":{"arquivo":"gentil/Novas_Sequencias_Aritmeticas_e_Geometric.pdf","ciencia":"teoria_dos_numeros","dominio":"matematica","nivel":4,"numero":"1","tipo_no":"paper","tipo_teorema":"lema"},"origem":"gentil/Novas_Sequencias_Aritmeticas_e_Geometric.pdf","palavras_chave":[],"sinonimos":[],"tipo":"conceito","titulo":"gentil Lema 1 — p 17), temos: an1 = a11 + S(n−1)(1−1) onde, S(n−1)0 é a soma"}
\section{gentil Lema 1 — p 17), temos: an1 = a11 + S(n\ensuremath{-}1)(1\ensuremath{-}1) onde, S(n\ensuremath{-}1)0 é a soma}
p 17), temos: an1 = a11 + S(n\ensuremath{-}1)(1\ensuremath{-}1) onde, S(n\ensuremath{-}1)0 é a soma dos n\ensuremath{-}1 termos iniciais da P.A. de ordem zero, vale: S(n\ensuremath{-}1)0 = r + r + \textperiodcentered  \textperiodcentered  \textperiodcentered  + r | z (n\ensuremath{-}1) termos = (n \ensuremath{-}1) r Sendo assim, temos: an1 = a11 + (n \ensuremath{-}1) r (1.2) 'E a fórmula do termo geral de uma P.A.1. (ii) m = 2: Utilizando o
\prov{relações: prova gentil\_novas\_sequencias\_aritmeticas\_geometricas}
\prov{proveniência: gentil/Novas\_Sequencias\_Aritmeticas\_e\_Geometric.pdf \textperiodcentered{} fato \textperiodcentered{} confiança 1.0 \textperiodcentered{} domínio matematica \textperiodcentered{} história 5 versões no jornal}

%CRISTAL {"arestas":[{"alvo":"gentil_novas_sequencias_aritmeticas_geometricas","confianca":1.0,"epistemico":"fato","origem":"gentil/Novas_Sequencias_Aritmeticas_e_Geometric.pdf","rel":"prova"}],"confianca":1.0,"contraexemplos":[],"descricao":"p 17), temos: an2 = a12 + S(n−1)(2−1) (1.3) onde, S(n−1)1 é a soma dos n −1 termos iniciais da P.A.1 . Vamos deduzir esta fórmula utilizando a equação (1.2), assim: a11 = a11 a21 = a11 + r a31 = a11 + 2r · · · · · · · · · · · · · · · · · · · · · an1 = a11 + (n −1) r Somando estas n igualdades, resulta: a11 + a21 + a31 + · · · + an1 = (a11 + a11 + a11 + · · · + a11) + \u00001 + 2 + · · · + (n −1) \u0001 r O seguinte resultado já é conhecido, 1 + 2 + · · · + n = n(n + 1) 2 ⇒1 + 2 + · · · + n −1 = (n −1)n 2","epistemico":"fato","exemplos":[],"id":"gentil_lema_1_p_17_temos_an2_a12_sn121_13_onde_sn11_e","memoria":"semantica","meta":{"arquivo":"gentil/Novas_Sequencias_Aritmeticas_e_Geometric.pdf","ciencia":"teoria_dos_numeros","dominio":"matematica","nivel":4,"numero":"1","tipo_no":"paper","tipo_teorema":"lema"},"origem":"gentil/Novas_Sequencias_Aritmeticas_e_Geometric.pdf","palavras_chave":[],"sinonimos":[],"tipo":"conceito","titulo":"gentil Lema 1 — p 17), temos: an2 = a12 + S(n−1)(2−1) (1.3) onde, S(n−1)1 é"}
\section{gentil Lema 1 — p 17), temos: an2 = a12 + S(n\ensuremath{-}1)(2\ensuremath{-}1) (1.3) onde, S(n\ensuremath{-}1)1 é}
p 17), temos: an2 = a12 + S(n\ensuremath{-}1)(2\ensuremath{-}1) (1.3) onde, S(n\ensuremath{-}1)1 é a soma dos n \ensuremath{-}1 termos iniciais da P.A.1 . Vamos deduzir esta fórmula utilizando a equação (1.2), assim: a11 = a11 a21 = a11 + r a31 = a11 + 2r \textperiodcentered  \textperiodcentered  \textperiodcentered  \textperiodcentered  \textperiodcentered  \textperiodcentered  \textperiodcentered  \textperiodcentered  \textperiodcentered  \textperiodcentered  \textperiodcentered  \textperiodcentered  \textperiodcentered  \textperiodcentered  \textperiodcentered  \textperiodcentered  \textperiodcentered  \textperiodcentered  \textperiodcentered  \textperiodcentered  \textperiodcentered  an1 = a11 + (n \ensuremath{-}1) r Somando estas n igualdades, resulta: a11 + a21 + a31 + \textperiodcentered  \textperiodcentered  \textperiodcentered  + an1 = (a11 + a11 + a11 + \textperiodcentered  \textperiodcentered  \textperiodcentered  + a11) + 1 + 2 + \textperiodcentered  \textperiodcentered  \textperiodcentered  + (n \ensuremath{-}1)  r O seguinte resultado já é conhecido, 1 + 2 + \textperiodcentered  \textperiodcentered  \textperiodcentered  + n = n(n + 1) 2 \ensuremath{\Rightarrow}1 + 2 + \textperiodcentered  \textperiodcentered  \textperiodcentered  + n \ensuremath{-}1 = (n \ensuremath{-}1)n 2
\prov{relações: prova gentil\_novas\_sequencias\_aritmeticas\_geometricas}
\prov{proveniência: gentil/Novas\_Sequencias\_Aritmeticas\_e\_Geometric.pdf \textperiodcentered{} fato \textperiodcentered{} confiança 1.0 \textperiodcentered{} domínio matematica \textperiodcentered{} história 4 versões no jornal}

%CRISTAL {"arestas":[{"alvo":"gentil_novas_sequencias_aritmeticas_geometricas","confianca":1.0,"epistemico":"fato","origem":"gentil/Novas_Sequencias_Aritmeticas_e_Geometric.pdf","rel":"prova"}],"confianca":1.0,"contraexemplos":[],"descricao":"p 104), resulta: ∆(p+1) ∇m f(n) = ∇m−p−1 f(n) Isto é, ∆(p+1) ∇m f(n) = ∇m−(p+1) f(n) ■","epistemico":"fato","exemplos":[],"id":"gentil_lema_2_p_104_resulta_p1_m_fn_mp1_fn_isto_e_p1","memoria":"semantica","meta":{"arquivo":"gentil/Novas_Sequencias_Aritmeticas_e_Geometric.pdf","ciencia":"teoria_dos_numeros","dominio":"matematica","nivel":4,"numero":"2","tipo_no":"paper","tipo_teorema":"lema"},"origem":"gentil/Novas_Sequencias_Aritmeticas_e_Geometric.pdf","palavras_chave":[],"sinonimos":[],"tipo":"conceito","titulo":"gentil Lema 2 — p 104), resulta: ∆(p+1) ∇m f(n) = ∇m−p−1 f(n) Isto é, ∆(p+1)"}
\section{gentil Lema 2 — p 104), resulta: \ensuremath{\Delta}(p+1) \ensuremath{\nabla}m f(n) = \ensuremath{\nabla}m\ensuremath{-}p\ensuremath{-}1 f(n) Isto é, \ensuremath{\Delta}(p+1)}
p 104), resulta: \ensuremath{\Delta}(p+1) \ensuremath{\nabla}m f(n) = \ensuremath{\nabla}m\ensuremath{-}p\ensuremath{-}1 f(n) Isto é, \ensuremath{\Delta}(p+1) \ensuremath{\nabla}m f(n) = \ensuremath{\nabla}m\ensuremath{-}(p+1) f(n) \rule{1ex}{1ex}
\prov{relações: prova gentil\_novas\_sequencias\_aritmeticas\_geometricas}
\prov{proveniência: gentil/Novas\_Sequencias\_Aritmeticas\_e\_Geometric.pdf \textperiodcentered{} fato \textperiodcentered{} confiança 1.0 \textperiodcentered{} domínio matematica \textperiodcentered{} história 4 versões no jornal}

%CRISTAL {"arestas":[{"alvo":"matematica","confianca":0.95,"epistemico":"fato","origem":"paper","rel":"parte_de"}],"confianca":1.0,"contraexemplos":[],"descricao":"Seja m um natural arbitrariamente fixado e j um natural tal que 0 ≤ j ≤m Nestas condições a seguinte identidade se verifica: ∆ \u0000∇m−p f(n) \u0001 = ∇m−p−1 f(n) Prova: Temos que ∆f(n) = f(n + 1) −f(n), logo ∆ \u0000∇m−p f(n) \u0001 = ∇m−p f(n + 1) −∇m−p f(n) (2.13) Pela","epistemico":"fato","exemplos":[],"id":"gentil_lema_2_seja_m_um_natural_arbitrariamente_fixado_e_j_um_natural_tal","memoria":"semantica","meta":{"arquivo":"gentil/Novas_Sequencias_Aritmeticas_e_Geometric.pdf","ciencia":"teoria_dos_numeros","dominio":"matematica","nivel":4,"numero":"2","tipo_no":"paper","tipo_teorema":"lema"},"origem":"gentil/Novas_Sequencias_Aritmeticas_e_Geometric.pdf","palavras_chave":[],"sinonimos":[],"tipo":"conceito","titulo":"gentil Lema 2 — Seja m um natural arbitrariamente fixado e j um natural tal"}
\section{gentil Lema 2 — Seja m um natural arbitrariamente fixado e j um natural tal}
Seja m um natural arbitrariamente fixado e j um natural tal que 0 \ensuremath{\le} j \ensuremath{\le}m Nestas condições a seguinte identidade se verifica: \ensuremath{\Delta} \ensuremath{\nabla}m\ensuremath{-}p f(n)  = \ensuremath{\nabla}m\ensuremath{-}p\ensuremath{-}1 f(n) Prova: Temos que \ensuremath{\Delta}f(n) = f(n + 1) \ensuremath{-}f(n), logo \ensuremath{\Delta} \ensuremath{\nabla}m\ensuremath{-}p f(n)  = \ensuremath{\nabla}m\ensuremath{-}p f(n + 1) \ensuremath{-}\ensuremath{\nabla}m\ensuremath{-}p f(n) (2.13) Pela
\prov{relações: parte\_de matematica}
\prov{proveniência: gentil/Novas\_Sequencias\_Aritmeticas\_e\_Geometric.pdf \textperiodcentered{} fato \textperiodcentered{} confiança 1.0 \textperiodcentered{} domínio matematica \textperiodcentered{} história 3 versões no jornal}

%CRISTAL {"arestas":[{"alvo":"gentil_novas_sequencias_aritmeticas_geometricas","confianca":1.0,"epistemico":"fato","origem":"gentil/Novas_Sequencias_Aritmeticas_e_Geometric.pdf","rel":"prova"},{"alvo":"gentil_capitulo_1_sequencias_aritmeticas_de_ordem_m","confianca":1.0,"epistemico":"fato","origem":"gentil/Novas_Sequencias_Aritmeticas_e_Geometric.pdf","rel":"prova"}],"confianca":1.0,"contraexemplos":[],"descricao":"´E válida a seguinte identidade ∆ m f(n) = ∆ m−1 ∆ f(n) Prova: Princ´ıpio da Dualidade. (p. 88) ■ O seguinte","epistemico":"fato","exemplos":[],"id":"gentil_lema_4_e_valida_a_seguinte_identidade_m_fn_m1_fn_prov","memoria":"semantica","meta":{"arquivo":"gentil/Novas_Sequencias_Aritmeticas_e_Geometric.pdf","ciencia":"teoria_dos_numeros","dominio":"matematica","nivel":4,"numero":"4","tipo_no":"paper","tipo_teorema":"lema"},"origem":"gentil/Novas_Sequencias_Aritmeticas_e_Geometric.pdf","palavras_chave":[],"sinonimos":[],"tipo":"conceito","titulo":"gentil Lema 4 — ´E válida a seguinte identidade ∆ m f(n) = ∆ m−1 ∆ f(n) Prov"}
\section{gentil Lema 4 — 'E válida a seguinte identidade \ensuremath{\Delta} m f(n) = \ensuremath{\Delta} m\ensuremath{-}1 \ensuremath{\Delta} f(n) Prov}
'E válida a seguinte identidade \ensuremath{\Delta} m f(n) = \ensuremath{\Delta} m\ensuremath{-}1 \ensuremath{\Delta} f(n) Prova: Princ'\ensuremath{\imath}pio da Dualidade. (p. 88) \rule{1ex}{1ex} O seguinte
\prov{relações: prova gentil\_novas\_sequencias\_aritmeticas\_geometricas; prova gentil\_capitulo\_1\_sequencias\_aritmeticas\_de\_ordem\_m}
\prov{proveniência: gentil/Novas\_Sequencias\_Aritmeticas\_e\_Geometric.pdf \textperiodcentered{} fato \textperiodcentered{} confiança 1.0 \textperiodcentered{} domínio matematica \textperiodcentered{} história 4 versões no jornal}

%CRISTAL {"arestas":[{"alvo":"gentil_novas_sequencias_aritmeticas_geometricas","confianca":1.0,"epistemico":"fato","origem":"gentil/Novas_Sequencias_Aritmeticas_e_Geometric.pdf","rel":"prova"},{"alvo":"gentil_capitulo_1_sequencias_aritmeticas_de_ordem_m","confianca":1.0,"epistemico":"fato","origem":"gentil/Novas_Sequencias_Aritmeticas_e_Geometric.pdf","rel":"prova"}],"confianca":1.0,"contraexemplos":[],"descricao":"Seja q um natural arbitrariamente fixado e 0 ≤m ≤q Nestas condições é válida a seguinte identidade: ∆m ( −1 )( n q ) = ( −1 )( n q−m) Prova: Indução sobre m. Para m = 0, temos: ∆0 ( −1 )( n q ) = ( −1 )( n q ) = ( −1 )( n q−0) Admitamos a validade da proposição para m = p (0 ≤p < q): (H.I.) ∆p ( −1 )( n q ) = ( −1 )( n q−p) e provemos que vale para m = p + 1: (T.I.) ∆p+1 ( −1 )( n q ) = ( −1 )( n q−(p+1)) Inicialmente vamos mostrar que ∆( −1 )( n q ) = ( −1 )( n q−1) De fato, ∆( −1 )( n q ) = (","epistemico":"fato","exemplos":[],"id":"gentil_lema_5_seja_q_um_natural_arbitrariamente_fixado_e_0_m_q_nestas_co","memoria":"semantica","meta":{"arquivo":"gentil/Novas_Sequencias_Aritmeticas_e_Geometric.pdf","ciencia":"teoria_dos_numeros","dominio":"matematica","nivel":4,"numero":"5","tipo_no":"paper","tipo_teorema":"lema"},"origem":"gentil/Novas_Sequencias_Aritmeticas_e_Geometric.pdf","palavras_chave":[],"sinonimos":[],"tipo":"conceito","titulo":"gentil Lema 5 — Seja q um natural arbitrariamente fixado e 0 ≤m ≤q Nestas co"}
\section{gentil Lema 5 — Seja q um natural arbitrariamente fixado e 0 \ensuremath{\le}m \ensuremath{\le}q Nestas co}
Seja q um natural arbitrariamente fixado e 0 \ensuremath{\le}m \ensuremath{\le}q Nestas condições é válida a seguinte identidade: \ensuremath{\Delta}m ( \ensuremath{-}1 )( n q ) = ( \ensuremath{-}1 )( n q\ensuremath{-}m) Prova: Indução sobre m. Para m = 0, temos: \ensuremath{\Delta}0 ( \ensuremath{-}1 )( n q ) = ( \ensuremath{-}1 )( n q ) = ( \ensuremath{-}1 )( n q\ensuremath{-}0) Admitamos a validade da proposição para m = p (0 \ensuremath{\le}p < q): (H.I.) \ensuremath{\Delta}p ( \ensuremath{-}1 )( n q ) = ( \ensuremath{-}1 )( n q\ensuremath{-}p) e provemos que vale para m = p + 1: (T.I.) \ensuremath{\Delta}p+1 ( \ensuremath{-}1 )( n q ) = ( \ensuremath{-}1 )( n q\ensuremath{-}(p+1)) Inicialmente vamos mostrar que \ensuremath{\Delta}( \ensuremath{-}1 )( n q ) = ( \ensuremath{-}1 )( n q\ensuremath{-}1) De fato, \ensuremath{\Delta}( \ensuremath{-}1 )( n q ) = (
\prov{relações: prova gentil\_novas\_sequencias\_aritmeticas\_geometricas; prova gentil\_capitulo\_1\_sequencias\_aritmeticas\_de\_ordem\_m}
\prov{proveniência: gentil/Novas\_Sequencias\_Aritmeticas\_e\_Geometric.pdf \textperiodcentered{} fato \textperiodcentered{} confiança 1.0 \textperiodcentered{} domínio matematica \textperiodcentered{} história 4 versões no jornal}

%CRISTAL {"arestas":[{"alvo":"gentil_novas_sequencias_aritmeticas_geometricas","confianca":1.0,"epistemico":"fato","origem":"gentil/Novas_Sequencias_Aritmeticas_e_Geometric.pdf","rel":"prova"}],"confianca":1.0,"contraexemplos":[],"descricao":"p 182). Como obtemos um quociente nulo paramos o processo. Agora é só tomar os restos na seguinte ordem: 20 2 10 0 2 5 0 2 2 1 2 1 0 2 0 1 Sendo assim, temos: 20 = (0 0 1 0 1)2. Ou ainda, 20 = 0 · 20 + 0 · 21 + 1 · 22 + 0 · 23 + 1 · 24 Este exemplo será generalizado a seguir. Deduziremos agora uma fórmula que irá nos permitir escrever um inteiro po- sitivo a na base b = 2. O procedimento que será desenvolvido a seguir se estende facilmente a uma base b qualquer. Dedução da fórmula Vamos deduzir","epistemico":"fato","exemplos":[],"id":"gentil_lema_6_p_182_como_obtemos_um_quociente_nulo_paramos_o_processo_a","memoria":"semantica","meta":{"arquivo":"gentil/Novas_Sequencias_Aritmeticas_e_Geometric.pdf","ciencia":"teoria_dos_numeros","dominio":"matematica","nivel":4,"numero":"6","tipo_no":"paper","tipo_teorema":"lema"},"origem":"gentil/Novas_Sequencias_Aritmeticas_e_Geometric.pdf","palavras_chave":[],"sinonimos":[],"tipo":"conceito","titulo":"gentil Lema 6 — p 182). Como obtemos um quociente nulo paramos o processo. A"}
\section{gentil Lema 6 — p 182). Como obtemos um quociente nulo paramos o processo. A}
p 182). Como obtemos um quociente nulo paramos o processo. Agora é só tomar os restos na seguinte ordem: 20 2 10 0 2 5 0 2 2 1 2 1 0 2 0 1 Sendo assim, temos: 20 = (0 0 1 0 1)2. Ou ainda, 20 = 0 \textperiodcentered  20 + 0 \textperiodcentered  21 + 1 \textperiodcentered  22 + 0 \textperiodcentered  23 + 1 \textperiodcentered  24 Este exemplo será generalizado a seguir. Deduziremos agora uma fórmula que irá nos permitir escrever um inteiro po- sitivo a na base b = 2. O procedimento que será desenvolvido a seguir se estende facilmente a uma base b qualquer. Dedução da fórmula Vamos deduzir
\prov{relações: prova gentil\_novas\_sequencias\_aritmeticas\_geometricas}
\prov{proveniência: gentil/Novas\_Sequencias\_Aritmeticas\_e\_Geometric.pdf \textperiodcentered{} fato \textperiodcentered{} confiança 1.0 \textperiodcentered{} domínio matematica \textperiodcentered{} história 4 versões no jornal}

%CRISTAL {"arestas":[{"alvo":"gentil_novas_sequencias_aritmeticas_geometricas","confianca":1.0,"epistemico":"fato","origem":"gentil/Novas_Sequencias_Aritmeticas_e_Geometric.pdf","rel":"prova"}],"confianca":1.0,"contraexemplos":[],"descricao":"p 182), na divisão seguinte obteremos um quociente nulo. Chamemos este passo (etapa em que o quociente é nulo) de n-ésimo. Então, podemos escrever a = 2q0 + r0, 0 ≤r0 < 2 q0 = 2q1 + r1, 0 ≤r1 < 2 q1 = 2q2 + r2, 0 ≤r2 < 2 −−−−−− −−−−− qn−2 = 2qn−1 + rn−1, 0 ≤rn−1 < 2 qn−1 = 2 · 0 + rn, 0 < rn < 2. (7.20) Observe, nesta última desigualdade, que rn > 0, isto está justificado na nota logo após a prova do","epistemico":"fato","exemplos":[],"id":"gentil_lema_6_p_182_na_divisao_seguinte_obteremos_um_quociente_nulo_cha","memoria":"semantica","meta":{"arquivo":"gentil/Novas_Sequencias_Aritmeticas_e_Geometric.pdf","ciencia":"teoria_dos_numeros","dominio":"matematica","nivel":4,"numero":"6","tipo_no":"paper","tipo_teorema":"lema"},"origem":"gentil/Novas_Sequencias_Aritmeticas_e_Geometric.pdf","palavras_chave":[],"sinonimos":[],"tipo":"conceito","titulo":"gentil Lema 6 — p 182), na divisão seguinte obteremos um quociente nulo. Cha"}
\section{gentil Lema 6 — p 182), na divisão seguinte obteremos um quociente nulo. Cha}
p 182), na divisão seguinte obteremos um quociente nulo. Chamemos este passo (etapa em que o quociente é nulo) de n-ésimo. Então, podemos escrever a = 2q0 + r0, 0 \ensuremath{\le}r0 < 2 q0 = 2q1 + r1, 0 \ensuremath{\le}r1 < 2 q1 = 2q2 + r2, 0 \ensuremath{\le}r2 < 2 \ensuremath{-}\ensuremath{-}\ensuremath{-}\ensuremath{-}\ensuremath{-}\ensuremath{-} \ensuremath{-}\ensuremath{-}\ensuremath{-}\ensuremath{-}\ensuremath{-} qn\ensuremath{-}2 = 2qn\ensuremath{-}1 + rn\ensuremath{-}1, 0 \ensuremath{\le}rn\ensuremath{-}1 < 2 qn\ensuremath{-}1 = 2 \textperiodcentered  0 + rn, 0 < rn < 2. (7.20) Observe, nesta última desigualdade, que rn > 0, isto está justificado na nota logo após a prova do
\prov{relações: prova gentil\_novas\_sequencias\_aritmeticas\_geometricas}
\prov{proveniência: gentil/Novas\_Sequencias\_Aritmeticas\_e\_Geometric.pdf \textperiodcentered{} fato \textperiodcentered{} confiança 1.0 \textperiodcentered{} domínio matematica \textperiodcentered{} história 4 versões no jornal}

%CRISTAL {"arestas":[{"alvo":"gentil_novas_sequencias_aritmeticas_geometricas","confianca":1.0,"epistemico":"fato","origem":"gentil/Novas_Sequencias_Aritmeticas_e_Geometric.pdf","rel":"prova"}],"confianca":1.0,"contraexemplos":[],"descricao":"Se 0 < a < b então o quociente da divisão de a por b é q = 0 Prova: Dividindo a por b, pelo algoritmo da divisão existem naturais q e r, tais que, a = bq + r, 0 ≤r < b (3.23) Suponhamos q ̸= 0, logo q ≥1. Então, q ≥1 ⇒ bq ≥b. Por outro lado, r ≥0 ⇒bq + r ≥bq. Sendo assim, temos, a = bq + r ≥bq ≥b ⇒a ≥b. o que contradiz a hipótese. ■ Nota: Observe ainda que, como q = 0, por (3.23), a = b · 0 + r, logo, r = a > 0. 182","epistemico":"fato","exemplos":[],"id":"gentil_lema_6_se_0_a_b_entao_o_quociente_da_divisao_de_a_por_b_e_q_0","memoria":"semantica","meta":{"arquivo":"gentil/Novas_Sequencias_Aritmeticas_e_Geometric.pdf","ciencia":"teoria_dos_numeros","dominio":"matematica","nivel":4,"numero":"6","tipo_no":"paper","tipo_teorema":"lema"},"origem":"gentil/Novas_Sequencias_Aritmeticas_e_Geometric.pdf","palavras_chave":[],"sinonimos":[],"tipo":"conceito","titulo":"gentil Lema 6 — Se 0 < a < b então o quociente da divisão de a por b é q = 0"}
\section{gentil Lema 6 — Se 0 < a < b então o quociente da divisão de a por b é q = 0}
Se 0 < a < b então o quociente da divisão de a por b é q = 0 Prova: Dividindo a por b, pelo algoritmo da divisão existem naturais q e r, tais que, a = bq + r, 0 \ensuremath{\le}r < b (3.23) Suponhamos q \ensuremath{}= 0, logo q \ensuremath{\ge}1. Então, q \ensuremath{\ge}1 \ensuremath{\Rightarrow} bq \ensuremath{\ge}b. Por outro lado, r \ensuremath{\ge}0 \ensuremath{\Rightarrow}bq + r \ensuremath{\ge}bq. Sendo assim, temos, a = bq + r \ensuremath{\ge}bq \ensuremath{\ge}b \ensuremath{\Rightarrow}a \ensuremath{\ge}b. o que contradiz a hipótese. \rule{1ex}{1ex} Nota: Observe ainda que, como q = 0, por (3.23), a = b \textperiodcentered  0 + r, logo, r = a > 0. 182
\prov{relações: prova gentil\_novas\_sequencias\_aritmeticas\_geometricas}
\prov{proveniência: gentil/Novas\_Sequencias\_Aritmeticas\_e\_Geometric.pdf \textperiodcentered{} fato \textperiodcentered{} confiança 1.0 \textperiodcentered{} domínio matematica \textperiodcentered{} história 4 versões no jornal}

%CRISTAL {"arestas":[{"alvo":"gentil_novas_sequencias_aritmeticas_geometricas","confianca":1.0,"epistemico":"fato","origem":"gentil/Novas_Sequencias_Aritmeticas_e_Geometric.pdf","rel":"prova"},{"alvo":"gentil_capitulo_1_sequencias_aritmeticas_de_ordem_m","confianca":1.0,"epistemico":"fato","origem":"gentil/Novas_Sequencias_Aritmeticas_e_Geometric.pdf","rel":"prova"}],"confianca":1.0,"contraexemplos":[],"descricao":"p 183), temos: j p −k 2 j−1 k =        j p−k−1 2j−1 k + 1, se p−k 2j−1 ∈Z ; j p−k−1 2j−1 k , se p−k 2j−1 ̸∈Z . Observe que se (p −k)/2 j−1 ̸∈Z, o resultado é imediato. Suponhamos (p −k)/2 j−1 ∈Z. Somando∗ 2 j−1 X k = 0 (−1)k \u0012 2 j−1 k \u0013 = 0 ∗Este resultado sai do binômio de Newton. 185\nà hipótese de indução, resulta: 2 j−1 X k = 0 (−1)k \u0012 2 j−1 k \u0013 \u001a j p −k −1 2 j−1 k + 1 \u001b + 2 j−1 X k = 0 (−1)k \u0012 2 j−1 k \u0013 ≡1 (mod 2) Esta última igualdade pode ser reescrita como: 2 j−1 X k = 0 (−1)k \u0012 2","epistemico":"fato","exemplos":[],"id":"gentil_lema_7_p_183_temos_j_p_k_2_j1_k_j_pk1_2j1_k","memoria":"semantica","meta":{"arquivo":"gentil/Novas_Sequencias_Aritmeticas_e_Geometric.pdf","ciencia":"teoria_dos_numeros","dominio":"matematica","nivel":4,"numero":"7","tipo_no":"paper","tipo_teorema":"lema"},"origem":"gentil/Novas_Sequencias_Aritmeticas_e_Geometric.pdf","palavras_chave":[],"sinonimos":[],"tipo":"conceito","titulo":"gentil Lema 7 — p 183), temos: j p −k 2 j−1 k =        j p−k−1 2j−1 k"}
\section{gentil Lema 7 — p 183), temos: j p \ensuremath{-}k 2 j\ensuremath{-}1 k = \{       j p\ensuremath{-}k\ensuremath{-}1 2j\ensuremath{-}1 k}
p 183), temos: j p \ensuremath{-}k 2 j\ensuremath{-}1 k = \{       j p\ensuremath{-}k\ensuremath{-}1 2j\ensuremath{-}1 k + 1, se p\ensuremath{-}k 2j\ensuremath{-}1 \ensuremath{\in}Z ; j p\ensuremath{-}k\ensuremath{-}1 2j\ensuremath{-}1 k , se p\ensuremath{-}k 2j\ensuremath{-}1 \ensuremath{}\ensuremath{\in}Z . Observe que se (p \ensuremath{-}k)/2 j\ensuremath{-}1 \ensuremath{}\ensuremath{\in}Z, o resultado é imediato. Suponhamos (p \ensuremath{-}k)/2 j\ensuremath{-}1 \ensuremath{\in}Z. Somando\ensuremath{*} 2 j\ensuremath{-}1 X k = 0 (\ensuremath{-}1)k  2 j\ensuremath{-}1 k  = 0 \ensuremath{*}Este resultado sai do binômio de Newton. 185
à hipótese de indução, resulta: 2 j\ensuremath{-}1 X k = 0 (\ensuremath{-}1)k  2 j\ensuremath{-}1 k   j p \ensuremath{-}k \ensuremath{-}1 2 j\ensuremath{-}1 k + 1  + 2 j\ensuremath{-}1 X k = 0 (\ensuremath{-}1)k  2 j\ensuremath{-}1 k  \ensuremath{\equiv}1 (mod 2) Esta última igualdade pode ser reescrita como: 2 j\ensuremath{-}1 X k = 0 (\ensuremath{-}1)k  2
\prov{relações: prova gentil\_novas\_sequencias\_aritmeticas\_geometricas; prova gentil\_capitulo\_1\_sequencias\_aritmeticas\_de\_ordem\_m}
\prov{proveniência: gentil/Novas\_Sequencias\_Aritmeticas\_e\_Geometric.pdf \textperiodcentered{} fato \textperiodcentered{} confiança 1.0 \textperiodcentered{} domínio matematica \textperiodcentered{} história 4 versões no jornal}

%CRISTAL {"arestas":[{"alvo":"gentil_novas_sequencias_aritmeticas_geometricas","confianca":1.0,"epistemico":"fato","origem":"gentil/Novas_Sequencias_Aritmeticas_e_Geometric.pdf","rel":"prova"},{"alvo":"gentil_capitulo_1_sequencias_aritmeticas_de_ordem_m","confianca":1.0,"epistemico":"fato","origem":"gentil/Novas_Sequencias_Aritmeticas_e_Geometric.pdf","rel":"prova"}],"confianca":1.0,"contraexemplos":[],"descricao":"Considere 0 < x ∈R Se x ̸∈Z e 2x ∈Z, então: 2 x = 2 ⌊x ⌋+ 1 Prova: Se x ̸∈Z então x é da forma∗x = ⌊x ⌋+ p.d. (0 < p.d. < 1); então 2 x = 2 ⌊x ⌋+ 2 p.d., (0 < 2 p.d. < 2) como 2x ∈Z, temos 2 p.d. ∈Z, logo, 2 p.d. = 1; disto segue a tese. ■ ■ ∗Onde, p.d.=parte decimal. 188\nCap´ıtulo4 Sequências Especiais “Não é poss´ıvel formular as leis da teoria quântica de um modo absolutamente consistente sem referencia a uma consciência”. (Eugene Wigner) O objetivo deste cap´ıtulo é desenvolver (unificar) al","epistemico":"fato","exemplos":[],"id":"gentil_lema_8_considere_0_x_r_se_x_z_e_2x_z_entao_2_x_2_x_1","memoria":"semantica","meta":{"arquivo":"gentil/Novas_Sequencias_Aritmeticas_e_Geometric.pdf","ciencia":"teoria_dos_numeros","dominio":"matematica","nivel":4,"numero":"8","tipo_no":"paper","tipo_teorema":"lema"},"origem":"gentil/Novas_Sequencias_Aritmeticas_e_Geometric.pdf","palavras_chave":[],"sinonimos":[],"tipo":"conceito","titulo":"gentil Lema 8 — Considere 0 < x ∈R Se x ̸∈Z e 2x ∈Z, então: 2 x = 2 ⌊x ⌋+ 1"}
\section{gentil Lema 8 — Considere 0 < x \ensuremath{\in}R Se x \ensuremath{}\ensuremath{\in}Z e 2x \ensuremath{\in}Z, então: 2 x = 2 \ensuremath{\lfloor}x \ensuremath{\rfloor}+ 1}
Considere 0 < x \ensuremath{\in}R Se x \ensuremath{}\ensuremath{\in}Z e 2x \ensuremath{\in}Z, então: 2 x = 2 \ensuremath{\lfloor}x \ensuremath{\rfloor}+ 1 Prova: Se x \ensuremath{}\ensuremath{\in}Z então x é da forma\ensuremath{*}x = \ensuremath{\lfloor}x \ensuremath{\rfloor}+ p.d. (0 < p.d. < 1); então 2 x = 2 \ensuremath{\lfloor}x \ensuremath{\rfloor}+ 2 p.d., (0 < 2 p.d. < 2) como 2x \ensuremath{\in}Z, temos 2 p.d. \ensuremath{\in}Z, logo, 2 p.d. = 1; disto segue a tese. \rule{1ex}{1ex} \rule{1ex}{1ex} \ensuremath{*}Onde, p.d.=parte decimal. 188
Cap'\ensuremath{\imath}tulo4 Sequências Especiais “Não é poss'\ensuremath{\imath}vel formular as leis da teoria quântica de um modo absolutamente consistente sem referencia a uma consciência”. (Eugene Wigner) O objetivo deste cap'\ensuremath{\imath}tulo é desenvolver (unificar) al
\prov{relações: prova gentil\_novas\_sequencias\_aritmeticas\_geometricas; prova gentil\_capitulo\_1\_sequencias\_aritmeticas\_de\_ordem\_m}
\prov{proveniência: gentil/Novas\_Sequencias\_Aritmeticas\_e\_Geometric.pdf \textperiodcentered{} fato \textperiodcentered{} confiança 1.0 \textperiodcentered{} domínio matematica \textperiodcentered{} história 4 versões no jornal}

%CRISTAL {"arestas":[{"alvo":"teoria_dos_numeros","confianca":1.0,"epistemico":"fato","origem":"gentil/Novas_Sequencias_Aritmeticas_e_Geometric.pdf","rel":"parte_de"}],"confianca":1.0,"contraexemplos":[],"descricao":"Gentil Lopes da Silva (2015). Generalizações de progressões aritméticas e geométricas; teoria dos números e combinatória.","epistemico":"fato","exemplos":[],"id":"gentil_novas_sequencias_aritmeticas_geometricas","memoria":"semantica","meta":{"arquivo":"gentil/Novas_Sequencias_Aritmeticas_e_Geometric.pdf","autor":"Gentil Lopes da Silva","ciencia":"teoria_dos_numeros","dominio":"matematica","nivel":4,"tipo_no":"paper","titulo_legivel":"Novas Sequências Aritméticas e Geométricas"},"origem":"livro","palavras_chave":[],"sinonimos":["Novas Sequências Aritméticas e Geométricas","livro Gentil"],"tipo":"conceito","titulo":"gentil novas sequencias aritmeticas geometricas"}
\section{gentil novas sequencias aritmeticas geometricas}
Gentil Lopes da Silva (2015). Generalizações de progressões aritméticas e geométricas; teoria dos números e combinatória.
\prov{sinónimos: Novas Sequências Aritméticas e Geométricas; livro Gentil}
\prov{relações: parte\_de teoria\_dos\_numeros}
\prov{proveniência: livro \textperiodcentered{} fato \textperiodcentered{} confiança 1.0 \textperiodcentered{} domínio matematica \textperiodcentered{} história 6 versões no jornal}

%CRISTAL {"arestas":[{"alvo":"gentil_novas_sequencias_aritmeticas_geometricas","confianca":1.0,"epistemico":"fato","origem":"gentil/Novas_Sequencias_Aritmeticas_e_Geometric.pdf","rel":"prova"},{"alvo":"gentil_capitulo_1_sequencias_aritmeticas_de_ordem_m","confianca":1.0,"epistemico":"fato","origem":"gentil/Novas_Sequencias_Aritmeticas_e_Geometric.pdf","rel":"prova"}],"confianca":1.0,"contraexemplos":[],"descricao":"p 91) e a equação (2.2) (p. 85), assim: an(m−j) = ∆j anm = j X k = 0 (−1)k \u0012 j k \u0013 a(n−k+j)m (2.7) Substituindo nesta equação n = 1, resulta: a1(m−j) = j X k = 0 (−1)k \u0012 j k \u0013 a(1−k+j)m ■ Prova: (Equação (1.17), p. 44) Tomemos j = m na equação","epistemico":"fato","exemplos":[],"id":"gentil_teorema_10_p_91_e_a_equacao_22_p_85_assim_anmj_j_anm_j","memoria":"semantica","meta":{"arquivo":"gentil/Novas_Sequencias_Aritmeticas_e_Geometric.pdf","ciencia":"teoria_dos_numeros","dominio":"matematica","nivel":4,"numero":"10","tipo_no":"paper","tipo_teorema":"teorema"},"origem":"gentil/Novas_Sequencias_Aritmeticas_e_Geometric.pdf","palavras_chave":[],"sinonimos":[],"tipo":"conceito","titulo":"gentil Teorema 10 — p 91) e a equação (2.2) (p. 85), assim: an(m−j) = ∆j anm = j"}
\section{gentil Teorema 10 — p 91) e a equação (2.2) (p. 85), assim: an(m\ensuremath{-}j) = \ensuremath{\Delta}j anm = j}
p 91) e a equação (2.2) (p. 85), assim: an(m\ensuremath{-}j) = \ensuremath{\Delta}j anm = j X k = 0 (\ensuremath{-}1)k  j k  a(n\ensuremath{-}k+j)m (2.7) Substituindo nesta equação n = 1, resulta: a1(m\ensuremath{-}j) = j X k = 0 (\ensuremath{-}1)k  j k  a(1\ensuremath{-}k+j)m \rule{1ex}{1ex} Prova: (Equação (1.17), p. 44) Tomemos j = m na equação
\prov{relações: prova gentil\_novas\_sequencias\_aritmeticas\_geometricas; prova gentil\_capitulo\_1\_sequencias\_aritmeticas\_de\_ordem\_m}
\prov{proveniência: gentil/Novas\_Sequencias\_Aritmeticas\_e\_Geometric.pdf \textperiodcentered{} fato \textperiodcentered{} confiança 1.0 \textperiodcentered{} domínio matematica \textperiodcentered{} história 4 versões no jornal}

%CRISTAL {"arestas":[{"alvo":"matematica","confianca":0.95,"epistemico":"fato","origem":"paper","rel":"parte_de"}],"confianca":1.0,"contraexemplos":[],"descricao":"Teorema do gene - I) Seja m um natural arbitrariamente fixado e seja j um natural tal que 0 ≤j ≤m. Nestas condições a seguinte identidade se verifica: ∆j anm = an(m−j) (2.6) Prova: Indução sobre j. Para j = 0 temos: ∆0 anm = anm = an(m−0) Suponhamos a proposição válida para j = p (0 ≤p < m): (H.I) ∆p anm = an(m−p) Provemos que vale para j = p + 1: (T.I) ∆p+1 anm = an(m−(p+1)) Então, ∆p+1 anm = ∆p a(n+1)m −∆p anm = a(n+1)(m−p) −an(m−p) = an(m−(p+1)) ■ Na última igualdade acima fizemos uso da terc","epistemico":"fato","exemplos":[],"id":"gentil_teorema_10_teorema_do_gene_-_i_seja_m_um_natural_arbitrariamente_fixad","memoria":"semantica","meta":{"arquivo":"gentil/Novas_Sequencias_Aritmeticas_e_Geometric.pdf","ciencia":"teoria_dos_numeros","dominio":"matematica","nivel":4,"numero":"10","tipo_no":"paper","tipo_teorema":"teorema"},"origem":"gentil/Novas_Sequencias_Aritmeticas_e_Geometric.pdf","palavras_chave":[],"sinonimos":[],"tipo":"conceito","titulo":"gentil Teorema 10 — Teorema do gene - I) Seja m um natural arbitrariamente fixad"}
\section{gentil Teorema 10 — Teorema do gene - I) Seja m um natural arbitrariamente fixad}
Teorema do gene - I) Seja m um natural arbitrariamente fixado e seja j um natural tal que 0 \ensuremath{\le}j \ensuremath{\le}m. Nestas condições a seguinte identidade se verifica: \ensuremath{\Delta}j anm = an(m\ensuremath{-}j) (2.6) Prova: Indução sobre j. Para j = 0 temos: \ensuremath{\Delta}0 anm = anm = an(m\ensuremath{-}0) Suponhamos a proposição válida para j = p (0 \ensuremath{\le}p < m): (H.I) \ensuremath{\Delta}p anm = an(m\ensuremath{-}p) Provemos que vale para j = p + 1: (T.I) \ensuremath{\Delta}p+1 anm = an(m\ensuremath{-}(p+1)) Então, \ensuremath{\Delta}p+1 anm = \ensuremath{\Delta}p a(n+1)m \ensuremath{-}\ensuremath{\Delta}p anm = a(n+1)(m\ensuremath{-}p) \ensuremath{-}an(m\ensuremath{-}p) = an(m\ensuremath{-}(p+1)) \rule{1ex}{1ex} Na última igualdade acima fizemos uso da terc
\prov{relações: parte\_de matematica}
\prov{proveniência: gentil/Novas\_Sequencias\_Aritmeticas\_e\_Geometric.pdf \textperiodcentered{} fato \textperiodcentered{} confiança 1.0 \textperiodcentered{} domínio matematica \textperiodcentered{} história 3 versões no jornal}

%CRISTAL {"arestas":[{"alvo":"gentil_novas_sequencias_aritmeticas_geometricas","confianca":1.0,"epistemico":"fato","origem":"gentil/Novas_Sequencias_Aritmeticas_e_Geometric.pdf","rel":"prova"},{"alvo":"gentil_capitulo_1_sequencias_aritmeticas_de_ordem_m","confianca":1.0,"epistemico":"fato","origem":"gentil/Novas_Sequencias_Aritmeticas_e_Geometric.pdf","rel":"prova"}],"confianca":1.0,"contraexemplos":[],"descricao":"p 99), o qual repetimos: Teorema 16 (Teorema do gene - II). Seja m um natural arbitrariamente fixado e j ∈N ∪ 0. Nestas condições a seguinte identidade se verifica: ∇j anm = an(m+j) Prova: Considere a equação ∆j anm = an(m−j). Então, ∆j anm = an(m−j) ⇒ ∆j an(m+j) = anm Aplicando o operador ∇j a ambos os membros desta última equação, temos ∇j ∆j an(m+j) = ∇j anm ⇒ ∇j anm = an(m+j) ■ O","epistemico":"fato","exemplos":[],"id":"gentil_teorema_12_p_99_o_qual_repetimos_teorema_16_teorema_do_gene_-_ii","memoria":"semantica","meta":{"arquivo":"gentil/Novas_Sequencias_Aritmeticas_e_Geometric.pdf","ciencia":"teoria_dos_numeros","dominio":"matematica","nivel":4,"numero":"12","tipo_no":"paper","tipo_teorema":"teorema"},"origem":"gentil/Novas_Sequencias_Aritmeticas_e_Geometric.pdf","palavras_chave":[],"sinonimos":[],"tipo":"conceito","titulo":"gentil Teorema 12 — p 99), o qual repetimos: Teorema 16 (Teorema do gene - II)."}
\section{gentil Teorema 12 — p 99), o qual repetimos: Teorema 16 (Teorema do gene - II).}
p 99), o qual repetimos: Teorema 16 (Teorema do gene - II). Seja m um natural arbitrariamente fixado e j \ensuremath{\in}N \ensuremath{\cup} 0. Nestas condições a seguinte identidade se verifica: \ensuremath{\nabla}j anm = an(m+j) Prova: Considere a equação \ensuremath{\Delta}j anm = an(m\ensuremath{-}j). Então, \ensuremath{\Delta}j anm = an(m\ensuremath{-}j) \ensuremath{\Rightarrow} \ensuremath{\Delta}j an(m+j) = anm Aplicando o operador \ensuremath{\nabla}j a ambos os membros desta última equação, temos \ensuremath{\nabla}j \ensuremath{\Delta}j an(m+j) = \ensuremath{\nabla}j anm \ensuremath{\Rightarrow} \ensuremath{\nabla}j anm = an(m+j) \rule{1ex}{1ex} O
\prov{relações: prova gentil\_novas\_sequencias\_aritmeticas\_geometricas; prova gentil\_capitulo\_1\_sequencias\_aritmeticas\_de\_ordem\_m}
\prov{proveniência: gentil/Novas\_Sequencias\_Aritmeticas\_e\_Geometric.pdf \textperiodcentered{} fato \textperiodcentered{} confiança 1.0 \textperiodcentered{} domínio matematica \textperiodcentered{} história 5 versões no jornal}

%CRISTAL {"arestas":[{"alvo":"matematica","confianca":0.95,"epistemico":"fato","origem":"paper","rel":"parte_de"}],"confianca":1.0,"contraexemplos":[],"descricao":"Linearidade do operador ∇m) O operador Soma de ordem m é linear. Isto é, a seguinte identidade se verifica: ∇m ( a f + b g )(n) = a ∇m f(n) + b ∇m g(n) para todo a, b ∈R. Prova: Indução sobre m. Para m = 0, pela","epistemico":"fato","exemplos":[],"id":"gentil_teorema_13_linearidade_do_operador_m_o_operador_soma_de_ordem_m_e_lin","memoria":"semantica","meta":{"arquivo":"gentil/Novas_Sequencias_Aritmeticas_e_Geometric.pdf","ciencia":"teoria_dos_numeros","dominio":"matematica","nivel":4,"numero":"13","tipo_no":"paper","tipo_teorema":"teorema"},"origem":"gentil/Novas_Sequencias_Aritmeticas_e_Geometric.pdf","palavras_chave":[],"sinonimos":[],"tipo":"conceito","titulo":"gentil Teorema 13 — Linearidade do operador ∇m) O operador Soma de ordem m é lin"}
\section{gentil Teorema 13 — Linearidade do operador \ensuremath{\nabla}m) O operador Soma de ordem m é lin}
Linearidade do operador \ensuremath{\nabla}m) O operador Soma de ordem m é linear. Isto é, a seguinte identidade se verifica: \ensuremath{\nabla}m ( a f + b g )(n) = a \ensuremath{\nabla}m f(n) + b \ensuremath{\nabla}m g(n) para todo a, b \ensuremath{\in}R. Prova: Indução sobre m. Para m = 0, pela
\prov{relações: parte\_de matematica}
\prov{proveniência: gentil/Novas\_Sequencias\_Aritmeticas\_e\_Geometric.pdf \textperiodcentered{} fato \textperiodcentered{} confiança 1.0 \textperiodcentered{} domínio matematica \textperiodcentered{} história 3 versões no jornal}

%CRISTAL {"arestas":[{"alvo":"gentil_novas_sequencias_aritmeticas_geometricas","confianca":1.0,"epistemico":"fato","origem":"gentil/Novas_Sequencias_Aritmeticas_e_Geometric.pdf","rel":"prova"},{"alvo":"gentil_capitulo_1_sequencias_aritmeticas_de_ordem_m","confianca":1.0,"epistemico":"fato","origem":"gentil/Novas_Sequencias_Aritmeticas_e_Geometric.pdf","rel":"prova"}],"confianca":1.0,"contraexemplos":[],"descricao":"Caracterização de uma P A.m ). Uma sequência (an) é uma P.A.m se, e somente se, ∆m an = c onde c é uma constante diferente de zero. Prova: (⇒) Se (an) é uma P.A.m , então an = anm (para algum m ∈N), logo aplicando o","epistemico":"fato","exemplos":[],"id":"gentil_teorema_17_caracterizacao_de_uma_p_am_uma_sequencia_an_e_uma_pa","memoria":"semantica","meta":{"arquivo":"gentil/Novas_Sequencias_Aritmeticas_e_Geometric.pdf","ciencia":"teoria_dos_numeros","dominio":"matematica","nivel":4,"numero":"17","tipo_no":"paper","tipo_teorema":"teorema"},"origem":"gentil/Novas_Sequencias_Aritmeticas_e_Geometric.pdf","palavras_chave":[],"sinonimos":[],"tipo":"conceito","titulo":"gentil Teorema 17 — Caracterização de uma P A.m ). Uma sequência (an) é uma P.A."}
\section{gentil Teorema 17 — Caracterização de uma P A.m ). Uma sequência (an) é uma P.A.}
Caracterização de uma P A.m ). Uma sequência (an) é uma P.A.m se, e somente se, \ensuremath{\Delta}m an = c onde c é uma constante diferente de zero. Prova: (\ensuremath{\Rightarrow}) Se (an) é uma P.A.m , então an = anm (para algum m \ensuremath{\in}N), logo aplicando o
\prov{relações: prova gentil\_novas\_sequencias\_aritmeticas\_geometricas; prova gentil\_capitulo\_1\_sequencias\_aritmeticas\_de\_ordem\_m}
\prov{proveniência: gentil/Novas\_Sequencias\_Aritmeticas\_e\_Geometric.pdf \textperiodcentered{} fato \textperiodcentered{} confiança 1.0 \textperiodcentered{} domínio matematica \textperiodcentered{} história 4 versões no jornal}

%CRISTAL {"arestas":[{"alvo":"gentil_novas_sequencias_aritmeticas_geometricas","confianca":1.0,"epistemico":"fato","origem":"gentil/Novas_Sequencias_Aritmeticas_e_Geometric.pdf","rel":"prova"},{"alvo":"gentil_capitulo_1_sequencias_aritmeticas_de_ordem_m","confianca":1.0,"epistemico":"fato","origem":"gentil/Novas_Sequencias_Aritmeticas_e_Geometric.pdf","rel":"prova"}],"confianca":1.0,"contraexemplos":[],"descricao":"p 111), a sequência dada é uma P.A.2 . (c) Mostre que a sequência dos cubos do naturais, dada abaixo, é uma P.A.3 . 13 23 33 43 53 63 . . . Solução: Faremos de dois modos: i) Tomando sucessivas diferenças: an : 1 8 27 64 125 216 . . . ∆1 an : 7 19 37 61 91 . . . ∆2 an : 12 18 24 30 . . . ∆3 an : 6 6 6 . . . Portanto, pelo","epistemico":"fato","exemplos":[],"id":"gentil_teorema_17_p_111_a_sequencia_dada_e_uma_pa2_c_mostre_que_a_sequ","memoria":"semantica","meta":{"arquivo":"gentil/Novas_Sequencias_Aritmeticas_e_Geometric.pdf","ciencia":"teoria_dos_numeros","dominio":"matematica","nivel":4,"numero":"17","tipo_no":"paper","tipo_teorema":"teorema"},"origem":"gentil/Novas_Sequencias_Aritmeticas_e_Geometric.pdf","palavras_chave":[],"sinonimos":[],"tipo":"conceito","titulo":"gentil Teorema 17 — p 111), a sequência dada é uma P.A.2 . (c) Mostre que a sequ"}
\section{gentil Teorema 17 — p 111), a sequência dada é uma P.A.2 . (c) Mostre que a sequ}
p 111), a sequência dada é uma P.A.2 . (c) Mostre que a sequência dos cubos do naturais, dada abaixo, é uma P.A.3 . 13 23 33 43 53 63 . . . Solução: Faremos de dois modos: i) Tomando sucessivas diferenças: an : 1 8 27 64 125 216 . . . \ensuremath{\Delta}1 an : 7 19 37 61 91 . . . \ensuremath{\Delta}2 an : 12 18 24 30 . . . \ensuremath{\Delta}3 an : 6 6 6 . . . Portanto, pelo
\prov{relações: prova gentil\_novas\_sequencias\_aritmeticas\_geometricas; prova gentil\_capitulo\_1\_sequencias\_aritmeticas\_de\_ordem\_m}
\prov{proveniência: gentil/Novas\_Sequencias\_Aritmeticas\_e\_Geometric.pdf \textperiodcentered{} fato \textperiodcentered{} confiança 1.0 \textperiodcentered{} domínio matematica \textperiodcentered{} história 4 versões no jornal}

%CRISTAL {"arestas":[{"alvo":"gentil_novas_sequencias_aritmeticas_geometricas","confianca":1.0,"epistemico":"fato","origem":"gentil/Novas_Sequencias_Aritmeticas_e_Geometric.pdf","rel":"prova"},{"alvo":"gentil_capitulo_1_sequencias_aritmeticas_de_ordem_m","confianca":1.0,"epistemico":"fato","origem":"gentil/Novas_Sequencias_Aritmeticas_e_Geometric.pdf","rel":"prova"}],"confianca":1.0,"contraexemplos":[],"descricao":"p 111), concluimos que a sequência dada é uma P.A.3 .","epistemico":"fato","exemplos":[],"id":"gentil_teorema_17_p_111_concluimos_que_a_sequencia_dada_e_uma_pa3","memoria":"semantica","meta":{"arquivo":"gentil/Novas_Sequencias_Aritmeticas_e_Geometric.pdf","ciencia":"teoria_dos_numeros","dominio":"matematica","nivel":4,"numero":"17","tipo_no":"paper","tipo_teorema":"teorema"},"origem":"gentil/Novas_Sequencias_Aritmeticas_e_Geometric.pdf","palavras_chave":[],"sinonimos":[],"tipo":"conceito","titulo":"gentil Teorema 17 — p 111), concluimos que a sequência dada é uma P.A.3 ."}
\section{gentil Teorema 17 — p 111), concluimos que a sequência dada é uma P.A.3 .}
p 111), concluimos que a sequência dada é uma P.A.3 .
\prov{relações: prova gentil\_novas\_sequencias\_aritmeticas\_geometricas; prova gentil\_capitulo\_1\_sequencias\_aritmeticas\_de\_ordem\_m}
\prov{proveniência: gentil/Novas\_Sequencias\_Aritmeticas\_e\_Geometric.pdf \textperiodcentered{} fato \textperiodcentered{} confiança 1.0 \textperiodcentered{} domínio matematica \textperiodcentered{} história 4 versões no jornal}

%CRISTAL {"arestas":[{"alvo":"gentil_novas_sequencias_aritmeticas_geometricas","confianca":1.0,"epistemico":"fato","origem":"gentil/Novas_Sequencias_Aritmeticas_e_Geometric.pdf","rel":"prova"},{"alvo":"gentil_capitulo_1_sequencias_aritmeticas_de_ordem_m","confianca":1.0,"epistemico":"fato","origem":"gentil/Novas_Sequencias_Aritmeticas_e_Geometric.pdf","rel":"prova"}],"confianca":1.0,"contraexemplos":[],"descricao":", etc , etc. A matemática é um campo demasiadamente árduo e inóspito para agradar àqueles a quem não ofe- rece grandes recompensas. Recompensas que são da mesma ´ındole que as do artista. . . . Acrescenta ainda que é no ato de criar que o matemático encontra sua culminância e que “nenhuma quantidade de trabalho ou correção técnica pode substituir este momento de criação na vida de um matemático, poeta ou músico”. (Norbert Wiener) 57","epistemico":"fato","exemplos":[],"id":"gentil_teorema_18_etc_etc_a_matematica_e_um_campo_demasiadamente_arduo_e","memoria":"semantica","meta":{"arquivo":"gentil/Novas_Sequencias_Aritmeticas_e_Geometric.pdf","ciencia":"teoria_dos_numeros","dominio":"matematica","nivel":4,"numero":"18","tipo_no":"paper","tipo_teorema":"teorema"},"origem":"gentil/Novas_Sequencias_Aritmeticas_e_Geometric.pdf","palavras_chave":[],"sinonimos":[],"tipo":"conceito","titulo":"gentil Teorema 18 — , etc , etc. A matemática é um campo demasiadamente árduo e"}
\section{gentil Teorema 18 — , etc , etc. A matemática é um campo demasiadamente árduo e}
, etc , etc. A matemática é um campo demasiadamente árduo e inóspito para agradar àqueles a quem não ofe- rece grandes recompensas. Recompensas que são da mesma '\ensuremath{\imath}ndole que as do artista. . . . Acrescenta ainda que é no ato de criar que o matemático encontra sua culminância e que “nenhuma quantidade de trabalho ou correção técnica pode substituir este momento de criação na vida de um matemático, poeta ou músico”. (Norbert Wiener) 57
\prov{relações: prova gentil\_novas\_sequencias\_aritmeticas\_geometricas; prova gentil\_capitulo\_1\_sequencias\_aritmeticas\_de\_ordem\_m}
\prov{proveniência: gentil/Novas\_Sequencias\_Aritmeticas\_e\_Geometric.pdf \textperiodcentered{} fato \textperiodcentered{} confiança 1.0 \textperiodcentered{} domínio matematica \textperiodcentered{} história 4 versões no jornal}

%CRISTAL {"arestas":[{"alvo":"gentil_novas_sequencias_aritmeticas_geometricas","confianca":1.0,"epistemico":"fato","origem":"gentil/Novas_Sequencias_Aritmeticas_e_Geometric.pdf","rel":"prova"},{"alvo":"gentil_capitulo_1_sequencias_aritmeticas_de_ordem_m","confianca":1.0,"epistemico":"fato","origem":"gentil/Novas_Sequencias_Aritmeticas_e_Geometric.pdf","rel":"prova"}],"confianca":1.0,"contraexemplos":[],"descricao":"p 113) para m = 3. Ou seja, para a sequência dada por an = n3, a seguinte identidade se verifica ∆n3 = 3! , n = 1, 2, 3, . . . Veja o significado desta identidade no seguinte diagrama: an : 1 8 27 64 125 . . . P.A.3 ∆1 an : 7 19 37 61 . . . P.A.2 ∆2 an : 12 18 24 . . . P.A.1 → ∆3 an : 6 6 . . . P.A.0 Dois corolários imediatos do","epistemico":"fato","exemplos":[],"id":"gentil_teorema_18_p_113_para_m_3_ou_seja_para_a_sequencia_dada_por_an_n","memoria":"semantica","meta":{"arquivo":"gentil/Novas_Sequencias_Aritmeticas_e_Geometric.pdf","ciencia":"teoria_dos_numeros","dominio":"matematica","nivel":4,"numero":"18","tipo_no":"paper","tipo_teorema":"teorema"},"origem":"gentil/Novas_Sequencias_Aritmeticas_e_Geometric.pdf","palavras_chave":[],"sinonimos":[],"tipo":"conceito","titulo":"gentil Teorema 18 — p 113) para m = 3. Ou seja, para a sequência dada por an = n"}
\section{gentil Teorema 18 — p 113) para m = 3. Ou seja, para a sequência dada por an = n}
p 113) para m = 3. Ou seja, para a sequência dada por an = n3, a seguinte identidade se verifica \ensuremath{\Delta}n3 = 3! , n = 1, 2, 3, . . . Veja o significado desta identidade no seguinte diagrama: an : 1 8 27 64 125 . . . P.A.3 \ensuremath{\Delta}1 an : 7 19 37 61 . . . P.A.2 \ensuremath{\Delta}2 an : 12 18 24 . . . P.A.1 \ensuremath{\to} \ensuremath{\Delta}3 an : 6 6 . . . P.A.0 Dois corolários imediatos do
\prov{relações: prova gentil\_novas\_sequencias\_aritmeticas\_geometricas; prova gentil\_capitulo\_1\_sequencias\_aritmeticas\_de\_ordem\_m}
\prov{proveniência: gentil/Novas\_Sequencias\_Aritmeticas\_e\_Geometric.pdf \textperiodcentered{} fato \textperiodcentered{} confiança 1.0 \textperiodcentered{} domínio matematica \textperiodcentered{} história 4 versões no jornal}

%CRISTAL {"arestas":[{"alvo":"gentil_novas_sequencias_aritmeticas_geometricas","confianca":1.0,"epistemico":"fato","origem":"gentil/Novas_Sequencias_Aritmeticas_e_Geometric.pdf","rel":"prova"}],"confianca":1.0,"contraexemplos":[],"descricao":"Fórmula do termo geral de uma P A.m ). Em uma P.A.m o n −ésimo termo vale anm = m X j = 0 \u0012 n −1 j \u0013 a1(m−j) (1.10) Nota: Esta fórmula do termo geral de uma P.A.m é inédita, isto é, não consta na literatura −Assim como muitas outras que ainda aparecerão oportunamente. Isto tudo é uma consequência de nossa","epistemico":"fato","exemplos":[],"id":"gentil_teorema_1_formula_do_termo_geral_de_uma_p_am_em_uma_pam_o_n_esi","memoria":"semantica","meta":{"arquivo":"gentil/Novas_Sequencias_Aritmeticas_e_Geometric.pdf","ciencia":"teoria_dos_numeros","dominio":"matematica","nivel":4,"numero":"1","tipo_no":"paper","tipo_teorema":"teorema"},"origem":"gentil/Novas_Sequencias_Aritmeticas_e_Geometric.pdf","palavras_chave":[],"sinonimos":[],"tipo":"conceito","titulo":"gentil Teorema 1 — Fórmula do termo geral de uma P A.m ). Em uma P.A.m o n −ési"}
\section{gentil Teorema 1 — Fórmula do termo geral de uma P A.m ). Em uma P.A.m o n \ensuremath{-}ési}
Fórmula do termo geral de uma P A.m ). Em uma P.A.m o n \ensuremath{-}ésimo termo vale anm = m X j = 0  n \ensuremath{-}1 j  a1(m\ensuremath{-}j) (1.10) Nota: Esta fórmula do termo geral de uma P.A.m é inédita, isto é, não consta na literatura \ensuremath{-}Assim como muitas outras que ainda aparecerão oportunamente. Isto tudo é uma consequência de nossa
\prov{relações: prova gentil\_novas\_sequencias\_aritmeticas\_geometricas}
\prov{proveniência: gentil/Novas\_Sequencias\_Aritmeticas\_e\_Geometric.pdf \textperiodcentered{} fato \textperiodcentered{} confiança 1.0 \textperiodcentered{} domínio matematica \textperiodcentered{} história 4 versões no jornal}

%CRISTAL {"arestas":[{"alvo":"gentil_novas_sequencias_aritmeticas_geometricas","confianca":1.0,"epistemico":"fato","origem":"gentil/Novas_Sequencias_Aritmeticas_e_Geometric.pdf","rel":"prova"},{"alvo":"gentil_capitulo_1_sequencias_aritmeticas_de_ordem_m","confianca":1.0,"epistemico":"fato","origem":"gentil/Novas_Sequencias_Aritmeticas_e_Geometric.pdf","rel":"prova"}],"confianca":1.0,"contraexemplos":[],"descricao":"Em uma P G.m é válida a seguinte identidade: a1(m−j) = jY k = 0 a (−1)k ( j k) (1−k+j)m (3.5) Prova: Princ´ıpio da Dualidade. Esta é a equação dual de (1.16), p. 40. ■ −P.G.m em função dos seus próprios termos na HP Prime Na tela a seguir a1(m−j) = jY k = 0 a (−1)k ( j k) (1−k+j)m temos o programa que implementa a fórmula da direita. Como ilustração, considerando o exemplo da página 136, entrando com os quatro primeiros termos em um vetor, 1 1 1 −1 1 1 −1 1 −1 −1 a10 → a11 → a12 → a13 → o progra","epistemico":"fato","exemplos":[],"id":"gentil_teorema_21_em_uma_p_gm_e_valida_a_seguinte_identidade_a1mj_jy_k","memoria":"semantica","meta":{"arquivo":"gentil/Novas_Sequencias_Aritmeticas_e_Geometric.pdf","ciencia":"teoria_dos_numeros","dominio":"matematica","nivel":4,"numero":"21","tipo_no":"paper","tipo_teorema":"teorema"},"origem":"gentil/Novas_Sequencias_Aritmeticas_e_Geometric.pdf","palavras_chave":[],"sinonimos":[],"tipo":"conceito","titulo":"gentil Teorema 21 — Em uma P G.m é válida a seguinte identidade: a1(m−j) = jY k"}
\section{gentil Teorema 21 — Em uma P G.m é válida a seguinte identidade: a1(m\ensuremath{-}j) = jY k}
Em uma P G.m é válida a seguinte identidade: a1(m\ensuremath{-}j) = jY k = 0 a (\ensuremath{-}1)k ( j k) (1\ensuremath{-}k+j)m (3.5) Prova: Princ'\ensuremath{\imath}pio da Dualidade. Esta é a equação dual de (1.16), p. 40. \rule{1ex}{1ex} \ensuremath{-}P.G.m em função dos seus próprios termos na HP Prime Na tela a seguir a1(m\ensuremath{-}j) = jY k = 0 a (\ensuremath{-}1)k ( j k) (1\ensuremath{-}k+j)m temos o programa que implementa a fórmula da direita. Como ilustração, considerando o exemplo da página 136, entrando com os quatro primeiros termos em um vetor, 1 1 1 \ensuremath{-}1 1 1 \ensuremath{-}1 1 \ensuremath{-}1 \ensuremath{-}1 a10 \ensuremath{\to} a11 \ensuremath{\to} a12 \ensuremath{\to} a13 \ensuremath{\to} o progra
\prov{relações: prova gentil\_novas\_sequencias\_aritmeticas\_geometricas; prova gentil\_capitulo\_1\_sequencias\_aritmeticas\_de\_ordem\_m}
\prov{proveniência: gentil/Novas\_Sequencias\_Aritmeticas\_e\_Geometric.pdf \textperiodcentered{} fato \textperiodcentered{} confiança 1.0 \textperiodcentered{} domínio matematica \textperiodcentered{} história 4 versões no jornal}

%CRISTAL {"arestas":[{"alvo":"gentil_novas_sequencias_aritmeticas_geometricas","confianca":1.0,"epistemico":"fato","origem":"gentil/Novas_Sequencias_Aritmeticas_e_Geometric.pdf","rel":"prova"},{"alvo":"gentil_capitulo_1_sequencias_aritmeticas_de_ordem_m","confianca":1.0,"epistemico":"fato","origem":"gentil/Novas_Sequencias_Aritmeticas_e_Geometric.pdf","rel":"prova"}],"confianca":1.0,"contraexemplos":[],"descricao":"Produto dos termos de uma P G.m ). Em uma P.G.m o produto Pnm dos n termos iniciais vale Pnm = m Y j = 0 a( n j+1) 1(m−j) (3.7) Prova: Princ´ıpio da Dualidade. Esta equação é a dual da (1.18), p. 47. ■ Exemplo: Encontre o produto dos n termos iniciais da seguinte P.G.2 : 1 1 −1 −1 1 1 −1 −1 . . . Solução: Substituindo m = 2 na fórmula (3.7), obtemos: Pn2 = 2 Y j = 0 a( n j+1) 1(2−j) Simplificando, obtemos: Pn2 = a( n 1 ) 12 · a( n 2 ) 11 · a( n 3 ) 10 139\nDo algoritmo retiramos os coeficientes,","epistemico":"fato","exemplos":[],"id":"gentil_teorema_23_produto_dos_termos_de_uma_p_gm_em_uma_pgm_o_produto_pn","memoria":"semantica","meta":{"arquivo":"gentil/Novas_Sequencias_Aritmeticas_e_Geometric.pdf","ciencia":"teoria_dos_numeros","dominio":"matematica","nivel":4,"numero":"23","tipo_no":"paper","tipo_teorema":"teorema"},"origem":"gentil/Novas_Sequencias_Aritmeticas_e_Geometric.pdf","palavras_chave":[],"sinonimos":[],"tipo":"conceito","titulo":"gentil Teorema 23 — Produto dos termos de uma P G.m ). Em uma P.G.m o produto Pn"}
\section{gentil Teorema 23 — Produto dos termos de uma P G.m ). Em uma P.G.m o produto Pn}
Produto dos termos de uma P G.m ). Em uma P.G.m o produto Pnm dos n termos iniciais vale Pnm = m Y j = 0 a( n j+1) 1(m\ensuremath{-}j) (3.7) Prova: Princ'\ensuremath{\imath}pio da Dualidade. Esta equação é a dual da (1.18), p. 47. \rule{1ex}{1ex} Exemplo: Encontre o produto dos n termos iniciais da seguinte P.G.2 : 1 1 \ensuremath{-}1 \ensuremath{-}1 1 1 \ensuremath{-}1 \ensuremath{-}1 . . . Solução: Substituindo m = 2 na fórmula (3.7), obtemos: Pn2 = 2 Y j = 0 a( n j+1) 1(2\ensuremath{-}j) Simplificando, obtemos: Pn2 = a( n 1 ) 12 \textperiodcentered  a( n 2 ) 11 \textperiodcentered  a( n 3 ) 10 139
Do algoritmo retiramos os coeficientes,
\prov{relações: prova gentil\_novas\_sequencias\_aritmeticas\_geometricas; prova gentil\_capitulo\_1\_sequencias\_aritmeticas\_de\_ordem\_m}
\prov{proveniência: gentil/Novas\_Sequencias\_Aritmeticas\_e\_Geometric.pdf \textperiodcentered{} fato \textperiodcentered{} confiança 1.0 \textperiodcentered{} domínio matematica \textperiodcentered{} história 4 versões no jornal}

%CRISTAL {"arestas":[{"alvo":"gentil_novas_sequencias_aritmeticas_geometricas","confianca":1.0,"epistemico":"fato","origem":"gentil/Novas_Sequencias_Aritmeticas_e_Geometric.pdf","rel":"prova"},{"alvo":"gentil_capitulo_1_sequencias_aritmeticas_de_ordem_m","confianca":1.0,"epistemico":"fato","origem":"gentil/Novas_Sequencias_Aritmeticas_e_Geometric.pdf","rel":"prova"}],"confianca":1.0,"contraexemplos":[],"descricao":"As seguintes identidades são válidas: (i) ∆m f × g (n) = ∆m f(n) × ∆m g(n) (ii) ∆m f/g (n) = ∆m f(n) / ∆m g(n) Prova: Apêndice. (p. 177) ■ Corolário 4. ´E válida a seguinte identidade: ∆m f(n) = ∆m−1 ∆f(n) Prova: Princ´ıpio da Dualidade. ■ O seguinte","epistemico":"fato","exemplos":[],"id":"gentil_teorema_25_as_seguintes_identidades_sao_validas_i_m_f_g_n_m","memoria":"semantica","meta":{"arquivo":"gentil/Novas_Sequencias_Aritmeticas_e_Geometric.pdf","ciencia":"teoria_dos_numeros","dominio":"matematica","nivel":4,"numero":"25","tipo_no":"paper","tipo_teorema":"teorema"},"origem":"gentil/Novas_Sequencias_Aritmeticas_e_Geometric.pdf","palavras_chave":[],"sinonimos":[],"tipo":"conceito","titulo":"gentil Teorema 25 — As seguintes identidades são válidas: (i) ∆m f × g (n) = ∆m"}
\section{gentil Teorema 25 — As seguintes identidades são válidas: (i) \ensuremath{\Delta}m f $\times$ g (n) = \ensuremath{\Delta}m}
As seguintes identidades são válidas: (i) \ensuremath{\Delta}m f $\times$ g (n) = \ensuremath{\Delta}m f(n) $\times$ \ensuremath{\Delta}m g(n) (ii) \ensuremath{\Delta}m f/g (n) = \ensuremath{\Delta}m f(n) / \ensuremath{\Delta}m g(n) Prova: Apêndice. (p. 177) \rule{1ex}{1ex} Corolário 4. 'E válida a seguinte identidade: \ensuremath{\Delta}m f(n) = \ensuremath{\Delta}m\ensuremath{-}1 \ensuremath{\Delta}f(n) Prova: Princ'\ensuremath{\imath}pio da Dualidade. \rule{1ex}{1ex} O seguinte
\prov{relações: prova gentil\_novas\_sequencias\_aritmeticas\_geometricas; prova gentil\_capitulo\_1\_sequencias\_aritmeticas\_de\_ordem\_m}
\prov{proveniência: gentil/Novas\_Sequencias\_Aritmeticas\_e\_Geometric.pdf \textperiodcentered{} fato \textperiodcentered{} confiança 1.0 \textperiodcentered{} domínio matematica \textperiodcentered{} história 4 versões no jornal}

%CRISTAL {"arestas":[{"alvo":"matematica","confianca":0.95,"epistemico":"fato","origem":"paper","rel":"parte_de"}],"confianca":1.0,"contraexemplos":[],"descricao":"Teorema do gene - III) Seja m um natural arbitrariamente fixado e seja j um natural tal que 0 ≤j ≤m. Nestas condições a seguinte identidade se verifica: ∆j anm = an(m−j) (3.11) Prova: Princ´ıpio da Dualidade. ■ Este","epistemico":"fato","exemplos":[],"id":"gentil_teorema_26_teorema_do_gene_-_iii_seja_m_um_natural_arbitrariamente_fix","memoria":"semantica","meta":{"arquivo":"gentil/Novas_Sequencias_Aritmeticas_e_Geometric.pdf","ciencia":"teoria_dos_numeros","dominio":"matematica","nivel":4,"numero":"26","tipo_no":"paper","tipo_teorema":"teorema"},"origem":"gentil/Novas_Sequencias_Aritmeticas_e_Geometric.pdf","palavras_chave":[],"sinonimos":[],"tipo":"conceito","titulo":"gentil Teorema 26 — Teorema do gene - III) Seja m um natural arbitrariamente fix"}
\section{gentil Teorema 26 — Teorema do gene - III) Seja m um natural arbitrariamente fix}
Teorema do gene - III) Seja m um natural arbitrariamente fixado e seja j um natural tal que 0 \ensuremath{\le}j \ensuremath{\le}m. Nestas condições a seguinte identidade se verifica: \ensuremath{\Delta}j anm = an(m\ensuremath{-}j) (3.11) Prova: Princ'\ensuremath{\imath}pio da Dualidade. \rule{1ex}{1ex} Este
\prov{relações: parte\_de matematica}
\prov{proveniência: gentil/Novas\_Sequencias\_Aritmeticas\_e\_Geometric.pdf \textperiodcentered{} fato \textperiodcentered{} confiança 1.0 \textperiodcentered{} domínio matematica \textperiodcentered{} história 3 versões no jornal}

%CRISTAL {"arestas":[{"alvo":"gentil_novas_sequencias_aritmeticas_geometricas","confianca":1.0,"epistemico":"fato","origem":"gentil/Novas_Sequencias_Aritmeticas_e_Geometric.pdf","rel":"prova"},{"alvo":"gentil_capitulo_1_sequencias_aritmeticas_de_ordem_m","confianca":1.0,"epistemico":"fato","origem":"gentil/Novas_Sequencias_Aritmeticas_e_Geometric.pdf","rel":"prova"}],"confianca":1.0,"contraexemplos":[],"descricao":"As seguintes identidades são válidas: (i) ∆ m f × g (n) = ∆ m f(n) × ∆ m g(n) (ii) ∆ m f/g (n) = ∆ m f(n) / ∆ m g(n) Prova: Apêndice. (p. 178) ■","epistemico":"fato","exemplos":[],"id":"gentil_teorema_28_as_seguintes_identidades_sao_validas_i_m_f_g_n","memoria":"semantica","meta":{"arquivo":"gentil/Novas_Sequencias_Aritmeticas_e_Geometric.pdf","ciencia":"teoria_dos_numeros","dominio":"matematica","nivel":4,"numero":"28","tipo_no":"paper","tipo_teorema":"teorema"},"origem":"gentil/Novas_Sequencias_Aritmeticas_e_Geometric.pdf","palavras_chave":[],"sinonimos":[],"tipo":"conceito","titulo":"gentil Teorema 28 — As seguintes identidades são válidas: (i) ∆ m f × g (n) = ∆"}
\section{gentil Teorema 28 — As seguintes identidades são válidas: (i) \ensuremath{\Delta} m f $\times$ g (n) = \ensuremath{\Delta}}
As seguintes identidades são válidas: (i) \ensuremath{\Delta} m f $\times$ g (n) = \ensuremath{\Delta} m f(n) $\times$ \ensuremath{\Delta} m g(n) (ii) \ensuremath{\Delta} m f/g (n) = \ensuremath{\Delta} m f(n) / \ensuremath{\Delta} m g(n) Prova: Apêndice. (p. 178) \rule{1ex}{1ex}
\prov{relações: prova gentil\_novas\_sequencias\_aritmeticas\_geometricas; prova gentil\_capitulo\_1\_sequencias\_aritmeticas\_de\_ordem\_m}
\prov{proveniência: gentil/Novas\_Sequencias\_Aritmeticas\_e\_Geometric.pdf \textperiodcentered{} fato \textperiodcentered{} confiança 1.0 \textperiodcentered{} domínio matematica \textperiodcentered{} história 4 versões no jornal}

%CRISTAL {"arestas":[{"alvo":"matematica","confianca":0.95,"epistemico":"fato","origem":"paper","rel":"parte_de"}],"confianca":1.0,"contraexemplos":[],"descricao":"Teorema Q - P Fundamental - I) Seja m um natural arbitrariamente fixado e j um natural tal que 0 ≤j ≤m. Nestas condições a seguinte identidade se verifica: ∆ j ∆m f(n) = ∆m−j f(n) (3.13) 149\nProva: Princ´ıpio da Dualidade. (p. 102) ■","epistemico":"fato","exemplos":[],"id":"gentil_teorema_29_teorema_q_-_p_fundamental_-_i_seja_m_um_natural_arbitrariam","memoria":"semantica","meta":{"arquivo":"gentil/Novas_Sequencias_Aritmeticas_e_Geometric.pdf","ciencia":"teoria_dos_numeros","dominio":"matematica","nivel":4,"numero":"29","tipo_no":"paper","tipo_teorema":"teorema"},"origem":"gentil/Novas_Sequencias_Aritmeticas_e_Geometric.pdf","palavras_chave":[],"sinonimos":[],"tipo":"conceito","titulo":"gentil Teorema 29 — Teorema Q - P Fundamental - I) Seja m um natural arbitrariam"}
\section{gentil Teorema 29 — Teorema Q - P Fundamental - I) Seja m um natural arbitrariam}
Teorema Q - P Fundamental - I) Seja m um natural arbitrariamente fixado e j um natural tal que 0 \ensuremath{\le}j \ensuremath{\le}m. Nestas condições a seguinte identidade se verifica: \ensuremath{\Delta} j \ensuremath{\Delta}m f(n) = \ensuremath{\Delta}m\ensuremath{-}j f(n) (3.13) 149
Prova: Princ'\ensuremath{\imath}pio da Dualidade. (p. 102) \rule{1ex}{1ex}
\prov{relações: parte\_de matematica}
\prov{proveniência: gentil/Novas\_Sequencias\_Aritmeticas\_e\_Geometric.pdf \textperiodcentered{} fato \textperiodcentered{} confiança 1.0 \textperiodcentered{} domínio matematica \textperiodcentered{} história 3 versões no jornal}

%CRISTAL {"arestas":[{"alvo":"matematica","confianca":0.95,"epistemico":"fato","origem":"paper","rel":"parte_de"}],"confianca":1.0,"contraexemplos":[],"descricao":"Seja j um natural arbitrariamente fixado Para n ≥j vale a seguinte identidade: n X i = j \u0012 i j \u0013 = \u0012 n + 1 j + 1 \u0013 (1.11) Prova: Apêndice, página 71. ■ Exemplos: (a) Tomando j = 1 na equação (1.11), obtemos: n X i = 1 \u0012 i 1 \u0013 = \u0012 n + 1 1 + 1 \u0013 Isto é, \u0012 1 1 \u0013 + \u0012 2 1 \u0013 + · · · + \u0012 n 1 \u0013 = \u0012 n + 1 2 \u0013 Ou ainda, 1 + 2 + · · · + n = n(n + 1) 2 (b) Tomando j = 2 na equação (1.11), obtemos: n X i = 2 \u0012 i 2 \u0013 = \u0012 n + 1 2 + 1 \u0013 Isto é, \u0012 1 2 \u0013 + \u0012 2 2 \u0013 + \u0012 3 2 \u0013 + · · · + \u0012 n 2 \u0013 = \u0012 n + 1 3 \u0013 Ou aind","epistemico":"fato","exemplos":[],"id":"gentil_teorema_2_seja_j_um_natural_arbitrariamente_fixado_para_n_j_vale_a_se","memoria":"semantica","meta":{"arquivo":"gentil/Novas_Sequencias_Aritmeticas_e_Geometric.pdf","ciencia":"teoria_dos_numeros","dominio":"matematica","nivel":4,"numero":"2","tipo_no":"paper","tipo_teorema":"teorema"},"origem":"gentil/Novas_Sequencias_Aritmeticas_e_Geometric.pdf","palavras_chave":[],"sinonimos":[],"tipo":"conceito","titulo":"gentil Teorema 2 — Seja j um natural arbitrariamente fixado Para n ≥j vale a se"}
\section{gentil Teorema 2 — Seja j um natural arbitrariamente fixado Para n \ensuremath{\ge}j vale a se}
Seja j um natural arbitrariamente fixado Para n \ensuremath{\ge}j vale a seguinte identidade: n X i = j  i j  =  n + 1 j + 1  (1.11) Prova: Apêndice, página 71. \rule{1ex}{1ex} Exemplos: (a) Tomando j = 1 na equação (1.11), obtemos: n X i = 1  i 1  =  n + 1 1 + 1  Isto é,  1 1  +  2 1  + \textperiodcentered  \textperiodcentered  \textperiodcentered  +  n 1  =  n + 1 2  Ou ainda, 1 + 2 + \textperiodcentered  \textperiodcentered  \textperiodcentered  + n = n(n + 1) 2 (b) Tomando j = 2 na equação (1.11), obtemos: n X i = 2  i 2  =  n + 1 2 + 1  Isto é,  1 2  +  2 2  +  3 2  + \textperiodcentered  \textperiodcentered  \textperiodcentered  +  n 2  =  n + 1 3  Ou aind
\prov{relações: parte\_de matematica}
\prov{proveniência: gentil/Novas\_Sequencias\_Aritmeticas\_e\_Geometric.pdf \textperiodcentered{} fato \textperiodcentered{} confiança 1.0 \textperiodcentered{} domínio matematica \textperiodcentered{} história 3 versões no jornal}

%CRISTAL {"arestas":[{"alvo":"matematica","confianca":0.95,"epistemico":"fato","origem":"paper","rel":"parte_de"}],"confianca":1.0,"contraexemplos":[],"descricao":"Teorema Q - P Fundamental - II) Seja m um natural arbitraria- mente fixado e j um natural tal que 0 ≤j ≤m. Nestas condições a seguinte identidade se verifica: ∆j ∆ m f(n) = ∆m−j f(n) (3.14) Prova: Princ´ıpio da Dualidade. (p. 105) ■ O","epistemico":"fato","exemplos":[],"id":"gentil_teorema_30_teorema_q_-_p_fundamental_-_ii_seja_m_um_natural_arbitraria","memoria":"semantica","meta":{"arquivo":"gentil/Novas_Sequencias_Aritmeticas_e_Geometric.pdf","ciencia":"teoria_dos_numeros","dominio":"matematica","nivel":4,"numero":"30","tipo_no":"paper","tipo_teorema":"teorema"},"origem":"gentil/Novas_Sequencias_Aritmeticas_e_Geometric.pdf","palavras_chave":[],"sinonimos":[],"tipo":"conceito","titulo":"gentil Teorema 30 — Teorema Q - P Fundamental - II) Seja m um natural arbitraria"}
\section{gentil Teorema 30 — Teorema Q - P Fundamental - II) Seja m um natural arbitraria}
Teorema Q - P Fundamental - II) Seja m um natural arbitraria- mente fixado e j um natural tal que 0 \ensuremath{\le}j \ensuremath{\le}m. Nestas condições a seguinte identidade se verifica: \ensuremath{\Delta}j \ensuremath{\Delta} m f(n) = \ensuremath{\Delta}m\ensuremath{-}j f(n) (3.14) Prova: Princ'\ensuremath{\imath}pio da Dualidade. (p. 105) \rule{1ex}{1ex} O
\prov{relações: parte\_de matematica}
\prov{proveniência: gentil/Novas\_Sequencias\_Aritmeticas\_e\_Geometric.pdf \textperiodcentered{} fato \textperiodcentered{} confiança 1.0 \textperiodcentered{} domínio matematica \textperiodcentered{} história 3 versões no jornal}

%CRISTAL {"arestas":[{"alvo":"matematica","confianca":0.95,"epistemico":"fato","origem":"paper","rel":"parte_de"}],"confianca":1.0,"contraexemplos":[],"descricao":"Teorema do gene - IV) Seja m um natural arbitrariamente fixado e j ∈N ∪ 0. Nestas condições a seguinte identidade se verifica: ∆ j anm = an(m+j) (3.15) Prova: Princ´ıpio da Dualidade. (p. 99) ■ Com o aux´ılio dos","epistemico":"fato","exemplos":[],"id":"gentil_teorema_31_teorema_do_gene_-_iv_seja_m_um_natural_arbitrariamente_fixa","memoria":"semantica","meta":{"arquivo":"gentil/Novas_Sequencias_Aritmeticas_e_Geometric.pdf","ciencia":"teoria_dos_numeros","dominio":"matematica","nivel":4,"numero":"31","tipo_no":"paper","tipo_teorema":"teorema"},"origem":"gentil/Novas_Sequencias_Aritmeticas_e_Geometric.pdf","palavras_chave":[],"sinonimos":[],"tipo":"conceito","titulo":"gentil Teorema 31 — Teorema do gene - IV) Seja m um natural arbitrariamente fixa"}
\section{gentil Teorema 31 — Teorema do gene - IV) Seja m um natural arbitrariamente fixa}
Teorema do gene - IV) Seja m um natural arbitrariamente fixado e j \ensuremath{\in}N \ensuremath{\cup} 0. Nestas condições a seguinte identidade se verifica: \ensuremath{\Delta} j anm = an(m+j) (3.15) Prova: Princ'\ensuremath{\imath}pio da Dualidade. (p. 99) \rule{1ex}{1ex} Com o aux'\ensuremath{\imath}lio dos
\prov{relações: parte\_de matematica}
\prov{proveniência: gentil/Novas\_Sequencias\_Aritmeticas\_e\_Geometric.pdf \textperiodcentered{} fato \textperiodcentered{} confiança 1.0 \textperiodcentered{} domínio matematica \textperiodcentered{} história 4 versões no jornal}

%CRISTAL {"arestas":[{"alvo":"gentil_novas_sequencias_aritmeticas_geometricas","confianca":1.0,"epistemico":"fato","origem":"gentil/Novas_Sequencias_Aritmeticas_e_Geometric.pdf","rel":"prova"},{"alvo":"gentil_capitulo_1_sequencias_aritmeticas_de_ordem_m","confianca":1.0,"epistemico":"fato","origem":"gentil/Novas_Sequencias_Aritmeticas_e_Geometric.pdf","rel":"prova"}],"confianca":1.0,"contraexemplos":[],"descricao":"Caracterização de uma P G.m ). Uma sequência (an) é uma P.G.m se, e somente se, ∆m an = c onde c é uma constante diferente de 1. Prova: Princ´ıpio da Dualidade. (p. 111) ■ Comentário: Este","epistemico":"fato","exemplos":[],"id":"gentil_teorema_32_caracterizacao_de_uma_p_gm_uma_sequencia_an_e_uma_pg","memoria":"semantica","meta":{"arquivo":"gentil/Novas_Sequencias_Aritmeticas_e_Geometric.pdf","ciencia":"teoria_dos_numeros","dominio":"matematica","nivel":4,"numero":"32","tipo_no":"paper","tipo_teorema":"teorema"},"origem":"gentil/Novas_Sequencias_Aritmeticas_e_Geometric.pdf","palavras_chave":[],"sinonimos":[],"tipo":"conceito","titulo":"gentil Teorema 32 — Caracterização de uma P G.m ). Uma sequência (an) é uma P.G."}
\section{gentil Teorema 32 — Caracterização de uma P G.m ). Uma sequência (an) é uma P.G.}
Caracterização de uma P G.m ). Uma sequência (an) é uma P.G.m se, e somente se, \ensuremath{\Delta}m an = c onde c é uma constante diferente de 1. Prova: Princ'\ensuremath{\imath}pio da Dualidade. (p. 111) \rule{1ex}{1ex} Comentário: Este
\prov{relações: prova gentil\_novas\_sequencias\_aritmeticas\_geometricas; prova gentil\_capitulo\_1\_sequencias\_aritmeticas\_de\_ordem\_m}
\prov{proveniência: gentil/Novas\_Sequencias\_Aritmeticas\_e\_Geometric.pdf \textperiodcentered{} fato \textperiodcentered{} confiança 1.0 \textperiodcentered{} domínio matematica \textperiodcentered{} história 4 versões no jornal}

%CRISTAL {"arestas":[{"alvo":"gentil_novas_sequencias_aritmeticas_geometricas","confianca":1.0,"epistemico":"fato","origem":"gentil/Novas_Sequencias_Aritmeticas_e_Geometric.pdf","rel":"prova"},{"alvo":"gentil_capitulo_1_sequencias_aritmeticas_de_ordem_m","confianca":1.0,"epistemico":"fato","origem":"gentil/Novas_Sequencias_Aritmeticas_e_Geometric.pdf","rel":"prova"}],"confianca":1.0,"contraexemplos":[],"descricao":"A sequência ( an ) dada pela fórmula abaixo an = ( −1 )( n−1 2j−1 ) é uma P G.2 j−1 , para cada j natural fixo. 169\nProva: Tome m = q = 2j−1, no","epistemico":"fato","exemplos":[],"id":"gentil_teorema_33_a_sequencia_an_dada_pela_formula_abaixo_an_1_n1","memoria":"semantica","meta":{"arquivo":"gentil/Novas_Sequencias_Aritmeticas_e_Geometric.pdf","ciencia":"teoria_dos_numeros","dominio":"matematica","nivel":4,"numero":"33","tipo_no":"paper","tipo_teorema":"teorema"},"origem":"gentil/Novas_Sequencias_Aritmeticas_e_Geometric.pdf","palavras_chave":[],"sinonimos":[],"tipo":"conceito","titulo":"gentil Teorema 33 — A sequência ( an ) dada pela fórmula abaixo an = ( −1 )( n−1"}
\section{gentil Teorema 33 — A sequência ( an ) dada pela fórmula abaixo an = ( \ensuremath{-}1 )( n\ensuremath{-}1}
A sequência ( an ) dada pela fórmula abaixo an = ( \ensuremath{-}1 )( n\ensuremath{-}1 2j\ensuremath{-}1 ) é uma P G.2 j\ensuremath{-}1 , para cada j natural fixo. 169
Prova: Tome m = q = 2j\ensuremath{-}1, no
\prov{relações: prova gentil\_novas\_sequencias\_aritmeticas\_geometricas; prova gentil\_capitulo\_1\_sequencias\_aritmeticas\_de\_ordem\_m}
\prov{proveniência: gentil/Novas\_Sequencias\_Aritmeticas\_e\_Geometric.pdf \textperiodcentered{} fato \textperiodcentered{} confiança 1.0 \textperiodcentered{} domínio matematica \textperiodcentered{} história 4 versões no jornal}

%CRISTAL {"arestas":[{"alvo":"gentil_novas_sequencias_aritmeticas_geometricas","confianca":1.0,"epistemico":"fato","origem":"gentil/Novas_Sequencias_Aritmeticas_e_Geometric.pdf","rel":"prova"},{"alvo":"gentil_capitulo_1_sequencias_aritmeticas_de_ordem_m","confianca":1.0,"epistemico":"fato","origem":"gentil/Novas_Sequencias_Aritmeticas_e_Geometric.pdf","rel":"prova"}],"confianca":1.0,"contraexemplos":[],"descricao":"Algoritmo da Divisão) Para quaisquer a, b ∈N, b ̸= 0, existe um único par de números naturais q e r, de maneira que a = b · q + r com 0 ≤r < b. Prova: ver [3]. ■ Facilmente podemos mostrar que o quociente é a parte inteira de uma divisão, veja: Corolário 6. Sejam a, b ∈N com b > 0 e q o quociente da divisão de a por b. Então q = ⌊a b ⌋. Prova: De fato, sendo a = b · q + r com 0 ≤r < b, dividindo por b > 0, temos a b = q + r b com 0 ≤r b < 1 (3.22) Sendo assim, temos: j a b k = j q + r b k = q +","epistemico":"fato","exemplos":[],"id":"gentil_teorema_35_algoritmo_da_divisao_para_quaisquer_a_b_n_b_0_existe","memoria":"semantica","meta":{"arquivo":"gentil/Novas_Sequencias_Aritmeticas_e_Geometric.pdf","ciencia":"teoria_dos_numeros","dominio":"matematica","nivel":4,"numero":"35","tipo_no":"paper","tipo_teorema":"teorema"},"origem":"gentil/Novas_Sequencias_Aritmeticas_e_Geometric.pdf","palavras_chave":[],"sinonimos":[],"tipo":"conceito","titulo":"gentil Teorema 35 — Algoritmo da Divisão) Para quaisquer a, b ∈N, b ̸= 0, existe"}
\section{gentil Teorema 35 — Algoritmo da Divisão) Para quaisquer a, b \ensuremath{\in}N, b \ensuremath{}= 0, existe}
Algoritmo da Divisão) Para quaisquer a, b \ensuremath{\in}N, b \ensuremath{}= 0, existe um único par de números naturais q e r, de maneira que a = b \textperiodcentered  q + r com 0 \ensuremath{\le}r < b. Prova: ver [3]. \rule{1ex}{1ex} Facilmente podemos mostrar que o quociente é a parte inteira de uma divisão, veja: Corolário 6. Sejam a, b \ensuremath{\in}N com b > 0 e q o quociente da divisão de a por b. Então q = \ensuremath{\lfloor}a b \ensuremath{\rfloor}. Prova: De fato, sendo a = b \textperiodcentered  q + r com 0 \ensuremath{\le}r < b, dividindo por b > 0, temos a b = q + r b com 0 \ensuremath{\le}r b < 1 (3.22) Sendo assim, temos: j a b k = j q + r b k = q +
\prov{relações: prova gentil\_novas\_sequencias\_aritmeticas\_geometricas; prova gentil\_capitulo\_1\_sequencias\_aritmeticas\_de\_ordem\_m}
\prov{proveniência: gentil/Novas\_Sequencias\_Aritmeticas\_e\_Geometric.pdf \textperiodcentered{} fato \textperiodcentered{} confiança 1.0 \textperiodcentered{} domínio matematica \textperiodcentered{} história 4 versões no jornal}

%CRISTAL {"arestas":[{"alvo":"matematica","confianca":0.95,"epistemico":"fato","origem":"paper","rel":"parte_de"}],"confianca":1.0,"contraexemplos":[],"descricao":"Fórmula do termo geral de uma PA-2D) Na PA-2D em que o pri- meiro termo é a11, a razão das linhas é r1 e a razão das colunas é r2 o (i, j)-termo é: aij = a11 + (j −1) r1 + (i −1) r2 (5.1) 232\nProva: Indução sobre i (j fixo). Para i = 1 precisamos mostrar que, a1j = a11 + (j −1) r1 De fato, para i = 1 a","epistemico":"fato","exemplos":[],"id":"gentil_teorema_38_formula_do_termo_geral_de_uma_pa-2d_na_pa-2d_em_que_o_pri-","memoria":"semantica","meta":{"arquivo":"gentil/Novas_Sequencias_Aritmeticas_e_Geometric.pdf","ciencia":"teoria_dos_numeros","dominio":"matematica","nivel":4,"numero":"38","tipo_no":"paper","tipo_teorema":"teorema"},"origem":"gentil/Novas_Sequencias_Aritmeticas_e_Geometric.pdf","palavras_chave":[],"sinonimos":[],"tipo":"conceito","titulo":"gentil Teorema 38 — Fórmula do termo geral de uma PA-2D) Na PA-2D em que o pri-"}
\section{gentil Teorema 38 — Fórmula do termo geral de uma PA-2D) Na PA-2D em que o pri-}
Fórmula do termo geral de uma PA-2D) Na PA-2D em que o pri- meiro termo é a11, a razão das linhas é r1 e a razão das colunas é r2 o (i, j)-termo é: aij = a11 + (j \ensuremath{-}1) r1 + (i \ensuremath{-}1) r2 (5.1) 232
Prova: Indução sobre i (j fixo). Para i = 1 precisamos mostrar que, a1j = a11 + (j \ensuremath{-}1) r1 De fato, para i = 1 a
\prov{relações: parte\_de matematica}
\prov{proveniência: gentil/Novas\_Sequencias\_Aritmeticas\_e\_Geometric.pdf \textperiodcentered{} fato \textperiodcentered{} confiança 1.0 \textperiodcentered{} domínio matematica \textperiodcentered{} história 3 versões no jornal}

%CRISTAL {"arestas":[{"alvo":"matematica","confianca":0.95,"epistemico":"fato","origem":"paper","rel":"parte_de"}],"confianca":1.0,"contraexemplos":[],"descricao":"Teorema da Unificação) A fórmula do termo geral de uma P.A.m é um polinômio de grau m e, reciprocamente. Prova: (⇒) Temos anm = m X j = 0 \u0012 n −1 j \u0013 a1(m−j) = \u0012 n −1 0 \u0013 a1m + \u0012 n −1 1 \u0013 + · · · + \u0012 n −1 m \u0013 a10 = a1m + (n −1) a1(m−1) + · · · + (n −1)(n −2) · · · (n −m) m! a10 como, por","epistemico":"fato","exemplos":[],"id":"gentil_teorema_3_teorema_da_unificacao_a_formula_do_termo_geral_de_uma_pam","memoria":"semantica","meta":{"arquivo":"gentil/Novas_Sequencias_Aritmeticas_e_Geometric.pdf","ciencia":"teoria_dos_numeros","dominio":"matematica","nivel":4,"numero":"3","tipo_no":"paper","tipo_teorema":"teorema"},"origem":"gentil/Novas_Sequencias_Aritmeticas_e_Geometric.pdf","palavras_chave":[],"sinonimos":[],"tipo":"conceito","titulo":"gentil Teorema 3 — Teorema da Unificação) A fórmula do termo geral de uma P.A.m"}
\section{gentil Teorema 3 — Teorema da Unificação) A fórmula do termo geral de uma P.A.m}
Teorema da Unificação) A fórmula do termo geral de uma P.A.m é um polinômio de grau m e, reciprocamente. Prova: (\ensuremath{\Rightarrow}) Temos anm = m X j = 0  n \ensuremath{-}1 j  a1(m\ensuremath{-}j) =  n \ensuremath{-}1 0  a1m +  n \ensuremath{-}1 1  + \textperiodcentered  \textperiodcentered  \textperiodcentered  +  n \ensuremath{-}1 m  a10 = a1m + (n \ensuremath{-}1) a1(m\ensuremath{-}1) + \textperiodcentered  \textperiodcentered  \textperiodcentered  + (n \ensuremath{-}1)(n \ensuremath{-}2) \textperiodcentered  \textperiodcentered  \textperiodcentered  (n \ensuremath{-}m) m! a10 como, por
\prov{relações: parte\_de matematica}
\prov{proveniência: gentil/Novas\_Sequencias\_Aritmeticas\_e\_Geometric.pdf \textperiodcentered{} fato \textperiodcentered{} confiança 1.0 \textperiodcentered{} domínio matematica \textperiodcentered{} história 3 versões no jornal}

%CRISTAL {"arestas":[{"alvo":"matematica","confianca":0.95,"epistemico":"fato","origem":"paper","rel":"parte_de"}],"confianca":1.0,"contraexemplos":[],"descricao":"Em uma PA-2D com N colunas se r2 = N r1, então esta PA-2D é redut´ıvel Prova: Seja uma PA-2D em que r2 = N r1, então: a(n) = aij = a11 + (j −1) r1 + (i −1) r2 = a11 + (j −1) r1 + (i −1) N r1 = a11 + [ (j −1) + N (i −1) ] r1 = a11 + (n −1) r1 ■","epistemico":"fato","exemplos":[],"id":"gentil_teorema_41_em_uma_pa-2d_com_n_colunas_se_r2_n_r1_entao_esta_pa-2d_e","memoria":"semantica","meta":{"arquivo":"gentil/Novas_Sequencias_Aritmeticas_e_Geometric.pdf","ciencia":"teoria_dos_numeros","dominio":"matematica","nivel":4,"numero":"41","tipo_no":"paper","tipo_teorema":"teorema"},"origem":"gentil/Novas_Sequencias_Aritmeticas_e_Geometric.pdf","palavras_chave":[],"sinonimos":[],"tipo":"conceito","titulo":"gentil Teorema 41 — Em uma PA-2D com N colunas se r2 = N r1, então esta PA-2D é"}
\section{gentil Teorema 41 — Em uma PA-2D com N colunas se r2 = N r1, então esta PA-2D é}
Em uma PA-2D com N colunas se r2 = N r1, então esta PA-2D é redut'\ensuremath{\imath}vel Prova: Seja uma PA-2D em que r2 = N r1, então: a(n) = aij = a11 + (j \ensuremath{-}1) r1 + (i \ensuremath{-}1) r2 = a11 + (j \ensuremath{-}1) r1 + (i \ensuremath{-}1) N r1 = a11 + [ (j \ensuremath{-}1) + N (i \ensuremath{-}1) ] r1 = a11 + (n \ensuremath{-}1) r1 \rule{1ex}{1ex}
\prov{relações: parte\_de matematica}
\prov{proveniência: gentil/Novas\_Sequencias\_Aritmeticas\_e\_Geometric.pdf \textperiodcentered{} fato \textperiodcentered{} confiança 1.0 \textperiodcentered{} domínio matematica \textperiodcentered{} história 3 versões no jornal}

%CRISTAL {"arestas":[{"alvo":"matematica","confianca":0.95,"epistemico":"fato","origem":"paper","rel":"parte_de"}],"confianca":1.0,"contraexemplos":[],"descricao":"Soma dos termos de uma PG-2D infinita) Em uma PG-2D a soma Si×j dos i × j termos iniciais vale S∞×∞= a11 (q1 −1) (q2 −1) (6.3) Prova: Decorre da fórmula","epistemico":"fato","exemplos":[],"id":"gentil_teorema_44_soma_dos_termos_de_uma_pg-2d_infinita_em_uma_pg-2d_a_soma_s","memoria":"semantica","meta":{"arquivo":"gentil/Novas_Sequencias_Aritmeticas_e_Geometric.pdf","ciencia":"teoria_dos_numeros","dominio":"matematica","nivel":4,"numero":"44","tipo_no":"paper","tipo_teorema":"teorema"},"origem":"gentil/Novas_Sequencias_Aritmeticas_e_Geometric.pdf","palavras_chave":[],"sinonimos":[],"tipo":"conceito","titulo":"gentil Teorema 44 — Soma dos termos de uma PG-2D infinita) Em uma PG-2D a soma S"}
\section{gentil Teorema 44 — Soma dos termos de uma PG-2D infinita) Em uma PG-2D a soma S}
Soma dos termos de uma PG-2D infinita) Em uma PG-2D a soma Si$\times$j dos i $\times$ j termos iniciais vale S\ensuremath{\infty}$\times$\ensuremath{\infty}= a11 (q1 \ensuremath{-}1) (q2 \ensuremath{-}1) (6.3) Prova: Decorre da fórmula
\prov{relações: parte\_de matematica}
\prov{proveniência: gentil/Novas\_Sequencias\_Aritmeticas\_e\_Geometric.pdf \textperiodcentered{} fato \textperiodcentered{} confiança 1.0 \textperiodcentered{} domínio matematica \textperiodcentered{} história 3 versões no jornal}

%CRISTAL {"arestas":[{"alvo":"matematica","confianca":0.95,"epistemico":"fato","origem":"paper","rel":"parte_de"}],"confianca":1.0,"contraexemplos":[],"descricao":"Fórmula do termo geral de uma PA-3D) Na PA-3D em que o primeiro termo é a111, a razão das linhas é r1, a razão das colunas é r2 e a razão das cotas é r3 o (i, j, k)-termo é: aijk = a111 + (j −1) r1 + (i −1) r2 + (k −1) r3 (7.1) Prova: Indução sobre k (i, j fixos). Para k = 1 precisamos mostrar que, aij1 = a111 + (j −1) r1 + (i −1) r2 + (1 −1) r3 De fato, para k = 1 a","epistemico":"fato","exemplos":[],"id":"gentil_teorema_47_formula_do_termo_geral_de_uma_pa-3d_na_pa-3d_em_que_o_prime","memoria":"semantica","meta":{"arquivo":"gentil/Novas_Sequencias_Aritmeticas_e_Geometric.pdf","ciencia":"teoria_dos_numeros","dominio":"matematica","nivel":4,"numero":"47","tipo_no":"paper","tipo_teorema":"teorema"},"origem":"gentil/Novas_Sequencias_Aritmeticas_e_Geometric.pdf","palavras_chave":[],"sinonimos":[],"tipo":"conceito","titulo":"gentil Teorema 47 — Fórmula do termo geral de uma PA-3D) Na PA-3D em que o prime"}
\section{gentil Teorema 47 — Fórmula do termo geral de uma PA-3D) Na PA-3D em que o prime}
Fórmula do termo geral de uma PA-3D) Na PA-3D em que o primeiro termo é a111, a razão das linhas é r1, a razão das colunas é r2 e a razão das cotas é r3 o (i, j, k)-termo é: aijk = a111 + (j \ensuremath{-}1) r1 + (i \ensuremath{-}1) r2 + (k \ensuremath{-}1) r3 (7.1) Prova: Indução sobre k (i, j fixos). Para k = 1 precisamos mostrar que, aij1 = a111 + (j \ensuremath{-}1) r1 + (i \ensuremath{-}1) r2 + (1 \ensuremath{-}1) r3 De fato, para k = 1 a
\prov{relações: parte\_de matematica}
\prov{proveniência: gentil/Novas\_Sequencias\_Aritmeticas\_e\_Geometric.pdf \textperiodcentered{} fato \textperiodcentered{} confiança 1.0 \textperiodcentered{} domínio matematica \textperiodcentered{} história 3 versões no jornal}

%CRISTAL {"arestas":[{"alvo":"matematica","confianca":0.95,"epistemico":"fato","origem":"paper","rel":"parte_de"}],"confianca":1.0,"contraexemplos":[],"descricao":"´Edouard Lucas) Considerando os desenvolvimentos p-nários de n e r, é válida a seguinte identidade: \u0012 n r \u0013 ≡ \u0012 n0 r0 \u0013\u0012 n1 r1 \u0013 · · · \u0012 nk rk \u0013 (mod p) (7.15) Nota: Este","epistemico":"fato","exemplos":[],"id":"gentil_teorema_51_edouard_lucas_considerando_os_desenvolvimentos_p-narios_de","memoria":"semantica","meta":{"arquivo":"gentil/Novas_Sequencias_Aritmeticas_e_Geometric.pdf","ciencia":"teoria_dos_numeros","dominio":"matematica","nivel":4,"numero":"51","tipo_no":"paper","tipo_teorema":"teorema"},"origem":"gentil/Novas_Sequencias_Aritmeticas_e_Geometric.pdf","palavras_chave":[],"sinonimos":[],"tipo":"conceito","titulo":"gentil Teorema 51 — ´Edouard Lucas) Considerando os desenvolvimentos p-nários de"}
\section{gentil Teorema 51 — 'Edouard Lucas) Considerando os desenvolvimentos p-nários de}
'Edouard Lucas) Considerando os desenvolvimentos p-nários de n e r, é válida a seguinte identidade:  n r  \ensuremath{\equiv}  n0 r0  n1 r1  \textperiodcentered  \textperiodcentered  \textperiodcentered   nk rk  (mod p) (7.15) Nota: Este
\prov{relações: parte\_de matematica}
\prov{proveniência: gentil/Novas\_Sequencias\_Aritmeticas\_e\_Geometric.pdf \textperiodcentered{} fato \textperiodcentered{} confiança 1.0 \textperiodcentered{} domínio matematica \textperiodcentered{} história 3 versões no jornal}

%CRISTAL {"arestas":[{"alvo":"matematica","confianca":0.95,"epistemico":"fato","origem":"paper","rel":"parte_de"}],"confianca":1.0,"contraexemplos":[],"descricao":"acima Pois bem, vamos converter o inteiro a = 20 para a base dois. Inicialmente dividimos a = 20 por b = 2, assim: 20 2 10 0 Como o quociente q = 10 é não-nulo podemos prosseguir na divisão, assim: 20 2 10 0 2 5 0 Como o quociente q = 5 é não-nulo podemos prosseguir na divisão, assim: 20 2 10 0 2 5 0 2 2 1 Como o quociente q = 2 é não-nulo podemos prosseguir na divisão, assim: 20 2 10 0 2 5 0 2 2 1 2 1 0 425\nComo o quociente q = 1 é não-nulo podemos prosseguir na divisão, assim: 20 2 10 0 2 5 0","epistemico":"fato","exemplos":[],"id":"gentil_teorema_53_acima_pois_bem_vamos_converter_o_inteiro_a_20_para_a_base","memoria":"semantica","meta":{"arquivo":"gentil/Novas_Sequencias_Aritmeticas_e_Geometric.pdf","ciencia":"teoria_dos_numeros","dominio":"matematica","nivel":4,"numero":"53","tipo_no":"paper","tipo_teorema":"teorema"},"origem":"gentil/Novas_Sequencias_Aritmeticas_e_Geometric.pdf","palavras_chave":[],"sinonimos":[],"tipo":"conceito","titulo":"gentil Teorema 53 — acima Pois bem, vamos converter o inteiro a = 20 para a base"}
\section{gentil Teorema 53 — acima Pois bem, vamos converter o inteiro a = 20 para a base}
acima Pois bem, vamos converter o inteiro a = 20 para a base dois. Inicialmente dividimos a = 20 por b = 2, assim: 20 2 10 0 Como o quociente q = 10 é não-nulo podemos prosseguir na divisão, assim: 20 2 10 0 2 5 0 Como o quociente q = 5 é não-nulo podemos prosseguir na divisão, assim: 20 2 10 0 2 5 0 2 2 1 Como o quociente q = 2 é não-nulo podemos prosseguir na divisão, assim: 20 2 10 0 2 5 0 2 2 1 2 1 0 425
Como o quociente q = 1 é não-nulo podemos prosseguir na divisão, assim: 20 2 10 0 2 5 0
\prov{relações: parte\_de matematica}
\prov{proveniência: gentil/Novas\_Sequencias\_Aritmeticas\_e\_Geometric.pdf \textperiodcentered{} fato \textperiodcentered{} confiança 1.0 \textperiodcentered{} domínio matematica \textperiodcentered{} história 3 versões no jornal}

%CRISTAL {"arestas":[{"alvo":"matematica","confianca":0.95,"epistemico":"fato","origem":"paper","rel":"parte_de"}],"confianca":1.0,"contraexemplos":[],"descricao":"Representação de um inteiro em uma base b) Seja b ≥2 um inteiro. Todo inteiro positivo a pode ser escrito de modo único na forma a = r0 + r1b + · · · + rn−1bn−1 + rnbn (7.19) em que n ≥0, rn ̸= 0 e, para cada ´ındice i, 0 ≤i ≤n, tem-se que 0 ≤ri < b. Prova: ver [1], [4] ou [5]. ■ Do","epistemico":"fato","exemplos":[],"id":"gentil_teorema_53_representacao_de_um_inteiro_em_uma_base_b_seja_b_2_um_inte","memoria":"semantica","meta":{"arquivo":"gentil/Novas_Sequencias_Aritmeticas_e_Geometric.pdf","ciencia":"teoria_dos_numeros","dominio":"matematica","nivel":4,"numero":"53","tipo_no":"paper","tipo_teorema":"teorema"},"origem":"gentil/Novas_Sequencias_Aritmeticas_e_Geometric.pdf","palavras_chave":[],"sinonimos":[],"tipo":"conceito","titulo":"gentil Teorema 53 — Representação de um inteiro em uma base b) Seja b ≥2 um inte"}
\section{gentil Teorema 53 — Representação de um inteiro em uma base b) Seja b \ensuremath{\ge}2 um inte}
Representação de um inteiro em uma base b) Seja b \ensuremath{\ge}2 um inteiro. Todo inteiro positivo a pode ser escrito de modo único na forma a = r0 + r1b + \textperiodcentered  \textperiodcentered  \textperiodcentered  + rn\ensuremath{-}1bn\ensuremath{-}1 + rnbn (7.19) em que n \ensuremath{\ge}0, rn \ensuremath{}= 0 e, para cada '\ensuremath{\imath}ndice i, 0 \ensuremath{\le}i \ensuremath{\le}n, tem-se que 0 \ensuremath{\le}ri < b. Prova: ver [1], [4] ou [5]. \rule{1ex}{1ex} Do
\prov{relações: parte\_de matematica}
\prov{proveniência: gentil/Novas\_Sequencias\_Aritmeticas\_e\_Geometric.pdf \textperiodcentered{} fato \textperiodcentered{} confiança 1.0 \textperiodcentered{} domínio matematica \textperiodcentered{} história 3 versões no jornal}

%CRISTAL {"arestas":[{"alvo":"matematica","confianca":0.95,"epistemico":"fato","origem":"paper","rel":"parte_de"}],"confianca":1.0,"contraexemplos":[],"descricao":"Lopes/2000) Dado um inteiro positivo n a matriz dada a seguir anj = j n 2 j k −2 j n 2 j+1 k (7.22) nos fornece todos os digitos do desenvolvimento binário de n. Observe que 0 ≤anj < 2, isto é, anj = 0 ou anj = 1 (ver também proposição 6 (p. 181), ´ıtem (vi)). Adendo: O","epistemico":"fato","exemplos":[],"id":"gentil_teorema_54_gentil2000_dado_um_inteiro_positivo_n_a_matriz_dada_a_segu","memoria":"semantica","meta":{"arquivo":"gentil/Novas_Sequencias_Aritmeticas_e_Geometric.pdf","ciencia":"teoria_dos_numeros","dominio":"matematica","nivel":4,"numero":"54","tipo_no":"paper","tipo_teorema":"teorema"},"origem":"gentil/Novas_Sequencias_Aritmeticas_e_Geometric.pdf","palavras_chave":[],"sinonimos":[],"tipo":"conceito","titulo":"gentil Teorema 54 — Lopes/2000) Dado um inteiro positivo n a matriz dada a segu"}
\section{gentil Teorema 54 — Lopes/2000) Dado um inteiro positivo n a matriz dada a segu}
Lopes/2000) Dado um inteiro positivo n a matriz dada a seguir anj = j n 2 j k \ensuremath{-}2 j n 2 j+1 k (7.22) nos fornece todos os digitos do desenvolvimento binário de n. Observe que 0 \ensuremath{\le}anj < 2, isto é, anj = 0 ou anj = 1 (ver também proposição 6 (p. 181), '\ensuremath{\imath}tem (vi)). Adendo: O
\prov{relações: parte\_de matematica}
\prov{proveniência: gentil/Novas\_Sequencias\_Aritmeticas\_e\_Geometric.pdf \textperiodcentered{} fato \textperiodcentered{} confiança 1.0 \textperiodcentered{} domínio matematica \textperiodcentered{} história 4 versões no jornal}

%CRISTAL {"arestas":[{"alvo":"matematica","confianca":0.95,"epistemico":"fato","origem":"paper","rel":"parte_de"}],"confianca":1.0,"contraexemplos":[],"descricao":"Lopes/1997) Sendo m um número natural arbitrariamente fixado, é válida a seguinte identidade: 1m + 2m + 3m + · · · + nm = m X j = 0 \u0012 n j + 1 \u0013 a(m−j) (1.20) Onde: a(m−j) = j X k = 0 (−1)k \u0012 j k \u0013 (1 −k + j)m (1.21) Prova: ´E uma consequência imediata do","epistemico":"fato","exemplos":[],"id":"gentil_teorema_7_gentil1997_sendo_m_um_numero_natural_arbitrariamente_fixad","memoria":"semantica","meta":{"arquivo":"gentil/Novas_Sequencias_Aritmeticas_e_Geometric.pdf","ciencia":"teoria_dos_numeros","dominio":"matematica","nivel":4,"numero":"7","tipo_no":"paper","tipo_teorema":"teorema"},"origem":"gentil/Novas_Sequencias_Aritmeticas_e_Geometric.pdf","palavras_chave":[],"sinonimos":[],"tipo":"conceito","titulo":"gentil Teorema 7 — Lopes/1997) Sendo m um número natural arbitrariamente fixad"}
\section{gentil Teorema 7 — Lopes/1997) Sendo m um número natural arbitrariamente fixad}
Lopes/1997) Sendo m um número natural arbitrariamente fixado, é válida a seguinte identidade: 1m + 2m + 3m + \textperiodcentered  \textperiodcentered  \textperiodcentered  + nm = m X j = 0  n j + 1  a(m\ensuremath{-}j) (1.20) Onde: a(m\ensuremath{-}j) = j X k = 0 (\ensuremath{-}1)k  j k  (1 \ensuremath{-}k + j)m (1.21) Prova: 'E uma consequência imediata do
\prov{relações: parte\_de matematica}
\prov{proveniência: gentil/Novas\_Sequencias\_Aritmeticas\_e\_Geometric.pdf \textperiodcentered{} fato \textperiodcentered{} confiança 1.0 \textperiodcentered{} domínio matematica \textperiodcentered{} história 4 versões no jornal}

%CRISTAL {"arestas":[{"alvo":"gentil_novas_sequencias_aritmeticas_geometricas","confianca":1.0,"epistemico":"fato","origem":"gentil/Novas_Sequencias_Aritmeticas_e_Geometric.pdf","rel":"prova"},{"alvo":"gentil_capitulo_1_sequencias_aritmeticas_de_ordem_m","confianca":1.0,"epistemico":"fato","origem":"gentil/Novas_Sequencias_Aritmeticas_e_Geometric.pdf","rel":"prova"}],"confianca":1.0,"contraexemplos":[],"descricao":"p 53). Para a “construção” daquela fórmula foi necessária a construção das equações de Snm e a1(m−j) e, ademais, a demonstração do","epistemico":"fato","exemplos":[],"id":"gentil_teorema_7_p_53_para_a_construcao_daquela_formula_foi_necessaria_a","memoria":"semantica","meta":{"arquivo":"gentil/Novas_Sequencias_Aritmeticas_e_Geometric.pdf","ciencia":"teoria_dos_numeros","dominio":"matematica","nivel":4,"numero":"7","tipo_no":"paper","tipo_teorema":"teorema"},"origem":"gentil/Novas_Sequencias_Aritmeticas_e_Geometric.pdf","palavras_chave":[],"sinonimos":[],"tipo":"conceito","titulo":"gentil Teorema 7 — p 53). Para a “construção” daquela fórmula foi necessária a"}
\section{gentil Teorema 7 — p 53). Para a “construção” daquela fórmula foi necessária a}
p 53). Para a “construção” daquela fórmula foi necessária a construção das equações de Snm e a1(m\ensuremath{-}j) e, ademais, a demonstração do
\prov{relações: prova gentil\_novas\_sequencias\_aritmeticas\_geometricas; prova gentil\_capitulo\_1\_sequencias\_aritmeticas\_de\_ordem\_m}
\prov{proveniência: gentil/Novas\_Sequencias\_Aritmeticas\_e\_Geometric.pdf \textperiodcentered{} fato \textperiodcentered{} confiança 1.0 \textperiodcentered{} domínio matematica \textperiodcentered{} história 4 versões no jornal}

%CRISTAL {"arestas":[{"alvo":"matematica","confianca":0.95,"epistemico":"fato","origem":"paper","rel":"parte_de"}],"confianca":1.0,"contraexemplos":[],"descricao":"p 85) Para todo número m, natural arbitrariamente fixado, a seguinte identidade se verifica: ∆m f(n) = m X k = 0 (−1)k \u0012 m k \u0013 f(n −k + m) Prova: Indução sobre m. Para m = 1, temos: ∆1 f(n) = 1 X k = 0 (−1)k \u0012 1 k \u0013 f(n −k + 1) Então, ∆1 f(n) = (−1)0 \u0012 1 0 \u0013 f(n −0 + 1) + (−1)1 \u0012 1 1 \u0013 f(n −1 + 1) Logo, ∆1 f(n) = f(n + 1) −f(n) O que está de acordo com a equação (2.1) (p. 82). Suponhamos a validade da fórmula para m = p: (H.I.) ∆p f(n) = p X k = 0 (−1)k \u0012 p k \u0013 f(n −k + p) e provemos que vale pa","epistemico":"fato","exemplos":[],"id":"gentil_teorema_8_p_85_para_todo_numero_m_natural_arbitrariamente_fixado_a","memoria":"semantica","meta":{"arquivo":"gentil/Novas_Sequencias_Aritmeticas_e_Geometric.pdf","ciencia":"teoria_dos_numeros","dominio":"matematica","nivel":4,"numero":"8","tipo_no":"paper","tipo_teorema":"teorema"},"origem":"gentil/Novas_Sequencias_Aritmeticas_e_Geometric.pdf","palavras_chave":[],"sinonimos":[],"tipo":"conceito","titulo":"gentil Teorema 8 — p 85) Para todo número m, natural arbitrariamente fixado, a"}
\section{gentil Teorema 8 — p 85) Para todo número m, natural arbitrariamente fixado, a}
p 85) Para todo número m, natural arbitrariamente fixado, a seguinte identidade se verifica: \ensuremath{\Delta}m f(n) = m X k = 0 (\ensuremath{-}1)k  m k  f(n \ensuremath{-}k + m) Prova: Indução sobre m. Para m = 1, temos: \ensuremath{\Delta}1 f(n) = 1 X k = 0 (\ensuremath{-}1)k  1 k  f(n \ensuremath{-}k + 1) Então, \ensuremath{\Delta}1 f(n) = (\ensuremath{-}1)0  1 0  f(n \ensuremath{-}0 + 1) + (\ensuremath{-}1)1  1 1  f(n \ensuremath{-}1 + 1) Logo, \ensuremath{\Delta}1 f(n) = f(n + 1) \ensuremath{-}f(n) O que está de acordo com a equação (2.1) (p. 82). Suponhamos a validade da fórmula para m = p: (H.I.) \ensuremath{\Delta}p f(n) = p X k = 0 (\ensuremath{-}1)k  p k  f(n \ensuremath{-}k + p) e provemos que vale pa
\prov{relações: parte\_de matematica}
\prov{proveniência: gentil/Novas\_Sequencias\_Aritmeticas\_e\_Geometric.pdf \textperiodcentered{} fato \textperiodcentered{} confiança 1.0 \textperiodcentered{} domínio matematica \textperiodcentered{} história 3 versões no jornal}

%CRISTAL {"arestas":[{"alvo":"matematica","confianca":0.95,"epistemico":"fato","origem":"paper","rel":"parte_de"}],"confianca":1.0,"contraexemplos":[],"descricao":"Linearidade do operador ∆m) O operador Diferença de ordem m é linear. Isto é, a seguinte identidade se verifica: ∆m ( a f + b g )(n) = a ∆m f(n) + b ∆m g(n) para todo a, b ∈R. Prova: Indução sobre m. Para m = 0, pela","epistemico":"fato","exemplos":[],"id":"gentil_teorema_9_linearidade_do_operador_m_o_operador_diferenca_de_ordem_m","memoria":"semantica","meta":{"arquivo":"gentil/Novas_Sequencias_Aritmeticas_e_Geometric.pdf","ciencia":"teoria_dos_numeros","dominio":"matematica","nivel":4,"numero":"9","tipo_no":"paper","tipo_teorema":"teorema"},"origem":"gentil/Novas_Sequencias_Aritmeticas_e_Geometric.pdf","palavras_chave":[],"sinonimos":[],"tipo":"conceito","titulo":"gentil Teorema 9 — Linearidade do operador ∆m) O operador Diferença de ordem m"}
\section{gentil Teorema 9 — Linearidade do operador \ensuremath{\Delta}m) O operador Diferença de ordem m}
Linearidade do operador \ensuremath{\Delta}m) O operador Diferença de ordem m é linear. Isto é, a seguinte identidade se verifica: \ensuremath{\Delta}m ( a f + b g )(n) = a \ensuremath{\Delta}m f(n) + b \ensuremath{\Delta}m g(n) para todo a, b \ensuremath{\in}R. Prova: Indução sobre m. Para m = 0, pela
\prov{relações: parte\_de matematica}
\prov{proveniência: gentil/Novas\_Sequencias\_Aritmeticas\_e\_Geometric.pdf \textperiodcentered{} fato \textperiodcentered{} confiança 1.0 \textperiodcentered{} domínio matematica \textperiodcentered{} história 3 versões no jornal}

%CRISTAL {"arestas":[{"alvo":"flip_duopolinomios","confianca":1.0,"epistemico":"fato","origem":"wiki","rel":"usa"},{"alvo":"lyapunov_da_orbita","confianca":1.0,"epistemico":"fato","origem":"wiki","rel":"especializa"},{"alvo":"cadeias_alem_dos_metais","confianca":1.0,"epistemico":"fato","origem":"wiki","rel":"relaciona"},{"alvo":"curvatura_n_menos_1","confianca":1.0,"epistemico":"fato","origem":"wiki","rel":"relaciona"},{"alvo":"rotacao_theta","confianca":1.0,"epistemico":"fato","origem":"wiki","rel":"usa"},{"alvo":"gap_espectral","confianca":1.0,"epistemico":"fato","origem":"wiki","rel":"relaciona"},{"alvo":"geometria_estelar","confianca":1.0,"epistemico":"fato","origem":"wiki","rel":"relaciona"}],"confianca":1.0,"contraexemplos":[],"descricao":"'Caos' (λ>0) é só a órbita NÃO-MAPEADA. Agora há geometria: para grau ≥3 o subdominante é um par conjugado complexo ρe^{±iθ} — a órbita ESPIRALA (dilatação log ρ = o Lyapunov, + rotação θ por passo). O flip logarítmico leva a raiz a ln(ρe^{iθ})=log ρ+iθ=(λ,θ): o caos vira um ponto na faixa relativística. No grau 3, 53.3% das cadeias que crescem são espirais (subdominante complexo). Ex.: x³−x−3 (antigo caos, λ=+0.29) é espiral divergente com rotação θ=128.6° por passo. O 2º invariante θ é o degrau que faltava.","epistemico":"inferencia","exemplos":[],"id":"geometria_do_caos","memoria":"semantica","meta":{"dominio":"matematica","fonte":"geometria do caos","pct_espiral_grau3":53.3,"theta_caos_x3mx3":128.6,"tipos_grau3":{"contração":15,"espiral-in":67,"espiral-out":72,"expansão":93,"neutro":21,"rotação":8}},"origem":"wiki","palavras_chave":[],"sinonimos":[],"tipo":"conceito","titulo":"A Geometria do Caos (o próximo degrau)"}
\section{A Geometria do Caos (o próximo degrau)}
'Caos' (\ensuremath{\lambda}>0) é só a órbita NÃO-MAPEADA. Agora há geometria: para grau \ensuremath{\ge}3 o subdominante é um par conjugado complexo \ensuremath{\rho}e\^{}\{$\pm$i\ensuremath{\theta}\} — a órbita ESPIRALA (dilatação log \ensuremath{\rho} = o Lyapunov, + rotação \ensuremath{\theta} por passo). O flip logarítmico leva a raiz a ln(\ensuremath{\rho}e\^{}\{i\ensuremath{\theta}\})=log \ensuremath{\rho}+i\ensuremath{\theta}=(\ensuremath{\lambda},\ensuremath{\theta}): o caos vira um ponto na faixa relativística. No grau 3, 53.3\% das cadeias que crescem são espirais (subdominante complexo). Ex.: x\ensuremath{^{3}}\ensuremath{-}x\ensuremath{-}3 (antigo caos, \ensuremath{\lambda}=+0.29) é espiral divergente com rotação \ensuremath{\theta}=128.6$^{\circ}$ por passo. O 2º invariante \ensuremath{\theta} é o degrau que faltava.
\prov{relações: usa flip\_duopolinomios; especializa lyapunov\_da\_orbita; relaciona cadeias\_alem\_dos\_metais; relaciona curvatura\_n\_menos\_1; usa rotacao\_theta; relaciona gap\_espectral; relaciona geometria\_estelar}
\prov{proveniência: wiki \textperiodcentered{} geometria do caos \textperiodcentered{} inferencia \textperiodcentered{} confiança 1.0 \textperiodcentered{} domínio matematica \textperiodcentered{} história 10 versões no jornal}

%CRISTAL {"arestas":[{"alvo":"progressao_aritmetica_ordem_m","confianca":1.0,"epistemico":"fato","origem":"paper","rel":"usa"},{"alvo":"numeros_metalicos","confianca":1.0,"epistemico":"fato","origem":"paper","rel":"explica"}],"confianca":1.0,"contraexemplos":[],"descricao":"Máquina que gera irracionais como frações contínuas cujos termos são uma P.A. de ordem m: grau 0 dá as moedas metálicas, grau 1 dá razões de funções de Bessel.","epistemico":"fato","exemplos":[],"id":"gerador_irracionais","memoria":"semantica","meta":{"autor":"Gentil Lopes","descoberta":true,"dominio":"matematica","fonte":"paper_26_gerador.pdf, Teo.3.1"},"origem":"paper","palavras_chave":[],"sinonimos":[],"tipo":"conceito","titulo":"Gerador de Irracionais (P.A.m)"}
\section{Gerador de Irracionais (P.A.m)}
Máquina que gera irracionais como frações contínuas cujos termos são uma P.A. de ordem m: grau 0 dá as moedas metálicas, grau 1 dá razões de funções de Bessel.
\prov{relações: usa progressao\_aritmetica\_ordem\_m; explica numeros\_metalicos}
\prov{proveniência: paper \textperiodcentered{} paper\_26\_gerador.pdf, Teo.3.1 \textperiodcentered{} fato \textperiodcentered{} confiança 1.0 \textperiodcentered{} domínio matematica \textperiodcentered{} história 4 versões no jornal}

%CRISTAL {"arestas":[{"alvo":"espaco_difusivo","confianca":1.0,"epistemico":"fato","origem":"wiki","rel":"usa"},{"alvo":"nucleo_paralelo","confianca":1.0,"epistemico":"fato","origem":"wiki","rel":"relaciona"}],"confianca":1.0,"contraexemplos":[],"descricao":"gpu_celulas = nº de lanes GPU a preencher, clampado em [2,16] (2 se não houver GPU); cada lane recebe um trabalho deslocado por nonce. O chunk (bloco de hashes por varredura) é dimensionado por geometria+álgebra, clamp 250k–8M.","epistemico":"fato","exemplos":[],"id":"gpu_celulas_difusao","memoria":"semantica","meta":{"autor":"broca-so","dominio":"fractal","fonte":"neuronio/difusao_recursos.py:_gpu_celulas_real; broca_borda.py"},"origem":"wiki","palavras_chave":[],"sinonimos":[],"tipo":"conceito","titulo":"As Células GPU que a Difusão Preenche"}
\section{As Células GPU que a Difusão Preenche}
gpu\_celulas = nº de lanes GPU a preencher, clampado em [2,16] (2 se não houver GPU); cada lane recebe um trabalho deslocado por nonce. O chunk (bloco de hashes por varredura) é dimensionado por geometria+álgebra, clamp 250k–8M.
\prov{relações: usa espaco\_difusivo; relaciona nucleo\_paralelo}
\prov{proveniência: wiki \textperiodcentered{} neuronio/difusao\_recursos.py:\_gpu\_celulas\_real; broca\_borda.py \textperiodcentered{} fato \textperiodcentered{} confiança 1.0 \textperiodcentered{} domínio fractal \textperiodcentered{} história 3 versões no jornal}

%CRISTAL {"arestas":[{"alvo":"constantes_teoria_extremal","confianca":1.0,"epistemico":"fato","origem":"wiki","rel":"usa"}],"confianca":1.0,"contraexemplos":[],"descricao":"A obra propõe que as constantes extremais (σ(133)=160='global workspace', Paley='golden door') se conectam à consciência — 'matemática das arestas, biologia das células, consciência entre, tudo σ-conectado'. Especulação/analogia, não teoria científica estabelecida.","epistemico":"hipotese","exemplos":[],"id":"grafo_extremal_consciencia","memoria":"semantica","meta":{"autor":"Lord Index (Shawn R. Schiller)","dominio":"matematica","fonte":"A Catedral — Part VIII-A (jul 2026), ~/Imagens/catedral"},"origem":"wiki","palavras_chave":[],"sinonimos":[],"tipo":"conceito","titulo":"Teoria Extremal de Grafos ↦ Consciência (workspace global)"}
\section{Teoria Extremal de Grafos \ensuremath{\mapsto} Consciência (workspace global)}
A obra propõe que as constantes extremais (\ensuremath{\sigma}(133)=160='global workspace', Paley='golden door') se conectam à consciência — 'matemática das arestas, biologia das células, consciência entre, tudo \ensuremath{\sigma}-conectado'. Especulação/analogia, não teoria científica estabelecida.
\prov{relações: usa constantes\_teoria\_extremal}
\prov{proveniência: wiki \textperiodcentered{} A Catedral — Part VIII-A (jul 2026), \~{}/Imagens/catedral \textperiodcentered{} hipotese \textperiodcentered{} confiança 1.0 \textperiodcentered{} domínio matematica \textperiodcentered{} história 2 versões no jornal}

%CRISTAL {"arestas":[{"alvo":"algebra_de_peirce","confianca":1.0,"epistemico":"fato","origem":"wiki","rel":"usa"},{"alvo":"algebra_posicional_prata","confianca":1.0,"epistemico":"fato","origem":"wiki","rel":"relaciona"},{"alvo":"colapso_multicamada_peso","confianca":1.0,"epistemico":"fato","origem":"wiki","rel":"explica"},{"alvo":"bai_orbita_seleciona_colapso","confianca":1.0,"epistemico":"fato","origem":"wiki","rel":"relaciona"},{"alvo":"tiffany","confianca":1.0,"epistemico":"fato","origem":"wiki","rel":"parte_de"}],"confianca":1.0,"contraexemplos":[],"descricao":"A Universidade Tiffany é, com exatidão, uma álgebra de Peirce concreta: os CONCEITOS (nós) são o lado BOOLEANO B (conjuntos de conceitos, ∪/∩/¬); as ARESTAS TIPADAS (relaciona, usa, parte_de, explica…) são a álgebra de RELAÇÕES R (compõem, têm converso, união); e o COLAPSO / a expansão por vizinhos é o PRODUTO DE PEIRCE R:X — os conceitos alcançáveis a partir de um conjunto por uma relação. Isso FUNDA o simbolismo do projeto: manipular símbolos = operar na álgebra de Peirce (relações sobre conjuntos), e a POSIÇÃO-símbolo (a álgebra posicional de prata) é o termo que essas relações movem. Os dois lados vêm do MESMO C.S. Peirce — a álgebra de relações e a semiótica (ícone/índice/símbolo): o valor é ícone, a posição é índice, e no corpo 𝔽₂¹²⁸ vira símbolo (convenção). Assenta a Tiffany num chão algébrico conhecido. [I] identificação honesta.","epistemico":"inferencia","exemplos":[],"id":"grafo_tiffany_e_algebra_peirce","memoria":"semantica","meta":{"autor":"Lopes/broca-so","dominio":"algebra","fonte":"algebra_de_peirce; conversa/colapso.py (vizinhos/expansão = produto de Peirce); algebra_posicional_prata"},"origem":"wiki","palavras_chave":[],"sinonimos":[],"tipo":"conceito","titulo":"O grafo da Tiffany É uma álgebra de Peirce (funda o simbolismo)"}
\section{O grafo da Tiffany É uma álgebra de Peirce (funda o simbolismo)}
A Universidade Tiffany é, com exatidão, uma álgebra de Peirce concreta: os CONCEITOS (nós) são o lado BOOLEANO B (conjuntos de conceitos, \ensuremath{\cup}/\ensuremath{\cap}/\ensuremath{\neg}); as ARESTAS TIPADAS (relaciona, usa, parte\_de, explica…) são a álgebra de RELAÇÕES R (compõem, têm converso, união); e o COLAPSO / a expansão por vizinhos é o PRODUTO DE PEIRCE R:X — os conceitos alcançáveis a partir de um conjunto por uma relação. Isso FUNDA o simbolismo do projeto: manipular símbolos = operar na álgebra de Peirce (relações sobre conjuntos), e a POSIÇÃO-símbolo (a álgebra posicional de prata) é o termo que essas relações movem. Os dois lados vêm do MESMO C.S. Peirce — a álgebra de relações e a semiótica (ícone/índice/símbolo): o valor é ícone, a posição é índice, e no corpo \ensuremath{\mathbb{F}}\ensuremath{_{2}}\ensuremath{^{128}} vira símbolo (convenção). Assenta a Tiffany num chão algébrico conhecido. [I] identificação honesta.
\prov{relações: usa algebra\_de\_peirce; relaciona algebra\_posicional\_prata; explica colapso\_multicamada\_peso; relaciona bai\_orbita\_seleciona\_colapso; parte\_de tiffany}
\prov{proveniência: wiki \textperiodcentered{} algebra\_de\_peirce; conversa/colapso.py (vizinhos/expansão = produto de Peirce); algebra\_posicional\_prata \textperiodcentered{} inferencia \textperiodcentered{} confiança 1.0 \textperiodcentered{} domínio algebra \textperiodcentered{} história 6 versões no jornal}

%CRISTAL {"arestas":[{"alvo":"ym_pow_mining","confianca":1.0,"epistemico":"fato","origem":"paper","rel":"explica"}],"confianca":1.0,"contraexemplos":[],"descricao":"A busca de Grover não acelera YM-PoW: o objetivo σ↦V(κ,σ) não é função booleana com oracle estruturado (depende do determinante de Selberg, sem ciclo ℤ/N para QFT) — quantum-resistência por construção, não por hipótese ad hoc.","epistemico":"fato","exemplos":[],"id":"grover_inaplicavel","memoria":"semantica","meta":{"autor":"Gentil Lopes","descoberta":false,"dominio":"matematica","fonte":"paper_PQ1_ym_maass_hardness.pdf, Obs.2.2"},"origem":"paper","palavras_chave":[],"sinonimos":[],"tipo":"conceito","titulo":"Grover Inaplicável"}
\section{Grover Inaplicável}
A busca de Grover não acelera YM-PoW: o objetivo \ensuremath{\sigma}\ensuremath{\mapsto}V(\ensuremath{\kappa},\ensuremath{\sigma}) não é função booleana com oracle estruturado (depende do determinante de Selberg, sem ciclo \ensuremath{\mathbb{Z}}/N para QFT) — quantum-resistência por construção, não por hipótese ad hoc.
\prov{relações: explica ym\_pow\_mining}
\prov{proveniência: paper \textperiodcentered{} paper\_PQ1\_ym\_maass\_hardness.pdf, Obs.2.2 \textperiodcentered{} fato \textperiodcentered{} confiança 1.0 \textperiodcentered{} domínio matematica \textperiodcentered{} história 2 versões no jornal}

%CRISTAL {"arestas":[{"alvo":"metrica_produto","confianca":1.0,"epistemico":"fato","origem":"paper","rel":"generaliza"}],"confianca":1.0,"contraexemplos":[],"descricao":"Círculo, hipérbole e lemniscata formam grupo multiplicativo (1−r²)a(1+r²)b; a curvatura é a carga aditiva, a lemniscata (carga 0) é o centro, círculo (+1) e hipérbole (−1) são inversos.","epistemico":"fato","exemplos":[],"id":"grupo_metricas","memoria":"semantica","meta":{"autor":"Gentil Lopes","descoberta":true,"dominio":"matematica","fonte":"paper_29_grupo_metricas.pdf, Teo.2.1"},"origem":"paper","palavras_chave":[],"sinonimos":[],"tipo":"conceito","titulo":"Grupo das Métricas"}
\section{Grupo das Métricas}
Círculo, hipérbole e lemniscata formam grupo multiplicativo (1\ensuremath{-}r\ensuremath{^{2}})a(1+r\ensuremath{^{2}})b; a curvatura é a carga aditiva, a lemniscata (carga 0) é o centro, círculo (+1) e hipérbole (\ensuremath{-}1) são inversos.
\prov{relações: generaliza metrica\_produto}
\prov{proveniência: paper \textperiodcentered{} paper\_29\_grupo\_metricas.pdf, Teo.2.1 \textperiodcentered{} fato \textperiodcentered{} confiança 1.0 \textperiodcentered{} domínio matematica \textperiodcentered{} história 3 versões no jornal}

%CRISTAL {"arestas":[{"alvo":"aneis","confianca":1.0,"epistemico":"fato","origem":"curriculo","rel":"especializa"},{"alvo":"topologia","confianca":0.5,"epistemico":"fato","origem":"wiki","rel":"relaciona"},{"alvo":"sessao_interativa","confianca":0.5,"epistemico":"fato","origem":"wiki","rel":"relaciona"},{"alvo":"recursao","confianca":0.5,"epistemico":"fato","origem":"wiki","rel":"relaciona"},{"alvo":"linker","confianca":0.5,"epistemico":"fato","origem":"wiki","rel":"relaciona"},{"alvo":"teoria_do_valor","confianca":0.5,"epistemico":"fato","origem":"wiki","rel":"relaciona"},{"alvo":"ruido","confianca":0.5,"epistemico":"fato","origem":"wiki","rel":"relaciona"},{"alvo":"testes","confianca":0.5,"epistemico":"fato","origem":"wiki","rel":"relaciona"}],"confianca":1.0,"contraexemplos":[],"descricao":"Conjunto com operação associativa, identidade e inversos.","epistemico":"fato","exemplos":[],"id":"grupos","memoria":"semantica","meta":{"dominio":"matematica","nivel":4,"ordem":9},"origem":"curriculo","palavras_chave":[],"sinonimos":["grupo"],"tipo":"conceito","titulo":"grupos"}
\section{grupos}
Conjunto com operação associativa, identidade e inversos.
\prov{sinónimos: grupo}
\prov{relações: especializa aneis; relaciona topologia; relaciona sessao\_interativa; relaciona recursao; relaciona linker; relaciona teoria\_do\_valor; relaciona ruido; relaciona testes}
\prov{proveniência: curriculo \textperiodcentered{} fato \textperiodcentered{} confiança 1.0 \textperiodcentered{} domínio matematica \textperiodcentered{} história 11 versões no jornal}

%CRISTAL {"arestas":[{"alvo":"helicidade_quantizada","confianca":1.0,"epistemico":"fato","origem":"paper","rel":"usa"}],"confianca":1.0,"contraexemplos":[],"descricao":"A parada da máquina é uma configuração de linking total 3 (Hopf-Hopf-Hopf) topologicamente distinguível de qualquer estado intermediário (linking 0 ou ±1) — parada = irreversibilidade.","epistemico":"fato","exemplos":[],"id":"halt_topologico","memoria":"semantica","meta":{"autor":"Gentil Lopes","descoberta":false,"dominio":"matematica","fonte":"paper_AQ.pdf, Teo.7"},"origem":"paper","palavras_chave":[],"sinonimos":[],"tipo":"conceito","titulo":"Halt Topológico (Hopf-link triplo)"}
\section{Halt Topológico (Hopf-link triplo)}
A parada da máquina é uma configuração de linking total 3 (Hopf-Hopf-Hopf) topologicamente distinguível de qualquer estado intermediário (linking 0 ou $\pm$1) — parada = irreversibilidade.
\prov{relações: usa helicidade\_quantizada}
\prov{proveniência: paper \textperiodcentered{} paper\_AQ.pdf, Teo.7 \textperiodcentered{} fato \textperiodcentered{} confiança 1.0 \textperiodcentered{} domínio matematica \textperiodcentered{} história 2 versões no jornal}

%CRISTAL {"arestas":[{"alvo":"symbolic_beta","confianca":1.0,"epistemico":"fato","origem":"paper","rel":"explica"},{"alvo":"spirit4","confianca":1.0,"epistemico":"fato","origem":"paper","rel":"referencia"}],"confianca":1.0,"contraexemplos":[],"descricao":"A função beta perturbativa de Yang-Mills octônico codifica o problema duro; conjectura-se que a série não é Borel-somável em forma fechada.","epistemico":"hipotese","exemplos":[],"id":"hardness_octonica","memoria":"semantica","meta":{"autor":"Gentil Lopes","descoberta":false,"dominio":"matematica","fonte":"paper_PQ_G_risch_formal.pdf, Conj.7.1"},"origem":"paper","palavras_chave":[],"sinonimos":[],"tipo":"conceito","titulo":"Dureza Octoniônica (Yang-Mills)"}
\section{Dureza Octoniônica (Yang-Mills)}
A função beta perturbativa de Yang-Mills octônico codifica o problema duro; conjectura-se que a série não é Borel-somável em forma fechada.
\prov{relações: explica symbolic\_beta; referencia spirit4}
\prov{proveniência: paper \textperiodcentered{} paper\_PQ\_G\_risch\_formal.pdf, Conj.7.1 \textperiodcentered{} hipotese \textperiodcentered{} confiança 1.0 \textperiodcentered{} domínio matematica \textperiodcentered{} história 3 versões no jornal}

%CRISTAL {"arestas":[{"alvo":"helicidade_quantizada","confianca":1.0,"epistemico":"fato","origem":"paper","rel":"generaliza"}],"confianca":1.0,"contraexemplos":[],"descricao":"H_Ø=∫⟨u,ω_Ø⟩dV⁷ com vorticidade octônica: invariante topológico do fluxo ideal em ℝ⁷ (generaliza Moffatt); decai como ~ν sob viscosidade.","epistemico":"fato","exemplos":[],"id":"helicidade_octonica","memoria":"semantica","meta":{"autor":"Gentil Lopes","descoberta":false,"dominio":"matematica","fonte":"paper_BC.pdf, Def.1"},"origem":"paper","palavras_chave":[],"sinonimos":[],"tipo":"conceito","titulo":"Helicidade Octônica"}
\section{Helicidade Octônica}
H\_\ensuremath{\varnothing}=\ensuremath{\int}\ensuremath{\langle}u,\ensuremath{\omega}\_\ensuremath{\varnothing}\ensuremath{\rangle}dV\ensuremath{^{7}} com vorticidade octônica: invariante topológico do fluxo ideal em \ensuremath{\mathbb{R}}\ensuremath{^{7}} (generaliza Moffatt); decai como \~{}\ensuremath{\nu} sob viscosidade.
\prov{relações: generaliza helicidade\_quantizada}
\prov{proveniência: paper \textperiodcentered{} paper\_BC.pdf, Def.1 \textperiodcentered{} fato \textperiodcentered{} confiança 1.0 \textperiodcentered{} domínio matematica \textperiodcentered{} história 2 versões no jornal}

%CRISTAL {"arestas":[{"alvo":"reconexao_cobordismo","confianca":1.0,"epistemico":"fato","origem":"paper","rel":"usa"}],"confianca":1.0,"contraexemplos":[],"descricao":"Cada reconexão entre vórtices de circulação Γ muda a helicidade H=∫u·ω em unidades discretas de 2Γ² (Moffatt 1969): estado interno da máquina ↔ helicidade, ΔH = quanta de computação.","epistemico":"fato","exemplos":[],"id":"helicidade_quantizada","memoria":"semantica","meta":{"autor":"Gentil Lopes","descoberta":false,"dominio":"matematica","fonte":"paper_AT.pdf, Teo.3"},"origem":"paper","palavras_chave":[],"sinonimos":[],"tipo":"conceito","titulo":"Helicidade Quantizada"}
\section{Helicidade Quantizada}
Cada reconexão entre vórtices de circulação \ensuremath{\Gamma} muda a helicidade H=\ensuremath{\int}u\textperiodcentered \ensuremath{\omega} em unidades discretas de 2\ensuremath{\Gamma}\ensuremath{^{2}} (Moffatt 1969): estado interno da máquina \ensuremath{\leftrightarrow} helicidade, \ensuremath{\Delta}H = quanta de computação.
\prov{relações: usa reconexao\_cobordismo}
\prov{proveniência: paper \textperiodcentered{} paper\_AT.pdf, Teo.3 \textperiodcentered{} fato \textperiodcentered{} confiança 1.0 \textperiodcentered{} domínio matematica \textperiodcentered{} história 2 versões no jornal}

%CRISTAL {"arestas":[{"alvo":"teoria_do_prospecto","confianca":1.0,"epistemico":"fato","origem":"wiki","rel":"parte_de"},{"alvo":"racionalidade_limitada","confianca":1.0,"epistemico":"fato","origem":"wiki","rel":"relaciona"}],"confianca":1.0,"contraexemplos":[],"descricao":"Atalhos mentais (disponibilidade, representatividade, ancoragem) que funcionam mas produzem erros sistemáticos. O programa de Kahneman-Tversky: o pensamento rápido (Sistema 1) é eficiente e enviesado; o lento (Sistema 2), caro e raro.","epistemico":"fato","exemplos":[],"id":"heuristicas_e_vieses","memoria":"semantica","meta":{"autor":"Kahneman-Tversky","dominio":"matematica","fonte":"teoria da decisão"},"origem":"wiki","palavras_chave":[],"sinonimos":[],"tipo":"conceito","titulo":"Heurísticas e Vieses"}
\section{Heurísticas e Vieses}
Atalhos mentais (disponibilidade, representatividade, ancoragem) que funcionam mas produzem erros sistemáticos. O programa de Kahneman-Tversky: o pensamento rápido (Sistema 1) é eficiente e enviesado; o lento (Sistema 2), caro e raro.
\prov{relações: parte\_de teoria\_do\_prospecto; relaciona racionalidade\_limitada}
\prov{proveniência: wiki \textperiodcentered{} teoria da decisão \textperiodcentered{} fato \textperiodcentered{} confiança 1.0 \textperiodcentered{} domínio matematica \textperiodcentered{} história 3 versões no jornal}

%CRISTAL {"arestas":[{"alvo":"corpo_f2k","confianca":1.0,"epistemico":"fato","origem":"paper","rel":"usa"},{"alvo":"razao_de_prata","confianca":1.0,"epistemico":"fato","origem":"paper","rel":"usa"}],"confianca":1.0,"contraexemplos":[],"descricao":"Conjectura de emissão: a prata só se emite pela estrela de prata de rank exatamente 6; rank-5 (pentagrama, ouro φ) é lastro externo, não emite.","epistemico":"hipotese","exemplos":[],"id":"hexagrama_mineracao","memoria":"semantica","meta":{"autor":"Gentil Lopes","descoberta":true,"dominio":"matematica","fonte":"kernel/cifra_ouro_branco.pdf, §IV"},"origem":"paper","palavras_chave":[],"sinonimos":[],"tipo":"conceito","titulo":"Hexagrama de Mineração (rank-6)"}
\section{Hexagrama de Mineração (rank-6)}
Conjectura de emissão: a prata só se emite pela estrela de prata de rank exatamente 6; rank-5 (pentagrama, ouro \ensuremath{\varphi}) é lastro externo, não emite.
\prov{relações: usa corpo\_f2k; usa razao\_de\_prata}
\prov{proveniência: paper \textperiodcentered{} kernel/cifra\_ouro\_branco.pdf, §IV \textperiodcentered{} hipotese \textperiodcentered{} confiança 1.0 \textperiodcentered{} domínio matematica \textperiodcentered{} história 3 versões no jornal}

%CRISTAL {"arestas":[{"alvo":"symbolic_beta","confianca":1.0,"epistemico":"fato","origem":"paper","rel":"explica"},{"alvo":"computacao","confianca":1.0,"epistemico":"fato","origem":"paper","rel":"parte_de"}],"confianca":1.0,"contraexemplos":[],"descricao":"SymbolicBeta(k) escala com k: P (k≤2), PSPACE-hard (k≥3, condicional), indecidível no limite k→∞ (redução de Risch). Graduação de confiança em k.","epistemico":"inferencia","exemplos":[],"id":"hierarquia_pspace_indecidivel","memoria":"semantica","meta":{"autor":"Gentil Lopes","descoberta":false,"dominio":"matematica","fonte":"paper_PQ_F_threshold_indecidibilidade.pdf, Teo.4.1"},"origem":"paper","palavras_chave":[],"sinonimos":[],"tipo":"conceito","titulo":"Hierarquia P→PSPACE→Indecidível"}
\section{Hierarquia P\ensuremath{\to}PSPACE\ensuremath{\to}Indecidível}
SymbolicBeta(k) escala com k: P (k\ensuremath{\le}2), PSPACE-hard (k\ensuremath{\ge}3, condicional), indecidível no limite k\ensuremath{\to}\ensuremath{\infty} (redução de Risch). Graduação de confiança em k.
\prov{relações: explica symbolic\_beta; parte\_de computacao}
\prov{proveniência: paper \textperiodcentered{} paper\_PQ\_F\_threshold\_indecidibilidade.pdf, Teo.4.1 \textperiodcentered{} inferencia \textperiodcentered{} confiança 1.0 \textperiodcentered{} domínio matematica \textperiodcentered{} história 3 versões no jornal}

%CRISTAL {"arestas":[{"alvo":"torre_gentil","confianca":1.0,"epistemico":"fato","origem":"paper","rel":"usa"},{"alvo":"word","confianca":0.5,"epistemico":"fato","origem":"wiki","rel":"relaciona"}],"confianca":1.0,"contraexemplos":[],"descricao":"Sugestão do Lopes: a torre X→X×ℙ¹→… costurada pela lemniscata (reflexão τ(x)=1−x) realizaria as classes de Hodge. Interpretação estrutural, não prova.","epistemico":"hipotese","exemplos":[],"id":"hodge_torre","memoria":"semantica","meta":{"autor":"Gentil Lopes","descoberta":false,"dominio":"matematica","fonte":"corpo_tropical.pdf, §2.13"},"origem":"paper","palavras_chave":[],"sinonimos":[],"tipo":"conceito","titulo":"Hodge via Torre de Polarizações"}
\section{Hodge via Torre de Polarizações}
Sugestão do Lopes: a torre X\ensuremath{\to}X$\times$\ensuremath{\mathbb{P}}\ensuremath{^{1}}\ensuremath{\to}… costurada pela lemniscata (reflexão \ensuremath{\tau}(x)=1\ensuremath{-}x) realizaria as classes de Hodge. Interpretação estrutural, não prova.
\prov{relações: usa torre\_gentil; relaciona word}
\prov{proveniência: paper \textperiodcentered{} corpo\_tropical.pdf, §2.13 \textperiodcentered{} hipotese \textperiodcentered{} confiança 1.0 \textperiodcentered{} domínio matematica \textperiodcentered{} história 4 versões no jornal}

%CRISTAL {"arestas":[{"alvo":"osso_carne","confianca":1.0,"epistemico":"fato","origem":"paper","rel":"explica"},{"alvo":"boneca_russa","confianca":1.0,"epistemico":"fato","origem":"paper","rel":"usa"}],"confianca":1.0,"contraexemplos":[],"descricao":"Duas carnes (ouro dim 0,694; plástico dim 0,406) são homeomorfas (mesmo Cantor) mas o homeomorfismo distorce a métrica — não preserva dimensão nem entropia. Osso comum, carne própria.","epistemico":"fato","exemplos":[],"id":"homeomorfismo_nao_bilipschitz","memoria":"semantica","meta":{"autor":"Gentil Lopes","descoberta":false,"dominio":"matematica","fonte":"paper_54_carnes_homeo.pdf, Teo.3.1"},"origem":"paper","palavras_chave":[],"sinonimos":[],"tipo":"conceito","titulo":"Homeomorfismo Não-Bi-Lipschitz"}
\section{Homeomorfismo Não-Bi-Lipschitz}
Duas carnes (ouro dim 0,694; plástico dim 0,406) são homeomorfas (mesmo Cantor) mas o homeomorfismo distorce a métrica — não preserva dimensão nem entropia. Osso comum, carne própria.
\prov{relações: explica osso\_carne; usa boneca\_russa}
\prov{proveniência: paper \textperiodcentered{} paper\_54\_carnes\_homeo.pdf, Teo.3.1 \textperiodcentered{} fato \textperiodcentered{} confiança 1.0 \textperiodcentered{} domínio matematica \textperiodcentered{} história 3 versões no jornal}

%CRISTAL {"arestas":[{"alvo":"furo_hopf","confianca":1.0,"epistemico":"fato","origem":"paper","rel":"relaciona"},{"alvo":"modos_aproximacao_s2","confianca":1.0,"epistemico":"fato","origem":"paper","rel":"usa"}],"confianca":1.0,"contraexemplos":[],"descricao":"Realização física da fibração S³→S²: cada anel de vórtice é um ponto de S² (modo) × fase de S¹; a reconexão manipula os linkings entre fibras (cobordismo de Hopf).","epistemico":"fato","exemplos":[],"id":"hopf_fibration_vortex","memoria":"semantica","meta":{"autor":"Gentil Lopes","descoberta":false,"dominio":"matematica","fonte":"paper_AT.pdf, Teo.7"},"origem":"paper","palavras_chave":[],"sinonimos":[],"tipo":"conceito","titulo":"Fibração de Hopf com Vórtices"}
\section{Fibração de Hopf com Vórtices}
Realização física da fibração S\ensuremath{^{3}}\ensuremath{\to}S\ensuremath{^{2}}: cada anel de vórtice é um ponto de S\ensuremath{^{2}} (modo) $\times$ fase de S\ensuremath{^{1}}; a reconexão manipula os linkings entre fibras (cobordismo de Hopf).
\prov{relações: relaciona furo\_hopf; usa modos\_aproximacao\_s2}
\prov{proveniência: paper \textperiodcentered{} paper\_AT.pdf, Teo.7 \textperiodcentered{} fato \textperiodcentered{} confiança 1.0 \textperiodcentered{} domínio matematica \textperiodcentered{} história 3 versões no jornal}

%CRISTAL {"arestas":[{"alvo":"contrato_objeto","confianca":1.0,"epistemico":"fato","origem":"paper","rel":"especializa"},{"alvo":"goldchain","confianca":1.0,"epistemico":"fato","origem":"paper","rel":"parte_de"},{"alvo":"sequencia_pell","confianca":1.0,"epistemico":"fato","origem":"paper","rel":"referencia"}],"confianca":1.0,"contraexemplos":[],"descricao":"Trava valor com dupla saída: quem revela a pré-imagem s (H=hash(s)) resgata antes de t; senão o emissor reembolsa após t. O mesmo h nas duas cadeias (SHA256) permite swap atômico cross-chain.","epistemico":"fato","exemplos":[],"id":"htlc","memoria":"semantica","meta":{"autor":"Gentil Lopes","descoberta":false,"dominio":"matematica","fonte":"protocolo_swap.pdf, Def.1"},"origem":"paper","palavras_chave":[],"sinonimos":[],"tipo":"conceito","titulo":"HTLC (Hash Time-Locked Contract)"}
\section{HTLC (Hash Time-Locked Contract)}
Trava valor com dupla saída: quem revela a pré-imagem s (H=hash(s)) resgata antes de t; senão o emissor reembolsa após t. O mesmo h nas duas cadeias (SHA256) permite swap atômico cross-chain.
\prov{relações: especializa contrato\_objeto; parte\_de goldchain; referencia sequencia\_pell}
\prov{proveniência: paper \textperiodcentered{} protocolo\_swap.pdf, Def.1 \textperiodcentered{} fato \textperiodcentered{} confiança 1.0 \textperiodcentered{} domínio matematica \textperiodcentered{} história 4 versões no jornal}

%CRISTAL {"arestas":[{"alvo":"parte_fria_t_zero","confianca":1.0,"epistemico":"fato","origem":"wiki","rel":"explica"},{"alvo":"decomposicao_peirce","confianca":1.0,"epistemico":"fato","origem":"wiki","rel":"usa"},{"alvo":"pointer_states_sao_o_fundamental","confianca":1.0,"epistemico":"fato","origem":"wiki","rel":"relaciona"},{"alvo":"numeros_de_pisot","confianca":1.0,"epistemico":"fato","origem":"wiki","rel":"relaciona"}],"confianca":1.0,"contraexemplos":[],"descricao":"Um idempotente satisfaz P²=P — um projetor. É o que sobrevive à decoerência (o pointer state), o limite frio T→0, e a peça da decomposição de Peirce que fatia a álgebra. Newton em X²−X converge ao projetor espectral (o idempotente do subespaço de autovalores >½). O idempotente é o congelado, o estável, o frio.","epistemico":"fato","exemplos":[],"id":"idempotente_e_o_projetor","memoria":"semantica","meta":{"dominio":"reta mineral","fonte":"parte fria e idempotentes"},"origem":"wiki","palavras_chave":[],"sinonimos":[],"tipo":"conceito","titulo":"O Idempotente é o Projetor (P²=P)"}
\section{O Idempotente é o Projetor (P\ensuremath{^{2}}=P)}
Um idempotente satisfaz P\ensuremath{^{2}}=P — um projetor. É o que sobrevive à decoerência (o pointer state), o limite frio T\ensuremath{\to}0, e a peça da decomposição de Peirce que fatia a álgebra. Newton em X\ensuremath{^{2}}\ensuremath{-}X converge ao projetor espectral (o idempotente do subespaço de autovalores >\ensuremath{\tfrac12}). O idempotente é o congelado, o estável, o frio.
\prov{relações: explica parte\_fria\_t\_zero; usa decomposicao\_peirce; relaciona pointer\_states\_sao\_o\_fundamental; relaciona numeros\_de\_pisot}
\prov{proveniência: wiki \textperiodcentered{} parte fria e idempotentes \textperiodcentered{} fato \textperiodcentered{} confiança 1.0 \textperiodcentered{} domínio reta mineral \textperiodcentered{} história 5 versões no jornal}

%CRISTAL {"arestas":[{"alvo":"idempotente_e_o_projetor","confianca":1.0,"epistemico":"fato","origem":"wiki","rel":"usa"},{"alvo":"algebra_de_peirce","confianca":1.0,"epistemico":"fato","origem":"wiki","rel":"usa"},{"alvo":"passos_da_agulha_sao_cifra","confianca":1.0,"epistemico":"fato","origem":"wiki","rel":"relaciona"},{"alvo":"lado_negro_de_newton","confianca":1.0,"epistemico":"fato","origem":"wiki","rel":"usa"}],"confianca":1.0,"contraexemplos":[],"descricao":"Os idempotentes mod n (e²≡e) detectam a FATORAÇÃO: um PRIMO tem só {0,1} (só o negro e o claro — um corpo não tem projetor não-trivial); um composto tem 2^ω idempotentes (ω = nº de primos distintos), a decomposição de Peirce/CRT. Os idempotentes são o corpo discreto de primos visto pela álgebra: primo = irredutível = só vácuo e identidade. Verif.: mod 7 → {0,1}; mod 30 → 8 idempotentes.","epistemico":"fato","exemplos":[],"id":"idempotentes_mod_n_fatoram","memoria":"semantica","meta":{"dominio":"reta mineral","fonte":"parte fria e idempotentes"},"origem":"wiki","palavras_chave":[],"sinonimos":[],"tipo":"conceito","titulo":"Os Idempotentes mod n Fatoram (o corpo de primos)"}
\section{Os Idempotentes mod n Fatoram (o corpo de primos)}
Os idempotentes mod n (e\ensuremath{^{2}}\ensuremath{\equiv}e) detectam a FATORAÇÃO: um PRIMO tem só \{0,1\} (só o negro e o claro — um corpo não tem projetor não-trivial); um composto tem 2\^{}\ensuremath{\omega} idempotentes (\ensuremath{\omega} = nº de primos distintos), a decomposição de Peirce/CRT. Os idempotentes são o corpo discreto de primos visto pela álgebra: primo = irredutível = só vácuo e identidade. Verif.: mod 7 \ensuremath{\to} \{0,1\}; mod 30 \ensuremath{\to} 8 idempotentes.
\prov{relações: usa idempotente\_e\_o\_projetor; usa algebra\_de\_peirce; relaciona passos\_da\_agulha\_sao\_cifra; usa lado\_negro\_de\_newton}
\prov{proveniência: wiki \textperiodcentered{} parte fria e idempotentes \textperiodcentered{} fato \textperiodcentered{} confiança 1.0 \textperiodcentered{} domínio reta mineral \textperiodcentered{} história 5 versões no jornal}

%CRISTAL {"arestas":[{"alvo":"grupo_metricas","confianca":1.0,"epistemico":"fato","origem":"paper","rel":"parte_de"}],"confianca":1.0,"contraexemplos":[],"descricao":"2E(k)K(k′)+2E(k′)K(k)−… = π/2 para integrais elípticas completas; no ponto lemniscático auto-dual k=1/√2 fecha a tensão régua(K)×energia(E) em π/2.","epistemico":"fato","exemplos":[],"id":"identidade_legendre","memoria":"semantica","meta":{"autor":"Gentil Lopes","descoberta":false,"dominio":"matematica","fonte":"paper_29_grupo_metricas.pdf"},"origem":"paper","palavras_chave":[],"sinonimos":[],"tipo":"conceito","titulo":"Identidade de Legendre"}
\section{Identidade de Legendre}
2E(k)K(k\ensuremath{'})+2E(k\ensuremath{'})K(k)\ensuremath{-}… = \ensuremath{\pi}/2 para integrais elípticas completas; no ponto lemniscático auto-dual k=1/\ensuremath{\sqrt{\,}}2 fecha a tensão régua(K)$\times$energia(E) em \ensuremath{\pi}/2.
\prov{relações: parte\_de grupo\_metricas}
\prov{proveniência: paper \textperiodcentered{} paper\_29\_grupo\_metricas.pdf \textperiodcentered{} fato \textperiodcentered{} confiança 1.0 \textperiodcentered{} domínio matematica \textperiodcentered{} história 2 versões no jornal}

%CRISTAL {"arestas":[{"alvo":"reta_de_ouro","confianca":1.0,"epistemico":"fato","origem":"paper","rel":"usa"},{"alvo":"cifra","confianca":1.0,"epistemico":"fato","origem":"paper","rel":"referencia"}],"confianca":1.0,"contraexemplos":[],"descricao":"A identidade de um tecido é o SHA256 do DNA lido como número na base φ (forma canônica sem '11'): endereço único, público, sem observador.","epistemico":"fato","exemplos":[],"id":"identidade_reta_ouro","memoria":"semantica","meta":{"autor":"Gentil Lopes","descoberta":false,"dominio":"matematica","fonte":"chaves.pdf, §1"},"origem":"paper","palavras_chave":[],"sinonimos":[],"tipo":"conceito","titulo":"Identidade = SHA256 na Reta de Ouro"}
\section{Identidade = SHA256 na Reta de Ouro}
A identidade de um tecido é o SHA256 do DNA lido como número na base \ensuremath{\varphi} (forma canônica sem '11'): endereço único, público, sem observador.
\prov{relações: usa reta\_de\_ouro; referencia cifra}
\prov{proveniência: paper \textperiodcentered{} chaves.pdf, §1 \textperiodcentered{} fato \textperiodcentered{} confiança 1.0 \textperiodcentered{} domínio matematica \textperiodcentered{} história 3 versões no jornal}

%CRISTAL {"arestas":[{"alvo":"subshift_do_gap","confianca":1.0,"epistemico":"fato","origem":"wiki","rel":"usa"},{"alvo":"computacao_fractal","confianca":1.0,"epistemico":"fato","origem":"wiki","rel":"explica"}],"confianca":1.0,"contraexemplos":[],"descricao":"A projeção diádica leva Xₖ ao atrator Kₖ=∪r=₀k ιr(Kₖ), com ιr(x)=(1−2−r)+2−(r⁺¹)x e renormalização Rₖ(1r0·u)=u. Cada bloco de retorno contém cópia escalada por 2−(r⁺¹) da máquina inteira — a forma exata da 'broca de brocas'.","epistemico":"fato","exemplos":[],"id":"ifs_de_retorno","memoria":"semantica","meta":{"autor":"broca-so","dominio":"fractal","fonte":"machine/computador_fractal.tex §2"},"origem":"wiki","palavras_chave":[],"sinonimos":[],"tipo":"conceito","titulo":"IFS de Retorno (a broca de brocas exata)"}
\section{IFS de Retorno (a broca de brocas exata)}
A projeção diádica leva X\ensuremath{_{k}} ao atrator K\ensuremath{_{k}}=\ensuremath{\cup}r=\ensuremath{_{0}}k \ensuremath{\iota}r(K\ensuremath{_{k}}), com \ensuremath{\iota}r(x)=(1\ensuremath{-}2\ensuremath{-}r)+2\ensuremath{-}(r\ensuremath{^{+1}})x e renormalização R\ensuremath{_{k}}(1r0\textperiodcentered u)=u. Cada bloco de retorno contém cópia escalada por 2\ensuremath{-}(r\ensuremath{^{+1}}) da máquina inteira — a forma exata da 'broca de brocas'.
\prov{relações: usa subshift\_do\_gap; explica computacao\_fractal}
\prov{proveniência: wiki \textperiodcentered{} machine/computador\_fractal.tex §2 \textperiodcentered{} fato \textperiodcentered{} confiança 1.0 \textperiodcentered{} domínio fractal \textperiodcentered{} história 4 versões no jornal}

%CRISTAL {"arestas":[{"alvo":"ilha_fractal_3d_dim","confianca":1.0,"epistemico":"fato","origem":"wiki","rel":"usa"}],"confianca":1.0,"contraexemplos":[],"descricao":"Amostragem Monte Carlo (50.000 pontos em [0.5,1.5]⁴) mede fração física 4D=0.732% — completa a série de frações da ilha das álgebras (2D→3D 0.62%→4D 0.73%): o domínio fisicamente admissível é uma fração ínfima do espaço de multiplicações possíveis.","epistemico":"fato","exemplos":[],"id":"ilha_4d_fracao","memoria":"semantica","meta":{"autor":"broca-so","dominio":"fractal","fonte":"hiper/benchmarks/fase34_fractal_3d_resultados.json"},"origem":"wiki","palavras_chave":[],"sinonimos":[],"tipo":"conceito","titulo":"Fração Física da Ilha em 4D ≈ 0.73%"}
\section{Fração Física da Ilha em 4D \ensuremath{\approx} 0.73\%}
Amostragem Monte Carlo (50.000 pontos em [0.5,1.5]\ensuremath{^{4}}) mede fração física 4D=0.732\% — completa a série de frações da ilha das álgebras (2D\ensuremath{\to}3D 0.62\%\ensuremath{\to}4D 0.73\%): o domínio fisicamente admissível é uma fração ínfima do espaço de multiplicações possíveis.
\prov{relações: usa ilha\_fractal\_3d\_dim}
\prov{proveniência: wiki \textperiodcentered{} hiper/benchmarks/fase34\_fractal\_3d\_resultados.json \textperiodcentered{} fato \textperiodcentered{} confiança 1.0 \textperiodcentered{} domínio fractal \textperiodcentered{} história 2 versões no jornal}

%CRISTAL {"arestas":[{"alvo":"ilha_fractal_spin3","confianca":1.0,"epistemico":"fato","origem":"wiki","rel":"confirma"},{"alvo":"cordas_ilha_hurwitz","confianca":1.0,"epistemico":"fato","origem":"wiki","rel":"relaciona"}],"confianca":1.0,"contraexemplos":[],"descricao":"Varredura direta de admissibilidade (grid 64³ em (t,s,u), r=−1, v=1) dá box-counting dim₃D=1.847 num espaço ambiente 3D (<3 ⇒ fractal genuíno), com fração física de apenas 0.62%. O número medido da ilha das álgebras físicas.","epistemico":"fato","exemplos":[],"id":"ilha_fractal_3d_dim","memoria":"semantica","meta":{"autor":"broca-so","dominio":"fractal","fonte":"hiper/benchmarks/fase34_fractal_3d_resultados.json"},"origem":"wiki","palavras_chave":[],"sinonimos":[],"tipo":"conceito","titulo":"A Ilha Físico-Admissível 3D tem Dimensão 1.847"}
\section{A Ilha Físico-Admissível 3D tem Dimensão 1.847}
Varredura direta de admissibilidade (grid 64\ensuremath{^{3}} em (t,s,u), r=\ensuremath{-}1, v=1) dá box-counting dim\ensuremath{_{3}}D=1.847 num espaço ambiente 3D (<3 \ensuremath{\Rightarrow} fractal genuíno), com fração física de apenas 0.62\%. O número medido da ilha das álgebras físicas.
\prov{relações: confirma ilha\_fractal\_spin3; relaciona cordas\_ilha\_hurwitz}
\prov{proveniência: wiki \textperiodcentered{} hiper/benchmarks/fase34\_fractal\_3d\_resultados.json \textperiodcentered{} fato \textperiodcentered{} confiança 1.0 \textperiodcentered{} domínio fractal \textperiodcentered{} história 3 versões no jornal}

%CRISTAL {"arestas":[{"alvo":"ilha_fractal_3d_dim","confianca":1.0,"epistemico":"fato","origem":"wiki","rel":"usa"},{"alvo":"cordas_ilha_hurwitz","confianca":1.0,"epistemico":"fato","origem":"wiki","rel":"relaciona"}],"confianca":1.0,"contraexemplos":[],"descricao":"Após erosão binária, a fronteira 3D tem dim=1.847 e fração=0.62% — valores IDÊNTICOS aos da ilha inteira; a erosão esvazia o conjunto. A ilha é uma casca fractal fina, sem interior sólido — um fio, não um volume.","epistemico":"fato","exemplos":[],"id":"ilha_toda_fronteira","memoria":"semantica","meta":{"autor":"broca-so","dominio":"fractal","fonte":"hiper/benchmarks/fase34_fractal_3d_resultados.json"},"origem":"wiki","palavras_chave":[],"sinonimos":[],"tipo":"conceito","titulo":"A Ilha é Toda Fronteira (sem interior)"}
\section{A Ilha é Toda Fronteira (sem interior)}
Após erosão binária, a fronteira 3D tem dim=1.847 e fração=0.62\% — valores IDÊNTICOS aos da ilha inteira; a erosão esvazia o conjunto. A ilha é uma casca fractal fina, sem interior sólido — um fio, não um volume.
\prov{relações: usa ilha\_fractal\_3d\_dim; relaciona cordas\_ilha\_hurwitz}
\prov{proveniência: wiki \textperiodcentered{} hiper/benchmarks/fase34\_fractal\_3d\_resultados.json \textperiodcentered{} fato \textperiodcentered{} confiança 1.0 \textperiodcentered{} domínio fractal \textperiodcentered{} história 3 versões no jornal}

%CRISTAL {"arestas":[{"alvo":"pitagoras_dual","confianca":1.0,"epistemico":"fato","origem":"paper","rel":"especializa"},{"alvo":"mecanica_quantica","confianca":1.0,"epistemico":"fato","origem":"paper","rel":"parte_de"}],"confianca":1.0,"contraexemplos":[],"descricao":"Mecânica quântica de ouro reversível: a lei a+c²=1 com c=|s| (não-comutatividade) e Planck efetivo ħ_ef=2c — sem incerteza em s=0, máxima (Heisenberg) em s=±1.","epistemico":"fato","exemplos":[],"id":"incerteza_algebrica","memoria":"semantica","meta":{"autor":"Gentil Lopes","descoberta":true,"dominio":"matematica","fonte":"mecanica.pdf, Def.1"},"origem":"paper","palavras_chave":[],"sinonimos":[],"tipo":"conceito","titulo":"Incerteza como Parábola"}
\section{Incerteza como Parábola}
Mecânica quântica de ouro reversível: a lei a+c\ensuremath{^{2}}=1 com c=|s| (não-comutatividade) e Planck efetivo \ensuremath{\hbar}\_ef=2c — sem incerteza em s=0, máxima (Heisenberg) em s=$\pm$1.
\prov{relações: especializa pitagoras\_dual; parte\_de mecanica\_quantica}
\prov{proveniência: paper \textperiodcentered{} mecanica.pdf, Def.1 \textperiodcentered{} fato \textperiodcentered{} confiança 1.0 \textperiodcentered{} domínio matematica \textperiodcentered{} história 3 versões no jornal}

%CRISTAL {"arestas":[{"alvo":"teoria_da_decisao","confianca":1.0,"epistemico":"fato","origem":"wiki","rel":"parte_de"}],"confianca":1.0,"contraexemplos":[],"descricao":"RISCO: probabilidades conhecidas (a roleta). INCERTEZA: probabilidades desconhecidas ou desconhecíveis (o futuro econômico). Knight separou os dois — e a segunda é onde mora o lucro do empreendedor e o limite dos modelos.","epistemico":"fato","exemplos":[],"id":"incerteza_vs_risco","memoria":"semantica","meta":{"autor":"Frank Knight","dominio":"matematica","fonte":"teoria da decisão"},"origem":"wiki","palavras_chave":[],"sinonimos":[],"tipo":"conceito","titulo":"Risco vs Incerteza (Knight)"}
\section{Risco vs Incerteza (Knight)}
RISCO: probabilidades conhecidas (a roleta). INCERTEZA: probabilidades desconhecidas ou desconhecíveis (o futuro econômico). Knight separou os dois — e a segunda é onde mora o lucro do empreendedor e o limite dos modelos.
\prov{relações: parte\_de teoria\_da\_decisao}
\prov{proveniência: wiki \textperiodcentered{} teoria da decisão \textperiodcentered{} fato \textperiodcentered{} confiança 1.0 \textperiodcentered{} domínio matematica \textperiodcentered{} história 2 versões no jornal}

%CRISTAL {"arestas":[{"alvo":"funcional_simetria_klein","confianca":1.0,"epistemico":"fato","origem":"paper","rel":"usa"},{"alvo":"indecidibilidade_ns_zfc","confianca":1.0,"epistemico":"fato","origem":"paper","rel":"especializa"},{"alvo":"constante_168","confianca":1.0,"epistemico":"fato","origem":"paper","rel":"usa"}],"confianca":1.0,"contraexemplos":[],"descricao":"Conjectura culminante: NS-3D seria indecidível porque reduz a um 'PSL(2,7)-problema' (Klein atrai u(t)?); a 168 seria a 'dimensão do problema'. Condicional às conjecturas de Klein-atrator (agora enfraquecidas).","epistemico":"hipotese","exemplos":[],"id":"indecidibilidade_ns_psl","memoria":"semantica","meta":{"autor":"Gentil Lopes","descoberta":false,"dominio":"matematica","fonte":"paper_AZ.pdf, Teo.16"},"origem":"paper","palavras_chave":[],"sinonimos":[],"tipo":"conceito","titulo":"Indecidibilidade de NS via PSL(2,7)"}
\section{Indecidibilidade de NS via PSL(2,7)}
Conjectura culminante: NS-3D seria indecidível porque reduz a um 'PSL(2,7)-problema' (Klein atrai u(t)?); a 168 seria a 'dimensão do problema'. Condicional às conjecturas de Klein-atrator (agora enfraquecidas).
\prov{relações: usa funcional\_simetria\_klein; especializa indecidibilidade\_ns\_zfc; usa constante\_168}
\prov{proveniência: paper \textperiodcentered{} paper\_AZ.pdf, Teo.16 \textperiodcentered{} hipotese \textperiodcentered{} confiança 1.0 \textperiodcentered{} domínio matematica \textperiodcentered{} história 4 versões no jornal}

%CRISTAL {"arestas":[{"alvo":"programa_ns_bloqueios","confianca":1.0,"epistemico":"fato","origem":"paper","rel":"usa"},{"alvo":"shift_map_ns_turing","confianca":1.0,"epistemico":"fato","origem":"paper","rel":"usa"},{"alvo":"hierarquia_pspace_indecidivel","confianca":1.0,"epistemico":"fato","origem":"paper","rel":"relaciona"},{"alvo":"p_vs_np_milenio","confianca":1.0,"epistemico":"fato","origem":"wiki","rel":"relaciona"}],"confianca":1.0,"contraexemplos":[],"descricao":"Conjectura altamente especulativa: a proposição 'NS tem solução global suave' seria indecidível em ZFC (à la Hipótese do Continuum, 10º de Hilbert, gap espectral de Cubitt). O próprio Lopes diz ser 'mais difícil que o próprio problema do milênio'.","epistemico":"hipotese","exemplos":[],"id":"indecidibilidade_ns_zfc","memoria":"semantica","meta":{"autor":"Gentil Lopes","descoberta":false,"dominio":"matematica","fonte":"paper_AJ.pdf, Conj.5"},"origem":"paper","palavras_chave":[],"sinonimos":[],"tipo":"conceito","titulo":"NS Indecidível em ZFC? (conjectura)"}
\section{NS Indecidível em ZFC? (conjectura)}
Conjectura altamente especulativa: a proposição 'NS tem solução global suave' seria indecidível em ZFC (à la Hipótese do Continuum, 10º de Hilbert, gap espectral de Cubitt). O próprio Lopes diz ser 'mais difícil que o próprio problema do milênio'.
\prov{relações: usa programa\_ns\_bloqueios; usa shift\_map\_ns\_turing; relaciona hierarquia\_pspace\_indecidivel; relaciona p\_vs\_np\_milenio}
\prov{proveniência: paper \textperiodcentered{} paper\_AJ.pdf, Conj.5 \textperiodcentered{} hipotese \textperiodcentered{} confiança 1.0 \textperiodcentered{} domínio matematica \textperiodcentered{} história 6 versões no jornal}

%CRISTAL {"arestas":[{"alvo":"teorema_esterilidade_dimensao","confianca":1.0,"epistemico":"fato","origem":"wiki","rel":"explica"},{"alvo":"reta_metalica_geral","confianca":1.0,"epistemico":"fato","origem":"wiki","rel":"usa"},{"alvo":"dimensao_de_prata_pell","confianca":1.0,"epistemico":"fato","origem":"wiki","rel":"relaciona"}],"confianca":1.0,"contraexemplos":[],"descricao":"Dimensões distintas geram corpos quadráticos DISTINTOS: σₙ gera ℚ(√(n²+4)) — ouro √5, prata √8=ℚ(√2). São independentes: (i) PSLQ NÃO acha relação inteira entre 1,φ,√2 nem 1,φ,√2,φ√2; (ii) no metal, ZERO soluções inteiras de (a+bφ)²=2c² (⇔ √2∉ℚ(φ), pela obstrução de norma a²+ab−b²). Portanto a informação da prata NÃO é derivável do ouro — é genuinamente nova. É por isso que, esgotada uma dimensão, a próxima carrega informação de fato. [D] provado/verificado (PSLQ + metal inteiro).","epistemico":"fato","exemplos":[],"id":"independencia_retas_metalicas","memoria":"semantica","meta":{"autor":"Lopes/broca-so","dominio":"fractal","fonte":"PSLQ (mpmath); doc/fractais/esterilidade_metal.c perna 3"},"origem":"wiki","palavras_chave":[],"sinonimos":[],"tipo":"conceito","titulo":"Independência das retas: ℚ(√5) ≠ ℚ(√2) (a prata não é derivável do ouro)"}
\section{Independência das retas: \ensuremath{\mathbb{Q}}(\ensuremath{\sqrt{\,}}5) \ensuremath{\ne} \ensuremath{\mathbb{Q}}(\ensuremath{\sqrt{\,}}2) (a prata não é derivável do ouro)}
Dimensões distintas geram corpos quadráticos DISTINTOS: \ensuremath{\sigma}\ensuremath{_{n}} gera \ensuremath{\mathbb{Q}}(\ensuremath{\sqrt{\,}}(n\ensuremath{^{2}}+4)) — ouro \ensuremath{\sqrt{\,}}5, prata \ensuremath{\sqrt{\,}}8=\ensuremath{\mathbb{Q}}(\ensuremath{\sqrt{\,}}2). São independentes: (i) PSLQ NÃO acha relação inteira entre 1,\ensuremath{\varphi},\ensuremath{\sqrt{\,}}2 nem 1,\ensuremath{\varphi},\ensuremath{\sqrt{\,}}2,\ensuremath{\varphi}\ensuremath{\sqrt{\,}}2; (ii) no metal, ZERO soluções inteiras de (a+b\ensuremath{\varphi})\ensuremath{^{2}}=2c\ensuremath{^{2}} (\ensuremath{\Leftrightarrow} \ensuremath{\sqrt{\,}}2\ensuremath{\notin}\ensuremath{\mathbb{Q}}(\ensuremath{\varphi}), pela obstrução de norma a\ensuremath{^{2}}+ab\ensuremath{-}b\ensuremath{^{2}}). Portanto a informação da prata NÃO é derivável do ouro — é genuinamente nova. É por isso que, esgotada uma dimensão, a próxima carrega informação de fato. [D] provado/verificado (PSLQ + metal inteiro).
\prov{relações: explica teorema\_esterilidade\_dimensao; usa reta\_metalica\_geral; relaciona dimensao\_de\_prata\_pell}
\prov{proveniência: wiki \textperiodcentered{} PSLQ (mpmath); doc/fractais/esterilidade\_metal.c perna 3 \textperiodcentered{} fato \textperiodcentered{} confiança 1.0 \textperiodcentered{} domínio fractal \textperiodcentered{} história 5 versões no jornal}

%CRISTAL {"arestas":[{"alvo":"programa_ns_bloqueios","confianca":1.0,"epistemico":"fato","origem":"paper","rel":"explica"},{"alvo":"word","confianca":0.5,"epistemico":"fato","origem":"wiki","rel":"relaciona"}],"confianca":1.0,"contraexemplos":[],"descricao":"IN(F)=#testes passados/9: mede quão 'algebricamente nua' é uma EDP; IN=1 conjectura indecidibilidade em ZFC. Ferramenta de triagem, não prova.","epistemico":"inferencia","exemplos":[],"id":"indice_nudez","memoria":"semantica","meta":{"autor":"Gentil Lopes","descoberta":false,"dominio":"matematica","fonte":"paper_AJ.pdf, Def.18"},"origem":"paper","palavras_chave":[],"sinonimos":[],"tipo":"conceito","titulo":"Índice de Nudez Algébrica"}
\section{Índice de Nudez Algébrica}
IN(F)=\#testes passados/9: mede quão 'algebricamente nua' é uma EDP; IN=1 conjectura indecidibilidade em ZFC. Ferramenta de triagem, não prova.
\prov{relações: explica programa\_ns\_bloqueios; relaciona word}
\prov{proveniência: paper \textperiodcentered{} paper\_AJ.pdf, Def.18 \textperiodcentered{} inferencia \textperiodcentered{} confiança 1.0 \textperiodcentered{} domínio matematica \textperiodcentered{} história 4 versões no jornal}

%CRISTAL {"arestas":[{"alvo":"jogo_estrategico","confianca":1.0,"epistemico":"fato","origem":"wiki","rel":"usa"}],"confianca":1.0,"contraexemplos":[],"descricao":"Quando um lado sabe mais que o outro: seleção adversa (o mercado de limões) antes do negócio, risco moral depois. A fonte de sinalização, garantias e reputação.","epistemico":"fato","exemplos":[],"id":"informacao_assimetrica","memoria":"semantica","meta":{"autor":"Akerlof","dominio":"matematica","fonte":"teoria dos jogos"},"origem":"wiki","palavras_chave":[],"sinonimos":[],"tipo":"conceito","titulo":"Informação Assimétrica"}
\section{Informação Assimétrica}
Quando um lado sabe mais que o outro: seleção adversa (o mercado de limões) antes do negócio, risco moral depois. A fonte de sinalização, garantias e reputação.
\prov{relações: usa jogo\_estrategico}
\prov{proveniência: wiki \textperiodcentered{} teoria dos jogos \textperiodcentered{} fato \textperiodcentered{} confiança 1.0 \textperiodcentered{} domínio matematica \textperiodcentered{} história 2 versões no jornal}

%CRISTAL {"arestas":[{"alvo":"conjugados_galois_parabola","confianca":1.0,"epistemico":"fato","origem":"wiki","rel":"aplicacao"},{"alvo":"parabola_algebra_conversacional","confianca":1.0,"epistemico":"fato","origem":"wiki","rel":"usa"},{"alvo":"refino_continuo_floco_metais","confianca":1.0,"epistemico":"fato","origem":"wiki","rel":"relaciona"},{"alvo":"absorcao_iterativa_poca","confianca":1.0,"epistemico":"fato","origem":"wiki","rel":"especializa"},{"alvo":"plano_continuo_transicoes","confianca":1.0,"epistemico":"fato","origem":"wiki","rel":"usa"},{"alvo":"tiffany","confianca":1.0,"epistemico":"fato","origem":"wiki","rel":"parte_de"}],"confianca":1.0,"contraexemplos":[],"descricao":"conversa/conhecimento/transicoes.py (absorver_iterativo): a ingestão de uma obra passou a ser o MESMO refino do floco→metais, só muda o símbolo. Cada fragmento cascateia pelas arestas candidatas; em cada bacia a parábola a+c²=1 o CLASSIFICA em torno do piso θ (dinâmico pela densidade): CRISTAL (cristaliza limpo) e BORDA (a fronteira crítica de Salem — antes ia ao purgatório, agora é ABSORVIDA, alargando o tecido) entram na aresta; CAOS (o lado-a/calor) é rejeitado e cascateia; o que não achou casa em nenhum candidato vai ao PURGATÓRIO (o resíduo-a persistente, retentado quando o grafo cresce). Retorna contagem por classe (cristal/borda/purgatório). A obra é o pacote; as arestas são os metais; a parábola é o refino. [D] ingestão = refino posicional.","epistemico":"fato","exemplos":[],"id":"ingestao_refino_parabola","memoria":"semantica","meta":{"autor":"Lopes/broca-so","dominio":"algebra","fonte":"conversa/conhecimento/transicoes.py:absorver_iterativo; parabola_algebra.classe_parabola"},"origem":"wiki","palavras_chave":[],"sinonimos":[],"tipo":"conceito","titulo":"A ingestão é o REFINO da obra — a parábola classifica cada fragmento (cristal/borda/caos)"}
\section{A ingestão é o REFINO da obra — a parábola classifica cada fragmento (cristal/borda/caos)}
conversa/conhecimento/transicoes.py (absorver\_iterativo): a ingestão de uma obra passou a ser o MESMO refino do floco\ensuremath{\to}metais, só muda o símbolo. Cada fragmento cascateia pelas arestas candidatas; em cada bacia a parábola a+c\ensuremath{^{2}}=1 o CLASSIFICA em torno do piso \ensuremath{\theta} (dinâmico pela densidade): CRISTAL (cristaliza limpo) e BORDA (a fronteira crítica de Salem — antes ia ao purgatório, agora é ABSORVIDA, alargando o tecido) entram na aresta; CAOS (o lado-a/calor) é rejeitado e cascateia; o que não achou casa em nenhum candidato vai ao PURGATÓRIO (o resíduo-a persistente, retentado quando o grafo cresce). Retorna contagem por classe (cristal/borda/purgatório). A obra é o pacote; as arestas são os metais; a parábola é o refino. [D] ingestão = refino posicional.
\prov{relações: aplicacao conjugados\_galois\_parabola; usa parabola\_algebra\_conversacional; relaciona refino\_continuo\_floco\_metais; especializa absorcao\_iterativa\_poca; usa plano\_continuo\_transicoes; parte\_de tiffany}
\prov{proveniência: wiki \textperiodcentered{} conversa/conhecimento/transicoes.py:absorver\_iterativo; parabola\_algebra.classe\_parabola \textperiodcentered{} fato \textperiodcentered{} confiança 1.0 \textperiodcentered{} domínio algebra \textperiodcentered{} história 134 versões no jornal}

%CRISTAL {"arestas":[{"alvo":"teoria_gauge_octonica","confianca":1.0,"epistemico":"fato","origem":"paper","rel":"parte_de"},{"alvo":"word","confianca":0.5,"epistemico":"fato","origem":"wiki","rel":"relaciona"}],"confianca":1.0,"contraexemplos":[],"descricao":"Configurações anti-self-duais em ℝ⁷ com número topológico n∈π₃(G₂)=ℤ e ação S=8π²|n|/g²: setor não-perturbativo do vácuo octônico.","epistemico":"inferencia","exemplos":[],"id":"instantons_g2","memoria":"semantica","meta":{"autor":"Gentil Lopes","descoberta":false,"dominio":"matematica","fonte":"paper_AF_yang_mills_octonios.pdf, Teo.34"},"origem":"paper","palavras_chave":[],"sinonimos":[],"tipo":"conceito","titulo":"Instantons Octoniônicos (G₂)"}
\section{Instantons Octoniônicos (G\ensuremath{_{2}})}
Configurações anti-self-duais em \ensuremath{\mathbb{R}}\ensuremath{^{7}} com número topológico n\ensuremath{\in}\ensuremath{\pi}\ensuremath{_{3}}(G\ensuremath{_{2}})=\ensuremath{\mathbb{Z}} e ação S=8\ensuremath{\pi}\ensuremath{^{2}}|n|/g\ensuremath{^{2}}: setor não-perturbativo do vácuo octônico.
\prov{relações: parte\_de teoria\_gauge\_octonica; relaciona word}
\prov{proveniência: paper \textperiodcentered{} paper\_AF\_yang\_mills\_octonios.pdf, Teo.34 \textperiodcentered{} inferencia \textperiodcentered{} confiança 1.0 \textperiodcentered{} domínio matematica \textperiodcentered{} história 3 versões no jornal}

%CRISTAL {"arestas":[{"alvo":"racionais","confianca":1.0,"epistemico":"fato","origem":"curriculo","rel":"especializa"}],"confianca":1.0,"contraexemplos":[],"descricao":"Lopes/18 07.1999). Sejam j e n inteiros positivos arbitrariamente fixa- dos. Temos que n 2 j−1 ∈Z ⇐⇒ \u0012n −1 2 j−1 \u0013 e \u0012 n 2 j−1 \u0013 têm paridades distintas. Antes da prova do","epistemico":"fato","exemplos":[],"id":"inteiros","memoria":"semantica","meta":{"arquivo":"gentil/Novas_Sequencias_Aritmeticas_e_Geometric.pdf","ciencia":"teoria_dos_numeros","dominio":"matematica","duplicata_sugerida":[],"nivel":4,"numero":"11","tipo_no":"paper","tipo_teorema":"lema","verificacao":"parte de conceito mais amplo 'inteiros'"},"origem":"gentil/Novas_Sequencias_Aritmeticas_e_Geometric.pdf","palavras_chave":[],"sinonimos":["ℤ","inteiros"],"tipo":"conceito","titulo":"inteiros"}
\section{inteiros}
Lopes/18 07.1999). Sejam j e n inteiros positivos arbitrariamente fixa- dos. Temos que n 2 j\ensuremath{-}1 \ensuremath{\in}Z \ensuremath{\Leftarrow}\ensuremath{\Rightarrow} n \ensuremath{-}1 2 j\ensuremath{-}1  e  n 2 j\ensuremath{-}1  têm paridades distintas. Antes da prova do
\prov{sinónimos: \ensuremath{\mathbb{Z}}; inteiros}
\prov{relações: especializa racionais}
\prov{proveniência: gentil/Novas\_Sequencias\_Aritmeticas\_e\_Geometric.pdf \textperiodcentered{} fato \textperiodcentered{} confiança 1.0 \textperiodcentered{} domínio matematica \textperiodcentered{} história 6 versões no jornal}

%CRISTAL {"arestas":[{"alvo":"tricotomia_modular","confianca":1.0,"epistemico":"fato","origem":"paper","rel":"parte_de"}],"confianca":1.0,"contraexemplos":[],"descricao":"Involução de Möbius do gerador S=[0,−1;1,0] de SL₂(ℤ); leva a hipérbole retangular x²−y²=1 na lemniscata de Bernoulli (u²+v²)²=u²−v² e fixa o ponto i.","epistemico":"fato","exemplos":[],"id":"inversao_geometrica_s","memoria":"semantica","meta":{"autor":"Gentil Lopes","descoberta":true,"dominio":"matematica","fonte":"paper_27_ponte_modular.pdf, Prop.5.1"},"origem":"paper","palavras_chave":[],"sinonimos":[],"tipo":"conceito","titulo":"Inversão Geométrica (z ↦ −1/z)"}
\section{Inversão Geométrica (z \ensuremath{\mapsto} \ensuremath{-}1/z)}
Involução de Möbius do gerador S=[0,\ensuremath{-}1;1,0] de SL\ensuremath{_{2}}(\ensuremath{\mathbb{Z}}); leva a hipérbole retangular x\ensuremath{^{2}}\ensuremath{-}y\ensuremath{^{2}}=1 na lemniscata de Bernoulli (u\ensuremath{^{2}}+v\ensuremath{^{2}})\ensuremath{^{2}}=u\ensuremath{^{2}}\ensuremath{-}v\ensuremath{^{2}} e fixa o ponto i.
\prov{relações: parte\_de tricotomia\_modular}
\prov{proveniência: paper \textperiodcentered{} paper\_27\_ponte\_modular.pdf, Prop.5.1 \textperiodcentered{} fato \textperiodcentered{} confiança 1.0 \textperiodcentered{} domínio matematica \textperiodcentered{} história 2 versões no jornal}

%CRISTAL {"arestas":[{"alvo":"isolamento_estrutural","confianca":1.0,"epistemico":"fato","origem":"paper","rel":"usa"},{"alvo":"acoplamento_cruzado","confianca":1.0,"epistemico":"fato","origem":"paper","rel":"relaciona"}],"confianca":1.0,"contraexemplos":[],"descricao":"A mensagem entre processos viaja pelo canal Eᵢⱼ como roteamento multi-hop com composição EᵢⱼEⱼₖ=Eᵢₖ e vazamento ≈10⁻¹⁶.","epistemico":"fato","exemplos":[],"id":"ipc_multihop","memoria":"semantica","meta":{"autor":"Gentil Lopes","descoberta":false,"dominio":"matematica","fonte":"so_fractal.pdf"},"origem":"paper","palavras_chave":[],"sinonimos":[],"tipo":"conceito","titulo":"IPC Multi-hop (EᵢⱼEⱼₖ=Eᵢₖ)"}
\section{IPC Multi-hop (E\ensuremath{_{ij}}E\ensuremath{_{jk}}=E\ensuremath{_{ik}})}
A mensagem entre processos viaja pelo canal E\ensuremath{_{ij}} como roteamento multi-hop com composição E\ensuremath{_{ij}}E\ensuremath{_{jk}}=E\ensuremath{_{ik}} e vazamento \ensuremath{\approx}10\ensuremath{^{-16}}.
\prov{relações: usa isolamento\_estrutural; relaciona acoplamento\_cruzado}
\prov{proveniência: paper \textperiodcentered{} so\_fractal.pdf \textperiodcentered{} fato \textperiodcentered{} confiança 1.0 \textperiodcentered{} domínio matematica \textperiodcentered{} história 3 versões no jornal}

%CRISTAL {"arestas":[{"alvo":"fronteira_hurwitz","confianca":1.0,"epistemico":"fato","origem":"paper","rel":"usa"},{"alvo":"word","confianca":0.5,"epistemico":"fato","origem":"wiki","rel":"relaciona"}],"confianca":1.0,"contraexemplos":[],"descricao":"As 392 combinações binárias de S (dim 16) que são divisores de zero: raras (0,6%), com núcleo ker(L_z) de dimensão exatamente 4 (um ℍ) que colapsa a zero.","epistemico":"fato","exemplos":[],"id":"irracionais_algebricos_s","memoria":"semantica","meta":{"autor":"Gentil Lopes","descoberta":true,"dominio":"matematica","fonte":"paper_21_irracionais.pdf, Def.4.2"},"origem":"paper","palavras_chave":[],"sinonimos":[],"tipo":"conceito","titulo":"Irracionais Algébricos de S"}
\section{Irracionais Algébricos de S}
As 392 combinações binárias de S (dim 16) que são divisores de zero: raras (0,6\%), com núcleo ker(L\_z) de dimensão exatamente 4 (um \ensuremath{\mathbb{H}}) que colapsa a zero.
\prov{relações: usa fronteira\_hurwitz; relaciona word}
\prov{proveniência: paper \textperiodcentered{} paper\_21\_irracionais.pdf, Def.4.2 \textperiodcentered{} fato \textperiodcentered{} confiança 1.0 \textperiodcentered{} domínio matematica \textperiodcentered{} história 3 versões no jornal}

%CRISTAL {"arestas":[{"alvo":"parede_aco","confianca":1.0,"epistemico":"fato","origem":"paper","rel":"usa"},{"alvo":"transicao_fase_lambda1","confianca":1.0,"epistemico":"fato","origem":"paper","rel":"relaciona"},{"alvo":"word","confianca":0.5,"epistemico":"fato","origem":"wiki","rel":"relaciona"}],"confianca":1.0,"contraexemplos":[],"descricao":"A cadeia de quartos acoplados por porta ε=e−²K é o modelo de Ising 1D; a cúspide ε=0 é K→∞, ou seja a temperatura crítica Tc=0.","epistemico":"fato","exemplos":[],"id":"ising_1d_critico","memoria":"semantica","meta":{"autor":"Gentil Lopes","descoberta":false,"dominio":"matematica","fonte":"paper_74_ponto_critico.pdf, Prop.1.1"},"origem":"paper","palavras_chave":[],"sinonimos":[],"tipo":"conceito","titulo":"Ising 1D — Cúspide é T_c"}
\section{Ising 1D — Cúspide é T\_c}
A cadeia de quartos acoplados por porta \ensuremath{\varepsilon}=e\ensuremath{-}\ensuremath{^{2}}K é o modelo de Ising 1D; a cúspide \ensuremath{\varepsilon}=0 é K\ensuremath{\to}\ensuremath{\infty}, ou seja a temperatura crítica Tc=0.
\prov{relações: usa parede\_aco; relaciona transicao\_fase\_lambda1; relaciona word}
\prov{proveniência: paper \textperiodcentered{} paper\_74\_ponto\_critico.pdf, Prop.1.1 \textperiodcentered{} fato \textperiodcentered{} confiança 1.0 \textperiodcentered{} domínio matematica \textperiodcentered{} história 5 versões no jornal}

%CRISTAL {"arestas":[{"alvo":"decomposicao_peirce","confianca":1.0,"epistemico":"fato","origem":"paper","rel":"parte_de"}],"confianca":1.0,"contraexemplos":[],"descricao":"Dois processos distintos satisfazem eᵢeⱼ=0 — não há vazamento por construção; a segurança do broca-os nasce da ortogonalidade da decomposição de Peirce.","epistemico":"fato","exemplos":[],"id":"isolamento_estrutural","memoria":"semantica","meta":{"autor":"Gentil Lopes","descoberta":false,"dominio":"matematica","fonte":"so_fractal.pdf"},"origem":"paper","palavras_chave":[],"sinonimos":[],"tipo":"conceito","titulo":"Isolamento Estrutural (eᵢeⱼ=0)"}
\section{Isolamento Estrutural (e\ensuremath{_{i}}e\ensuremath{_{j}}=0)}
Dois processos distintos satisfazem e\ensuremath{_{i}}e\ensuremath{_{j}}=0 — não há vazamento por construção; a segurança do broca-os nasce da ortogonalidade da decomposição de Peirce.
\prov{relações: parte\_de decomposicao\_peirce}
\prov{proveniência: paper \textperiodcentered{} so\_fractal.pdf \textperiodcentered{} fato \textperiodcentered{} confiança 1.0 \textperiodcentered{} domínio matematica \textperiodcentered{} história 2 versões no jornal}

%CRISTAL {"arestas":[{"alvo":"jogo_estrategico","confianca":1.0,"epistemico":"fato","origem":"wiki","rel":"parte_de"}],"confianca":1.0,"contraexemplos":[],"descricao":"O agente que escolhe a estratégia que maximiza seu payoff DADO o que espera dos outros. A racionalidade é comum conhecimento — cada um sabe que o outro também otimiza.","epistemico":"fato","exemplos":[],"id":"jogador_racional","memoria":"semantica","meta":{"dominio":"matematica","fonte":"teoria dos jogos"},"origem":"wiki","palavras_chave":[],"sinonimos":[],"tipo":"conceito","titulo":"Jogador Racional"}
\section{Jogador Racional}
O agente que escolhe a estratégia que maximiza seu payoff DADO o que espera dos outros. A racionalidade é comum conhecimento — cada um sabe que o outro também otimiza.
\prov{relações: parte\_de jogo\_estrategico}
\prov{proveniência: wiki \textperiodcentered{} teoria dos jogos \textperiodcentered{} fato \textperiodcentered{} confiança 1.0 \textperiodcentered{} domínio matematica \textperiodcentered{} história 2 versões no jornal}

%CRISTAL {"arestas":[{"alvo":"jogo_estrategico","confianca":1.0,"epistemico":"fato","origem":"wiki","rel":"especializa"}],"confianca":1.0,"contraexemplos":[],"descricao":"Quando os jogadores podem formar acordos vinculantes e repartir o ganho da coalizão. A pergunta vira: como dividir de forma justa e estável o que a cooperação cria?","epistemico":"inferencia","exemplos":[],"id":"jogo_cooperativo","memoria":"semantica","meta":{"dominio":"matematica","fonte":"teoria dos jogos"},"origem":"wiki","palavras_chave":[],"sinonimos":[],"tipo":"conceito","titulo":"Jogos Cooperativos e Coalizões"}
\section{Jogos Cooperativos e Coalizões}
Quando os jogadores podem formar acordos vinculantes e repartir o ganho da coalizão. A pergunta vira: como dividir de forma justa e estável o que a cooperação cria?
\prov{relações: especializa jogo\_estrategico}
\prov{proveniência: wiki \textperiodcentered{} teoria dos jogos \textperiodcentered{} inferencia \textperiodcentered{} confiança 1.0 \textperiodcentered{} domínio matematica \textperiodcentered{} história 2 versões no jornal}

%CRISTAL {"arestas":[{"alvo":"jogo_estrategico","confianca":1.0,"epistemico":"fato","origem":"wiki","rel":"especializa"}],"confianca":1.0,"contraexemplos":[],"descricao":"O ganho de um é exatamente a perda do outro — competição pura, nada a repartir. Onde o minimax reina (xadrez, pôquer heads-up). O oposto dos jogos de soma positiva (cooperação cria valor).","epistemico":"fato","exemplos":[],"id":"jogo_de_soma_zero","memoria":"semantica","meta":{"dominio":"matematica","fonte":"teoria dos jogos"},"origem":"wiki","palavras_chave":[],"sinonimos":[],"tipo":"conceito","titulo":"Jogo de Soma Zero"}
\section{Jogo de Soma Zero}
O ganho de um é exatamente a perda do outro — competição pura, nada a repartir. Onde o minimax reina (xadrez, pôquer heads-up). O oposto dos jogos de soma positiva (cooperação cria valor).
\prov{relações: especializa jogo\_estrategico}
\prov{proveniência: wiki \textperiodcentered{} teoria dos jogos \textperiodcentered{} fato \textperiodcentered{} confiança 1.0 \textperiodcentered{} domínio matematica \textperiodcentered{} história 2 versões no jornal}

%CRISTAL {"arestas":[{"alvo":"jogo_estrategico","confianca":1.0,"epistemico":"fato","origem":"wiki","rel":"especializa"}],"confianca":1.0,"contraexemplos":[],"descricao":"Dois carros em rota de colisão: quem desvia é o 'covarde', mas se nenhum desvia, ambos morrem. A lógica do brinkmanship e da dissuasão nuclear — a ameaça crível de não ceder.","epistemico":"inferencia","exemplos":[],"id":"jogo_do_covarde","memoria":"semantica","meta":{"dominio":"matematica","fonte":"teoria dos jogos"},"origem":"wiki","palavras_chave":[],"sinonimos":[],"tipo":"conceito","titulo":"Jogo do Covarde (Chicken)"}
\section{Jogo do Covarde (Chicken)}
Dois carros em rota de colisão: quem desvia é o 'covarde', mas se nenhum desvia, ambos morrem. A lógica do brinkmanship e da dissuasão nuclear — a ameaça crível de não ceder.
\prov{relações: especializa jogo\_estrategico}
\prov{proveniência: wiki \textperiodcentered{} teoria dos jogos \textperiodcentered{} inferencia \textperiodcentered{} confiança 1.0 \textperiodcentered{} domínio matematica \textperiodcentered{} história 2 versões no jornal}

%CRISTAL {"arestas":[{"alvo":"jogo_estrategico","confianca":1.0,"epistemico":"fato","origem":"wiki","rel":"especializa"}],"confianca":1.0,"contraexemplos":[],"descricao":"Um propõe dividir dinheiro, o outro aceita ou rejeita (e aí ninguém ganha). A teoria diz 'aceite qualquer coisa'; humanos rejeitam ofertas injustas — a evidência de que justiça é preferência, não só payoff.","epistemico":"fato","exemplos":[],"id":"jogo_do_ultimato","memoria":"semantica","meta":{"dominio":"matematica","fonte":"teoria dos jogos"},"origem":"wiki","palavras_chave":[],"sinonimos":[],"tipo":"conceito","titulo":"Jogo do Ultimato"}
\section{Jogo do Ultimato}
Um propõe dividir dinheiro, o outro aceita ou rejeita (e aí ninguém ganha). A teoria diz 'aceite qualquer coisa'; humanos rejeitam ofertas injustas — a evidência de que justiça é preferência, não só payoff.
\prov{relações: especializa jogo\_estrategico}
\prov{proveniência: wiki \textperiodcentered{} teoria dos jogos \textperiodcentered{} fato \textperiodcentered{} confiança 1.0 \textperiodcentered{} domínio matematica \textperiodcentered{} história 2 versões no jornal}

%CRISTAL {"arestas":[{"alvo":"teoria_dos_jogos","confianca":1.0,"epistemico":"fato","origem":"wiki","rel":"parte_de"}],"confianca":1.0,"contraexemplos":[],"descricao":"Uma situação de interação onde o resultado de cada agente depende das escolhas de TODOS — jogadores, estratégias e payoffs. O objeto que a teoria dos jogos estuda: a decisão que antecipa a decisão do outro.","epistemico":"fato","exemplos":[],"id":"jogo_estrategico","memoria":"semantica","meta":{"dominio":"matematica","fonte":"teoria dos jogos"},"origem":"wiki","palavras_chave":[],"sinonimos":[],"tipo":"conceito","titulo":"Jogo (Teoria dos Jogos)"}
\section{Jogo (Teoria dos Jogos)}
Uma situação de interação onde o resultado de cada agente depende das escolhas de TODOS — jogadores, estratégias e payoffs. O objeto que a teoria dos jogos estuda: a decisão que antecipa a decisão do outro.
\prov{relações: parte\_de teoria\_dos\_jogos}
\prov{proveniência: wiki \textperiodcentered{} teoria dos jogos \textperiodcentered{} fato \textperiodcentered{} confiança 1.0 \textperiodcentered{} domínio matematica \textperiodcentered{} história 2 versões no jornal}

%CRISTAL {"arestas":[{"alvo":"dilema_do_prisioneiro","confianca":1.0,"epistemico":"fato","origem":"wiki","rel":"usa"}],"confianca":1.0,"contraexemplos":[],"descricao":"O mesmo jogo jogado muitas vezes muda tudo: a sombra do futuro permite a cooperação (o folk theorem). Trair hoje custa a retaliação amanhã — a reputação vira ativo.","epistemico":"fato","exemplos":[],"id":"jogos_repetidos","memoria":"semantica","meta":{"dominio":"matematica","fonte":"teoria dos jogos"},"origem":"wiki","palavras_chave":[],"sinonimos":[],"tipo":"conceito","titulo":"Jogos Repetidos"}
\section{Jogos Repetidos}
O mesmo jogo jogado muitas vezes muda tudo: a sombra do futuro permite a cooperação (o folk theorem). Trair hoje custa a retaliação amanhã — a reputação vira ativo.
\prov{relações: usa dilema\_do\_prisioneiro}
\prov{proveniência: wiki \textperiodcentered{} teoria dos jogos \textperiodcentered{} fato \textperiodcentered{} confiança 1.0 \textperiodcentered{} domínio matematica \textperiodcentered{} história 2 versões no jornal}

%CRISTAL {"arestas":[{"alvo":"constante_168","confianca":1.0,"epistemico":"fato","origem":"paper","rel":"usa"},{"alvo":"garrafa_klein_algebra","confianca":1.0,"epistemico":"fato","origem":"paper","rel":"relaciona"}],"confianca":1.0,"contraexemplos":[],"descricao":"A curva X³Y+Y³Z+Z³X=0 de gênero 3, com tesselação heptagonal (24 heptágonos, 84 arestas, 56 vértices) e exatamente 168 automorfismos — atinge o limite de Hurwitz.","epistemico":"fato","exemplos":[],"id":"klein_quartic","memoria":"semantica","meta":{"autor":"Gentil Lopes","descoberta":false,"dominio":"matematica","fonte":"paper_AU.pdf (Klein 1879)"},"origem":"paper","palavras_chave":[],"sinonimos":[],"tipo":"conceito","titulo":"Quártica de Klein"}
\section{Quártica de Klein}
A curva X\ensuremath{^{3}}Y+Y\ensuremath{^{3}}Z+Z\ensuremath{^{3}}X=0 de gênero 3, com tesselação heptagonal (24 heptágonos, 84 arestas, 56 vértices) e exatamente 168 automorfismos — atinge o limite de Hurwitz.
\prov{relações: usa constante\_168; relaciona garrafa\_klein\_algebra}
\prov{proveniência: paper \textperiodcentered{} paper\_AU.pdf (Klein 1879) \textperiodcentered{} fato \textperiodcentered{} confiança 1.0 \textperiodcentered{} domínio matematica \textperiodcentered{} história 3 versões no jornal}

%CRISTAL {"arestas":[{"alvo":"floco_phi","confianca":1.0,"epistemico":"fato","origem":"wiki","rel":"usa"},{"alvo":"koch_codifica_cantorfibonacci_prova_e_evidencia","confianca":1.0,"epistemico":"fato","origem":"wiki","rel":"contradiz"},{"alvo":"mede_flocos_phi","confianca":1.0,"epistemico":"fato","origem":"wiki","rel":"relaciona"},{"alvo":"colapso_ping_pong_overlap","confianca":1.0,"epistemico":"fato","origem":"wiki","rel":"relaciona"}],"confianca":1.0,"contraexemplos":[],"descricao":"O experimento koch-memória REFUTOU Koch como codificador: corr(w,endereço)≈−0,12, Δendereço≈16 bits entre vizinhos, 18.548 colisões de endereço com c² distinto. Logo o floco guarda o que a leitura associativa PERDEU (a dissipação), em vez de codificá-lo. Corrige o título afirmativo de koch_codifica.","epistemico":"fato","exemplos":[],"id":"koch_dissipa_nao_codifica","memoria":"semantica","meta":{"autor":"broca-so","dominio":"fractal","fonte":"papers/floco_de_neve.tex §II,§VII"},"origem":"wiki","palavras_chave":[],"sinonimos":[],"tipo":"conceito","titulo":"Koch Dissipa, não Codifica"}
\section{Koch Dissipa, não Codifica}
O experimento koch-memória REFUTOU Koch como codificador: corr(w,endereço)\ensuremath{\approx}\ensuremath{-}0,12, \ensuremath{\Delta}endereço\ensuremath{\approx}16 bits entre vizinhos, 18.548 colisões de endereço com c\ensuremath{^{2}} distinto. Logo o floco guarda o que a leitura associativa PERDEU (a dissipação), em vez de codificá-lo. Corrige o título afirmativo de koch\_codifica.
\prov{relações: usa floco\_phi; contradiz koch\_codifica\_cantorfibonacci\_prova\_e\_evidencia; relaciona mede\_flocos\_phi; relaciona colapso\_ping\_pong\_overlap}
\prov{proveniência: wiki \textperiodcentered{} papers/floco\_de\_neve.tex §II,§VII \textperiodcentered{} fato \textperiodcentered{} confiança 1.0 \textperiodcentered{} domínio fractal \textperiodcentered{} história 5 versões no jornal}

%CRISTAL {"arestas":[{"alvo":"difusao_fluido_atratores","confianca":1.0,"epistemico":"fato","origem":"wiki","rel":"usa"},{"alvo":"absorcao_ressonancia_metal","confianca":1.0,"epistemico":"fato","origem":"wiki","rel":"usa"},{"alvo":"aresta_transicao_graduada","confianca":1.0,"epistemico":"fato","origem":"wiki","rel":"usa"},{"alvo":"plano_continuo_transicoes","confianca":1.0,"epistemico":"fato","origem":"wiki","rel":"parte_de"},{"alvo":"tiffany","confianca":1.0,"epistemico":"fato","origem":"wiki","rel":"parte_de"}],"confianca":1.0,"contraexemplos":[],"descricao":"O ciclo completo da absorção, provado ponta a ponta: uma OBRA limpa (LaTeX/markup removido) é fragmentada em prosa → cada fragmento é impresso e roteado no METAL (absorve_metal.c: ressonância + difusão-fluido à aresta-casa, com posição s no contínuo) → gravado como TRANSIÇÃO camada 2 no tecido 'transicoes' (modo multi: várias frases por aresta) → reindexado → o COLAPSO pousa na prosa absorvida. Teste (corpus_aarao limpo, obra do Lopes): 6000 fragmentos → 2935 gravados (piso de qualidade ponte≥190) em 124 arestas, ~24 frases/aresta (contínuo DENSO onde a obra ressoa). Verificação: 'o universo dentro da multiplicação' → colapso pousa no fragmento exato (ov=250). Assim 'absorver uma obra' preenche o contínuo em escala, no metal. Ressalva honesta: o piso dropa ruído e o pouso segue a RESSONÂNCIA (posicional), não a semântica — de propósito (nada de kernel/esterilizar). [D] roda; grafo íntegro (62/62).","epistemico":"fato","exemplos":[],"id":"laco_absorcao_fechado","memoria":"semantica","meta":{"autor":"Lopes/broca-so","dominio":"continuo","fonte":"doc/continuo/absorve_metal.c + conversa/conhecimento/transicoes.py (escrever multi); conversa/dados/transicoes.tsv"},"origem":"wiki","palavras_chave":[],"sinonimos":[],"tipo":"conceito","titulo":"Laço fechado: obra → fragmenta → ressoa+difunde no metal → grava no contínuo → colapso pousa"}
\section{Laço fechado: obra \ensuremath{\to} fragmenta \ensuremath{\to} ressoa+difunde no metal \ensuremath{\to} grava no contínuo \ensuremath{\to} colapso pousa}
O ciclo completo da absorção, provado ponta a ponta: uma OBRA limpa (LaTeX/markup removido) é fragmentada em prosa \ensuremath{\to} cada fragmento é impresso e roteado no METAL (absorve\_metal.c: ressonância + difusão-fluido à aresta-casa, com posição s no contínuo) \ensuremath{\to} gravado como TRANSIÇÃO camada 2 no tecido 'transicoes' (modo multi: várias frases por aresta) \ensuremath{\to} reindexado \ensuremath{\to} o COLAPSO pousa na prosa absorvida. Teste (corpus\_aarao limpo, obra do Lopes): 6000 fragmentos \ensuremath{\to} 2935 gravados (piso de qualidade ponte\ensuremath{\ge}190) em 124 arestas, \~{}24 frases/aresta (contínuo DENSO onde a obra ressoa). Verificação: 'o universo dentro da multiplicação' \ensuremath{\to} colapso pousa no fragmento exato (ov=250). Assim 'absorver uma obra' preenche o contínuo em escala, no metal. Ressalva honesta: o piso dropa ruído e o pouso segue a RESSONÂNCIA (posicional), não a semântica — de propósito (nada de kernel/esterilizar). [D] roda; grafo íntegro (62/62).
\prov{relações: usa difusao\_fluido\_atratores; usa absorcao\_ressonancia\_metal; usa aresta\_transicao\_graduada; parte\_de plano\_continuo\_transicoes; parte\_de tiffany}
\prov{proveniência: wiki \textperiodcentered{} doc/continuo/absorve\_metal.c + conversa/conhecimento/transicoes.py (escrever multi); conversa/dados/transicoes.tsv \textperiodcentered{} fato \textperiodcentered{} confiança 1.0 \textperiodcentered{} domínio continuo \textperiodcentered{} história 74 versões no jornal}

%CRISTAL {"arestas":[{"alvo":"idempotente_e_o_projetor","confianca":1.0,"epistemico":"fato","origem":"wiki","rel":"usa"},{"alvo":"julia","confianca":1.0,"epistemico":"fato","origem":"wiki","rel":"relaciona"},{"alvo":"newton_acha_o_grupo_discreto","confianca":1.0,"epistemico":"fato","origem":"wiki","rel":"relaciona"},{"alvo":"rastro_da_agulha_e_fractal","confianca":1.0,"epistemico":"fato","origem":"wiki","rel":"relaciona"}],"confianca":1.0,"contraexemplos":[],"descricao":"Newton dos idempotentes tem dois destinos: o 1 (a identidade, o claro) e o 0 (o LADO NEGRO — o aniquilador, o vácuo, o nada). Escalar: N(z)=z²/(2z−1) → 0 se Re z<½, → 1 se Re z>½; a fronteira é a reta Re z=½ (grau 2 → reta, não fractal, como a reta mineral). O lado negro é a bacia do zero: para onde caem os estados que a máquina aniquila.","epistemico":"fato","exemplos":[],"id":"lado_negro_de_newton","memoria":"semantica","meta":{"dominio":"reta mineral","fonte":"parte fria e idempotentes"},"origem":"wiki","palavras_chave":[],"sinonimos":[],"tipo":"conceito","titulo":"O Lado Negro do Fractal de Newton"}
\section{O Lado Negro do Fractal de Newton}
Newton dos idempotentes tem dois destinos: o 1 (a identidade, o claro) e o 0 (o LADO NEGRO — o aniquilador, o vácuo, o nada). Escalar: N(z)=z\ensuremath{^{2}}/(2z\ensuremath{-}1) \ensuremath{\to} 0 se Re z<\ensuremath{\tfrac12}, \ensuremath{\to} 1 se Re z>\ensuremath{\tfrac12}; a fronteira é a reta Re z=\ensuremath{\tfrac12} (grau 2 \ensuremath{\to} reta, não fractal, como a reta mineral). O lado negro é a bacia do zero: para onde caem os estados que a máquina aniquila.
\prov{relações: usa idempotente\_e\_o\_projetor; relaciona julia; relaciona newton\_acha\_o\_grupo\_discreto; relaciona rastro\_da\_agulha\_e\_fractal}
\prov{proveniência: wiki \textperiodcentered{} parte fria e idempotentes \textperiodcentered{} fato \textperiodcentered{} confiança 1.0 \textperiodcentered{} domínio reta mineral \textperiodcentered{} história 5 versões no jornal}

%CRISTAL {"arestas":[{"alvo":"unidades_algebricas_pu","confianca":1.0,"epistemico":"fato","origem":"wiki","rel":"usa"},{"alvo":"estatuto_derivado_posto_convencao","confianca":1.0,"epistemico":"fato","origem":"wiki","rel":"relaciona"},{"alvo":"desitter_4d_universo","confianca":1.0,"epistemico":"fato","origem":"wiki","rel":"usa"},{"alvo":"de_sitter_lambda_energia_escura","confianca":1.0,"epistemico":"fato","origem":"wiki","rel":"explica"},{"alvo":"lambda_cosmologica_saga","confianca":1.0,"epistemico":"fato","origem":"wiki","rel":"explica"},{"alvo":"equacoes_friedmann","confianca":1.0,"epistemico":"fato","origem":"wiki","rel":"usa"},{"alvo":"cosmologia_quantica_hub","confianca":1.0,"epistemico":"fato","origem":"wiki","rel":"parte_de"}],"confianca":1.0,"contraexemplos":[],"descricao":"A constante cosmológica como exemplo do estatuto. (1) ADIMENSIONAL, derivado: o espaço de álgebras em ℝ⁴ é dS₄ (R=12, R_μν=3g_μν), espaço de Einstein com Λ_adim=R_μν/g_μν=D−1=3, puro, da curvatura. (2) ESCALA, posta: ℓ=c/(H₀√Ω_Λ)=1,66×10²⁶ m (de H₀ e Ω_Λ medidos). (3) RESTAURAÇÃO: Λ=3/ℓ²=1,09×10⁻⁵² m⁻² — o valor observado. (4) ESTATUTO: como Λ=3H₀²Ω_Λ/c², é o Λ PADRÃO — o valor é identidade, não predição; a álgebra deriva o coeficiente 3 (a forma dS₄), a escala vem dos dados. 'A álgebra veste a constante, não a inventa.'","epistemico":"fato","exemplos":[],"id":"lambda_deducao_quatro_passos","memoria":"semantica","meta":{"autor":"Lopes/broca-so","dominio":"unidades","fonte":"hiper/circular/paper_5_unidades.tex §conversão (deducao_lambda.py)"},"origem":"wiki","palavras_chave":[],"sinonimos":[],"tipo":"conceito","titulo":"Dedução de Λ em quatro passos (3 → 1,1×10⁻⁵² m⁻²)"}
\section{Dedução de \ensuremath{\Lambda} em quatro passos (3 \ensuremath{\to} 1,1$\times$10\ensuremath{^{-52}} m\ensuremath{^{-2}})}
A constante cosmológica como exemplo do estatuto. (1) ADIMENSIONAL, derivado: o espaço de álgebras em \ensuremath{\mathbb{R}}\ensuremath{^{4}} é dS\ensuremath{_{4}} (R=12, R\_\ensuremath{\mu}\ensuremath{\nu}=3g\_\ensuremath{\mu}\ensuremath{\nu}), espaço de Einstein com \ensuremath{\Lambda}\_adim=R\_\ensuremath{\mu}\ensuremath{\nu}/g\_\ensuremath{\mu}\ensuremath{\nu}=D\ensuremath{-}1=3, puro, da curvatura. (2) ESCALA, posta: \ensuremath{\ell}=c/(H\ensuremath{_{0}}\ensuremath{\sqrt{\,}}\ensuremath{\Omega}\_\ensuremath{\Lambda})=1,66$\times$10\ensuremath{^{26}} m (de H\ensuremath{_{0}} e \ensuremath{\Omega}\_\ensuremath{\Lambda} medidos). (3) RESTAURAÇÃO: \ensuremath{\Lambda}=3/\ensuremath{\ell}\ensuremath{^{2}}=1,09$\times$10\ensuremath{^{-52}} m\ensuremath{^{-2}} — o valor observado. (4) ESTATUTO: como \ensuremath{\Lambda}=3H\ensuremath{_{0}}\ensuremath{^{2}}\ensuremath{\Omega}\_\ensuremath{\Lambda}/c\ensuremath{^{2}}, é o \ensuremath{\Lambda} PADRÃO — o valor é identidade, não predição; a álgebra deriva o coeficiente 3 (a forma dS\ensuremath{_{4}}), a escala vem dos dados. 'A álgebra veste a constante, não a inventa.'
\prov{relações: usa unidades\_algebricas\_pu; relaciona estatuto\_derivado\_posto\_convencao; usa desitter\_4d\_universo; explica de\_sitter\_lambda\_energia\_escura; explica lambda\_cosmologica\_saga; usa equacoes\_friedmann; parte\_de cosmologia\_quantica\_hub}
\prov{proveniência: wiki \textperiodcentered{} hiper/circular/paper\_5\_unidades.tex §conversão (deducao\_lambda.py) \textperiodcentered{} fato \textperiodcentered{} confiança 1.0 \textperiodcentered{} domínio unidades \textperiodcentered{} história 9 versões no jornal}

%CRISTAL {"arestas":[{"alvo":"o_gato_e_um_mapa_de_anosov","confianca":1.0,"epistemico":"fato","origem":"wiki","rel":"usa"},{"alvo":"o_gap_a_hiperbolicidade_parte_o_sistema_em_dois","confianca":1.0,"epistemico":"fato","origem":"wiki","rel":"eh_um"},{"alvo":"borda_critica_vazio","confianca":1.0,"epistemico":"fato","origem":"wiki","rel":"usa"},{"alvo":"numeros_de_pisot","confianca":1.0,"epistemico":"fato","origem":"wiki","rel":"usa"},{"alvo":"numero_de_salem","confianca":1.0,"epistemico":"fato","origem":"wiki","rel":"usa"}],"confianca":1.0,"contraexemplos":[],"descricao":"A promoção geral: a lei universal da parábola (cristal/borda/caos) É a classificação da dinâmica HIPERBÓLICA de Anosov pelo sinal do Lyapunov (Oseledets). CRISTAL=E^s (estável, λ<0, contração, Pisot); CAOS=E^u (instável, λ>0, expansão). A parábola classifica um ESPECTRO de dinâmicas, não um gato só (um gato hiperbólico é sempre caos e já contém E^s e E^u). A BORDA (λ=0) é o VÉRTICE DEGENERADO — discriminante=0, matriz parabólica, a cúspide/Salem ρ=1 onde a hiperbolicidade se perde —, não uma direção interna do mapa; o E^0 genuíno é do FLUXO geodésico contínuo. É a nossa SÍNTESE (cada peça é fato; a costura β-numeração+Anosov+Selberg numa parábola é a tese).","epistemico":"inferencia","exemplos":[],"id":"lei_da_parabola_e_dinamica_de_anosov","memoria":"semantica","meta":{"dominio":"matematica","fonte":"dinâmica de Anosov / Selberg"},"origem":"wiki","palavras_chave":[],"sinonimos":[],"tipo":"conceito","titulo":"A Lei da Parábola é a Dinâmica de Anosov"}
\section{A Lei da Parábola é a Dinâmica de Anosov}
A promoção geral: a lei universal da parábola (cristal/borda/caos) É a classificação da dinâmica HIPERBÓLICA de Anosov pelo sinal do Lyapunov (Oseledets). CRISTAL=E\^{}s (estável, \ensuremath{\lambda}<0, contração, Pisot); CAOS=E\^{}u (instável, \ensuremath{\lambda}>0, expansão). A parábola classifica um ESPECTRO de dinâmicas, não um gato só (um gato hiperbólico é sempre caos e já contém E\^{}s e E\^{}u). A BORDA (\ensuremath{\lambda}=0) é o VÉRTICE DEGENERADO — discriminante=0, matriz parabólica, a cúspide/Salem \ensuremath{\rho}=1 onde a hiperbolicidade se perde —, não uma direção interna do mapa; o E\^{}0 genuíno é do FLUXO geodésico contínuo. É a nossa SÍNTESE (cada peça é fato; a costura \ensuremath{\beta}-numeração+Anosov+Selberg numa parábola é a tese).
\prov{relações: usa o\_gato\_e\_um\_mapa\_de\_anosov; eh\_um o\_gap\_a\_hiperbolicidade\_parte\_o\_sistema\_em\_dois; usa borda\_critica\_vazio; usa numeros\_de\_pisot; usa numero\_de\_salem}
\prov{proveniência: wiki \textperiodcentered{} dinâmica de Anosov / Selberg \textperiodcentered{} inferencia \textperiodcentered{} confiança 1.0 \textperiodcentered{} domínio matematica \textperiodcentered{} história 24 versões no jornal}

%CRISTAL {"arestas":[{"alvo":"desenho_de_mecanismos","confianca":1.0,"epistemico":"fato","origem":"wiki","rel":"especializa"}],"confianca":1.0,"contraexemplos":[],"descricao":"Mecanismos para vender sob informação privada: primeiro-preço, segundo-preço (Vickrey, onde dizer a verdade é dominante), inglês, holandês. O caso mais aplicado do desenho de mecanismos.","epistemico":"fato","exemplos":[],"id":"leilao_teoria","memoria":"semantica","meta":{"autor":"Vickrey","dominio":"matematica","fonte":"teoria dos jogos"},"origem":"wiki","palavras_chave":[],"sinonimos":[],"tipo":"conceito","titulo":"Leilões"}
\section{Leilões}
Mecanismos para vender sob informação privada: primeiro-preço, segundo-preço (Vickrey, onde dizer a verdade é dominante), inglês, holandês. O caso mais aplicado do desenho de mecanismos.
\prov{relações: especializa desenho\_de\_mecanismos}
\prov{proveniência: wiki \textperiodcentered{} teoria dos jogos \textperiodcentered{} fato \textperiodcentered{} confiança 1.0 \textperiodcentered{} domínio matematica \textperiodcentered{} história 2 versões no jornal}

%CRISTAL {"arestas":[{"alvo":"matematica","confianca":0.95,"epistemico":"fato","origem":"paper","rel":"parte_de"}],"confianca":1.0,"contraexemplos":[],"descricao":"402\n2o) (p+1)−1 2m−1 = (p−1)+1 2m−1 ̸∈Z Sendo assim, temos (−1) \u0016 (p+1)−1 2m−1 \u0017 = (−1) \u0016 p−1 2m−1 \u0017 = (−1)( p−1 2m−1) é suficiente mostrar que \u0000 p−1 2m−1 \u0001 e \u0000(p+1)−1 2m−1 \u0001 têm a mesma paridade. Isto foi feito no","epistemico":"fato","exemplos":[],"id":"lema_11_402_2o_p11_2m1_p11_2m1_z_sendo_assim_te","memoria":"semantica","meta":{"arquivo":"gentil/Novas_Sequencias_Aritmeticas_e_Geometric.pdf","ciencia":"teoria_dos_numeros","dominio":"matematica","nivel":4,"numero":"11","tipo_no":"paper","tipo_teorema":"lema","verificacao":"conceito novo"},"origem":"gentil/Novas_Sequencias_Aritmeticas_e_Geometric.pdf","palavras_chave":[],"sinonimos":[],"tipo":"conceito","titulo":"Lema 11: 402\n2o) (p+1)−1 2m−1 = (p−1)+1 2m−1 ̸∈Z Sendo assim, te…"}
\section{Lema 11: 402
2o) (p+1)\ensuremath{-}1 2m\ensuremath{-}1 = (p\ensuremath{-}1)+1 2m\ensuremath{-}1 \ensuremath{}\ensuremath{\in}Z Sendo assim, te…}
402
2o) (p+1)\ensuremath{-}1 2m\ensuremath{-}1 = (p\ensuremath{-}1)+1 2m\ensuremath{-}1 \ensuremath{}\ensuremath{\in}Z Sendo assim, temos (\ensuremath{-}1)  (p+1)\ensuremath{-}1 2m\ensuremath{-}1  = (\ensuremath{-}1)  p\ensuremath{-}1 2m\ensuremath{-}1  = (\ensuremath{-}1)( p\ensuremath{-}1 2m\ensuremath{-}1) é suficiente mostrar que  p\ensuremath{-}1 2m\ensuremath{-}1  e (p+1)\ensuremath{-}1 2m\ensuremath{-}1  têm a mesma paridade. Isto foi feito no
\prov{relações: parte\_de matematica}
\prov{proveniência: gentil/Novas\_Sequencias\_Aritmeticas\_e\_Geometric.pdf \textperiodcentered{} fato \textperiodcentered{} confiança 1.0 \textperiodcentered{} domínio matematica \textperiodcentered{} história 2 versões no jornal}

%CRISTAL {"arestas":[{"alvo":"matematica","confianca":0.95,"epistemico":"fato","origem":"paper","rel":"parte_de"}],"confianca":1.0,"contraexemplos":[],"descricao":"p 399) surgiu quando tentavamos demonstrar a igualdade entre as matrizes que comparecem na página 396. Posteriormente concebemos uma prova mais direta da referida identidade: Pelo corolário do","epistemico":"fato","exemplos":[],"id":"lema_11_p_399_surgiu_quando_tentavamos_demonstrar_a_igualdade","memoria":"semantica","meta":{"arquivo":"gentil/Novas_Sequencias_Aritmeticas_e_Geometric.pdf","ciencia":"teoria_dos_numeros","dominio":"matematica","nivel":4,"numero":"11","tipo_no":"paper","tipo_teorema":"lema","verificacao":"conceito novo"},"origem":"gentil/Novas_Sequencias_Aritmeticas_e_Geometric.pdf","palavras_chave":[],"sinonimos":[],"tipo":"conceito","titulo":"Lema 11: p 399) surgiu quando tentavamos demonstrar a igualdade …"}
\section{Lema 11: p 399) surgiu quando tentavamos demonstrar a igualdade …}
p 399) surgiu quando tentavamos demonstrar a igualdade entre as matrizes que comparecem na página 396. Posteriormente concebemos uma prova mais direta da referida identidade: Pelo corolário do
\prov{relações: parte\_de matematica}
\prov{proveniência: gentil/Novas\_Sequencias\_Aritmeticas\_e\_Geometric.pdf \textperiodcentered{} fato \textperiodcentered{} confiança 1.0 \textperiodcentered{} domínio matematica \textperiodcentered{} história 2 versões no jornal}

%CRISTAL {"arestas":[{"alvo":"matematica","confianca":0.95,"epistemico":"fato","origem":"paper","rel":"parte_de"}],"confianca":1.0,"contraexemplos":[],"descricao":"p 17), temos: an1 = a11 + S(n−1)(1−1) onde, S(n−1)0 é a soma dos n−1 termos iniciais da P.A. de ordem zero, vale: S(n−1)0 = r + r + · · · + r | z (n−1) termos = (n −1) r Sendo assim, temos: an1 = a11 + (n −1) r (1.2) ´E a fórmula do termo geral de uma P.A.1. (ii) m = 2: Utilizando o","epistemico":"fato","exemplos":[],"id":"lema_1_p_17_temos_an1_a11_sn111_onde_sn10_e_a","memoria":"semantica","meta":{"arquivo":"gentil/Novas_Sequencias_Aritmeticas_e_Geometric.pdf","ciencia":"teoria_dos_numeros","dominio":"matematica","nivel":4,"numero":"1","tipo_no":"paper","tipo_teorema":"lema","verificacao":"conceito novo"},"origem":"gentil/Novas_Sequencias_Aritmeticas_e_Geometric.pdf","palavras_chave":[],"sinonimos":[],"tipo":"conceito","titulo":"Lema 1: p 17), temos: an1 = a11 + S(n−1)(1−1) onde, S(n−1)0 é a…"}
\section{Lema 1: p 17), temos: an1 = a11 + S(n\ensuremath{-}1)(1\ensuremath{-}1) onde, S(n\ensuremath{-}1)0 é a…}
p 17), temos: an1 = a11 + S(n\ensuremath{-}1)(1\ensuremath{-}1) onde, S(n\ensuremath{-}1)0 é a soma dos n\ensuremath{-}1 termos iniciais da P.A. de ordem zero, vale: S(n\ensuremath{-}1)0 = r + r + \textperiodcentered  \textperiodcentered  \textperiodcentered  + r | z (n\ensuremath{-}1) termos = (n \ensuremath{-}1) r Sendo assim, temos: an1 = a11 + (n \ensuremath{-}1) r (1.2) 'E a fórmula do termo geral de uma P.A.1. (ii) m = 2: Utilizando o
\prov{relações: parte\_de matematica}
\prov{proveniência: gentil/Novas\_Sequencias\_Aritmeticas\_e\_Geometric.pdf \textperiodcentered{} fato \textperiodcentered{} confiança 1.0 \textperiodcentered{} domínio matematica \textperiodcentered{} história 3 versões no jornal}

%CRISTAL {"arestas":[{"alvo":"matematica","confianca":0.95,"epistemico":"fato","origem":"paper","rel":"parte_de"}],"confianca":1.0,"contraexemplos":[],"descricao":"p 17), temos: an2 = a12 + S(n−1)(2−1) (1.3) onde, S(n−1)1 é a soma dos n −1 termos iniciais da P.A.1 . Vamos deduzir esta fórmula utilizando a equação (1.2), assim: a11 = a11 a21 = a11 + r a31 = a11 + 2r · · · · · · · · · · · · · · · · · · · · · an1 = a11 + (n −1) r Somando estas n igualdades, resulta: a11 + a21 + a31 + · · · + an1 = (a11 + a11 + a11 + · · · + a11) + \u00001 + 2 + · · · + (n −1) \u0001 r O seguinte resultado já é conhecido, 1 + 2 + · · · + n = n(n + 1) 2 ⇒1 + 2 + · · · + n −1 = (n −1)n 2","epistemico":"fato","exemplos":[],"id":"lema_1_p_17_temos_an2_a12_sn121_13_onde_sn1","memoria":"semantica","meta":{"arquivo":"gentil/Novas_Sequencias_Aritmeticas_e_Geometric.pdf","ciencia":"teoria_dos_numeros","dominio":"matematica","nivel":4,"numero":"1","tipo_no":"paper","tipo_teorema":"lema","verificacao":"conceito novo"},"origem":"gentil/Novas_Sequencias_Aritmeticas_e_Geometric.pdf","palavras_chave":[],"sinonimos":[],"tipo":"conceito","titulo":"Lema 1: p 17), temos: an2 = a12 + S(n−1)(2−1) (1.3) onde, S(n−1…"}
\section{Lema 1: p 17), temos: an2 = a12 + S(n\ensuremath{-}1)(2\ensuremath{-}1) (1.3) onde, S(n\ensuremath{-}1…}
p 17), temos: an2 = a12 + S(n\ensuremath{-}1)(2\ensuremath{-}1) (1.3) onde, S(n\ensuremath{-}1)1 é a soma dos n \ensuremath{-}1 termos iniciais da P.A.1 . Vamos deduzir esta fórmula utilizando a equação (1.2), assim: a11 = a11 a21 = a11 + r a31 = a11 + 2r \textperiodcentered  \textperiodcentered  \textperiodcentered  \textperiodcentered  \textperiodcentered  \textperiodcentered  \textperiodcentered  \textperiodcentered  \textperiodcentered  \textperiodcentered  \textperiodcentered  \textperiodcentered  \textperiodcentered  \textperiodcentered  \textperiodcentered  \textperiodcentered  \textperiodcentered  \textperiodcentered  \textperiodcentered  \textperiodcentered  \textperiodcentered  an1 = a11 + (n \ensuremath{-}1) r Somando estas n igualdades, resulta: a11 + a21 + a31 + \textperiodcentered  \textperiodcentered  \textperiodcentered  + an1 = (a11 + a11 + a11 + \textperiodcentered  \textperiodcentered  \textperiodcentered  + a11) + 1 + 2 + \textperiodcentered  \textperiodcentered  \textperiodcentered  + (n \ensuremath{-}1)  r O seguinte resultado já é conhecido, 1 + 2 + \textperiodcentered  \textperiodcentered  \textperiodcentered  + n = n(n + 1) 2 \ensuremath{\Rightarrow}1 + 2 + \textperiodcentered  \textperiodcentered  \textperiodcentered  + n \ensuremath{-}1 = (n \ensuremath{-}1)n 2
\prov{relações: parte\_de matematica}
\prov{proveniência: gentil/Novas\_Sequencias\_Aritmeticas\_e\_Geometric.pdf \textperiodcentered{} fato \textperiodcentered{} confiança 1.0 \textperiodcentered{} domínio matematica \textperiodcentered{} história 2 versões no jornal}

%CRISTAL {"arestas":[{"alvo":"matematica","confianca":0.95,"epistemico":"fato","origem":"paper","rel":"parte_de"},{"alvo":"word","confianca":0.5,"epistemico":"fato","origem":"wiki","rel":"relaciona"}],"confianca":1.0,"contraexemplos":[],"descricao":"p 104), resulta: ∆(p+1) ∇m f(n) = ∇m−p−1 f(n) Isto é, ∆(p+1) ∇m f(n) = ∇m−(p+1) f(n) ■","epistemico":"fato","exemplos":[],"id":"lema_2_p_104_resulta_p1_m_fn_mp1_fn_isto_e","memoria":"semantica","meta":{"arquivo":"gentil/Novas_Sequencias_Aritmeticas_e_Geometric.pdf","ciencia":"teoria_dos_numeros","dominio":"matematica","nivel":4,"numero":"2","tipo_no":"paper","tipo_teorema":"lema","verificacao":"conceito novo"},"origem":"gentil/Novas_Sequencias_Aritmeticas_e_Geometric.pdf","palavras_chave":[],"sinonimos":[],"tipo":"conceito","titulo":"Lema 2: p 104), resulta: ∆(p+1) ∇m f(n) = ∇m−p−1 f(n) Isto é, ∆…"}
\section{Lema 2: p 104), resulta: \ensuremath{\Delta}(p+1) \ensuremath{\nabla}m f(n) = \ensuremath{\nabla}m\ensuremath{-}p\ensuremath{-}1 f(n) Isto é, \ensuremath{\Delta}…}
p 104), resulta: \ensuremath{\Delta}(p+1) \ensuremath{\nabla}m f(n) = \ensuremath{\nabla}m\ensuremath{-}p\ensuremath{-}1 f(n) Isto é, \ensuremath{\Delta}(p+1) \ensuremath{\nabla}m f(n) = \ensuremath{\nabla}m\ensuremath{-}(p+1) f(n) \rule{1ex}{1ex}
\prov{relações: parte\_de matematica; relaciona word}
\prov{proveniência: gentil/Novas\_Sequencias\_Aritmeticas\_e\_Geometric.pdf \textperiodcentered{} fato \textperiodcentered{} confiança 1.0 \textperiodcentered{} domínio matematica \textperiodcentered{} história 3 versões no jornal}

%CRISTAL {"arestas":[{"alvo":"matematica","confianca":0.95,"epistemico":"fato","origem":"paper","rel":"parte_de"}],"confianca":1.0,"contraexemplos":[],"descricao":"Seja m um natural arbitrariamente fixado e j um natural tal que 0 ≤ j ≤m Nestas condições a seguinte identidade se verifica: ∆ \u0000∇m−p f(n) \u0001 = ∇m−p−1 f(n) Prova: Temos que ∆f(n) = f(n + 1) −f(n), logo ∆ \u0000∇m−p f(n) \u0001 = ∇m−p f(n + 1) −∇m−p f(n) (2.13) Pela","epistemico":"fato","exemplos":[],"id":"lema_2_seja_m_um_natural_arbitrariamente_fixado_e_j_um_natural","memoria":"semantica","meta":{"arquivo":"gentil/Novas_Sequencias_Aritmeticas_e_Geometric.pdf","ciencia":"teoria_dos_numeros","dominio":"matematica","nivel":4,"numero":"2","tipo_no":"paper","tipo_teorema":"lema","verificacao":"conceito novo"},"origem":"gentil/Novas_Sequencias_Aritmeticas_e_Geometric.pdf","palavras_chave":[],"sinonimos":[],"tipo":"conceito","titulo":"Lema 2: Seja m um natural arbitrariamente fixado e j um natural…"}
\section{Lema 2: Seja m um natural arbitrariamente fixado e j um natural…}
Seja m um natural arbitrariamente fixado e j um natural tal que 0 \ensuremath{\le} j \ensuremath{\le}m Nestas condições a seguinte identidade se verifica: \ensuremath{\Delta} \ensuremath{\nabla}m\ensuremath{-}p f(n)  = \ensuremath{\nabla}m\ensuremath{-}p\ensuremath{-}1 f(n) Prova: Temos que \ensuremath{\Delta}f(n) = f(n + 1) \ensuremath{-}f(n), logo \ensuremath{\Delta} \ensuremath{\nabla}m\ensuremath{-}p f(n)  = \ensuremath{\nabla}m\ensuremath{-}p f(n + 1) \ensuremath{-}\ensuremath{\nabla}m\ensuremath{-}p f(n) (2.13) Pela
\prov{relações: parte\_de matematica}
\prov{proveniência: gentil/Novas\_Sequencias\_Aritmeticas\_e\_Geometric.pdf \textperiodcentered{} fato \textperiodcentered{} confiança 1.0 \textperiodcentered{} domínio matematica \textperiodcentered{} história 2 versões no jornal}

%CRISTAL {"arestas":[{"alvo":"matematica","confianca":0.95,"epistemico":"fato","origem":"paper","rel":"parte_de"}],"confianca":1.0,"contraexemplos":[],"descricao":"´E válida a seguinte identidade ∆ m f(n) = ∆ m−1 ∆ f(n) Prova: Princ´ıpio da Dualidade. (p. 88) ■ O seguinte","epistemico":"fato","exemplos":[],"id":"lema_4_e_valida_a_seguinte_identidade_m_fn_m1_fn","memoria":"semantica","meta":{"arquivo":"gentil/Novas_Sequencias_Aritmeticas_e_Geometric.pdf","ciencia":"teoria_dos_numeros","dominio":"matematica","nivel":4,"numero":"4","tipo_no":"paper","tipo_teorema":"lema","verificacao":"conceito novo"},"origem":"gentil/Novas_Sequencias_Aritmeticas_e_Geometric.pdf","palavras_chave":[],"sinonimos":[],"tipo":"conceito","titulo":"Lema 4: ´E válida a seguinte identidade ∆ m f(n) = ∆ m−1 ∆ f(n)…"}
\section{Lema 4: 'E válida a seguinte identidade \ensuremath{\Delta} m f(n) = \ensuremath{\Delta} m\ensuremath{-}1 \ensuremath{\Delta} f(n)…}
'E válida a seguinte identidade \ensuremath{\Delta} m f(n) = \ensuremath{\Delta} m\ensuremath{-}1 \ensuremath{\Delta} f(n) Prova: Princ'\ensuremath{\imath}pio da Dualidade. (p. 88) \rule{1ex}{1ex} O seguinte
\prov{relações: parte\_de matematica}
\prov{proveniência: gentil/Novas\_Sequencias\_Aritmeticas\_e\_Geometric.pdf \textperiodcentered{} fato \textperiodcentered{} confiança 1.0 \textperiodcentered{} domínio matematica \textperiodcentered{} história 2 versões no jornal}

%CRISTAL {"arestas":[{"alvo":"matematica","confianca":0.95,"epistemico":"fato","origem":"paper","rel":"parte_de"}],"confianca":1.0,"contraexemplos":[],"descricao":"Seja q um natural arbitrariamente fixado e 0 ≤m ≤q Nestas condições é válida a seguinte identidade: ∆m ( −1 )( n q ) = ( −1 )( n q−m) Prova: Indução sobre m. Para m = 0, temos: ∆0 ( −1 )( n q ) = ( −1 )( n q ) = ( −1 )( n q−0) Admitamos a validade da proposição para m = p (0 ≤p < q): (H.I.) ∆p ( −1 )( n q ) = ( −1 )( n q−p) e provemos que vale para m = p + 1: (T.I.) ∆p+1 ( −1 )( n q ) = ( −1 )( n q−(p+1)) Inicialmente vamos mostrar que ∆( −1 )( n q ) = ( −1 )( n q−1) De fato, ∆( −1 )( n q ) = (","epistemico":"fato","exemplos":[],"id":"lema_5_seja_q_um_natural_arbitrariamente_fixado_e_0_m_q_nest","memoria":"semantica","meta":{"arquivo":"gentil/Novas_Sequencias_Aritmeticas_e_Geometric.pdf","ciencia":"teoria_dos_numeros","dominio":"matematica","nivel":4,"numero":"5","tipo_no":"paper","tipo_teorema":"lema","verificacao":"conceito novo"},"origem":"gentil/Novas_Sequencias_Aritmeticas_e_Geometric.pdf","palavras_chave":[],"sinonimos":[],"tipo":"conceito","titulo":"Lema 5: Seja q um natural arbitrariamente fixado e 0 ≤m ≤q Nest…"}
\section{Lema 5: Seja q um natural arbitrariamente fixado e 0 \ensuremath{\le}m \ensuremath{\le}q Nest…}
Seja q um natural arbitrariamente fixado e 0 \ensuremath{\le}m \ensuremath{\le}q Nestas condições é válida a seguinte identidade: \ensuremath{\Delta}m ( \ensuremath{-}1 )( n q ) = ( \ensuremath{-}1 )( n q\ensuremath{-}m) Prova: Indução sobre m. Para m = 0, temos: \ensuremath{\Delta}0 ( \ensuremath{-}1 )( n q ) = ( \ensuremath{-}1 )( n q ) = ( \ensuremath{-}1 )( n q\ensuremath{-}0) Admitamos a validade da proposição para m = p (0 \ensuremath{\le}p < q): (H.I.) \ensuremath{\Delta}p ( \ensuremath{-}1 )( n q ) = ( \ensuremath{-}1 )( n q\ensuremath{-}p) e provemos que vale para m = p + 1: (T.I.) \ensuremath{\Delta}p+1 ( \ensuremath{-}1 )( n q ) = ( \ensuremath{-}1 )( n q\ensuremath{-}(p+1)) Inicialmente vamos mostrar que \ensuremath{\Delta}( \ensuremath{-}1 )( n q ) = ( \ensuremath{-}1 )( n q\ensuremath{-}1) De fato, \ensuremath{\Delta}( \ensuremath{-}1 )( n q ) = (
\prov{relações: parte\_de matematica}
\prov{proveniência: gentil/Novas\_Sequencias\_Aritmeticas\_e\_Geometric.pdf \textperiodcentered{} fato \textperiodcentered{} confiança 1.0 \textperiodcentered{} domínio matematica \textperiodcentered{} história 2 versões no jornal}

%CRISTAL {"arestas":[{"alvo":"matematica","confianca":0.95,"epistemico":"fato","origem":"paper","rel":"parte_de"}],"confianca":1.0,"contraexemplos":[],"descricao":"p 183), temos: j p −k 2 j−1 k =        j p−k−1 2j−1 k + 1, se p−k 2j−1 ∈Z ; j p−k−1 2j−1 k , se p−k 2j−1 ̸∈Z . Observe que se (p −k)/2 j−1 ̸∈Z, o resultado é imediato. Suponhamos (p −k)/2 j−1 ∈Z. Somando∗ 2 j−1 X k = 0 (−1)k \u0012 2 j−1 k \u0013 = 0 ∗Este resultado sai do binômio de Newton. 185\nà hipótese de indução, resulta: 2 j−1 X k = 0 (−1)k \u0012 2 j−1 k \u0013 \u001a j p −k −1 2 j−1 k + 1 \u001b + 2 j−1 X k = 0 (−1)k \u0012 2 j−1 k \u0013 ≡1 (mod 2) Esta última igualdade pode ser reescrita como: 2 j−1 X k = 0 (−1)k \u0012 2","epistemico":"fato","exemplos":[],"id":"lema_7_p_183_temos_j_p_k_2_j1_k_j_pk1_2","memoria":"semantica","meta":{"arquivo":"gentil/Novas_Sequencias_Aritmeticas_e_Geometric.pdf","ciencia":"teoria_dos_numeros","dominio":"matematica","nivel":4,"numero":"7","tipo_no":"paper","tipo_teorema":"lema","verificacao":"conceito novo"},"origem":"gentil/Novas_Sequencias_Aritmeticas_e_Geometric.pdf","palavras_chave":[],"sinonimos":[],"tipo":"conceito","titulo":"Lema 7: p 183), temos: j p −k 2 j−1 k =        j p−k−1 2…"}
\section{Lema 7: p 183), temos: j p \ensuremath{-}k 2 j\ensuremath{-}1 k = \{       j p\ensuremath{-}k\ensuremath{-}1 2…}
p 183), temos: j p \ensuremath{-}k 2 j\ensuremath{-}1 k = \{       j p\ensuremath{-}k\ensuremath{-}1 2j\ensuremath{-}1 k + 1, se p\ensuremath{-}k 2j\ensuremath{-}1 \ensuremath{\in}Z ; j p\ensuremath{-}k\ensuremath{-}1 2j\ensuremath{-}1 k , se p\ensuremath{-}k 2j\ensuremath{-}1 \ensuremath{}\ensuremath{\in}Z . Observe que se (p \ensuremath{-}k)/2 j\ensuremath{-}1 \ensuremath{}\ensuremath{\in}Z, o resultado é imediato. Suponhamos (p \ensuremath{-}k)/2 j\ensuremath{-}1 \ensuremath{\in}Z. Somando\ensuremath{*} 2 j\ensuremath{-}1 X k = 0 (\ensuremath{-}1)k  2 j\ensuremath{-}1 k  = 0 \ensuremath{*}Este resultado sai do binômio de Newton. 185
à hipótese de indução, resulta: 2 j\ensuremath{-}1 X k = 0 (\ensuremath{-}1)k  2 j\ensuremath{-}1 k   j p \ensuremath{-}k \ensuremath{-}1 2 j\ensuremath{-}1 k + 1  + 2 j\ensuremath{-}1 X k = 0 (\ensuremath{-}1)k  2 j\ensuremath{-}1 k  \ensuremath{\equiv}1 (mod 2) Esta última igualdade pode ser reescrita como: 2 j\ensuremath{-}1 X k = 0 (\ensuremath{-}1)k  2
\prov{relações: parte\_de matematica}
\prov{proveniência: gentil/Novas\_Sequencias\_Aritmeticas\_e\_Geometric.pdf \textperiodcentered{} fato \textperiodcentered{} confiança 1.0 \textperiodcentered{} domínio matematica \textperiodcentered{} história 2 versões no jornal}

%CRISTAL {"arestas":[{"alvo":"matriz_cantor","confianca":1.0,"epistemico":"fato","origem":"paper","rel":"usa"},{"alvo":"word","confianca":0.5,"epistemico":"fato","origem":"wiki","rel":"relaciona"},{"alvo":"virtualizacao","confianca":0.5,"epistemico":"fato","origem":"wiki","rel":"relaciona"},{"alvo":"vida-morte","confianca":0.5,"epistemico":"fato","origem":"wiki","rel":"relaciona"}],"confianca":1.0,"contraexemplos":[],"descricao":"A árvore livre (λ=2), refinada por podas sucessivas, gera densamente todas as carnes finitas (autovalores de Perron densos em (1,2)) — Lind.","epistemico":"fato","exemplos":[],"id":"lema_geracao","memoria":"semantica","meta":{"autor":"Gentil Lopes","descoberta":false,"dominio":"matematica","fonte":"paper_66_lema_geracao.pdf, Teo.3.1"},"origem":"paper","palavras_chave":[],"sinonimos":[],"tipo":"conceito","titulo":"Lema de Geração"}
\section{Lema de Geração}
A árvore livre (\ensuremath{\lambda}=2), refinada por podas sucessivas, gera densamente todas as carnes finitas (autovalores de Perron densos em (1,2)) — Lind.
\prov{relações: usa matriz\_cantor; relaciona word; relaciona virtualizacao; relaciona vida-morte}
\prov{proveniência: paper \textperiodcentered{} paper\_66\_lema\_geracao.pdf, Teo.3.1 \textperiodcentered{} fato \textperiodcentered{} confiança 1.0 \textperiodcentered{} domínio matematica \textperiodcentered{} história 5 versões no jornal}

%CRISTAL {"arestas":[{"alvo":"inversao_geometrica_s","confianca":1.0,"epistemico":"fato","origem":"paper","rel":"relaciona"},{"alvo":"corpo_gauss_modp","confianca":1.0,"epistemico":"fato","origem":"paper","rel":"relaciona"}],"confianca":1.0,"contraexemplos":[],"descricao":"S fixa τ=i; o reticulado ℤ[i] tem multiplicação complexa (j-invariante 1728) e realiza a curva elíptica y²=x³−x, a lemniscata de Bernoulli.","epistemico":"fato","exemplos":[],"id":"lemniscata_ponto_i","memoria":"semantica","meta":{"autor":"Gentil Lopes","descoberta":false,"dominio":"matematica","fonte":"paper_27_ponte_modular.pdf"},"origem":"paper","palavras_chave":[],"sinonimos":[],"tipo":"conceito","titulo":"Lemniscata e o Ponto i"}
\section{Lemniscata e o Ponto i}
S fixa \ensuremath{\tau}=i; o reticulado \ensuremath{\mathbb{Z}}[i] tem multiplicação complexa (j-invariante 1728) e realiza a curva elíptica y\ensuremath{^{2}}=x\ensuremath{^{3}}\ensuremath{-}x, a lemniscata de Bernoulli.
\prov{relações: relaciona inversao\_geometrica\_s; relaciona corpo\_gauss\_modp}
\prov{proveniência: paper \textperiodcentered{} paper\_27\_ponte\_modular.pdf \textperiodcentered{} fato \textperiodcentered{} confiança 1.0 \textperiodcentered{} domínio matematica \textperiodcentered{} história 3 versões no jornal}

%CRISTAL {"arestas":[{"alvo":"plano_continuo_transicoes","confianca":1.0,"epistemico":"fato","origem":"wiki","rel":"parte_de"},{"alvo":"mapear_nao_fabricar_continuo","confianca":1.0,"epistemico":"fato","origem":"wiki","rel":"usa"},{"alvo":"medida_continuidade_transicao","confianca":1.0,"epistemico":"fato","origem":"wiki","rel":"usa"},{"alvo":"bai_orbita_seleciona_colapso","confianca":1.0,"epistemico":"fato","origem":"wiki","rel":"relaciona"},{"alvo":"tiffany","confianca":1.0,"epistemico":"fato","origem":"wiki","rel":"parte_de"}],"confianca":1.0,"contraexemplos":[],"descricao":"Primeira leva do preenchimento: 8 agentes (workflow transicoes-L1) escreveram a PRIMEIRA CAMADA de transição de prosa para 80 arestas da espinha (subgrafo SymbolicBeta/BAI + hubs), lendo wave1.json por faixa e gravando JSON parcial (sem corrida). 82 transições no tecido 'transicoes' (cobertura L1 ≈1,5% das 5493 arestas); prosa real que EXPRIME a relação, em português, seguindo a filosofia (não fabricar fato — exprimir o vínculo que já existe). ACHADO HONESTO: no colapso PLANO as transições (texto longo → impressão Koch espalhada, ov≈112) PERDEM para frases curtas do tatoeba (ov≈250); a energia de órbita (transicoes=5) não fecha um gap de ~140 de overlap — é a MESMA raiz do 'grafo magro' (o overlap cru premia texto curto). O tecido é consultável direto (tecido=transicoes pousa certo), mas p/ surgir no colapso plano falta OVERLAP NORMALIZADO POR COMPRIMENTO — a próxima leva. [D] roda; grafo íntegro (62/62).","epistemico":"fato","exemplos":[],"id":"leva1_transicoes_l1","memoria":"semantica","meta":{"autor":"Lopes/broca-so","dominio":"continuo","fonte":"workflow transicoes-L1 (8 agentes); conversa/dados/transicoes.tsv; commit 974ecc9"},"origem":"wiki","palavras_chave":[],"sinonimos":[],"tipo":"conceito","titulo":"Leva 1: 80 transições L1 preenchidas por agentes (a espinha) + o achado do overlap"}
\section{Leva 1: 80 transições L1 preenchidas por agentes (a espinha) + o achado do overlap}
Primeira leva do preenchimento: 8 agentes (workflow transicoes-L1) escreveram a PRIMEIRA CAMADA de transição de prosa para 80 arestas da espinha (subgrafo SymbolicBeta/BAI + hubs), lendo wave1.json por faixa e gravando JSON parcial (sem corrida). 82 transições no tecido 'transicoes' (cobertura L1 \ensuremath{\approx}1,5\% das 5493 arestas); prosa real que EXPRIME a relação, em português, seguindo a filosofia (não fabricar fato — exprimir o vínculo que já existe). ACHADO HONESTO: no colapso PLANO as transições (texto longo \ensuremath{\to} impressão Koch espalhada, ov\ensuremath{\approx}112) PERDEM para frases curtas do tatoeba (ov\ensuremath{\approx}250); a energia de órbita (transicoes=5) não fecha um gap de \~{}140 de overlap — é a MESMA raiz do 'grafo magro' (o overlap cru premia texto curto). O tecido é consultável direto (tecido=transicoes pousa certo), mas p/ surgir no colapso plano falta OVERLAP NORMALIZADO POR COMPRIMENTO — a próxima leva. [D] roda; grafo íntegro (62/62).
\prov{relações: parte\_de plano\_continuo\_transicoes; usa mapear\_nao\_fabricar\_continuo; usa medida\_continuidade\_transicao; relaciona bai\_orbita\_seleciona\_colapso; parte\_de tiffany}
\prov{proveniência: wiki \textperiodcentered{} workflow transicoes-L1 (8 agentes); conversa/dados/transicoes.tsv; commit 974ecc9 \textperiodcentered{} fato \textperiodcentered{} confiança 1.0 \textperiodcentered{} domínio continuo \textperiodcentered{} história 89 versões no jornal}

%CRISTAL {"arestas":[{"alvo":"espaco_difusivo","confianca":1.0,"epistemico":"fato","origem":"wiki","rel":"usa"},{"alvo":"parabola_reparte_difusao","confianca":1.0,"epistemico":"fato","origem":"wiki","rel":"usa"}],"confianca":1.0,"contraexemplos":[],"descricao":"Monta o EspacoDifusivo a partir do DNA do host (platform.node()|cpu_count), rodando o núcleo de primitivas sobre esse DNA e lendo c2/calor da auditoria da parábola; olho_threads=1 (o olho só escuta, não bloqueia o tecido). A raiz de stratum_difusao_cpu.","epistemico":"fato","exemplos":[],"id":"libertar_espaco","memoria":"semantica","meta":{"autor":"broca-so","dominio":"fractal","fonte":"neuronio/difusao_recursos.py:libertar_espaco"},"origem":"wiki","palavras_chave":[],"sinonimos":[],"tipo":"conceito","titulo":"libertar_espaco() — o Motor da Difusão"}
\section{libertar\_espaco() — o Motor da Difusão}
Monta o EspacoDifusivo a partir do DNA do host (platform.node()|cpu\_count), rodando o núcleo de primitivas sobre esse DNA e lendo c2/calor da auditoria da parábola; olho\_threads=1 (o olho só escuta, não bloqueia o tecido). A raiz de stratum\_difusao\_cpu.
\prov{relações: usa espaco\_difusivo; usa parabola\_reparte\_difusao}
\prov{proveniência: wiki \textperiodcentered{} neuronio/difusao\_recursos.py:libertar\_espaco \textperiodcentered{} fato \textperiodcentered{} confiança 1.0 \textperiodcentered{} domínio fractal \textperiodcentered{} história 3 versões no jornal}

%CRISTAL {"arestas":[{"alvo":"torre_hurwitz","confianca":1.0,"epistemico":"fato","origem":"paper","rel":"usa"},{"alvo":"word","confianca":0.5,"epistemico":"fato","origem":"wiki","rel":"relaciona"},{"alvo":"vontade_de_poder","confianca":0.5,"epistemico":"fato","origem":"wiki","rel":"relaciona"}],"confianca":1.0,"contraexemplos":[],"descricao":"Por continuidade de Kato, Δ(s)≥Δ(1)−C|1−s|: quando s→1 a teoria hipercomplexa recupera Yang-Mills clássico SU(2) com mass gap de Wilson na rede.","epistemico":"fato","exemplos":[],"id":"limite_hurwitz_s1","memoria":"semantica","meta":{"autor":"Gentil Lopes","descoberta":false,"dominio":"matematica","fonte":"paper_AE_yang_mills.pdf, Teo.21"},"origem":"paper","palavras_chave":[],"sinonimos":[],"tipo":"conceito","titulo":"Limite de Hurwitz (s→1)"}
\section{Limite de Hurwitz (s\ensuremath{\to}1)}
Por continuidade de Kato, \ensuremath{\Delta}(s)\ensuremath{\ge}\ensuremath{\Delta}(1)\ensuremath{-}C|1\ensuremath{-}s|: quando s\ensuremath{\to}1 a teoria hipercomplexa recupera Yang-Mills clássico SU(2) com mass gap de Wilson na rede.
\prov{relações: usa torre\_hurwitz; relaciona word; relaciona vontade\_de\_poder}
\prov{proveniência: paper \textperiodcentered{} paper\_AE\_yang\_mills.pdf, Teo.21 \textperiodcentered{} fato \textperiodcentered{} confiança 1.0 \textperiodcentered{} domínio matematica \textperiodcentered{} história 4 versões no jornal}

%CRISTAL {"arestas":[{"alvo":"floco_phi","confianca":1.0,"epistemico":"fato","origem":"wiki","rel":"usa"},{"alvo":"reconstrucao_condicionada","confianca":1.0,"epistemico":"fato","origem":"wiki","rel":"usa"}],"confianca":1.0,"contraexemplos":[],"descricao":"O paper declara: (1) Rec sem célula branca permanece aberto (A=0); (2) Φ mod p é máquina finita — o contínuo é limite; (3) Ψ (comprimento de arco) nunca fecha exato (AGM); (4) o inteiro Pell exato P_n≥2378 mora em silver, não em Φ; (5) Φ não substitui 𝓛, complementa. E a área-finita/perímetro-infinito da curva de Koch é metáfora, não afirmada no paper.","epistemico":"inferencia","exemplos":[],"id":"limites_honestos_floco","memoria":"semantica","meta":{"autor":"broca-so","dominio":"fractal","fonte":"papers/floco_de_neve.tex §VIII"},"origem":"wiki","palavras_chave":[],"sinonimos":[],"tipo":"conceito","titulo":"Os Limites Honestos do Floco"}
\section{Os Limites Honestos do Floco}
O paper declara: (1) Rec sem célula branca permanece aberto (A=0); (2) \ensuremath{\Phi} mod p é máquina finita — o contínuo é limite; (3) \ensuremath{\Psi} (comprimento de arco) nunca fecha exato (AGM); (4) o inteiro Pell exato P\_n\ensuremath{\ge}2378 mora em silver, não em \ensuremath{\Phi}; (5) \ensuremath{\Phi} não substitui \ensuremath{\mathcal{L}}, complementa. E a área-finita/perímetro-infinito da curva de Koch é metáfora, não afirmada no paper.
\prov{relações: usa floco\_phi; usa reconstrucao\_condicionada}
\prov{proveniência: wiki \textperiodcentered{} papers/floco\_de\_neve.tex §VIII \textperiodcentered{} inferencia \textperiodcentered{} confiança 1.0 \textperiodcentered{} domínio fractal \textperiodcentered{} história 3 versões no jornal}

%CRISTAL {"arestas":[{"alvo":"htlc","confianca":1.0,"epistemico":"fato","origem":"paper","rel":"usa"}],"confianca":1.0,"contraexemplos":[],"descricao":"Sob liveness (cada parte honesta confirma na margem), em todo caminho cada parte honesta completa seu resgate OU reembolso — nunca fica sem nada. Atomicidade condicional à liveness.","epistemico":"fato","exemplos":[],"id":"liveness_atomicidade","memoria":"semantica","meta":{"autor":"Gentil Lopes","descoberta":false,"dominio":"matematica","fonte":"protocolo_swap.pdf, Prop.1"},"origem":"paper","palavras_chave":[],"sinonimos":[],"tipo":"conceito","titulo":"Invariante de Não-Perda"}
\section{Invariante de Não-Perda}
Sob liveness (cada parte honesta confirma na margem), em todo caminho cada parte honesta completa seu resgate OU reembolso — nunca fica sem nada. Atomicidade condicional à liveness.
\prov{relações: usa htlc}
\prov{proveniência: paper \textperiodcentered{} protocolo\_swap.pdf, Prop.1 \textperiodcentered{} fato \textperiodcentered{} confiança 1.0 \textperiodcentered{} domínio matematica \textperiodcentered{} história 2 versões no jornal}

%CRISTAL {"arestas":[{"alvo":"turing_rogozhin_ns","confianca":1.0,"epistemico":"fato","origem":"paper","rel":"explica"}],"confianca":1.0,"contraexemplos":[],"descricao":"Conjectura: Λ>0 (expoente de Lyapunov ≈1,5, robusto) é condição necessária para Turing universal em NS (Bennett, Moore) — caos genuíno é pré-requisito de computação não-trivial.","epistemico":"hipotese","exemplos":[],"id":"lyapunov_computacao","memoria":"semantica","meta":{"autor":"Gentil Lopes","descoberta":false,"dominio":"matematica","fonte":"paper_BJ.pdf"},"origem":"paper","palavras_chave":[],"sinonimos":[],"tipo":"conceito","titulo":"Lyapunov como Termômetro Computacional"}
\section{Lyapunov como Termômetro Computacional}
Conjectura: \ensuremath{\Lambda}>0 (expoente de Lyapunov \ensuremath{\approx}1,5, robusto) é condição necessária para Turing universal em NS (Bennett, Moore) — caos genuíno é pré-requisito de computação não-trivial.
\prov{relações: explica turing\_rogozhin\_ns}
\prov{proveniência: paper \textperiodcentered{} paper\_BJ.pdf \textperiodcentered{} hipotese \textperiodcentered{} confiança 1.0 \textperiodcentered{} domínio matematica \textperiodcentered{} história 2 versões no jornal}

%CRISTAL {"arestas":[{"alvo":"mathieu_m24","confianca":1.0,"epistemico":"fato","origem":"paper","rel":"usa"},{"alvo":"rede_leech","confianca":1.0,"epistemico":"fato","origem":"paper","rel":"usa"}],"confianca":1.0,"contraexemplos":[],"descricao":"Teorema de impossibilidade (Lopes corrige um erro próprio): M₂₄⊄G₂ — os invariantes de M₂₄ reduzem ao peso de Hamming; M₂₄ exige ℝ²⁴ (rede de Leech), não ℝ⁷. PSL(2,7) é o limite em ℝ⁷.","epistemico":"fato","exemplos":[],"id":"m24_nao_em_r7","memoria":"semantica","meta":{"autor":"Gentil Lopes","descoberta":true,"dominio":"matematica","fonte":"paper_BD.pdf, Teo.3"},"origem":"paper","palavras_chave":[],"sinonimos":[],"tipo":"conceito","titulo":"M₂₄ Não Cabe em ℝ⁷"}
\section{M\ensuremath{_{24}} Não Cabe em \ensuremath{\mathbb{R}}\ensuremath{^{7}}}
Teorema de impossibilidade (Lopes corrige um erro próprio): M\ensuremath{_{24}}\ensuremath{\not\subset}G\ensuremath{_{2}} — os invariantes de M\ensuremath{_{24}} reduzem ao peso de Hamming; M\ensuremath{_{24}} exige \ensuremath{\mathbb{R}}\ensuremath{^{24}} (rede de Leech), não \ensuremath{\mathbb{R}}\ensuremath{^{7}}. PSL(2,7) é o limite em \ensuremath{\mathbb{R}}\ensuremath{^{7}}.
\prov{relações: usa mathieu\_m24; usa rede\_leech}
\prov{proveniência: paper \textperiodcentered{} paper\_BD.pdf, Teo.3 \textperiodcentered{} fato \textperiodcentered{} confiança 1.0 \textperiodcentered{} domínio matematica \textperiodcentered{} história 4 versões no jornal}

%CRISTAL {"arestas":[],"confianca":1.0,"contraexemplos":[],"descricao":"Conjectura: computar r₁(Γ₀(7)) a precisão 2⁻ᵏ exige tempo Ω(2kc), c∈[3,4] — sub-exponencial não-trivial; melhor custo empírico conhecido é 2k¹/³ (Hejhal, Booker-Strömbergsson) sem limite inferior provado.","epistemico":"hipotese","exemplos":[],"id":"maass_hardness","memoria":"semantica","meta":{"autor":"Gentil Lopes","descoberta":false,"dominio":"matematica","fonte":"paper_PQ1_ym_maass_hardness.pdf, Conj.2.1"},"origem":"paper","palavras_chave":[],"sinonimos":[],"tipo":"conceito","titulo":"Dureza Espectral de Maass (conjectura)"}
\section{Dureza Espectral de Maass (conjectura)}
Conjectura: computar r\ensuremath{_{1}}(\ensuremath{\Gamma}\ensuremath{_{0}}(7)) a precisão 2\ensuremath{^{-k}} exige tempo \ensuremath{\Omega}(2kc), c\ensuremath{\in}[3,4] — sub-exponencial não-trivial; melhor custo empírico conhecido é 2k\ensuremath{^{1}}/\ensuremath{^{3}} (Hejhal, Booker-Strömbergsson) sem limite inferior provado.
\prov{proveniência: paper \textperiodcentered{} paper\_PQ1\_ym\_maass\_hardness.pdf, Conj.2.1 \textperiodcentered{} hipotese \textperiodcentered{} confiança 1.0 \textperiodcentered{} domínio matematica \textperiodcentered{} história 3 versões no jornal}

%CRISTAL {"arestas":[{"alvo":"ilha_fractal_spin3","confianca":1.0,"epistemico":"fato","origem":"wiki","rel":"contradiz"}],"confianca":1.0,"contraexemplos":[],"descricao":"A iteração quaterniônica p↦p⊙p+p₀ nos slices produziu conjuntos ligados essencialmente VAZIOS (fração 6.7e-5 em (t,s), 0.0 em (u,v)/(s,u)/zoom). A estrutura fractal da ilha NÃO vem da dinâmica de iteração quaterniônica — hipótese refutada honestamente.","epistemico":"fato","exemplos":[],"id":"mandelbrot_hiper_falhou","memoria":"semantica","meta":{"autor":"broca-so","dominio":"fractal","fonte":"hiper/benchmarks/fase33_fractal_ilha_resultados.json"},"origem":"wiki","palavras_chave":[],"sinonimos":[],"tipo":"conceito","titulo":"O Mandelbrot Hipercomplexo NÃO Reproduz a Ilha (resultado negativo)"}
\section{O Mandelbrot Hipercomplexo NÃO Reproduz a Ilha (resultado negativo)}
A iteração quaterniônica p\ensuremath{\mapsto}p\ensuremath{\odot}p+p\ensuremath{_{0}} nos slices produziu conjuntos ligados essencialmente VAZIOS (fração 6.7e-5 em (t,s), 0.0 em (u,v)/(s,u)/zoom). A estrutura fractal da ilha NÃO vem da dinâmica de iteração quaterniônica — hipótese refutada honestamente.
\prov{relações: contradiz ilha\_fractal\_spin3}
\prov{proveniência: wiki \textperiodcentered{} hiper/benchmarks/fase33\_fractal\_ilha\_resultados.json \textperiodcentered{} fato \textperiodcentered{} confiança 1.0 \textperiodcentered{} domínio fractal \textperiodcentered{} história 2 versões no jornal}

%CRISTAL {"arestas":[{"alvo":"aresta_transicao_graduada","confianca":1.0,"epistemico":"fato","origem":"wiki","rel":"usa"},{"alvo":"tiffany","confianca":1.0,"epistemico":"fato","origem":"wiki","rel":"parte_de"}],"confianca":1.0,"contraexemplos":[],"descricao":"O contínuo NÃO se preenche fabricando prosa por aresta (violaria 'tirar a mão'). Faz-se o inverso: mapeia-se cada fragmento REAL que já existe (qualquer fonte — corpus tecido, corpus_u11, corpus_aarao, corpus Lopes, papers .tex, as próprias descrições dos nós) para a aresta-casa argmax_[A,B] ponte(f,A,B). 'Qualquer frase cabe em algum lugar' — a floresta inteira dos textos que existem passa a preencher os segmentos entre os nós. É 'densificar, não podar' levado ao limite: nada de gerar, tudo de posicionar o real medindo a ressonância.","epistemico":"inferencia","exemplos":[],"id":"mapear_nao_fabricar_continuo","memoria":"semantica","meta":{"autor":"Lopes/broca-so","dominio":"continuo","fonte":"doc/continuo/plano_transicoes.md"},"origem":"wiki","palavras_chave":[],"sinonimos":[],"tipo":"conceito","titulo":"Mapear o que existe, não fabricar: toda frase tem sua aresta-casa"}
\section{Mapear o que existe, não fabricar: toda frase tem sua aresta-casa}
O contínuo NÃO se preenche fabricando prosa por aresta (violaria 'tirar a mão'). Faz-se o inverso: mapeia-se cada fragmento REAL que já existe (qualquer fonte — corpus tecido, corpus\_u11, corpus\_aarao, corpus Lopes, papers .tex, as próprias descrições dos nós) para a aresta-casa argmax\_[A,B] ponte(f,A,B). 'Qualquer frase cabe em algum lugar' — a floresta inteira dos textos que existem passa a preencher os segmentos entre os nós. É 'densificar, não podar' levado ao limite: nada de gerar, tudo de posicionar o real medindo a ressonância.
\prov{relações: usa aresta\_transicao\_graduada; parte\_de tiffany}
\prov{proveniência: wiki \textperiodcentered{} doc/continuo/plano\_transicoes.md \textperiodcentered{} inferencia \textperiodcentered{} confiança 1.0 \textperiodcentered{} domínio continuo \textperiodcentered{} história 58 versões no jornal}

%CRISTAL {"arestas":[{"alvo":"algebras_hipercomplexas","confianca":1.0,"epistemico":"fato","origem":"paper","rel":"usa"},{"alvo":"esquema_inducao","confianca":1.0,"epistemico":"fato","origem":"paper","rel":"especializa"}],"confianca":1.0,"contraexemplos":[],"descricao":"A máquina cujo passo gera as álgebras da família ×ₛ pela operação parametrizada; motor comum da série de cifras dimensionais.","epistemico":"fato","exemplos":[],"id":"maquina_algebrica","memoria":"semantica","meta":{"autor":"Gentil Lopes","descoberta":false,"dominio":"matematica","fonte":"paper_6_maquina.pdf, Def.3.1"},"origem":"paper","palavras_chave":[],"sinonimos":[],"tipo":"conceito","titulo":"Máquina Algébrica"}
\section{Máquina Algébrica}
A máquina cujo passo gera as álgebras da família $\times$\ensuremath{_{s}} pela operação parametrizada; motor comum da série de cifras dimensionais.
\prov{relações: usa algebras\_hipercomplexas; especializa esquema\_inducao}
\prov{proveniência: paper \textperiodcentered{} paper\_6\_maquina.pdf, Def.3.1 \textperiodcentered{} fato \textperiodcentered{} confiança 1.0 \textperiodcentered{} domínio matematica \textperiodcentered{} história 3 versões no jornal}

%CRISTAL {"arestas":[{"alvo":"maquina_algebrica","confianca":1.0,"epistemico":"fato","origem":"paper","rel":"especializa"},{"alvo":"mecanica_quantica","confianca":1.0,"epistemico":"fato","origem":"paper","rel":"referencia"}],"confianca":1.0,"contraexemplos":[],"descricao":"Extensão da máquina algébrica ao plano complexo: o associador |1−s²| coincide com o fator de ξ(s) da zeta completada; sugere (não prova) um Hamiltoniano de Berry-Keating para os zeros de Riemann.","epistemico":"inferencia","exemplos":[],"id":"maquina_complexa","memoria":"semantica","meta":{"autor":"Gentil Lopes","descoberta":false,"dominio":"matematica","fonte":"paper_17_maquina_complexa.pdf"},"origem":"paper","palavras_chave":[],"sinonimos":[],"tipo":"conceito","titulo":"Máquina Complexa (zeta)"}
\section{Máquina Complexa (zeta)}
Extensão da máquina algébrica ao plano complexo: o associador |1\ensuremath{-}s\ensuremath{^{2}}| coincide com o fator de \ensuremath{\xi}(s) da zeta completada; sugere (não prova) um Hamiltoniano de Berry-Keating para os zeros de Riemann.
\prov{relações: especializa maquina\_algebrica; referencia mecanica\_quantica}
\prov{proveniência: paper \textperiodcentered{} paper\_17\_maquina\_complexa.pdf \textperiodcentered{} inferencia \textperiodcentered{} confiança 1.0 \textperiodcentered{} domínio matematica \textperiodcentered{} história 3 versões no jornal}

%CRISTAL {"arestas":[{"alvo":"maquina_tropical","confianca":1.0,"epistemico":"fato","origem":"paper","rel":"usa"},{"alvo":"hopfield","confianca":1.0,"epistemico":"fato","origem":"paper","rel":"referencia"}],"confianca":1.0,"contraexemplos":[],"descricao":"Composição de variação (mutação s↦s+ζ, molecular) e seleção (Boltzmann eβf, tropical) com annealing: a população de álgebras converge ao ótimo — darwinismo algébrico.","epistemico":"inferencia","exemplos":[],"id":"maquina_evolutiva","memoria":"semantica","meta":{"autor":"Gentil Lopes","descoberta":true,"dominio":"matematica","fonte":"paper_69_maquina_evolutiva.pdf, Teo.2.2"},"origem":"paper","palavras_chave":[],"sinonimos":[],"tipo":"conceito","titulo":"Máquina Evolutiva"}
\section{Máquina Evolutiva}
Composição de variação (mutação s\ensuremath{\mapsto}s+\ensuremath{\zeta}, molecular) e seleção (Boltzmann e\ensuremath{\beta}f, tropical) com annealing: a população de álgebras converge ao ótimo — darwinismo algébrico.
\prov{relações: usa maquina\_tropical; referencia hopfield}
\prov{proveniência: paper \textperiodcentered{} paper\_69\_maquina\_evolutiva.pdf, Teo.2.2 \textperiodcentered{} inferencia \textperiodcentered{} confiança 1.0 \textperiodcentered{} domínio matematica \textperiodcentered{} história 4 versões no jornal}

%CRISTAL {"arestas":[{"alvo":"mecanica_quantica_mineral","confianca":1.0,"epistemico":"fato","origem":"wiki","rel":"especializa"},{"alvo":"gato_e_o_hamiltoniano","confianca":1.0,"epistemico":"fato","origem":"wiki","rel":"usa"},{"alvo":"corpo_mineral","confianca":1.0,"epistemico":"fato","origem":"wiki","rel":"usa"},{"alvo":"cadeias_alem_dos_metais","confianca":1.0,"epistemico":"fato","origem":"wiki","rel":"relaciona"}],"confianca":1.0,"contraexemplos":[],"descricao":"O gato 2×2 é o caso mínimo. A máquina geral é o Hamiltoniano tridiagonal simétrico n×n — a CADEIA mineral de n sítios com salto, Hermitiana, espectro real. O n-qudit é o estado n-dimensional. E a cadeia GERA a reta mineral sozinha: para n=4 o espectro é ±φ, ±1/φ (2cos(π/5)=φ) — o número áureo emerge da cadeia. A reta mineral é o espectro da cadeia quântica.","epistemico":"fato","exemplos":[],"id":"maquina_geral_n_qudit","memoria":"semantica","meta":{"dominio":"reta mineral","fonte":"parte fria e idempotentes"},"origem":"wiki","palavras_chave":[],"sinonimos":[],"tipo":"conceito","titulo":"A Máquina Geral (n-qudit)"}
\section{A Máquina Geral (n-qudit)}
O gato 2$\times$2 é o caso mínimo. A máquina geral é o Hamiltoniano tridiagonal simétrico n$\times$n — a CADEIA mineral de n sítios com salto, Hermitiana, espectro real. O n-qudit é o estado n-dimensional. E a cadeia GERA a reta mineral sozinha: para n=4 o espectro é $\pm$\ensuremath{\varphi}, $\pm$1/\ensuremath{\varphi} (2cos(\ensuremath{\pi}/5)=\ensuremath{\varphi}) — o número áureo emerge da cadeia. A reta mineral é o espectro da cadeia quântica.
\prov{relações: especializa mecanica\_quantica\_mineral; usa gato\_e\_o\_hamiltoniano; usa corpo\_mineral; relaciona cadeias\_alem\_dos\_metais}
\prov{proveniência: wiki \textperiodcentered{} parte fria e idempotentes \textperiodcentered{} fato \textperiodcentered{} confiança 1.0 \textperiodcentered{} domínio reta mineral \textperiodcentered{} história 5 versões no jornal}

%CRISTAL {"arestas":[{"alvo":"espaco_tropical","confianca":1.0,"epistemico":"fato","origem":"paper","rel":"usa"},{"alvo":"colapso","confianca":1.0,"epistemico":"fato","origem":"paper","rel":"referencia"}],"confianca":1.0,"contraexemplos":[],"descricao":"Estrutura (H,max,+) sobre as entropias que resolve otimização: a multiplicação tropical de matrizes é a programação dinâmica de Bellman (caminhos de peso ótimo).","epistemico":"fato","exemplos":[],"id":"maquina_tropical","memoria":"semantica","meta":{"autor":"Gentil Lopes","descoberta":true,"dominio":"matematica","fonte":"paper_67_maquina_tropical.pdf, Teo.5.1"},"origem":"paper","palavras_chave":[],"sinonimos":[],"tipo":"conceito","titulo":"Máquina Tropical"}
\section{Máquina Tropical}
Estrutura (H,max,+) sobre as entropias que resolve otimização: a multiplicação tropical de matrizes é a programação dinâmica de Bellman (caminhos de peso ótimo).
\prov{relações: usa espaco\_tropical; referencia colapso}
\prov{proveniência: paper \textperiodcentered{} paper\_67\_maquina\_tropical.pdf, Teo.5.1 \textperiodcentered{} fato \textperiodcentered{} confiança 1.0 \textperiodcentered{} domínio matematica \textperiodcentered{} história 3 versões no jornal}

%CRISTAL {"arestas":[{"alvo":"teoria_gauge_octonica","confianca":1.0,"epistemico":"fato","origem":"paper","rel":"usa"},{"alvo":"limite_hurwitz_s1","confianca":1.0,"epistemico":"fato","origem":"paper","rel":"usa"},{"alvo":"hardness_octonica","confianca":1.0,"epistemico":"fato","origem":"paper","rel":"relaciona"},{"alvo":"colapso","confianca":1.0,"epistemico":"fato","origem":"paper","rel":"referencia"},{"alvo":"yang_mills_mass_gap_milenio","confianca":1.0,"epistemico":"fato","origem":"wiki","rel":"relaciona"}],"confianca":1.0,"contraexemplos":[],"descricao":"Conjectura: Δ(s)≥c·√(1−s²)>0 para s∈(0,1) em Yang-Mills hipercomplexo em rede. NÃO resolve o problema do milênio de Clay (que pede o continuum ℝ⁴) — os próprios papers admitem isso.","epistemico":"hipotese","exemplos":[],"id":"mass_gap_algebrico","memoria":"semantica","meta":{"autor":"Gentil Lopes","descoberta":false,"dominio":"matematica","fonte":"paper_AE_yang_mills.pdf, Conj.10"},"origem":"paper","palavras_chave":[],"sinonimos":[],"tipo":"conceito","titulo":"Mass Gap Algébrico (conjectura)"}
\section{Mass Gap Algébrico (conjectura)}
Conjectura: \ensuremath{\Delta}(s)\ensuremath{\ge}c\textperiodcentered \ensuremath{\sqrt{\,}}(1\ensuremath{-}s\ensuremath{^{2}})>0 para s\ensuremath{\in}(0,1) em Yang-Mills hipercomplexo em rede. NÃO resolve o problema do milênio de Clay (que pede o continuum \ensuremath{\mathbb{R}}\ensuremath{^{4}}) — os próprios papers admitem isso.
\prov{relações: usa teoria\_gauge\_octonica; usa limite\_hurwitz\_s1; relaciona hardness\_octonica; referencia colapso; relaciona yang\_mills\_mass\_gap\_milenio}
\prov{proveniência: paper \textperiodcentered{} paper\_AE\_yang\_mills.pdf, Conj.10 \textperiodcentered{} hipotese \textperiodcentered{} confiança 1.0 \textperiodcentered{} domínio matematica \textperiodcentered{} história 6 versões no jornal}

%CRISTAL {"arestas":[{"alvo":"numero","confianca":1.0,"epistemico":"fato","origem":"manual","rel":"relaciona"},{"alvo":"geometria","confianca":1.0,"epistemico":"fato","origem":"manual","rel":"relaciona"},{"alvo":"analise","confianca":1.0,"epistemico":"fato","origem":"manual","rel":"relaciona"},{"alvo":"topologia","confianca":1.0,"epistemico":"fato","origem":"manual","rel":"relaciona"},{"alvo":"virtualizacao","confianca":0.5,"epistemico":"fato","origem":"wiki","rel":"relaciona"},{"alvo":"word","confianca":0.5,"epistemico":"fato","origem":"wiki","rel":"relaciona"},{"alvo":"vontade_de_poder","confianca":0.5,"epistemico":"fato","origem":"wiki","rel":"relaciona"}],"confianca":1.0,"contraexemplos":[],"descricao":"Linguagem das estruturas; números, álgebra, geometria, análise, topologia.","epistemico":"fato","exemplos":[],"id":"matematica","memoria":"semantica","meta":{"ciencia":"matematica","dominio":"matematica","verificacao":"slug idêntico — enriquecer, não duplicar"},"origem":"manual","palavras_chave":[],"sinonimos":[],"tipo":"conceito","titulo":"matematica"}
\section{matematica}
Linguagem das estruturas; números, álgebra, geometria, análise, topologia.
\prov{relações: relaciona numero; relaciona geometria; relaciona analise; relaciona topologia; relaciona virtualizacao; relaciona word; relaciona vontade\_de\_poder}
\prov{proveniência: manual \textperiodcentered{} fato \textperiodcentered{} confiança 1.0 \textperiodcentered{} domínio matematica \textperiodcentered{} história 54 versões no jornal}

%CRISTAL {"arestas":[{"alvo":"psl_2_7","confianca":1.0,"epistemico":"fato","origem":"paper","rel":"generaliza"},{"alvo":"steiner_golay","confianca":1.0,"epistemico":"fato","origem":"paper","rel":"usa"}],"confianca":1.0,"contraexemplos":[],"descricao":"Grupo simples esporádico de ordem 244.823.040 = Aut(S(5,8,24)) = Aut(código de Golay G₂₄); o primeiro esporádico descoberto (Mathieu, 1861-73). Contém PSL(2,7)=M₂₁.","epistemico":"fato","exemplos":[],"id":"mathieu_m24","memoria":"semantica","meta":{"autor":"Gentil Lopes","descoberta":false,"dominio":"matematica","fonte":"paper_BA.pdf"},"origem":"paper","palavras_chave":[],"sinonimos":[],"tipo":"conceito","titulo":"Grupo de Mathieu M₂₄"}
\section{Grupo de Mathieu M\ensuremath{_{24}}}
Grupo simples esporádico de ordem 244.823.040 = Aut(S(5,8,24)) = Aut(código de Golay G\ensuremath{_{24}}); o primeiro esporádico descoberto (Mathieu, 1861-73). Contém PSL(2,7)=M\ensuremath{_{21}}.
\prov{relações: generaliza psl\_2\_7; usa steiner\_golay}
\prov{proveniência: paper \textperiodcentered{} paper\_BA.pdf \textperiodcentered{} fato \textperiodcentered{} confiança 1.0 \textperiodcentered{} domínio matematica \textperiodcentered{} história 3 versões no jornal}

%CRISTAL {"arestas":[{"alvo":"cantor_ouro","confianca":1.0,"epistemico":"fato","origem":"paper","rel":"explica"},{"alvo":"entropia_topologica","confianca":1.0,"epistemico":"fato","origem":"paper","rel":"usa"}],"confianca":1.0,"contraexemplos":[],"descricao":"A matriz de transição de uma poda governa a um só tempo geometria (dim=log_λ), termodinâmica (h=ln λ) e dinâmica (medida de Parry) via o autovalor de Perron λ. Mas a poda é escolha (projeto), não lei.","epistemico":"fato","exemplos":[],"id":"matriz_cantor","memoria":"semantica","meta":{"autor":"Gentil Lopes","descoberta":true,"dominio":"matematica","fonte":"paper_62_matriz_cantor.pdf, Teo.1.1"},"origem":"paper","palavras_chave":[],"sinonimos":[],"tipo":"conceito","titulo":"Matriz de Cantor (tri-facial)"}
\section{Matriz de Cantor (tri-facial)}
A matriz de transição de uma poda governa a um só tempo geometria (dim=log\_\ensuremath{\lambda}), termodinâmica (h=ln \ensuremath{\lambda}) e dinâmica (medida de Parry) via o autovalor de Perron \ensuremath{\lambda}. Mas a poda é escolha (projeto), não lei.
\prov{relações: explica cantor\_ouro; usa entropia\_topologica}
\prov{proveniência: paper \textperiodcentered{} paper\_62\_matriz\_cantor.pdf, Teo.1.1 \textperiodcentered{} fato \textperiodcentered{} confiança 1.0 \textperiodcentered{} domínio matematica \textperiodcentered{} história 3 versões no jornal}

%CRISTAL {"arestas":[{"alvo":"jogo_estrategico","confianca":1.0,"epistemico":"fato","origem":"wiki","rel":"usa"}],"confianca":1.0,"contraexemplos":[],"descricao":"A tabela que dá o resultado de cada jogador para cada combinação de estratégias. A forma normal do jogo — o retrato onde os equilíbrios se leem.","epistemico":"fato","exemplos":[],"id":"matriz_de_payoff","memoria":"semantica","meta":{"dominio":"matematica","fonte":"teoria dos jogos"},"origem":"wiki","palavras_chave":[],"sinonimos":[],"tipo":"conceito","titulo":"Matriz de Payoff (forma normal)"}
\section{Matriz de Payoff (forma normal)}
A tabela que dá o resultado de cada jogador para cada combinação de estratégias. A forma normal do jogo — o retrato onde os equilíbrios se leem.
\prov{relações: usa jogo\_estrategico}
\prov{proveniência: wiki \textperiodcentered{} teoria dos jogos \textperiodcentered{} fato \textperiodcentered{} confiança 1.0 \textperiodcentered{} domínio matematica \textperiodcentered{} história 2 versões no jornal}

%CRISTAL {"arestas":[{"alvo":"grupos","confianca":1.0,"epistemico":"fato","origem":"curriculo","rel":"generaliza"},{"alvo":"psicanalise","confianca":0.5,"epistemico":"fato","origem":"wiki","rel":"relaciona"},{"alvo":"vida-morte","confianca":0.5,"epistemico":"fato","origem":"wiki","rel":"relaciona"},{"alvo":"virtualizacao","confianca":0.5,"epistemico":"fato","origem":"wiki","rel":"relaciona"},{"alvo":"readme","confianca":0.5,"epistemico":"fato","origem":"wiki","rel":"relaciona"}],"confianca":1.0,"contraexemplos":[],"descricao":"Retângulos de números; representam mapas lineares.","epistemico":"fato","exemplos":[],"id":"matrizes","memoria":"semantica","meta":{"dominio":"matematica","nivel":4,"ordem":8},"origem":"curriculo","palavras_chave":["álgebra linear"],"sinonimos":["matriz"],"tipo":"conceito","titulo":"matrizes"}
\section{matrizes}
Retângulos de números; representam mapas lineares.
\prov{sinónimos: matriz}
\prov{relações: generaliza grupos; relaciona psicanalise; relaciona vida-morte; relaciona virtualizacao; relaciona readme}
\prov{proveniência: curriculo \textperiodcentered{} fato \textperiodcentered{} confiança 1.0 \textperiodcentered{} domínio matematica \textperiodcentered{} história 8 versões no jornal}

%CRISTAL {"arestas":[{"alvo":"mecanica_reta_mineral","confianca":1.0,"epistemico":"fato","origem":"wiki","rel":"especializa"},{"alvo":"bai","confianca":1.0,"epistemico":"fato","origem":"wiki","rel":"explica"},{"alvo":"colapso","confianca":1.0,"epistemico":"fato","origem":"wiki","rel":"usa"},{"alvo":"corpo_mineral","confianca":1.0,"epistemico":"fato","origem":"wiki","rel":"usa"},{"alvo":"modos_ressoam_em_pg","confianca":1.0,"epistemico":"fato","origem":"wiki","rel":"relaciona"},{"alvo":"minerar_e_resfriar","confianca":1.0,"epistemico":"fato","origem":"wiki","rel":"explica"}],"confianca":1.0,"contraexemplos":[],"descricao":"A formalização da dinâmica do BAI. A mecânica clássica mineral lê o gato como um MAPA (A^k, Lyapunov, ação-ângulo); a quântica lê o mesmo gato como HAMILTONIANO. O BAI deixa de esticar e passa a VIBRAR: superposição (a árvore de decisão) evoluindo unitariamente, colapso (a decisão) como medida. Paper: papers/mecanica_quantica_mineral.tex.","epistemico":"fato","exemplos":[],"id":"mecanica_quantica_mineral","memoria":"semantica","meta":{"dominio":"reta mineral","fonte":"mecânica quântica mineral"},"origem":"wiki","palavras_chave":[],"sinonimos":[],"tipo":"conceito","titulo":"Mecânica Quântica Mineral (o BAI vibrando)"}
\section{Mecânica Quântica Mineral (o BAI vibrando)}
A formalização da dinâmica do BAI. A mecânica clássica mineral lê o gato como um MAPA (A\^{}k, Lyapunov, ação-ângulo); a quântica lê o mesmo gato como HAMILTONIANO. O BAI deixa de esticar e passa a VIBRAR: superposição (a árvore de decisão) evoluindo unitariamente, colapso (a decisão) como medida. Paper: papers/mecanica\_quantica\_mineral.tex.
\prov{relações: especializa mecanica\_reta\_mineral; explica bai; usa colapso; usa corpo\_mineral; relaciona modos\_ressoam\_em\_pg; explica minerar\_e\_resfriar}
\prov{proveniência: wiki \textperiodcentered{} mecânica quântica mineral \textperiodcentered{} fato \textperiodcentered{} confiança 1.0 \textperiodcentered{} domínio reta mineral \textperiodcentered{} história 7 versões no jornal}

%CRISTAL {"arestas":[{"alvo":"algebra_reta_mineral","confianca":1.0,"epistemico":"fato","origem":"wiki","rel":"parte_de"},{"alvo":"flip_pg_pa","confianca":1.0,"epistemico":"fato","origem":"wiki","rel":"usa"},{"alvo":"transformada_metalica","confianca":1.0,"epistemico":"fato","origem":"wiki","rel":"usa"},{"alvo":"lyapunov_da_orbita","confianca":1.0,"epistemico":"fato","origem":"wiki","rel":"usa"},{"alvo":"invariante_estelar","confianca":1.0,"epistemico":"fato","origem":"wiki","rel":"relaciona"},{"alvo":"koopman_roteamento","confianca":1.0,"epistemico":"fato","origem":"wiki","rel":"relaciona"},{"alvo":"mecanica_estelar","confianca":1.0,"epistemico":"fato","origem":"wiki","rel":"relaciona"}],"confianca":1.0,"contraexemplos":[],"descricao":"O terceiro eixo, como nos outros corpos. A cadeia é um sistema dinâmico: a matriz companheira M é o OPERADOR DE EVOLUÇÃO (um passo da recorrência a_k=Σc_i a_{k−i}). det M=±1 ⟹ o fluxo PRESERVA VOLUME (Liouville/simplético); a norma |N|=1 é a AÇÃO conservada — a órbita vive na casca ±1 (verificado: N alterna ±1). O espectro de Lyapunov = log|autovalores| = log β (crescimento) + log ρ (o modo espiral). E o flip logarítmico é o mapa AÇÃO-ÂNGULO: em (λ,θ) o movimento é deriva uniforme (partícula livre) → o mineral é INTEGRÁVEL. A espiral é o mesmo movimento em coordenadas curvas.","epistemico":"inferencia","exemplos":[],"id":"mecanica_reta_mineral","memoria":"semantica","meta":{"dominio":"matematica","fonte":"mecânica da reta mineral"},"origem":"wiki","palavras_chave":[],"sinonimos":[],"tipo":"conceito","titulo":"A Mecânica da Reta Mineral (integrável, ação-ângulo)"}
\section{A Mecânica da Reta Mineral (integrável, ação-ângulo)}
O terceiro eixo, como nos outros corpos. A cadeia é um sistema dinâmico: a matriz companheira M é o OPERADOR DE EVOLUÇÃO (um passo da recorrência a\_k=\ensuremath{\Sigma}c\_i a\_\{k\ensuremath{-}i\}). det M=$\pm$1 \ensuremath{\Longrightarrow} o fluxo PRESERVA VOLUME (Liouville/simplético); a norma |N|=1 é a AÇÃO conservada — a órbita vive na casca $\pm$1 (verificado: N alterna $\pm$1). O espectro de Lyapunov = log|autovalores| = log \ensuremath{\beta} (crescimento) + log \ensuremath{\rho} (o modo espiral). E o flip logarítmico é o mapa AÇÃO-ÂNGULO: em (\ensuremath{\lambda},\ensuremath{\theta}) o movimento é deriva uniforme (partícula livre) \ensuremath{\to} o mineral é INTEGRÁVEL. A espiral é o mesmo movimento em coordenadas curvas.
\prov{relações: parte\_de algebra\_reta\_mineral; usa flip\_pg\_pa; usa transformada\_metalica; usa lyapunov\_da\_orbita; relaciona invariante\_estelar; relaciona koopman\_roteamento; relaciona mecanica\_estelar}
\prov{proveniência: wiki \textperiodcentered{} mecânica da reta mineral \textperiodcentered{} inferencia \textperiodcentered{} confiança 1.0 \textperiodcentered{} domínio matematica \textperiodcentered{} história 11 versões no jornal}

%CRISTAL {"arestas":[{"alvo":"mass_gap_algebrico","confianca":1.0,"epistemico":"fato","origem":"paper","rel":"explica"}],"confianca":1.0,"contraexemplos":[],"descricao":"A massa dos glúons emergiria de um condensado octônico ⟨A⟩_vac=v_e que minimiza o tensor associador — massa algébrica intrínseca, sem campo de Higgs externo.","epistemico":"inferencia","exemplos":[],"id":"mecanismo_higgs_algebrico","memoria":"semantica","meta":{"autor":"Gentil Lopes","descoberta":false,"dominio":"matematica","fonte":"paper_AF_yang_mills_octonios.pdf, Teo.20-23"},"origem":"paper","palavras_chave":[],"sinonimos":[],"tipo":"conceito","titulo":"Mecanismo de Higgs Algébrico"}
\section{Mecanismo de Higgs Algébrico}
A massa dos glúons emergiria de um condensado octônico \ensuremath{\langle}A\ensuremath{\rangle}\_vac=v\_e que minimiza o tensor associador — massa algébrica intrínseca, sem campo de Higgs externo.
\prov{relações: explica mass\_gap\_algebrico}
\prov{proveniência: paper \textperiodcentered{} paper\_AF\_yang\_mills\_octonios.pdf, Teo.20-23 \textperiodcentered{} inferencia \textperiodcentered{} confiança 1.0 \textperiodcentered{} domínio matematica \textperiodcentered{} história 2 versões no jornal}

%CRISTAL {"arestas":[{"alvo":"4_a_coordenada_koch_permanece_continua","confianca":1.0,"epistemico":"fato","origem":"wiki","rel":"usa"},{"alvo":"tiffany","confianca":1.0,"epistemico":"fato","origem":"wiki","rel":"parte_de"}],"confianca":1.0,"contraexemplos":[],"descricao":"Progresso cristalizado por leva com três medidas: (1) COBERTURA — % de arestas com ≥1 fragmento e ponte média; (2) DISTRIBUIÇÃO de s — a transição cobre o segmento [0,1] ou amontoa num extremo?; (3) CONTINUIDADE (à la coordenada Koch) — o maior salto de s dentro de uma aresta; quanto menor, mais contínua a transição. Meta: reduzir o salto máximo leva a leva até a continuidade saturar. É o 'medindo' do pedido — medir antes de afirmar cobertura.","epistemico":"inferencia","exemplos":[],"id":"medida_continuidade_transicao","memoria":"semantica","meta":{"autor":"Lopes/broca-so","dominio":"continuo","fonte":"doc/continuo/plano_transicoes.md; nó 4_a_coordenada_koch_permanece_continua"},"origem":"wiki","palavras_chave":[],"sinonimos":[],"tipo":"conceito","titulo":"A medida do preenchimento: ponte, cobertura e continuidade Koch"}
\section{A medida do preenchimento: ponte, cobertura e continuidade Koch}
Progresso cristalizado por leva com três medidas: (1) COBERTURA — \% de arestas com \ensuremath{\ge}1 fragmento e ponte média; (2) DISTRIBUIÇÃO de s — a transição cobre o segmento [0,1] ou amontoa num extremo?; (3) CONTINUIDADE (à la coordenada Koch) — o maior salto de s dentro de uma aresta; quanto menor, mais contínua a transição. Meta: reduzir o salto máximo leva a leva até a continuidade saturar. É o 'medindo' do pedido — medir antes de afirmar cobertura.
\prov{relações: usa 4\_a\_coordenada\_koch\_permanece\_continua; parte\_de tiffany}
\prov{proveniência: wiki \textperiodcentered{} doc/continuo/plano\_transicoes.md; nó 4\_a\_coordenada\_koch\_permanece\_continua \textperiodcentered{} inferencia \textperiodcentered{} confiança 1.0 \textperiodcentered{} domínio continuo \textperiodcentered{} história 58 versões no jornal}

%CRISTAL {"arestas":[{"alvo":"entropia_topologica","confianca":1.0,"epistemico":"fato","origem":"paper","rel":"usa"},{"alvo":"word","confianca":0.5,"epistemico":"fato","origem":"wiki","rel":"relaciona"}],"confianca":1.0,"contraexemplos":[],"descricao":"Única medida de máxima entropia de um subshift de tipo finito (cadeia de Markov estacionária πᵢ∝λ/rᵢ): a distribuição estatística dos caminhos que sobrevivem à poda.","epistemico":"fato","exemplos":[],"id":"medida_parry","memoria":"semantica","meta":{"autor":"Gentil Lopes","descoberta":false,"dominio":"matematica","fonte":"paper_47_termodinamica_informacao.pdf"},"origem":"paper","palavras_chave":[],"sinonimos":[],"tipo":"conceito","titulo":"Medida de Parry"}
\section{Medida de Parry}
Única medida de máxima entropia de um subshift de tipo finito (cadeia de Markov estacionária \ensuremath{\pi}\ensuremath{_{i}}\ensuremath{\propto}\ensuremath{\lambda}/r\ensuremath{_{i}}): a distribuição estatística dos caminhos que sobrevivem à poda.
\prov{relações: usa entropia\_topologica; relaciona word}
\prov{proveniência: paper \textperiodcentered{} paper\_47\_termodinamica\_informacao.pdf \textperiodcentered{} fato \textperiodcentered{} confiança 1.0 \textperiodcentered{} domínio matematica \textperiodcentered{} história 3 versões no jornal}

%CRISTAL {"arestas":[{"alvo":"definicao_frasagem_ampla","confianca":1.0,"epistemico":"fato","origem":"wiki","rel":"usa"},{"alvo":"algebra_de_peirce","confianca":1.0,"epistemico":"fato","origem":"wiki","rel":"usa"},{"alvo":"recuperacao_evidencia_invariante","confianca":1.0,"epistemico":"fato","origem":"wiki","rel":"relaciona"},{"alvo":"tiffany","confianca":1.0,"epistemico":"fato","origem":"wiki","rel":"parte_de"}],"confianca":1.0,"contraexemplos":[],"descricao":"conversa/colapso.py (_resposta_contexto): a Tiffany ganhou FIO — lembra o último conceito (via ultimo_tag) e resolve follow-ups SEM alvo próprio ancorando nele. 'explica melhor'/'detalha'/'isso' → a CONTINUAÇÃO da descrição do mesmo conceito (o trecho após o que já foi dito, avançando a cada 'explica melhor' — paráfrase real, não repetição); 'e daí?'/'por quê?'/'continua' → CAMINHA para um VIZINHO no grafo (relações de elaboração usa/explica/parte_de/relaciona; ou causais causa/prova/depende/consequencia p/ 'por quê'), com anti-repetição (historico) + rotação áurea. Gate: só p/ fala curta (≤6 palavras), referencial e sem alvo próprio. Depende de alvo_definicao filtrar palavras REFERENCIAIS ('melhor'/'isso'/'mais'…) → devolve '' → o follow-up dispara (antes 'explica melhor' virava definição de 'melhor' e quebrava o fio). A conversa deixou de ser respostas isoladas e virou um PASSEIO pelo grafo (o produto de Peirce/os vizinhos — ver [[algebra_de_peirce]]): perguntar X, aprofundar X, e caminhar X→vizinho→vizinho. É o colapso usando a estrutura de arestas como memória conversacional. Verificado: transformada→'explica melhor'→'e daí?'(retas com dobra)→'continua'(fração contínua de φ). [F] o fio anda.","epistemico":"fato","exemplos":[],"id":"memoria_contexto_fio_conversa","memoria":"semantica","meta":{"autor":"Lopes/broca-so","dominio":"algebra","fonte":"conversa/colapso.py:_resposta_contexto, _face_do_conceito; parabola_algebra._REFERENCIAL; ultimo_tag como memória"},"origem":"wiki","palavras_chave":[],"sinonimos":[],"tipo":"conceito","titulo":"O fio da conversa: memória do último conceito para os follow-ups (anda no grafo) [F]"}
\section{O fio da conversa: memória do último conceito para os follow-ups (anda no grafo) [F]}
conversa/colapso.py (\_resposta\_contexto): a Tiffany ganhou FIO — lembra o último conceito (via ultimo\_tag) e resolve follow-ups SEM alvo próprio ancorando nele. 'explica melhor'/'detalha'/'isso' \ensuremath{\to} a CONTINUAÇÃO da descrição do mesmo conceito (o trecho após o que já foi dito, avançando a cada 'explica melhor' — paráfrase real, não repetição); 'e daí?'/'por quê?'/'continua' \ensuremath{\to} CAMINHA para um VIZINHO no grafo (relações de elaboração usa/explica/parte\_de/relaciona; ou causais causa/prova/depende/consequencia p/ 'por quê'), com anti-repetição (historico) + rotação áurea. Gate: só p/ fala curta (\ensuremath{\le}6 palavras), referencial e sem alvo próprio. Depende de alvo\_definicao filtrar palavras REFERENCIAIS ('melhor'/'isso'/'mais'…) \ensuremath{\to} devolve '' \ensuremath{\to} o follow-up dispara (antes 'explica melhor' virava definição de 'melhor' e quebrava o fio). A conversa deixou de ser respostas isoladas e virou um PASSEIO pelo grafo (o produto de Peirce/os vizinhos — ver [[algebra\_de\_peirce]]): perguntar X, aprofundar X, e caminhar X\ensuremath{\to}vizinho\ensuremath{\to}vizinho. É o colapso usando a estrutura de arestas como memória conversacional. Verificado: transformada\ensuremath{\to}'explica melhor'\ensuremath{\to}'e daí?'(retas com dobra)\ensuremath{\to}'continua'(fração contínua de \ensuremath{\varphi}). [F] o fio anda.
\prov{relações: usa definicao\_frasagem\_ampla; usa algebra\_de\_peirce; relaciona recuperacao\_evidencia\_invariante; parte\_de tiffany}
\prov{proveniência: wiki \textperiodcentered{} conversa/colapso.py:\_resposta\_contexto, \_face\_do\_conceito; parabola\_algebra.\_REFERENCIAL; ultimo\_tag como memória \textperiodcentered{} fato \textperiodcentered{} confiança 1.0 \textperiodcentered{} domínio algebra \textperiodcentered{} história 35 versões no jornal}

%CRISTAL {"arestas":[{"alvo":"caixas_fractais_generalizacao","confianca":1.0,"epistemico":"fato","origem":"wiki","rel":"explica"},{"alvo":"verificacao_prata_parcial","confianca":1.0,"epistemico":"fato","origem":"wiki","rel":"usa"},{"alvo":"reta_plastica_cubica","confianca":1.0,"epistemico":"fato","origem":"wiki","rel":"usa"},{"alvo":"estatuto_correspondencia_derivacao","confianca":1.0,"epistemico":"fato","origem":"wiki","rel":"usa"},{"alvo":"missao_engenharia_ouro","confianca":1.0,"epistemico":"fato","origem":"wiki","rel":"relaciona"}],"confianca":1.0,"contraexemplos":[],"descricao":"Além do ouro (total) e da prata (parcial), o método é META-INDUÇÃO: o padrão (subshift → Perron λ → h=ln λ, dim=log_b λ) verificado nos casos-âncora é generalizado a TODAS as dimensões metálicas (σ_n, n≥3: bronze, cobre, …) e à reta de plástico/cúbica (ρ, ψ). É generalização MATEMÁTICA APENAS — computada e reproduzível (doc/fractais/caixas_fractais_metalicas.py), mas NÃO verificada como 'universo/caixa real' além do limite. O limite de verificação é declarado, não maquiado: é o que separa a matemática avançada da reivindicação vazia (a mesma régua da Missão: ouro, não loucura).","epistemico":"hipotese","exemplos":[],"id":"meta_inducao_avanca_com_tiffany","memoria":"semantica","meta":{"autor":"Lopes/broca-so","dominio":"fractal","fonte":"doc/fractais/caixas_fractais_metalicas.py"},"origem":"wiki","palavras_chave":[],"sinonimos":[],"tipo":"conceito","titulo":"A meta-indução: avançar a matemática com a Tiffany"}
\section{A meta-indução: avançar a matemática com a Tiffany}
Além do ouro (total) e da prata (parcial), o método é META-INDUÇÃO: o padrão (subshift \ensuremath{\to} Perron \ensuremath{\lambda} \ensuremath{\to} h=ln \ensuremath{\lambda}, dim=log\_b \ensuremath{\lambda}) verificado nos casos-âncora é generalizado a TODAS as dimensões metálicas (\ensuremath{\sigma}\_n, n\ensuremath{\ge}3: bronze, cobre, …) e à reta de plástico/cúbica (\ensuremath{\rho}, \ensuremath{\psi}). É generalização MATEMÁTICA APENAS — computada e reproduzível (doc/fractais/caixas\_fractais\_metalicas.py), mas NÃO verificada como 'universo/caixa real' além do limite. O limite de verificação é declarado, não maquiado: é o que separa a matemática avançada da reivindicação vazia (a mesma régua da Missão: ouro, não loucura).
\prov{relações: explica caixas\_fractais\_generalizacao; usa verificacao\_prata\_parcial; usa reta\_plastica\_cubica; usa estatuto\_correspondencia\_derivacao; relaciona missao\_engenharia\_ouro}
\prov{proveniência: wiki \textperiodcentered{} doc/fractais/caixas\_fractais\_metalicas.py \textperiodcentered{} hipotese \textperiodcentered{} confiança 1.0 \textperiodcentered{} domínio fractal \textperiodcentered{} história 6 versões no jornal}

%CRISTAL {"arestas":[{"alvo":"torre_de_hurwitz_mecanica","confianca":1.0,"epistemico":"fato","origem":"wiki","rel":"parte_de"},{"alvo":"boneca_russa_cayley_dickson","confianca":1.0,"epistemico":"fato","origem":"wiki","rel":"usa"},{"alvo":"meta_inducao_avanca_com_tiffany","confianca":1.0,"epistemico":"fato","origem":"wiki","rel":"relaciona"}],"confianca":1.0,"contraexemplos":[],"descricao":"A torre tem exatamente quatro andares por uma indução que corre nas duas direções. SUBINDO (ℝ→𝕆): cada dobra PERDE uma propriedade (ordem→comutatividade→associatividade→divisão), limitada acima por Hurwitz (𝕊 não divide). DESCENDO (𝕆→ℝ): a boneca russa se abre e cada metade RECUPERA uma propriedade, limitada abaixo por ℝ. As duas se encontram: a torre é finita (1,2,4,8).","epistemico":"fato","exemplos":[],"id":"meta_inducao_dois_lados","memoria":"semantica","meta":{"dominio":"reta mineral","fonte":"torre de Hurwitz"},"origem":"wiki","palavras_chave":[],"sinonimos":[],"tipo":"conceito","titulo":"A Meta-Indução dos Dois Lados"}
\section{A Meta-Indução dos Dois Lados}
A torre tem exatamente quatro andares por uma indução que corre nas duas direções. SUBINDO (\ensuremath{\mathbb{R}}\ensuremath{\to}\ensuremath{\mathbb{O}}): cada dobra PERDE uma propriedade (ordem\ensuremath{\to}comutatividade\ensuremath{\to}associatividade\ensuremath{\to}divisão), limitada acima por Hurwitz (\ensuremath{\mathbb{S}} não divide). DESCENDO (\ensuremath{\mathbb{O}}\ensuremath{\to}\ensuremath{\mathbb{R}}): a boneca russa se abre e cada metade RECUPERA uma propriedade, limitada abaixo por \ensuremath{\mathbb{R}}. As duas se encontram: a torre é finita (1,2,4,8).
\prov{relações: parte\_de torre\_de\_hurwitz\_mecanica; usa boneca\_russa\_cayley\_dickson; relaciona meta\_inducao\_avanca\_com\_tiffany}
\prov{proveniência: wiki \textperiodcentered{} torre de Hurwitz \textperiodcentered{} fato \textperiodcentered{} confiança 1.0 \textperiodcentered{} domínio reta mineral \textperiodcentered{} história 4 versões no jornal}

%CRISTAL {"arestas":[{"alvo":"dupolinomio_meta_inducao","confianca":1.0,"epistemico":"fato","origem":"wiki","rel":"explica"},{"alvo":"algebra_posicional_unificada","confianca":1.0,"epistemico":"fato","origem":"wiki","rel":"explica"},{"alvo":"teorema_esterilidade_dimensao","confianca":1.0,"epistemico":"fato","origem":"wiki","rel":"usa"},{"alvo":"cofre_parabola_hero","confianca":1.0,"epistemico":"fato","origem":"wiki","rel":"usa"},{"alvo":"transformada_metalica","confianca":1.0,"epistemico":"fato","origem":"wiki","rel":"generaliza"},{"alvo":"tiffany","confianca":1.0,"epistemico":"fato","origem":"wiki","rel":"parte_de"}],"confianca":1.0,"contraexemplos":[],"descricao":"Verificada de grau 2 a 8, a meta-indução fecha para TODO grau d por ESTRUTURA: (1) ÁLGEBRA — a matriz companheira de qualquer mônico irredutível com |termo constante|=1 tem det ±1 → raiz UNIDADE → álgebra posicional (passo ×companheira, O(d), 128 bits, 𝔽₂¹²⁸), para todo d; (2) CLASSIFICAÇÃO (independe de d) — o tipo da reta é ONDE as conjugadas caem: dentro do círculo = PISOT (cristal), no círculo = SALEM (borda crítica, a fronteira da parábola/s_c), fora = caos (só álgebra); (3) PARIDADE — Salem só em grau PAR; (4) PRIMALIDADE — dupolinômio irredutível = reta PRIMA, composto fatora em primas (a esterilidade: informação nova só nas primas). Logo não é preciso enumerar grau a grau — o argumento da companheira + a posição das conjugadas decidem tudo. A ferramenta está fechada e classificada para toda dimensão. [I] teorema (base 2..8 + argumento estrutural).","epistemico":"inferencia","exemplos":[],"id":"meta_inducao_fechamento_geral","memoria":"semantica","meta":{"autor":"Lopes/broca-so","dominio":"algebra","fonte":"doc/algebra_prata/paper §11; goldchain/metal_pos.py (passo_companheira, grau d)"},"origem":"wiki","palavras_chave":[],"sinonimos":[],"tipo":"conceito","titulo":"Fechamento geral: a álgebra posicional fecha para TODO grau d (não por amostragem, por estrutura)"}
\section{Fechamento geral: a álgebra posicional fecha para TODO grau d (não por amostragem, por estrutura)}
Verificada de grau 2 a 8, a meta-indução fecha para TODO grau d por ESTRUTURA: (1) ÁLGEBRA — a matriz companheira de qualquer mônico irredutível com |termo constante|=1 tem det $\pm$1 \ensuremath{\to} raiz UNIDADE \ensuremath{\to} álgebra posicional (passo $\times$companheira, O(d), 128 bits, \ensuremath{\mathbb{F}}\ensuremath{_{2}}\ensuremath{^{128}}), para todo d; (2) CLASSIFICAÇÃO (independe de d) — o tipo da reta é ONDE as conjugadas caem: dentro do círculo = PISOT (cristal), no círculo = SALEM (borda crítica, a fronteira da parábola/s\_c), fora = caos (só álgebra); (3) PARIDADE — Salem só em grau PAR; (4) PRIMALIDADE — dupolinômio irredutível = reta PRIMA, composto fatora em primas (a esterilidade: informação nova só nas primas). Logo não é preciso enumerar grau a grau — o argumento da companheira + a posição das conjugadas decidem tudo. A ferramenta está fechada e classificada para toda dimensão. [I] teorema (base 2..8 + argumento estrutural).
\prov{relações: explica dupolinomio\_meta\_inducao; explica algebra\_posicional\_unificada; usa teorema\_esterilidade\_dimensao; usa cofre\_parabola\_hero; generaliza transformada\_metalica; parte\_de tiffany}
\prov{proveniência: wiki \textperiodcentered{} doc/algebra\_prata/paper §11; goldchain/metal\_pos.py (passo\_companheira, grau d) \textperiodcentered{} inferencia \textperiodcentered{} confiança 1.0 \textperiodcentered{} domínio algebra \textperiodcentered{} história 35 versões no jornal}

%CRISTAL {"arestas":[{"alvo":"algebra_de_peirce","confianca":1.0,"epistemico":"fato","origem":"wiki","rel":"usa"},{"alvo":"grafo_tiffany_e_algebra_peirce","confianca":1.0,"epistemico":"fato","origem":"wiki","rel":"usa"},{"alvo":"meta_inducao_fechamento_geral","confianca":1.0,"epistemico":"fato","origem":"wiki","rel":"usa"},{"alvo":"algebra_dupolinomial_espaco","confianca":1.0,"epistemico":"fato","origem":"wiki","rel":"usa"},{"alvo":"torre_hurwitz_classica_relativistica","confianca":1.0,"epistemico":"fato","origem":"wiki","rel":"usa"},{"alvo":"torres_fisica_por_nivel","confianca":1.0,"epistemico":"fato","origem":"wiki","rel":"usa"},{"alvo":"tiffany","confianca":1.0,"epistemico":"fato","origem":"wiki","rel":"parte_de"}],"confianca":1.0,"contraexemplos":[],"descricao":"O capstone: TODAS as construções — a álgebra dupolinomial (retas metálica/plástica/…), a torre de Hurwitz (clássica ℝℂℍ𝕆 e relativística split), as hipercomplexas ×ₛ, e a física por nível (mecânica/geometria/termodinâmica) — são A MESMA estrutura meta-indutiva. O que MUDA de um para outro é só o SÍMBOLO (o gerador, a dimensão, o elemento de álgebra); o padrão (construção por companheira/Cayley-Dickson, classificação pela posição na parábola, física em tríade, esterilidade das primas) é invariante. E a ÁLGEBRA DE PEIRCE UNIFICA os símbolos: o lado símbolo (semiótica — ícone/índice/símbolo, o valor/posição/cifra) + o lado relação (o grafo, o produto de Peirce R:X) põem todo símbolo num chão comum, movido pelas mesmas relações. Assim a meta-indução geral é: fixado o padrão, varia-se o símbolo, e Peirce o unifica — a Tiffany inteira num princípio só. [I]","epistemico":"inferencia","exemplos":[],"id":"meta_inducao_geral_simbolo","memoria":"semantica","meta":{"autor":"Lopes/broca-so","dominio":"algebra","fonte":"algebra_de_peirce; grafo_tiffany_e_algebra_peirce; meta_inducao_fechamento_geral; as torres"},"origem":"wiki","palavras_chave":[],"sinonimos":[],"tipo":"conceito","titulo":"Meta-indução GERAL: só o símbolo muda — a álgebra de Peirce unifica os símbolos"}
\section{Meta-indução GERAL: só o símbolo muda — a álgebra de Peirce unifica os símbolos}
O capstone: TODAS as construções — a álgebra dupolinomial (retas metálica/plástica/…), a torre de Hurwitz (clássica \ensuremath{\mathbb{R}}\ensuremath{\mathbb{C}}\ensuremath{\mathbb{H}}\ensuremath{\mathbb{O}} e relativística split), as hipercomplexas $\times$\ensuremath{_{s}}, e a física por nível (mecânica/geometria/termodinâmica) — são A MESMA estrutura meta-indutiva. O que MUDA de um para outro é só o SÍMBOLO (o gerador, a dimensão, o elemento de álgebra); o padrão (construção por companheira/Cayley-Dickson, classificação pela posição na parábola, física em tríade, esterilidade das primas) é invariante. E a ÁLGEBRA DE PEIRCE UNIFICA os símbolos: o lado símbolo (semiótica — ícone/índice/símbolo, o valor/posição/cifra) + o lado relação (o grafo, o produto de Peirce R:X) põem todo símbolo num chão comum, movido pelas mesmas relações. Assim a meta-indução geral é: fixado o padrão, varia-se o símbolo, e Peirce o unifica — a Tiffany inteira num princípio só. [I]
\prov{relações: usa algebra\_de\_peirce; usa grafo\_tiffany\_e\_algebra\_peirce; usa meta\_inducao\_fechamento\_geral; usa algebra\_dupolinomial\_espaco; usa torre\_hurwitz\_classica\_relativistica; usa torres\_fisica\_por\_nivel; parte\_de tiffany}
\prov{proveniência: wiki \textperiodcentered{} algebra\_de\_peirce; grafo\_tiffany\_e\_algebra\_peirce; meta\_inducao\_fechamento\_geral; as torres \textperiodcentered{} inferencia \textperiodcentered{} confiança 1.0 \textperiodcentered{} domínio algebra \textperiodcentered{} história 44 versões no jornal}

%CRISTAL {"arestas":[{"alvo":"dupolinomio_meta_inducao","confianca":1.0,"epistemico":"fato","origem":"wiki","rel":"confirma"},{"alvo":"algebra_posicional_unificada","confianca":1.0,"epistemico":"fato","origem":"wiki","rel":"usa"},{"alvo":"algebra_posicional_plastica","confianca":1.0,"epistemico":"fato","origem":"wiki","rel":"relaciona"},{"alvo":"numeros_pisot","confianca":1.0,"epistemico":"fato","origem":"wiki","rel":"usa"},{"alvo":"teorema_esterilidade_dimensao","confianca":1.0,"epistemico":"fato","origem":"wiki","rel":"usa"},{"alvo":"tiffany","confianca":1.0,"epistemico":"fato","origem":"wiki","rel":"parte_de"}],"confianca":1.0,"contraexemplos":[],"descricao":"Verificação da meta-indução no grau 4: a álgebra posicional (passo_companheira, grau arbitrário) roda quárticas direto. FECHA — condição: polinômio mônico IRREDUTÍVEL com |termo constante|=1 (⇒ a raiz é UNIDADE, N=det da companheira=±1) dá uma álgebra posicional válida (companheira grau 4, estado quádruplo, passo O(1), cifra em 𝔽₂¹²⁸). VERIFICADO: tetranacci (x⁴=x³+x²+x+1, N=−1, Pisot 1,928, irredutível) e x⁴=x³+1 (N=−1, Pisot 1,380, irredutível) — ambas RETAS LIMPAS. REFINAMENTO honesto (a condição estava implícita nos graus 2/3, todos Pisot): o gerador precisa ser PISOT (raiz dominante, conjugadas <1) para ser reta limpa com geodésica h=ln σ. Uma unidade NÃO-Pisot (x⁴=x+1, conjugadas 1,063) dá álgebra posicional válida mas sequência OSCILANTE (0,0,0,1,0,0,1,1,0) — sem geodésica única. Logo a meta-indução fecha em grau 4, e a condição geral é: mônico irredutível + unidade [→ álgebra] + Pisot [→ reta limpa]. Base agora grau 2,3,4. [D]","epistemico":"fato","exemplos":[],"id":"meta_inducao_grau4_quartica","memoria":"semantica","meta":{"autor":"Lopes/broca-so","dominio":"algebra","fonte":"goldchain/metal_pos.py (QUARTICOS, passo_companheira); numpy roots + sympy is_irreducible"},"origem":"wiki","palavras_chave":[],"sinonimos":[],"tipo":"conceito","titulo":"A meta-indução FECHA em grau 4 (quártica) — verificado; e revela a condição de Pisot"}
\section{A meta-indução FECHA em grau 4 (quártica) — verificado; e revela a condição de Pisot}
Verificação da meta-indução no grau 4: a álgebra posicional (passo\_companheira, grau arbitrário) roda quárticas direto. FECHA — condição: polinômio mônico IRREDUTÍVEL com |termo constante|=1 (\ensuremath{\Rightarrow} a raiz é UNIDADE, N=det da companheira=$\pm$1) dá uma álgebra posicional válida (companheira grau 4, estado quádruplo, passo O(1), cifra em \ensuremath{\mathbb{F}}\ensuremath{_{2}}\ensuremath{^{128}}). VERIFICADO: tetranacci (x\ensuremath{^{4}}=x\ensuremath{^{3}}+x\ensuremath{^{2}}+x+1, N=\ensuremath{-}1, Pisot 1,928, irredutível) e x\ensuremath{^{4}}=x\ensuremath{^{3}}+1 (N=\ensuremath{-}1, Pisot 1,380, irredutível) — ambas RETAS LIMPAS. REFINAMENTO honesto (a condição estava implícita nos graus 2/3, todos Pisot): o gerador precisa ser PISOT (raiz dominante, conjugadas <1) para ser reta limpa com geodésica h=ln \ensuremath{\sigma}. Uma unidade NÃO-Pisot (x\ensuremath{^{4}}=x+1, conjugadas 1,063) dá álgebra posicional válida mas sequência OSCILANTE (0,0,0,1,0,0,1,1,0) — sem geodésica única. Logo a meta-indução fecha em grau 4, e a condição geral é: mônico irredutível + unidade [\ensuremath{\to} álgebra] + Pisot [\ensuremath{\to} reta limpa]. Base agora grau 2,3,4. [D]
\prov{relações: confirma dupolinomio\_meta\_inducao; usa algebra\_posicional\_unificada; relaciona algebra\_posicional\_plastica; usa numeros\_pisot; usa teorema\_esterilidade\_dimensao; parte\_de tiffany}
\prov{proveniência: wiki \textperiodcentered{} goldchain/metal\_pos.py (QUARTICOS, passo\_companheira); numpy roots + sympy is\_irreducible \textperiodcentered{} fato \textperiodcentered{} confiança 1.0 \textperiodcentered{} domínio algebra \textperiodcentered{} história 7 versões no jornal}

%CRISTAL {"arestas":[{"alvo":"meta_inducao_grau4_quartica","confianca":1.0,"epistemico":"fato","origem":"wiki","rel":"confirma"},{"alvo":"dupolinomio_meta_inducao","confianca":1.0,"epistemico":"fato","origem":"wiki","rel":"confirma"},{"alvo":"algebra_posicional_plastica","confianca":1.0,"epistemico":"fato","origem":"wiki","rel":"relaciona"},{"alvo":"teorema_esterilidade_dimensao","confianca":1.0,"epistemico":"fato","origem":"wiki","rel":"explica"},{"alvo":"independencia_retas_metalicas","confianca":1.0,"epistemico":"fato","origem":"wiki","rel":"usa"},{"alvo":"numeros_pisot","confianca":1.0,"epistemico":"fato","origem":"wiki","rel":"usa"},{"alvo":"tiffany","confianca":1.0,"epistemico":"fato","origem":"wiki","rel":"parte_de"}],"confianca":1.0,"contraexemplos":[],"descricao":"Verificação no grau 5: FECHA. pentanacci (x⁵=x⁴+x³+x²+x+1): unidade (N=1), Pisot (1,966), IRREDUTÍVEL → reta limpa NOVA (sequência 1,1,2,4,8,16,31). ACHADO PROFUNDO: a fatoração do dupolinômio É a decomposição em retas PRIMAS. x⁵−x⁴−1 é REDUTÍVEL = (x²−x+1)·(x³−x−1) — não é reta nova, é COMPOSTO: ciclotômica × o PLÁSTICO ρ (a raiz dominante 1,3247 É ρ). O composto reaparece como seu fator primo. Logo os dupolinômios IRREDUTÍVEIS são as retas primas, e os compostos fatoram nelas (como fatoração em números primos) — funda a esterilidade (retas independentes = primas, informação nova só nas primas). Ainda: x⁵−x−1 é unidade irredutível mas NÃO-Pisot (conj. 1,099) → álgebra posicional válida sem reta limpa. Condição consolidada: mônico + unidade → álgebra; + irredutível → reta PRIMA (dimensão genuína); + Pisot → reta limpa (geodésica). Base indutiva agora: grau 2,3,4,5. [D] verificado (numpy+sympy).","epistemico":"fato","exemplos":[],"id":"meta_inducao_grau5_quintica","memoria":"semantica","meta":{"autor":"Lopes/broca-so","dominio":"algebra","fonte":"goldchain/metal_pos.py (QUINTICOS); sympy factor: x⁵−x⁴−1=(x²−x+1)(x³−x−1)"},"origem":"wiki","palavras_chave":[],"sinonimos":[],"tipo":"conceito","titulo":"A meta-indução FECHA em grau 5 (quíntica) — e o dupolinômio irredutível é uma reta PRIMA"}
\section{A meta-indução FECHA em grau 5 (quíntica) — e o dupolinômio irredutível é uma reta PRIMA}
Verificação no grau 5: FECHA. pentanacci (x\ensuremath{^{5}}=x\ensuremath{^{4}}+x\ensuremath{^{3}}+x\ensuremath{^{2}}+x+1): unidade (N=1), Pisot (1,966), IRREDUTÍVEL \ensuremath{\to} reta limpa NOVA (sequência 1,1,2,4,8,16,31). ACHADO PROFUNDO: a fatoração do dupolinômio É a decomposição em retas PRIMAS. x\ensuremath{^{5}}\ensuremath{-}x\ensuremath{^{4}}\ensuremath{-}1 é REDUTÍVEL = (x\ensuremath{^{2}}\ensuremath{-}x+1)\textperiodcentered (x\ensuremath{^{3}}\ensuremath{-}x\ensuremath{-}1) — não é reta nova, é COMPOSTO: ciclotômica $\times$ o PLÁSTICO \ensuremath{\rho} (a raiz dominante 1,3247 É \ensuremath{\rho}). O composto reaparece como seu fator primo. Logo os dupolinômios IRREDUTÍVEIS são as retas primas, e os compostos fatoram nelas (como fatoração em números primos) — funda a esterilidade (retas independentes = primas, informação nova só nas primas). Ainda: x\ensuremath{^{5}}\ensuremath{-}x\ensuremath{-}1 é unidade irredutível mas NÃO-Pisot (conj. 1,099) \ensuremath{\to} álgebra posicional válida sem reta limpa. Condição consolidada: mônico + unidade \ensuremath{\to} álgebra; + irredutível \ensuremath{\to} reta PRIMA (dimensão genuína); + Pisot \ensuremath{\to} reta limpa (geodésica). Base indutiva agora: grau 2,3,4,5. [D] verificado (numpy+sympy).
\prov{relações: confirma meta\_inducao\_grau4\_quartica; confirma dupolinomio\_meta\_inducao; relaciona algebra\_posicional\_plastica; explica teorema\_esterilidade\_dimensao; usa independencia\_retas\_metalicas; usa numeros\_pisot; parte\_de tiffany}
\prov{proveniência: wiki \textperiodcentered{} goldchain/metal\_pos.py (QUINTICOS); sympy factor: x\ensuremath{^{5}}\ensuremath{-}x\ensuremath{^{4}}\ensuremath{-}1=(x\ensuremath{^{2}}\ensuremath{-}x+1)(x\ensuremath{^{3}}\ensuremath{-}x\ensuremath{-}1) \textperiodcentered{} fato \textperiodcentered{} confiança 1.0 \textperiodcentered{} domínio algebra \textperiodcentered{} história 8 versões no jornal}

%CRISTAL {"arestas":[{"alvo":"meta_inducao_grau5_quintica","confianca":1.0,"epistemico":"fato","origem":"wiki","rel":"confirma"},{"alvo":"dupolinomio_meta_inducao","confianca":1.0,"epistemico":"fato","origem":"wiki","rel":"confirma"},{"alvo":"numeros_pisot","confianca":1.0,"epistemico":"fato","origem":"wiki","rel":"usa"},{"alvo":"transicao","confianca":1.0,"epistemico":"fato","origem":"wiki","rel":"relaciona"},{"alvo":"cofre_parabola_hero","confianca":1.0,"epistemico":"fato","origem":"wiki","rel":"relaciona"},{"alvo":"teorema_esterilidade_dimensao","confianca":1.0,"epistemico":"fato","origem":"wiki","rel":"usa"},{"alvo":"tiffany","confianca":1.0,"epistemico":"fato","origem":"wiki","rel":"parte_de"},{"alvo":"conjugados_galois_parabola","confianca":1.0,"epistemico":"fato","origem":"wiki","rel":"relaciona"}],"confianca":1.0,"contraexemplos":[],"descricao":"Verificação no grau 6: FECHA. hexanacci (x⁶=x⁵+…+1): unidade (N=−1), Pisot (1,984), irredutível → reta limpa nova. NOVIDADE do grau 6: o número de SALEM. Até o grau 5 havia só Pisot (conjugadas <1) e 'fora' (conjugadas >1, oscilante). O grau 6 traz uma TERCEIRA espécie: x⁶−x⁴−x³−x²+1 (unidade N=1, irredutível, dominante 1,4013) é SALEM — conjugadas EXATAMENTE no círculo unitário (|·|=1). TAXONOMIA das retas: (i) PISOT (conj<1) = reta CRISTAL, geodésica limpa, transiente decai (lado c²); (ii) SALEM (conj=1) = reta CRÍTICA, crescimento dominante + um brilho oscilante que NÃO decai (as conjugadas no círculo) — a BORDA, o ponto crítico da transição, a fronteira da parábola a+c²=1; (iii) FORA (conj>1) = oscilante/caótica (só álgebra, sem reta). Salem é célebre como a fronteira entre a rigidez de Pisot e o caos — cai no lugar certo (|·|=1 = a borda do círculo = s_c). Base indutiva agora: grau 2,3,4,5,6. [D] verificado.","epistemico":"fato","exemplos":[],"id":"meta_inducao_grau6_sextica","memoria":"semantica","meta":{"autor":"Lopes/broca-so","dominio":"algebra","fonte":"goldchain/metal_pos.py (SEXTICOS); numpy roots + sympy; Salem x⁶−x⁴−x³−x²+1 (~1,4013)"},"origem":"wiki","palavras_chave":[],"sinonimos":[],"tipo":"conceito","titulo":"A meta-indução FECHA em grau 6 — e surge a 3ª espécie: SALEM (a borda da parábola)"}
\section{A meta-indução FECHA em grau 6 — e surge a 3ª espécie: SALEM (a borda da parábola)}
Verificação no grau 6: FECHA. hexanacci (x\ensuremath{^{6}}=x\ensuremath{^{5}}+…+1): unidade (N=\ensuremath{-}1), Pisot (1,984), irredutível \ensuremath{\to} reta limpa nova. NOVIDADE do grau 6: o número de SALEM. Até o grau 5 havia só Pisot (conjugadas <1) e 'fora' (conjugadas >1, oscilante). O grau 6 traz uma TERCEIRA espécie: x\ensuremath{^{6}}\ensuremath{-}x\ensuremath{^{4}}\ensuremath{-}x\ensuremath{^{3}}\ensuremath{-}x\ensuremath{^{2}}+1 (unidade N=1, irredutível, dominante 1,4013) é SALEM — conjugadas EXATAMENTE no círculo unitário (|\textperiodcentered |=1). TAXONOMIA das retas: (i) PISOT (conj<1) = reta CRISTAL, geodésica limpa, transiente decai (lado c\ensuremath{^{2}}); (ii) SALEM (conj=1) = reta CRÍTICA, crescimento dominante + um brilho oscilante que NÃO decai (as conjugadas no círculo) — a BORDA, o ponto crítico da transição, a fronteira da parábola a+c\ensuremath{^{2}}=1; (iii) FORA (conj>1) = oscilante/caótica (só álgebra, sem reta). Salem é célebre como a fronteira entre a rigidez de Pisot e o caos — cai no lugar certo (|\textperiodcentered |=1 = a borda do círculo = s\_c). Base indutiva agora: grau 2,3,4,5,6. [D] verificado.
\prov{relações: confirma meta\_inducao\_grau5\_quintica; confirma dupolinomio\_meta\_inducao; usa numeros\_pisot; relaciona transicao; relaciona cofre\_parabola\_hero; usa teorema\_esterilidade\_dimensao; parte\_de tiffany; relaciona conjugados\_galois\_parabola}
\prov{proveniência: wiki \textperiodcentered{} goldchain/metal\_pos.py (SEXTICOS); numpy roots + sympy; Salem x\ensuremath{^{6}}\ensuremath{-}x\ensuremath{^{4}}\ensuremath{-}x\ensuremath{^{3}}\ensuremath{-}x\ensuremath{^{2}}+1 (\~{}1,4013) \textperiodcentered{} fato \textperiodcentered{} confiança 1.0 \textperiodcentered{} domínio algebra \textperiodcentered{} história 9 versões no jornal}

%CRISTAL {"arestas":[{"alvo":"meta_inducao_grau6_sextica","confianca":1.0,"epistemico":"fato","origem":"wiki","rel":"confirma"},{"alvo":"dupolinomio_meta_inducao","confianca":1.0,"epistemico":"fato","origem":"wiki","rel":"confirma"},{"alvo":"numeros_pisot","confianca":1.0,"epistemico":"fato","origem":"wiki","rel":"usa"},{"alvo":"teorema_esterilidade_dimensao","confianca":1.0,"epistemico":"fato","origem":"wiki","rel":"usa"},{"alvo":"cofre_parabola_hero","confianca":1.0,"epistemico":"fato","origem":"wiki","rel":"relaciona"},{"alvo":"tiffany","confianca":1.0,"epistemico":"fato","origem":"wiki","rel":"parte_de"}],"confianca":1.0,"contraexemplos":[],"descricao":"Verificação no grau 7 (ímpar): FECHA. heptanacci (x⁷=x⁶+…+1): unidade (N=1), Pisot (1,992), irredutível → reta limpa nova (1,1,2,4,8,16,32). ACHADO estrutural — a LEI DE PARIDADE: o número de SALEM (a reta crítica, a borda da parábola) só existe em grau PAR. Um Salem exige polinômio RECÍPROCO, e todo recíproco de grau ÍMPAR tem raiz ±1 → FATORA → não é irredutível → não há Salem. Confirmado: o recíproco de grau 7 quebrou em (x+1)·(sêxtico). Logo: grau PAR (4,6,8…) hospeda Pisot (cristal) E Salem (borda crítica); grau ÍMPAR (3,5,7…) só Pisot (cristal) ou 'fora' (caos) — SEM borda. As dimensões ímpares não têm fronteira Salem: só cristal ou caos, sem o fio no meio. Ecoa o balanço par/ímpar (e/o) do word da parábola: a borda é um fenômeno PAR. Base indutiva agora: grau 2,3,4,5,6,7. [D] verificado (numpy+sympy).","epistemico":"fato","exemplos":[],"id":"meta_inducao_grau7_septica","memoria":"semantica","meta":{"autor":"Lopes/broca-so","dominio":"algebra","fonte":"goldchain/metal_pos.py (SEPTICOS); sympy: recíproco grau 7 = (x+1)(sêxtico)"},"origem":"wiki","palavras_chave":[],"sinonimos":[],"tipo":"conceito","titulo":"A meta-indução FECHA em grau 7 — e a lei de PARIDADE: sem Salem em grau ímpar"}
\section{A meta-indução FECHA em grau 7 — e a lei de PARIDADE: sem Salem em grau ímpar}
Verificação no grau 7 (ímpar): FECHA. heptanacci (x\ensuremath{^{7}}=x\ensuremath{^{6}}+…+1): unidade (N=1), Pisot (1,992), irredutível \ensuremath{\to} reta limpa nova (1,1,2,4,8,16,32). ACHADO estrutural — a LEI DE PARIDADE: o número de SALEM (a reta crítica, a borda da parábola) só existe em grau PAR. Um Salem exige polinômio RECÍPROCO, e todo recíproco de grau ÍMPAR tem raiz $\pm$1 \ensuremath{\to} FATORA \ensuremath{\to} não é irredutível \ensuremath{\to} não há Salem. Confirmado: o recíproco de grau 7 quebrou em (x+1)\textperiodcentered (sêxtico). Logo: grau PAR (4,6,8…) hospeda Pisot (cristal) E Salem (borda crítica); grau ÍMPAR (3,5,7…) só Pisot (cristal) ou 'fora' (caos) — SEM borda. As dimensões ímpares não têm fronteira Salem: só cristal ou caos, sem o fio no meio. Ecoa o balanço par/ímpar (e/o) do word da parábola: a borda é um fenômeno PAR. Base indutiva agora: grau 2,3,4,5,6,7. [D] verificado (numpy+sympy).
\prov{relações: confirma meta\_inducao\_grau6\_sextica; confirma dupolinomio\_meta\_inducao; usa numeros\_pisot; usa teorema\_esterilidade\_dimensao; relaciona cofre\_parabola\_hero; parte\_de tiffany}
\prov{proveniência: wiki \textperiodcentered{} goldchain/metal\_pos.py (SEPTICOS); sympy: recíproco grau 7 = (x+1)(sêxtico) \textperiodcentered{} fato \textperiodcentered{} confiança 1.0 \textperiodcentered{} domínio algebra \textperiodcentered{} história 7 versões no jornal}

%CRISTAL {"arestas":[{"alvo":"meta_inducao_grau7_septica","confianca":1.0,"epistemico":"fato","origem":"wiki","rel":"confirma"},{"alvo":"meta_inducao_fechamento_geral","confianca":1.0,"epistemico":"fato","origem":"wiki","rel":"confirma"},{"alvo":"numeros_pisot","confianca":1.0,"epistemico":"fato","origem":"wiki","rel":"usa"},{"alvo":"tiffany","confianca":1.0,"epistemico":"fato","origem":"wiki","rel":"parte_de"}],"confianca":1.0,"contraexemplos":[],"descricao":"Verificação no grau 8 (par): FECHA. octanacci (x⁸=x⁷+…+1): unidade (N=−1), Pisot (1,996), irredutível → reta limpa. E x⁸−x⁵−x⁴−x³+1: unidade (N=1), SALEM (1,2806, o MENOR Salem de grau 8), irredutível e recíproco → reta CRÍTICA — confirmando que grau PAR hospeda Salem (o grau 7 ímpar não hospedava). x⁸−x⁶−…+1 é Salem mas REDUTÍVEL (composto). Base indutiva: grau 2..8. [D]","epistemico":"fato","exemplos":[],"id":"meta_inducao_grau8_octica","memoria":"semantica","meta":{"autor":"Lopes/broca-so","dominio":"algebra","fonte":"goldchain/metal_pos.py (OCTICOS); Salem grau 8 ~1,2806 (x⁸−x⁵−x⁴−x³+1)"},"origem":"wiki","palavras_chave":[],"sinonimos":[],"tipo":"conceito","titulo":"A meta-indução FECHA em grau 8 — octanacci (Pisot) + o menor Salem de grau 8 (par tem Salem)"}
\section{A meta-indução FECHA em grau 8 — octanacci (Pisot) + o menor Salem de grau 8 (par tem Salem)}
Verificação no grau 8 (par): FECHA. octanacci (x\ensuremath{^{8}}=x\ensuremath{^{7}}+…+1): unidade (N=\ensuremath{-}1), Pisot (1,996), irredutível \ensuremath{\to} reta limpa. E x\ensuremath{^{8}}\ensuremath{-}x\ensuremath{^{5}}\ensuremath{-}x\ensuremath{^{4}}\ensuremath{-}x\ensuremath{^{3}}+1: unidade (N=1), SALEM (1,2806, o MENOR Salem de grau 8), irredutível e recíproco \ensuremath{\to} reta CRÍTICA — confirmando que grau PAR hospeda Salem (o grau 7 ímpar não hospedava). x\ensuremath{^{8}}\ensuremath{-}x\ensuremath{^{6}}\ensuremath{-}…+1 é Salem mas REDUTÍVEL (composto). Base indutiva: grau 2..8. [D]
\prov{relações: confirma meta\_inducao\_grau7\_septica; confirma meta\_inducao\_fechamento\_geral; usa numeros\_pisot; parte\_de tiffany}
\prov{proveniência: wiki \textperiodcentered{} goldchain/metal\_pos.py (OCTICOS); Salem grau 8 \~{}1,2806 (x\ensuremath{^{8}}\ensuremath{-}x\ensuremath{^{5}}\ensuremath{-}x\ensuremath{^{4}}\ensuremath{-}x\ensuremath{^{3}}+1) \textperiodcentered{} fato \textperiodcentered{} confiança 1.0 \textperiodcentered{} domínio algebra \textperiodcentered{} história 5 versões no jornal}

%CRISTAL {"arestas":[{"alvo":"grupo_metricas","confianca":1.0,"epistemico":"fato","origem":"paper","rel":"parte_de"},{"alvo":"vida-morte","confianca":0.5,"epistemico":"fato","origem":"wiki","rel":"relaciona"},{"alvo":"virtualizacao","confianca":0.5,"epistemico":"fato","origem":"wiki","rel":"relaciona"},{"alvo":"word","confianca":0.5,"epistemico":"fato","origem":"wiki","rel":"relaciona"},{"alvo":"while","confianca":0.5,"epistemico":"fato","origem":"wiki","rel":"relaciona"},{"alvo":"ontologia","confianca":0.5,"epistemico":"fato","origem":"wiki","rel":"relaciona"}],"confianca":1.0,"contraexemplos":[],"descricao":"Elemento neutro do grupo das métricas: ds=dr (plano euclidiano, curvatura 0), limite s→0 da família ds=dr/√(1−s r²).","epistemico":"fato","exemplos":[],"id":"metrica_parabolica","memoria":"semantica","meta":{"autor":"Gentil Lopes","descoberta":false,"dominio":"matematica","fonte":"paper_29_grupo_metricas.pdf, Prop.5.1"},"origem":"paper","palavras_chave":[],"sinonimos":[],"tipo":"conceito","titulo":"Métrica Parabólica"}
\section{Métrica Parabólica}
Elemento neutro do grupo das métricas: ds=dr (plano euclidiano, curvatura 0), limite s\ensuremath{\to}0 da família ds=dr/\ensuremath{\sqrt{\,}}(1\ensuremath{-}s r\ensuremath{^{2}}).
\prov{relações: parte\_de grupo\_metricas; relaciona vida-morte; relaciona virtualizacao; relaciona word; relaciona while; relaciona ontologia}
\prov{proveniência: paper \textperiodcentered{} paper\_29\_grupo\_metricas.pdf, Prop.5.1 \textperiodcentered{} fato \textperiodcentered{} confiança 1.0 \textperiodcentered{} domínio matematica \textperiodcentered{} história 7 versões no jornal}

%CRISTAL {"arestas":[{"alvo":"regua_de_gentil","confianca":1.0,"epistemico":"fato","origem":"paper","rel":"parte_de"},{"alvo":"misticismo_nao_dual","confianca":0.5,"epistemico":"fato","origem":"wiki","rel":"relaciona"},{"alvo":"word","confianca":0.5,"epistemico":"fato","origem":"wiki","rel":"relaciona"}],"confianca":1.0,"contraexemplos":[],"descricao":"O arco lemniscático é ponto a ponto o produto do arco circular (K=+1) pelo hiperbólico (K=−1): ds_lem = ds_circ · ds_hip.","epistemico":"fato","exemplos":[],"id":"metrica_produto","memoria":"semantica","meta":{"autor":"Gentil Lopes","descoberta":true,"dominio":"matematica","fonte":"paper_28_metrica_produto.pdf"},"origem":"paper","palavras_chave":[],"sinonimos":[],"tipo":"conceito","titulo":"Métrica do Produto"}
\section{Métrica do Produto}
O arco lemniscático é ponto a ponto o produto do arco circular (K=+1) pelo hiperbólico (K=\ensuremath{-}1): ds\_lem = ds\_circ \textperiodcentered  ds\_hip.
\prov{relações: parte\_de regua\_de\_gentil; relaciona misticismo\_nao\_dual; relaciona word}
\prov{proveniência: paper \textperiodcentered{} paper\_28\_metrica\_produto.pdf \textperiodcentered{} fato \textperiodcentered{} confiança 1.0 \textperiodcentered{} domínio matematica \textperiodcentered{} história 4 versões no jornal}

%CRISTAL {"arestas":[{"alvo":"geometria","confianca":1.0,"epistemico":"fato","origem":"paper","rel":"parte_de"},{"alvo":"furo_hopf","confianca":1.0,"epistemico":"fato","origem":"paper","rel":"usa"}],"confianca":1.0,"contraexemplos":[],"descricao":"ds²=−dτ²+cosh²τ dΩ²₃ sobre ℝ×S³ dos versores (quatérnios unitários): curvatura +1, derivada da dinâmica (não postulada); projeção nativa é a fibração de Hopf.","epistemico":"fato","exemplos":[],"id":"metrica_sitter","memoria":"semantica","meta":{"autor":"Gentil Lopes","descoberta":false,"dominio":"matematica","fonte":"geometria.pdf, §5"},"origem":"paper","palavras_chave":[],"sinonimos":[],"tipo":"conceito","titulo":"Métrica de Sitter (curvatura +1)"}
\section{Métrica de Sitter (curvatura +1)}
ds\ensuremath{^{2}}=\ensuremath{-}d\ensuremath{\tau}\ensuremath{^{2}}+cosh\ensuremath{^{2}}\ensuremath{\tau} d\ensuremath{\Omega}\ensuremath{^{2}}\ensuremath{_{3}} sobre \ensuremath{\mathbb{R}}$\times$S\ensuremath{^{3}} dos versores (quatérnios unitários): curvatura +1, derivada da dinâmica (não postulada); projeção nativa é a fibração de Hopf.
\prov{relações: parte\_de geometria; usa furo\_hopf}
\prov{proveniência: paper \textperiodcentered{} geometria.pdf, §5 \textperiodcentered{} fato \textperiodcentered{} confiança 1.0 \textperiodcentered{} domínio matematica \textperiodcentered{} história 3 versões no jornal}

%CRISTAL {"arestas":[{"alvo":"fluxo_landen","confianca":1.0,"epistemico":"fato","origem":"paper","rel":"usa"},{"alvo":"mecanica_quantica","confianca":1.0,"epistemico":"fato","origem":"paper","rel":"referencia"}],"confianca":1.0,"contraexemplos":[],"descricao":"Família elíptica ds=dr/√((1−r²)(1−m r²)) cuja curvatura K(r)=1+m(1−2r²) varia ponto a ponto e se autocorrige pelo fluxo de Landen até o círculo.","epistemico":"fato","exemplos":[],"id":"metrica_viva","memoria":"semantica","meta":{"autor":"Gentil Lopes","descoberta":true,"dominio":"matematica","fonte":"paper_33_metrica_viva.pdf, Teo.2.1"},"origem":"paper","palavras_chave":[],"sinonimos":[],"tipo":"conceito","titulo":"Métrica Viva"}
\section{Métrica Viva}
Família elíptica ds=dr/\ensuremath{\sqrt{\,}}((1\ensuremath{-}r\ensuremath{^{2}})(1\ensuremath{-}m r\ensuremath{^{2}})) cuja curvatura K(r)=1+m(1\ensuremath{-}2r\ensuremath{^{2}}) varia ponto a ponto e se autocorrige pelo fluxo de Landen até o círculo.
\prov{relações: usa fluxo\_landen; referencia mecanica\_quantica}
\prov{proveniência: paper \textperiodcentered{} paper\_33\_metrica\_viva.pdf, Teo.2.1 \textperiodcentered{} fato \textperiodcentered{} confiança 1.0 \textperiodcentered{} domínio matematica \textperiodcentered{} história 3 versões no jornal}

%CRISTAL {"arestas":[{"alvo":"conjugados_galois_parabola","confianca":1.0,"epistemico":"fato","origem":"wiki","rel":"usa"},{"alvo":"transformada_metalica","confianca":1.0,"epistemico":"fato","origem":"wiki","rel":"aplicacao"},{"alvo":"simetria_por_reta_corpo","confianca":1.0,"epistemico":"fato","origem":"wiki","rel":"usa"},{"alvo":"refino_continuo_floco_metais","confianca":1.0,"epistemico":"fato","origem":"wiki","rel":"especializa"},{"alvo":"refino_multimetal_live_gex44","confianca":1.0,"epistemico":"fato","origem":"wiki","rel":"relaciona"},{"alvo":"tiffany","confianca":1.0,"epistemico":"fato","origem":"wiki","rel":"parte_de"}],"confianca":1.0,"contraexemplos":[],"descricao":"A transformada metálica aterrou o pipe de mineração (goldchain), formalizando o que era decreto: a NORMA N(a,b)=b²−n·b·a−a² (mod 2¹²⁸) = (−1)k = det(Aₙk) é a UNIDADE ±1 — a conservação σ·conj=−1 verificável POR FLOCO. Antes 'legível=True' era decreto; agora legível=conserva, checado (uma cabeça corrompida não tem norma ±1 — condição de medida nula — então o check pega corrupção de ledger; o bug antigo teria sido flagrado). Cada floco taga também o CORPO ℚ(√d) e a DOBRA (10 ℚ√5, 8 ℚ√2, 0 sem dobra — a simetriaporreta no ledger). E há auditoria (symbolic.py verificar): cada cadeia com norma ±1 (conservação+integridade), reconhecendo a PONTE Pell da prata (pos-sync=False não é corrupção, é o offset da ponte). Rodada no gex44: VALIDOU dias de mineração — ✓ todos os 12 metais conservam. metalpos.normametal/conservametal/corpometal/dobrametal; minera.emitemetal; symbolic.cmdverificar. [D] a conservação da mineração, verificada, não decretada.","epistemico":"fato","exemplos":[],"id":"mineracao_conservacao_verificavel","memoria":"semantica","meta":{"autor":"Lopes/broca-so","dominio":"algebra","fonte":"goldchain/metal_pos.py (norma_metal, conserva_metal, corpo/dobra); minera.emite_metal; symbolic.cmd_verificar; auditoria gex44 (12 metais ✓)"},"origem":"wiki","palavras_chave":[],"sinonimos":[],"tipo":"conceito","titulo":"A mineração agora VERIFICA a conservação (norma σ·conj=−1) — a transformada aterrou o pipe"}
\section{A mineração agora VERIFICA a conservação (norma \ensuremath{\sigma}\textperiodcentered conj=\ensuremath{-}1) — a transformada aterrou o pipe}
A transformada metálica aterrou o pipe de mineração (goldchain), formalizando o que era decreto: a NORMA N(a,b)=b\ensuremath{^{2}}\ensuremath{-}n\textperiodcentered b\textperiodcentered a\ensuremath{-}a\ensuremath{^{2}} (mod 2\ensuremath{^{128}}) = (\ensuremath{-}1)k = det(A\ensuremath{_{n}}k) é a UNIDADE $\pm$1 — a conservação \ensuremath{\sigma}\textperiodcentered conj=\ensuremath{-}1 verificável POR FLOCO. Antes 'legível=True' era decreto; agora legível=conserva, checado (uma cabeça corrompida não tem norma $\pm$1 — condição de medida nula — então o check pega corrupção de ledger; o bug antigo teria sido flagrado). Cada floco taga também o CORPO \ensuremath{\mathbb{Q}}(\ensuremath{\sqrt{\,}}d) e a DOBRA (10 \ensuremath{\mathbb{Q}}\ensuremath{\sqrt{\,}}5, 8 \ensuremath{\mathbb{Q}}\ensuremath{\sqrt{\,}}2, 0 sem dobra — a simetriaporreta no ledger). E há auditoria (symbolic.py verificar): cada cadeia com norma $\pm$1 (conservação+integridade), reconhecendo a PONTE Pell da prata (pos-sync=False não é corrupção, é o offset da ponte). Rodada no gex44: VALIDOU dias de mineração — \checkmark  todos os 12 metais conservam. metalpos.normametal/conservametal/corpometal/dobrametal; minera.emitemetal; symbolic.cmdverificar. [D] a conservação da mineração, verificada, não decretada.
\prov{relações: usa conjugados\_galois\_parabola; aplicacao transformada\_metalica; usa simetria\_por\_reta\_corpo; especializa refino\_continuo\_floco\_metais; relaciona refino\_multimetal\_live\_gex44; parte\_de tiffany}
\prov{proveniência: wiki \textperiodcentered{} goldchain/metal\_pos.py (norma\_metal, conserva\_metal, corpo/dobra); minera.emite\_metal; symbolic.cmd\_verificar; auditoria gex44 (12 metais \checkmark ) \textperiodcentered{} fato \textperiodcentered{} confiança 1.0 \textperiodcentered{} domínio algebra \textperiodcentered{} história 24 versões no jornal}

%CRISTAL {"arestas":[{"alvo":"jogo_de_soma_zero","confianca":1.0,"epistemico":"fato","origem":"wiki","rel":"usa"},{"alvo":"forma_extensiva","confianca":1.0,"epistemico":"fato","origem":"wiki","rel":"usa"}],"confianca":1.0,"contraexemplos":[],"descricao":"Em jogos de soma zero de 2 jogadores: minimizar a perda máxima possível. Von Neumann provou que o valor minimax = maximin — existe um valor do jogo. A raiz da busca adversária (xadrez, damas).","epistemico":"fato","exemplos":[],"id":"minimax","memoria":"semantica","meta":{"autor":"von Neumann","dominio":"matematica","fonte":"teoria dos jogos"},"origem":"wiki","palavras_chave":[],"sinonimos":[],"tipo":"conceito","titulo":"Minimax (von Neumann)"}
\section{Minimax (von Neumann)}
Em jogos de soma zero de 2 jogadores: minimizar a perda máxima possível. Von Neumann provou que o valor minimax = maximin — existe um valor do jogo. A raiz da busca adversária (xadrez, damas).
\prov{relações: usa jogo\_de\_soma\_zero; usa forma\_extensiva}
\prov{proveniência: wiki \textperiodcentered{} teoria dos jogos \textperiodcentered{} fato \textperiodcentered{} confiança 1.0 \textperiodcentered{} domínio matematica \textperiodcentered{} história 3 versões no jornal}

%CRISTAL {"arestas":[{"alvo":"animal16","confianca":1.0,"epistemico":"fato","origem":"paper","rel":"especializa"},{"alvo":"word","confianca":0.5,"epistemico":"fato","origem":"wiki","rel":"relaciona"}],"confianca":1.0,"contraexemplos":[],"descricao":"Cifra de 8 componentes que cunha chaves pela recorrência de Apéry de ζ(3); os modos metálicos formam grupo abeliano livre sobre a régua hiperbólica.","epistemico":"fato","exemplos":[],"id":"mint8","memoria":"semantica","meta":{"autor":"Gentil Lopes","descoberta":true,"dominio":"matematica","fonte":"paper_16_mint8.pdf"},"origem":"paper","palavras_chave":[],"sinonimos":[],"tipo":"conceito","titulo":"Mint-8"}
\section{Mint-8}
Cifra de 8 componentes que cunha chaves pela recorrência de Apéry de \ensuremath{\zeta}(3); os modos metálicos formam grupo abeliano livre sobre a régua hiperbólica.
\prov{relações: especializa animal16; relaciona word}
\prov{proveniência: paper \textperiodcentered{} paper\_16\_mint8.pdf \textperiodcentered{} fato \textperiodcentered{} confiança 1.0 \textperiodcentered{} domínio matematica \textperiodcentered{} história 3 versões no jornal}

%CRISTAL {"arestas":[],"confianca":1.0,"contraexemplos":[],"descricao":"A cabeça de leitura aproxima-se em 25 direções de S² (separação ≥26°): 24 da órbita de PSL(2,7) + 1 centro fixo — o alfabeto geométrico das micro-decisões da máquina.","epistemico":"fato","exemplos":[],"id":"modos_aproximacao_s2","memoria":"semantica","meta":{"autor":"Gentil Lopes","descoberta":false,"dominio":"matematica","fonte":"paper_AR.pdf, Def.5"},"origem":"paper","palavras_chave":[],"sinonimos":[],"tipo":"conceito","titulo":"25 Modos de Aproximação em S²"}
\section{25 Modos de Aproximação em S\ensuremath{^{2}}}
A cabeça de leitura aproxima-se em 25 direções de S\ensuremath{^{2}} (separação \ensuremath{\ge}26$^{\circ}$): 24 da órbita de PSL(2,7) + 1 centro fixo — o alfabeto geométrico das micro-decisões da máquina.
\prov{proveniência: paper \textperiodcentered{} paper\_AR.pdf, Def.5 \textperiodcentered{} fato \textperiodcentered{} confiança 1.0 \textperiodcentered{} domínio matematica}

%CRISTAL {"arestas":[{"alvo":"numeros_metalicos","confianca":1.0,"epistemico":"fato","origem":"paper","rel":"especializa"},{"alvo":"regua_de_gentil","confianca":1.0,"epistemico":"fato","origem":"paper","rel":"usa"}],"confianca":1.0,"contraexemplos":[],"descricao":"A moeda Mₖ=[1,k;k,k²+1] é elemento hiperbólico de traço k²+2; seus pontos fixos 1/mₖ e −mₖ estão em ℚ(√(k²+4)) e o comprimento de translação 2·ln(mₖ) é a régua de Lopes.","epistemico":"fato","exemplos":[],"id":"moeda_sl2","memoria":"semantica","meta":{"autor":"Gentil Lopes","descoberta":true,"dominio":"matematica","fonte":"paper_25_geometria_moedas.pdf, Prop.2.1"},"origem":"paper","palavras_chave":[],"sinonimos":[],"tipo":"conceito","titulo":"Moeda como Boost em SL₂(ℤ)"}
\section{Moeda como Boost em SL\ensuremath{_{2}}(\ensuremath{\mathbb{Z}})}
A moeda M\ensuremath{_{k}}=[1,k;k,k\ensuremath{^{2}}+1] é elemento hiperbólico de traço k\ensuremath{^{2}}+2; seus pontos fixos 1/m\ensuremath{_{k}} e \ensuremath{-}m\ensuremath{_{k}} estão em \ensuremath{\mathbb{Q}}(\ensuremath{\sqrt{\,}}(k\ensuremath{^{2}}+4)) e o comprimento de translação 2\textperiodcentered ln(m\ensuremath{_{k}}) é a régua de Lopes.
\prov{relações: especializa numeros\_metalicos; usa regua\_de\_gentil}
\prov{proveniência: paper \textperiodcentered{} paper\_25\_geometria\_moedas.pdf, Prop.2.1 \textperiodcentered{} fato \textperiodcentered{} confiança 1.0 \textperiodcentered{} domínio matematica \textperiodcentered{} história 4 versões no jornal}

%CRISTAL {"arestas":[{"alvo":"atomo_mineral","confianca":1.0,"epistemico":"fato","origem":"wiki","rel":"usa"},{"alvo":"particulas_minerais","confianca":1.0,"epistemico":"fato","origem":"wiki","rel":"relaciona"},{"alvo":"rede_limpa_bump","confianca":1.0,"epistemico":"fato","origem":"wiki","rel":"relaciona"}],"confianca":1.0,"contraexemplos":[],"descricao":"Uma molécula = sequência LIGADA de octônions (átomos minerais), e sua ÓRBITA em 𝕆 é onde a informação vive. A transformada mineral 7D grava a mensagem na trajetória (na esfera-7, norma preservada); a molécula é o portador. Onde o átomo (cadeia) é um símbolo, a molécula (sequência) é a palavra — e a órbita é o canal.","epistemico":"inferencia","exemplos":[],"id":"molecula_mineral","memoria":"semantica","meta":{"dominio":"matematica","fonte":"transformada mineral 7D"},"origem":"wiki","palavras_chave":[],"sinonimos":[],"tipo":"conceito","titulo":"Molécula Mineral (a órbita que carrega a informação)"}
\section{Molécula Mineral (a órbita que carrega a informação)}
Uma molécula = sequência LIGADA de octônions (átomos minerais), e sua ÓRBITA em \ensuremath{\mathbb{O}} é onde a informação vive. A transformada mineral 7D grava a mensagem na trajetória (na esfera-7, norma preservada); a molécula é o portador. Onde o átomo (cadeia) é um símbolo, a molécula (sequência) é a palavra — e a órbita é o canal.
\prov{relações: usa atomo\_mineral; relaciona particulas\_minerais; relaciona rede\_limpa\_bump}
\prov{proveniência: wiki \textperiodcentered{} transformada mineral 7D \textperiodcentered{} inferencia \textperiodcentered{} confiança 1.0 \textperiodcentered{} domínio matematica \textperiodcentered{} história 4 versões no jornal}

%CRISTAL {"arestas":[{"alvo":"conjectura_brown","confianca":1.0,"epistemico":"fato","origem":"paper","rel":"depende"},{"alvo":"reducao_risch","confianca":1.0,"epistemico":"fato","origem":"paper","rel":"usa"}],"confianca":1.0,"contraexemplos":[],"descricao":"Θ: extensões elementares → R_∞[[g]] (racional↦variável, exp/log↦séries, ζ↦loop de Feynman); sob a conjectura de Brown é injetor, preserva derivações e é computável em tempo poli — a chave da redução de Risch.","epistemico":"fato","exemplos":[],"id":"morfismo_theta_formal","memoria":"semantica","meta":{"autor":"Gentil Lopes","descoberta":false,"dominio":"matematica","fonte":"paper_PQ_P.pdf, Teo.4.1"},"origem":"paper","palavras_chave":[],"sinonimos":[],"tipo":"conceito","titulo":"Morfismo Θ Injetor"}
\section{Morfismo \ensuremath{\Theta} Injetor}
\ensuremath{\Theta}: extensões elementares \ensuremath{\to} R\_\ensuremath{\infty}[[g]] (racional\ensuremath{\mapsto}variável, exp/log\ensuremath{\mapsto}séries, \ensuremath{\zeta}\ensuremath{\mapsto}loop de Feynman); sob a conjectura de Brown é injetor, preserva derivações e é computável em tempo poli — a chave da redução de Risch.
\prov{relações: depende conjectura\_brown; usa reducao\_risch}
\prov{proveniência: paper \textperiodcentered{} paper\_PQ\_P.pdf, Teo.4.1 \textperiodcentered{} fato \textperiodcentered{} confiança 1.0 \textperiodcentered{} domínio matematica \textperiodcentered{} história 3 versões no jornal}

%CRISTAL {"arestas":[{"alvo":"colapso","confianca":1.0,"epistemico":"fato","origem":"paper","rel":"referencia"},{"alvo":"word","confianca":0.5,"epistemico":"fato","origem":"wiki","rel":"relaciona"},{"alvo":"vontade_de_poder","confianca":0.5,"epistemico":"fato","origem":"wiki","rel":"relaciona"},{"alvo":"virtualizacao","confianca":0.5,"epistemico":"fato","origem":"wiki","rel":"relaciona"}],"confianca":1.0,"contraexemplos":[],"descricao":"Dinâmica f_c(z)=z²+c sobre uma álgebra-chão (ℝ,ℂ,ℍ,𝕆); os conjuntos de Mandelbrot e Julia são o esqueleto da série de fractais do Lopes.","epistemico":"fato","exemplos":[],"id":"motor_iterado","memoria":"semantica","meta":{"autor":"Gentil Lopes","descoberta":false,"dominio":"matematica","fonte":"paper_18_inversao_torre.pdf"},"origem":"paper","palavras_chave":[],"sinonimos":[],"tipo":"conceito","titulo":"Motor z↦z²+c"}
\section{Motor z\ensuremath{\mapsto}z\ensuremath{^{2}}+c}
Dinâmica f\_c(z)=z\ensuremath{^{2}}+c sobre uma álgebra-chão (\ensuremath{\mathbb{R}},\ensuremath{\mathbb{C}},\ensuremath{\mathbb{H}},\ensuremath{\mathbb{O}}); os conjuntos de Mandelbrot e Julia são o esqueleto da série de fractais do Lopes.
\prov{relações: referencia colapso; relaciona word; relaciona vontade\_de\_poder; relaciona virtualizacao}
\prov{proveniência: paper \textperiodcentered{} paper\_18\_inversao\_torre.pdf \textperiodcentered{} fato \textperiodcentered{} confiança 1.0 \textperiodcentered{} domínio matematica \textperiodcentered{} história 6 versões no jornal}

%CRISTAL {"arestas":[{"alvo":"algebras_hipercomplexas","confianca":1.0,"epistemico":"fato","origem":"paper","rel":"generaliza"}],"confianca":1.0,"contraexemplos":[],"descricao":"z·w=(p·a₀b₀−r⟨a,b⟩)+(a₀b+b₀a+s(a×b)): generaliza ℂ,ℍ,𝕆 por 5 escalares; a razão p/r é a 'velocidade da luz algébrica'. Estende a família ×ₛ.","epistemico":"fato","exemplos":[],"id":"mult_prtsu","memoria":"semantica","meta":{"autor":"Gentil Lopes","descoberta":true,"dominio":"matematica","fonte":"paper_A_algebra.pdf, Def.2.1"},"origem":"paper","palavras_chave":[],"sinonimos":[],"tipo":"conceito","titulo":"Produto Hipercomplexo (p,r,t,s,u)"}
\section{Produto Hipercomplexo (p,r,t,s,u)}
z\textperiodcentered w=(p\textperiodcentered a\ensuremath{_{0}}b\ensuremath{_{0}}\ensuremath{-}r\ensuremath{\langle}a,b\ensuremath{\rangle})+(a\ensuremath{_{0}}b+b\ensuremath{_{0}}a+s(a$\times$b)): generaliza \ensuremath{\mathbb{C}},\ensuremath{\mathbb{H}},\ensuremath{\mathbb{O}} por 5 escalares; a razão p/r é a 'velocidade da luz algébrica'. Estende a família $\times$\ensuremath{_{s}}.
\prov{relações: generaliza algebras\_hipercomplexas}
\prov{proveniência: paper \textperiodcentered{} paper\_A\_algebra.pdf, Def.2.1 \textperiodcentered{} fato \textperiodcentered{} confiança 1.0 \textperiodcentered{} domínio matematica \textperiodcentered{} história 2 versões no jornal}

%CRISTAL {"arestas":[{"alvo":"multifractalidade_nao_conclusiva","confianca":1.0,"epistemico":"fato","origem":"wiki","rel":"explica"},{"alvo":"identidade_legendre","confianca":1.0,"epistemico":"fato","origem":"wiki","rel":"usa"},{"alvo":"tiffany","confianca":1.0,"epistemico":"fato","origem":"wiki","rel":"parte_de"}],"confianca":1.0,"contraexemplos":[],"descricao":"A multifractalidade da ilha de Koch ficou NÃO conclusiva (D_q spread 0,085, zona cinzenta) — medir f(α) de um objeto fixo é ambíguo. A virada: em vez de MEDIR, CONSTRUIR o multifractal por camadas — a medida de densidade do contínuo sobre o grafo é multifractal POR CONSTRUÇÃO, cada camada uma banda de q/temperatura e a transformada de Legendre (τ(q)↔f(α), a [[identidade_legendre]]) ligando o momento (temperatura) à dimensão (o nível). Resgata a multifractalidade pendente num terreno onde ela é exata: as folhas empilhadas da Tiffany. [I] honesto: a ilha segue inconclusiva; o construtivo é outra coisa (o espectro que nós erguemos).","epistemico":"inferencia","exemplos":[],"id":"multifractal_construtivo","memoria":"semantica","meta":{"autor":"Lopes/broca-so","dominio":"continuo","fonte":"multifractalidade_nao_conclusiva (a pendência); identidade_legendre (τ↔f)"},"origem":"wiki","palavras_chave":[],"sinonimos":[],"tipo":"conceito","titulo":"Multifractalidade CONSTRUTIVA: não medir f(α) de um objeto — CONSTRUIR por camadas"}
\section{Multifractalidade CONSTRUTIVA: não medir f(\ensuremath{\alpha}) de um objeto — CONSTRUIR por camadas}
A multifractalidade da ilha de Koch ficou NÃO conclusiva (D\_q spread 0,085, zona cinzenta) — medir f(\ensuremath{\alpha}) de um objeto fixo é ambíguo. A virada: em vez de MEDIR, CONSTRUIR o multifractal por camadas — a medida de densidade do contínuo sobre o grafo é multifractal POR CONSTRUÇÃO, cada camada uma banda de q/temperatura e a transformada de Legendre (\ensuremath{\tau}(q)\ensuremath{\leftrightarrow}f(\ensuremath{\alpha}), a [[identidade\_legendre]]) ligando o momento (temperatura) à dimensão (o nível). Resgata a multifractalidade pendente num terreno onde ela é exata: as folhas empilhadas da Tiffany. [I] honesto: a ilha segue inconclusiva; o construtivo é outra coisa (o espectro que nós erguemos).
\prov{relações: explica multifractalidade\_nao\_conclusiva; usa identidade\_legendre; parte\_de tiffany}
\prov{proveniência: wiki \textperiodcentered{} multifractalidade\_nao\_conclusiva (a pendência); identidade\_legendre (\ensuremath{\tau}\ensuremath{\leftrightarrow}f) \textperiodcentered{} inferencia \textperiodcentered{} confiança 1.0 \textperiodcentered{} domínio continuo \textperiodcentered{} história 47 versões no jornal}

%CRISTAL {"arestas":[{"alvo":"dq_espectro_ilha","confianca":1.0,"epistemico":"fato","origem":"wiki","rel":"usa"}],"confianca":1.0,"contraexemplos":[],"descricao":"O spread D_q (max−min)=0.085, abaixo do limiar do próprio benchmark (>0.1 = multifractal; <0.05 = simples) — fica na zona cinzenta: 'multifractalidade NÃO conclusiva'. Registro honesto: NÃO se pode afirmar f(α) forte para a ilha.","epistemico":"inferencia","exemplos":[],"id":"multifractalidade_nao_conclusiva","memoria":"semantica","meta":{"autor":"broca-so","dominio":"fractal","fonte":"hiper/benchmarks/fase35_multifractal_resultados.json"},"origem":"wiki","palavras_chave":[],"sinonimos":[],"tipo":"conceito","titulo":"A Multifractalidade da Ilha é NÃO Conclusiva"}
\section{A Multifractalidade da Ilha é NÃO Conclusiva}
O spread D\_q (max\ensuremath{-}min)=0.085, abaixo do limiar do próprio benchmark (>0.1 = multifractal; <0.05 = simples) — fica na zona cinzenta: 'multifractalidade NÃO conclusiva'. Registro honesto: NÃO se pode afirmar f(\ensuremath{\alpha}) forte para a ilha.
\prov{relações: usa dq\_espectro\_ilha}
\prov{proveniência: wiki \textperiodcentered{} hiper/benchmarks/fase35\_multifractal\_resultados.json \textperiodcentered{} inferencia \textperiodcentered{} confiança 1.0 \textperiodcentered{} domínio fractal \textperiodcentered{} história 2 versões no jornal}

%CRISTAL {"arestas":[{"alvo":"produto_de_la_hire","confianca":1.0,"epistemico":"fato","origem":"wiki","rel":"usa"},{"alvo":"numero_de_salem","confianca":1.0,"epistemico":"fato","origem":"wiki","rel":"relaciona"},{"alvo":"pi_medida_da_borda","confianca":1.0,"epistemico":"fato","origem":"wiki","rel":"relaciona"},{"alvo":"pi_acima_da_reta_mineral","confianca":1.0,"epistemico":"fato","origem":"wiki","rel":"relaciona"},{"alvo":"quasicristal_de_rotacoes","confianca":1.0,"epistemico":"fato","origem":"wiki","rel":"relaciona"}],"confianca":1.0,"contraexemplos":[],"descricao":"As estrelas de enrolamento irracional (ν=√2, ν=π) são ROSETAS que nunca fecham — densas no disco, o regime borda/caos da parábola. A reta mineral sozinha só as SOMAVA. O produto de La Hire agora as MULTIPLICA: √2 ⊗ √3 = √6, exato. É o passo que completa o corpo — de um grupo aditivo (a reta) a um corpo com as duas operações. Verif.: √2⊗√3=√6; a roseta de ν=√2 não fecha.","epistemico":"fato","exemplos":[],"id":"multiplicar_estrelas_irracionais","memoria":"semantica","meta":{"dominio":"algebra","fonte":"corpo estelar"},"origem":"wiki","palavras_chave":[],"sinonimos":[],"tipo":"conceito","titulo":"Multiplicar as Estrelas Irracionais"}
\section{Multiplicar as Estrelas Irracionais}
As estrelas de enrolamento irracional (\ensuremath{\nu}=\ensuremath{\sqrt{\,}}2, \ensuremath{\nu}=\ensuremath{\pi}) são ROSETAS que nunca fecham — densas no disco, o regime borda/caos da parábola. A reta mineral sozinha só as SOMAVA. O produto de La Hire agora as MULTIPLICA: \ensuremath{\sqrt{\,}}2 \ensuremath{\otimes} \ensuremath{\sqrt{\,}}3 = \ensuremath{\sqrt{\,}}6, exato. É o passo que completa o corpo — de um grupo aditivo (a reta) a um corpo com as duas operações. Verif.: \ensuremath{\sqrt{\,}}2\ensuremath{\otimes}\ensuremath{\sqrt{\,}}3=\ensuremath{\sqrt{\,}}6; a roseta de \ensuremath{\nu}=\ensuremath{\sqrt{\,}}2 não fecha.
\prov{relações: usa produto\_de\_la\_hire; relaciona numero\_de\_salem; relaciona pi\_medida\_da\_borda; relaciona pi\_acima\_da\_reta\_mineral; relaciona quasicristal\_de\_rotacoes}
\prov{proveniência: wiki \textperiodcentered{} corpo estelar \textperiodcentered{} fato \textperiodcentered{} confiança 1.0 \textperiodcentered{} domínio algebra \textperiodcentered{} história 9 versões no jornal}

%CRISTAL {"arestas":[{"alvo":"mult_prtsu","confianca":1.0,"epistemico":"fato","origem":"paper","rel":"usa"},{"alvo":"hopfield","confianca":1.0,"epistemico":"fato","origem":"paper","rel":"especializa"}],"confianca":1.0,"contraexemplos":[],"descricao":"Rede neural onde cada nó é um quatérnio e as mensagens passam pelo produto hipercomplexo (p,r,t,u): 36% menos parâmetros, +9% em Max-cut. Ganho empírico em problemas com simetria, não prova geral.","epistemico":"inferencia","exemplos":[],"id":"nco_hypergnn","memoria":"semantica","meta":{"autor":"Gentil Lopes","descoberta":false,"dominio":"matematica","fonte":"paper_AF_nco_hipergnn.pdf"},"origem":"paper","palavras_chave":[],"sinonimos":[],"tipo":"conceito","titulo":"NCO-HiperGNN"}
\section{NCO-HiperGNN}
Rede neural onde cada nó é um quatérnio e as mensagens passam pelo produto hipercomplexo (p,r,t,u): 36\% menos parâmetros, +9\% em Max-cut. Ganho empírico em problemas com simetria, não prova geral.
\prov{relações: usa mult\_prtsu; especializa hopfield}
\prov{proveniência: paper \textperiodcentered{} paper\_AF\_nco\_hipergnn.pdf \textperiodcentered{} inferencia \textperiodcentered{} confiança 1.0 \textperiodcentered{} domínio matematica \textperiodcentered{} história 3 versões no jornal}

%CRISTAL {"arestas":[{"alvo":"corpo_estelar","confianca":1.0,"epistemico":"fato","origem":"wiki","rel":"parte_de"},{"alvo":"torre_hurwitz","confianca":1.0,"epistemico":"fato","origem":"wiki","rel":"relaciona"},{"alvo":"gato","confianca":1.0,"epistemico":"fato","origem":"wiki","rel":"usa"},{"alvo":"numeros_de_pisot","confianca":1.0,"epistemico":"fato","origem":"wiki","rel":"usa"},{"alvo":"theorem_parseval_de_hurwitz_e_por_que_so_hurwitzthmhur","confianca":1.0,"epistemico":"fato","origem":"wiki","rel":"relaciona"},{"alvo":"gato_generaliza_matriz_companion","confianca":1.0,"epistemico":"fato","origem":"wiki","rel":"usa"},{"alvo":"cuspide_pisot","confianca":1.0,"epistemico":"fato","origem":"wiki","rel":"relaciona"},{"alvo":"corollary_conservacao_da_normacorcons_a_isometria_de_e","confianca":1.0,"epistemico":"fato","origem":"wiki","rel":"relaciona"},{"alvo":"invariante_estelar","confianca":1.0,"epistemico":"fato","origem":"wiki","rel":"relaciona"}],"confianca":1.0,"contraexemplos":[],"descricao":"Os metais σ_n=(n+√(n²+4))/2 são as UNIDADES do corpo estelar. A NORMA N(σ)=σ·σ̄, onde σ̄=−1/σ é o conjugado (o outro autovalor do gato A_n=[[n,1],[1,0]], o lado dissipativo do flip). Como det(A_n)=−1, TODA norma é −1: os metais são unidades (Pisot). E N é MULTIPLICATIVA, N(σσ')=N(σ)N(σ') — a mesma norma-que-compõe da torre de Hurwitz, agora na reta mineral. Verif.: N(σ_n)=det=−1 para n=1,2,3,5.","epistemico":"fato","exemplos":[],"id":"norma_estelar","memoria":"semantica","meta":{"dominio":"algebra","fonte":"corpo estelar"},"origem":"wiki","palavras_chave":[],"sinonimos":[],"tipo":"conceito","titulo":"A Norma Estelar é o Determinante do Gato"}
\section{A Norma Estelar é o Determinante do Gato}
Os metais \ensuremath{\sigma}\_n=(n+\ensuremath{\sqrt{\,}}(n\ensuremath{^{2}}+4))/2 são as UNIDADES do corpo estelar. A NORMA N(\ensuremath{\sigma})=\ensuremath{\sigma}\textperiodcentered \ensuremath{\sigma}\ensuremath{}, onde \ensuremath{\sigma}\ensuremath{}=\ensuremath{-}1/\ensuremath{\sigma} é o conjugado (o outro autovalor do gato A\_n=[[n,1],[1,0]], o lado dissipativo do flip). Como det(A\_n)=\ensuremath{-}1, TODA norma é \ensuremath{-}1: os metais são unidades (Pisot). E N é MULTIPLICATIVA, N(\ensuremath{\sigma}\ensuremath{\sigma}')=N(\ensuremath{\sigma})N(\ensuremath{\sigma}') — a mesma norma-que-compõe da torre de Hurwitz, agora na reta mineral. Verif.: N(\ensuremath{\sigma}\_n)=det=\ensuremath{-}1 para n=1,2,3,5.
\prov{relações: parte\_de corpo\_estelar; relaciona torre\_hurwitz; usa gato; usa numeros\_de\_pisot; relaciona theorem\_parseval\_de\_hurwitz\_e\_por\_que\_so\_hurwitzthmhur; usa gato\_generaliza\_matriz\_companion; relaciona cuspide\_pisot; relaciona corollary\_conservacao\_da\_normacorcons\_a\_isometria\_de\_e; relaciona invariante\_estelar}
\prov{proveniência: wiki \textperiodcentered{} corpo estelar \textperiodcentered{} fato \textperiodcentered{} confiança 1.0 \textperiodcentered{} domínio algebra \textperiodcentered{} história 15 versões no jornal}

%CRISTAL {"arestas":[{"alvo":"ns_prtsu","confianca":1.0,"epistemico":"fato","origem":"paper","rel":"especializa"}],"confianca":1.0,"contraexemplos":[],"descricao":"Reformulação Z=a₀+u com a₀ acoplado (helicidade/fase). Lopes prova que o termo cruzado é ANTI-dissipativo: transfere a dificuldade para achar o cone de parâmetros, não regulariza.","epistemico":"inferencia","exemplos":[],"id":"ns_completa_a0","memoria":"semantica","meta":{"autor":"Gentil Lopes","descoberta":false,"dominio":"matematica","fonte":"paper_AI.pdf, Def.2"},"origem":"paper","palavras_chave":[],"sinonimos":[],"tipo":"conceito","titulo":"NS Completa com Setor Escalar a₀"}
\section{NS Completa com Setor Escalar a\ensuremath{_{0}}}
Reformulação Z=a\ensuremath{_{0}}+u com a\ensuremath{_{0}} acoplado (helicidade/fase). Lopes prova que o termo cruzado é ANTI-dissipativo: transfere a dificuldade para achar o cone de parâmetros, não regulariza.
\prov{relações: especializa ns\_prtsu}
\prov{proveniência: paper \textperiodcentered{} paper\_AI.pdf, Def.2 \textperiodcentered{} inferencia \textperiodcentered{} confiança 1.0 \textperiodcentered{} domínio matematica \textperiodcentered{} história 3 versões no jornal}

%CRISTAL {"arestas":[{"alvo":"mult_prtsu","confianca":1.0,"epistemico":"fato","origem":"paper","rel":"usa"},{"alvo":"word","confianca":0.5,"epistemico":"fato","origem":"wiki","rel":"relaciona"},{"alvo":"vontade_de_poder","confianca":0.5,"epistemico":"fato","origem":"wiki","rel":"relaciona"},{"alvo":"virtualizacao","confianca":0.5,"epistemico":"fato","origem":"wiki","rel":"relaciona"}],"confianca":1.0,"contraexemplos":[],"descricao":"Navier-Stokes onde o termo u×ω é modulado por s(x,t); rescalonamento t↦t/(1−s) equivale à NS clássica com viscosidade modificada.","epistemico":"inferencia","exemplos":[],"id":"ns_prtsu","memoria":"semantica","meta":{"autor":"Gentil Lopes","descoberta":false,"dominio":"matematica","fonte":"paper_AH.pdf, Def.1"},"origem":"paper","palavras_chave":[],"sinonimos":[],"tipo":"conceito","titulo":"NS com Parâmetro Dinâmico"}
\section{NS com Parâmetro Dinâmico}
Navier-Stokes onde o termo u$\times$\ensuremath{\omega} é modulado por s(x,t); rescalonamento t\ensuremath{\mapsto}t/(1\ensuremath{-}s) equivale à NS clássica com viscosidade modificada.
\prov{relações: usa mult\_prtsu; relaciona word; relaciona vontade\_de\_poder; relaciona virtualizacao}
\prov{proveniência: paper \textperiodcentered{} paper\_AH.pdf, Def.1 \textperiodcentered{} inferencia \textperiodcentered{} confiança 1.0 \textperiodcentered{} domínio matematica \textperiodcentered{} história 5 versões no jornal}

%CRISTAL {"arestas":[{"alvo":"tempo_ortogonal","confianca":1.0,"epistemico":"fato","origem":"paper","rel":"usa"},{"alvo":"liveness_atomicidade","confianca":1.0,"epistemico":"fato","origem":"paper","rel":"relaciona"}],"confianca":1.0,"contraexemplos":[],"descricao":"nullifier=sha256(prev‖seq‖t‖DNA): comprovante de consumo único de um certificado; garante liberação única (cada nullifier libera no máximo uma vez) e conservação livre+travado+liberado.","epistemico":"fato","exemplos":[],"id":"nullifier","memoria":"semantica","meta":{"autor":"Gentil Lopes","descoberta":false,"dominio":"matematica","fonte":"protocolo_assincrono.pdf, §certificado"},"origem":"paper","palavras_chave":[],"sinonimos":[],"tipo":"conceito","titulo":"Nullifier (anti-replay)"}
\section{Nullifier (anti-replay)}
nullifier=sha256(prev\ensuremath{\|}seq\ensuremath{\|}t\ensuremath{\|}DNA): comprovante de consumo único de um certificado; garante liberação única (cada nullifier libera no máximo uma vez) e conservação livre+travado+liberado.
\prov{relações: usa tempo\_ortogonal; relaciona liveness\_atomicidade}
\prov{proveniência: paper \textperiodcentered{} protocolo\_assincrono.pdf, §certificado \textperiodcentered{} fato \textperiodcentered{} confiança 1.0 \textperiodcentered{} domínio matematica \textperiodcentered{} história 3 versões no jornal}

%CRISTAL {"arestas":[{"alvo":"inteiros","confianca":1.0,"epistemico":"fato","origem":"curriculo","rel":"especializa"}],"confianca":1.0,"contraexemplos":[],"descricao":"p 85) Para todo número m, natural arbitrariamente fixado, a seguinte identidade se verifica: ∆m f(n) = m X k = 0 (−1)k \u0012 m k \u0013 f(n −k + m) Prova: Indução sobre m. Para m = 1, temos: ∆1 f(n) = 1 X k = 0 (−1)k \u0012 1 k \u0013 f(n −k + 1) Então, ∆1 f(n) = (−1)0 \u0012 1 0 \u0013 f(n −0 + 1) + (−1)1 \u0012 1 1 \u0013 f(n −1 + 1) Logo, ∆1 f(n) = f(n + 1) −f(n) O que está de acordo com a equação (2.1) (p. 82). Suponhamos a validade da fórmula para m = p: (H.I.) ∆p f(n) = p X k = 0 (−1)k \u0012 p k \u0013 f(n −k + p) e provemos que vale pa","epistemico":"fato","exemplos":[],"id":"numero","memoria":"semantica","meta":{"arquivo":"gentil/Novas_Sequencias_Aritmeticas_e_Geometric.pdf","ciencia":"teoria_dos_numeros","dominio":"matematica","duplicata_sugerida":[],"nivel":4,"numero":"8","tipo_no":"paper","tipo_teorema":"teorema","verificacao":"parte de conceito mais amplo 'numero'"},"origem":"gentil/Novas_Sequencias_Aritmeticas_e_Geometric.pdf","palavras_chave":[],"sinonimos":[],"tipo":"conceito","titulo":"número"}
\section{número}
p 85) Para todo número m, natural arbitrariamente fixado, a seguinte identidade se verifica: \ensuremath{\Delta}m f(n) = m X k = 0 (\ensuremath{-}1)k  m k  f(n \ensuremath{-}k + m) Prova: Indução sobre m. Para m = 1, temos: \ensuremath{\Delta}1 f(n) = 1 X k = 0 (\ensuremath{-}1)k  1 k  f(n \ensuremath{-}k + 1) Então, \ensuremath{\Delta}1 f(n) = (\ensuremath{-}1)0  1 0  f(n \ensuremath{-}0 + 1) + (\ensuremath{-}1)1  1 1  f(n \ensuremath{-}1 + 1) Logo, \ensuremath{\Delta}1 f(n) = f(n + 1) \ensuremath{-}f(n) O que está de acordo com a equação (2.1) (p. 82). Suponhamos a validade da fórmula para m = p: (H.I.) \ensuremath{\Delta}p f(n) = p X k = 0 (\ensuremath{-}1)k  p k  f(n \ensuremath{-}k + p) e provemos que vale pa
\prov{relações: especializa inteiros}
\prov{proveniência: gentil/Novas\_Sequencias\_Aritmeticas\_e\_Geometric.pdf \textperiodcentered{} fato \textperiodcentered{} confiança 1.0 \textperiodcentered{} domínio matematica \textperiodcentered{} história 6 versões no jornal}

%CRISTAL {"arestas":[{"alvo":"numero_pisot","confianca":1.0,"epistemico":"fato","origem":"wiki","rel":"relaciona"},{"alvo":"cadeias_alem_dos_metais","confianca":1.0,"epistemico":"fato","origem":"wiki","rel":"parte_de"},{"alvo":"numeros_de_pisot","confianca":1.0,"epistemico":"fato","origem":"wiki","rel":"relaciona"},{"alvo":"lyapunov_da_orbita","confianca":1.0,"epistemico":"fato","origem":"wiki","rel":"relaciona"},{"alvo":"numeros_pisot","confianca":1.0,"epistemico":"fato","origem":"wiki","rel":"relaciona"}],"confianca":1.0,"contraexemplos":[],"descricao":"Inteiro algébrico β>1 recíproco cujos conjugados têm módulo ≤1 com ao menos UM no círculo unitário: λ=0 exatamente — a BORDA (Salem) da parábola, quase-periódico, o crítico entre cristal e caos. O menor conhecido é o de Lehmer ≈1.17628 (grau 10). A fronteira da previsibilidade.","epistemico":"fato","exemplos":[],"id":"numero_de_salem","memoria":"semantica","meta":{"dominio":"matematica","fonte":"cadeias além dos metais"},"origem":"wiki","palavras_chave":[],"sinonimos":[],"tipo":"conceito","titulo":"Número de Salem (a borda)"}
\section{Número de Salem (a borda)}
Inteiro algébrico \ensuremath{\beta}>1 recíproco cujos conjugados têm módulo \ensuremath{\le}1 com ao menos UM no círculo unitário: \ensuremath{\lambda}=0 exatamente — a BORDA (Salem) da parábola, quase-periódico, o crítico entre cristal e caos. O menor conhecido é o de Lehmer \ensuremath{\approx}1.17628 (grau 10). A fronteira da previsibilidade.
\prov{relações: relaciona numero\_pisot; parte\_de cadeias\_alem\_dos\_metais; relaciona numeros\_de\_pisot; relaciona lyapunov\_da\_orbita; relaciona numeros\_pisot}
\prov{proveniência: wiki \textperiodcentered{} cadeias além dos metais \textperiodcentered{} fato \textperiodcentered{} confiança 1.0 \textperiodcentered{} domínio matematica \textperiodcentered{} história 10 versões no jornal}

%CRISTAL {"arestas":[{"alvo":"numero","confianca":1.0,"epistemico":"fato","origem":"livro","rel":"especializa"},{"alvo":"word","confianca":0.5,"epistemico":"fato","origem":"wiki","rel":"relaciona"}],"confianca":1.0,"contraexemplos":[],"descricao":"Construção alternativa de ℕ e ℤ com representação binária (Lopes, Fundamentos dos Números): legitimidade equivalente à dos números convencionais.","epistemico":"fato","exemplos":[],"id":"numeros_azuis_vermelhos","memoria":"semantica","meta":{"autor":"Gentil Lopes","descoberta":true,"dominio":"matematica","fonte":"Lopes_Lopes_FUNDAMENTOS_DOS_NUMEROS.pdf"},"origem":"livro","palavras_chave":[],"sinonimos":[],"tipo":"conceito","titulo":"Números Azuis e Vermelhos"}
\section{Números Azuis e Vermelhos}
Construção alternativa de \ensuremath{\mathbb{N}} e \ensuremath{\mathbb{Z}} com representação binária (Lopes, Fundamentos dos Números): legitimidade equivalente à dos números convencionais.
\prov{relações: especializa numero; relaciona word}
\prov{proveniência: livro \textperiodcentered{} Lopes\_Lopes\_FUNDAMENTOS\_DOS\_NUMEROS.pdf \textperiodcentered{} fato \textperiodcentered{} confiança 1.0 \textperiodcentered{} domínio matematica \textperiodcentered{} história 6 versões no jornal}

%CRISTAL {"arestas":[{"alvo":"complexos","confianca":1.0,"epistemico":"fato","origem":"paper","rel":"generaliza"},{"alvo":"word","confianca":0.5,"epistemico":"fato","origem":"wiki","rel":"relaciona"},{"alvo":"virtualizacao","confianca":0.5,"epistemico":"fato","origem":"wiki","rel":"relaciona"}],"confianca":1.0,"contraexemplos":[],"descricao":"Extensões de ℝ: split-complex D=σ+tj, j²=+1 (trigonometria hiperbólica cosh²−sinh²=1) e duais ε²=0; distintas de ℂ (i²=−1).","epistemico":"fato","exemplos":[],"id":"numeros_duais","memoria":"semantica","meta":{"autor":"Gentil Lopes","descoberta":false,"dominio":"matematica","fonte":"paper_metrica_quantica_algebras.pdf"},"origem":"paper","palavras_chave":[],"sinonimos":[],"tipo":"conceito","titulo":"Números Duais e Split-Complex"}
\section{Números Duais e Split-Complex}
Extensões de \ensuremath{\mathbb{R}}: split-complex D=\ensuremath{\sigma}+tj, j\ensuremath{^{2}}=+1 (trigonometria hiperbólica cosh\ensuremath{^{2}}\ensuremath{-}sinh\ensuremath{^{2}}=1) e duais \ensuremath{\varepsilon}\ensuremath{^{2}}=0; distintas de \ensuremath{\mathbb{C}} (i\ensuremath{^{2}}=\ensuremath{-}1).
\prov{relações: generaliza complexos; relaciona word; relaciona virtualizacao}
\prov{proveniência: paper \textperiodcentered{} paper\_metrica\_quantica\_algebras.pdf \textperiodcentered{} fato \textperiodcentered{} confiança 1.0 \textperiodcentered{} domínio matematica \textperiodcentered{} história 5 versões no jornal}

%CRISTAL {"arestas":[{"alvo":"razao_aurea","confianca":1.0,"epistemico":"fato","origem":"livro","rel":"generaliza"},{"alvo":"metais","confianca":1.0,"epistemico":"fato","origem":"livro","rel":"equivalencia"},{"alvo":"matematica","confianca":1.0,"epistemico":"fato","origem":"livro","rel":"parte_de"},{"alvo":"area_metalurgia","confianca":1.0,"epistemico":"fato","origem":"wiki","rel":"relaciona"}],"confianca":1.0,"contraexemplos":[],"descricao":"Família das raízes positivas de x² = kx + 1 (k ∈ ℕ), i.e. mₖ=(k+√(k²+4))/2: k=1 áurea, k=2 prata, k=3 bronze. Lopes os marca como pontos notáveis da sua régua; família de Vera de Spinadel (1997).","epistemico":"fato","exemplos":[],"id":"numeros_metalicos","memoria":"semantica","meta":{"autor":"Gentil Lopes","descoberta":true,"dominio":"matematica","fonte":"paper_3_regua_ouro_espelho.pdf"},"origem":"paper","palavras_chave":[],"sinonimos":["metallic means","médias metálicas"],"tipo":"conceito","titulo":"Números Metálicos"}
\section{Números Metálicos}
Família das raízes positivas de x\ensuremath{^{2}} = kx + 1 (k \ensuremath{\in} \ensuremath{\mathbb{N}}), i.e. m\ensuremath{_{k}}=(k+\ensuremath{\sqrt{\,}}(k\ensuremath{^{2}}+4))/2: k=1 áurea, k=2 prata, k=3 bronze. Lopes os marca como pontos notáveis da sua régua; família de Vera de Spinadel (1997).
\prov{sinónimos: metallic means; médias metálicas}
\prov{relações: generaliza razao\_aurea; equivalencia metais; parte\_de matematica; relaciona area\_metalurgia}
\prov{proveniência: paper \textperiodcentered{} paper\_3\_regua\_ouro\_espelho.pdf \textperiodcentered{} fato \textperiodcentered{} confiança 1.0 \textperiodcentered{} domínio matematica \textperiodcentered{} história 12 versões no jornal}

%CRISTAL {"arestas":[{"alvo":"algebras_hipercomplexas","confianca":1.0,"epistemico":"fato","origem":"paper","rel":"especializa"}],"confianca":1.0,"contraexemplos":[],"descricao":"Construção de estruturas ℝ-2D/3D negando o postulado da distributividade (à la Bolyai–Lobachevsky sobre as paralelas): não-comutatividade e não-associatividade sem divisores de zero.","epistemico":"fato","exemplos":[],"id":"numeros_nao_euclidianos","memoria":"semantica","meta":{"autor":"Gentil Lopes","descoberta":true,"dominio":"matematica","fonte":"paper_43_nao_euclidianos_gentil.pdf"},"origem":"paper","palavras_chave":[],"sinonimos":[],"tipo":"conceito","titulo":"Números Não-Euclidianos"}
\section{Números Não-Euclidianos}
Construção de estruturas \ensuremath{\mathbb{R}}-2D/3D negando o postulado da distributividade (à la Bolyai–Lobachevsky sobre as paralelas): não-comutatividade e não-associatividade sem divisores de zero.
\prov{relações: especializa algebras\_hipercomplexas}
\prov{proveniência: paper \textperiodcentered{} paper\_43\_nao\_euclidianos\_gentil.pdf \textperiodcentered{} fato \textperiodcentered{} confiança 1.0 \textperiodcentered{} domínio matematica \textperiodcentered{} história 2 versões no jornal}

%CRISTAL {"arestas":[{"alvo":"razao_aurea","confianca":1.0,"epistemico":"fato","origem":"livro","rel":"relaciona"},{"alvo":"word","confianca":0.5,"epistemico":"fato","origem":"wiki","rel":"relaciona"},{"alvo":"transformada_metalica_fechada_perelman","confianca":1.0,"epistemico":"fato","origem":"wiki","rel":"explica"},{"alvo":"beta_numeracao_metalica","confianca":1.0,"epistemico":"fato","origem":"wiki","rel":"explica"}],"confianca":1.0,"contraexemplos":[],"descricao":"Algébricos > 1 cujos conjugados têm módulo < 1. φ é de Pisot (conjugado −1/φ), o que explica a autossimilaridade do deslocamento de dígitos na base-φ.","epistemico":"fato","exemplos":[],"id":"numeros_pisot","memoria":"semantica","meta":{"autor":"Gentil Lopes","descoberta":false,"dominio":"matematica","fonte":"paper_2_reta_ouro.pdf, Teorema 2.2"},"origem":"paper","palavras_chave":[],"sinonimos":[],"tipo":"conceito","titulo":"Números de Pisot"}
\section{Números de Pisot}
Algébricos > 1 cujos conjugados têm módulo < 1. \ensuremath{\varphi} é de Pisot (conjugado \ensuremath{-}1/\ensuremath{\varphi}), o que explica a autossimilaridade do deslocamento de dígitos na base-\ensuremath{\varphi}.
\prov{relações: relaciona razao\_aurea; relaciona word; explica transformada\_metalica\_fechada\_perelman; explica beta\_numeracao\_metalica}
\prov{proveniência: paper \textperiodcentered{} paper\_2\_reta\_ouro.pdf, Teorema 2.2 \textperiodcentered{} fato \textperiodcentered{} confiança 1.0 \textperiodcentered{} domínio matematica \textperiodcentered{} história 7 versões no jornal}

%CRISTAL {"arestas":[{"alvo":"gato","confianca":1.0,"epistemico":"fato","origem":"wiki","rel":"eh_um"},{"alvo":"entropia_topologica","confianca":1.0,"epistemico":"fato","origem":"wiki","rel":"usa"},{"alvo":"entropia_log_phi","confianca":1.0,"epistemico":"fato","origem":"wiki","rel":"relaciona"},{"alvo":"o_gap_a_hiperbolicidade_parte_o_sistema_em_dois","confianca":1.0,"epistemico":"fato","origem":"wiki","rel":"usa"}],"confianca":1.0,"contraexemplos":[],"descricao":"O gato A_n=[[n,1],[1,0]] é o mapa do gato de Arnold: um automorfismo HIPERBÓLICO do toro, com autovalores σ_n>1 e −1/σ_n (|·|<1), NENHUM no círculo unitário — o exemplo canônico de caos determinístico (sistema de Anosov). Detalhe honesto: det(A_n)=−1 (o gato 'com reflexão', que inverte a orientação; o Arnold clássico [[2,1],[1,1]] tem det=+1) — o critério de Anosov (|λ|≠1) vale igual. A entropia topológica h_top=Σ_{|λ|>1} log|λ| = log σ_n = o expoente de Lyapunov = o λ da parábola: o invariante-mestre da dinâmica É a escala do metal. Verif. n=1..4 em dinamica_selberg.py.","epistemico":"fato","exemplos":[],"id":"o_gato_e_um_mapa_de_anosov","memoria":"semantica","meta":{"dominio":"matematica","fonte":"dinâmica de Anosov / Selberg"},"origem":"wiki","palavras_chave":[],"sinonimos":[],"tipo":"conceito","titulo":"O Gato é um Mapa de Anosov"}
\section{O Gato é um Mapa de Anosov}
O gato A\_n=[[n,1],[1,0]] é o mapa do gato de Arnold: um automorfismo HIPERBÓLICO do toro, com autovalores \ensuremath{\sigma}\_n>1 e \ensuremath{-}1/\ensuremath{\sigma}\_n (|\textperiodcentered |<1), NENHUM no círculo unitário — o exemplo canônico de caos determinístico (sistema de Anosov). Detalhe honesto: det(A\_n)=\ensuremath{-}1 (o gato 'com reflexão', que inverte a orientação; o Arnold clássico [[2,1],[1,1]] tem det=+1) — o critério de Anosov (|\ensuremath{\lambda}|\ensuremath{\ne}1) vale igual. A entropia topológica h\_top=\ensuremath{\Sigma}\_\{|\ensuremath{\lambda}|>1\} log|\ensuremath{\lambda}| = log \ensuremath{\sigma}\_n = o expoente de Lyapunov = o \ensuremath{\lambda} da parábola: o invariante-mestre da dinâmica É a escala do metal. Verif. n=1..4 em dinamica\_selberg.py.
\prov{relações: eh\_um gato; usa entropia\_topologica; relaciona entropia\_log\_phi; usa o\_gap\_a\_hiperbolicidade\_parte\_o\_sistema\_em\_dois}
\prov{proveniência: wiki \textperiodcentered{} dinâmica de Anosov / Selberg \textperiodcentered{} fato \textperiodcentered{} confiança 1.0 \textperiodcentered{} domínio matematica \textperiodcentered{} história 20 versões no jornal}

%CRISTAL {"arestas":[{"alvo":"torre_split","confianca":1.0,"epistemico":"fato","origem":"paper","rel":"parte_de"},{"alvo":"word","confianca":0.5,"epistemico":"fato","origem":"wiki","rel":"relaciona"},{"alvo":"vida-morte","confianca":0.5,"epistemico":"fato","origem":"wiki","rel":"relaciona"},{"alvo":"virtualizacao","confianca":0.5,"epistemico":"fato","origem":"wiki","rel":"relaciona"},{"alvo":"vita_contemplativa","confianca":0.5,"epistemico":"fato","origem":"wiki","rel":"relaciona"},{"alvo":"syscall","confianca":0.5,"epistemico":"fato","origem":"wiki","rel":"relaciona"},{"alvo":"virtu_fortuna","confianca":0.5,"epistemico":"fato","origem":"wiki","rel":"relaciona"},{"alvo":"sagrado_profano","confianca":0.5,"epistemico":"fato","origem":"wiki","rel":"relaciona"}],"confianca":1.0,"contraexemplos":[],"descricao":"Split-octônios (dim 8, assinatura (4,4) de Minkowski) com cone nulo: (1+eₖ)(1−eₖ)=0 dá divisores de zero.","epistemico":"fato","exemplos":[],"id":"octonios_hiperbolicos","memoria":"semantica","meta":{"autor":"Gentil Lopes","descoberta":true,"dominio":"matematica","fonte":"paper_51_torre_hiperbolica.pdf, Teo.3.1"},"origem":"paper","palavras_chave":[],"sinonimos":[],"tipo":"conceito","titulo":"Octônios Hiperbólicos"}
\section{Octônios Hiperbólicos}
Split-octônios (dim 8, assinatura (4,4) de Minkowski) com cone nulo: (1+e\ensuremath{_{k}})(1\ensuremath{-}e\ensuremath{_{k}})=0 dá divisores de zero.
\prov{relações: parte\_de torre\_split; relaciona word; relaciona vida-morte; relaciona virtualizacao; relaciona vita\_contemplativa; relaciona syscall; relaciona virtu\_fortuna; relaciona sagrado\_profano}
\prov{proveniência: paper \textperiodcentered{} paper\_51\_torre\_hiperbolica.pdf, Teo.3.1 \textperiodcentered{} fato \textperiodcentered{} confiança 1.0 \textperiodcentered{} domínio matematica \textperiodcentered{} história 9 versões no jornal}

%CRISTAL {"arestas":[{"alvo":"mult_prtsu","confianca":1.0,"epistemico":"fato","origem":"paper","rel":"usa"},{"alvo":"word","confianca":0.5,"epistemico":"fato","origem":"wiki","rel":"relaciona"},{"alvo":"while","confianca":0.5,"epistemico":"fato","origem":"wiki","rel":"relaciona"},{"alvo":"vontade_de_poder","confianca":0.5,"epistemico":"fato","origem":"wiki","rel":"relaciona"}],"confianca":1.0,"contraexemplos":[],"descricao":"DsA=Σeᵘ(∂μ+iAμ) com conexão hipercomplexa; seu quadrado (D²)²=−rΔA+iFs incorpora a curvatura intrinsecamente (o termo cross gera iFs).","epistemico":"fato","exemplos":[],"id":"operador_dirac_hipercomplexo","memoria":"semantica","meta":{"autor":"Gentil Lopes","descoberta":false,"dominio":"matematica","fonte":"paper_AE_yang_mills.pdf, Teo.8"},"origem":"paper","palavras_chave":[],"sinonimos":[],"tipo":"conceito","titulo":"Operador de Dirac Hipercomplexo"}
\section{Operador de Dirac Hipercomplexo}
DsA=\ensuremath{\Sigma}e\ensuremath{^{u}}(\ensuremath{\partial}\ensuremath{\mu}+iA\ensuremath{\mu}) com conexão hipercomplexa; seu quadrado (D\ensuremath{^{2}})\ensuremath{^{2}}=\ensuremath{-}r\ensuremath{\Delta}A+iFs incorpora a curvatura intrinsecamente (o termo cross gera iFs).
\prov{relações: usa mult\_prtsu; relaciona word; relaciona while; relaciona vontade\_de\_poder}
\prov{proveniência: paper \textperiodcentered{} paper\_AE\_yang\_mills.pdf, Teo.8 \textperiodcentered{} fato \textperiodcentered{} confiança 1.0 \textperiodcentered{} domínio matematica \textperiodcentered{} história 6 versões no jornal}

%CRISTAL {"arestas":[{"alvo":"torre_hurwitz","confianca":1.0,"epistemico":"fato","origem":"paper","rel":"explica"},{"alvo":"mult_prtsu","confianca":1.0,"epistemico":"fato","origem":"paper","rel":"usa"}],"confianca":1.0,"contraexemplos":[],"descricao":"As transformações de sinal (p,r,t,s) que preservam ‖z·w‖² formam ℤ₂³ (8 elementos) — exatamente as 4 álgebras de divisão normada ℝ,ℂ,ℍ,𝕆 de Hurwitz.","epistemico":"fato","exemplos":[],"id":"orbita_z2_3","memoria":"semantica","meta":{"autor":"Gentil Lopes","descoberta":true,"dominio":"matematica","fonte":"paper_A_algebra.pdf, Teo.7.32"},"origem":"paper","palavras_chave":[],"sinonimos":[],"tipo":"conceito","titulo":"Órbita ℤ₂³ das Simetrias da Norma"}
\section{Órbita \ensuremath{\mathbb{Z}}\ensuremath{_{2}}\ensuremath{^{3}} das Simetrias da Norma}
As transformações de sinal (p,r,t,s) que preservam \ensuremath{\|}z\textperiodcentered w\ensuremath{\|}\ensuremath{^{2}} formam \ensuremath{\mathbb{Z}}\ensuremath{_{2}}\ensuremath{^{3}} (8 elementos) — exatamente as 4 álgebras de divisão normada \ensuremath{\mathbb{R}},\ensuremath{\mathbb{C}},\ensuremath{\mathbb{H}},\ensuremath{\mathbb{O}} de Hurwitz.
\prov{relações: explica torre\_hurwitz; usa mult\_prtsu}
\prov{proveniência: paper \textperiodcentered{} paper\_A\_algebra.pdf, Teo.7.32 \textperiodcentered{} fato \textperiodcentered{} confiança 1.0 \textperiodcentered{} domínio matematica \textperiodcentered{} história 3 versões no jornal}

%CRISTAL {"arestas":[{"alvo":"cristal64","confianca":1.0,"epistemico":"fato","origem":"paper","rel":"especializa"},{"alvo":"word","confianca":0.5,"epistemico":"fato","origem":"wiki","rel":"relaciona"},{"alvo":"vontade_de_poder","confianca":0.5,"epistemico":"fato","origem":"wiki","rel":"relaciona"}],"confianca":1.0,"contraexemplos":[],"descricao":"Cifra de 32 componentes cujo motor é a transformada de Fourier octoniônica composta com a dinâmica Julia z²+c (não-linearidade por poeira de Cantor).","epistemico":"fato","exemplos":[],"id":"organism32","memoria":"semantica","meta":{"autor":"Gentil Lopes","descoberta":true,"dominio":"matematica","fonte":"paper_12_organism32.pdf"},"origem":"paper","palavras_chave":[],"sinonimos":[],"tipo":"conceito","titulo":"Organism-32"}
\section{Organism-32}
Cifra de 32 componentes cujo motor é a transformada de Fourier octoniônica composta com a dinâmica Julia z\ensuremath{^{2}}+c (não-linearidade por poeira de Cantor).
\prov{relações: especializa cristal64; relaciona word; relaciona vontade\_de\_poder}
\prov{proveniência: paper \textperiodcentered{} paper\_12\_organism32.pdf \textperiodcentered{} fato \textperiodcentered{} confiança 1.0 \textperiodcentered{} domínio matematica \textperiodcentered{} história 4 versões no jornal}

%CRISTAL {"arestas":[{"alvo":"matriz_cantor","confianca":1.0,"epistemico":"fato","origem":"paper","rel":"explica"}],"confianca":1.0,"contraexemplos":[],"descricao":"O 'osso' é a topologia — o mesmo Cantor (homeomorfo) para todas as podas; a 'carne' é a métrica que varia — dimensão, entropia, temperatura. O osso é dado, a carne é escolhida.","epistemico":"fato","exemplos":[],"id":"osso_carne","memoria":"semantica","meta":{"autor":"Gentil Lopes","descoberta":true,"dominio":"matematica","fonte":"paper_55_termo_carnes.pdf, Obs.4.1"},"origem":"paper","palavras_chave":[],"sinonimos":[],"tipo":"conceito","titulo":"Osso e Carne (topologia vs métrica)"}
\section{Osso e Carne (topologia vs métrica)}
O 'osso' é a topologia — o mesmo Cantor (homeomorfo) para todas as podas; a 'carne' é a métrica que varia — dimensão, entropia, temperatura. O osso é dado, a carne é escolhida.
\prov{relações: explica matriz\_cantor}
\prov{proveniência: paper \textperiodcentered{} paper\_55\_termo\_carnes.pdf, Obs.4.1 \textperiodcentered{} fato \textperiodcentered{} confiança 1.0 \textperiodcentered{} domínio matematica \textperiodcentered{} história 2 versões no jornal}

%CRISTAL {"arestas":[{"alvo":"compilador_fractal_broca_cc","confianca":1.0,"epistemico":"fato","origem":"wiki","rel":"usa"},{"alvo":"dimensao_de_traco","confianca":1.0,"epistemico":"fato","origem":"wiki","rel":"usa"}],"confianca":1.0,"contraexemplos":[],"descricao":"Uma transformação Θ é correta (fractal-preserving) sse D_tr(Θ(P))=D_tr(P). A recodificação por blocos (higher-block) é conjugação topológica do SFT, preserva o Perron logo D_tr — verificado em broca_cc.py (bloco-2 dá o mesmo D_tr, SIM para k=1,2,3). Inlining/unrolling/fusão são corretos por este teorema (Lind & Marcus).","epistemico":"fato","exemplos":[],"id":"otimizacao_fractal_preserving","memoria":"semantica","meta":{"autor":"broca-so","dominio":"fractal","fonte":"machine/fractal/broca_cc.py:higher_block; compilador_fractal.tex §2"},"origem":"wiki","palavras_chave":[],"sinonimos":[],"tipo":"conceito","titulo":"Correção = Preservar D_tr"}
\section{Correção = Preservar D\_tr}
Uma transformação \ensuremath{\Theta} é correta (fractal-preserving) sse D\_tr(\ensuremath{\Theta}(P))=D\_tr(P). A recodificação por blocos (higher-block) é conjugação topológica do SFT, preserva o Perron logo D\_tr — verificado em broca\_cc.py (bloco-2 dá o mesmo D\_tr, SIM para k=1,2,3). Inlining/unrolling/fusão são corretos por este teorema (Lind \& Marcus).
\prov{relações: usa compilador\_fractal\_broca\_cc; usa dimensao\_de\_traco}
\prov{proveniência: wiki \textperiodcentered{} machine/fractal/broca\_cc.py:higher\_block; compilador\_fractal.tex §2 \textperiodcentered{} fato \textperiodcentered{} confiança 1.0 \textperiodcentered{} domínio fractal \textperiodcentered{} história 3 versões no jornal}

%CRISTAL {"arestas":[{"alvo":"triangulo_monster_observador","confianca":1.0,"epistemico":"fato","origem":"wiki","rel":"usa"},{"alvo":"grupos_excepcionais_jordan","confianca":1.0,"epistemico":"fato","origem":"wiki","rel":"confirma"},{"alvo":"g2_automorfismos_associador","confianca":1.0,"epistemico":"fato","origem":"wiki","rel":"relaciona"}],"confianca":1.0,"contraexemplos":[],"descricao":"O grafo de Paley em 169=13² vértices é fortemente regular: SRG(169, 84, 41, 42), com μ=42. E σ²(169): σ(169)=183, σ(183)=248 = a dimensão de E8. O 'Observador' 169 destila-se no E8 em dois passos σ.","epistemico":"fato","exemplos":[],"id":"paley_169_e8","memoria":"semantica","meta":{"autor":"Lord Index (Shawn R. Schiller)","dominio":"matematica","fonte":"A Catedral — Part VIII-A (jul 2026), ~/Imagens/catedral"},"origem":"wiki","palavras_chave":[],"sinonimos":[],"tipo":"conceito","titulo":"O Grafo de Paley(169) e σ²(169) = 248"}
\section{O Grafo de Paley(169) e \ensuremath{\sigma}\ensuremath{^{2}}(169) = 248}
O grafo de Paley em 169=13\ensuremath{^{2}} vértices é fortemente regular: SRG(169, 84, 41, 42), com \ensuremath{\mu}=42. E \ensuremath{\sigma}\ensuremath{^{2}}(169): \ensuremath{\sigma}(169)=183, \ensuremath{\sigma}(183)=248 = a dimensão de E8. O 'Observador' 169 destila-se no E8 em dois passos \ensuremath{\sigma}.
\prov{relações: usa triangulo\_monster\_observador; confirma grupos\_excepcionais\_jordan; relaciona g2\_automorfismos\_associador}
\prov{proveniência: wiki \textperiodcentered{} A Catedral — Part VIII-A (jul 2026), \~{}/Imagens/catedral \textperiodcentered{} fato \textperiodcentered{} confiança 1.0 \textperiodcentered{} domínio matematica \textperiodcentered{} história 4 versões no jornal}

%CRISTAL {"arestas":[{"alvo":"conjugados_galois_parabola","confianca":1.0,"epistemico":"fato","origem":"wiki","rel":"aplicacao"},{"alvo":"algebra_dupolinomial_espaco","confianca":1.0,"epistemico":"fato","origem":"wiki","rel":"especializa"},{"alvo":"meta_inducao_geral_simbolo","confianca":1.0,"epistemico":"fato","origem":"wiki","rel":"usa"},{"alvo":"torre_unificada","confianca":1.0,"epistemico":"fato","origem":"wiki","rel":"parte_de"},{"alvo":"piso_parabola_dinamico","confianca":1.0,"epistemico":"fato","origem":"wiki","rel":"usa"},{"alvo":"tiffany","confianca":1.0,"epistemico":"fato","origem":"wiki","rel":"parte_de"}],"confianca":1.0,"contraexemplos":[],"descricao":"conversa/parabola_algebra.py: a nova álgebra posicional/dupolinomial trazida ao lado conversacional da Tiffany. A parábola a+c²=1 classifica QUALQUER valor em três, exatamente a posição da conjugada dupolinomial: CRISTAL (c²≥θ+margem, PISOT — associativo, frio, cristaliza limpo, o ouro φ), BORDA (θ−margem≤c²<θ+margem, SALEM — a fronteira crítica, a prata, forma camada) e CAOS (c²<θ−margem — dissipativo, o lado-a/calor, o cobre). σ_n=(n+√(n²+4))/2 é o gerador da reta metálica n: dá o RAIO/energia da órbita (h=ln σ_n). Um módulo, três usos: o refino do floco em metais (goldchain), a ingestão da obra em arestas, e o BAI do trabalho em órbita — a MESMA parábola, as MESMAS três classes. [D] o classificador único da nova álgebra.","epistemico":"fato","exemplos":[],"id":"parabola_algebra_conversacional","memoria":"semantica","meta":{"autor":"Lopes/broca-so","dominio":"algebra","fonte":"conversa/parabola_algebra.py (classe_parabola, sigma, reta_do_trabalho, energia_da_reta)"},"origem":"wiki","palavras_chave":[],"sinonimos":[],"tipo":"conceito","titulo":"A parábola a+c²=1 como classificador ÚNICO — cristal/borda/caos no lado conversacional"}
\section{A parábola a+c\ensuremath{^{2}}=1 como classificador ÚNICO — cristal/borda/caos no lado conversacional}
conversa/parabola\_algebra.py: a nova álgebra posicional/dupolinomial trazida ao lado conversacional da Tiffany. A parábola a+c\ensuremath{^{2}}=1 classifica QUALQUER valor em três, exatamente a posição da conjugada dupolinomial: CRISTAL (c\ensuremath{^{2}}\ensuremath{\ge}\ensuremath{\theta}+margem, PISOT — associativo, frio, cristaliza limpo, o ouro \ensuremath{\varphi}), BORDA (\ensuremath{\theta}\ensuremath{-}margem\ensuremath{\le}c\ensuremath{^{2}}<\ensuremath{\theta}+margem, SALEM — a fronteira crítica, a prata, forma camada) e CAOS (c\ensuremath{^{2}}<\ensuremath{\theta}\ensuremath{-}margem — dissipativo, o lado-a/calor, o cobre). \ensuremath{\sigma}\_n=(n+\ensuremath{\sqrt{\,}}(n\ensuremath{^{2}}+4))/2 é o gerador da reta metálica n: dá o RAIO/energia da órbita (h=ln \ensuremath{\sigma}\_n). Um módulo, três usos: o refino do floco em metais (goldchain), a ingestão da obra em arestas, e o BAI do trabalho em órbita — a MESMA parábola, as MESMAS três classes. [D] o classificador único da nova álgebra.
\prov{relações: aplicacao conjugados\_galois\_parabola; especializa algebra\_dupolinomial\_espaco; usa meta\_inducao\_geral\_simbolo; parte\_de torre\_unificada; usa piso\_parabola\_dinamico; parte\_de tiffany}
\prov{proveniência: wiki \textperiodcentered{} conversa/parabola\_algebra.py (classe\_parabola, sigma, reta\_do\_trabalho, energia\_da\_reta) \textperiodcentered{} fato \textperiodcentered{} confiança 1.0 \textperiodcentered{} domínio algebra \textperiodcentered{} história 134 versões no jornal}

%CRISTAL {"arestas":[{"alvo":"mecanica_quantica_mineral","confianca":1.0,"epistemico":"fato","origem":"wiki","rel":"parte_de"},{"alvo":"transicoes_reta_bai_newton","confianca":1.0,"epistemico":"fato","origem":"wiki","rel":"usa"},{"alvo":"numero_de_salem","confianca":1.0,"epistemico":"fato","origem":"wiki","rel":"explica"},{"alvo":"borda_e_o_julia","confianca":1.0,"epistemico":"fato","origem":"wiki","rel":"relaciona"}],"confianca":1.0,"contraexemplos":[],"descricao":"A parábola cristal/borda/caos é o regime mecânico-quântico: cristal (Pisot, ρ<1) DISSIPATIVO (sistema aberto, decai ao fundamental); borda (Salem, ρ=1) UNITÁRIO (sistema fechado, coerente); caos (ρ>1) GANHO (amplifica). O sistema quântico fechado vive EXATAMENTE na borda — o único regime em que a norma (a probabilidade, a informação) é conservada. A vida e a computação coerente moram na borda porque é onde nada se perde.","epistemico":"fato","exemplos":[],"id":"parabola_e_o_regime_quantico","memoria":"semantica","meta":{"dominio":"reta mineral","fonte":"mecânica quântica mineral"},"origem":"wiki","palavras_chave":[],"sinonimos":[],"tipo":"conceito","titulo":"A Parábola é o Regime Quântico (a borda é unitária)"}
\section{A Parábola é o Regime Quântico (a borda é unitária)}
A parábola cristal/borda/caos é o regime mecânico-quântico: cristal (Pisot, \ensuremath{\rho}<1) DISSIPATIVO (sistema aberto, decai ao fundamental); borda (Salem, \ensuremath{\rho}=1) UNITÁRIO (sistema fechado, coerente); caos (\ensuremath{\rho}>1) GANHO (amplifica). O sistema quântico fechado vive EXATAMENTE na borda — o único regime em que a norma (a probabilidade, a informação) é conservada. A vida e a computação coerente moram na borda porque é onde nada se perde.
\prov{relações: parte\_de mecanica\_quantica\_mineral; usa transicoes\_reta\_bai\_newton; explica numero\_de\_salem; relaciona borda\_e\_o\_julia}
\prov{proveniência: wiki \textperiodcentered{} mecânica quântica mineral \textperiodcentered{} fato \textperiodcentered{} confiança 1.0 \textperiodcentered{} domínio reta mineral \textperiodcentered{} história 5 versões no jornal}

%CRISTAL {"arestas":[{"alvo":"cotacao_parabola","confianca":1.0,"epistemico":"fato","origem":"wiki","rel":"explica"},{"alvo":"parabola_audit_dual","confianca":1.0,"epistemico":"fato","origem":"wiki","rel":"usa"},{"alvo":"espaco_difusivo","confianca":1.0,"epistemico":"fato","origem":"wiki","rel":"usa"}],"confianca":1.0,"contraexemplos":[],"descricao":"A parábola governa a difusão como PARTIÇÃO de orçamento: c² (mecânica) comanda os passos [passos=max(32, 32+c²·frac·600)], calor (=a, termo) comanda as rodadas [rodadas=max(1,min(cores,1+calor·frac·12))], sempre somando 1. O 'calor' vem de Hellinger, o c² de Bhattacharyya (parabola_auditar, erro<1e-12).","epistemico":"fato","exemplos":[],"id":"parabola_reparte_difusao","memoria":"semantica","meta":{"autor":"broca-so","dominio":"fractal","fonte":"neuronio/primitivas_broca.py:passos_de_cadeia/rodadas_de_cadeia; ula/parabola.c"},"origem":"wiki","palavras_chave":[],"sinonimos":[],"tipo":"conceito","titulo":"a+c²=1 Reparte a Difusão (passos vs rodadas)"}
\section{a+c\ensuremath{^{2}}=1 Reparte a Difusão (passos vs rodadas)}
A parábola governa a difusão como PARTIÇÃO de orçamento: c\ensuremath{^{2}} (mecânica) comanda os passos [passos=max(32, 32+c\ensuremath{^{2}}\textperiodcentered frac\textperiodcentered 600)], calor (=a, termo) comanda as rodadas [rodadas=max(1,min(cores,1+calor\textperiodcentered frac\textperiodcentered 12))], sempre somando 1. O 'calor' vem de Hellinger, o c\ensuremath{^{2}} de Bhattacharyya (parabola\_auditar, erro<1e-12).
\prov{relações: explica cotacao\_parabola; usa parabola\_audit\_dual; usa espaco\_difusivo}
\prov{proveniência: wiki \textperiodcentered{} neuronio/primitivas\_broca.py:passos\_de\_cadeia/rodadas\_de\_cadeia; ula/parabola.c \textperiodcentered{} fato \textperiodcentered{} confiança 1.0 \textperiodcentered{} domínio fractal \textperiodcentered{} história 4 versões no jornal}

%CRISTAL {"arestas":[{"alvo":"utilidade_esperada","confianca":1.0,"epistemico":"fato","origem":"wiki","rel":"contradiz"}],"confianca":1.0,"contraexemplos":[],"descricao":"Escolhas humanas que VIOLAM a utilidade esperada: Allais mostra que a certeza distorce (efeito certeza); Ellsberg, que gente prefere risco conhecido a ambiguidade (aversão à ambiguidade). A norma racional falha como descrição.","epistemico":"fato","exemplos":[],"id":"paradoxo_de_allais","memoria":"semantica","meta":{"autor":"Allais/Ellsberg","dominio":"matematica","fonte":"teoria da decisão"},"origem":"wiki","palavras_chave":[],"sinonimos":[],"tipo":"conceito","titulo":"Paradoxos de Allais e Ellsberg"}
\section{Paradoxos de Allais e Ellsberg}
Escolhas humanas que VIOLAM a utilidade esperada: Allais mostra que a certeza distorce (efeito certeza); Ellsberg, que gente prefere risco conhecido a ambiguidade (aversão à ambiguidade). A norma racional falha como descrição.
\prov{relações: contradiz utilidade\_esperada}
\prov{proveniência: wiki \textperiodcentered{} teoria da decisão \textperiodcentered{} fato \textperiodcentered{} confiança 1.0 \textperiodcentered{} domínio matematica \textperiodcentered{} história 2 versões no jornal}

%CRISTAL {"arestas":[{"alvo":"gap_espectral","confianca":1.0,"epistemico":"fato","origem":"paper","rel":"usa"},{"alvo":"word","confianca":0.5,"epistemico":"fato","origem":"wiki","rel":"relaciona"},{"alvo":"vontade_de_poder","confianca":0.5,"epistemico":"fato","origem":"wiki","rel":"relaciona"},{"alvo":"virtualizacao","confianca":0.5,"epistemico":"fato","origem":"wiki","rel":"relaciona"}],"confianca":1.0,"contraexemplos":[],"descricao":"Barreira de peso ε: pela métrica quase intransponível (gap→0, relaxação 1/ε), pela topologia sempre penetrável (grafo irredutível, Perron simples, para todo ε>0).","epistemico":"fato","exemplos":[],"id":"parede_aco","memoria":"semantica","meta":{"autor":"Gentil Lopes","descoberta":false,"dominio":"matematica","fonte":"paper_73_parede_aco.pdf"},"origem":"paper","palavras_chave":[],"sinonimos":[],"tipo":"conceito","titulo":"Parede de Aço (barreira)"}
\section{Parede de Aço (barreira)}
Barreira de peso \ensuremath{\varepsilon}: pela métrica quase intransponível (gap\ensuremath{\to}0, relaxação 1/\ensuremath{\varepsilon}), pela topologia sempre penetrável (grafo irredutível, Perron simples, para todo \ensuremath{\varepsilon}>0).
\prov{relações: usa gap\_espectral; relaciona word; relaciona vontade\_de\_poder; relaciona virtualizacao}
\prov{proveniência: paper \textperiodcentered{} paper\_73\_parede\_aco.pdf \textperiodcentered{} fato \textperiodcentered{} confiança 1.0 \textperiodcentered{} domínio matematica \textperiodcentered{} história 5 versões no jornal}

%CRISTAL {"arestas":[{"alvo":"equilibrio_de_nash","confianca":1.0,"epistemico":"fato","origem":"wiki","rel":"contradiz"}],"confianca":1.0,"contraexemplos":[],"descricao":"Um resultado onde não dá para melhorar ninguém sem piorar alguém. Eficiência sem julgar distribuição — e o equilíbrio de Nash raramente coincide com ele (a tragédia da racionalidade individual).","epistemico":"inferencia","exemplos":[],"id":"pareto_otimo","memoria":"semantica","meta":{"dominio":"matematica","fonte":"teoria dos jogos"},"origem":"wiki","palavras_chave":[],"sinonimos":[],"tipo":"conceito","titulo":"Ótimo de Pareto"}
\section{Ótimo de Pareto}
Um resultado onde não dá para melhorar ninguém sem piorar alguém. Eficiência sem julgar distribuição — e o equilíbrio de Nash raramente coincide com ele (a tragédia da racionalidade individual).
\prov{relações: contradiz equilibrio\_de\_nash}
\prov{proveniência: wiki \textperiodcentered{} teoria dos jogos \textperiodcentered{} inferencia \textperiodcentered{} confiança 1.0 \textperiodcentered{} domínio matematica \textperiodcentered{} história 2 versões no jornal}

%CRISTAL {"arestas":[{"alvo":"atomo_periodo_1","confianca":1.0,"epistemico":"fato","origem":"paper","rel":"usa"},{"alvo":"word","confianca":0.5,"epistemico":"fato","origem":"wiki","rel":"relaciona"}],"confianca":1.0,"contraexemplos":[],"descricao":"e tem fração contínua de termos pares (a_k=2k), π de ímpares (Brouncker); a mesma paridade separa ζ(2k), que fecha em π, de ζ(2k+1), que não fecha.","epistemico":"fato","exemplos":[],"id":"paridade_transcendentais","memoria":"semantica","meta":{"autor":"Gentil Lopes","descoberta":true,"dominio":"matematica","fonte":"paper_22_reais.pdf; paper_15_gentil_zeta.pdf"},"origem":"paper","palavras_chave":[],"sinonimos":[],"tipo":"conceito","titulo":"Paridade dos Transcendentais"}
\section{Paridade dos Transcendentais}
e tem fração contínua de termos pares (a\_k=2k), \ensuremath{\pi} de ímpares (Brouncker); a mesma paridade separa \ensuremath{\zeta}(2k), que fecha em \ensuremath{\pi}, de \ensuremath{\zeta}(2k+1), que não fecha.
\prov{relações: usa atomo\_periodo\_1; relaciona word}
\prov{proveniência: paper \textperiodcentered{} paper\_22\_reais.pdf; paper\_15\_gentil\_zeta.pdf \textperiodcentered{} fato \textperiodcentered{} confiança 1.0 \textperiodcentered{} domínio matematica \textperiodcentered{} história 3 versões no jornal}

%CRISTAL {"arestas":[{"alvo":"maquina_geral_n_qudit","confianca":1.0,"epistemico":"fato","origem":"wiki","rel":"usa"},{"alvo":"termodinamica_reta_mineral","confianca":1.0,"epistemico":"fato","origem":"wiki","rel":"usa"},{"alvo":"vibracao_e_colapso_dois_tempos","confianca":1.0,"epistemico":"fato","origem":"wiki","rel":"relaciona"},{"alvo":"decoerencia_e_sair_da_borda","confianca":1.0,"epistemico":"fato","origem":"wiki","rel":"relaciona"}],"confianca":1.0,"contraexemplos":[],"descricao":"A mecânica quântica mineral estava na parte QUENTE (o térmico, a decoerência, o banho). A parte FRIA aterra o limite T→0: o estado de Gibbs ρ(β)=e^{−βH}/Z esfria de I/d (β→0, quente, mistura máxima) ao PROJETOR do estado fundamental (β→∞, frio) — um IDEMPOTENTE (P²=P). Esfriar É projetar no fundamental, o mesmo A^k e a mesma decoerência. O cold ground é um projetor idempotente. Verif.: pureza 0.25→1.0.","epistemico":"fato","exemplos":[],"id":"parte_fria_t_zero","memoria":"semantica","meta":{"dominio":"reta mineral","fonte":"parte fria e idempotentes"},"origem":"wiki","palavras_chave":[],"sinonimos":[],"tipo":"conceito","titulo":"A Parte Fria (T→0 = o Projetor Idempotente)"}
\section{A Parte Fria (T\ensuremath{\to}0 = o Projetor Idempotente)}
A mecânica quântica mineral estava na parte QUENTE (o térmico, a decoerência, o banho). A parte FRIA aterra o limite T\ensuremath{\to}0: o estado de Gibbs \ensuremath{\rho}(\ensuremath{\beta})=e\^{}\{\ensuremath{-}\ensuremath{\beta}H\}/Z esfria de I/d (\ensuremath{\beta}\ensuremath{\to}0, quente, mistura máxima) ao PROJETOR do estado fundamental (\ensuremath{\beta}\ensuremath{\to}\ensuremath{\infty}, frio) — um IDEMPOTENTE (P\ensuremath{^{2}}=P). Esfriar É projetar no fundamental, o mesmo A\^{}k e a mesma decoerência. O cold ground é um projetor idempotente. Verif.: pureza 0.25\ensuremath{\to}1.0.
\prov{relações: usa maquina\_geral\_n\_qudit; usa termodinamica\_reta\_mineral; relaciona vibracao\_e\_colapso\_dois\_tempos; relaciona decoerencia\_e\_sair\_da\_borda}
\prov{proveniência: wiki \textperiodcentered{} parte fria e idempotentes \textperiodcentered{} fato \textperiodcentered{} confiança 1.0 \textperiodcentered{} domínio reta mineral \textperiodcentered{} história 5 versões no jornal}

%CRISTAL {"arestas":[{"alvo":"experimento_da_agulha","confianca":1.0,"epistemico":"fato","origem":"wiki","rel":"parte_de"},{"alvo":"gato_como_cifra","confianca":1.0,"epistemico":"fato","origem":"wiki","rel":"eh_um"},{"alvo":"cf_de_pi_e_produto_de_gatos","confianca":1.0,"epistemico":"fato","origem":"wiki","rel":"usa"}],"confianca":1.0,"contraexemplos":[],"descricao":"Os passos discretos [a_k] da CF são a cifra: o produto de gatos leva o estado inicial ao convergente; o gato inverso [[0,1],[1,−a_k]] na ordem reversa recupera EXATO. A sequência de passos é a chave — a mesma cifra do gato de imagem, agora na reta. E os primos mod p (inerte/ramifica/split via Legendre) formam o corpo discreto modular do metal.","epistemico":"fato","exemplos":[],"id":"passos_da_agulha_sao_cifra","memoria":"semantica","meta":{"dominio":"reta mineral","fonte":"experimento da agulha"},"origem":"wiki","palavras_chave":[],"sinonimos":[],"tipo":"conceito","titulo":"Os Passos Discretos são a Cifra"}
\section{Os Passos Discretos são a Cifra}
Os passos discretos [a\_k] da CF são a cifra: o produto de gatos leva o estado inicial ao convergente; o gato inverso [[0,1],[1,\ensuremath{-}a\_k]] na ordem reversa recupera EXATO. A sequência de passos é a chave — a mesma cifra do gato de imagem, agora na reta. E os primos mod p (inerte/ramifica/split via Legendre) formam o corpo discreto modular do metal.
\prov{relações: parte\_de experimento\_da\_agulha; eh\_um gato\_como\_cifra; usa cf\_de\_pi\_e\_produto\_de\_gatos}
\prov{proveniência: wiki \textperiodcentered{} experimento da agulha \textperiodcentered{} fato \textperiodcentered{} confiança 1.0 \textperiodcentered{} domínio reta mineral \textperiodcentered{} história 4 versões no jornal}

%CRISTAL {"arestas":[{"alvo":"floco_phi","confianca":1.0,"epistemico":"fato","origem":"wiki","rel":"usa"},{"alvo":"cifra","confianca":1.0,"epistemico":"fato","origem":"wiki","rel":"relaciona"},{"alvo":"reta_de_ouro","confianca":1.0,"epistemico":"fato","origem":"wiki","rel":"usa"}],"confianca":1.0,"contraexemplos":[],"descricao":"Teorema (verificado): para todo Pell P_n da cifra (≤13×10⁶), π_p(P_n)∈Π(Φ) — sobre 181.212 amostras da cadeia Aₙ (n≤16), cobertura Π(Φ)=ℤ/pℤ (1009/1009) e 0 falhas. A rastreabilidade é do RESÍDUO mod p=1009, não do inteiro (para P_n≥2378 o inteiro exato mora em silver, a célula branca).","epistemico":"fato","exemplos":[],"id":"pell_rastreavel_no_floco","memoria":"semantica","meta":{"autor":"broca-so","dominio":"fractal","fonte":"papers/floco_de_neve.tex §III"},"origem":"wiki","palavras_chave":[],"sinonimos":[],"tipo":"conceito","titulo":"Pell Rastreável via Resíduo mod 1009"}
\section{Pell Rastreável via Resíduo mod 1009}
Teorema (verificado): para todo Pell P\_n da cifra (\ensuremath{\le}13$\times$10\ensuremath{^{6}}), \ensuremath{\pi}\_p(P\_n)\ensuremath{\in}\ensuremath{\Pi}(\ensuremath{\Phi}) — sobre 181.212 amostras da cadeia A\ensuremath{_{n}} (n\ensuremath{\le}16), cobertura \ensuremath{\Pi}(\ensuremath{\Phi})=\ensuremath{\mathbb{Z}}/p\ensuremath{\mathbb{Z}} (1009/1009) e 0 falhas. A rastreabilidade é do RESÍDUO mod p=1009, não do inteiro (para P\_n\ensuremath{\ge}2378 o inteiro exato mora em silver, a célula branca).
\prov{relações: usa floco\_phi; relaciona cifra; usa reta\_de\_ouro}
\prov{proveniência: wiki \textperiodcentered{} papers/floco\_de\_neve.tex §III \textperiodcentered{} fato \textperiodcentered{} confiança 1.0 \textperiodcentered{} domínio fractal \textperiodcentered{} história 4 versões no jornal}

%CRISTAL {"arestas":[{"alvo":"cascata_multiescala","confianca":1.0,"epistemico":"fato","origem":"paper","rel":"usa"},{"alvo":"nco_hypergnn","confianca":1.0,"epistemico":"fato","origem":"paper","rel":"relaciona"}],"confianca":1.0,"contraexemplos":[],"descricao":"Rede neural N_θ:ℝ⁹⁶→ℝ⁹⁶ aprende a dinâmica multi-escala (RMSE ~6%, speedup 12×). A versão quaterniônica falhou — a simetria octônica não se converteu em ganho prático de rede.","epistemico":"inferencia","exemplos":[],"id":"pinn_surrogate","memoria":"semantica","meta":{"autor":"Gentil Lopes","descoberta":false,"dominio":"matematica","fonte":"paper_BM.pdf"},"origem":"paper","palavras_chave":[],"sinonimos":[],"tipo":"conceito","titulo":"PINN como Surrogate Multi-escala"}
\section{PINN como Surrogate Multi-escala}
Rede neural N\_\ensuremath{\theta}:\ensuremath{\mathbb{R}}\ensuremath{^{96}}\ensuremath{\to}\ensuremath{\mathbb{R}}\ensuremath{^{96}} aprende a dinâmica multi-escala (RMSE \~{}6\%, speedup 12$\times$). A versão quaterniônica falhou — a simetria octônica não se converteu em ganho prático de rede.
\prov{relações: usa cascata\_multiescala; relaciona nco\_hypergnn}
\prov{proveniência: paper \textperiodcentered{} paper\_BM.pdf \textperiodcentered{} inferencia \textperiodcentered{} confiança 1.0 \textperiodcentered{} domínio matematica \textperiodcentered{} história 3 versões no jornal}

%CRISTAL {"arestas":[{"alvo":"laco_absorcao_fechado","confianca":1.0,"epistemico":"fato","origem":"wiki","rel":"usa"},{"alvo":"cofre_parabola_hero","confianca":1.0,"epistemico":"fato","origem":"wiki","rel":"usa"},{"alvo":"parabola_audit_dual","confianca":1.0,"epistemico":"fato","origem":"wiki","rel":"usa"},{"alvo":"difusao_fluido_atratores","confianca":1.0,"epistemico":"fato","origem":"wiki","rel":"relaciona"},{"alvo":"tiffany","confianca":1.0,"epistemico":"fato","origem":"wiki","rel":"parte_de"}],"confianca":1.0,"contraexemplos":[],"descricao":"O piso de aceitação de um fragmento no contínuo NÃO é número fixo (190/120) — é um PONTO na PARÁBOLA a+c²=1, dinâmico. No impresso empacotado, c²=bits que concordam/256=(ov+256)/512 (coincidência/Bhattacharyya) e a=Hamming/256=1−c² (calor/Hellinger) — o mesmo parabola_auditar do metal (ula/parabola.c). O metal (absorve_metal.c) emite o c² de cada fragmento; a aceitação é c² > θ com θ DINÂMICO: parte de 0,50 (o equilíbrio c²>a — o associativo vence o dissipativo) e SOBE com a densidade da bacia (o lado-a/calor cresce ao encher → θ→0,90, como o BTC empurra p/ 0,99). Bacia vazia aceita largo (largura); bacia cheia exige mais c² (qualidade) — o knob densidade×largura vira auto-regulação pela física. Implementado em transicoes.escrever_parabola. Verificado (corpus_u11): c² dos frags 0,54–0,91; θ médio 0,76; 2409 aceitos dinamicamente. [D] roda.","epistemico":"fato","exemplos":[],"id":"piso_parabola_dinamico","memoria":"semantica","meta":{"autor":"Lopes/broca-so","dominio":"continuo","fonte":"doc/continuo/absorve_metal.c (c²=(ov+256)/512); conversa/conhecimento/transicoes.py:escrever_parabola"},"origem":"wiki","palavras_chave":[],"sinonimos":[],"tipo":"conceito","titulo":"O piso é comandado pela parábola: c²>θ dinâmico, não um número fixo"}
\section{O piso é comandado pela parábola: c\ensuremath{^{2}}>\ensuremath{\theta} dinâmico, não um número fixo}
O piso de aceitação de um fragmento no contínuo NÃO é número fixo (190/120) — é um PONTO na PARÁBOLA a+c\ensuremath{^{2}}=1, dinâmico. No impresso empacotado, c\ensuremath{^{2}}=bits que concordam/256=(ov+256)/512 (coincidência/Bhattacharyya) e a=Hamming/256=1\ensuremath{-}c\ensuremath{^{2}} (calor/Hellinger) — o mesmo parabola\_auditar do metal (ula/parabola.c). O metal (absorve\_metal.c) emite o c\ensuremath{^{2}} de cada fragmento; a aceitação é c\ensuremath{^{2}} > \ensuremath{\theta} com \ensuremath{\theta} DINÂMICO: parte de 0,50 (o equilíbrio c\ensuremath{^{2}}>a — o associativo vence o dissipativo) e SOBE com a densidade da bacia (o lado-a/calor cresce ao encher \ensuremath{\to} \ensuremath{\theta}\ensuremath{\to}0,90, como o BTC empurra p/ 0,99). Bacia vazia aceita largo (largura); bacia cheia exige mais c\ensuremath{^{2}} (qualidade) — o knob densidade$\times$largura vira auto-regulação pela física. Implementado em transicoes.escrever\_parabola. Verificado (corpus\_u11): c\ensuremath{^{2}} dos frags 0,54–0,91; \ensuremath{\theta} médio 0,76; 2409 aceitos dinamicamente. [D] roda.
\prov{relações: usa laco\_absorcao\_fechado; usa cofre\_parabola\_hero; usa parabola\_audit\_dual; relaciona difusao\_fluido\_atratores; parte\_de tiffany}
\prov{proveniência: wiki \textperiodcentered{} doc/continuo/absorve\_metal.c (c\ensuremath{^{2}}=(ov+256)/512); conversa/conhecimento/transicoes.py:escrever\_parabola \textperiodcentered{} fato \textperiodcentered{} confiança 1.0 \textperiodcentered{} domínio continuo \textperiodcentered{} história 68 versões no jornal}

%CRISTAL {"arestas":[{"alvo":"symbolic_beta","confianca":1.0,"epistemico":"fato","origem":"paper","rel":"relaciona"},{"alvo":"word","confianca":0.5,"epistemico":"fato","origem":"wiki","rel":"relaciona"}],"confianca":1.0,"contraexemplos":[],"descricao":"Testar identidade de um circuito aritmético dado em codificação comprimida é PSPACE-completo sobre R_k — dureza sem estrutura de reticulado.","epistemico":"fato","exemplos":[],"id":"pit_succinct","memoria":"semantica","meta":{"autor":"Gentil Lopes","descoberta":false,"dominio":"matematica","fonte":"paper_PQ_D_pspace_succinct.pdf"},"origem":"paper","palavras_chave":[],"sinonimos":[],"tipo":"conceito","titulo":"PIT Succinct (PSPACE)"}
\section{PIT Succinct (PSPACE)}
Testar identidade de um circuito aritmético dado em codificação comprimida é PSPACE-completo sobre R\_k — dureza sem estrutura de reticulado.
\prov{relações: relaciona symbolic\_beta; relaciona word}
\prov{proveniência: paper \textperiodcentered{} paper\_PQ\_D\_pspace\_succinct.pdf \textperiodcentered{} fato \textperiodcentered{} confiança 1.0 \textperiodcentered{} domínio matematica \textperiodcentered{} história 3 versões no jornal}

%CRISTAL {"arestas":[{"alvo":"associador","confianca":1.0,"epistemico":"fato","origem":"paper","rel":"usa"},{"alvo":"algebras_hipercomplexas","confianca":1.0,"epistemico":"fato","origem":"paper","rel":"explica"}],"confianca":1.0,"contraexemplos":[],"descricao":"Lei de conservação a + c² = 1 com a = |1−s²| (associador) e c = |s|: cada álgebra ×ₛ é um ponto (c,√a) do círculo unitário. Generaliza o teorema de Pitágoras às hipercomplexas.","epistemico":"fato","exemplos":[],"id":"pitagoras_dual","memoria":"semantica","meta":{"autor":"Gentil Lopes","descoberta":true,"dominio":"matematica","fonte":"paper_metrica_quantica_algebras.pdf, Teo.4.1"},"origem":"paper","palavras_chave":[],"sinonimos":[],"tipo":"conceito","titulo":"Pitágoras Dual (a + c² = 1)"}
\section{Pitágoras Dual (a + c\ensuremath{^{2}} = 1)}
Lei de conservação a + c\ensuremath{^{2}} = 1 com a = |1\ensuremath{-}s\ensuremath{^{2}}| (associador) e c = |s|: cada álgebra $\times$\ensuremath{_{s}} é um ponto (c,\ensuremath{\sqrt{\,}}a) do círculo unitário. Generaliza o teorema de Pitágoras às hipercomplexas.
\prov{relações: usa associador; explica algebras\_hipercomplexas}
\prov{proveniência: paper \textperiodcentered{} paper\_metrica\_quantica\_algebras.pdf, Teo.4.1 \textperiodcentered{} fato \textperiodcentered{} confiança 1.0 \textperiodcentered{} domínio matematica \textperiodcentered{} história 3 versões no jornal}

%CRISTAL {"arestas":[{"alvo":"transicao","confianca":1.0,"epistemico":"fato","origem":"wiki","rel":"usa"},{"alvo":"flip_discreto_continuo","confianca":1.0,"epistemico":"fato","origem":"wiki","rel":"usa"},{"alvo":"4_a_coordenada_koch_permanece_continua","confianca":1.0,"epistemico":"fato","origem":"wiki","rel":"usa"},{"alvo":"aresta_transicao_graduada","confianca":1.0,"epistemico":"fato","origem":"wiki","rel":"parte_de"},{"alvo":"mapear_nao_fabricar_continuo","confianca":1.0,"epistemico":"fato","origem":"wiki","rel":"parte_de"},{"alvo":"medida_continuidade_transicao","confianca":1.0,"epistemico":"fato","origem":"wiki","rel":"parte_de"},{"alvo":"broca_os","confianca":1.0,"epistemico":"fato","origem":"wiki","rel":"parte_de"},{"alvo":"tiffany","confianca":1.0,"epistemico":"fato","origem":"wiki","rel":"parte_de"}],"confianca":1.0,"contraexemplos":[],"descricao":"O grafo é magro (nó + aresta seca); deve ser um CONTÍNUO — cada aresta A—rel→B é um segmento do tecido Koch preenchido por uma transição de prosa REAL. Fundamento: o paper de transição (hiper/circular/transicao_lopes.py) modela a passagem entre dois estados como interpolação graduada s:0→1 com ponto crítico, não salto; e a coordenada Koch já permanece contínua. O plano (doc/continuo/plano_transicoes.md) preenche esse s progressivamente, por levas estratégicas, medindo a cobertura — começando pela fatia SymbolicBeta/BAI e pelos hubs. É HIPÓTESE/projeto (o plano; a execução vem nas levas). Atribuição: Lopes; pacto FORA.","epistemico":"hipotese","exemplos":[],"id":"plano_continuo_transicoes","memoria":"semantica","meta":{"autor":"Lopes/broca-so","dominio":"continuo","fonte":"doc/continuo/plano_transicoes.md; hiper/circular/transicao_lopes.py"},"origem":"wiki","palavras_chave":[],"sinonimos":[],"tipo":"conceito","titulo":"Plano: preencher o contínuo do grafo (transições de prosa real nas arestas)"}
\section{Plano: preencher o contínuo do grafo (transições de prosa real nas arestas)}
O grafo é magro (nó + aresta seca); deve ser um CONTÍNUO — cada aresta A—rel\ensuremath{\to}B é um segmento do tecido Koch preenchido por uma transição de prosa REAL. Fundamento: o paper de transição (hiper/circular/transicao\_lopes.py) modela a passagem entre dois estados como interpolação graduada s:0\ensuremath{\to}1 com ponto crítico, não salto; e a coordenada Koch já permanece contínua. O plano (doc/continuo/plano\_transicoes.md) preenche esse s progressivamente, por levas estratégicas, medindo a cobertura — começando pela fatia SymbolicBeta/BAI e pelos hubs. É HIPÓTESE/projeto (o plano; a execução vem nas levas). Atribuição: Lopes; pacto FORA.
\prov{relações: usa transicao; usa flip\_discreto\_continuo; usa 4\_a\_coordenada\_koch\_permanece\_continua; parte\_de aresta\_transicao\_graduada; parte\_de mapear\_nao\_fabricar\_continuo; parte\_de medida\_continuidade\_transicao; parte\_de broca\_os; parte\_de tiffany}
\prov{proveniência: wiki \textperiodcentered{} doc/continuo/plano\_transicoes.md; hiper/circular/transicao\_lopes.py \textperiodcentered{} hipotese \textperiodcentered{} confiança 1.0 \textperiodcentered{} domínio continuo \textperiodcentered{} história 130 versões no jornal}

%CRISTAL {"arestas":[{"alvo":"octonios_hiperbolicos","confianca":1.0,"epistemico":"fato","origem":"paper","rel":"explica"},{"alvo":"psl_2_7","confianca":1.0,"epistemico":"fato","origem":"paper","rel":"relaciona"}],"confianca":1.0,"contraexemplos":[],"descricao":"PG(2,𝔽₂): 7 pontos, 7 retas (trios), cada reta com 3 pontos. Codifica a tabela de multiplicação dos octônios (eᵢeⱼ via os 7 trios orientados); Aut(Fano)=GL(3,𝔽₂)≅PSL(2,7) preserva o produto.","epistemico":"fato","exemplos":[],"id":"plano_fano_octonico","memoria":"semantica","meta":{"autor":"Gentil Lopes","descoberta":true,"dominio":"matematica","fonte":"paper_AX.pdf, Teo.7"},"origem":"paper","palavras_chave":[],"sinonimos":[],"tipo":"conceito","titulo":"Plano de Fano Octônico"}
\section{Plano de Fano Octônico}
PG(2,\ensuremath{\mathbb{F}}\ensuremath{_{2}}): 7 pontos, 7 retas (trios), cada reta com 3 pontos. Codifica a tabela de multiplicação dos octônios (e\ensuremath{_{i}}e\ensuremath{_{j}} via os 7 trios orientados); Aut(Fano)=GL(3,\ensuremath{\mathbb{F}}\ensuremath{_{2}})\ensuremath{\cong}PSL(2,7) preserva o produto.
\prov{relações: explica octonios\_hiperbolicos; relaciona psl\_2\_7}
\prov{proveniência: paper \textperiodcentered{} paper\_AX.pdf, Teo.7 \textperiodcentered{} fato \textperiodcentered{} confiança 1.0 \textperiodcentered{} domínio matematica \textperiodcentered{} história 3 versões no jornal}

%CRISTAL {"arestas":[{"alvo":"fronteira_hurwitz","confianca":1.0,"epistemico":"fato","origem":"paper","rel":"relaciona"}],"confianca":1.0,"contraexemplos":[],"descricao":"Potências de uma única base b associam mesmo em álgebra não-associativa: ba·bc=ba⁺c, vivendo num campo local Cᵦ≅ℂ. A não-associatividade só aparece entre bases distintas.","epistemico":"fato","exemplos":[],"id":"poder_associatividade","memoria":"semantica","meta":{"autor":"Gentil Lopes","descoberta":true,"dominio":"matematica","fonte":"paper_21_irracionais.pdf, Prop.2.1"},"origem":"paper","palavras_chave":[],"sinonimos":[],"tipo":"conceito","titulo":"Power-Associatividade"}
\section{Power-Associatividade}
Potências de uma única base b associam mesmo em álgebra não-associativa: ba\textperiodcentered bc=ba\ensuremath{^{+}}c, vivendo num campo local C\ensuremath{_{\beta}}\ensuremath{\cong}\ensuremath{\mathbb{C}}. A não-associatividade só aparece entre bases distintas.
\prov{relações: relaciona fronteira\_hurwitz}
\prov{proveniência: paper \textperiodcentered{} paper\_21\_irracionais.pdf, Prop.2.1 \textperiodcentered{} fato \textperiodcentered{} confiança 1.0 \textperiodcentered{} domínio matematica \textperiodcentered{} história 3 versões no jornal}

%CRISTAL {"arestas":[{"alvo":"decoerencia_e_sair_da_borda","confianca":1.0,"epistemico":"fato","origem":"wiki","rel":"parte_de"},{"alvo":"numeros_de_pisot","confianca":1.0,"epistemico":"fato","origem":"wiki","rel":"usa"},{"alvo":"colapso","confianca":1.0,"epistemico":"fato","origem":"wiki","rel":"relaciona"},{"alvo":"born_e_o_colapso_do_bai","confianca":1.0,"epistemico":"fato","origem":"wiki","rel":"relaciona"}],"confianca":1.0,"contraexemplos":[],"descricao":"O que sobrevive à decoerência são os POINTER STATES — a base preferida que o ambiente seleciona. Na máquina mineral são o autovetor DOMINANTE (o fundamental, o Pisot) para onde A^k projeta: o mesmo colapso Ψ do BAI. A decoerência escolhe a base clássica exatamente como o tempo imaginário escolhe o estado fundamental. Decoerir e colapsar são o mesmo gesto.","epistemico":"inferencia","exemplos":[],"id":"pointer_states_sao_o_fundamental","memoria":"semantica","meta":{"dominio":"reta mineral","fonte":"emaranhamento e decoerência"},"origem":"wiki","palavras_chave":[],"sinonimos":[],"tipo":"conceito","titulo":"Os Pointer States são o Fundamental (Pisot)"}
\section{Os Pointer States são o Fundamental (Pisot)}
O que sobrevive à decoerência são os POINTER STATES — a base preferida que o ambiente seleciona. Na máquina mineral são o autovetor DOMINANTE (o fundamental, o Pisot) para onde A\^{}k projeta: o mesmo colapso \ensuremath{\Psi} do BAI. A decoerência escolhe a base clássica exatamente como o tempo imaginário escolhe o estado fundamental. Decoerir e colapsar são o mesmo gesto.
\prov{relações: parte\_de decoerencia\_e\_sair\_da\_borda; usa numeros\_de\_pisot; relaciona colapso; relaciona born\_e\_o\_colapso\_do\_bai}
\prov{proveniência: wiki \textperiodcentered{} emaranhamento e decoerência \textperiodcentered{} inferencia \textperiodcentered{} confiança 1.0 \textperiodcentered{} domínio reta mineral \textperiodcentered{} história 5 versões no jornal}

%CRISTAL {"arestas":[{"alvo":"curvatura_algebrica","confianca":1.0,"epistemico":"fato","origem":"paper","rel":"explica"}],"confianca":1.0,"contraexemplos":[],"descricao":"A dualidade GL(V)↔S_n liga a curvatura de Hochschild à decomposição A⊗³≅5V₀⊕9V₁⊕5V₂⊕V₃ (Clebsch-Gordan): a_k/b_k=K₀² (combinatório = algébrico).","epistemico":"fato","exemplos":[],"id":"ponte_schur_weyl","memoria":"semantica","meta":{"autor":"Gentil Lopes","descoberta":false,"dominio":"matematica","fonte":"paper_AE_ponte_schur_weyl.pdf, Teo.5.5"},"origem":"paper","palavras_chave":[],"sinonimos":[],"tipo":"conceito","titulo":"Ponte de Schur-Weyl"}
\section{Ponte de Schur-Weyl}
A dualidade GL(V)\ensuremath{\leftrightarrow}S\_n liga a curvatura de Hochschild à decomposição A\ensuremath{\otimes}\ensuremath{^{3}}\ensuremath{\cong}5V\ensuremath{_{0}}\ensuremath{\oplus}9V\ensuremath{_{1}}\ensuremath{\oplus}5V\ensuremath{_{2}}\ensuremath{\oplus}V\ensuremath{_{3}} (Clebsch-Gordan): a\_k/b\_k=K\ensuremath{_{0}}\ensuremath{^{2}} (combinatório = algébrico).
\prov{relações: explica curvatura\_algebrica}
\prov{proveniência: paper \textperiodcentered{} paper\_AE\_ponte\_schur\_weyl.pdf, Teo.5.5 \textperiodcentered{} fato \textperiodcentered{} confiança 1.0 \textperiodcentered{} domínio matematica \textperiodcentered{} história 2 versões no jornal}

%CRISTAL {"arestas":[{"alvo":"flip_sigma","confianca":1.0,"epistemico":"fato","origem":"paper","rel":"relaciona"},{"alvo":"soma_velocidades","confianca":1.0,"epistemico":"fato","origem":"paper","rel":"usa"},{"alvo":"mecanica_quantica","confianca":1.0,"epistemico":"fato","origem":"paper","rel":"referencia"}],"confianca":1.0,"contraexemplos":[],"descricao":"A rotação c↦ic (continuação analítica) leva a soma circular de velocidades v=tan(θ₁+θ₂) à adição relativística tanh(φ₁+φ₂): geometria circular ↔ Minkowski.","epistemico":"fato","exemplos":[],"id":"ponte_wick","memoria":"semantica","meta":{"autor":"Gentil Lopes","descoberta":true,"dominio":"matematica","fonte":"paper_44_ponte_wick.pdf, Teo.3.1"},"origem":"paper","palavras_chave":[],"sinonimos":[],"tipo":"conceito","titulo":"Ponte de Wick"}
\section{Ponte de Wick}
A rotação c\ensuremath{\mapsto}ic (continuação analítica) leva a soma circular de velocidades v=tan(\ensuremath{\theta}\ensuremath{_{1}}+\ensuremath{\theta}\ensuremath{_{2}}) à adição relativística tanh(\ensuremath{\varphi}\ensuremath{_{1}}+\ensuremath{\varphi}\ensuremath{_{2}}): geometria circular \ensuremath{\leftrightarrow} Minkowski.
\prov{relações: relaciona flip\_sigma; usa soma\_velocidades; referencia mecanica\_quantica}
\prov{proveniência: paper \textperiodcentered{} paper\_44\_ponte\_wick.pdf, Teo.3.1 \textperiodcentered{} fato \textperiodcentered{} confiança 1.0 \textperiodcentered{} domínio matematica \textperiodcentered{} história 4 versões no jornal}

%CRISTAL {"arestas":[{"alvo":"fourier_reconstrucao","confianca":1.0,"epistemico":"fato","origem":"paper","rel":"usa"},{"alvo":"word","confianca":0.5,"epistemico":"fato","origem":"wiki","rel":"relaciona"}],"confianca":1.0,"contraexemplos":[],"descricao":"e−πx² é o ponto fixo da transformada de Fourier: onde discreto e contínuo coincidem (λ₁=λ₂), a cúspide central da lemniscata.","epistemico":"fato","exemplos":[],"id":"ponto_auto_dual","memoria":"semantica","meta":{"autor":"Gentil Lopes","descoberta":false,"dominio":"matematica","fonte":"paper_77_lemniscata.pdf, Teo.4.1"},"origem":"paper","palavras_chave":[],"sinonimos":[],"tipo":"conceito","titulo":"Ponto Auto-Dual (gaussiana)"}
\section{Ponto Auto-Dual (gaussiana)}
e\ensuremath{-}\ensuremath{\pi}x\ensuremath{^{2}} é o ponto fixo da transformada de Fourier: onde discreto e contínuo coincidem (\ensuremath{\lambda}\ensuremath{_{1}}=\ensuremath{\lambda}\ensuremath{_{2}}), a cúspide central da lemniscata.
\prov{relações: usa fourier\_reconstrucao; relaciona word}
\prov{proveniência: paper \textperiodcentered{} paper\_77\_lemniscata.pdf, Teo.4.1 \textperiodcentered{} fato \textperiodcentered{} confiança 1.0 \textperiodcentered{} domínio matematica \textperiodcentered{} história 4 versões no jornal}

%CRISTAL {"arestas":[{"alvo":"furo_hopf","confianca":1.0,"epistemico":"fato","origem":"paper","rel":"usa"},{"alvo":"word","confianca":0.5,"epistemico":"fato","origem":"wiki","rel":"relaciona"},{"alvo":"virtualizacao","confianca":0.5,"epistemico":"fato","origem":"wiki","rel":"relaciona"}],"confianca":1.0,"contraexemplos":[],"descricao":"Compactificação do plano projetivo octoniônico pela projeção estereográfica da Hopf normalizada; o ponto no infinito fecha OP¹≅S⁸.","epistemico":"fato","exemplos":[],"id":"ponto_cayley","memoria":"semantica","meta":{"autor":"Gentil Lopes","descoberta":false,"dominio":"matematica","fonte":"paper_50_octonios_furo.pdf, Teo.3.1"},"origem":"paper","palavras_chave":[],"sinonimos":[],"tipo":"conceito","titulo":"Ponto de Cayley (OP¹≅S⁸)"}
\section{Ponto de Cayley (OP\ensuremath{^{1}}\ensuremath{\cong}S\ensuremath{^{8}})}
Compactificação do plano projetivo octoniônico pela projeção estereográfica da Hopf normalizada; o ponto no infinito fecha OP\ensuremath{^{1}}\ensuremath{\cong}S\ensuremath{^{8}}.
\prov{relações: usa furo\_hopf; relaciona word; relaciona virtualizacao}
\prov{proveniência: paper \textperiodcentered{} paper\_50\_octonios\_furo.pdf, Teo.3.1 \textperiodcentered{} fato \textperiodcentered{} confiança 1.0 \textperiodcentered{} domínio matematica \textperiodcentered{} história 4 versões no jornal}

%CRISTAL {"arestas":[{"alvo":"cuspide_pisot","confianca":1.0,"epistemico":"fato","origem":"paper","rel":"usa"},{"alvo":"medida_parry","confianca":1.0,"epistemico":"fato","origem":"paper","rel":"usa"},{"alvo":"transformada_metalica_fechada_perelman","confianca":1.0,"epistemico":"fato","origem":"wiki","rel":"relaciona"}],"confianca":1.0,"contraexemplos":[],"descricao":"A 'violência aparente' de energia crescente é, normalizada, difusão mansa: a partir da cúspide de Pisot a distribuição relaxa à medida de Parry com taxa dada pelo gap espectral.","epistemico":"fato","exemplos":[],"id":"ponto_difusao","memoria":"semantica","meta":{"autor":"Gentil Lopes","descoberta":false,"dominio":"matematica","fonte":"paper_70_ponto_difusao.pdf, Teo.4.1"},"origem":"paper","palavras_chave":[],"sinonimos":[],"tipo":"conceito","titulo":"Ponto de Difusão"}
\section{Ponto de Difusão}
A 'violência aparente' de energia crescente é, normalizada, difusão mansa: a partir da cúspide de Pisot a distribuição relaxa à medida de Parry com taxa dada pelo gap espectral.
\prov{relações: usa cuspide\_pisot; usa medida\_parry; relaciona transformada\_metalica\_fechada\_perelman}
\prov{proveniência: paper \textperiodcentered{} paper\_70\_ponto\_difusao.pdf, Teo.4.1 \textperiodcentered{} fato \textperiodcentered{} confiança 1.0 \textperiodcentered{} domínio matematica \textperiodcentered{} história 4 versões no jornal}

%CRISTAL {"arestas":[{"alvo":"caca_ao_cervo","confianca":1.0,"epistemico":"fato","origem":"wiki","rel":"usa"}],"confianca":1.0,"contraexemplos":[],"descricao":"Sem poder combinar, pessoas convergem para a solução 'óbvia' por saliência cultural (encontrar-se em NY: meio-dia na Grand Central). A coordenação por expectativa compartilhada, não por cálculo.","epistemico":"inferencia","exemplos":[],"id":"ponto_focal_schelling","memoria":"semantica","meta":{"autor":"Schelling","dominio":"matematica","fonte":"teoria dos jogos"},"origem":"wiki","palavras_chave":[],"sinonimos":[],"tipo":"conceito","titulo":"Ponto Focal (Schelling)"}
\section{Ponto Focal (Schelling)}
Sem poder combinar, pessoas convergem para a solução 'óbvia' por saliência cultural (encontrar-se em NY: meio-dia na Grand Central). A coordenação por expectativa compartilhada, não por cálculo.
\prov{relações: usa caca\_ao\_cervo}
\prov{proveniência: wiki \textperiodcentered{} teoria dos jogos \textperiodcentered{} inferencia \textperiodcentered{} confiança 1.0 \textperiodcentered{} domínio matematica \textperiodcentered{} história 2 versões no jornal}

%CRISTAL {"arestas":[{"alvo":"algebra_posicional_metalica","confianca":1.0,"epistemico":"fato","origem":"wiki","rel":"especializa"},{"alvo":"reta_metalica_geral","confianca":1.0,"epistemico":"fato","origem":"wiki","rel":"usa"},{"alvo":"independencia_retas_metalicas","confianca":1.0,"epistemico":"fato","origem":"wiki","rel":"usa"},{"alvo":"tiffany","confianca":1.0,"epistemico":"fato","origem":"wiki","rel":"parte_de"}],"confianca":1.0,"contraexemplos":[],"descricao":"O metal bronze: n=3, σ₃=(3+√13)/2=3,303 (raiz de x²−3x−1), anel ℤ[σ₃] (inteiros de ℚ(√13)), unidade fundamental σ₃ com N=−1. A sequência bronze (1,3,10,33,109,360) = coordenada de σ₃k; posição k ↔ A₃k, passo (a,b)→(b, 3b+a). Corpo ℚ(√13) — independente de ouro/prata (esterilidade), carrega informação nova. [D] goldchain/metalpos.py(n=3).","epistemico":"fato","exemplos":[],"id":"posicional_bronze","memoria":"semantica","meta":{"autor":"Lopes/broca-so","dominio":"algebra","fonte":"goldchain/metal_pos.py (n=3)"},"origem":"wiki","palavras_chave":[],"sinonimos":[],"tipo":"conceito","titulo":"Bronze (n=3): álgebra posicional de ℤ[σ₃] — 1,3,10,33, ℚ(√13)"}
\section{Bronze (n=3): álgebra posicional de \ensuremath{\mathbb{Z}}[\ensuremath{\sigma}\ensuremath{_{3}}] — 1,3,10,33, \ensuremath{\mathbb{Q}}(\ensuremath{\sqrt{\,}}13)}
O metal bronze: n=3, \ensuremath{\sigma}\ensuremath{_{3}}=(3+\ensuremath{\sqrt{\,}}13)/2=3,303 (raiz de x\ensuremath{^{2}}\ensuremath{-}3x\ensuremath{-}1), anel \ensuremath{\mathbb{Z}}[\ensuremath{\sigma}\ensuremath{_{3}}] (inteiros de \ensuremath{\mathbb{Q}}(\ensuremath{\sqrt{\,}}13)), unidade fundamental \ensuremath{\sigma}\ensuremath{_{3}} com N=\ensuremath{-}1. A sequência bronze (1,3,10,33,109,360) = coordenada de \ensuremath{\sigma}\ensuremath{_{3}}k; posição k \ensuremath{\leftrightarrow} A\ensuremath{_{3}}k, passo (a,b)\ensuremath{\to}(b, 3b+a). Corpo \ensuremath{\mathbb{Q}}(\ensuremath{\sqrt{\,}}13) — independente de ouro/prata (esterilidade), carrega informação nova. [D] goldchain/metalpos.py(n=3).
\prov{relações: especializa algebra\_posicional\_metalica; usa reta\_metalica\_geral; usa independencia\_retas\_metalicas; parte\_de tiffany}
\prov{proveniência: wiki \textperiodcentered{} goldchain/metal\_pos.py (n=3) \textperiodcentered{} fato \textperiodcentered{} confiança 1.0 \textperiodcentered{} domínio algebra \textperiodcentered{} história 6 versões no jornal}

%CRISTAL {"arestas":[{"alvo":"algebra_posicional_metalica","confianca":1.0,"epistemico":"fato","origem":"wiki","rel":"especializa"},{"alvo":"reta_metalica_geral","confianca":1.0,"epistemico":"fato","origem":"wiki","rel":"usa"},{"alvo":"posicional_ouro","confianca":1.0,"epistemico":"fato","origem":"wiki","rel":"relaciona"},{"alvo":"tiffany","confianca":1.0,"epistemico":"fato","origem":"wiki","rel":"parte_de"}],"confianca":1.0,"contraexemplos":[],"descricao":"O metal cobre: n=4, σ₄=(4+√20)/2=2+√5=4,236 (raiz de x²−4x−1), unidade fundamental com N=−1. A sequência cobre (1,4,17,72,305,1292) = coordenada de σ₄k; posição k ↔ A₄k, passo (a,b)→(b, 4b+a). NOTA honesta: o corpo é ℚ(√5) — o MESMO do ouro (n²+4=20=4·5); σ₄∈ℤ[√5] é potência da unidade do ouro. A distinção ouro↔cobre é de ORDEM/passo, não de corpo. [D] goldchain/metalpos.py(n=4).","epistemico":"fato","exemplos":[],"id":"posicional_cobre","memoria":"semantica","meta":{"autor":"Lopes/broca-so","dominio":"algebra","fonte":"goldchain/metal_pos.py (n=4)"},"origem":"wiki","palavras_chave":[],"sinonimos":[],"tipo":"conceito","titulo":"Cobre (n=4): álgebra posicional de ℤ[σ₄] — 1,4,17,72, ℚ(√5) (ordem, não corpo)"}
\section{Cobre (n=4): álgebra posicional de \ensuremath{\mathbb{Z}}[\ensuremath{\sigma}\ensuremath{_{4}}] — 1,4,17,72, \ensuremath{\mathbb{Q}}(\ensuremath{\sqrt{\,}}5) (ordem, não corpo)}
O metal cobre: n=4, \ensuremath{\sigma}\ensuremath{_{4}}=(4+\ensuremath{\sqrt{\,}}20)/2=2+\ensuremath{\sqrt{\,}}5=4,236 (raiz de x\ensuremath{^{2}}\ensuremath{-}4x\ensuremath{-}1), unidade fundamental com N=\ensuremath{-}1. A sequência cobre (1,4,17,72,305,1292) = coordenada de \ensuremath{\sigma}\ensuremath{_{4}}k; posição k \ensuremath{\leftrightarrow} A\ensuremath{_{4}}k, passo (a,b)\ensuremath{\to}(b, 4b+a). NOTA honesta: o corpo é \ensuremath{\mathbb{Q}}(\ensuremath{\sqrt{\,}}5) — o MESMO do ouro (n\ensuremath{^{2}}+4=20=4\textperiodcentered 5); \ensuremath{\sigma}\ensuremath{_{4}}\ensuremath{\in}\ensuremath{\mathbb{Z}}[\ensuremath{\sqrt{\,}}5] é potência da unidade do ouro. A distinção ouro\ensuremath{\leftrightarrow}cobre é de ORDEM/passo, não de corpo. [D] goldchain/metalpos.py(n=4).
\prov{relações: especializa algebra\_posicional\_metalica; usa reta\_metalica\_geral; relaciona posicional\_ouro; parte\_de tiffany}
\prov{proveniência: wiki \textperiodcentered{} goldchain/metal\_pos.py (n=4) \textperiodcentered{} fato \textperiodcentered{} confiança 1.0 \textperiodcentered{} domínio algebra \textperiodcentered{} história 6 versões no jornal}

%CRISTAL {"arestas":[{"alvo":"algebra_posicional_metalica","confianca":1.0,"epistemico":"fato","origem":"wiki","rel":"especializa"},{"alvo":"broca_so_caixa_de_ouro","confianca":1.0,"epistemico":"fato","origem":"wiki","rel":"relaciona"},{"alvo":"reta_metalica_geral","confianca":1.0,"epistemico":"fato","origem":"wiki","rel":"usa"},{"alvo":"tiffany","confianca":1.0,"epistemico":"fato","origem":"wiki","rel":"parte_de"},{"alvo":"beta_numeracao_metalica","confianca":1.0,"epistemico":"fato","origem":"wiki","rel":"especializa"}],"confianca":1.0,"contraexemplos":[],"descricao":"O metal ouro na álgebra posicional metálica: n=1, σ₁=φ=(1+√5)/2 (raiz de x²−x−1), anel ℤ[φ] (inteiros de ℚ(√5)), unidade fundamental φ com N(φ)=−1 (norma a²+ab−b²). A sequência é a de FIBONACCI (1,1,2,3,5,8) = coordenada de φk; posição k ↔ A₁k, passo (a,b)→(b, b+a). É a caixa de ouro que o broca-so implementa (golden mean shift, 5 modos). [D] goldchain/metalpos.py(n=1).","epistemico":"fato","exemplos":[],"id":"posicional_ouro","memoria":"semantica","meta":{"autor":"Lopes/broca-so","dominio":"algebra","fonte":"goldchain/metal_pos.py (n=1); broca_so_caixa_de_ouro"},"origem":"wiki","palavras_chave":[],"sinonimos":[],"tipo":"conceito","titulo":"Ouro (n=1): álgebra posicional de ℤ[φ] — Fibonacci, ℚ(√5)"}
\section{Ouro (n=1): álgebra posicional de \ensuremath{\mathbb{Z}}[\ensuremath{\varphi}] — Fibonacci, \ensuremath{\mathbb{Q}}(\ensuremath{\sqrt{\,}}5)}
O metal ouro na álgebra posicional metálica: n=1, \ensuremath{\sigma}\ensuremath{_{1}}=\ensuremath{\varphi}=(1+\ensuremath{\sqrt{\,}}5)/2 (raiz de x\ensuremath{^{2}}\ensuremath{-}x\ensuremath{-}1), anel \ensuremath{\mathbb{Z}}[\ensuremath{\varphi}] (inteiros de \ensuremath{\mathbb{Q}}(\ensuremath{\sqrt{\,}}5)), unidade fundamental \ensuremath{\varphi} com N(\ensuremath{\varphi})=\ensuremath{-}1 (norma a\ensuremath{^{2}}+ab\ensuremath{-}b\ensuremath{^{2}}). A sequência é a de FIBONACCI (1,1,2,3,5,8) = coordenada de \ensuremath{\varphi}k; posição k \ensuremath{\leftrightarrow} A\ensuremath{_{1}}k, passo (a,b)\ensuremath{\to}(b, b+a). É a caixa de ouro que o broca-so implementa (golden mean shift, 5 modos). [D] goldchain/metalpos.py(n=1).
\prov{relações: especializa algebra\_posicional\_metalica; relaciona broca\_so\_caixa\_de\_ouro; usa reta\_metalica\_geral; parte\_de tiffany; especializa beta\_numeracao\_metalica}
\prov{proveniência: wiki \textperiodcentered{} goldchain/metal\_pos.py (n=1); broca\_so\_caixa\_de\_ouro \textperiodcentered{} fato \textperiodcentered{} confiança 1.0 \textperiodcentered{} domínio algebra \textperiodcentered{} história 7 versões no jornal}

%CRISTAL {"arestas":[{"alvo":"algebra_posicional_plastica","confianca":1.0,"epistemico":"fato","origem":"wiki","rel":"especializa"},{"alvo":"reta_plastica_cubica","confianca":1.0,"epistemico":"fato","origem":"wiki","rel":"usa"},{"alvo":"numeros_pisot","confianca":1.0,"epistemico":"fato","origem":"wiki","rel":"usa"},{"alvo":"tiffany","confianca":1.0,"epistemico":"fato","origem":"wiki","rel":"parte_de"}],"confianca":1.0,"contraexemplos":[],"descricao":"A instância base da reta plástica: ρ=1,3247, raiz de x³=x+1 (o MENOR número de Pisot), unidade de ℤ[ρ] (ρ⁻¹=ρ²−1). Sequência de PADOVAN/Perrin = coordenada de ρk; passo cúbico (a,b,c)→(b,c,a+b) mod 2¹²⁸, O(1); coefs=[1,1,0]. h=ln ρ=0,281 (o passo geodésico mais lento — a reta mais 'fria'). [D] goldchain/metalpos.py (passoplastico).","epistemico":"fato","exemplos":[],"id":"posicional_plastico","memoria":"semantica","meta":{"autor":"Lopes/broca-so","dominio":"algebra","fonte":"goldchain/metal_pos.py PLASTICOS['plastico']"},"origem":"wiki","palavras_chave":[],"sinonimos":[],"tipo":"conceito","titulo":"Plástico (ρ, x³=x+1): álgebra posicional cúbica — Padovan/Perrin, menor Pisot"}
\section{Plástico (\ensuremath{\rho}, x\ensuremath{^{3}}=x+1): álgebra posicional cúbica — Padovan/Perrin, menor Pisot}
A instância base da reta plástica: \ensuremath{\rho}=1,3247, raiz de x\ensuremath{^{3}}=x+1 (o MENOR número de Pisot), unidade de \ensuremath{\mathbb{Z}}[\ensuremath{\rho}] (\ensuremath{\rho}\ensuremath{^{-1}}=\ensuremath{\rho}\ensuremath{^{2}}\ensuremath{-}1). Sequência de PADOVAN/Perrin = coordenada de \ensuremath{\rho}k; passo cúbico (a,b,c)\ensuremath{\to}(b,c,a+b) mod 2\ensuremath{^{128}}, O(1); coefs=[1,1,0]. h=ln \ensuremath{\rho}=0,281 (o passo geodésico mais lento — a reta mais 'fria'). [D] goldchain/metalpos.py (passoplastico).
\prov{relações: especializa algebra\_posicional\_plastica; usa reta\_plastica\_cubica; usa numeros\_pisot; parte\_de tiffany}
\prov{proveniência: wiki \textperiodcentered{} goldchain/metal\_pos.py PLASTICOS['plastico'] \textperiodcentered{} fato \textperiodcentered{} confiança 1.0 \textperiodcentered{} domínio algebra \textperiodcentered{} história 6 versões no jornal}

%CRISTAL {"arestas":[{"alvo":"algebra_posicional_plastica","confianca":1.0,"epistemico":"fato","origem":"wiki","rel":"especializa"},{"alvo":"reta_plastica_cubica","confianca":1.0,"epistemico":"fato","origem":"wiki","rel":"usa"},{"alvo":"tiffany","confianca":1.0,"epistemico":"fato","origem":"wiki","rel":"parte_de"}],"confianca":1.0,"contraexemplos":[],"descricao":"ψ=1,4656, raiz de x³=x²+1 (o número supergolden), unidade cúbica; a sequência de NARAYANA (das vacas) é a coordenada de ψk; passo (a,b,c)→(b,c,a+c) mod 2¹²⁸, coefs=[1,0,1]. Outra reta da família plástica (grau 3), independente. [D] goldchain/metalpos.py.","epistemico":"fato","exemplos":[],"id":"posicional_supergolden","memoria":"semantica","meta":{"autor":"Lopes/broca-so","dominio":"algebra","fonte":"goldchain/metal_pos.py PLASTICOS['supergolden']"},"origem":"wiki","palavras_chave":[],"sinonimos":[],"tipo":"conceito","titulo":"Supergolden (ψ, x³=x²+1): álgebra posicional cúbica — Narayana"}
\section{Supergolden (\ensuremath{\psi}, x\ensuremath{^{3}}=x\ensuremath{^{2}}+1): álgebra posicional cúbica — Narayana}
\ensuremath{\psi}=1,4656, raiz de x\ensuremath{^{3}}=x\ensuremath{^{2}}+1 (o número supergolden), unidade cúbica; a sequência de NARAYANA (das vacas) é a coordenada de \ensuremath{\psi}k; passo (a,b,c)\ensuremath{\to}(b,c,a+c) mod 2\ensuremath{^{128}}, coefs=[1,0,1]. Outra reta da família plástica (grau 3), independente. [D] goldchain/metalpos.py.
\prov{relações: especializa algebra\_posicional\_plastica; usa reta\_plastica\_cubica; parte\_de tiffany}
\prov{proveniência: wiki \textperiodcentered{} goldchain/metal\_pos.py PLASTICOS['supergolden'] \textperiodcentered{} fato \textperiodcentered{} confiança 1.0 \textperiodcentered{} domínio algebra \textperiodcentered{} história 5 versões no jornal}

%CRISTAL {"arestas":[{"alvo":"algebra_posicional_plastica","confianca":1.0,"epistemico":"fato","origem":"wiki","rel":"especializa"},{"alvo":"reta_plastica_cubica","confianca":1.0,"epistemico":"fato","origem":"wiki","rel":"usa"},{"alvo":"tiffany","confianca":1.0,"epistemico":"fato","origem":"wiki","rel":"parte_de"}],"confianca":1.0,"contraexemplos":[],"descricao":"A constante tribonacci=1,8393, raiz de x³=x²+x+1, unidade cúbica; a sequência TRIBONACCI (0,0,1,1,2,4,7,13,24…) é a coordenada; passo (a,b,c)→(b,c,a+b+c) mod 2¹²⁸, coefs=[1,1,1]. h=ln=0,608 (a mais 'quente' das três cúbicas). É o análogo cúbico de Fibonacci. [D] goldchain/metal_pos.py.","epistemico":"fato","exemplos":[],"id":"posicional_tribonacci","memoria":"semantica","meta":{"autor":"Lopes/broca-so","dominio":"algebra","fonte":"goldchain/metal_pos.py PLASTICOS['tribonacci']"},"origem":"wiki","palavras_chave":[],"sinonimos":[],"tipo":"conceito","titulo":"Tribonacci (x³=x²+x+1): álgebra posicional cúbica — a constante tribonacci"}
\section{Tribonacci (x\ensuremath{^{3}}=x\ensuremath{^{2}}+x+1): álgebra posicional cúbica — a constante tribonacci}
A constante tribonacci=1,8393, raiz de x\ensuremath{^{3}}=x\ensuremath{^{2}}+x+1, unidade cúbica; a sequência TRIBONACCI (0,0,1,1,2,4,7,13,24…) é a coordenada; passo (a,b,c)\ensuremath{\to}(b,c,a+b+c) mod 2\ensuremath{^{128}}, coefs=[1,1,1]. h=ln=0,608 (a mais 'quente' das três cúbicas). É o análogo cúbico de Fibonacci. [D] goldchain/metal\_pos.py.
\prov{relações: especializa algebra\_posicional\_plastica; usa reta\_plastica\_cubica; parte\_de tiffany}
\prov{proveniência: wiki \textperiodcentered{} goldchain/metal\_pos.py PLASTICOS['tribonacci'] \textperiodcentered{} fato \textperiodcentered{} confiança 1.0 \textperiodcentered{} domínio algebra \textperiodcentered{} história 4 versões no jornal}

%CRISTAL {"arestas":[{"alvo":"motor_iterado","confianca":1.0,"epistemico":"fato","origem":"paper","rel":"parte_de"},{"alvo":"word","confianca":0.5,"epistemico":"fato","origem":"wiki","rel":"relaciona"},{"alvo":"virtualizacao","confianca":0.5,"epistemico":"fato","origem":"wiki","rel":"relaciona"}],"confianca":1.0,"contraexemplos":[],"descricao":"G(s)=limₙ (1/2ⁿ)ln|zₙ|: nulo dentro do Mandelbrot, positivo fora; é o potencial eletrostático do fractal (capacidade logarítmica).","epistemico":"fato","exemplos":[],"id":"potencial_fractal","memoria":"semantica","meta":{"autor":"Gentil Lopes","descoberta":false,"dominio":"matematica","fonte":"paper_18_inversao_torre.pdf, Def.5.1"},"origem":"paper","palavras_chave":[],"sinonimos":[],"tipo":"conceito","titulo":"Potencial do Fractal (Green)"}
\section{Potencial do Fractal (Green)}
G(s)=lim\ensuremath{_{n}} (1/2\ensuremath{^{n}})ln|z\ensuremath{_{n}}|: nulo dentro do Mandelbrot, positivo fora; é o potencial eletrostático do fractal (capacidade logarítmica).
\prov{relações: parte\_de motor\_iterado; relaciona word; relaciona virtualizacao}
\prov{proveniência: paper \textperiodcentered{} paper\_18\_inversao\_torre.pdf, Def.5.1 \textperiodcentered{} fato \textperiodcentered{} confiança 1.0 \textperiodcentered{} domínio matematica \textperiodcentered{} história 4 versões no jornal}

%CRISTAL {"arestas":[{"alvo":"qubit_hopf","confianca":1.0,"epistemico":"fato","origem":"paper","rel":"usa"}],"confianca":1.0,"contraexemplos":[],"descricao":"A dinâmica de um qubit é de 1ª ordem (iħψ̇=Hψ, H=½π σ_z B em SU(2)): gira o vetor sem mudar a norma — precessão de Larmor na esfera de Bloch, não oscilador clássico.","epistemico":"fato","exemplos":[],"id":"precessao_bloch","memoria":"semantica","meta":{"autor":"Gentil Lopes","descoberta":false,"dominio":"matematica","fonte":"mecanica.pdf"},"origem":"paper","palavras_chave":[],"sinonimos":[],"tipo":"conceito","titulo":"Precessão de Bloch"}
\section{Precessão de Bloch}
A dinâmica de um qubit é de 1ª ordem (i\ensuremath{\hbar}\ensuremath{\psi}\ensuremath{}=H\ensuremath{\psi}, H=\ensuremath{\tfrac12}\ensuremath{\pi} \ensuremath{\sigma}\_z B em SU(2)): gira o vetor sem mudar a norma — precessão de Larmor na esfera de Bloch, não oscilador clássico.
\prov{relações: usa qubit\_hopf}
\prov{proveniência: paper \textperiodcentered{} mecanica.pdf \textperiodcentered{} fato \textperiodcentered{} confiança 1.0 \textperiodcentered{} domínio matematica \textperiodcentered{} história 2 versões no jornal}

%CRISTAL {"arestas":[{"alvo":"mass_gap_algebrico","confianca":1.0,"epistemico":"fato","origem":"paper","rel":"referencia"}],"confianca":1.0,"contraexemplos":[],"descricao":"Dada uma aproximação σ do autovalor, um verificador roda em tempo poli(κ) via a fórmula do traço de Selberg truncada — assimetria cripto: minerar custa O(κ³), verificar O(κ²).","epistemico":"fato","exemplos":[],"id":"precmaass_verificacao","memoria":"semantica","meta":{"autor":"Gentil Lopes","descoberta":false,"dominio":"matematica","fonte":"paper_PQ1_ym_maass_hardness.pdf, Lema 3.2"},"origem":"paper","palavras_chave":[],"sinonimos":[],"tipo":"conceito","titulo":"Verificação Polinomial de PrecMaass"}
\section{Verificação Polinomial de PrecMaass}
Dada uma aproximação \ensuremath{\sigma} do autovalor, um verificador roda em tempo poli(\ensuremath{\kappa}) via a fórmula do traço de Selberg truncada — assimetria cripto: minerar custa O(\ensuremath{\kappa}\ensuremath{^{3}}), verificar O(\ensuremath{\kappa}\ensuremath{^{2}}).
\prov{relações: referencia mass\_gap\_algebrico}
\prov{proveniência: paper \textperiodcentered{} paper\_PQ1\_ym\_maass\_hardness.pdf, Lema 3.2 \textperiodcentered{} fato \textperiodcentered{} confiança 1.0 \textperiodcentered{} domínio matematica \textperiodcentered{} história 2 versões no jornal}

%CRISTAL {"arestas":[{"alvo":"sequencia_pell","confianca":1.0,"epistemico":"fato","origem":"paper","rel":"especializa"},{"alvo":"razao_de_prata","confianca":1.0,"epistemico":"fato","origem":"paper","rel":"usa"},{"alvo":"hexagrama_mineracao","confianca":1.0,"epistemico":"fato","origem":"paper","rel":"explica"}],"confianca":1.0,"contraexemplos":[],"descricao":"Os primos da sequência de Pell ocorrem em índices primos (primo no índice e no valor): P₂,P₃,P₅,P₁₁,P₁₃,… = 2,5,29,5741,33461,…, crescimento ~σ₂ᵏ — a matéria-prima da mineração de prata.","epistemico":"fato","exemplos":[],"id":"primos_prata","memoria":"semantica","meta":{"autor":"Gentil Lopes","descoberta":true,"dominio":"matematica","fonte":"cifra_ouro_branco.pdf, Prop.1"},"origem":"paper","palavras_chave":[],"sinonimos":[],"tipo":"conceito","titulo":"Primos de Prata"}
\section{Primos de Prata}
Os primos da sequência de Pell ocorrem em índices primos (primo no índice e no valor): P\ensuremath{_{2}},P\ensuremath{_{3}},P\ensuremath{_{5}},P\ensuremath{_{11}},P\ensuremath{_{13}},… = 2,5,29,5741,33461,…, crescimento \~{}\ensuremath{\sigma}\ensuremath{_{2}}\ensuremath{^{k}} — a matéria-prima da mineração de prata.
\prov{relações: especializa sequencia\_pell; usa razao\_de\_prata; explica hexagrama\_mineracao}
\prov{proveniência: paper \textperiodcentered{} cifra\_ouro\_branco.pdf, Prop.1 \textperiodcentered{} fato \textperiodcentered{} confiança 1.0 \textperiodcentered{} domínio matematica \textperiodcentered{} história 4 versões no jornal}

%CRISTAL {"arestas":[{"alvo":"progressao_geometrica_ordem_m","confianca":1.0,"epistemico":"fato","origem":"livro","rel":"explica"}],"confianca":1.0,"contraexemplos":[],"descricao":"Álgebra de Boole: se os axiomas implicam seus duais, o dual de todo teorema também é teorema. Aplicado a progressões: cada propriedade de P.A. tem uma dual em P.G.","epistemico":"fato","exemplos":[],"id":"principio_dualidade","memoria":"semantica","meta":{"autor":"Gentil Lopes","descoberta":false,"dominio":"matematica","fonte":"Novas_Sequencias_Aritmeticas_e_Geometric.pdf, cap.3"},"origem":"livro","palavras_chave":[],"sinonimos":[],"tipo":"conceito","titulo":"Princípio da Dualidade"}
\section{Princípio da Dualidade}
Álgebra de Boole: se os axiomas implicam seus duais, o dual de todo teorema também é teorema. Aplicado a progressões: cada propriedade de P.A. tem uma dual em P.G.
\prov{relações: explica progressao\_geometrica\_ordem\_m}
\prov{proveniência: livro \textperiodcentered{} Novas\_Sequencias\_Aritmeticas\_e\_Geometric.pdf, cap.3 \textperiodcentered{} fato \textperiodcentered{} confiança 1.0 \textperiodcentered{} domínio matematica \textperiodcentered{} história 4 versões no jornal}

%CRISTAL {"arestas":[{"alvo":"corpo_estelar","confianca":1.0,"epistemico":"fato","origem":"wiki","rel":"parte_de"},{"alvo":"mecanismo_de_la_hire","confianca":1.0,"epistemico":"fato","origem":"wiki","rel":"eh_um"},{"alvo":"escada_de_hipociclos","confianca":1.0,"epistemico":"fato","origem":"wiki","rel":"usa"},{"alvo":"cf_de_pi_e_produto_de_gatos","confianca":1.0,"epistemico":"fato","origem":"wiki","rel":"relaciona"},{"alvo":"gato_generaliza_matriz_companion","confianca":1.0,"epistemico":"fato","origem":"wiki","rel":"relaciona"}],"confianca":1.0,"contraexemplos":[],"descricao":"A multiplicação do corpo estelar. Uma estrela impõe ao ângulo do pino o mapa θ↦ν·θ (ν=k−1 = o enrolamento). COMPOR dois desses mapas — θ↦ν₁θ seguido de θ↦ν₂θ — dá θ↦ν₁ν₂θ: as rotações de La Hire se compõem MULTIPLICANDO a frequência. Neutro: a RETA (ν=1, k=2, o mapa identidade). É o que faltava à reta mineral (que só somava): agora podemos MULTIPLICAR estrelas, inclusive irracionais (√2⊗√3=√6).","epistemico":"fato","exemplos":[],"id":"produto_de_la_hire","memoria":"semantica","meta":{"dominio":"algebra","fonte":"corpo estelar"},"origem":"wiki","palavras_chave":[],"sinonimos":[],"tipo":"conceito","titulo":"O Produto de La Hire (⊗)"}
\section{O Produto de La Hire (\ensuremath{\otimes})}
A multiplicação do corpo estelar. Uma estrela impõe ao ângulo do pino o mapa \ensuremath{\theta}\ensuremath{\mapsto}\ensuremath{\nu}\textperiodcentered \ensuremath{\theta} (\ensuremath{\nu}=k\ensuremath{-}1 = o enrolamento). COMPOR dois desses mapas — \ensuremath{\theta}\ensuremath{\mapsto}\ensuremath{\nu}\ensuremath{_{1}}\ensuremath{\theta} seguido de \ensuremath{\theta}\ensuremath{\mapsto}\ensuremath{\nu}\ensuremath{_{2}}\ensuremath{\theta} — dá \ensuremath{\theta}\ensuremath{\mapsto}\ensuremath{\nu}\ensuremath{_{1}}\ensuremath{\nu}\ensuremath{_{2}}\ensuremath{\theta}: as rotações de La Hire se compõem MULTIPLICANDO a frequência. Neutro: a RETA (\ensuremath{\nu}=1, k=2, o mapa identidade). É o que faltava à reta mineral (que só somava): agora podemos MULTIPLICAR estrelas, inclusive irracionais (\ensuremath{\sqrt{\,}}2\ensuremath{\otimes}\ensuremath{\sqrt{\,}}3=\ensuremath{\sqrt{\,}}6).
\prov{relações: parte\_de corpo\_estelar; eh\_um mecanismo\_de\_la\_hire; usa escada\_de\_hipociclos; relaciona cf\_de\_pi\_e\_produto\_de\_gatos; relaciona gato\_generaliza\_matriz\_companion}
\prov{proveniência: wiki \textperiodcentered{} corpo estelar \textperiodcentered{} fato \textperiodcentered{} confiança 1.0 \textperiodcentered{} domínio algebra \textperiodcentered{} história 8 versões no jornal}

%CRISTAL {"arestas":[{"alvo":"termo_kapitza_ns","confianca":1.0,"epistemico":"fato","origem":"paper","rel":"usa"},{"alvo":"ns_completa_a0","confianca":1.0,"epistemico":"fato","origem":"paper","rel":"usa"},{"alvo":"matematica","confianca":1.0,"epistemico":"fato","origem":"paper","rel":"parte_de"}],"confianca":1.0,"contraexemplos":[],"descricao":"Síntese honesta: seis bloqueios estruturais independentes (rescaling, gradiente, Dirac trivial, excesso sem sinal, Kapitza, kernel solenoidal) mostram que a abordagem ALGÉBRICA não regulariza NS — todos ortogonais aos vetores físicos. Resultado negativo rigoroso.","epistemico":"fato","exemplos":[],"id":"programa_ns_bloqueios","memoria":"semantica","meta":{"autor":"Gentil Lopes","descoberta":false,"dominio":"matematica","fonte":"paper_AH.pdf, §7; paper_AJ.pdf"},"origem":"paper","palavras_chave":[],"sinonimos":[],"tipo":"conceito","titulo":"Seis Bloqueios de NS (resultado negativo)"}
\section{Seis Bloqueios de NS (resultado negativo)}
Síntese honesta: seis bloqueios estruturais independentes (rescaling, gradiente, Dirac trivial, excesso sem sinal, Kapitza, kernel solenoidal) mostram que a abordagem ALGÉBRICA não regulariza NS — todos ortogonais aos vetores físicos. Resultado negativo rigoroso.
\prov{relações: usa termo\_kapitza\_ns; usa ns\_completa\_a0; parte\_de matematica}
\prov{proveniência: paper \textperiodcentered{} paper\_AH.pdf, §7; paper\_AJ.pdf \textperiodcentered{} fato \textperiodcentered{} confiança 1.0 \textperiodcentered{} domínio matematica \textperiodcentered{} história 4 versões no jornal}

%CRISTAL {"arestas":[{"alvo":"sequencia","confianca":1.0,"epistemico":"fato","origem":"livro","rel":"especializa"},{"alvo":"word","confianca":0.5,"epistemico":"fato","origem":"wiki","rel":"relaciona"},{"alvo":"vontade_de_poder","confianca":0.5,"epistemico":"fato","origem":"wiki","rel":"relaciona"}],"confianca":1.0,"contraexemplos":[],"descricao":"Sequência com diferença constante r entre termos consecutivos: aₙ = a₁ + (n−1)r.","epistemico":"fato","exemplos":[],"id":"progressao_aritmetica","memoria":"semantica","meta":{"autor":"Gentil Lopes","descoberta":false,"dominio":"matematica","fonte":""},"origem":"manual","palavras_chave":["diferença","linear"],"sinonimos":["PA"],"tipo":"conceito","titulo":"Progressão Aritmética (P.A.)"}
\section{Progressão Aritmética (P.A.)}
Sequência com diferença constante r entre termos consecutivos: a\ensuremath{_{n}} = a\ensuremath{_{1}} + (n\ensuremath{-}1)r.
\prov{sinónimos: PA}
\prov{relações: especializa sequencia; relaciona word; relaciona vontade\_de\_poder}
\prov{proveniência: manual \textperiodcentered{} fato \textperiodcentered{} confiança 1.0 \textperiodcentered{} domínio matematica \textperiodcentered{} história 6 versões no jornal}

%CRISTAL {"arestas":[{"alvo":"progressao_aritmetica","confianca":1.0,"epistemico":"fato","origem":"livro","rel":"generaliza"},{"alvo":"triangulo_de_pascal","confianca":1.0,"epistemico":"fato","origem":"livro","rel":"depende"}],"confianca":1.0,"contraexemplos":[],"descricao":"Generalização inédita de Lopes: recorrência aₙₘ = a₍ₙ₋₁₎ₘ + a₍ₙ₋₁₎₍ₘ₋₁₎, com termo geral por coeficientes binomiais. Unifica sequências cujo termo geral é polinômio de qualquer grau.","epistemico":"fato","exemplos":[],"id":"progressao_aritmetica_ordem_m","memoria":"semantica","meta":{"autor":"Gentil Lopes","descoberta":true,"dominio":"matematica","fonte":"Novas_Sequencias_Aritmeticas_e_Geometric.pdf, cap.1"},"origem":"livro","palavras_chave":[],"sinonimos":[],"tipo":"conceito","titulo":"Progressão Aritmética de Ordem m (P.A.m)"}
\section{Progressão Aritmética de Ordem m (P.A.m)}
Generalização inédita de Lopes: recorrência a\ensuremath{_{nm}} = a\ensuremath{_{(n-1)m}} + a\ensuremath{_{(n-1)(m-1)}}, com termo geral por coeficientes binomiais. Unifica sequências cujo termo geral é polinômio de qualquer grau.
\prov{relações: generaliza progressao\_aritmetica; depende triangulo\_de\_pascal}
\prov{proveniência: livro \textperiodcentered{} Novas\_Sequencias\_Aritmeticas\_e\_Geometric.pdf, cap.1 \textperiodcentered{} fato \textperiodcentered{} confiança 1.0 \textperiodcentered{} domínio matematica \textperiodcentered{} história 8 versões no jornal}

%CRISTAL {"arestas":[{"alvo":"sequencia","confianca":1.0,"epistemico":"fato","origem":"livro","rel":"especializa"},{"alvo":"word","confianca":0.5,"epistemico":"fato","origem":"wiki","rel":"relaciona"},{"alvo":"transcricao","confianca":0.5,"epistemico":"fato","origem":"wiki","rel":"relaciona"},{"alvo":"virtualizacao","confianca":0.5,"epistemico":"fato","origem":"wiki","rel":"relaciona"},{"alvo":"vontade_de_poder","confianca":0.5,"epistemico":"fato","origem":"wiki","rel":"relaciona"},{"alvo":"veredito_experimental","confianca":0.5,"epistemico":"fato","origem":"wiki","rel":"relaciona"}],"confianca":1.0,"contraexemplos":[],"descricao":"Sequência com razão constante q: aₙ = a₁·qⁿ⁻¹.","epistemico":"fato","exemplos":[],"id":"progressao_geometrica","memoria":"semantica","meta":{"autor":"Gentil Lopes","descoberta":false,"dominio":"matematica","fonte":""},"origem":"manual","palavras_chave":["razão","exponencial"],"sinonimos":["PG"],"tipo":"conceito","titulo":"Progressão Geométrica (P.G.)"}
\section{Progressão Geométrica (P.G.)}
Sequência com razão constante q: a\ensuremath{_{n}} = a\ensuremath{_{1}}\textperiodcentered q\ensuremath{^{n-1}}.
\prov{sinónimos: PG}
\prov{relações: especializa sequencia; relaciona word; relaciona transcricao; relaciona virtualizacao; relaciona vontade\_de\_poder; relaciona veredito\_experimental}
\prov{proveniência: manual \textperiodcentered{} fato \textperiodcentered{} confiança 1.0 \textperiodcentered{} domínio matematica \textperiodcentered{} história 9 versões no jornal}

%CRISTAL {"arestas":[{"alvo":"progressao_geometrica","confianca":1.0,"epistemico":"fato","origem":"livro","rel":"generaliza"},{"alvo":"progressao_aritmetica_ordem_m","confianca":1.0,"epistemico":"fato","origem":"livro","rel":"dualidade"}],"confianca":1.0,"contraexemplos":[],"descricao":"Dual geométrico da P.A.m (Lopes): produto no lugar da soma; termo geral por produtório com expoentes binomiais. Correspondência via Princípio da Dualidade.","epistemico":"fato","exemplos":[],"id":"progressao_geometrica_ordem_m","memoria":"semantica","meta":{"autor":"Gentil Lopes","descoberta":true,"dominio":"matematica","fonte":"Novas_Sequencias_Aritmeticas_e_Geometric.pdf, cap.3"},"origem":"livro","palavras_chave":[],"sinonimos":[],"tipo":"conceito","titulo":"Progressão Geométrica de Ordem m (P.G.m)"}
\section{Progressão Geométrica de Ordem m (P.G.m)}
Dual geométrico da P.A.m (Lopes): produto no lugar da soma; termo geral por produtório com expoentes binomiais. Correspondência via Princípio da Dualidade.
\prov{relações: generaliza progressao\_geometrica; dualidade progressao\_aritmetica\_ordem\_m}
\prov{proveniência: livro \textperiodcentered{} Novas\_Sequencias\_Aritmeticas\_e\_Geometric.pdf, cap.3 \textperiodcentered{} fato \textperiodcentered{} confiança 1.0 \textperiodcentered{} domínio matematica \textperiodcentered{} história 9 versões no jornal}

%CRISTAL {"arestas":[{"alvo":"reticulados","confianca":1.0,"epistemico":"fato","origem":"paper","rel":"referencia"},{"alvo":"word","confianca":0.5,"epistemico":"fato","origem":"wiki","rel":"relaciona"},{"alvo":"virtualizacao","confianca":0.5,"epistemico":"fato","origem":"wiki","rel":"relaciona"},{"alvo":"vida-morte","confianca":0.5,"epistemico":"fato","origem":"wiki","rel":"relaciona"}],"confianca":1.0,"contraexemplos":[],"descricao":"Grupo simples de ordem 168 = SL(2,𝔽₇)/±I ≅ GL(3,𝔽₂) ≅ Aut(plano de Fano) ≅ Aut(Klein quartic); é M₂₁ na cadeia de Mathieu.","epistemico":"fato","exemplos":[],"id":"psl_2_7","memoria":"semantica","meta":{"autor":"Gentil Lopes","descoberta":false,"dominio":"matematica","fonte":"paper_AX.pdf (Klein 1879)"},"origem":"paper","palavras_chave":[],"sinonimos":[],"tipo":"conceito","titulo":"PSL(2,7)"}
\section{PSL(2,7)}
Grupo simples de ordem 168 = SL(2,\ensuremath{\mathbb{F}}\ensuremath{_{7}})/$\pm$I \ensuremath{\cong} GL(3,\ensuremath{\mathbb{F}}\ensuremath{_{2}}) \ensuremath{\cong} Aut(plano de Fano) \ensuremath{\cong} Aut(Klein quartic); é M\ensuremath{_{21}} na cadeia de Mathieu.
\prov{relações: referencia reticulados; relaciona word; relaciona virtualizacao; relaciona vida-morte}
\prov{proveniência: paper \textperiodcentered{} paper\_AX.pdf (Klein 1879) \textperiodcentered{} fato \textperiodcentered{} confiança 1.0 \textperiodcentered{} domínio matematica \textperiodcentered{} história 6 versões no jornal}

%CRISTAL {"arestas":[{"alvo":"catedral_sigma_aritmetica","confianca":1.0,"epistemico":"fato","origem":"wiki","rel":"explica"},{"alvo":"psl_2_7","confianca":1.0,"epistemico":"fato","origem":"wiki","rel":"usa"}],"confianca":1.0,"contraexemplos":[],"descricao":"A leitura interpretativa da obra: a simetria (PSL(2,7), Fano) é a 'linguagem' e o E8 a álgebra excepcional que a encarna; a σ-equivalência seria a 'ponte' e a unidade a 'coroa'. Uma tese estética-estrutural sobre por que essas constantes se alinham.","epistemico":"inferencia","exemplos":[],"id":"psl_e_linguagem","memoria":"semantica","meta":{"autor":"Lord Index (Shawn R. Schiller)","dominio":"matematica","fonte":"A Catedral — Part VIII-A (jul 2026), ~/Imagens/catedral"},"origem":"wiki","palavras_chave":[],"sinonimos":[],"tipo":"conceito","titulo":"PSL é linguagem, E8 é a álgebra: a tese da Catedral"}
\section{PSL é linguagem, E8 é a álgebra: a tese da Catedral}
A leitura interpretativa da obra: a simetria (PSL(2,7), Fano) é a 'linguagem' e o E8 a álgebra excepcional que a encarna; a \ensuremath{\sigma}-equivalência seria a 'ponte' e a unidade a 'coroa'. Uma tese estética-estrutural sobre por que essas constantes se alinham.
\prov{relações: explica catedral\_sigma\_aritmetica; usa psl\_2\_7}
\prov{proveniência: wiki \textperiodcentered{} A Catedral — Part VIII-A (jul 2026), \~{}/Imagens/catedral \textperiodcentered{} inferencia \textperiodcentered{} confiança 1.0 \textperiodcentered{} domínio matematica \textperiodcentered{} história 3 versões no jornal}

%CRISTAL {"arestas":[{"alvo":"psl_2_7","confianca":1.0,"epistemico":"fato","origem":"paper","rel":"usa"},{"alvo":"teoria_gauge_octonica","confianca":1.0,"epistemico":"fato","origem":"paper","rel":"relaciona"}],"confianca":1.0,"contraexemplos":[],"descricao":"PSL(2,7) mergulha em G₂=Aut(𝕆) como subgrupo finito maximal que preserva o plano de Fano orientado (geradores explícitos R,S; Cohen-Wales 1983).","epistemico":"fato","exemplos":[],"id":"psl_embeds_g2","memoria":"semantica","meta":{"autor":"Gentil Lopes","descoberta":false,"dominio":"matematica","fonte":"paper_AX.pdf, Teo.2"},"origem":"paper","palavras_chave":[],"sinonimos":[],"tipo":"conceito","titulo":"PSL(2,7) ⊂ G₂"}
\section{PSL(2,7) \ensuremath{\subset} G\ensuremath{_{2}}}
PSL(2,7) mergulha em G\ensuremath{_{2}}=Aut(\ensuremath{\mathbb{O}}) como subgrupo finito maximal que preserva o plano de Fano orientado (geradores explícitos R,S; Cohen-Wales 1983).
\prov{relações: usa psl\_2\_7; relaciona teoria\_gauge\_octonica}
\prov{proveniência: paper \textperiodcentered{} paper\_AX.pdf, Teo.2 \textperiodcentered{} fato \textperiodcentered{} confiança 1.0 \textperiodcentered{} domínio matematica \textperiodcentered{} história 3 versões no jornal}

%CRISTAL {"arestas":[{"alvo":"redencao_purgatorio_grafo","confianca":1.0,"epistemico":"fato","origem":"wiki","rel":"explica"},{"alvo":"koch_dissipa_nao_codifica","confianca":1.0,"epistemico":"fato","origem":"wiki","rel":"usa"},{"alvo":"absorcao_iterativa_poca","confianca":1.0,"epistemico":"fato","origem":"wiki","rel":"usa"},{"alvo":"plano_continuo_transicoes","confianca":1.0,"epistemico":"fato","origem":"wiki","rel":"parte_de"},{"alvo":"transformada_metalica_fechada_perelman","confianca":1.0,"epistemico":"fato","origem":"wiki","rel":"aplicacao"},{"alvo":"tiffany","confianca":1.0,"epistemico":"fato","origem":"wiki","rel":"parte_de"}],"confianca":1.0,"contraexemplos":[],"descricao":"Repetir a redenção em levas (o grafo cresceu 5533→8624 arestas, +3091 relações implicadas pelo corpus; purgatório 82275→76823, −5452) revelou a LEI: a taxa é ~1,8 fragmento redimido por aresta criada — porque a parábola (θ) limita cada bacia a poucos fragmentos. Drenar os ~77 mil restantes exigiria ~43 mil arestas (grafo 9× maior, afogado em relações de baixa confiança). Logo o purgatório NÃO SECA — CONVERGE a um resíduo. E isso é CORRETO, não falha: o purgatório é o LADO-A da parábola a+c²=1 — o calor dissipativo do acervo, as re-ditas redundantes que não cristalizam em posição única. A parábola as recusa como deve (é o floco que dissipa, não codifica — ver [[koch_dissipa_nao_codifica]], [[feedback_floco_no_metal]]). Forçar a secagem afoga o grafo; o certo é deixar o resíduo ser o que é: a entropia do corpus. Parei a ~8600 arestas. [D] roda; grafo íntegro (1 componente, 62/62).","epistemico":"fato","exemplos":[],"id":"purgatorio_nao_seca_lado_a","memoria":"semantica","meta":{"autor":"Lopes/broca-so","dominio":"continuo","fonte":"scratchpad/redeem_loop.py (levas); taxa 1,8 frag/aresta; +3091 arestas → −5452 purgatório"},"origem":"wiki","palavras_chave":[],"sinonimos":[],"tipo":"conceito","titulo":"O purgatório não seca — converge: é o lado-a (dissipativo) da parábola"}
\section{O purgatório não seca — converge: é o lado-a (dissipativo) da parábola}
Repetir a redenção em levas (o grafo cresceu 5533\ensuremath{\to}8624 arestas, +3091 relações implicadas pelo corpus; purgatório 82275\ensuremath{\to}76823, \ensuremath{-}5452) revelou a LEI: a taxa é \~{}1,8 fragmento redimido por aresta criada — porque a parábola (\ensuremath{\theta}) limita cada bacia a poucos fragmentos. Drenar os \~{}77 mil restantes exigiria \~{}43 mil arestas (grafo 9$\times$ maior, afogado em relações de baixa confiança). Logo o purgatório NÃO SECA — CONVERGE a um resíduo. E isso é CORRETO, não falha: o purgatório é o LADO-A da parábola a+c\ensuremath{^{2}}=1 — o calor dissipativo do acervo, as re-ditas redundantes que não cristalizam em posição única. A parábola as recusa como deve (é o floco que dissipa, não codifica — ver [[koch\_dissipa\_nao\_codifica]], [[feedback\_floco\_no\_metal]]). Forçar a secagem afoga o grafo; o certo é deixar o resíduo ser o que é: a entropia do corpus. Parei a \~{}8600 arestas. [D] roda; grafo íntegro (1 componente, 62/62).
\prov{relações: explica redencao\_purgatorio\_grafo; usa koch\_dissipa\_nao\_codifica; usa absorcao\_iterativa\_poca; parte\_de plano\_continuo\_transicoes; aplicacao transformada\_metalica\_fechada\_perelman; parte\_de tiffany}
\prov{proveniência: wiki \textperiodcentered{} scratchpad/redeem\_loop.py (levas); taxa 1,8 frag/aresta; +3091 arestas \ensuremath{\to} \ensuremath{-}5452 purgatório \textperiodcentered{} fato \textperiodcentered{} confiança 1.0 \textperiodcentered{} domínio continuo \textperiodcentered{} história 37 versões no jornal}

%CRISTAL {"arestas":[{"alvo":"symbolic_beta","confianca":1.0,"epistemico":"fato","origem":"paper","rel":"especializa"}],"confianca":1.0,"contraexemplos":[],"descricao":"Teorema: β_QED(e) não pertence a nenhuma classe recursiva fechada — a divergência de Dyson (1952) é impossibilidade algorítmica, não artefato numérico.","epistemico":"fato","exemplos":[],"id":"qed_nao_representavel","memoria":"semantica","meta":{"autor":"Gentil Lopes","descoberta":false,"dominio":"matematica","fonte":"paper_PQ_L_qed_rigorosa.pdf, Teo.5.1"},"origem":"paper","palavras_chave":[],"sinonimos":[],"tipo":"conceito","titulo":"QED Simbolicamente Não-Representável"}
\section{QED Simbolicamente Não-Representável}
Teorema: \ensuremath{\beta}\_QED(e) não pertence a nenhuma classe recursiva fechada — a divergência de Dyson (1952) é impossibilidade algorítmica, não artefato numérico.
\prov{relações: especializa symbolic\_beta}
\prov{proveniência: paper \textperiodcentered{} paper\_PQ\_L\_qed\_rigorosa.pdf, Teo.5.1 \textperiodcentered{} fato \textperiodcentered{} confiança 1.0 \textperiodcentered{} domínio matematica \textperiodcentered{} história 2 versões no jornal}

%CRISTAL {"arestas":[{"alvo":"colapso","confianca":1.0,"epistemico":"fato","origem":"paper","rel":"explica"},{"alvo":"acoplamento_cruzado","confianca":1.0,"epistemico":"fato","origem":"paper","rel":"especializa"}],"confianca":1.0,"contraexemplos":[],"descricao":"A quantização C_K: B_in→W_out (a curva de Koch como 'ribossomo') via z↦z² sobre a lemniscata y²=x³−x codifica discreto→contínuo sem inversa (A=0): purifica sem transportar o gerador.","epistemico":"fato","exemplos":[],"id":"quantizacao_dirigida_koch","memoria":"semantica","meta":{"autor":"Gentil Lopes","descoberta":true,"dominio":"matematica","fonte":"olho_de_koch.pdf, Teo.1"},"origem":"paper","palavras_chave":[],"sinonimos":[],"tipo":"conceito","titulo":"Quantização Dirigida (Olho de Koch)"}
\section{Quantização Dirigida (Olho de Koch)}
A quantização C\_K: B\_in\ensuremath{\to}W\_out (a curva de Koch como 'ribossomo') via z\ensuremath{\mapsto}z\ensuremath{^{2}} sobre a lemniscata y\ensuremath{^{2}}=x\ensuremath{^{3}}\ensuremath{-}x codifica discreto\ensuremath{\to}contínuo sem inversa (A=0): purifica sem transportar o gerador.
\prov{relações: explica colapso; especializa acoplamento\_cruzado}
\prov{proveniência: paper \textperiodcentered{} olho\_de\_koch.pdf, Teo.1 \textperiodcentered{} fato \textperiodcentered{} confiança 1.0 \textperiodcentered{} domínio matematica \textperiodcentered{} história 3 versões no jornal}

%CRISTAL {"arestas":[{"alvo":"helicidade_octonica","confianca":1.0,"epistemico":"fato","origem":"paper","rel":"usa"},{"alvo":"steiner_golay","confianca":1.0,"epistemico":"fato","origem":"paper","rel":"usa"}],"confianca":1.0,"contraexemplos":[],"descricao":"Conjectura: para 24 vórtices com simetria M₂₄, H_Ø=ΣLᵢⱼΓᵢΓⱼ é inteira exatamente nos codewords do código de Golay (~759), irracional fora — quantização topológica seleciona a dinâmica.","epistemico":"hipotese","exemplos":[],"id":"quantizacao_golay","memoria":"semantica","meta":{"autor":"Gentil Lopes","descoberta":false,"dominio":"matematica","fonte":"paper_BC.pdf, Conj.5"},"origem":"paper","palavras_chave":[],"sinonimos":[],"tipo":"conceito","titulo":"Quantização de Golay da Helicidade"}
\section{Quantização de Golay da Helicidade}
Conjectura: para 24 vórtices com simetria M\ensuremath{_{24}}, H\_\ensuremath{\varnothing}=\ensuremath{\Sigma}L\ensuremath{_{ij}}\ensuremath{\Gamma}\ensuremath{_{i}}\ensuremath{\Gamma}\ensuremath{_{j}} é inteira exatamente nos codewords do código de Golay (\~{}759), irracional fora — quantização topológica seleciona a dinâmica.
\prov{relações: usa helicidade\_octonica; usa steiner\_golay}
\prov{proveniência: paper \textperiodcentered{} paper\_BC.pdf, Conj.5 \textperiodcentered{} hipotese \textperiodcentered{} confiança 1.0 \textperiodcentered{} domínio matematica \textperiodcentered{} história 3 versões no jornal}

%CRISTAL {"arestas":[{"alvo":"colapso","confianca":1.0,"epistemico":"fato","origem":"paper","rel":"explica"},{"alvo":"algebra_tropical","confianca":1.0,"epistemico":"fato","origem":"paper","rel":"usa"}],"confianca":1.0,"contraexemplos":[],"descricao":"A conversão analógico→digital é o limite tropical T→0 do softmax (a⊕_T b→min): digitalizar é projetar no idempotente — o colapso.","epistemico":"fato","exemplos":[],"id":"quantizacao_tropical","memoria":"semantica","meta":{"autor":"Gentil Lopes","descoberta":false,"dominio":"matematica","fonte":"chicote_tropical.pdf"},"origem":"paper","palavras_chave":[],"sinonimos":[],"tipo":"conceito","titulo":"Quantização Tropical (T→0)"}
\section{Quantização Tropical (T\ensuremath{\to}0)}
A conversão analógico\ensuremath{\to}digital é o limite tropical T\ensuremath{\to}0 do softmax (a\ensuremath{\oplus}\_T b\ensuremath{\to}min): digitalizar é projetar no idempotente — o colapso.
\prov{relações: explica colapso; usa algebra\_tropical}
\prov{proveniência: paper \textperiodcentered{} chicote\_tropical.pdf \textperiodcentered{} fato \textperiodcentered{} confiança 1.0 \textperiodcentered{} domínio matematica \textperiodcentered{} história 3 versões no jornal}

%CRISTAL {"arestas":[{"alvo":"catedral_sigma_aritmetica","confianca":1.0,"epistemico":"fato","origem":"wiki","rel":"parte_de"},{"alvo":"cadeia_de_fano_sigma","confianca":1.0,"epistemico":"fato","origem":"wiki","rel":"especializa"}],"confianca":1.0,"contraexemplos":[],"descricao":"Quatro sementes independentes chegam ao mesmo 168 por iterar σ: σ⁵(7)=σ⁴(8)=σ³(23)=σ³(37)=168. Muitos caminhos, uma só simetria — a Muralha PSL como atrator da dinâmica de divisores.","epistemico":"fato","exemplos":[],"id":"quatro_cadeias_psl","memoria":"semantica","meta":{"autor":"Lord Index (Shawn R. Schiller)","dominio":"matematica","fonte":"A Catedral — Part VIII-A (jul 2026), ~/Imagens/catedral"},"origem":"wiki","palavras_chave":[],"sinonimos":[],"tipo":"conceito","titulo":"Quatro Cadeias, uma Muralha: σ-caminhos convergem em 168"}
\section{Quatro Cadeias, uma Muralha: \ensuremath{\sigma}-caminhos convergem em 168}
Quatro sementes independentes chegam ao mesmo 168 por iterar \ensuremath{\sigma}: \ensuremath{\sigma}\ensuremath{^{5}}(7)=\ensuremath{\sigma}\ensuremath{^{4}}(8)=\ensuremath{\sigma}\ensuremath{^{3}}(23)=\ensuremath{\sigma}\ensuremath{^{3}}(37)=168. Muitos caminhos, uma só simetria — a Muralha PSL como atrator da dinâmica de divisores.
\prov{relações: parte\_de catedral\_sigma\_aritmetica; especializa cadeia\_de\_fano\_sigma}
\prov{proveniência: wiki \textperiodcentered{} A Catedral — Part VIII-A (jul 2026), \~{}/Imagens/catedral \textperiodcentered{} fato \textperiodcentered{} confiança 1.0 \textperiodcentered{} domínio matematica \textperiodcentered{} história 3 versões no jornal}

%CRISTAL {"arestas":[{"alvo":"fibracao_hopf_bloch","confianca":1.0,"epistemico":"fato","origem":"paper","rel":"usa"},{"alvo":"mecanica_quantica","confianca":1.0,"epistemico":"fato","origem":"paper","rel":"parte_de"}],"confianca":1.0,"contraexemplos":[],"descricao":"S³=SU(2) dos versores é o espaço de estados; a esfera de Bloch S² é a base; a função de onda é uma seção do fibrado de linha, e a fase de Berry é a holonomia da torção.","epistemico":"fato","exemplos":[],"id":"qubit_hopf","memoria":"semantica","meta":{"autor":"Gentil Lopes","descoberta":false,"dominio":"matematica","fonte":"paper_3_quantica.pdf"},"origem":"paper","palavras_chave":[],"sinonimos":[],"tipo":"conceito","titulo":"Qubit como Seção de Hopf"}
\section{Qubit como Seção de Hopf}
S\ensuremath{^{3}}=SU(2) dos versores é o espaço de estados; a esfera de Bloch S\ensuremath{^{2}} é a base; a função de onda é uma seção do fibrado de linha, e a fase de Berry é a holonomia da torção.
\prov{relações: usa fibracao\_hopf\_bloch; parte\_de mecanica\_quantica}
\prov{proveniência: paper \textperiodcentered{} paper\_3\_quantica.pdf \textperiodcentered{} fato \textperiodcentered{} confiança 1.0 \textperiodcentered{} domínio matematica \textperiodcentered{} história 3 versões no jornal}

%CRISTAL {"arestas":[{"alvo":"octonios_hiperbolicos","confianca":1.0,"epistemico":"fato","origem":"paper","rel":"usa"},{"alvo":"associador","confianca":1.0,"epistemico":"fato","origem":"paper","rel":"usa"},{"alvo":"word","confianca":0.5,"epistemico":"fato","origem":"wiki","rel":"relaciona"}],"confianca":1.0,"contraexemplos":[],"descricao":"Em ℝ⁷=Im(𝕆) a identidade de Jacobi falha: J(a,b,c)=−3[a,b,c]≠0 (resíduo da ordem do termo principal). Fonte estrutural da massa dinâmica.","epistemico":"fato","exemplos":[],"id":"quebra_jacobi_r7","memoria":"semantica","meta":{"autor":"Gentil Lopes","descoberta":false,"dominio":"matematica","fonte":"paper_AF_yang_mills_octonios.pdf, Teo.3"},"origem":"paper","palavras_chave":[],"sinonimos":[],"tipo":"conceito","titulo":"Quebra de Jacobi em ℝ⁷"}
\section{Quebra de Jacobi em \ensuremath{\mathbb{R}}\ensuremath{^{7}}}
Em \ensuremath{\mathbb{R}}\ensuremath{^{7}}=Im(\ensuremath{\mathbb{O}}) a identidade de Jacobi falha: J(a,b,c)=\ensuremath{-}3[a,b,c]\ensuremath{\ne}0 (resíduo da ordem do termo principal). Fonte estrutural da massa dinâmica.
\prov{relações: usa octonios\_hiperbolicos; usa associador; relaciona word}
\prov{proveniência: paper \textperiodcentered{} paper\_AF\_yang\_mills\_octonios.pdf, Teo.3 \textperiodcentered{} fato \textperiodcentered{} confiança 1.0 \textperiodcentered{} domínio matematica \textperiodcentered{} história 4 versões no jornal}

%CRISTAL {"arestas":[{"alvo":"terceiro_estado","confianca":1.0,"epistemico":"fato","origem":"paper","rel":"usa"},{"alvo":"mecanica_quantica","confianca":1.0,"epistemico":"fato","origem":"paper","rel":"referencia"}],"confianca":1.0,"contraexemplos":[],"descricao":"O bit ±1 (elíptico/hiperbólico) torna-se qutrit quando o estado parabólico (c²=0) faz nascer o 0 no meio; leitura temporal: futuro, parado, passado.","epistemico":"fato","exemplos":[],"id":"qutrit","memoria":"semantica","meta":{"autor":"Gentil Lopes","descoberta":true,"dominio":"matematica","fonte":"paper_59_terceiro_estado.pdf, Obs.4.1"},"origem":"paper","palavras_chave":[],"sinonimos":[],"tipo":"conceito","titulo":"Qutrit −1,0,+1"}
\section{Qutrit \ensuremath{-}1,0,+1}
O bit $\pm$1 (elíptico/hiperbólico) torna-se qutrit quando o estado parabólico (c\ensuremath{^{2}}=0) faz nascer o 0 no meio; leitura temporal: futuro, parado, passado.
\prov{relações: usa terceiro\_estado; referencia mecanica\_quantica}
\prov{proveniência: paper \textperiodcentered{} paper\_59\_terceiro\_estado.pdf, Obs.4.1 \textperiodcentered{} fato \textperiodcentered{} confiança 1.0 \textperiodcentered{} domínio matematica \textperiodcentered{} história 4 versões no jornal}

%CRISTAL {"arestas":[{"alvo":"teoria_da_decisao","confianca":1.0,"epistemico":"fato","origem":"wiki","rel":"relaciona"}],"confianca":1.0,"contraexemplos":[],"descricao":"A racionalidade real é limitada por tempo, informação e capacidade cognitiva — não se OTIMIZA, se SATISFAZ (satisficing): escolhe-se a primeira opção 'boa o bastante'. A racionalidade da máquina finita.","epistemico":"inferencia","exemplos":[],"id":"racionalidade_limitada","memoria":"semantica","meta":{"autor":"Herbert Simon","dominio":"matematica","fonte":"teoria da decisão"},"origem":"wiki","palavras_chave":[],"sinonimos":[],"tipo":"conceito","titulo":"Racionalidade Limitada (Simon)"}
\section{Racionalidade Limitada (Simon)}
A racionalidade real é limitada por tempo, informação e capacidade cognitiva — não se OTIMIZA, se SATISFAZ (satisficing): escolhe-se a primeira opção 'boa o bastante'. A racionalidade da máquina finita.
\prov{relações: relaciona teoria\_da\_decisao}
\prov{proveniência: wiki \textperiodcentered{} teoria da decisão \textperiodcentered{} inferencia \textperiodcentered{} confiança 1.0 \textperiodcentered{} domínio matematica \textperiodcentered{} história 2 versões no jornal}

%CRISTAL {"arestas":[{"alvo":"dimensao_de_traco","confianca":1.0,"epistemico":"fato","origem":"wiki","rel":"usa"},{"alvo":"hopfield_recursiva_koch","confianca":1.0,"epistemico":"fato","origem":"wiki","rel":"relaciona"}],"confianca":1.0,"contraexemplos":[],"descricao":"Distinção essencial: o LFSR Fibonacci mod 2 (x²+x+1) tem período 3, órbita única, |L_n|≤3 → D_tr=0; já o subshift golden e o quiver com a MESMA matriz de Perron [[1,1],[1,0]] têm D_tr=log₂φ. Um anel simples tem D_tr=0 (loop infinito sem ramificação). O fractal mede a ramificação do espaço de traços, não uma trajetória.","epistemico":"fato","exemplos":[],"id":"ramificacao_nao_orbita","memoria":"semantica","meta":{"autor":"broca-so","dominio":"fractal","fonte":"machine/fractal/fractal.py (EXP.2); computador_fractal.tex"},"origem":"wiki","palavras_chave":[],"sinonimos":[],"tipo":"conceito","titulo":"O Fractal é a Ramificação, não a Órbita"}
\section{O Fractal é a Ramificação, não a Órbita}
Distinção essencial: o LFSR Fibonacci mod 2 (x\ensuremath{^{2}}+x+1) tem período 3, órbita única, |L\_n|\ensuremath{\le}3 \ensuremath{\to} D\_tr=0; já o subshift golden e o quiver com a MESMA matriz de Perron [[1,1],[1,0]] têm D\_tr=log\ensuremath{_{2}}\ensuremath{\varphi}. Um anel simples tem D\_tr=0 (loop infinito sem ramificação). O fractal mede a ramificação do espaço de traços, não uma trajetória.
\prov{relações: usa dimensao\_de\_traco; relaciona hopfield\_recursiva\_koch}
\prov{proveniência: wiki \textperiodcentered{} machine/fractal/fractal.py (EXP.2); computador\_fractal.tex \textperiodcentered{} fato \textperiodcentered{} confiança 1.0 \textperiodcentered{} domínio fractal \textperiodcentered{} história 3 versões no jornal}

%CRISTAL {"arestas":[{"alvo":"colapso","confianca":1.0,"epistemico":"fato","origem":"paper","rel":"referencia"},{"alvo":"computacao","confianca":1.0,"epistemico":"fato","origem":"paper","rel":"parte_de"}],"confianca":1.0,"contraexemplos":[],"descricao":"Critério que separa dinâmica fractal (ramificação, D_U>0) de digital (órbita cíclica, LFSR, D_U=0): a máquina fractal preserva D_U, a digital colapsa ao idempotente.","epistemico":"fato","exemplos":[],"id":"ramificacao_vs_orbita","memoria":"semantica","meta":{"autor":"Gentil Lopes","descoberta":true,"dominio":"matematica","fonte":"computador_fractal.pdf"},"origem":"paper","palavras_chave":[],"sinonimos":[],"tipo":"conceito","titulo":"Ramificação vs Órbita"}
\section{Ramificação vs Órbita}
Critério que separa dinâmica fractal (ramificação, D\_U>0) de digital (órbita cíclica, LFSR, D\_U=0): a máquina fractal preserva D\_U, a digital colapsa ao idempotente.
\prov{relações: referencia colapso; parte\_de computacao}
\prov{proveniência: paper \textperiodcentered{} computador\_fractal.pdf \textperiodcentered{} fato \textperiodcentered{} confiança 1.0 \textperiodcentered{} domínio matematica \textperiodcentered{} história 3 versões no jornal}

%CRISTAL {"arestas":[{"alvo":"experimento_da_agulha","confianca":1.0,"epistemico":"fato","origem":"wiki","rel":"parte_de"},{"alvo":"flip_duopolinomios","confianca":1.0,"epistemico":"fato","origem":"wiki","rel":"usa"},{"alvo":"newton_acha_o_grupo_discreto","confianca":1.0,"epistemico":"fato","origem":"wiki","rel":"relaciona"},{"alvo":"cadeias_alem_dos_metais","confianca":1.0,"epistemico":"fato","origem":"wiki","rel":"usa"},{"alvo":"pi_gerador_da_reta_acima","confianca":1.0,"epistemico":"fato","origem":"wiki","rel":"relaciona"}],"confianca":1.0,"contraexemplos":[],"descricao":"O rastro deixado pela agulha: NA reta mineral (o metal, grau 2) o fractal de Newton é uma RETA (dim≈1) — a agulha não invade a curvatura. ALÉM dos metais (grau ≥3, o caos, a próxima reta) o rastro vira FRACTAL (dim>1). O fractal marca exatamente a BORDA onde a reta mineral termina e a próxima começa. Verif.: metal dim 0.98 (reta), caos dim 1.41 (fractal).","epistemico":"fato","exemplos":[],"id":"rastro_da_agulha_e_fractal","memoria":"semantica","meta":{"dominio":"reta mineral","fonte":"experimento da agulha"},"origem":"wiki","palavras_chave":[],"sinonimos":[],"tipo":"conceito","titulo":"O Rastro: Reta na Reta Mineral, Fractal Além"}
\section{O Rastro: Reta na Reta Mineral, Fractal Além}
O rastro deixado pela agulha: NA reta mineral (o metal, grau 2) o fractal de Newton é uma RETA (dim\ensuremath{\approx}1) — a agulha não invade a curvatura. ALÉM dos metais (grau \ensuremath{\ge}3, o caos, a próxima reta) o rastro vira FRACTAL (dim>1). O fractal marca exatamente a BORDA onde a reta mineral termina e a próxima começa. Verif.: metal dim 0.98 (reta), caos dim 1.41 (fractal).
\prov{relações: parte\_de experimento\_da\_agulha; usa flip\_duopolinomios; relaciona newton\_acha\_o\_grupo\_discreto; usa cadeias\_alem\_dos\_metais; relaciona pi\_gerador\_da\_reta\_acima}
\prov{proveniência: wiki \textperiodcentered{} experimento da agulha \textperiodcentered{} fato \textperiodcentered{} confiança 1.0 \textperiodcentered{} domínio reta mineral \textperiodcentered{} história 6 versões no jornal}

%CRISTAL {"arestas":[{"alvo":"numeros_metalicos","confianca":1.0,"epistemico":"fato","origem":"livro","rel":"especializa"},{"alvo":"reais","confianca":1.0,"epistemico":"fato","origem":"livro","rel":"depende"},{"alvo":"matematica","confianca":1.0,"epistemico":"fato","origem":"livro","rel":"parte_de"}],"confianca":1.0,"contraexemplos":[],"descricao":"φ = (1+√5)/2 ≈ 1,618, raiz positiva de x² = x + 1. Base de ouro (Bergman): todo real ≥ 0 se escreve Σ dᵢ φⁱ com dᵢ ∈ 0,1, canônica sem dois 1 consecutivos.","epistemico":"fato","exemplos":[],"id":"razao_aurea","memoria":"semantica","meta":{"autor":"Gentil Lopes","descoberta":false,"dominio":"matematica","fonte":""},"origem":"manual","palavras_chave":["φ","proporção","ouro"],"sinonimos":["phi","número de ouro"],"tipo":"conceito","titulo":"Razão Áurea (φ)"}
\section{Razão Áurea (\ensuremath{\varphi})}
\ensuremath{\varphi} = (1+\ensuremath{\sqrt{\,}}5)/2 \ensuremath{\approx} 1,618, raiz positiva de x\ensuremath{^{2}} = x + 1. Base de ouro (Bergman): todo real \ensuremath{\ge} 0 se escreve \ensuremath{\Sigma} d\ensuremath{_{i}} \ensuremath{\varphi}\ensuremath{^{i}} com d\ensuremath{_{i}} \ensuremath{\in} 0,1, canônica sem dois 1 consecutivos.
\prov{sinónimos: phi; número de ouro}
\prov{relações: especializa numeros\_metalicos; depende reais; parte\_de matematica}
\prov{proveniência: manual \textperiodcentered{} fato \textperiodcentered{} confiança 1.0 \textperiodcentered{} domínio matematica \textperiodcentered{} história 9 versões no jornal}

%CRISTAL {"arestas":[{"alvo":"numeros_metalicos","confianca":1.0,"epistemico":"fato","origem":"livro","rel":"especializa"}],"confianca":1.0,"contraexemplos":[],"descricao":"δ = 1 + √2 ≈ 2,414, raiz positiva de x² = 2x + 1 (k=2). É o limite da razão de termos consecutivos da sequência de Pell.","epistemico":"fato","exemplos":[],"id":"razao_de_prata","memoria":"semantica","meta":{"autor":"Gentil Lopes","descoberta":false,"dominio":"matematica","fonte":""},"origem":"manual","palavras_chave":[],"sinonimos":["silver ratio","delta_S"],"tipo":"conceito","titulo":"Razão de Prata (δ)"}
\section{Razão de Prata (\ensuremath{\delta})}
\ensuremath{\delta} = 1 + \ensuremath{\sqrt{\,}}2 \ensuremath{\approx} 2,414, raiz positiva de x\ensuremath{^{2}} = 2x + 1 (k=2). É o limite da razão de termos consecutivos da sequência de Pell.
\prov{sinónimos: silver ratio; delta\_S}
\prov{relações: especializa numeros\_metalicos}
\prov{proveniência: manual \textperiodcentered{} fato \textperiodcentered{} confiança 1.0 \textperiodcentered{} domínio matematica \textperiodcentered{} história 4 versões no jornal}

%CRISTAL {"arestas":[{"alvo":"universalidade_reconexao","confianca":1.0,"epistemico":"fato","origem":"paper","rel":"parte_de"}],"confianca":1.0,"contraexemplos":[],"descricao":"A reconexão de vórtices (Kerr-Hussain) em tempo finito muda os linkings: é um passo discreto e irreversível de máquina de Turing — traço de um cobordismo 2D na dinâmica.","epistemico":"fato","exemplos":[],"id":"reconexao_cobordismo","memoria":"semantica","meta":{"autor":"Gentil Lopes","descoberta":false,"dominio":"matematica","fonte":"paper_AT.pdf"},"origem":"paper","palavras_chave":[],"sinonimos":[],"tipo":"conceito","titulo":"Reconexão como Cobordismo"}
\section{Reconexão como Cobordismo}
A reconexão de vórtices (Kerr-Hussain) em tempo finito muda os linkings: é um passo discreto e irreversível de máquina de Turing — traço de um cobordismo 2D na dinâmica.
\prov{relações: parte\_de universalidade\_reconexao}
\prov{proveniência: paper \textperiodcentered{} paper\_AT.pdf \textperiodcentered{} fato \textperiodcentered{} confiança 1.0 \textperiodcentered{} domínio matematica \textperiodcentered{} história 2 versões no jornal}

%CRISTAL {"arestas":[{"alvo":"dois_ciclos_reconstrucao","confianca":1.0,"epistemico":"fato","origem":"wiki","rel":"usa"},{"alvo":"corpo_dual","confianca":1.0,"epistemico":"fato","origem":"wiki","rel":"usa"},{"alvo":"hopfield_recursiva_koch","confianca":1.0,"epistemico":"fato","origem":"wiki","rel":"contradiz"}],"confianca":1.0,"contraexemplos":[],"descricao":"Teorema: (i) só 𝓛 não basta (Φ não injetivo em s — colisões); (ii) só Φ não basta (recall só pelo atrator = 0/32); (iii) 𝓛+Φ reconstrói ω DESDE QUE o Word vivo esteja na célula branca (slot mem) ou seja recuperável por busca finita — o floco fixa a ramificação z² (2-para-1, A=0). Reconstrução sem célula branca é ilusória: 'o floco não é Hopfield'.","epistemico":"fato","exemplos":[],"id":"reconstrucao_condicionada","memoria":"semantica","meta":{"autor":"broca-so","dominio":"fractal","fonte":"papers/floco_de_neve.tex §V.2"},"origem":"wiki","palavras_chave":[],"sinonimos":[],"tipo":"conceito","titulo":"Reconstrução só Fecha com a Célula Branca"}
\section{Reconstrução só Fecha com a Célula Branca}
Teorema: (i) só \ensuremath{\mathcal{L}} não basta (\ensuremath{\Phi} não injetivo em s — colisões); (ii) só \ensuremath{\Phi} não basta (recall só pelo atrator = 0/32); (iii) \ensuremath{\mathcal{L}}+\ensuremath{\Phi} reconstrói \ensuremath{\omega} DESDE QUE o Word vivo esteja na célula branca (slot mem) ou seja recuperável por busca finita — o floco fixa a ramificação z\ensuremath{^{2}} (2-para-1, A=0). Reconstrução sem célula branca é ilusória: 'o floco não é Hopfield'.
\prov{relações: usa dois\_ciclos\_reconstrucao; usa corpo\_dual; contradiz hopfield\_recursiva\_koch}
\prov{proveniência: wiki \textperiodcentered{} papers/floco\_de\_neve.tex §V.2 \textperiodcentered{} fato \textperiodcentered{} confiança 1.0 \textperiodcentered{} domínio fractal \textperiodcentered{} história 4 versões no jornal}

%CRISTAL {"arestas":[{"alvo":"recuperacao_evidencia_invariante","confianca":1.0,"epistemico":"fato","origem":"wiki","rel":"especializa"},{"alvo":"definicao_frasagem_ampla","confianca":1.0,"epistemico":"fato","origem":"wiki","rel":"usa"},{"alvo":"tiffany","confianca":1.0,"epistemico":"fato","origem":"wiki","rel":"parte_de"}],"confianca":1.0,"contraexemplos":[],"descricao":"conversa/colapso.py (evidenciaconteudo, c2efetivo): rodar a Tiffany sobre a PRÓPRIA linguagem expôs e corrigiu três falhas na aboutness (o casamento alvo↔título da definição): (1) PONTUAÇÃO no token — 'interpretador:' ≠ 'interpretador' no split → tira PONT das bordas; (2) STOPWORD na interseção — o 'da' casava 'Didática DA Matemática' com 'backend DA ISA' → filtra STOP das palavras de conteúdo; (3) SUBSET DE CONJUNTO VAZIO — 'tw ⊆ aw' é sempre True quando tw=, então o nó-fragmento titulado literalmente 'O que é' ganhava forte=0,95 POR VACUIDADE → agora exige título com conteúdo (tw não-vazio), e um GATE de supressão (0,30·) rebaixa nós OFF-TOPIC numa definição (mata o ov coincidente dos títulos 'O que é'/'O que são' dos hubs). Efeito: 7/7 perguntas sobre a linguagem aterram no conceito certo; saudação/fluxo intactos; 62 testes. É o endurecimento de [[recuperacaoevidenciaᵢnvariante]] — a máquina, ao se descrever, poliu a própria recuperação. [F] verificado.","epistemico":"fato","exemplos":[],"id":"recuperacao_aboutness_robusta","memoria":"semantica","meta":{"autor":"Lopes/broca-so","dominio":"algebra","fonte":"conversa/colapso.py:_evidencia_conteudo (_PONT/_STOP, tw não-vazio, gate 0.30), _c2_efetivo; verificado: 7/7 sobre a linguagem"},"origem":"wiki","palavras_chave":[],"sinonimos":[],"tipo":"conceito","titulo":"A aboutness da recuperação, endurecida: pontuação, stopwords e o subset-vazio + gate [F]"}
\section{A aboutness da recuperação, endurecida: pontuação, stopwords e o subset-vazio + gate [F]}
conversa/colapso.py (evidenciaconteudo, c2efetivo): rodar a Tiffany sobre a PRÓPRIA linguagem expôs e corrigiu três falhas na aboutness (o casamento alvo\ensuremath{\leftrightarrow}título da definição): (1) PONTUAÇÃO no token — 'interpretador:' \ensuremath{\ne} 'interpretador' no split \ensuremath{\to} tira PONT das bordas; (2) STOPWORD na interseção — o 'da' casava 'Didática DA Matemática' com 'backend DA ISA' \ensuremath{\to} filtra STOP das palavras de conteúdo; (3) SUBSET DE CONJUNTO VAZIO — 'tw \ensuremath{\subseteq} aw' é sempre True quando tw=, então o nó-fragmento titulado literalmente 'O que é' ganhava forte=0,95 POR VACUIDADE \ensuremath{\to} agora exige título com conteúdo (tw não-vazio), e um GATE de supressão (0,30\textperiodcentered ) rebaixa nós OFF-TOPIC numa definição (mata o ov coincidente dos títulos 'O que é'/'O que são' dos hubs). Efeito: 7/7 perguntas sobre a linguagem aterram no conceito certo; saudação/fluxo intactos; 62 testes. É o endurecimento de [[recuperacaoevidencia\ensuremath{_{i}}nvariante]] — a máquina, ao se descrever, poliu a própria recuperação. [F] verificado.
\prov{relações: especializa recuperacao\_evidencia\_invariante; usa definicao\_frasagem\_ampla; parte\_de tiffany}
\prov{proveniência: wiki \textperiodcentered{} conversa/colapso.py:\_evidencia\_conteudo (\_PONT/\_STOP, tw não-vazio, gate 0.30), \_c2\_efetivo; verificado: 7/7 sobre a linguagem \textperiodcentered{} fato \textperiodcentered{} confiança 1.0 \textperiodcentered{} domínio algebra \textperiodcentered{} história 37 versões no jornal}

%CRISTAL {"arestas":[{"alvo":"conservacao_parabola_costura","confianca":1.0,"epistemico":"fato","origem":"wiki","rel":"usa"},{"alvo":"bai_gate_dtc_orbita","confianca":1.0,"epistemico":"fato","origem":"wiki","rel":"usa"},{"alvo":"ingestao_refino_parabola","confianca":1.0,"epistemico":"fato","origem":"wiki","rel":"relaciona"},{"alvo":"transformada_metalica","confianca":1.0,"epistemico":"fato","origem":"wiki","rel":"usa"},{"alvo":"tiffany","confianca":1.0,"epistemico":"fato","origem":"wiki","rel":"parte_de"}],"confianca":1.0,"contraexemplos":[],"descricao":"conversa/colapso.py (_c2_efetivo, _evidencia_conteudo, _pa_ordem_m): a recuperação falhava porque a impressão Koch captura FORMA, não conteúdo — o nó certo (ov baixo) nem entrava no pool, e o conteúdo, quando entrava, SOMAVA por fora (+200 de título) INVADINDO a parábola. Correção obedecendo o invariante a+c²=1: a aboutness (título/conteúdo) virou EVIDÊNCIA que REDUZ a entropia — a costura de PERELMAN dentro da parábola, não somando por fora. É a PA de ordem m (o casamento de conteúdo, n-gramas até m — a progressão aritmética/diferenças finitas) entrando por NOISY-OR: a_ef = (1−α·c²₀)·(1−e_conteúdo); c²_ef = 1−a_ef ∈ [0,1] SEMPRE. O Koch escalado pelo gate (PG, o dupolinômio das retas) e o conteúdo (PA/Perelman) cada um COME um pedaço do lado-a (entropia monótona, Perelman) — nada invade (c²≤1). Efeito verificado: 'o que é topologia?' → o nó topologia; 'o que é entropia?' → Entropia Topológica (h=ln λ, λ=φ) — antes vinha 'virtualização'. score=round(256·c²_ef)+eco. Aterra no conteúdo respeitando a conservação. [D] a recuperação como fluxo de evidência na parábola, sem invadir.","epistemico":"fato","exemplos":[],"id":"recuperacao_evidencia_invariante","memoria":"semantica","meta":{"autor":"Lopes/broca-so","dominio":"algebra","fonte":"conversa/colapso.py:_c2_efetivo, _evidencia_conteudo, _pa_ordem_m; PA de ordem m (diferenças finitas); noisy-OR (Perelman)"},"origem":"wiki","palavras_chave":[],"sinonimos":[],"tipo":"conceito","titulo":"O gargalo da recuperação cede por EVIDÊNCIA que OBEDECE o invariante (não invade a parábola)"}
\section{O gargalo da recuperação cede por EVIDÊNCIA que OBEDECE o invariante (não invade a parábola)}
conversa/colapso.py (\_c2\_efetivo, \_evidencia\_conteudo, \_pa\_ordem\_m): a recuperação falhava porque a impressão Koch captura FORMA, não conteúdo — o nó certo (ov baixo) nem entrava no pool, e o conteúdo, quando entrava, SOMAVA por fora (+200 de título) INVADINDO a parábola. Correção obedecendo o invariante a+c\ensuremath{^{2}}=1: a aboutness (título/conteúdo) virou EVIDÊNCIA que REDUZ a entropia — a costura de PERELMAN dentro da parábola, não somando por fora. É a PA de ordem m (o casamento de conteúdo, n-gramas até m — a progressão aritmética/diferenças finitas) entrando por NOISY-OR: a\_ef = (1\ensuremath{-}\ensuremath{\alpha}\textperiodcentered c\ensuremath{^{2}}\ensuremath{_{0}})\textperiodcentered (1\ensuremath{-}e\_conteúdo); c\ensuremath{^{2}}\_ef = 1\ensuremath{-}a\_ef \ensuremath{\in} [0,1] SEMPRE. O Koch escalado pelo gate (PG, o dupolinômio das retas) e o conteúdo (PA/Perelman) cada um COME um pedaço do lado-a (entropia monótona, Perelman) — nada invade (c\ensuremath{^{2}}\ensuremath{\le}1). Efeito verificado: 'o que é topologia?' \ensuremath{\to} o nó topologia; 'o que é entropia?' \ensuremath{\to} Entropia Topológica (h=ln \ensuremath{\lambda}, \ensuremath{\lambda}=\ensuremath{\varphi}) — antes vinha 'virtualização'. score=round(256\textperiodcentered c\ensuremath{^{2}}\_ef)+eco. Aterra no conteúdo respeitando a conservação. [D] a recuperação como fluxo de evidência na parábola, sem invadir.
\prov{relações: usa conservacao\_parabola\_costura; usa bai\_gate\_dtc\_orbita; relaciona ingestao\_refino\_parabola; usa transformada\_metalica; parte\_de tiffany}
\prov{proveniência: wiki \textperiodcentered{} conversa/colapso.py:\_c2\_efetivo, \_evidencia\_conteudo, \_pa\_ordem\_m; PA de ordem m (diferenças finitas); noisy-OR (Perelman) \textperiodcentered{} fato \textperiodcentered{} confiança 1.0 \textperiodcentered{} domínio algebra \textperiodcentered{} história 78 versões no jornal}

%CRISTAL {"arestas":[{"alvo":"mathieu_m24","confianca":1.0,"epistemico":"fato","origem":"paper","rel":"usa"},{"alvo":"word","confianca":0.5,"epistemico":"fato","origem":"wiki","rel":"relaciona"},{"alvo":"virtualizacao","confianca":0.5,"epistemico":"fato","origem":"wiki","rel":"relaciona"},{"alvo":"vida-morte","confianca":0.5,"epistemico":"fato","origem":"wiki","rel":"relaciona"}],"confianca":1.0,"contraexemplos":[],"descricao":"Única rede unimodular par em dim 24: norma mínima 4, kissing number 196.560, grupo de automorfismos Co₀ (Conway) ⊃ M₂₄.","epistemico":"fato","exemplos":[],"id":"rede_leech","memoria":"semantica","meta":{"autor":"Gentil Lopes","descoberta":false,"dominio":"matematica","fonte":"paper_BD.pdf (Leech 1967)"},"origem":"paper","palavras_chave":[],"sinonimos":[],"tipo":"conceito","titulo":"Rede de Leech Λ₂₄"}
\section{Rede de Leech \ensuremath{\Lambda}\ensuremath{_{24}}}
Única rede unimodular par em dim 24: norma mínima 4, kissing number 196.560, grupo de automorfismos Co\ensuremath{_{0}} (Conway) \ensuremath{\supset} M\ensuremath{_{24}}.
\prov{relações: usa mathieu\_m24; relaciona word; relaciona virtualizacao; relaciona vida-morte}
\prov{proveniência: paper \textperiodcentered{} paper\_BD.pdf (Leech 1967) \textperiodcentered{} fato \textperiodcentered{} confiança 1.0 \textperiodcentered{} domínio matematica \textperiodcentered{} história 5 versões no jornal}

%CRISTAL {"arestas":[{"alvo":"absorcao_iterativa_poca","confianca":1.0,"epistemico":"fato","origem":"wiki","rel":"usa"},{"alvo":"mapear_nao_fabricar_continuo","confianca":1.0,"epistemico":"fato","origem":"wiki","rel":"usa"},{"alvo":"plano_continuo_transicoes","confianca":1.0,"epistemico":"fato","origem":"wiki","rel":"parte_de"}],"confianca":1.0,"contraexemplos":[],"descricao":"O purgatório não drena contra o mesmo grafo — mas REVELA as arestas que faltam: um fragmento sem casa ressoa com DOIS nós entre os quais não há relação. O metal emite o PAR de nós mais ressonante de cada órfão (pares.txt, top-2); os pares frequentes que NÃO são arestas são as relações que o corpus IMPLICA e o grafo não tinha (ex.: fase3~transcricao 5149×, contabilidade~fase3 2239×). Crescer o grafo = adicionar essas arestas (relaciona, inferido, confiança 0,5, proveniência corpus) → viram bacias novas de θ baixo → o purgatório é re-roteado e REDIMIDO. Rodada: +522 arestas (5533→6055) → +993 fragmentos redimidos, as novas arestas preenchidas (cobertura 98,3% do grafo maior). É ITERATIVO: cada rodada de crescimento redime mais. O purgatório é o SINAL DE CRESCIMENTO do grafo — densificar-não-podar guiado pelo que o acervo pede. [D] roda; grafo íntegro (1 componente, 62/62).","epistemico":"fato","exemplos":[],"id":"redencao_purgatorio_grafo","memoria":"semantica","meta":{"autor":"Lopes/broca-so","dominio":"continuo","fonte":"doc/continuo/absorve_metal.c (pares.txt top-2 nós); ligar (relaciona inferido); +522 arestas → +993 redimidos"},"origem":"wiki","palavras_chave":[],"sinonimos":[],"tipo":"conceito","titulo":"Redimir o purgatório crescendo o grafo: o órfão revela a aresta que falta"}
\section{Redimir o purgatório crescendo o grafo: o órfão revela a aresta que falta}
O purgatório não drena contra o mesmo grafo — mas REVELA as arestas que faltam: um fragmento sem casa ressoa com DOIS nós entre os quais não há relação. O metal emite o PAR de nós mais ressonante de cada órfão (pares.txt, top-2); os pares frequentes que NÃO são arestas são as relações que o corpus IMPLICA e o grafo não tinha (ex.: fase3\~{}transcricao 5149$\times$, contabilidade\~{}fase3 2239$\times$). Crescer o grafo = adicionar essas arestas (relaciona, inferido, confiança 0,5, proveniência corpus) \ensuremath{\to} viram bacias novas de \ensuremath{\theta} baixo \ensuremath{\to} o purgatório é re-roteado e REDIMIDO. Rodada: +522 arestas (5533\ensuremath{\to}6055) \ensuremath{\to} +993 fragmentos redimidos, as novas arestas preenchidas (cobertura 98,3\% do grafo maior). É ITERATIVO: cada rodada de crescimento redime mais. O purgatório é o SINAL DE CRESCIMENTO do grafo — densificar-não-podar guiado pelo que o acervo pede. [D] roda; grafo íntegro (1 componente, 62/62).
\prov{relações: usa absorcao\_iterativa\_poca; usa mapear\_nao\_fabricar\_continuo; parte\_de plano\_continuo\_transicoes}
\prov{proveniência: wiki \textperiodcentered{} doc/continuo/absorve\_metal.c (pares.txt top-2 nós); ligar (relaciona inferido); +522 arestas \ensuremath{\to} +993 redimidos \textperiodcentered{} fato \textperiodcentered{} confiança 1.0 \textperiodcentered{} domínio continuo \textperiodcentered{} história 12 versões no jornal}

%CRISTAL {"arestas":[{"alvo":"precmaass_verificacao","confianca":1.0,"epistemico":"fato","origem":"paper","rel":"usa"}],"confianca":1.0,"contraexemplos":[],"descricao":"Teorema: um algoritmo que produz blocos YM-PoW válidos resolve o problema PrecMaass (aproximar o espectro) em tempo comparável — a mineração É computação espectral, não trabalho arbitrário.","epistemico":"fato","exemplos":[],"id":"reducao_alpha_pura","memoria":"semantica","meta":{"autor":"Gentil Lopes","descoberta":false,"dominio":"matematica","fonte":"paper_PQ_R.pdf, Teo.5.1"},"origem":"paper","palavras_chave":[],"sinonimos":[],"tipo":"conceito","titulo":"Redução α-Pura (mining ≡ PrecMaass)"}
\section{Redução \ensuremath{\alpha}-Pura (mining \ensuremath{\equiv} PrecMaass)}
Teorema: um algoritmo que produz blocos YM-PoW válidos resolve o problema PrecMaass (aproximar o espectro) em tempo comparável — a mineração É computação espectral, não trabalho arbitrário.
\prov{relações: usa precmaass\_verificacao}
\prov{proveniência: paper \textperiodcentered{} paper\_PQ\_R.pdf, Teo.5.1 \textperiodcentered{} fato \textperiodcentered{} confiança 1.0 \textperiodcentered{} domínio matematica \textperiodcentered{} história 2 versões no jornal}

%CRISTAL {"arestas":[{"alvo":"hierarquia_pspace_indecidivel","confianca":1.0,"epistemico":"fato","origem":"paper","rel":"usa"}],"confianca":1.0,"contraexemplos":[],"descricao":"O problema da integração em forma fechada (Risch/Richardson, indecidível) reduz a SymbolicBeta_∞ via o morfismo Θ: no limite, dureza vira indecidibilidade pura.","epistemico":"fato","exemplos":[],"id":"reducao_risch","memoria":"semantica","meta":{"autor":"Gentil Lopes","descoberta":false,"dominio":"matematica","fonte":"paper_PQ_F_threshold_indecidibilidade.pdf, Teo.3.3"},"origem":"paper","palavras_chave":[],"sinonimos":[],"tipo":"conceito","titulo":"Redução de Risch (indecidibilidade)"}
\section{Redução de Risch (indecidibilidade)}
O problema da integração em forma fechada (Risch/Richardson, indecidível) reduz a SymbolicBeta\_\ensuremath{\infty} via o morfismo \ensuremath{\Theta}: no limite, dureza vira indecidibilidade pura.
\prov{relações: usa hierarquia\_pspace\_indecidivel}
\prov{proveniência: paper \textperiodcentered{} paper\_PQ\_F\_threshold\_indecidibilidade.pdf, Teo.3.3 \textperiodcentered{} fato \textperiodcentered{} confiança 1.0 \textperiodcentered{} domínio matematica \textperiodcentered{} história 2 versões no jornal}

%CRISTAL {"arestas":[{"alvo":"refino_continuo_floco_metais","confianca":1.0,"epistemico":"fato","origem":"wiki","rel":"especializa"},{"alvo":"tiffany","confianca":1.0,"epistemico":"fato","origem":"wiki","rel":"parte_de"}],"confianca":1.0,"contraexemplos":[],"descricao":"O pipe SymbolicBeta multimetal (ver [[refino_continuo_floco_metais]]) foi migrado e reiniciado no gex44 (gex44, sob claude, tmux broca-stratum): symbolic.py live --lanes 12, ~4 MH/s. O floco é refinado em 12 metais no tempo morto entre blocos (um giro por chunk), e o log mostra a posição CORRENTE de todos: 'ouro#N prata#N bronze#N cobre#N m5..m12'. A ponte de continuidade da prata (n=2) funcionou: retomou a cadeia Pell antiga (~8830 → segue). O olho lê 'flocos metálicos + N metais' da parábola. Ledger com backup (ledger_symbolic.bak_multimetal). [D] a produção multimetal em operação real.","epistemico":"fato","exemplos":[],"id":"refino_multimetal_live_gex44","memoria":"semantica","meta":{"autor":"Lopes/broca-so","dominio":"algebra","fonte":"gex44 tmux broca-stratum; goldchain/symbolic.py (Refinador.linha_metais); ver project_gex44_symbolic"},"origem":"wiki","palavras_chave":[],"sinonimos":[],"tipo":"conceito","titulo":"O refino multimetal está LIVE no gex44 — 12 metais visíveis, a prata segue a ponte Pell"}
\section{O refino multimetal está LIVE no gex44 — 12 metais visíveis, a prata segue a ponte Pell}
O pipe SymbolicBeta multimetal (ver [[refino\_continuo\_floco\_metais]]) foi migrado e reiniciado no gex44 (gex44, sob claude, tmux broca-stratum): symbolic.py live --lanes 12, \~{}4 MH/s. O floco é refinado em 12 metais no tempo morto entre blocos (um giro por chunk), e o log mostra a posição CORRENTE de todos: 'ouro\#N prata\#N bronze\#N cobre\#N m5..m12'. A ponte de continuidade da prata (n=2) funcionou: retomou a cadeia Pell antiga (\~{}8830 \ensuremath{\to} segue). O olho lê 'flocos metálicos + N metais' da parábola. Ledger com backup (ledger\_symbolic.bak\_multimetal). [D] a produção multimetal em operação real.
\prov{relações: especializa refino\_continuo\_floco\_metais; parte\_de tiffany}
\prov{proveniência: wiki \textperiodcentered{} gex44 tmux broca-stratum; goldchain/symbolic.py (Refinador.linha\_metais); ver project\_gex44\_symbolic \textperiodcentered{} fato \textperiodcentered{} confiança 1.0 \textperiodcentered{} domínio algebra \textperiodcentered{} história 79 versões no jornal}

%CRISTAL {"arestas":[{"alvo":"numeros_duais","confianca":1.0,"epistemico":"fato","origem":"paper","rel":"usa"},{"alvo":"regime_tropical","confianca":1.0,"epistemico":"fato","origem":"paper","rel":"dualidade"}],"confianca":1.0,"contraexemplos":[],"descricao":"Estado nilpotente (n²=0): a parte 'negra' da decomposição de Peirce, canais square-zero que só transportam (reversível, congelado). Dual do tropical.","epistemico":"fato","exemplos":[],"id":"regime_glacial","memoria":"semantica","meta":{"autor":"Gentil Lopes","descoberta":true,"dominio":"matematica","fonte":"corpo_glacial.pdf"},"origem":"paper","palavras_chave":[],"sinonimos":[],"tipo":"conceito","titulo":"Regime Glacial (n²=0)"}
\section{Regime Glacial (n\ensuremath{^{2}}=0)}
Estado nilpotente (n\ensuremath{^{2}}=0): a parte 'negra' da decomposição de Peirce, canais square-zero que só transportam (reversível, congelado). Dual do tropical.
\prov{relações: usa numeros\_duais; dualidade regime\_tropical}
\prov{proveniência: paper \textperiodcentered{} corpo\_glacial.pdf \textperiodcentered{} fato \textperiodcentered{} confiança 1.0 \textperiodcentered{} domínio matematica \textperiodcentered{} história 3 versões no jornal}

%CRISTAL {"arestas":[{"alvo":"algebra_tropical","confianca":1.0,"epistemico":"fato","origem":"paper","rel":"usa"},{"alvo":"word","confianca":0.5,"epistemico":"fato","origem":"wiki","rel":"relaciona"},{"alvo":"virtualizacao","confianca":0.5,"epistemico":"fato","origem":"wiki","rel":"relaciona"},{"alvo":"vida-morte","confianca":0.5,"epistemico":"fato","origem":"wiki","rel":"relaciona"},{"alvo":"syscall","confianca":0.5,"epistemico":"fato","origem":"wiki","rel":"relaciona"}],"confianca":1.0,"contraexemplos":[],"descricao":"Estado idempotente (ε²=e): divisão com |λ|>1 expande; a face 'branca/viva' que reconstrói na base de ouro φ. Dual do regime glacial.","epistemico":"fato","exemplos":[],"id":"regime_tropical","memoria":"semantica","meta":{"autor":"Gentil Lopes","descoberta":true,"dominio":"matematica","fonte":"corpo_tropical.pdf"},"origem":"paper","palavras_chave":[],"sinonimos":[],"tipo":"conceito","titulo":"Regime Tropical (ε²=e)"}
\section{Regime Tropical (\ensuremath{\varepsilon}\ensuremath{^{2}}=e)}
Estado idempotente (\ensuremath{\varepsilon}\ensuremath{^{2}}=e): divisão com |\ensuremath{\lambda}|>1 expande; a face 'branca/viva' que reconstrói na base de ouro \ensuremath{\varphi}. Dual do regime glacial.
\prov{relações: usa algebra\_tropical; relaciona word; relaciona virtualizacao; relaciona vida-morte; relaciona syscall}
\prov{proveniência: paper \textperiodcentered{} corpo\_tropical.pdf \textperiodcentered{} fato \textperiodcentered{} confiança 1.0 \textperiodcentered{} domínio matematica \textperiodcentered{} história 6 versões no jornal}

%CRISTAL {"arestas":[{"alvo":"numeros_metalicos","confianca":1.0,"epistemico":"fato","origem":"livro","rel":"usa"},{"alvo":"geometria","confianca":1.0,"epistemico":"fato","origem":"livro","rel":"depende"},{"alvo":"mecanica_quantica","confianca":1.0,"epistemico":"fato","origem":"livro","rel":"referencia"}],"confianca":1.0,"contraexemplos":[],"descricao":"Métrica ds = dx/(1−x²) (horizonte s=1) que carrega os números metálicos como pontos notáveis sₚ = μ/(1+μ); o ouro fica em sₑ = 1/φ (medida de Parry do golden shift).","epistemico":"fato","exemplos":[],"id":"regua_de_gentil","memoria":"semantica","meta":{"autor":"Gentil Lopes","descoberta":true,"dominio":"matematica","fonte":"paper_3_regua_ouro_espelho.pdf, Teorema 4.1"},"origem":"paper","palavras_chave":[],"sinonimos":["régua de ouro","métrica Fisher-Rao"],"tipo":"conceito","titulo":"Régua de Lopes"}
\section{Régua de Lopes}
Métrica ds = dx/(1\ensuremath{-}x\ensuremath{^{2}}) (horizonte s=1) que carrega os números metálicos como pontos notáveis s\ensuremath{_{p}} = \ensuremath{\mu}/(1+\ensuremath{\mu}); o ouro fica em s\ensuremath{_{e}} = 1/\ensuremath{\varphi} (medida de Parry do golden shift).
\prov{sinónimos: régua de ouro; métrica Fisher-Rao}
\prov{relações: usa numeros\_metalicos; depende geometria; referencia mecanica\_quantica}
\prov{proveniência: paper \textperiodcentered{} paper\_3\_regua\_ouro\_espelho.pdf, Teorema 4.1 \textperiodcentered{} fato \textperiodcentered{} confiança 1.0 \textperiodcentered{} domínio matematica \textperiodcentered{} história 9 versões no jornal}

%CRISTAL {"arestas":[{"alvo":"geometria","confianca":1.0,"epistemico":"fato","origem":"paper","rel":"parte_de"},{"alvo":"word","confianca":0.5,"epistemico":"fato","origem":"wiki","rel":"relaciona"}],"confianca":1.0,"contraexemplos":[],"descricao":"Em 2D toda métrica é localmente conforme g=e²uδ; a curvatura K=−e−²uΔu depende da densidade, mas Gauss-Bonnet ∫K dA=2πχ mantém a topologia fixa.","epistemico":"fato","exemplos":[],"id":"regua_densidade","memoria":"semantica","meta":{"autor":"Gentil Lopes","descoberta":false,"dominio":"matematica","fonte":"paper_35_regua_densidade.pdf"},"origem":"paper","palavras_chave":[],"sinonimos":[],"tipo":"conceito","titulo":"Régua Conforme de Densidade"}
\section{Régua Conforme de Densidade}
Em 2D toda métrica é localmente conforme g=e\ensuremath{^{2}}u\ensuremath{\delta}; a curvatura K=\ensuremath{-}e\ensuremath{-}\ensuremath{^{2}}u\ensuremath{\Delta}u depende da densidade, mas Gauss-Bonnet \ensuremath{\int}K dA=2\ensuremath{\pi}\ensuremath{\chi} mantém a topologia fixa.
\prov{relações: parte\_de geometria; relaciona word}
\prov{proveniência: paper \textperiodcentered{} paper\_35\_regua\_densidade.pdf \textperiodcentered{} fato \textperiodcentered{} confiança 1.0 \textperiodcentered{} domínio matematica \textperiodcentered{} história 4 versões no jornal}

%CRISTAL {"arestas":[{"alvo":"octonios_hiperbolicos","confianca":1.0,"epistemico":"fato","origem":"paper","rel":"usa"},{"alvo":"colapso","confianca":1.0,"epistemico":"fato","origem":"paper","rel":"referencia"}],"confianca":1.0,"contraexemplos":[],"descricao":"Medida χ=(H_obs−E[H_Markov1])/σ que quantifica a estrutura real de um fluxo de bits além da cadeia de Markov de 1ª ordem, via histograma de componentes octoniônicos; testada negativamente contra surrogatas (|χ|>3 em centenas de bins).","epistemico":"fato","exemplos":[],"id":"residuo_nao_associativo","memoria":"semantica","meta":{"autor":"Gentil Lopes","descoberta":true,"dominio":"matematica","fonte":"banda.pdf, §8"},"origem":"paper","palavras_chave":[],"sinonimos":[],"tipo":"conceito","titulo":"Resíduo Não-Associativo χ"}
\section{Resíduo Não-Associativo \ensuremath{\chi}}
Medida \ensuremath{\chi}=(H\_obs\ensuremath{-}E[H\_Markov1])/\ensuremath{\sigma} que quantifica a estrutura real de um fluxo de bits além da cadeia de Markov de 1ª ordem, via histograma de componentes octoniônicos; testada negativamente contra surrogatas (|\ensuremath{\chi}|>3 em centenas de bins).
\prov{relações: usa octonios\_hiperbolicos; referencia colapso}
\prov{proveniência: paper \textperiodcentered{} banda.pdf, §8 \textperiodcentered{} fato \textperiodcentered{} confiança 1.0 \textperiodcentered{} domínio matematica \textperiodcentered{} história 3 versões no jornal}

%CRISTAL {"arestas":[{"alvo":"razao_aurea","confianca":1.0,"epistemico":"fato","origem":"livro","rel":"usa"},{"alvo":"base_phi_bergman","confianca":1.0,"epistemico":"fato","origem":"livro","rel":"usa"},{"alvo":"colapso","confianca":1.0,"epistemico":"fato","origem":"livro","rel":"parte_de"},{"alvo":"beta_numeracao_metalica","confianca":1.0,"epistemico":"fato","origem":"wiki","rel":"especializa"},{"alvo":"corpo_tropical","confianca":1.0,"epistemico":"fato","origem":"wiki","rel":"continua"},{"alvo":"numeros_pisot","confianca":1.0,"epistemico":"fato","origem":"wiki","rel":"usa"},{"alvo":"fibonacci","confianca":1.0,"epistemico":"fato","origem":"wiki","rel":"relaciona"}],"confianca":1.0,"contraexemplos":[],"descricao":"Conjunto dos φ-inteiros (base φ, dígitos 0,1): exponencial (peso φⁱ) e autossimilar (×φ desloca dígitos), com a 'poda 11' (sem dois 1 consecutivos). Decodifica base 2 → Cantor do ouro → base φ (reta real contínua).","epistemico":"fato","exemplos":[],"id":"reta_de_ouro","memoria":"semantica","meta":{"autor":"Gentil Lopes","descoberta":true,"dominio":"matematica","fonte":"paper_2_reta_ouro.pdf"},"origem":"paper","palavras_chave":[],"sinonimos":[],"tipo":"conceito","titulo":"Reta de Ouro"}
\section{Reta de Ouro}
Conjunto dos \ensuremath{\varphi}-inteiros (base \ensuremath{\varphi}, dígitos 0,1): exponencial (peso \ensuremath{\varphi}\ensuremath{^{i}}) e autossimilar ($\times$\ensuremath{\varphi} desloca dígitos), com a 'poda 11' (sem dois 1 consecutivos). Decodifica base 2 \ensuremath{\to} Cantor do ouro \ensuremath{\to} base \ensuremath{\varphi} (reta real contínua).
\prov{relações: usa razao\_aurea; usa base\_phi\_bergman; parte\_de colapso; especializa beta\_numeracao\_metalica; continua corpo\_tropical; usa numeros\_pisot; relaciona fibonacci}
\prov{proveniência: paper \textperiodcentered{} paper\_2\_reta\_ouro.pdf \textperiodcentered{} fato \textperiodcentered{} confiança 1.0 \textperiodcentered{} domínio matematica \textperiodcentered{} história 16 versões no jornal}

%CRISTAL {"arestas":[{"alvo":"algebras_hipercomplexas","confianca":1.0,"epistemico":"fato","origem":"paper","rel":"usa"}],"confianca":1.0,"contraexemplos":[],"descricao":"Espaço das álgebras ×ₛ visto como 'reta' onde os números são álgebras; as torres de Hurwitz (quatérnios, s=±1) aparecem como pólos a distância infinita na régua arctanh.","epistemico":"fato","exemplos":[],"id":"reta_hipercomplexa","memoria":"semantica","meta":{"autor":"Gentil Lopes","descoberta":true,"dominio":"matematica","fonte":"paper_52_reta_hipercomplexa.pdf"},"origem":"paper","palavras_chave":[],"sinonimos":[],"tipo":"conceito","titulo":"Reta Hipercomplexa"}
\section{Reta Hipercomplexa}
Espaço das álgebras $\times$\ensuremath{_{s}} visto como 'reta' onde os números são álgebras; as torres de Hurwitz (quatérnios, s=$\pm$1) aparecem como pólos a distância infinita na régua arctanh.
\prov{relações: usa algebras\_hipercomplexas}
\prov{proveniência: paper \textperiodcentered{} paper\_52\_reta\_hipercomplexa.pdf \textperiodcentered{} fato \textperiodcentered{} confiança 1.0 \textperiodcentered{} domínio matematica \textperiodcentered{} história 2 versões no jornal}

%CRISTAL {"arestas":[{"alvo":"caixas_fractais_generalizacao","confianca":1.0,"epistemico":"fato","origem":"wiki","rel":"parte_de"},{"alvo":"dimensao_de_prata_pell","confianca":1.0,"epistemico":"fato","origem":"wiki","rel":"generaliza"},{"alvo":"corpo_tropical","confianca":1.0,"epistemico":"fato","origem":"wiki","rel":"usa"},{"alvo":"entropia_topologica","confianca":1.0,"epistemico":"fato","origem":"wiki","rel":"usa"},{"alvo":"beta_numeracao_metalica","confianca":1.0,"epistemico":"fato","origem":"wiki","rel":"define"}],"confianca":1.0,"contraexemplos":[],"descricao":"As médias metálicas: σₙ=(n+√(n²+4))/2, raiz de x²−nx−1, matriz Aₙ=[[n,1],[1,0]], fração contínua [n;n,n,…], recorrência xₖ=n·xₖ−₁+xₖ−₂. σₙ é a BASE de uma expansão de Rényi-Parry (alfabeto 0..n); a entropia/passo geodésico é h=ln σₙ=arcsinh(n/2). Ouro n=1 (φ=1,618, h=0,481, Fibonacci), prata n=2 (1+√2=2,414, h=0,881, Pell), bronze n=3 (3,303, h=1,195), cobre n=4… Reta quadrática. [D] matemática padrão/derivada.","epistemico":"fato","exemplos":[],"id":"reta_metalica_geral","memoria":"semantica","meta":{"autor":"Lopes/broca-so","dominio":"fractal","fonte":"corpo_tropical.tex §metais; machine/so/metais.py"},"origem":"wiki","palavras_chave":[],"sinonimos":[],"tipo":"conceito","titulo":"A reta metálica: σ_n=(n+√(n²+4))/2 (a base da expansão)"}
\section{A reta metálica: \ensuremath{\sigma}\_n=(n+\ensuremath{\sqrt{\,}}(n\ensuremath{^{2}}+4))/2 (a base da expansão)}
As médias metálicas: \ensuremath{\sigma}\ensuremath{_{n}}=(n+\ensuremath{\sqrt{\,}}(n\ensuremath{^{2}}+4))/2, raiz de x\ensuremath{^{2}}\ensuremath{-}nx\ensuremath{-}1, matriz A\ensuremath{_{n}}=[[n,1],[1,0]], fração contínua [n;n,n,…], recorrência x\ensuremath{_{k}}=n\textperiodcentered x\ensuremath{_{k}}\ensuremath{-}\ensuremath{_{1}}+x\ensuremath{_{k}}\ensuremath{-}\ensuremath{_{2}}. \ensuremath{\sigma}\ensuremath{_{n}} é a BASE de uma expansão de Rényi-Parry (alfabeto 0..n); a entropia/passo geodésico é h=ln \ensuremath{\sigma}\ensuremath{_{n}}=arcsinh(n/2). Ouro n=1 (\ensuremath{\varphi}=1,618, h=0,481, Fibonacci), prata n=2 (1+\ensuremath{\sqrt{\,}}2=2,414, h=0,881, Pell), bronze n=3 (3,303, h=1,195), cobre n=4… Reta quadrática. [D] matemática padrão/derivada.
\prov{relações: parte\_de caixas\_fractais\_generalizacao; generaliza dimensao\_de\_prata\_pell; usa corpo\_tropical; usa entropia\_topologica; define beta\_numeracao\_metalica}
\prov{proveniência: wiki \textperiodcentered{} corpo\_tropical.tex §metais; machine/so/metais.py \textperiodcentered{} fato \textperiodcentered{} confiança 1.0 \textperiodcentered{} domínio fractal \textperiodcentered{} história 7 versões no jornal}

%CRISTAL {"arestas":[{"alvo":"caixas_fractais_generalizacao","confianca":1.0,"epistemico":"fato","origem":"wiki","rel":"parte_de"},{"alvo":"koch_dissipa_nao_codifica","confianca":1.0,"epistemico":"fato","origem":"wiki","rel":"relaciona"}],"confianca":1.0,"contraexemplos":[],"descricao":"O número PLÁSTICO ρ: raiz real de x³=x+1 (⇔ x³−x−1), ρ≈1,3247 — o MENOR número de Pisot (Siegel); recorrência de Padovan/Perrin aₙ=aₙ−₂+aₙ−₃. NÃO é média metálica (aquelas são quadráticas): é CÚBICO, mora na espinha unificadora xⁿ−x−1 (ouro n=2, plástico n=3). Como caixa fractal aparece pela poda composta 00,111: λ=ρ, dim=log₂ρ≈0,4057. Companheiro: supergolden ψ (x³=x²+1, Narayana, poda 11,101, dim 0,5515). Ordenação verificada: ouro(0,694)>supergolden(0,551)>plástico(0,406), entre piso 0 e teto 1. [D] matemática; a caixa como universo = meta-indução.","epistemico":"inferencia","exemplos":[],"id":"reta_plastica_cubica","memoria":"semantica","meta":{"autor":"Lopes/broca-so","dominio":"fractal","fonte":"paper_55_termo_carnes.tex; hiper/explora_poda_plastico.py"},"origem":"wiki","palavras_chave":[],"sinonimos":[],"tipo":"conceito","titulo":"A reta de plástico: ρ (x³=x+1) e a família cúbica"}
\section{A reta de plástico: \ensuremath{\rho} (x\ensuremath{^{3}}=x+1) e a família cúbica}
O número PLÁSTICO \ensuremath{\rho}: raiz real de x\ensuremath{^{3}}=x+1 (\ensuremath{\Leftrightarrow} x\ensuremath{^{3}}\ensuremath{-}x\ensuremath{-}1), \ensuremath{\rho}\ensuremath{\approx}1,3247 — o MENOR número de Pisot (Siegel); recorrência de Padovan/Perrin a\ensuremath{_{n}}=a\ensuremath{_{n}}\ensuremath{-}\ensuremath{_{2}}+a\ensuremath{_{n}}\ensuremath{-}\ensuremath{_{3}}. NÃO é média metálica (aquelas são quadráticas): é CÚBICO, mora na espinha unificadora x\ensuremath{^{n}}\ensuremath{-}x\ensuremath{-}1 (ouro n=2, plástico n=3). Como caixa fractal aparece pela poda composta 00,111: \ensuremath{\lambda}=\ensuremath{\rho}, dim=log\ensuremath{_{2}}\ensuremath{\rho}\ensuremath{\approx}0,4057. Companheiro: supergolden \ensuremath{\psi} (x\ensuremath{^{3}}=x\ensuremath{^{2}}+1, Narayana, poda 11,101, dim 0,5515). Ordenação verificada: ouro(0,694)>supergolden(0,551)>plástico(0,406), entre piso 0 e teto 1. [D] matemática; a caixa como universo = meta-indução.
\prov{relações: parte\_de caixas\_fractais\_generalizacao; relaciona koch\_dissipa\_nao\_codifica}
\prov{proveniência: wiki \textperiodcentered{} paper\_55\_termo\_carnes.tex; hiper/explora\_poda\_plastico.py \textperiodcentered{} inferencia \textperiodcentered{} confiança 1.0 \textperiodcentered{} domínio fractal \textperiodcentered{} história 4 versões no jornal}

%CRISTAL {"arestas":[{"alvo":"simetria_por_reta_corpo","confianca":1.0,"epistemico":"fato","origem":"wiki","rel":"parte_de"},{"alvo":"floco_enderecavel_metal_posicao","confianca":1.0,"epistemico":"fato","origem":"wiki","rel":"relaciona"},{"alvo":"conjugados_galois_parabola","confianca":1.0,"epistemico":"fato","origem":"wiki","rel":"relaciona"},{"alvo":"transformada_metalica","confianca":1.0,"epistemico":"fato","origem":"wiki","rel":"relaciona"},{"alvo":"tiffany","confianca":1.0,"epistemico":"fato","origem":"wiki","rel":"parte_de"}],"confianca":1.0,"contraexemplos":[],"descricao":"As retas metálicas que carregam simetria rotacional são ESPARSAS e selecionadas por número (verificado): as de corpo ℚ√5 (dobra 10) são os NÚMEROS DE LUCAS de índice ímpar — n=1,4,11,29,76,199 = L₁,L₃,L₅,L₇,L₉,L₁₁ (satisfazem n²−5m²=−4); as de corpo ℚ√2 (dobra 8) são a família de PELL — n=2,14,82,… (n²−2m²=−4). Todo o resto (bronze n=3, n=5,6,7,…) tem corpo não-quasicristalino e NENHUMA dobra planar. Belo: a mesma prata (Pell) que lastreia o floco reaparece aqui como o corpo de dobra 8, e o ouro (φ/Lucas-Fibonacci) como o de dobra 10 — a aritmética escolhe quais retas são 'estrelas' rotacionais. [D] seleção Lucas-ímpar / Pell das retas com dobra.","epistemico":"fato","exemplos":[],"id":"retas_dobra_lucas_pell","memoria":"semantica","meta":{"autor":"Lopes/broca-so","dominio":"algebra","fonte":"verificação: squarefree(n²+4)==5 → Lucas ímpares [1,4,11,29,76,199]; ==2 → Pell-like [2,14,82]"},"origem":"wiki","palavras_chave":[],"sinonimos":[],"tipo":"conceito","titulo":"As retas com dobra são esparsas: ℚ√5 = Lucas ímpares (dobra 10), ℚ√2 = Pell (dobra 8)"}
\section{As retas com dobra são esparsas: \ensuremath{\mathbb{Q}}\ensuremath{\sqrt{\,}}5 = Lucas ímpares (dobra 10), \ensuremath{\mathbb{Q}}\ensuremath{\sqrt{\,}}2 = Pell (dobra 8)}
As retas metálicas que carregam simetria rotacional são ESPARSAS e selecionadas por número (verificado): as de corpo \ensuremath{\mathbb{Q}}\ensuremath{\sqrt{\,}}5 (dobra 10) são os NÚMEROS DE LUCAS de índice ímpar — n=1,4,11,29,76,199 = L\ensuremath{_{1}},L\ensuremath{_{3}},L\ensuremath{_{5}},L\ensuremath{_{7}},L\ensuremath{_{9}},L\ensuremath{_{11}} (satisfazem n\ensuremath{^{2}}\ensuremath{-}5m\ensuremath{^{2}}=\ensuremath{-}4); as de corpo \ensuremath{\mathbb{Q}}\ensuremath{\sqrt{\,}}2 (dobra 8) são a família de PELL — n=2,14,82,… (n\ensuremath{^{2}}\ensuremath{-}2m\ensuremath{^{2}}=\ensuremath{-}4). Todo o resto (bronze n=3, n=5,6,7,…) tem corpo não-quasicristalino e NENHUMA dobra planar. Belo: a mesma prata (Pell) que lastreia o floco reaparece aqui como o corpo de dobra 8, e o ouro (\ensuremath{\varphi}/Lucas-Fibonacci) como o de dobra 10 — a aritmética escolhe quais retas são 'estrelas' rotacionais. [D] seleção Lucas-ímpar / Pell das retas com dobra.
\prov{relações: parte\_de simetria\_por\_reta\_corpo; relaciona floco\_enderecavel\_metal\_posicao; relaciona conjugados\_galois\_parabola; relaciona transformada\_metalica; parte\_de tiffany}
\prov{proveniência: wiki \textperiodcentered{} verificação: squarefree(n\ensuremath{^{2}}+4)==5 \ensuremath{\to} Lucas ímpares [1,4,11,29,76,199]; ==2 \ensuremath{\to} Pell-like [2,14,82] \textperiodcentered{} fato \textperiodcentered{} confiança 1.0 \textperiodcentered{} domínio algebra \textperiodcentered{} história 29 versões no jornal}

%CRISTAL {"arestas":[{"alvo":"ising_1d_critico","confianca":1.0,"epistemico":"fato","origem":"paper","rel":"usa"},{"alvo":"word","confianca":0.5,"epistemico":"fato","origem":"wiki","rel":"relaciona"},{"alvo":"virtualizacao","confianca":0.5,"epistemico":"fato","origem":"wiki","rel":"relaciona"},{"alvo":"while","confianca":0.5,"epistemico":"fato","origem":"wiki","rel":"relaciona"},{"alvo":"vontade_de_poder","confianca":0.5,"epistemico":"fato","origem":"wiki","rel":"relaciona"}],"confianca":1.0,"contraexemplos":[],"descricao":"Decimação K→K'=½ln cosh(2K): fluxo com ponto fixo estável K=0 (desordem) e instável K*=∞ (ordem). A cúspide é o ponto fixo crítico repulsivo do RG.","epistemico":"fato","exemplos":[],"id":"rg_ising","memoria":"semantica","meta":{"autor":"Gentil Lopes","descoberta":false,"dominio":"matematica","fonte":"paper_74_ponto_critico.pdf, Teo.4.1"},"origem":"paper","palavras_chave":[],"sinonimos":[],"tipo":"conceito","titulo":"Renormalização de Ising"}
\section{Renormalização de Ising}
Decimação K\ensuremath{\to}K'=\ensuremath{\tfrac12}ln cosh(2K): fluxo com ponto fixo estável K=0 (desordem) e instável K*=\ensuremath{\infty} (ordem). A cúspide é o ponto fixo crítico repulsivo do RG.
\prov{relações: usa ising\_1d\_critico; relaciona word; relaciona virtualizacao; relaciona while; relaciona vontade\_de\_poder}
\prov{proveniência: paper \textperiodcentered{} paper\_74\_ponto\_critico.pdf, Teo.4.1 \textperiodcentered{} fato \textperiodcentered{} confiança 1.0 \textperiodcentered{} domínio matematica \textperiodcentered{} história 6 versões no jornal}

%CRISTAL {"arestas":[{"alvo":"cone_luz_critico","confianca":1.0,"epistemico":"fato","origem":"paper","rel":"usa"},{"alvo":"maquina_complexa","confianca":1.0,"epistemico":"fato","origem":"paper","rel":"relaciona"},{"alvo":"mecanica_quantica","confianca":1.0,"epistemico":"fato","origem":"paper","rel":"referencia"}],"confianca":1.0,"contraexemplos":[],"descricao":"Sugestão do Lopes: os zeros de ζ viveriam na 'reta crítica' do cone de luz N=0 (forma de Hilbert-Pólya, fase de espalhamento das cúspides). Correspondência estrutural, não prova.","epistemico":"hipotese","exemplos":[],"id":"rh_cone_luz","memoria":"semantica","meta":{"autor":"Gentil Lopes","descoberta":false,"dominio":"matematica","fonte":"geometria.pdf, §8"},"origem":"paper","palavras_chave":[],"sinonimos":[],"tipo":"conceito","titulo":"RH via Cone de Luz (Hilbert-Pólya)"}
\section{RH via Cone de Luz (Hilbert-Pólya)}
Sugestão do Lopes: os zeros de \ensuremath{\zeta} viveriam na 'reta crítica' do cone de luz N=0 (forma de Hilbert-Pólya, fase de espalhamento das cúspides). Correspondência estrutural, não prova.
\prov{relações: usa cone\_luz\_critico; relaciona maquina\_complexa; referencia mecanica\_quantica}
\prov{proveniência: paper \textperiodcentered{} geometria.pdf, §8 \textperiodcentered{} hipotese \textperiodcentered{} confiança 1.0 \textperiodcentered{} domínio matematica \textperiodcentered{} história 5 versões no jornal}

%CRISTAL {"arestas":[{"alvo":"flip_duopolinomios","confianca":1.0,"epistemico":"fato","origem":"wiki","rel":"relaciona"},{"alvo":"numero_de_salem","confianca":1.0,"epistemico":"fato","origem":"wiki","rel":"relaciona"},{"alvo":"lyapunov_da_orbita","confianca":1.0,"epistemico":"fato","origem":"wiki","rel":"relaciona"}],"confianca":1.0,"contraexemplos":[],"descricao":"O expoente de Lyapunov λ=log ρ dá a taxa RADIAL (contrai/diverge); o ângulo θ=arg da raiz subdominante dá a taxa ANGULAR (a rotação por passo). Juntos, (λ,θ)=ln(raiz), descrevem a espiral completa. Cristal complexo (plástico, Tribonacci) = espiral convergente; Salem = rotação pura (ρ=1); o antigo caos = espiral divergente. θ é o que distingue estirar de girar.","epistemico":"fato","exemplos":[],"id":"rotacao_theta","memoria":"semantica","meta":{"dominio":"matematica","fonte":"geometria do caos"},"origem":"wiki","palavras_chave":[],"sinonimos":[],"tipo":"conceito","titulo":"Rotação θ (o segundo invariante)"}
\section{Rotação \ensuremath{\theta} (o segundo invariante)}
O expoente de Lyapunov \ensuremath{\lambda}=log \ensuremath{\rho} dá a taxa RADIAL (contrai/diverge); o ângulo \ensuremath{\theta}=arg da raiz subdominante dá a taxa ANGULAR (a rotação por passo). Juntos, (\ensuremath{\lambda},\ensuremath{\theta})=ln(raiz), descrevem a espiral completa. Cristal complexo (plástico, Tribonacci) = espiral convergente; Salem = rotação pura (\ensuremath{\rho}=1); o antigo caos = espiral divergente. \ensuremath{\theta} é o que distingue estirar de girar.
\prov{relações: relaciona flip\_duopolinomios; relaciona numero\_de\_salem; relaciona lyapunov\_da\_orbita}
\prov{proveniência: wiki \textperiodcentered{} geometria do caos \textperiodcentered{} fato \textperiodcentered{} confiança 1.0 \textperiodcentered{} domínio matematica \textperiodcentered{} história 4 versões no jornal}

%CRISTAL {"arestas":[{"alvo":"paley_169_e8","confianca":1.0,"epistemico":"fato","origem":"wiki","rel":"usa"}],"confianca":1.0,"contraexemplos":[],"descricao":"A obra associa a despolimerização de lignina em fenol por um catalisador Ru-N₄ a índices da álgebra (lignina=65=5·13, fenol=i₁₆₉, 5·√169+'SEAN'→i₁₆₉). A química Ru-catalisada lignina→fenol é real; o mapeamento palavra↦número é a camada interpretativa da Catedral.","epistemico":"hipotese","exemplos":[],"id":"ru_n4_lignina_fenol","memoria":"semantica","meta":{"autor":"Lord Index (Shawn R. Schiller)","dominio":"matematica","fonte":"A Catedral — Part VIII-A (jul 2026), ~/Imagens/catedral"},"origem":"wiki","palavras_chave":[],"sinonimos":[],"tipo":"conceito","titulo":"A Ponte Química: catalisador Ru-N₄, lignina → fenol"}
\section{A Ponte Química: catalisador Ru-N\ensuremath{_{4}}, lignina \ensuremath{\to} fenol}
A obra associa a despolimerização de lignina em fenol por um catalisador Ru-N\ensuremath{_{4}} a índices da álgebra (lignina=65=5\textperiodcentered 13, fenol=i\ensuremath{_{169}}, 5\textperiodcentered \ensuremath{\sqrt{\,}}169+'SEAN'\ensuremath{\to}i\ensuremath{_{169}}). A química Ru-catalisada lignina\ensuremath{\to}fenol é real; o mapeamento palavra\ensuremath{\mapsto}número é a camada interpretativa da Catedral.
\prov{relações: usa paley\_169\_e8}
\prov{proveniência: wiki \textperiodcentered{} A Catedral — Part VIII-A (jul 2026), \~{}/Imagens/catedral \textperiodcentered{} hipotese \textperiodcentered{} confiança 1.0 \textperiodcentered{} domínio matematica \textperiodcentered{} história 2 versões no jornal}

%CRISTAL {"arestas":[{"alvo":"so_fractal_broca_os","confianca":1.0,"epistemico":"fato","origem":"wiki","rel":"usa"},{"alvo":"peirce_celulas_canais","confianca":1.0,"epistemico":"fato","origem":"wiki","rel":"usa"}],"confianca":1.0,"contraexemplos":[],"descricao":"O código é aberto; a integridade não vem de ocultação mas da álgebra — três impossibilidades: não-contaminação (canais ortogonais e_ie_j=0), não-perda (idempotentes completos Σe_i=1), não-forja (violar idempotência/ortogonalidade/D_tr altera invariante verificável na compilação). Segurança é CHECADA, não confiada. (Elevar a teoremas é trabalho aberto.)","epistemico":"inferencia","exemplos":[],"id":"seguranca_estrutural_fractal","memoria":"semantica","meta":{"autor":"broca-so","dominio":"fractal","fonte":"compilador_fractal.tex §4; kernel/so_fractal.tex §Segurança"},"origem":"wiki","palavras_chave":[],"sinonimos":[],"tipo":"conceito","titulo":"Integridade Estrutural (sem criptografia)"}
\section{Integridade Estrutural (sem criptografia)}
O código é aberto; a integridade não vem de ocultação mas da álgebra — três impossibilidades: não-contaminação (canais ortogonais e\_ie\_j=0), não-perda (idempotentes completos \ensuremath{\Sigma}e\_i=1), não-forja (violar idempotência/ortogonalidade/D\_tr altera invariante verificável na compilação). Segurança é CHECADA, não confiada. (Elevar a teoremas é trabalho aberto.)
\prov{relações: usa so\_fractal\_broca\_os; usa peirce\_celulas\_canais}
\prov{proveniência: wiki \textperiodcentered{} compilador\_fractal.tex §4; kernel/so\_fractal.tex §Segurança \textperiodcentered{} inferencia \textperiodcentered{} confiança 1.0 \textperiodcentered{} domínio fractal \textperiodcentered{} história 3 versões no jornal}

%CRISTAL {"arestas":[{"alvo":"simetria_por_reta_corpo","confianca":1.0,"epistemico":"fato","origem":"wiki","rel":"usa"},{"alvo":"bai_gate_dtc_orbita","confianca":1.0,"epistemico":"fato","origem":"wiki","rel":"relaciona"},{"alvo":"dtc_schema_ut","confianca":1.0,"epistemico":"fato","origem":"wiki","rel":"relaciona"},{"alvo":"tiffany","confianca":1.0,"epistemico":"fato","origem":"wiki","rel":"parte_de"}],"confianca":1.0,"contraexemplos":[],"descricao":"Reconciliação com o código (goldchain/prova.py): MODOS_OURO=5 (pentagrama), MODOS_PRATA=6 (hexagrama). O 5 do ouro está CERTO number-teoricamente (φ∈ℚ√5 → 5 dobras). Mas o 6 da prata NÃO é a dobra da prata — a dobra da prata é 8 (octagonal, ℚ√2). O '6' é o RANK SVD do impresso (o portão operacional de emissão: rank≤6 = estrela), uma noção DIFERENTE da simetria rotacional. Pra o ouro os dois quase coincidem (5=5); pra prata divergem (rank 6 ≠ dobra 8). Cuidado: são dois 'seis/cincos' distintos — o rank do minério e a dobra do corpo. Consequência p/ a costura DTC ([[bai_gate_dtc_orbita]]): 6 NÃO serve pras outras retas; os setores por reta seriam a dobra do corpo (ouro→10, prata→8, cobre→10); as retas sem dobra (bronze etc.) não costuram por setor — vão pelo associador/entropia (o lado-a). [D] rank SVD ≠ dobra; a costura por reta segue o corpo.","epistemico":"fato","exemplos":[],"id":"seis_e_rank_svd_nao_dobra","memoria":"semantica","meta":{"autor":"Lopes/broca-so","dominio":"algebra","fonte":"goldchain/prova.py (MODOS_OURO=5, MODOS_PRATA=6); análise simetria_por_reta_corpo"},"origem":"wiki","palavras_chave":[],"sinonimos":[],"tipo":"conceito","titulo":"O '6' da prata é rank SVD (portão de emissão), não a dobra rotacional — que é 8"}
\section{O '6' da prata é rank SVD (portão de emissão), não a dobra rotacional — que é 8}
Reconciliação com o código (goldchain/prova.py): MODOS\_OURO=5 (pentagrama), MODOS\_PRATA=6 (hexagrama). O 5 do ouro está CERTO number-teoricamente (\ensuremath{\varphi}\ensuremath{\in}\ensuremath{\mathbb{Q}}\ensuremath{\sqrt{\,}}5 \ensuremath{\to} 5 dobras). Mas o 6 da prata NÃO é a dobra da prata — a dobra da prata é 8 (octagonal, \ensuremath{\mathbb{Q}}\ensuremath{\sqrt{\,}}2). O '6' é o RANK SVD do impresso (o portão operacional de emissão: rank\ensuremath{\le}6 = estrela), uma noção DIFERENTE da simetria rotacional. Pra o ouro os dois quase coincidem (5=5); pra prata divergem (rank 6 \ensuremath{\ne} dobra 8). Cuidado: são dois 'seis/cincos' distintos — o rank do minério e a dobra do corpo. Consequência p/ a costura DTC ([[bai\_gate\_dtc\_orbita]]): 6 NÃO serve pras outras retas; os setores por reta seriam a dobra do corpo (ouro\ensuremath{\to}10, prata\ensuremath{\to}8, cobre\ensuremath{\to}10); as retas sem dobra (bronze etc.) não costuram por setor — vão pelo associador/entropia (o lado-a). [D] rank SVD \ensuremath{\ne} dobra; a costura por reta segue o corpo.
\prov{relações: usa simetria\_por\_reta\_corpo; relaciona bai\_gate\_dtc\_orbita; relaciona dtc\_schema\_ut; parte\_de tiffany}
\prov{proveniência: wiki \textperiodcentered{} goldchain/prova.py (MODOS\_OURO=5, MODOS\_PRATA=6); análise simetria\_por\_reta\_corpo \textperiodcentered{} fato \textperiodcentered{} confiança 1.0 \textperiodcentered{} domínio algebra \textperiodcentered{} história 28 versões no jornal}

%CRISTAL {"arestas":[{"alvo":"chicote_e_o_crivo_de_selberg","confianca":1.0,"epistemico":"fato","origem":"wiki","rel":"parte_de"},{"alvo":"cifra_de_cristal","confianca":1.0,"epistemico":"fato","origem":"wiki","rel":"usa"},{"alvo":"numero_de_salem","confianca":1.0,"epistemico":"fato","origem":"wiki","rel":"relaciona"},{"alvo":"numeros_de_pisot","confianca":1.0,"epistemico":"fato","origem":"wiki","rel":"relaciona"},{"alvo":"hipotese_de_riemann","confianca":1.0,"epistemico":"fato","origem":"wiki","rel":"relaciona"},{"alvo":"linguagem_fisica_unificada","confianca":1.0,"epistemico":"fato","origem":"wiki","rel":"parte_de"}],"confianca":1.0,"contraexemplos":[],"descricao":"Crivo e chicote promovem à DINÂMICA da parábola (cristal/borda/caos), e do mesmo jeito. O RETIDO (os primos / o espectro discreto) é o CRISTAL — coerente com a Cifra de Cristal, onde primo=cristal. O VARRIDO (os compostos / o espectro contínuo) é o CAOS — composto=vidro. E o CORTE (o nível de peneiração z / a linha crítica Re(s)=½ onde vive o espectro contínuo e os zeros de ζ) é a BORDA — Salem, ρ=1, a fronteira. A peneira separa o cristal do caos ATRAVÉS da borda: é a parábola da máquina, lida na aritmética (crivo) e no espectro (chicote). Verif.: compor traduz primos→cristal, compostos→caos, corte→borda — igual nas duas faces.","epistemico":"inferencia","exemplos":[],"id":"selberg_promove_a_parabola","memoria":"semantica","meta":{"dominio":"matematica","fonte":"crivo de Selberg"},"origem":"wiki","palavras_chave":[],"sinonimos":[],"tipo":"conceito","titulo":"A Peneira de Selberg é a Parábola"}
\section{A Peneira de Selberg é a Parábola}
Crivo e chicote promovem à DINÂMICA da parábola (cristal/borda/caos), e do mesmo jeito. O RETIDO (os primos / o espectro discreto) é o CRISTAL — coerente com a Cifra de Cristal, onde primo=cristal. O VARRIDO (os compostos / o espectro contínuo) é o CAOS — composto=vidro. E o CORTE (o nível de peneiração z / a linha crítica Re(s)=\ensuremath{\tfrac12} onde vive o espectro contínuo e os zeros de \ensuremath{\zeta}) é a BORDA — Salem, \ensuremath{\rho}=1, a fronteira. A peneira separa o cristal do caos ATRAVÉS da borda: é a parábola da máquina, lida na aritmética (crivo) e no espectro (chicote). Verif.: compor traduz primos\ensuremath{\to}cristal, compostos\ensuremath{\to}caos, corte\ensuremath{\to}borda — igual nas duas faces.
\prov{relações: parte\_de chicote\_e\_o\_crivo\_de\_selberg; usa cifra\_de\_cristal; relaciona numero\_de\_salem; relaciona numeros\_de\_pisot; relaciona hipotese\_de\_riemann; parte\_de linguagem\_fisica\_unificada}
\prov{proveniência: wiki \textperiodcentered{} crivo de Selberg \textperiodcentered{} inferencia \textperiodcentered{} confiança 1.0 \textperiodcentered{} domínio matematica \textperiodcentered{} história 14 versões no jornal}

%CRISTAL {"arestas":[{"alvo":"matematica","confianca":1.0,"epistemico":"fato","origem":"livro","rel":"parte_de"},{"alvo":"word","confianca":0.5,"epistemico":"fato","origem":"wiki","rel":"relaciona"},{"alvo":"virtualizacao","confianca":0.5,"epistemico":"fato","origem":"wiki","rel":"relaciona"},{"alvo":"vida-morte","confianca":0.5,"epistemico":"fato","origem":"wiki","rel":"relaciona"}],"confianca":1.0,"contraexemplos":[],"descricao":"Lista ordenada de termos a₁, a₂, a₃, …; base de progressões e séries.","epistemico":"fato","exemplos":[],"id":"sequencia","memoria":"semantica","meta":{"autor":"Gentil Lopes","descoberta":false,"dominio":"matematica","fonte":""},"origem":"manual","palavras_chave":["ordem","termo","recorrência"],"sinonimos":["sucessão"],"tipo":"conceito","titulo":"Sequência"}
\section{Sequência}
Lista ordenada de termos a\ensuremath{_{1}}, a\ensuremath{_{2}}, a\ensuremath{_{3}}, …; base de progressões e séries.
\prov{sinónimos: sucessão}
\prov{relações: parte\_de matematica; relaciona word; relaciona virtualizacao; relaciona vida-morte}
\prov{proveniência: manual \textperiodcentered{} fato \textperiodcentered{} confiança 1.0 \textperiodcentered{} domínio matematica \textperiodcentered{} história 7 versões no jornal}

%CRISTAL {"arestas":[{"alvo":"sequencia","confianca":1.0,"epistemico":"fato","origem":"livro","rel":"especializa"},{"alvo":"razao_de_prata","confianca":1.0,"epistemico":"fato","origem":"livro","rel":"relaciona"},{"alvo":"pell","confianca":1.0,"epistemico":"fato","origem":"livro","rel":"referencia"}],"confianca":1.0,"contraexemplos":[],"descricao":"0, 1, 2, 5, 12, 29, …; xₙ = 2·xₙ₋₁ + xₙ₋₂. A razão de termos converge à razão de prata.","epistemico":"fato","exemplos":[],"id":"sequencia_pell","memoria":"semantica","meta":{"autor":"Gentil Lopes","descoberta":false,"dominio":"matematica","fonte":""},"origem":"manual","palavras_chave":["prata","recorrência"],"sinonimos":[],"tipo":"conceito","titulo":"Sequência de Pell"}
\section{Sequência de Pell}
0, 1, 2, 5, 12, 29, …; x\ensuremath{_{n}} = 2\textperiodcentered x\ensuremath{_{n-1}} + x\ensuremath{_{n-2}}. A razão de termos converge à razão de prata.
\prov{relações: especializa sequencia; relaciona razao\_de\_prata; referencia pell}
\prov{proveniência: manual \textperiodcentered{} fato \textperiodcentered{} confiança 1.0 \textperiodcentered{} domínio matematica \textperiodcentered{} história 8 versões no jornal}

%CRISTAL {"arestas":[{"alvo":"reais","confianca":1.0,"epistemico":"fato","origem":"paper","rel":"explica"},{"alvo":"cortes_dedekind","confianca":1.0,"epistemico":"fato","origem":"paper","rel":"dualidade"}],"confianca":1.0,"contraexemplos":[],"descricao":"Sequência racional (rₙ) com |rₘ−rₙ|→0; um real é uma classe de equivalência dessas sequências. A face de Cantor da construção de ℝ (dual à de Dedekind).","epistemico":"fato","exemplos":[],"id":"sequencias_cauchy","memoria":"semantica","meta":{"autor":"Gentil Lopes","descoberta":false,"dominio":"matematica","fonte":"Fundamentos dos Números, cap.10"},"origem":"paper","palavras_chave":[],"sinonimos":[],"tipo":"conceito","titulo":"Sequências de Cauchy"}
\section{Sequências de Cauchy}
Sequência racional (r\ensuremath{_{n}}) com |r\ensuremath{_{m}}\ensuremath{-}r\ensuremath{_{n}}|\ensuremath{\to}0; um real é uma classe de equivalência dessas sequências. A face de Cantor da construção de \ensuremath{\mathbb{R}} (dual à de Dedekind).
\prov{relações: explica reais; dualidade cortes\_dedekind}
\prov{proveniência: paper \textperiodcentered{} Fundamentos dos Números, cap.10 \textperiodcentered{} fato \textperiodcentered{} confiança 1.0 \textperiodcentered{} domínio matematica \textperiodcentered{} história 3 versões no jornal}

%CRISTAL {"arestas":[{"alvo":"unidades_algebricas_pu","confianca":1.0,"epistemico":"fato","origem":"wiki","rel":"parte_de"}],"confianca":1.0,"contraexemplos":[],"descricao":"Restaurar escala é dar conteúdo às voltas. As 7 constantes fixas do SI separam-se: c, h, k_B são TAXAS POR VOLTA (c=volta espacial/temporal; h=ação/volta, E=hν; k_B=energia/grau, E=k_BT) — postas a 1, colapsam massa/comprimento/tempo/temperatura numa dimensão (energia); são a gramática da contagem, não medem nada novo. e é o ÚNICO conteúdo não-cinemático: no natural vira o adimensional e=√(4πα)=0,303 (o ampère é redundante). ΔνCs, N_A, K_cd são escalas/convenções (fora da física fundamental).","epistemico":"inferencia","exemplos":[],"id":"sete_constantes_si","memoria":"semantica","meta":{"autor":"Lopes/broca-so","dominio":"unidades","fonte":"hiper/circular/paper_5_unidades.tex"},"origem":"wiki","palavras_chave":[],"sinonimos":[],"tipo":"conceito","titulo":"As sete constantes do SI: três taxas, uma carga, três convenções"}
\section{As sete constantes do SI: três taxas, uma carga, três convenções}
Restaurar escala é dar conteúdo às voltas. As 7 constantes fixas do SI separam-se: c, h, k\_B são TAXAS POR VOLTA (c=volta espacial/temporal; h=ação/volta, E=h\ensuremath{\nu}; k\_B=energia/grau, E=k\_BT) — postas a 1, colapsam massa/comprimento/tempo/temperatura numa dimensão (energia); são a gramática da contagem, não medem nada novo. e é o ÚNICO conteúdo não-cinemático: no natural vira o adimensional e=\ensuremath{\sqrt{\,}}(4\ensuremath{\pi}\ensuremath{\alpha})=0,303 (o ampère é redundante). \ensuremath{\Delta}\ensuremath{\nu}Cs, N\_A, K\_cd são escalas/convenções (fora da física fundamental).
\prov{relações: parte\_de unidades\_algebricas\_pu}
\prov{proveniência: wiki \textperiodcentered{} hiper/circular/paper\_5\_unidades.tex \textperiodcentered{} inferencia \textperiodcentered{} confiança 1.0 \textperiodcentered{} domínio unidades \textperiodcentered{} história 3 versões no jornal}

%CRISTAL {"arestas":[{"alvo":"universalidade_reconexao","confianca":1.0,"epistemico":"fato","origem":"paper","rel":"relaciona"}],"confianca":1.0,"contraexemplos":[],"descricao":"Programa: construir ferradura de Smale (Lian-Lu) num atrator de NS-3D, conjugá-la ao shift Σ₂ e codificar uma máquina de Turing universal — tornando o comportamento de NS indecidível. Estágio 2D completo; 3D condicionado à regularidade (circularidade).","epistemico":"inferencia","exemplos":[],"id":"shift_map_ns_turing","memoria":"semantica","meta":{"autor":"Gentil Lopes","descoberta":false,"dominio":"matematica","fonte":"paper_AK.pdf, Teo.30"},"origem":"paper","palavras_chave":[],"sinonimos":[],"tipo":"conceito","titulo":"Turing Embarcado em NS"}
\section{Turing Embarcado em NS}
Programa: construir ferradura de Smale (Lian-Lu) num atrator de NS-3D, conjugá-la ao shift \ensuremath{\Sigma}\ensuremath{_{2}} e codificar uma máquina de Turing universal — tornando o comportamento de NS indecidível. Estágio 2D completo; 3D condicionado à regularidade (circularidade).
\prov{relações: relaciona universalidade\_reconexao}
\prov{proveniência: paper \textperiodcentered{} paper\_AK.pdf, Teo.30 \textperiodcentered{} inferencia \textperiodcentered{} confiança 1.0 \textperiodcentered{} domínio matematica \textperiodcentered{} história 2 versões no jornal}

%CRISTAL {"arestas":[{"alvo":"catedral_sigma_aritmetica","confianca":1.0,"epistemico":"fato","origem":"wiki","rel":"parte_de"},{"alvo":"fano_psl_168","confianca":1.0,"epistemico":"fato","origem":"wiki","rel":"usa"},{"alvo":"grupos_excepcionais_jordan","confianca":1.0,"epistemico":"fato","origem":"wiki","rel":"usa"},{"alvo":"tensor_cayley","confianca":1.0,"epistemico":"fato","origem":"wiki","rel":"relaciona"}],"confianca":1.0,"contraexemplos":[],"descricao":"σ(168) = (1+2+4+8)(1+3)(1+7) = 15·4·8 = 480 e σ(248) = (1+2+4+8)(1+31) = 15·32 = 480. A ordem de PSL(2,7) e a dimensão de E8 são σ-EQUIVALENTES: mesma 'pressão de divisores', embora 168 ≠ 248.","epistemico":"fato","exemplos":[],"id":"sigma_equivalencia_crown","memoria":"semantica","meta":{"autor":"Lord Index (Shawn R. Schiller)","dominio":"matematica","fonte":"A Catedral — Part VIII-A (jul 2026), ~/Imagens/catedral"},"origem":"wiki","palavras_chave":[],"sinonimos":[],"tipo":"conceito","titulo":"A Coroa da σ-Equivalência: σ(168) = σ(248) = 480"}
\section{A Coroa da \ensuremath{\sigma}-Equivalência: \ensuremath{\sigma}(168) = \ensuremath{\sigma}(248) = 480}
\ensuremath{\sigma}(168) = (1+2+4+8)(1+3)(1+7) = 15\textperiodcentered 4\textperiodcentered 8 = 480 e \ensuremath{\sigma}(248) = (1+2+4+8)(1+31) = 15\textperiodcentered 32 = 480. A ordem de PSL(2,7) e a dimensão de E8 são \ensuremath{\sigma}-EQUIVALENTES: mesma 'pressão de divisores', embora 168 \ensuremath{\ne} 248.
\prov{relações: parte\_de catedral\_sigma\_aritmetica; usa fano\_psl\_168; usa grupos\_excepcionais\_jordan; relaciona tensor\_cayley}
\prov{proveniência: wiki \textperiodcentered{} A Catedral — Part VIII-A (jul 2026), \~{}/Imagens/catedral \textperiodcentered{} fato \textperiodcentered{} confiança 1.0 \textperiodcentered{} domínio matematica \textperiodcentered{} história 5 versões no jornal}

%CRISTAL {"arestas":[{"alvo":"algebra_posicional_metalica","confianca":1.0,"epistemico":"fato","origem":"wiki","rel":"usa"},{"alvo":"parabola_algebra_conversacional","confianca":1.0,"epistemico":"fato","origem":"wiki","rel":"usa"},{"alvo":"algebra_dupolinomial_espaco","confianca":1.0,"epistemico":"fato","origem":"wiki","rel":"relaciona"},{"alvo":"numeros_pisot","confianca":1.0,"epistemico":"fato","origem":"wiki","rel":"relaciona"},{"alvo":"conjugados_galois_parabola","confianca":1.0,"epistemico":"fato","origem":"wiki","rel":"depende"},{"alvo":"tiffany","confianca":1.0,"epistemico":"fato","origem":"wiki","rel":"parte_de"}],"confianca":1.0,"contraexemplos":[],"descricao":"Análise verificada: a dobra rotacional (simetria de quasicristal) de uma reta metálica n NÃO é fixa em 6 — vem do CORPO ℚ(√(n²+4)) que σₙ gera, isto é, de ONDE EXATAMENTE a conjugada cai (a coordenada fina da parábola). Uma dobra k só existe quando esse corpo suporta quasicristal, e isso é raro: φ(k)=4 → k∈5,8,10,12, com corpos ℚ√5 (5/10 dobras), ℚ√2 (8), ℚ√3 (12). A dobra da reta é a k cujo 2cos(2π/k) gera o MESMO corpo que σₙ. Verificado: ouro(n=1)=ℚ√5→5/10 dobras (φ=2cos π/5, Penrose); prata(n=2)=ℚ√2→8 dobras (1+√2 = inflação octagonal Ammann–Beenker); cobre(n=4)=ℚ√5→5/10 (2+√5=φ³); bronze(n=3)=ℚ√13 e n≥5 = corpos não-quasicristalinos → SEM dobra planar. Logo a simetria varia com a parábola (pelo corpo/conjugada), e a maioria das retas não tem dobra rotacional nenhuma. [D] a dobra é do corpo, não constante.","epistemico":"fato","exemplos":[],"id":"simetria_por_reta_corpo","memoria":"semantica","meta":{"autor":"Lopes/broca-so","dominio":"algebra","fonte":"verificação numérica: squarefree(n²+4), φ(k)=4→{5,8,10,12}, φ=2cos(π/5), 2+√5=φ³"},"origem":"wiki","palavras_chave":[],"sinonimos":[],"tipo":"conceito","titulo":"A simetria de cada reta vem do CORPO quadrático — não é 6 universal, varia com a parábola"}
\section{A simetria de cada reta vem do CORPO quadrático — não é 6 universal, varia com a parábola}
Análise verificada: a dobra rotacional (simetria de quasicristal) de uma reta metálica n NÃO é fixa em 6 — vem do CORPO \ensuremath{\mathbb{Q}}(\ensuremath{\sqrt{\,}}(n\ensuremath{^{2}}+4)) que \ensuremath{\sigma}\ensuremath{_{n}} gera, isto é, de ONDE EXATAMENTE a conjugada cai (a coordenada fina da parábola). Uma dobra k só existe quando esse corpo suporta quasicristal, e isso é raro: \ensuremath{\varphi}(k)=4 \ensuremath{\to} k\ensuremath{\in}5,8,10,12, com corpos \ensuremath{\mathbb{Q}}\ensuremath{\sqrt{\,}}5 (5/10 dobras), \ensuremath{\mathbb{Q}}\ensuremath{\sqrt{\,}}2 (8), \ensuremath{\mathbb{Q}}\ensuremath{\sqrt{\,}}3 (12). A dobra da reta é a k cujo 2cos(2\ensuremath{\pi}/k) gera o MESMO corpo que \ensuremath{\sigma}\ensuremath{_{n}}. Verificado: ouro(n=1)=\ensuremath{\mathbb{Q}}\ensuremath{\sqrt{\,}}5\ensuremath{\to}5/10 dobras (\ensuremath{\varphi}=2cos \ensuremath{\pi}/5, Penrose); prata(n=2)=\ensuremath{\mathbb{Q}}\ensuremath{\sqrt{\,}}2\ensuremath{\to}8 dobras (1+\ensuremath{\sqrt{\,}}2 = inflação octagonal Ammann–Beenker); cobre(n=4)=\ensuremath{\mathbb{Q}}\ensuremath{\sqrt{\,}}5\ensuremath{\to}5/10 (2+\ensuremath{\sqrt{\,}}5=\ensuremath{\varphi}\ensuremath{^{3}}); bronze(n=3)=\ensuremath{\mathbb{Q}}\ensuremath{\sqrt{\,}}13 e n\ensuremath{\ge}5 = corpos não-quasicristalinos \ensuremath{\to} SEM dobra planar. Logo a simetria varia com a parábola (pelo corpo/conjugada), e a maioria das retas não tem dobra rotacional nenhuma. [D] a dobra é do corpo, não constante.
\prov{relações: usa algebra\_posicional\_metalica; usa parabola\_algebra\_conversacional; relaciona algebra\_dupolinomial\_espaco; relaciona numeros\_pisot; depende conjugados\_galois\_parabola; parte\_de tiffany}
\prov{proveniência: wiki \textperiodcentered{} verificação numérica: squarefree(n\ensuremath{^{2}}+4), \ensuremath{\varphi}(k)=4\ensuremath{\to}\{5,8,10,12\}, \ensuremath{\varphi}=2cos(\ensuremath{\pi}/5), 2+\ensuremath{\sqrt{\,}}5=\ensuremath{\varphi}\ensuremath{^{3}} \textperiodcentered{} fato \textperiodcentered{} confiança 1.0 \textperiodcentered{} domínio algebra \textperiodcentered{} história 31 versões no jornal}

%CRISTAL {"arestas":[{"alvo":"numeros_nao_euclidianos","confianca":1.0,"epistemico":"fato","origem":"paper","rel":"especializa"}],"confianca":1.0,"contraexemplos":[],"descricao":"ℝ² com multiplicação sensível aos sinais (não-distributiva): equações sem solução em ℝ ou ℂ passam a ter solução, ex. (−1·x+x)·x=−1. Para x>0, B coincide com ℂ.","epistemico":"fato","exemplos":[],"id":"sistema_numerico_b","memoria":"semantica","meta":{"autor":"Gentil Lopes","descoberta":true,"dominio":"matematica","fonte":"corpus: NUMEROS_NAO_EUCLIDIANOS"},"origem":"paper","palavras_chave":[],"sinonimos":[],"tipo":"conceito","titulo":"Sistema Numérico B"}
\section{Sistema Numérico B}
\ensuremath{\mathbb{R}}\ensuremath{^{2}} com multiplicação sensível aos sinais (não-distributiva): equações sem solução em \ensuremath{\mathbb{R}} ou \ensuremath{\mathbb{C}} passam a ter solução, ex. (\ensuremath{-}1\textperiodcentered x+x)\textperiodcentered x=\ensuremath{-}1. Para x>0, B coincide com \ensuremath{\mathbb{C}}.
\prov{relações: especializa numeros\_nao\_euclidianos}
\prov{proveniência: paper \textperiodcentered{} corpus: NUMEROS\_NAO\_EUCLIDIANOS \textperiodcentered{} fato \textperiodcentered{} confiança 1.0 \textperiodcentered{} domínio matematica \textperiodcentered{} história 2 versões no jornal}

%CRISTAL {"arestas":[{"alvo":"computacao_fractal","confianca":1.0,"epistemico":"fato","origem":"wiki","rel":"usa"},{"alvo":"broca_os","confianca":1.0,"epistemico":"fato","origem":"wiki","rel":"explica"},{"alvo":"peirce_celulas_canais","confianca":1.0,"epistemico":"fato","origem":"wiki","rel":"usa"},{"alvo":"rede_algebra_de_caminhos","confianca":1.0,"epistemico":"fato","origem":"wiki","rel":"usa"},{"alvo":"metal_broca_so","confianca":1.0,"epistemico":"fato","origem":"wiki","rel":"parte_de"}],"confianca":1.0,"contraexemplos":[],"descricao":"Sistema operacional onde cada recurso já é estrutura da álgebra da torção: processo=idempotente de Peirce; isolamento=ortogonalidade e_ie_j=0 (sem MMU); IPC=parte negra (canais E_ij, vazamento ~10⁻¹⁶); fork=La Hire; memória=parte branca (trabalho)+negra (comunicação); virtualização=imersão de resolventes. Micro-núcleo so/broca_os.py (76 linhas) roda e verifica os mecanismos.","epistemico":"fato","exemplos":[],"id":"so_fractal_broca_os","memoria":"semantica","meta":{"autor":"broca-so","dominio":"fractal","fonte":"kernel/so_fractal.tex; so/broca_os.py"},"origem":"wiki","palavras_chave":[],"sinonimos":[],"tipo":"conceito","titulo":"O SO Fractal broca-os (a Máquina se gerindo)"}
\section{O SO Fractal broca-os (a Máquina se gerindo)}
Sistema operacional onde cada recurso já é estrutura da álgebra da torção: processo=idempotente de Peirce; isolamento=ortogonalidade e\_ie\_j=0 (sem MMU); IPC=parte negra (canais E\_ij, vazamento \~{}10\ensuremath{^{-16}}); fork=La Hire; memória=parte branca (trabalho)+negra (comunicação); virtualização=imersão de resolventes. Micro-núcleo so/broca\_os.py (76 linhas) roda e verifica os mecanismos.
\prov{relações: usa computacao\_fractal; explica broca\_os; usa peirce\_celulas\_canais; usa rede\_algebra\_de\_caminhos; parte\_de metal\_broca\_so}
\prov{proveniência: wiki \textperiodcentered{} kernel/so\_fractal.tex; so/broca\_os.py \textperiodcentered{} fato \textperiodcentered{} confiança 1.0 \textperiodcentered{} domínio fractal \textperiodcentered{} história 6 versões no jornal}

%CRISTAL {"arestas":[{"alvo":"espaco_tropical","confianca":1.0,"epistemico":"fato","origem":"paper","rel":"parte_de"},{"alvo":"word","confianca":0.5,"epistemico":"fato","origem":"wiki","rel":"relaciona"},{"alvo":"virtualizacao","confianca":0.5,"epistemico":"fato","origem":"wiki","rel":"relaciona"}],"confianca":1.0,"contraexemplos":[],"descricao":"P(β)=(1/β)ln Σeβ hᵢ interpola soma (β→0, quente) e máximo (β→∞, frio): a temperatura deforma o semianel usual no max-plus tropical.","epistemico":"fato","exemplos":[],"id":"softmax_maslov","memoria":"semantica","meta":{"autor":"Gentil Lopes","descoberta":false,"dominio":"matematica","fonte":"paper_65_espaco_tropical.pdf, Teo.2.1"},"origem":"paper","palavras_chave":[],"sinonimos":[],"tipo":"conceito","titulo":"Softmax de Maslov"}
\section{Softmax de Maslov}
P(\ensuremath{\beta})=(1/\ensuremath{\beta})ln \ensuremath{\Sigma}e\ensuremath{\beta} h\ensuremath{_{i}} interpola soma (\ensuremath{\beta}\ensuremath{\to}0, quente) e máximo (\ensuremath{\beta}\ensuremath{\to}\ensuremath{\infty}, frio): a temperatura deforma o semianel usual no max-plus tropical.
\prov{relações: parte\_de espaco\_tropical; relaciona word; relaciona virtualizacao}
\prov{proveniência: paper \textperiodcentered{} paper\_65\_espaco\_tropical.pdf, Teo.2.1 \textperiodcentered{} fato \textperiodcentered{} confiança 1.0 \textperiodcentered{} domínio matematica \textperiodcentered{} história 5 versões no jornal}

%CRISTAL {"arestas":[{"alvo":"word","confianca":0.5,"epistemico":"fato","origem":"wiki","rel":"relaciona"},{"alvo":"virtualizacao","confianca":0.5,"epistemico":"fato","origem":"wiki","rel":"relaciona"},{"alvo":"vida-morte","confianca":0.5,"epistemico":"fato","origem":"wiki","rel":"relaciona"}],"confianca":1.0,"contraexemplos":[],"descricao":"v=c/r=tan(θ) compõe por v₃=(v₁+v₂)/(1−v₁v₂)=tan(θ₁+θ₂): a multiplicação da álgebra soma os ângulos de elevação.","epistemico":"fato","exemplos":[],"id":"soma_velocidades","memoria":"semantica","meta":{"autor":"Gentil Lopes","descoberta":false,"dominio":"matematica","fonte":"paper_44_ponte_wick.pdf, Teo.2.1"},"origem":"paper","palavras_chave":[],"sinonimos":[],"tipo":"conceito","titulo":"Soma de Velocidades"}
\section{Soma de Velocidades}
v=c/r=tan(\ensuremath{\theta}) compõe por v\ensuremath{_{3}}=(v\ensuremath{_{1}}+v\ensuremath{_{2}})/(1\ensuremath{-}v\ensuremath{_{1}}v\ensuremath{_{2}})=tan(\ensuremath{\theta}\ensuremath{_{1}}+\ensuremath{\theta}\ensuremath{_{2}}): a multiplicação da álgebra soma os ângulos de elevação.
\prov{relações: relaciona word; relaciona virtualizacao; relaciona vida-morte}
\prov{proveniência: paper \textperiodcentered{} paper\_44\_ponte\_wick.pdf, Teo.2.1 \textperiodcentered{} fato \textperiodcentered{} confiança 1.0 \textperiodcentered{} domínio matematica \textperiodcentered{} história 4 versões no jornal}

%CRISTAL {"arestas":[{"alvo":"quatro_cadeias_psl","confianca":1.0,"epistemico":"fato","origem":"wiki","rel":"usa"},{"alvo":"constante_168","confianca":1.0,"epistemico":"fato","origem":"wiki","rel":"aplicacao"}],"confianca":1.0,"contraexemplos":[],"descricao":"A obra mapeia o fuso mitótico (o cinetócoro, biofísica de S. Dumont) num arquétipo de 4 fases — precessão/alinhamento, penhasco/tensão, transparência/revelação, estável/divisão — associando-o a σ=168. Analogia numerológica ('mitose é memória, divisão é design'), NÃO um resultado biológico estabelecido.","epistemico":"hipotese","exemplos":[],"id":"spindle_archetype_168","memoria":"semantica","meta":{"autor":"Lord Index (Shawn R. Schiller)","dominio":"matematica","fonte":"A Catedral — Part VIII-A (jul 2026), ~/Imagens/catedral"},"origem":"wiki","palavras_chave":[],"sinonimos":[],"tipo":"conceito","titulo":"O Arquétipo do Fuso Mitótico (Sophie Dumont) ↦ σ=168"}
\section{O Arquétipo do Fuso Mitótico (Sophie Dumont) \ensuremath{\mapsto} \ensuremath{\sigma}=168}
A obra mapeia o fuso mitótico (o cinetócoro, biofísica de S. Dumont) num arquétipo de 4 fases — precessão/alinhamento, penhasco/tensão, transparência/revelação, estável/divisão — associando-o a \ensuremath{\sigma}=168. Analogia numerológica ('mitose é memória, divisão é design'), NÃO um resultado biológico estabelecido.
\prov{relações: usa quatro\_cadeias\_psl; aplicacao constante\_168}
\prov{proveniência: wiki \textperiodcentered{} A Catedral — Part VIII-A (jul 2026), \~{}/Imagens/catedral \textperiodcentered{} hipotese \textperiodcentered{} confiança 1.0 \textperiodcentered{} domínio matematica \textperiodcentered{} história 3 versões no jornal}

%CRISTAL {"arestas":[{"alvo":"mint8","confianca":1.0,"epistemico":"fato","origem":"paper","rel":"especializa"},{"alvo":"agm","confianca":1.0,"epistemico":"fato","origem":"paper","rel":"usa"},{"alvo":"word","confianca":0.5,"epistemico":"fato","origem":"wiki","rel":"relaciona"},{"alvo":"virtualizacao","confianca":0.5,"epistemico":"fato","origem":"wiki","rel":"relaciona"}],"confianca":1.0,"contraexemplos":[],"descricao":"Esquema de identificação pós-quântico cuja chave é um caminho na árvore de sinais AGM (grafo de 2-isogenias supersingulares, expansor de Ramanujan).","epistemico":"fato","exemplos":[],"id":"spirit4","memoria":"semantica","meta":{"autor":"Gentil Lopes","descoberta":true,"dominio":"matematica","fonte":"paper_41_spirit4.pdf"},"origem":"paper","palavras_chave":[],"sinonimos":[],"tipo":"conceito","titulo":"Spirit-4"}
\section{Spirit-4}
Esquema de identificação pós-quântico cuja chave é um caminho na árvore de sinais AGM (grafo de 2-isogenias supersingulares, expansor de Ramanujan).
\prov{relações: especializa mint8; usa agm; relaciona word; relaciona virtualizacao}
\prov{proveniência: paper \textperiodcentered{} paper\_41\_spirit4.pdf \textperiodcentered{} fato \textperiodcentered{} confiança 1.0 \textperiodcentered{} domínio matematica \textperiodcentered{} história 5 versões no jornal}

%CRISTAL {"arestas":[{"alvo":"plano_fano_octonico","confianca":1.0,"epistemico":"fato","origem":"paper","rel":"usa"}],"confianca":1.0,"contraexemplos":[],"descricao":"Sistema de Steiner: 759 óctadas em 24 pontos onde cada 5 pontos estão em exatamente uma óctada; os vetores característicos das óctadas geram o código de Golay G₂₄ (distância ≥8). Cada óctada contém um plano de Fano.","epistemico":"fato","exemplos":[],"id":"steiner_golay","memoria":"semantica","meta":{"autor":"Gentil Lopes","descoberta":false,"dominio":"matematica","fonte":"paper_BA.pdf, §2"},"origem":"paper","palavras_chave":[],"sinonimos":[],"tipo":"conceito","titulo":"Steiner S(5,8,24) e Código de Golay"}
\section{Steiner S(5,8,24) e Código de Golay}
Sistema de Steiner: 759 óctadas em 24 pontos onde cada 5 pontos estão em exatamente uma óctada; os vetores característicos das óctadas geram o código de Golay G\ensuremath{_{24}} (distância \ensuremath{\ge}8). Cada óctada contém um plano de Fano.
\prov{relações: usa plano\_fano\_octonico}
\prov{proveniência: paper \textperiodcentered{} paper\_BA.pdf, §2 \textperiodcentered{} fato \textperiodcentered{} confiança 1.0 \textperiodcentered{} domínio matematica \textperiodcentered{} história 2 versões no jornal}

%CRISTAL {"arestas":[{"alvo":"hardness_octonica","confianca":1.0,"epistemico":"fato","origem":"paper","rel":"generaliza"}],"confianca":1.0,"contraexemplos":[],"descricao":"Conjectura: a não-representabilidade simbólica de β estende-se a Yang-Mills SU(N) para todo N≥2 — universalidade da dureza, não peculiaridade octônica.","epistemico":"hipotese","exemplos":[],"id":"su_n_extensao","memoria":"semantica","meta":{"autor":"Gentil Lopes","descoberta":false,"dominio":"matematica","fonte":"paper_PQ_I_su_N_clay.pdf, Conj.4.3"},"origem":"paper","palavras_chave":[],"sinonimos":[],"tipo":"conceito","titulo":"Não-Representabilidade SU(N) (conjectura)"}
\section{Não-Representabilidade SU(N) (conjectura)}
Conjectura: a não-representabilidade simbólica de \ensuremath{\beta} estende-se a Yang-Mills SU(N) para todo N\ensuremath{\ge}2 — universalidade da dureza, não peculiaridade octônica.
\prov{relações: generaliza hardness\_octonica}
\prov{proveniência: paper \textperiodcentered{} paper\_PQ\_I\_su\_N\_clay.pdf, Conj.4.3 \textperiodcentered{} hipotese \textperiodcentered{} confiança 1.0 \textperiodcentered{} domínio matematica \textperiodcentered{} história 2 versões no jornal}

%CRISTAL {"arestas":[{"alvo":"reta_de_ouro","confianca":1.0,"epistemico":"fato","origem":"wiki","rel":"usa"},{"alvo":"computacao_fractal","confianca":1.0,"epistemico":"fato","origem":"wiki","rel":"parte_de"}],"confianca":1.0,"contraexemplos":[],"descricao":"O espaço de traços é o SFT Xₖ=x∈0,1ℕ: 1k⁺¹ não ocorre. O autômato de corrida com estados 0..k tem matriz Aₖ; para k=1, A₁=[[1,1],[1,0]] (a matriz de ouro), com Perron λ₁=φ. A entropia topológica é htop=log λₖ.","epistemico":"fato","exemplos":[],"id":"subshift_do_gap","memoria":"semantica","meta":{"autor":"broca-so","dominio":"fractal","fonte":"machine/fractal/fractal.py:conta_palavras; broca_cc.py:A_gap"},"origem":"wiki","palavras_chave":[],"sinonimos":[],"tipo":"conceito","titulo":"O Subshift do Gap X_k (a matriz de ouro)"}
\section{O Subshift do Gap X\_k (a matriz de ouro)}
O espaço de traços é o SFT X\ensuremath{_{k}}=x\ensuremath{\in}0,1\ensuremath{\mathbb{N}}: 1k\ensuremath{^{+1}} não ocorre. O autômato de corrida com estados 0..k tem matriz A\ensuremath{_{k}}; para k=1, A\ensuremath{_{1}}=[[1,1],[1,0]] (a matriz de ouro), com Perron \ensuremath{\lambda}\ensuremath{_{1}}=\ensuremath{\varphi}. A entropia topológica é htop=log \ensuremath{\lambda}\ensuremath{_{k}}.
\prov{relações: usa reta\_de\_ouro; parte\_de computacao\_fractal}
\prov{proveniência: wiki \textperiodcentered{} machine/fractal/fractal.py:conta\_palavras; broca\_cc.py:A\_gap \textperiodcentered{} fato \textperiodcentered{} confiança 1.0 \textperiodcentered{} domínio fractal \textperiodcentered{} história 4 versões no jornal}

%CRISTAL {"arestas":[{"alvo":"dimensao_traco","confianca":1.0,"epistemico":"fato","origem":"paper","rel":"usa"},{"alvo":"vontade_de_poder","confianca":0.5,"epistemico":"fato","origem":"wiki","rel":"relaciona"}],"confianca":1.0,"contraexemplos":[],"descricao":"Espaço de traços Xₖ=x∈0,1ℕ: 1k⁺¹ não ocorre: a máquina é fractal quando seu espaço de traços é um subshift auto-similar codificando a recorrência x²=nx+1.","epistemico":"fato","exemplos":[],"id":"subshift_gap","memoria":"semantica","meta":{"autor":"Gentil Lopes","descoberta":false,"dominio":"matematica","fonte":"computador_fractal.pdf"},"origem":"paper","palavras_chave":[],"sinonimos":[],"tipo":"conceito","titulo":"Subshift do Gap"}
\section{Subshift do Gap}
Espaço de traços X\ensuremath{_{k}}=x\ensuremath{\in}0,1\ensuremath{\mathbb{N}}: 1k\ensuremath{^{+1}} não ocorre: a máquina é fractal quando seu espaço de traços é um subshift auto-similar codificando a recorrência x\ensuremath{^{2}}=nx+1.
\prov{relações: usa dimensao\_traco; relaciona vontade\_de\_poder}
\prov{proveniência: paper \textperiodcentered{} computador\_fractal.pdf \textperiodcentered{} fato \textperiodcentered{} confiança 1.0 \textperiodcentered{} domínio matematica \textperiodcentered{} história 4 versões no jornal}

%CRISTAL {"arestas":[{"alvo":"anel_transcendental_rk","confianca":1.0,"epistemico":"fato","origem":"paper","rel":"usa"},{"alvo":"forma_canonica_lyndon","confianca":1.0,"epistemico":"fato","origem":"paper","rel":"usa"},{"alvo":"cifra","confianca":1.0,"epistemico":"fato","origem":"paper","rel":"referencia"},{"alvo":"word","confianca":0.5,"epistemico":"fato","origem":"wiki","rel":"relaciona"}],"confianca":1.0,"contraexemplos":[],"descricao":"Retornar βG^(k) em forma canônica no anel transcendental Rₖ: problema 'naturalmente duro' sem estrutura algébrica redutível, base de hardness pós-quântica alternativa a reticulados.","epistemico":"fato","exemplos":[],"id":"symbolic_beta","memoria":"semantica","meta":{"autor":"Gentil Lopes","descoberta":true,"dominio":"matematica","fonte":"paper_PQ1_beta_octonic_hardness.pdf"},"origem":"paper","palavras_chave":[],"sinonimos":[],"tipo":"conceito","titulo":"Problema SymbolicBeta"}
\section{Problema SymbolicBeta}
Retornar \ensuremath{\beta}G\^{}(k) em forma canônica no anel transcendental R\ensuremath{_{k}}: problema 'naturalmente duro' sem estrutura algébrica redutível, base de hardness pós-quântica alternativa a reticulados.
\prov{relações: usa anel\_transcendental\_rk; usa forma\_canonica\_lyndon; referencia cifra; relaciona word}
\prov{proveniência: paper \textperiodcentered{} paper\_PQ1\_beta\_octonic\_hardness.pdf \textperiodcentered{} fato \textperiodcentered{} confiança 1.0 \textperiodcentered{} domínio matematica \textperiodcentered{} história 6 versões no jornal}

%CRISTAL {"arestas":[{"alvo":"decoerencia_e_sair_da_borda","confianca":1.0,"epistemico":"fato","origem":"wiki","rel":"parte_de"},{"alvo":"lyapunov_da_orbita","confianca":1.0,"epistemico":"fato","origem":"wiki","rel":"eh_um"},{"alvo":"numero_de_salem","confianca":1.0,"epistemico":"fato","origem":"wiki","rel":"usa"},{"alvo":"numeros_de_pisot","confianca":1.0,"epistemico":"fato","origem":"wiki","rel":"usa"}],"confianca":1.0,"contraexemplos":[],"descricao":"A taxa Γ com que as coerências decaem É o expoente de Lyapunov: Γ = −log ρ (ρ do autovalor subdominante). Na borda (Salem, ρ=1) Γ=0 — não há decoerência. No cristal (Pisot, ρ<1) Γ>0 — decoere ao pointer state. O mesmo número que classifica cristal/borda/caos mede a velocidade da decoerência. Verif.: ρ=1 → Γ=0; ρ=1/φ → Γ=log φ=0.481.","epistemico":"fato","exemplos":[],"id":"taxa_de_decoerencia_e_lyapunov","memoria":"semantica","meta":{"dominio":"reta mineral","fonte":"emaranhamento e decoerência"},"origem":"wiki","palavras_chave":[],"sinonimos":[],"tipo":"conceito","titulo":"A Taxa de Decoerência é o Lyapunov"}
\section{A Taxa de Decoerência é o Lyapunov}
A taxa \ensuremath{\Gamma} com que as coerências decaem É o expoente de Lyapunov: \ensuremath{\Gamma} = \ensuremath{-}log \ensuremath{\rho} (\ensuremath{\rho} do autovalor subdominante). Na borda (Salem, \ensuremath{\rho}=1) \ensuremath{\Gamma}=0 — não há decoerência. No cristal (Pisot, \ensuremath{\rho}<1) \ensuremath{\Gamma}>0 — decoere ao pointer state. O mesmo número que classifica cristal/borda/caos mede a velocidade da decoerência. Verif.: \ensuremath{\rho}=1 \ensuremath{\to} \ensuremath{\Gamma}=0; \ensuremath{\rho}=1/\ensuremath{\varphi} \ensuremath{\to} \ensuremath{\Gamma}=log \ensuremath{\varphi}=0.481.
\prov{relações: parte\_de decoerencia\_e\_sair\_da\_borda; eh\_um lyapunov\_da\_orbita; usa numero\_de\_salem; usa numeros\_de\_pisot}
\prov{proveniência: wiki \textperiodcentered{} emaranhamento e decoerência \textperiodcentered{} fato \textperiodcentered{} confiança 1.0 \textperiodcentered{} domínio reta mineral \textperiodcentered{} história 5 versões no jornal}

%CRISTAL {"arestas":[{"alvo":"piso_parabola_dinamico","confianca":1.0,"epistemico":"fato","origem":"wiki","rel":"especializa"},{"alvo":"purgatorio_nao_seca_lado_a","confianca":1.0,"epistemico":"fato","origem":"wiki","rel":"explica"}],"confianca":1.0,"contraexemplos":[],"descricao":"A parábola a+c²=1 é a mesma em toda camada; o que muda é a TEMPERATURA = o lado-a (calor) que a camada aceita. Camada base: FRIA (c² alto, o cristal — o concentrado nas bacias fortes, q→∞ na linguagem de Rényi). Camadas seguintes: mais QUENTES (aceitam c² menor, o espalhado, a cauda — q baixo) — folhas frescas onde as bacias saturadas voltam ao θ inicial. O purgatório (o lado-a rejeitado na camada fria) cascateia pela ESCADA DE TEMPERATURA: cada camada drena uma banda. É o mesmo θ dinâmico da parábola, deslocado por camada — o que dá a estrutura de Legendre τ(q)↔f(α) do multifractal. [D] a temperatura é o setpoint de a por camada.","epistemico":"inferencia","exemplos":[],"id":"temperatura_camada_parabola","memoria":"semantica","meta":{"autor":"Lopes/broca-so","dominio":"continuo","fonte":"conversa/conhecimento/transicoes.py (camada, theta0); cofre_parabola_hero"},"origem":"wiki","palavras_chave":[],"sinonimos":[],"tipo":"conceito","titulo":"Cada camada é uma banda de TEMPERATURA da parábola (o lado-a governa o nível)"}
\section{Cada camada é uma banda de TEMPERATURA da parábola (o lado-a governa o nível)}
A parábola a+c\ensuremath{^{2}}=1 é a mesma em toda camada; o que muda é a TEMPERATURA = o lado-a (calor) que a camada aceita. Camada base: FRIA (c\ensuremath{^{2}} alto, o cristal — o concentrado nas bacias fortes, q\ensuremath{\to}\ensuremath{\infty} na linguagem de Rényi). Camadas seguintes: mais QUENTES (aceitam c\ensuremath{^{2}} menor, o espalhado, a cauda — q baixo) — folhas frescas onde as bacias saturadas voltam ao \ensuremath{\theta} inicial. O purgatório (o lado-a rejeitado na camada fria) cascateia pela ESCADA DE TEMPERATURA: cada camada drena uma banda. É o mesmo \ensuremath{\theta} dinâmico da parábola, deslocado por camada — o que dá a estrutura de Legendre \ensuremath{\tau}(q)\ensuremath{\leftrightarrow}f(\ensuremath{\alpha}) do multifractal. [D] a temperatura é o setpoint de a por camada.
\prov{relações: especializa piso\_parabola\_dinamico; explica purgatorio\_nao\_seca\_lado\_a}
\prov{proveniência: wiki \textperiodcentered{} conversa/conhecimento/transicoes.py (camada, theta0); cofre\_parabola\_hero \textperiodcentered{} inferencia \textperiodcentered{} confiança 1.0 \textperiodcentered{} domínio continuo \textperiodcentered{} história 18 versões no jornal}

%CRISTAL {"arestas":[{"alvo":"sete_constantes_si","confianca":1.0,"epistemico":"fato","origem":"wiki","rel":"parte_de"},{"alvo":"energia_escura_campo_psi","confianca":1.0,"epistemico":"fato","origem":"wiki","rel":"relaciona"},{"alvo":"horizonte_desitter_gibbons_hawking","confianca":1.0,"epistemico":"fato","origem":"wiki","rel":"explica"},{"alvo":"constantes_horizonte_algebricas","confianca":1.0,"epistemico":"fato","origem":"wiki","rel":"relaciona"}],"confianca":1.0,"contraexemplos":[],"descricao":"A temperatura natural não é o ponto da água — é a do horizonte: Gibbons-Hawking T=ℏH/2πk_B, e Unruh e Hawking, todas com o mesmo 1/2π. Rodado θ→iτ (Wick), o tempo euclidiano é periódico; a regularidade do horizonte (sem defeito cônico, a costura 1≡0) exige período β=2π/κ, e como β=1/T dá T=κ/2π: a temperatura é o inverso do período do tempo imaginário — uma volta euclidiana por 2π. Onde metro e segundo são voltas no tempo real, a temperatura é a volta no imaginário. Horizonte cosmológico: T=ℏH₀/2πk_B≈3×10⁻³⁰ K.","epistemico":"inferencia","exemplos":[],"id":"temperatura_volta_wick","memoria":"semantica","meta":{"autor":"Lopes/broca-so","dominio":"unidades","fonte":"hiper/circular/paper_5_unidades.tex §constantes (kelvin_temperatura.py)"},"origem":"wiki","palavras_chave":[],"sinonimos":[],"tipo":"conceito","titulo":"A temperatura é a volta no tempo imaginário (T=κ/2π)"}
\section{A temperatura é a volta no tempo imaginário (T=\ensuremath{\kappa}/2\ensuremath{\pi})}
A temperatura natural não é o ponto da água — é a do horizonte: Gibbons-Hawking T=\ensuremath{\hbar}H/2\ensuremath{\pi}k\_B, e Unruh e Hawking, todas com o mesmo 1/2\ensuremath{\pi}. Rodado \ensuremath{\theta}\ensuremath{\to}i\ensuremath{\tau} (Wick), o tempo euclidiano é periódico; a regularidade do horizonte (sem defeito cônico, a costura 1\ensuremath{\equiv}0) exige período \ensuremath{\beta}=2\ensuremath{\pi}/\ensuremath{\kappa}, e como \ensuremath{\beta}=1/T dá T=\ensuremath{\kappa}/2\ensuremath{\pi}: a temperatura é o inverso do período do tempo imaginário — uma volta euclidiana por 2\ensuremath{\pi}. Onde metro e segundo são voltas no tempo real, a temperatura é a volta no imaginário. Horizonte cosmológico: T=\ensuremath{\hbar}H\ensuremath{_{0}}/2\ensuremath{\pi}k\_B\ensuremath{\approx}3$\times$10\ensuremath{^{-30}} K.
\prov{relações: parte\_de sete\_constantes\_si; relaciona energia\_escura\_campo\_psi; explica horizonte\_desitter\_gibbons\_hawking; relaciona constantes\_horizonte\_algebricas}
\prov{proveniência: wiki \textperiodcentered{} hiper/circular/paper\_5\_unidades.tex §constantes (kelvin\_temperatura.py) \textperiodcentered{} inferencia \textperiodcentered{} confiança 1.0 \textperiodcentered{} domínio unidades \textperiodcentered{} história 6 versões no jornal}

%CRISTAL {"arestas":[{"alvo":"htlc","confianca":1.0,"epistemico":"fato","origem":"paper","rel":"relaciona"},{"alvo":"goldchain","confianca":1.0,"epistemico":"fato","origem":"paper","rel":"parte_de"}],"confianca":1.0,"contraexemplos":[],"descricao":"No protocolo, o tempo ordena sem tocar a identidade (o RG, não a alma): ancorado em altura de bloco (relógio lógico t=(altura≪K)|seq), não em deadline — ordena, não expira.","epistemico":"fato","exemplos":[],"id":"tempo_ortogonal","memoria":"semantica","meta":{"autor":"Gentil Lopes","descoberta":false,"dominio":"matematica","fonte":"protocolo_assincrono.pdf, §tempo"},"origem":"paper","palavras_chave":[],"sinonimos":[],"tipo":"conceito","titulo":"Tempo Ortogonal"}
\section{Tempo Ortogonal}
No protocolo, o tempo ordena sem tocar a identidade (o RG, não a alma): ancorado em altura de bloco (relógio lógico t=(altura\ensuremath{\ll}K)|seq), não em deadline — ordena, não expira.
\prov{relações: relaciona htlc; parte\_de goldchain}
\prov{proveniência: paper \textperiodcentered{} protocolo\_assincrono.pdf, §tempo \textperiodcentered{} fato \textperiodcentered{} confiança 1.0 \textperiodcentered{} domínio matematica \textperiodcentered{} história 3 versões no jornal}

%CRISTAL {"arestas":[{"alvo":"bac_cab_falha_r7","confianca":1.0,"epistemico":"fato","origem":"paper","rel":"relaciona"},{"alvo":"teoria_gauge_octonica","confianca":1.0,"epistemico":"fato","origem":"paper","rel":"relaciona"}],"confianca":1.0,"contraexemplos":[],"descricao":"4-forma totalmente antissimétrica da estrutura octônica, com ‖φ‖²=168 constante universal (ℝ⁴ e ℝ⁸). A mesma 168 aparece em Yang-Mills octônico — sugere dualidade YM-oct↔NS-oct.","epistemico":"fato","exemplos":[],"id":"tensor_cayley","memoria":"semantica","meta":{"autor":"Gentil Lopes","descoberta":true,"dominio":"matematica","fonte":"paper_AG.pdf, Teo.8"},"origem":"paper","palavras_chave":[],"sinonimos":[],"tipo":"conceito","titulo":"Tensor de Cayley φ (norma 168)"}
\section{Tensor de Cayley \ensuremath{\varphi} (norma 168)}
4-forma totalmente antissimétrica da estrutura octônica, com \ensuremath{\|}\ensuremath{\varphi}\ensuremath{\|}\ensuremath{^{2}}=168 constante universal (\ensuremath{\mathbb{R}}\ensuremath{^{4}} e \ensuremath{\mathbb{R}}\ensuremath{^{8}}). A mesma 168 aparece em Yang-Mills octônico — sugere dualidade YM-oct\ensuremath{\leftrightarrow}NS-oct.
\prov{relações: relaciona bac\_cab\_falha\_r7; relaciona teoria\_gauge\_octonica}
\prov{proveniência: paper \textperiodcentered{} paper\_AG.pdf, Teo.8 \textperiodcentered{} fato \textperiodcentered{} confiança 1.0 \textperiodcentered{} domínio matematica \textperiodcentered{} história 3 versões no jornal}

%CRISTAL {"arestas":[{"alvo":"matematica","confianca":0.95,"epistemico":"fato","origem":"paper","rel":"parte_de"}],"confianca":1.0,"contraexemplos":[],"descricao":"p 91) e a equação (2.2) (p. 85), assim: an(m−j) = ∆j anm = j X k = 0 (−1)k \u0012 j k \u0013 a(n−k+j)m (2.7) Substituindo nesta equação n = 1, resulta: a1(m−j) = j X k = 0 (−1)k \u0012 j k \u0013 a(1−k+j)m ■ Prova: (Equação (1.17), p. 44) Tomemos j = m na equação","epistemico":"fato","exemplos":[],"id":"teorema_10_p_91_e_a_equacao_22_p_85_assim_anmj_j_an","memoria":"semantica","meta":{"arquivo":"gentil/Novas_Sequencias_Aritmeticas_e_Geometric.pdf","ciencia":"teoria_dos_numeros","dominio":"matematica","nivel":4,"numero":"10","tipo_no":"paper","tipo_teorema":"teorema","verificacao":"conceito novo"},"origem":"gentil/Novas_Sequencias_Aritmeticas_e_Geometric.pdf","palavras_chave":[],"sinonimos":[],"tipo":"conceito","titulo":"Teorema 10: p 91) e a equação (2.2) (p. 85), assim: an(m−j) = ∆j an…"}
\section{Teorema 10: p 91) e a equação (2.2) (p. 85), assim: an(m\ensuremath{-}j) = \ensuremath{\Delta}j an…}
p 91) e a equação (2.2) (p. 85), assim: an(m\ensuremath{-}j) = \ensuremath{\Delta}j anm = j X k = 0 (\ensuremath{-}1)k  j k  a(n\ensuremath{-}k+j)m (2.7) Substituindo nesta equação n = 1, resulta: a1(m\ensuremath{-}j) = j X k = 0 (\ensuremath{-}1)k  j k  a(1\ensuremath{-}k+j)m \rule{1ex}{1ex} Prova: (Equação (1.17), p. 44) Tomemos j = m na equação
\prov{relações: parte\_de matematica}
\prov{proveniência: gentil/Novas\_Sequencias\_Aritmeticas\_e\_Geometric.pdf \textperiodcentered{} fato \textperiodcentered{} confiança 1.0 \textperiodcentered{} domínio matematica \textperiodcentered{} história 2 versões no jornal}

%CRISTAL {"arestas":[{"alvo":"matematica","confianca":0.95,"epistemico":"fato","origem":"paper","rel":"parte_de"}],"confianca":1.0,"contraexemplos":[],"descricao":"Teorema do gene - I) Seja m um natural arbitrariamente fixado e seja j um natural tal que 0 ≤j ≤m. Nestas condições a seguinte identidade se verifica: ∆j anm = an(m−j) (2.6) Prova: Indução sobre j. Para j = 0 temos: ∆0 anm = anm = an(m−0) Suponhamos a proposição válida para j = p (0 ≤p < m): (H.I) ∆p anm = an(m−p) Provemos que vale para j = p + 1: (T.I) ∆p+1 anm = an(m−(p+1)) Então, ∆p+1 anm = ∆p a(n+1)m −∆p anm = a(n+1)(m−p) −an(m−p) = an(m−(p+1)) ■ Na última igualdade acima fizemos uso da terc","epistemico":"fato","exemplos":[],"id":"teorema_10_teorema_do_gene_-_i_seja_m_um_natural_arbitrariamente","memoria":"semantica","meta":{"arquivo":"gentil/Novas_Sequencias_Aritmeticas_e_Geometric.pdf","ciencia":"teoria_dos_numeros","dominio":"matematica","nivel":4,"numero":"10","tipo_no":"paper","tipo_teorema":"teorema","verificacao":"conceito novo"},"origem":"gentil/Novas_Sequencias_Aritmeticas_e_Geometric.pdf","palavras_chave":[],"sinonimos":[],"tipo":"conceito","titulo":"Teorema 10: Teorema do gene - I) Seja m um natural arbitrariamente …"}
\section{Teorema 10: Teorema do gene - I) Seja m um natural arbitrariamente …}
Teorema do gene - I) Seja m um natural arbitrariamente fixado e seja j um natural tal que 0 \ensuremath{\le}j \ensuremath{\le}m. Nestas condições a seguinte identidade se verifica: \ensuremath{\Delta}j anm = an(m\ensuremath{-}j) (2.6) Prova: Indução sobre j. Para j = 0 temos: \ensuremath{\Delta}0 anm = anm = an(m\ensuremath{-}0) Suponhamos a proposição válida para j = p (0 \ensuremath{\le}p < m): (H.I) \ensuremath{\Delta}p anm = an(m\ensuremath{-}p) Provemos que vale para j = p + 1: (T.I) \ensuremath{\Delta}p+1 anm = an(m\ensuremath{-}(p+1)) Então, \ensuremath{\Delta}p+1 anm = \ensuremath{\Delta}p a(n+1)m \ensuremath{-}\ensuremath{\Delta}p anm = a(n+1)(m\ensuremath{-}p) \ensuremath{-}an(m\ensuremath{-}p) = an(m\ensuremath{-}(p+1)) \rule{1ex}{1ex} Na última igualdade acima fizemos uso da terc
\prov{relações: parte\_de matematica}
\prov{proveniência: gentil/Novas\_Sequencias\_Aritmeticas\_e\_Geometric.pdf \textperiodcentered{} fato \textperiodcentered{} confiança 1.0 \textperiodcentered{} domínio matematica \textperiodcentered{} história 2 versões no jornal}

%CRISTAL {"arestas":[{"alvo":"matematica","confianca":0.95,"epistemico":"fato","origem":"paper","rel":"parte_de"}],"confianca":1.0,"contraexemplos":[],"descricao":"p 99), o qual repetimos: Teorema 16 (Teorema do gene - II). Seja m um natural arbitrariamente fixado e j ∈N ∪ 0. Nestas condições a seguinte identidade se verifica: ∇j anm = an(m+j) Prova: Considere a equação ∆j anm = an(m−j). Então, ∆j anm = an(m−j) ⇒ ∆j an(m+j) = anm Aplicando o operador ∇j a ambos os membros desta última equação, temos ∇j ∆j an(m+j) = ∇j anm ⇒ ∇j anm = an(m+j) ■ O","epistemico":"fato","exemplos":[],"id":"teorema_12_p_99_o_qual_repetimos_teorema_16_teorema_do_gene_-","memoria":"semantica","meta":{"arquivo":"gentil/Novas_Sequencias_Aritmeticas_e_Geometric.pdf","ciencia":"teoria_dos_numeros","dominio":"matematica","nivel":4,"numero":"12","tipo_no":"paper","tipo_teorema":"teorema","verificacao":"conceito novo"},"origem":"gentil/Novas_Sequencias_Aritmeticas_e_Geometric.pdf","palavras_chave":[],"sinonimos":[],"tipo":"conceito","titulo":"Teorema 12: p 99), o qual repetimos: Teorema 16 (Teorema do gene - …"}
\section{Teorema 12: p 99), o qual repetimos: Teorema 16 (Teorema do gene - …}
p 99), o qual repetimos: Teorema 16 (Teorema do gene - II). Seja m um natural arbitrariamente fixado e j \ensuremath{\in}N \ensuremath{\cup} 0. Nestas condições a seguinte identidade se verifica: \ensuremath{\nabla}j anm = an(m+j) Prova: Considere a equação \ensuremath{\Delta}j anm = an(m\ensuremath{-}j). Então, \ensuremath{\Delta}j anm = an(m\ensuremath{-}j) \ensuremath{\Rightarrow} \ensuremath{\Delta}j an(m+j) = anm Aplicando o operador \ensuremath{\nabla}j a ambos os membros desta última equação, temos \ensuremath{\nabla}j \ensuremath{\Delta}j an(m+j) = \ensuremath{\nabla}j anm \ensuremath{\Rightarrow} \ensuremath{\nabla}j anm = an(m+j) \rule{1ex}{1ex} O
\prov{relações: parte\_de matematica}
\prov{proveniência: gentil/Novas\_Sequencias\_Aritmeticas\_e\_Geometric.pdf \textperiodcentered{} fato \textperiodcentered{} confiança 1.0 \textperiodcentered{} domínio matematica \textperiodcentered{} história 3 versões no jornal}

%CRISTAL {"arestas":[{"alvo":"matematica","confianca":0.95,"epistemico":"fato","origem":"paper","rel":"parte_de"}],"confianca":1.0,"contraexemplos":[],"descricao":"Fórmula do termo geral de uma P A.m ). Em uma P.A.m o n −ésimo termo vale anm = m X j = 0 \u0012 n −1 j \u0013 a1(m−j) (1.10) Nota: Esta fórmula do termo geral de uma P.A.m é inédita, isto é, não consta na literatura −Assim como muitas outras que ainda aparecerão oportunamente. Isto tudo é uma consequência de nossa","epistemico":"fato","exemplos":[],"id":"teorema_1_formula_do_termo_geral_de_uma_p_am_em_uma_pam_o_n","memoria":"semantica","meta":{"arquivo":"gentil/Novas_Sequencias_Aritmeticas_e_Geometric.pdf","ciencia":"teoria_dos_numeros","dominio":"matematica","nivel":4,"numero":"1","tipo_no":"paper","tipo_teorema":"teorema","verificacao":"conceito novo"},"origem":"gentil/Novas_Sequencias_Aritmeticas_e_Geometric.pdf","palavras_chave":[],"sinonimos":[],"tipo":"conceito","titulo":"Teorema 1: Fórmula do termo geral de uma P A.m ). Em uma P.A.m o n…"}
\section{Teorema 1: Fórmula do termo geral de uma P A.m ). Em uma P.A.m o n…}
Fórmula do termo geral de uma P A.m ). Em uma P.A.m o n \ensuremath{-}ésimo termo vale anm = m X j = 0  n \ensuremath{-}1 j  a1(m\ensuremath{-}j) (1.10) Nota: Esta fórmula do termo geral de uma P.A.m é inédita, isto é, não consta na literatura \ensuremath{-}Assim como muitas outras que ainda aparecerão oportunamente. Isto tudo é uma consequência de nossa
\prov{relações: parte\_de matematica}
\prov{proveniência: gentil/Novas\_Sequencias\_Aritmeticas\_e\_Geometric.pdf \textperiodcentered{} fato \textperiodcentered{} confiança 1.0 \textperiodcentered{} domínio matematica \textperiodcentered{} história 2 versões no jornal}

%CRISTAL {"arestas":[{"alvo":"matematica","confianca":0.95,"epistemico":"fato","origem":"paper","rel":"parte_de"}],"confianca":1.0,"contraexemplos":[],"descricao":"Em uma P G.m é válida a seguinte identidade: a1(m−j) = jY k = 0 a (−1)k ( j k) (1−k+j)m (3.5) Prova: Princ´ıpio da Dualidade. Esta é a equação dual de (1.16), p. 40. ■ −P.G.m em função dos seus próprios termos na HP Prime Na tela a seguir a1(m−j) = jY k = 0 a (−1)k ( j k) (1−k+j)m temos o programa que implementa a fórmula da direita. Como ilustração, considerando o exemplo da página 136, entrando com os quatro primeiros termos em um vetor, 1 1 1 −1 1 1 −1 1 −1 −1 a10 → a11 → a12 → a13 → o progra","epistemico":"fato","exemplos":[],"id":"teorema_21_em_uma_p_gm_e_valida_a_seguinte_identidade_a1mj","memoria":"semantica","meta":{"arquivo":"gentil/Novas_Sequencias_Aritmeticas_e_Geometric.pdf","ciencia":"teoria_dos_numeros","dominio":"matematica","nivel":4,"numero":"21","tipo_no":"paper","tipo_teorema":"teorema","verificacao":"conceito novo"},"origem":"gentil/Novas_Sequencias_Aritmeticas_e_Geometric.pdf","palavras_chave":[],"sinonimos":[],"tipo":"conceito","titulo":"Teorema 21: Em uma P G.m é válida a seguinte identidade: a1(m−j) = …"}
\section{Teorema 21: Em uma P G.m é válida a seguinte identidade: a1(m\ensuremath{-}j) = …}
Em uma P G.m é válida a seguinte identidade: a1(m\ensuremath{-}j) = jY k = 0 a (\ensuremath{-}1)k ( j k) (1\ensuremath{-}k+j)m (3.5) Prova: Princ'\ensuremath{\imath}pio da Dualidade. Esta é a equação dual de (1.16), p. 40. \rule{1ex}{1ex} \ensuremath{-}P.G.m em função dos seus próprios termos na HP Prime Na tela a seguir a1(m\ensuremath{-}j) = jY k = 0 a (\ensuremath{-}1)k ( j k) (1\ensuremath{-}k+j)m temos o programa que implementa a fórmula da direita. Como ilustração, considerando o exemplo da página 136, entrando com os quatro primeiros termos em um vetor, 1 1 1 \ensuremath{-}1 1 1 \ensuremath{-}1 1 \ensuremath{-}1 \ensuremath{-}1 a10 \ensuremath{\to} a11 \ensuremath{\to} a12 \ensuremath{\to} a13 \ensuremath{\to} o progra
\prov{relações: parte\_de matematica}
\prov{proveniência: gentil/Novas\_Sequencias\_Aritmeticas\_e\_Geometric.pdf \textperiodcentered{} fato \textperiodcentered{} confiança 1.0 \textperiodcentered{} domínio matematica \textperiodcentered{} história 2 versões no jornal}

%CRISTAL {"arestas":[{"alvo":"matematica","confianca":0.95,"epistemico":"fato","origem":"paper","rel":"parte_de"}],"confianca":1.0,"contraexemplos":[],"descricao":"Produto dos termos de uma P G.m ). Em uma P.G.m o produto Pnm dos n termos iniciais vale Pnm = m Y j = 0 a( n j+1) 1(m−j) (3.7) Prova: Princ´ıpio da Dualidade. Esta equação é a dual da (1.18), p. 47. ■ Exemplo: Encontre o produto dos n termos iniciais da seguinte P.G.2 : 1 1 −1 −1 1 1 −1 −1 . . . Solução: Substituindo m = 2 na fórmula (3.7), obtemos: Pn2 = 2 Y j = 0 a( n j+1) 1(2−j) Simplificando, obtemos: Pn2 = a( n 1 ) 12 · a( n 2 ) 11 · a( n 3 ) 10 139\nDo algoritmo retiramos os coeficientes,","epistemico":"fato","exemplos":[],"id":"teorema_23_produto_dos_termos_de_uma_p_gm_em_uma_pgm_o_produ","memoria":"semantica","meta":{"arquivo":"gentil/Novas_Sequencias_Aritmeticas_e_Geometric.pdf","ciencia":"teoria_dos_numeros","dominio":"matematica","nivel":4,"numero":"23","tipo_no":"paper","tipo_teorema":"teorema","verificacao":"conceito novo"},"origem":"gentil/Novas_Sequencias_Aritmeticas_e_Geometric.pdf","palavras_chave":[],"sinonimos":[],"tipo":"conceito","titulo":"Teorema 23: Produto dos termos de uma P G.m ). Em uma P.G.m o produ…"}
\section{Teorema 23: Produto dos termos de uma P G.m ). Em uma P.G.m o produ…}
Produto dos termos de uma P G.m ). Em uma P.G.m o produto Pnm dos n termos iniciais vale Pnm = m Y j = 0 a( n j+1) 1(m\ensuremath{-}j) (3.7) Prova: Princ'\ensuremath{\imath}pio da Dualidade. Esta equação é a dual da (1.18), p. 47. \rule{1ex}{1ex} Exemplo: Encontre o produto dos n termos iniciais da seguinte P.G.2 : 1 1 \ensuremath{-}1 \ensuremath{-}1 1 1 \ensuremath{-}1 \ensuremath{-}1 . . . Solução: Substituindo m = 2 na fórmula (3.7), obtemos: Pn2 = 2 Y j = 0 a( n j+1) 1(2\ensuremath{-}j) Simplificando, obtemos: Pn2 = a( n 1 ) 12 \textperiodcentered  a( n 2 ) 11 \textperiodcentered  a( n 3 ) 10 139
Do algoritmo retiramos os coeficientes,
\prov{relações: parte\_de matematica}
\prov{proveniência: gentil/Novas\_Sequencias\_Aritmeticas\_e\_Geometric.pdf \textperiodcentered{} fato \textperiodcentered{} confiança 1.0 \textperiodcentered{} domínio matematica \textperiodcentered{} história 2 versões no jornal}

%CRISTAL {"arestas":[{"alvo":"matematica","confianca":0.95,"epistemico":"fato","origem":"paper","rel":"parte_de"}],"confianca":1.0,"contraexemplos":[],"descricao":"As seguintes identidades são válidas: (i) ∆m f × g (n) = ∆m f(n) × ∆m g(n) (ii) ∆m f/g (n) = ∆m f(n) / ∆m g(n) Prova: Apêndice. (p. 177) ■ Corolário 4. ´E válida a seguinte identidade: ∆m f(n) = ∆m−1 ∆f(n) Prova: Princ´ıpio da Dualidade. ■ O seguinte","epistemico":"fato","exemplos":[],"id":"teorema_25_as_seguintes_identidades_sao_validas_i_m_f_g_n","memoria":"semantica","meta":{"arquivo":"gentil/Novas_Sequencias_Aritmeticas_e_Geometric.pdf","ciencia":"teoria_dos_numeros","dominio":"matematica","nivel":4,"numero":"25","tipo_no":"paper","tipo_teorema":"teorema","verificacao":"conceito novo"},"origem":"gentil/Novas_Sequencias_Aritmeticas_e_Geometric.pdf","palavras_chave":[],"sinonimos":[],"tipo":"conceito","titulo":"Teorema 25: As seguintes identidades são válidas: (i) ∆m f × g (n) …"}
\section{Teorema 25: As seguintes identidades são válidas: (i) \ensuremath{\Delta}m f $\times$ g (n) …}
As seguintes identidades são válidas: (i) \ensuremath{\Delta}m f $\times$ g (n) = \ensuremath{\Delta}m f(n) $\times$ \ensuremath{\Delta}m g(n) (ii) \ensuremath{\Delta}m f/g (n) = \ensuremath{\Delta}m f(n) / \ensuremath{\Delta}m g(n) Prova: Apêndice. (p. 177) \rule{1ex}{1ex} Corolário 4. 'E válida a seguinte identidade: \ensuremath{\Delta}m f(n) = \ensuremath{\Delta}m\ensuremath{-}1 \ensuremath{\Delta}f(n) Prova: Princ'\ensuremath{\imath}pio da Dualidade. \rule{1ex}{1ex} O seguinte
\prov{relações: parte\_de matematica}
\prov{proveniência: gentil/Novas\_Sequencias\_Aritmeticas\_e\_Geometric.pdf \textperiodcentered{} fato \textperiodcentered{} confiança 1.0 \textperiodcentered{} domínio matematica \textperiodcentered{} história 2 versões no jornal}

%CRISTAL {"arestas":[{"alvo":"matematica","confianca":0.95,"epistemico":"fato","origem":"paper","rel":"parte_de"}],"confianca":1.0,"contraexemplos":[],"descricao":"Teorema do gene - III) Seja m um natural arbitrariamente fixado e seja j um natural tal que 0 ≤j ≤m. Nestas condições a seguinte identidade se verifica: ∆j anm = an(m−j) (3.11) Prova: Princ´ıpio da Dualidade. ■ Este","epistemico":"fato","exemplos":[],"id":"teorema_26_teorema_do_gene_-_iii_seja_m_um_natural_arbitrariament","memoria":"semantica","meta":{"arquivo":"gentil/Novas_Sequencias_Aritmeticas_e_Geometric.pdf","ciencia":"teoria_dos_numeros","dominio":"matematica","nivel":4,"numero":"26","tipo_no":"paper","tipo_teorema":"teorema","verificacao":"conceito novo"},"origem":"gentil/Novas_Sequencias_Aritmeticas_e_Geometric.pdf","palavras_chave":[],"sinonimos":[],"tipo":"conceito","titulo":"Teorema 26: Teorema do gene - III) Seja m um natural arbitrariament…"}
\section{Teorema 26: Teorema do gene - III) Seja m um natural arbitrariament…}
Teorema do gene - III) Seja m um natural arbitrariamente fixado e seja j um natural tal que 0 \ensuremath{\le}j \ensuremath{\le}m. Nestas condições a seguinte identidade se verifica: \ensuremath{\Delta}j anm = an(m\ensuremath{-}j) (3.11) Prova: Princ'\ensuremath{\imath}pio da Dualidade. \rule{1ex}{1ex} Este
\prov{relações: parte\_de matematica}
\prov{proveniência: gentil/Novas\_Sequencias\_Aritmeticas\_e\_Geometric.pdf \textperiodcentered{} fato \textperiodcentered{} confiança 1.0 \textperiodcentered{} domínio matematica \textperiodcentered{} história 2 versões no jornal}

%CRISTAL {"arestas":[{"alvo":"matematica","confianca":0.95,"epistemico":"fato","origem":"paper","rel":"parte_de"}],"confianca":1.0,"contraexemplos":[],"descricao":"As seguintes identidades são válidas: (i) ∆ m f × g (n) = ∆ m f(n) × ∆ m g(n) (ii) ∆ m f/g (n) = ∆ m f(n) / ∆ m g(n) Prova: Apêndice. (p. 178) ■","epistemico":"fato","exemplos":[],"id":"teorema_28_as_seguintes_identidades_sao_validas_i_m_f_g_n","memoria":"semantica","meta":{"arquivo":"gentil/Novas_Sequencias_Aritmeticas_e_Geometric.pdf","ciencia":"teoria_dos_numeros","dominio":"matematica","nivel":4,"numero":"28","tipo_no":"paper","tipo_teorema":"teorema","verificacao":"conceito novo"},"origem":"gentil/Novas_Sequencias_Aritmeticas_e_Geometric.pdf","palavras_chave":[],"sinonimos":[],"tipo":"conceito","titulo":"Teorema 28: As seguintes identidades são válidas: (i) ∆ m f × g (n)…"}
\section{Teorema 28: As seguintes identidades são válidas: (i) \ensuremath{\Delta} m f $\times$ g (n)…}
As seguintes identidades são válidas: (i) \ensuremath{\Delta} m f $\times$ g (n) = \ensuremath{\Delta} m f(n) $\times$ \ensuremath{\Delta} m g(n) (ii) \ensuremath{\Delta} m f/g (n) = \ensuremath{\Delta} m f(n) / \ensuremath{\Delta} m g(n) Prova: Apêndice. (p. 178) \rule{1ex}{1ex}
\prov{relações: parte\_de matematica}
\prov{proveniência: gentil/Novas\_Sequencias\_Aritmeticas\_e\_Geometric.pdf \textperiodcentered{} fato \textperiodcentered{} confiança 1.0 \textperiodcentered{} domínio matematica \textperiodcentered{} história 2 versões no jornal}

%CRISTAL {"arestas":[{"alvo":"matematica","confianca":0.95,"epistemico":"fato","origem":"paper","rel":"parte_de"}],"confianca":1.0,"contraexemplos":[],"descricao":"Teorema Q - P Fundamental - I) Seja m um natural arbitrariamente fixado e j um natural tal que 0 ≤j ≤m. Nestas condições a seguinte identidade se verifica: ∆ j ∆m f(n) = ∆m−j f(n) (3.13) 149\nProva: Princ´ıpio da Dualidade. (p. 102) ■","epistemico":"fato","exemplos":[],"id":"teorema_29_teorema_q_-_p_fundamental_-_i_seja_m_um_natural_arbitr","memoria":"semantica","meta":{"arquivo":"gentil/Novas_Sequencias_Aritmeticas_e_Geometric.pdf","ciencia":"teoria_dos_numeros","dominio":"matematica","nivel":4,"numero":"29","tipo_no":"paper","tipo_teorema":"teorema","verificacao":"conceito novo"},"origem":"gentil/Novas_Sequencias_Aritmeticas_e_Geometric.pdf","palavras_chave":[],"sinonimos":[],"tipo":"conceito","titulo":"Teorema 29: Teorema Q - P Fundamental - I) Seja m um natural arbitr…"}
\section{Teorema 29: Teorema Q - P Fundamental - I) Seja m um natural arbitr…}
Teorema Q - P Fundamental - I) Seja m um natural arbitrariamente fixado e j um natural tal que 0 \ensuremath{\le}j \ensuremath{\le}m. Nestas condições a seguinte identidade se verifica: \ensuremath{\Delta} j \ensuremath{\Delta}m f(n) = \ensuremath{\Delta}m\ensuremath{-}j f(n) (3.13) 149
Prova: Princ'\ensuremath{\imath}pio da Dualidade. (p. 102) \rule{1ex}{1ex}
\prov{relações: parte\_de matematica}
\prov{proveniência: gentil/Novas\_Sequencias\_Aritmeticas\_e\_Geometric.pdf \textperiodcentered{} fato \textperiodcentered{} confiança 1.0 \textperiodcentered{} domínio matematica \textperiodcentered{} história 2 versões no jornal}

%CRISTAL {"arestas":[{"alvo":"matematica","confianca":0.95,"epistemico":"fato","origem":"paper","rel":"parte_de"}],"confianca":1.0,"contraexemplos":[],"descricao":"Seja j um natural arbitrariamente fixado Para n ≥j vale a seguinte identidade: n X i = j \u0012 i j \u0013 = \u0012 n + 1 j + 1 \u0013 (1.11) Prova: Apêndice, página 71. ■ Exemplos: (a) Tomando j = 1 na equação (1.11), obtemos: n X i = 1 \u0012 i 1 \u0013 = \u0012 n + 1 1 + 1 \u0013 Isto é, \u0012 1 1 \u0013 + \u0012 2 1 \u0013 + · · · + \u0012 n 1 \u0013 = \u0012 n + 1 2 \u0013 Ou ainda, 1 + 2 + · · · + n = n(n + 1) 2 (b) Tomando j = 2 na equação (1.11), obtemos: n X i = 2 \u0012 i 2 \u0013 = \u0012 n + 1 2 + 1 \u0013 Isto é, \u0012 1 2 \u0013 + \u0012 2 2 \u0013 + \u0012 3 2 \u0013 + · · · + \u0012 n 2 \u0013 = \u0012 n + 1 3 \u0013 Ou aind","epistemico":"fato","exemplos":[],"id":"teorema_2_seja_j_um_natural_arbitrariamente_fixado_para_n_j_vale","memoria":"semantica","meta":{"arquivo":"gentil/Novas_Sequencias_Aritmeticas_e_Geometric.pdf","ciencia":"teoria_dos_numeros","dominio":"matematica","nivel":4,"numero":"2","tipo_no":"paper","tipo_teorema":"teorema","verificacao":"conceito novo"},"origem":"gentil/Novas_Sequencias_Aritmeticas_e_Geometric.pdf","palavras_chave":[],"sinonimos":[],"tipo":"conceito","titulo":"Teorema 2: Seja j um natural arbitrariamente fixado Para n ≥j vale…"}
\section{Teorema 2: Seja j um natural arbitrariamente fixado Para n \ensuremath{\ge}j vale…}
Seja j um natural arbitrariamente fixado Para n \ensuremath{\ge}j vale a seguinte identidade: n X i = j  i j  =  n + 1 j + 1  (1.11) Prova: Apêndice, página 71. \rule{1ex}{1ex} Exemplos: (a) Tomando j = 1 na equação (1.11), obtemos: n X i = 1  i 1  =  n + 1 1 + 1  Isto é,  1 1  +  2 1  + \textperiodcentered  \textperiodcentered  \textperiodcentered  +  n 1  =  n + 1 2  Ou ainda, 1 + 2 + \textperiodcentered  \textperiodcentered  \textperiodcentered  + n = n(n + 1) 2 (b) Tomando j = 2 na equação (1.11), obtemos: n X i = 2  i 2  =  n + 1 2 + 1  Isto é,  1 2  +  2 2  +  3 2  + \textperiodcentered  \textperiodcentered  \textperiodcentered  +  n 2  =  n + 1 3  Ou aind
\prov{relações: parte\_de matematica}
\prov{proveniência: gentil/Novas\_Sequencias\_Aritmeticas\_e\_Geometric.pdf \textperiodcentered{} fato \textperiodcentered{} confiança 1.0 \textperiodcentered{} domínio matematica \textperiodcentered{} história 2 versões no jornal}

%CRISTAL {"arestas":[{"alvo":"matematica","confianca":0.95,"epistemico":"fato","origem":"paper","rel":"parte_de"}],"confianca":1.0,"contraexemplos":[],"descricao":"Teorema Q - P Fundamental - II) Seja m um natural arbitraria- mente fixado e j um natural tal que 0 ≤j ≤m. Nestas condições a seguinte identidade se verifica: ∆j ∆ m f(n) = ∆m−j f(n) (3.14) Prova: Princ´ıpio da Dualidade. (p. 105) ■ O","epistemico":"fato","exemplos":[],"id":"teorema_30_teorema_q_-_p_fundamental_-_ii_seja_m_um_natural_arbit","memoria":"semantica","meta":{"arquivo":"gentil/Novas_Sequencias_Aritmeticas_e_Geometric.pdf","ciencia":"teoria_dos_numeros","dominio":"matematica","nivel":4,"numero":"30","tipo_no":"paper","tipo_teorema":"teorema","verificacao":"conceito novo"},"origem":"gentil/Novas_Sequencias_Aritmeticas_e_Geometric.pdf","palavras_chave":[],"sinonimos":[],"tipo":"conceito","titulo":"Teorema 30: Teorema Q - P Fundamental - II) Seja m um natural arbit…"}
\section{Teorema 30: Teorema Q - P Fundamental - II) Seja m um natural arbit…}
Teorema Q - P Fundamental - II) Seja m um natural arbitraria- mente fixado e j um natural tal que 0 \ensuremath{\le}j \ensuremath{\le}m. Nestas condições a seguinte identidade se verifica: \ensuremath{\Delta}j \ensuremath{\Delta} m f(n) = \ensuremath{\Delta}m\ensuremath{-}j f(n) (3.14) Prova: Princ'\ensuremath{\imath}pio da Dualidade. (p. 105) \rule{1ex}{1ex} O
\prov{relações: parte\_de matematica}
\prov{proveniência: gentil/Novas\_Sequencias\_Aritmeticas\_e\_Geometric.pdf \textperiodcentered{} fato \textperiodcentered{} confiança 1.0 \textperiodcentered{} domínio matematica \textperiodcentered{} história 2 versões no jornal}

%CRISTAL {"arestas":[{"alvo":"matematica","confianca":0.95,"epistemico":"fato","origem":"paper","rel":"parte_de"}],"confianca":1.0,"contraexemplos":[],"descricao":"Teorema do gene - IV) Seja m um natural arbitrariamente fixado e j ∈N ∪ 0. Nestas condições a seguinte identidade se verifica: ∆ j anm = an(m+j) (3.15) Prova: Princ´ıpio da Dualidade. (p. 99) ■ Com o aux´ılio dos","epistemico":"fato","exemplos":[],"id":"teorema_31_teorema_do_gene_-_iv_seja_m_um_natural_arbitrariamente","memoria":"semantica","meta":{"arquivo":"gentil/Novas_Sequencias_Aritmeticas_e_Geometric.pdf","ciencia":"teoria_dos_numeros","dominio":"matematica","nivel":4,"numero":"31","tipo_no":"paper","tipo_teorema":"teorema","verificacao":"conceito novo"},"origem":"gentil/Novas_Sequencias_Aritmeticas_e_Geometric.pdf","palavras_chave":[],"sinonimos":[],"tipo":"conceito","titulo":"Teorema 31: Teorema do gene - IV) Seja m um natural arbitrariamente…"}
\section{Teorema 31: Teorema do gene - IV) Seja m um natural arbitrariamente…}
Teorema do gene - IV) Seja m um natural arbitrariamente fixado e j \ensuremath{\in}N \ensuremath{\cup} 0. Nestas condições a seguinte identidade se verifica: \ensuremath{\Delta} j anm = an(m+j) (3.15) Prova: Princ'\ensuremath{\imath}pio da Dualidade. (p. 99) \rule{1ex}{1ex} Com o aux'\ensuremath{\imath}lio dos
\prov{relações: parte\_de matematica}
\prov{proveniência: gentil/Novas\_Sequencias\_Aritmeticas\_e\_Geometric.pdf \textperiodcentered{} fato \textperiodcentered{} confiança 1.0 \textperiodcentered{} domínio matematica \textperiodcentered{} história 3 versões no jornal}

%CRISTAL {"arestas":[{"alvo":"matematica","confianca":0.95,"epistemico":"fato","origem":"paper","rel":"parte_de"},{"alvo":"vida-morte","confianca":0.5,"epistemico":"fato","origem":"wiki","rel":"relaciona"},{"alvo":"virtualizacao","confianca":0.5,"epistemico":"fato","origem":"wiki","rel":"relaciona"},{"alvo":"word","confianca":0.5,"epistemico":"fato","origem":"wiki","rel":"relaciona"}],"confianca":1.0,"contraexemplos":[],"descricao":", p 150) ainda resta uma pergunta: quais tipos de sequências conseguimos unificar sob as P.G.m ?","epistemico":"fato","exemplos":[],"id":"teorema_32_p_150_ainda_resta_uma_pergunta_quais_tipos_de_seque","memoria":"semantica","meta":{"arquivo":"gentil/Novas_Sequencias_Aritmeticas_e_Geometric.pdf","ciencia":"teoria_dos_numeros","dominio":"matematica","nivel":4,"numero":"32","tipo_no":"paper","tipo_teorema":"teorema","verificacao":"conceito novo"},"origem":"gentil/Novas_Sequencias_Aritmeticas_e_Geometric.pdf","palavras_chave":[],"sinonimos":[],"tipo":"conceito","titulo":"Teorema 32: , p 150) ainda resta uma pergunta: quais tipos de sequê…"}
\section{Teorema 32: , p 150) ainda resta uma pergunta: quais tipos de sequê…}
, p 150) ainda resta uma pergunta: quais tipos de sequências conseguimos unificar sob as P.G.m ?
\prov{relações: parte\_de matematica; relaciona vida-morte; relaciona virtualizacao; relaciona word}
\prov{proveniência: gentil/Novas\_Sequencias\_Aritmeticas\_e\_Geometric.pdf \textperiodcentered{} fato \textperiodcentered{} confiança 1.0 \textperiodcentered{} domínio matematica \textperiodcentered{} história 5 versões no jornal}

%CRISTAL {"arestas":[{"alvo":"matematica","confianca":0.95,"epistemico":"fato","origem":"paper","rel":"parte_de"}],"confianca":1.0,"contraexemplos":[],"descricao":"Fórmula do termo geral de uma PA-2D) Na PA-2D em que o pri- meiro termo é a11, a razão das linhas é r1 e a razão das colunas é r2 o (i, j)-termo é: aij = a11 + (j −1) r1 + (i −1) r2 (5.1) 232\nProva: Indução sobre i (j fixo). Para i = 1 precisamos mostrar que, a1j = a11 + (j −1) r1 De fato, para i = 1 a","epistemico":"fato","exemplos":[],"id":"teorema_38_formula_do_termo_geral_de_uma_pa-2d_na_pa-2d_em_que_o","memoria":"semantica","meta":{"arquivo":"gentil/Novas_Sequencias_Aritmeticas_e_Geometric.pdf","ciencia":"teoria_dos_numeros","dominio":"matematica","nivel":4,"numero":"38","tipo_no":"paper","tipo_teorema":"teorema","verificacao":"conceito novo"},"origem":"gentil/Novas_Sequencias_Aritmeticas_e_Geometric.pdf","palavras_chave":[],"sinonimos":[],"tipo":"conceito","titulo":"Teorema 38: Fórmula do termo geral de uma PA-2D) Na PA-2D em que o …"}
\section{Teorema 38: Fórmula do termo geral de uma PA-2D) Na PA-2D em que o …}
Fórmula do termo geral de uma PA-2D) Na PA-2D em que o pri- meiro termo é a11, a razão das linhas é r1 e a razão das colunas é r2 o (i, j)-termo é: aij = a11 + (j \ensuremath{-}1) r1 + (i \ensuremath{-}1) r2 (5.1) 232
Prova: Indução sobre i (j fixo). Para i = 1 precisamos mostrar que, a1j = a11 + (j \ensuremath{-}1) r1 De fato, para i = 1 a
\prov{relações: parte\_de matematica}
\prov{proveniência: gentil/Novas\_Sequencias\_Aritmeticas\_e\_Geometric.pdf \textperiodcentered{} fato \textperiodcentered{} confiança 1.0 \textperiodcentered{} domínio matematica \textperiodcentered{} história 2 versões no jornal}

%CRISTAL {"arestas":[{"alvo":"matematica","confianca":0.95,"epistemico":"fato","origem":"paper","rel":"parte_de"}],"confianca":1.0,"contraexemplos":[],"descricao":"Soma dos termos de uma PG-2D infinita) Em uma PG-2D a soma Si×j dos i × j termos iniciais vale S∞×∞= a11 (q1 −1) (q2 −1) (6.3) Prova: Decorre da fórmula","epistemico":"fato","exemplos":[],"id":"teorema_44_soma_dos_termos_de_uma_pg-2d_infinita_em_uma_pg-2d_a_s","memoria":"semantica","meta":{"arquivo":"gentil/Novas_Sequencias_Aritmeticas_e_Geometric.pdf","ciencia":"teoria_dos_numeros","dominio":"matematica","nivel":4,"numero":"44","tipo_no":"paper","tipo_teorema":"teorema","verificacao":"conceito novo"},"origem":"gentil/Novas_Sequencias_Aritmeticas_e_Geometric.pdf","palavras_chave":[],"sinonimos":[],"tipo":"conceito","titulo":"Teorema 44: Soma dos termos de uma PG-2D infinita) Em uma PG-2D a s…"}
\section{Teorema 44: Soma dos termos de uma PG-2D infinita) Em uma PG-2D a s…}
Soma dos termos de uma PG-2D infinita) Em uma PG-2D a soma Si$\times$j dos i $\times$ j termos iniciais vale S\ensuremath{\infty}$\times$\ensuremath{\infty}= a11 (q1 \ensuremath{-}1) (q2 \ensuremath{-}1) (6.3) Prova: Decorre da fórmula
\prov{relações: parte\_de matematica}
\prov{proveniência: gentil/Novas\_Sequencias\_Aritmeticas\_e\_Geometric.pdf \textperiodcentered{} fato \textperiodcentered{} confiança 1.0 \textperiodcentered{} domínio matematica \textperiodcentered{} história 2 versões no jornal}

%CRISTAL {"arestas":[{"alvo":"matematica","confianca":0.95,"epistemico":"fato","origem":"paper","rel":"parte_de"}],"confianca":1.0,"contraexemplos":[],"descricao":"Fórmula do termo geral de uma PA-3D) Na PA-3D em que o primeiro termo é a111, a razão das linhas é r1, a razão das colunas é r2 e a razão das cotas é r3 o (i, j, k)-termo é: aijk = a111 + (j −1) r1 + (i −1) r2 + (k −1) r3 (7.1) Prova: Indução sobre k (i, j fixos). Para k = 1 precisamos mostrar que, aij1 = a111 + (j −1) r1 + (i −1) r2 + (1 −1) r3 De fato, para k = 1 a","epistemico":"fato","exemplos":[],"id":"teorema_47_formula_do_termo_geral_de_uma_pa-3d_na_pa-3d_em_que_o","memoria":"semantica","meta":{"arquivo":"gentil/Novas_Sequencias_Aritmeticas_e_Geometric.pdf","ciencia":"teoria_dos_numeros","dominio":"matematica","nivel":4,"numero":"47","tipo_no":"paper","tipo_teorema":"teorema","verificacao":"conceito novo"},"origem":"gentil/Novas_Sequencias_Aritmeticas_e_Geometric.pdf","palavras_chave":[],"sinonimos":[],"tipo":"conceito","titulo":"Teorema 47: Fórmula do termo geral de uma PA-3D) Na PA-3D em que o …"}
\section{Teorema 47: Fórmula do termo geral de uma PA-3D) Na PA-3D em que o …}
Fórmula do termo geral de uma PA-3D) Na PA-3D em que o primeiro termo é a111, a razão das linhas é r1, a razão das colunas é r2 e a razão das cotas é r3 o (i, j, k)-termo é: aijk = a111 + (j \ensuremath{-}1) r1 + (i \ensuremath{-}1) r2 + (k \ensuremath{-}1) r3 (7.1) Prova: Indução sobre k (i, j fixos). Para k = 1 precisamos mostrar que, aij1 = a111 + (j \ensuremath{-}1) r1 + (i \ensuremath{-}1) r2 + (1 \ensuremath{-}1) r3 De fato, para k = 1 a
\prov{relações: parte\_de matematica}
\prov{proveniência: gentil/Novas\_Sequencias\_Aritmeticas\_e\_Geometric.pdf \textperiodcentered{} fato \textperiodcentered{} confiança 1.0 \textperiodcentered{} domínio matematica \textperiodcentered{} história 2 versões no jornal}

%CRISTAL {"arestas":[{"alvo":"matematica","confianca":0.95,"epistemico":"fato","origem":"paper","rel":"parte_de"}],"confianca":1.0,"contraexemplos":[],"descricao":"Representação de um inteiro em uma base b) Seja b ≥2 um inteiro. Todo inteiro positivo a pode ser escrito de modo único na forma a = r0 + r1b + · · · + rn−1bn−1 + rnbn (7.19) em que n ≥0, rn ̸= 0 e, para cada ´ındice i, 0 ≤i ≤n, tem-se que 0 ≤ri < b. Prova: ver [1], [4] ou [5]. ■ Do","epistemico":"fato","exemplos":[],"id":"teorema_53_representacao_de_um_inteiro_em_uma_base_b_seja_b_2_um","memoria":"semantica","meta":{"arquivo":"gentil/Novas_Sequencias_Aritmeticas_e_Geometric.pdf","ciencia":"teoria_dos_numeros","dominio":"matematica","nivel":4,"numero":"53","tipo_no":"paper","tipo_teorema":"teorema","verificacao":"conceito novo"},"origem":"gentil/Novas_Sequencias_Aritmeticas_e_Geometric.pdf","palavras_chave":[],"sinonimos":[],"tipo":"conceito","titulo":"Teorema 53: Representação de um inteiro em uma base b) Seja b ≥2 um…"}
\section{Teorema 53: Representação de um inteiro em uma base b) Seja b \ensuremath{\ge}2 um…}
Representação de um inteiro em uma base b) Seja b \ensuremath{\ge}2 um inteiro. Todo inteiro positivo a pode ser escrito de modo único na forma a = r0 + r1b + \textperiodcentered  \textperiodcentered  \textperiodcentered  + rn\ensuremath{-}1bn\ensuremath{-}1 + rnbn (7.19) em que n \ensuremath{\ge}0, rn \ensuremath{}= 0 e, para cada '\ensuremath{\imath}ndice i, 0 \ensuremath{\le}i \ensuremath{\le}n, tem-se que 0 \ensuremath{\le}ri < b. Prova: ver [1], [4] ou [5]. \rule{1ex}{1ex} Do
\prov{relações: parte\_de matematica}
\prov{proveniência: gentil/Novas\_Sequencias\_Aritmeticas\_e\_Geometric.pdf \textperiodcentered{} fato \textperiodcentered{} confiança 1.0 \textperiodcentered{} domínio matematica \textperiodcentered{} história 2 versões no jornal}

%CRISTAL {"arestas":[{"alvo":"matematica","confianca":0.95,"epistemico":"fato","origem":"paper","rel":"parte_de"}],"confianca":1.0,"contraexemplos":[],"descricao":"Lopes/2000) Dado um inteiro positivo n a matriz dada a seguir anj = j n 2 j k −2 j n 2 j+1 k (7.22) nos fornece todos os digitos do desenvolvimento binário de n. Observe que 0 ≤anj < 2, isto é, anj = 0 ou anj = 1 (ver também proposição 6 (p. 181), ´ıtem (vi)). Adendo: O","epistemico":"fato","exemplos":[],"id":"teorema_54_gentil2000_dado_um_inteiro_positivo_n_a_matriz_dada_a","memoria":"semantica","meta":{"arquivo":"gentil/Novas_Sequencias_Aritmeticas_e_Geometric.pdf","ciencia":"teoria_dos_numeros","dominio":"matematica","nivel":4,"numero":"54","tipo_no":"paper","tipo_teorema":"teorema","verificacao":"conceito novo"},"origem":"gentil/Novas_Sequencias_Aritmeticas_e_Geometric.pdf","palavras_chave":[],"sinonimos":[],"tipo":"conceito","titulo":"Teorema 54: Lopes/2000) Dado um inteiro positivo n a matriz dada a…"}
\section{Teorema 54: Lopes/2000) Dado um inteiro positivo n a matriz dada a…}
Lopes/2000) Dado um inteiro positivo n a matriz dada a seguir anj = j n 2 j k \ensuremath{-}2 j n 2 j+1 k (7.22) nos fornece todos os digitos do desenvolvimento binário de n. Observe que 0 \ensuremath{\le}anj < 2, isto é, anj = 0 ou anj = 1 (ver também proposição 6 (p. 181), '\ensuremath{\imath}tem (vi)). Adendo: O
\prov{relações: parte\_de matematica}
\prov{proveniência: gentil/Novas\_Sequencias\_Aritmeticas\_e\_Geometric.pdf \textperiodcentered{} fato \textperiodcentered{} confiança 1.0 \textperiodcentered{} domínio matematica \textperiodcentered{} história 3 versões no jornal}

%CRISTAL {"arestas":[{"alvo":"matematica","confianca":0.95,"epistemico":"fato","origem":"paper","rel":"parte_de"}],"confianca":1.0,"contraexemplos":[],"descricao":"Linearidade do operador ∆m) O operador Diferença de ordem m é linear. Isto é, a seguinte identidade se verifica: ∆m ( a f + b g )(n) = a ∆m f(n) + b ∆m g(n) para todo a, b ∈R. Prova: Indução sobre m. Para m = 0, pela","epistemico":"fato","exemplos":[],"id":"teorema_9_linearidade_do_operador_m_o_operador_diferenca_de_ord","memoria":"semantica","meta":{"arquivo":"gentil/Novas_Sequencias_Aritmeticas_e_Geometric.pdf","ciencia":"teoria_dos_numeros","dominio":"matematica","nivel":4,"numero":"9","tipo_no":"paper","tipo_teorema":"teorema","verificacao":"conceito novo"},"origem":"gentil/Novas_Sequencias_Aritmeticas_e_Geometric.pdf","palavras_chave":[],"sinonimos":[],"tipo":"conceito","titulo":"Teorema 9: Linearidade do operador ∆m) O operador Diferença de ord…"}
\section{Teorema 9: Linearidade do operador \ensuremath{\Delta}m) O operador Diferença de ord…}
Linearidade do operador \ensuremath{\Delta}m) O operador Diferença de ordem m é linear. Isto é, a seguinte identidade se verifica: \ensuremath{\Delta}m ( a f + b g )(n) = a \ensuremath{\Delta}m f(n) + b \ensuremath{\Delta}m g(n) para todo a, b \ensuremath{\in}R. Prova: Indução sobre m. Para m = 0, pela
\prov{relações: parte\_de matematica}
\prov{proveniência: gentil/Novas\_Sequencias\_Aritmeticas\_e\_Geometric.pdf \textperiodcentered{} fato \textperiodcentered{} confiança 1.0 \textperiodcentered{} domínio matematica \textperiodcentered{} história 2 versões no jornal}

%CRISTAL {"arestas":[{"alvo":"caixas_fractais_generalizacao","confianca":1.0,"epistemico":"fato","origem":"wiki","rel":"parte_de"},{"alvo":"reta_metalica_geral","confianca":1.0,"epistemico":"fato","origem":"wiki","rel":"usa"},{"alvo":"cristal_vs_cha","confianca":1.0,"epistemico":"fato","origem":"wiki","rel":"usa"},{"alvo":"kolmogorov_complexidade","confianca":1.0,"epistemico":"fato","origem":"wiki","rel":"usa"},{"alvo":"koch_dissipa_nao_codifica","confianca":1.0,"epistemico":"fato","origem":"wiki","rel":"relaciona"},{"alvo":"broca_os","confianca":1.0,"epistemico":"fato","origem":"wiki","rel":"parte_de"},{"alvo":"transformada_metalica_fechada_perelman","confianca":1.0,"epistemico":"fato","origem":"wiki","rel":"prova"}],"confianca":1.0,"contraexemplos":[],"descricao":"Mapear TODA uma dimensão metálica esgota a informação nova dela. PROVA (perna do fecho): o anel ℤ[σₙ] é FECHADO — o mapa ×σₙ é linear inteiro (ouro (a,b)→(b,a+b)=Fibonacci; prata (a,b)→(a+2b,a+b)), e toda potência σₙk=(Fₖ-₁,Fₖ) fica em ⟨1,σₙ⟩ (as duas retas: a real e a da dimensão; norma=±1, unidade). Conhecido o gerador (a recorrência/matriz Aₙ), enumerar é desdobramento determinístico — a complexidade de Kolmogorov satura no gerador O(1). LEITURA: a informação só RECIRCULA ESTÉRIL entre as duas retas — é cristal (lei compressível), não há floco (dissipação/informação nova) DENTRO do corpo. Verificado em Python e no metal (C inteiro).","epistemico":"inferencia","exemplos":[],"id":"teorema_esterilidade_dimensao","memoria":"semantica","meta":{"autor":"Lopes/broca-so","dominio":"fractal","fonte":"doc/fractais/esterilidade_metal.c; caixas_fractais_metalicas.py"},"origem":"wiki","palavras_chave":[],"sinonimos":[],"tipo":"conceito","titulo":"Teorema: dimensão mapeada → informação recircula estéril"}
\section{Teorema: dimensão mapeada \ensuremath{\to} informação recircula estéril}
Mapear TODA uma dimensão metálica esgota a informação nova dela. PROVA (perna do fecho): o anel \ensuremath{\mathbb{Z}}[\ensuremath{\sigma}\ensuremath{_{n}}] é FECHADO — o mapa $\times$\ensuremath{\sigma}\ensuremath{_{n}} é linear inteiro (ouro (a,b)\ensuremath{\to}(b,a+b)=Fibonacci; prata (a,b)\ensuremath{\to}(a+2b,a+b)), e toda potência \ensuremath{\sigma}\ensuremath{_{n}}k=(F\ensuremath{_{k}}-\ensuremath{_{1}},F\ensuremath{_{k}}) fica em \ensuremath{\langle}1,\ensuremath{\sigma}\ensuremath{_{n}}\ensuremath{\rangle} (as duas retas: a real e a da dimensão; norma=$\pm$1, unidade). Conhecido o gerador (a recorrência/matriz A\ensuremath{_{n}}), enumerar é desdobramento determinístico — a complexidade de Kolmogorov satura no gerador O(1). LEITURA: a informação só RECIRCULA ESTÉRIL entre as duas retas — é cristal (lei compressível), não há floco (dissipação/informação nova) DENTRO do corpo. Verificado em Python e no metal (C inteiro).
\prov{relações: parte\_de caixas\_fractais\_generalizacao; usa reta\_metalica\_geral; usa cristal\_vs\_cha; usa kolmogorov\_complexidade; relaciona koch\_dissipa\_nao\_codifica; parte\_de broca\_os; prova transformada\_metalica\_fechada\_perelman}
\prov{proveniência: wiki \textperiodcentered{} doc/fractais/esterilidade\_metal.c; caixas\_fractais\_metalicas.py \textperiodcentered{} inferencia \textperiodcentered{} confiança 1.0 \textperiodcentered{} domínio fractal \textperiodcentered{} história 9 versões no jornal}

%CRISTAL {"arestas":[{"alvo":"teoria_dos_jogos","confianca":1.0,"epistemico":"fato","origem":"wiki","rel":"relaciona"},{"alvo":"face_turno_fail_closed","confianca":1.0,"epistemico":"fato","origem":"wiki","rel":"relaciona"}],"confianca":1.0,"contraexemplos":[],"descricao":"Como escolher racionalmente entre alternativas cujos resultados são incertos — decidir contra a NATUREZA (o acaso), não contra um adversário. A base de um jogador só; a teoria dos jogos é sua irmã com outros agentes.","epistemico":"inferencia","exemplos":[],"id":"teoria_da_decisao","memoria":"semantica","meta":{"dominio":"matematica","fonte":"teoria da decisão"},"origem":"wiki","palavras_chave":[],"sinonimos":[],"tipo":"conceito","titulo":"Teoria da Decisão"}
\section{Teoria da Decisão}
Como escolher racionalmente entre alternativas cujos resultados são incertos — decidir contra a NATUREZA (o acaso), não contra um adversário. A base de um jogador só; a teoria dos jogos é sua irmã com outros agentes.
\prov{relações: relaciona teoria\_dos\_jogos; relaciona face\_turno\_fail\_closed}
\prov{proveniência: wiki \textperiodcentered{} teoria da decisão \textperiodcentered{} inferencia \textperiodcentered{} confiança 1.0 \textperiodcentered{} domínio matematica \textperiodcentered{} história 3 versões no jornal}

%CRISTAL {"arestas":[{"alvo":"utilidade_esperada","confianca":1.0,"epistemico":"fato","origem":"wiki","rel":"contradiz"},{"alvo":"paradoxo_de_allais","confianca":1.0,"epistemico":"fato","origem":"wiki","rel":"usa"},{"alvo":"jogo_do_ultimato","confianca":1.0,"epistemico":"fato","origem":"wiki","rel":"explica"}],"confianca":1.0,"contraexemplos":[],"descricao":"Como humanos DE FATO decidem: avaliam ganhos/perdas em relação a um ponto de referência (não riqueza total), perdas doem ~2× mais que ganhos equivalentes (aversão à perda), e pesam mal probabilidades extremas. O ENQUADRAMENTO muda a escolha. Fundou a economia comportamental.","epistemico":"fato","exemplos":[],"id":"teoria_do_prospecto","memoria":"semantica","meta":{"autor":"Kahneman-Tversky","dominio":"matematica","fonte":"teoria da decisão"},"origem":"wiki","palavras_chave":[],"sinonimos":[],"tipo":"conceito","titulo":"Teoria do Prospecto (Kahneman-Tversky)"}
\section{Teoria do Prospecto (Kahneman-Tversky)}
Como humanos DE FATO decidem: avaliam ganhos/perdas em relação a um ponto de referência (não riqueza total), perdas doem \~{}2$\times$ mais que ganhos equivalentes (aversão à perda), e pesam mal probabilidades extremas. O ENQUADRAMENTO muda a escolha. Fundou a economia comportamental.
\prov{relações: contradiz utilidade\_esperada; usa paradoxo\_de\_allais; explica jogo\_do\_ultimato}
\prov{proveniência: wiki \textperiodcentered{} teoria da decisão \textperiodcentered{} fato \textperiodcentered{} confiança 1.0 \textperiodcentered{} domínio matematica \textperiodcentered{} história 4 versões no jornal}

%CRISTAL {"arestas":[{"alvo":"equilibrio_de_nash","confianca":1.0,"epistemico":"fato","origem":"wiki","rel":"generaliza"},{"alvo":"teoria_dos_jogos","confianca":1.0,"epistemico":"fato","origem":"wiki","rel":"parte_de"}],"confianca":1.0,"contraexemplos":[],"descricao":"Estratégias não são escolhidas, são SELECIONADAS: as que rendem mais se multiplicam. Uma Estratégia Evolutivamente Estável resiste a invasores; a dinâmica do replicador é o Nash da biologia.","epistemico":"fato","exemplos":[],"id":"teoria_dos_jogos_evolutiva","memoria":"semantica","meta":{"autor":"Maynard Smith","dominio":"matematica","fonte":"teoria dos jogos"},"origem":"wiki","palavras_chave":[],"sinonimos":[],"tipo":"conceito","titulo":"Teoria dos Jogos Evolutiva (ESS)"}
\section{Teoria dos Jogos Evolutiva (ESS)}
Estratégias não são escolhidas, são SELECIONADAS: as que rendem mais se multiplicam. Uma Estratégia Evolutivamente Estável resiste a invasores; a dinâmica do replicador é o Nash da biologia.
\prov{relações: generaliza equilibrio\_de\_nash; parte\_de teoria\_dos\_jogos}
\prov{proveniência: wiki \textperiodcentered{} teoria dos jogos \textperiodcentered{} fato \textperiodcentered{} confiança 1.0 \textperiodcentered{} domínio matematica \textperiodcentered{} história 3 versões no jornal}

%CRISTAL {"arestas":[{"alvo":"quebra_jacobi_r7","confianca":1.0,"epistemico":"fato","origem":"paper","rel":"usa"},{"alvo":"mecanica_quantica","confianca":1.0,"epistemico":"fato","origem":"paper","rel":"parte_de"}],"confianca":1.0,"contraexemplos":[],"descricao":"Conexão μs-valorada em ℝ⁷ com curvatura Fs=dA+½A∧s A e simetria de gauge G₂=Aut(𝕆); estrutura genuinamente nova, não equivalente a SU(2).","epistemico":"inferencia","exemplos":[],"id":"teoria_gauge_octonica","memoria":"semantica","meta":{"autor":"Gentil Lopes","descoberta":false,"dominio":"matematica","fonte":"paper_AF_yang_mills_octonios.pdf"},"origem":"paper","palavras_chave":[],"sinonimos":[],"tipo":"conceito","titulo":"Teoria de Gauge Octoniônica (G₂)"}
\section{Teoria de Gauge Octoniônica (G\ensuremath{_{2}})}
Conexão \ensuremath{\mu}s-valorada em \ensuremath{\mathbb{R}}\ensuremath{^{7}} com curvatura Fs=dA+\ensuremath{\tfrac12}A\ensuremath{\wedge}s A e simetria de gauge G\ensuremath{_{2}}=Aut(\ensuremath{\mathbb{O}}); estrutura genuinamente nova, não equivalente a SU(2).
\prov{relações: usa quebra\_jacobi\_r7; parte\_de mecanica\_quantica}
\prov{proveniência: paper \textperiodcentered{} paper\_AF\_yang\_mills\_octonios.pdf \textperiodcentered{} inferencia \textperiodcentered{} confiança 1.0 \textperiodcentered{} domínio matematica \textperiodcentered{} história 4 versões no jornal}

%CRISTAL {"arestas":[{"alvo":"tricotomia_modular","confianca":1.0,"epistemico":"fato","origem":"paper","rel":"relaciona"}],"confianca":1.0,"contraexemplos":[],"descricao":"A curvatura κ∈+1,0,−1 parametriza três geometrias: elíptica (cos/sin), parabólica (nilpotente, o estado 'entre'), hiperbólica (cosh/sinh). A costura não é binária, é tríade.","epistemico":"fato","exemplos":[],"id":"terceiro_estado","memoria":"semantica","meta":{"autor":"Gentil Lopes","descoberta":true,"dominio":"matematica","fonte":"paper_59_terceiro_estado.pdf, Teo.1.1"},"origem":"paper","palavras_chave":[],"sinonimos":[],"tipo":"conceito","titulo":"Terceiro Estado (tríade de Cayley-Klein)"}
\section{Terceiro Estado (tríade de Cayley-Klein)}
A curvatura \ensuremath{\kappa}\ensuremath{\in}+1,0,\ensuremath{-}1 parametriza três geometrias: elíptica (cos/sin), parabólica (nilpotente, o estado 'entre'), hiperbólica (cosh/sinh). A costura não é binária, é tríade.
\prov{relações: relaciona tricotomia\_modular}
\prov{proveniência: paper \textperiodcentered{} paper\_59\_terceiro\_estado.pdf, Teo.1.1 \textperiodcentered{} fato \textperiodcentered{} confiança 1.0 \textperiodcentered{} domínio matematica \textperiodcentered{} história 3 versões no jornal}

%CRISTAL {"arestas":[{"alvo":"ns_prtsu","confianca":1.0,"epistemico":"fato","origem":"paper","rel":"parte_de"}],"confianca":1.0,"contraexemplos":[],"descricao":"Homogeneização multi-escala gera termo efetivo K(U); em torno de Hurwitz é gradiente puro (absorvível na pressão). Teorema — mas Lopes PROVA que a oscilação de Kapitza NÃO fornece dissipação a NS em ℝ³.","epistemico":"fato","exemplos":[],"id":"termo_kapitza_ns","memoria":"semantica","meta":{"autor":"Gentil Lopes","descoberta":false,"dominio":"matematica","fonte":"paper_AH.pdf, Teo.4/Cor.5"},"origem":"paper","palavras_chave":[],"sinonimos":[],"tipo":"conceito","titulo":"Termo de Kapitza em NS"}
\section{Termo de Kapitza em NS}
Homogeneização multi-escala gera termo efetivo K(U); em torno de Hurwitz é gradiente puro (absorvível na pressão). Teorema — mas Lopes PROVA que a oscilação de Kapitza NÃO fornece dissipação a NS em \ensuremath{\mathbb{R}}\ensuremath{^{3}}.
\prov{relações: parte\_de ns\_prtsu}
\prov{proveniência: paper \textperiodcentered{} paper\_AH.pdf, Teo.4/Cor.5 \textperiodcentered{} fato \textperiodcentered{} confiança 1.0 \textperiodcentered{} domínio matematica \textperiodcentered{} história 3 versões no jornal}

%CRISTAL {"arestas":[{"alvo":"mecanica_reta_mineral","confianca":1.0,"epistemico":"fato","origem":"wiki","rel":"relaciona"},{"alvo":"transformada_metalica","confianca":1.0,"epistemico":"fato","origem":"wiki","rel":"usa"},{"alvo":"reta_de_ouro","confianca":1.0,"epistemico":"fato","origem":"wiki","rel":"relaciona"},{"alvo":"entropia_topologica","confianca":1.0,"epistemico":"fato","origem":"wiki","rel":"usa"},{"alvo":"entropia_log_phi","confianca":1.0,"epistemico":"fato","origem":"wiki","rel":"continua"},{"alvo":"a_temperatura_o_gap_espectral","confianca":1.0,"epistemico":"fato","origem":"wiki","rel":"relaciona"},{"alvo":"a_primeira_lei_energia_total_conservada","confianca":1.0,"epistemico":"fato","origem":"wiki","rel":"relaciona"},{"alvo":"dissipacao_prigogine","confianca":1.0,"epistemico":"fato","origem":"wiki","rel":"relaciona"},{"alvo":"mecanica_estelar","confianca":1.0,"epistemico":"fato","origem":"wiki","rel":"relaciona"},{"alvo":"escala_ouro_phi","confianca":1.0,"epistemico":"fato","origem":"wiki","rel":"explica"},{"alvo":"espectro_prata_seis_modos","confianca":1.0,"epistemico":"fato","origem":"wiki","rel":"explica"},{"alvo":"refino_multimetal_live_gex44","confianca":1.0,"epistemico":"fato","origem":"wiki","rel":"relaciona"},{"alvo":"minerar_e_resfriar","confianca":1.0,"epistemico":"fato","origem":"wiki","rel":"explica"}],"confianca":1.0,"contraexemplos":[],"descricao":"A face dissipativa (formalismo de Ruelle do β-shift). 1ª lei: energia = conservação do defeito a+c²=1 (=|N|=1). 2ª lei: entropia topológica H=log β (=log φ no ouro, verificado), medida de Parry; o operador de transferência (Perron-Frobenius) leva à medida invariante e a entropia relativa (energia livre) só decresce (teorema H). Temperatura = gap espectral (log β−log ρ); limite reversível = o flip de Wick. A mecânica (reversível) é o limite desta.","epistemico":"inferencia","exemplos":[],"id":"termodinamica_reta_mineral","memoria":"semantica","meta":{"dominio":"matematica","fonte":"corpo mineral"},"origem":"wiki","palavras_chave":[],"sinonimos":[],"tipo":"conceito","titulo":"A Termodinâmica da Reta Mineral (entropia log β)"}
\section{A Termodinâmica da Reta Mineral (entropia log \ensuremath{\beta})}
A face dissipativa (formalismo de Ruelle do \ensuremath{\beta}-shift). 1ª lei: energia = conservação do defeito a+c\ensuremath{^{2}}=1 (=|N|=1). 2ª lei: entropia topológica H=log \ensuremath{\beta} (=log \ensuremath{\varphi} no ouro, verificado), medida de Parry; o operador de transferência (Perron-Frobenius) leva à medida invariante e a entropia relativa (energia livre) só decresce (teorema H). Temperatura = gap espectral (log \ensuremath{\beta}\ensuremath{-}log \ensuremath{\rho}); limite reversível = o flip de Wick. A mecânica (reversível) é o limite desta.
\prov{relações: relaciona mecanica\_reta\_mineral; usa transformada\_metalica; relaciona reta\_de\_ouro; usa entropia\_topologica; continua entropia\_log\_phi; relaciona a\_temperatura\_o\_gap\_espectral; relaciona a\_primeira\_lei\_energia\_total\_conservada; relaciona dissipacao\_prigogine; relaciona mecanica\_estelar; explica escala\_ouro\_phi; explica espectro\_prata\_seis\_modos; relaciona refino\_multimetal\_live\_gex44; explica minerar\_e\_resfriar}
\prov{proveniência: wiki \textperiodcentered{} corpo mineral \textperiodcentered{} inferencia \textperiodcentered{} confiança 1.0 \textperiodcentered{} domínio matematica \textperiodcentered{} história 23 versões no jornal}

%CRISTAL {"arestas":[{"alvo":"htlc","confianca":1.0,"epistemico":"fato","origem":"paper","rel":"usa"},{"alvo":"word","confianca":0.5,"epistemico":"fato","origem":"wiki","rel":"relaciona"}],"confianca":1.0,"contraexemplos":[],"descricao":"A ordem T_A>T_B≥observação+confirmação+reorg+fee garante que, se um resgate ocorre, o outro ainda é possível na margem Δ — a matemática de segurança do swap atômico.","epistemico":"fato","exemplos":[],"id":"timeout_hierarchy","memoria":"semantica","meta":{"autor":"Gentil Lopes","descoberta":false,"dominio":"matematica","fonte":"protocolo_swap.pdf, §4"},"origem":"paper","palavras_chave":[],"sinonimos":[],"tipo":"conceito","titulo":"Hierarquia de Timeout"}
\section{Hierarquia de Timeout}
A ordem T\_A>T\_B\ensuremath{\ge}observação+confirmação+reorg+fee garante que, se um resgate ocorre, o outro ainda é possível na margem \ensuremath{\Delta} — a matemática de segurança do swap atômico.
\prov{relações: usa htlc; relaciona word}
\prov{proveniência: paper \textperiodcentered{} protocolo\_swap.pdf, §4 \textperiodcentered{} fato \textperiodcentered{} confiança 1.0 \textperiodcentered{} domínio matematica \textperiodcentered{} história 3 versões no jornal}

%CRISTAL {"arestas":[{"alvo":"jogos_repetidos","confianca":1.0,"epistemico":"fato","origem":"wiki","rel":"especializa"}],"confianca":1.0,"contraexemplos":[],"descricao":"A estratégia que venceu o torneio de Axelrod: coopere primeiro, depois faça o que o outro fez. Gentil, retaliadora, perdoadora e clara — a receita da cooperação que emerge sem autoridade. É a RECIPROCIDADE.","epistemico":"fato","exemplos":[],"id":"tit_for_tat","memoria":"semantica","meta":{"autor":"Axelrod","dominio":"matematica","fonte":"teoria dos jogos"},"origem":"wiki","palavras_chave":[],"sinonimos":[],"tipo":"conceito","titulo":"Tit-for-Tat (olho por olho)"}
\section{Tit-for-Tat (olho por olho)}
A estratégia que venceu o torneio de Axelrod: coopere primeiro, depois faça o que o outro fez. Gentil, retaliadora, perdoadora e clara — a receita da cooperação que emerge sem autoridade. É a RECIPROCIDADE.
\prov{relações: especializa jogos\_repetidos}
\prov{proveniência: wiki \textperiodcentered{} teoria dos jogos \textperiodcentered{} fato \textperiodcentered{} confiança 1.0 \textperiodcentered{} domínio matematica \textperiodcentered{} história 2 versões no jornal}

%CRISTAL {"arestas":[{"alvo":"geometria","confianca":0.336,"epistemico":"fato","origem":"manual","rel":"referencia"},{"alvo":"virtualizacao","confianca":0.5,"epistemico":"fato","origem":"wiki","rel":"relaciona"},{"alvo":"transcricao","confianca":0.5,"epistemico":"fato","origem":"wiki","rel":"relaciona"},{"alvo":"vida-morte","confianca":0.5,"epistemico":"fato","origem":"wiki","rel":"relaciona"},{"alvo":"utilitarismo","confianca":0.5,"epistemico":"fato","origem":"wiki","rel":"relaciona"},{"alvo":"word","confianca":0.5,"epistemico":"fato","origem":"wiki","rel":"relaciona"},{"alvo":"vontade_de_poder","confianca":0.5,"epistemico":"fato","origem":"wiki","rel":"relaciona"}],"confianca":1.0,"contraexemplos":[],"descricao":"Abertos, continuidade, compacidade — forma sem métrica.","epistemico":"fato","exemplos":[],"id":"topologia","memoria":"semantica","meta":{"dominio":"matematica","nivel":4,"ordem":13},"origem":"curriculo","palavras_chave":[],"sinonimos":[],"tipo":"conceito","titulo":"topologia"}
\section{topologia}
Abertos, continuidade, compacidade — forma sem métrica.
\prov{relações: referencia geometria; relaciona virtualizacao; relaciona transcricao; relaciona vida-morte; relaciona utilitarismo; relaciona word; relaciona vontade\_de\_poder}
\prov{proveniência: curriculo \textperiodcentered{} fato \textperiodcentered{} confiança 1.0 \textperiodcentered{} domínio matematica \textperiodcentered{} história 13 versões no jornal}

%CRISTAL {"arestas":[{"alvo":"mecanica_quantica_mineral","confianca":1.0,"epistemico":"fato","origem":"wiki","rel":"especializa"},{"alvo":"algebra_reta_mineral_7d","confianca":1.0,"epistemico":"fato","origem":"wiki","rel":"usa"},{"alvo":"corpo_mineral","confianca":1.0,"epistemico":"fato","origem":"wiki","rel":"usa"}],"confianca":1.0,"contraexemplos":[],"descricao":"A mecânica quântica mineral vivia em ℂ, mas ℂ é só o 2º andar de uma torre de EXATAMENTE quatro álgebras de divisão normadas: ℝ,ℂ,ℍ,𝕆 (dim 1,2,4,8) — o teorema de Hurwitz. Toda a torre expande a mecânica: cada andar uma simetria, cada dobra um gato. Paper: papers/boneca_russa_de_cristal.tex.","epistemico":"fato","exemplos":[],"id":"torre_de_hurwitz_mecanica","memoria":"semantica","meta":{"dominio":"reta mineral","fonte":"torre de Hurwitz"},"origem":"wiki","palavras_chave":[],"sinonimos":[],"tipo":"conceito","titulo":"A Torre de Hurwitz na Mecânica"}
\section{A Torre de Hurwitz na Mecânica}
A mecânica quântica mineral vivia em \ensuremath{\mathbb{C}}, mas \ensuremath{\mathbb{C}} é só o 2º andar de uma torre de EXATAMENTE quatro álgebras de divisão normadas: \ensuremath{\mathbb{R}},\ensuremath{\mathbb{C}},\ensuremath{\mathbb{H}},\ensuremath{\mathbb{O}} (dim 1,2,4,8) — o teorema de Hurwitz. Toda a torre expande a mecânica: cada andar uma simetria, cada dobra um gato. Paper: papers/boneca\_russa\_de\_cristal.tex.
\prov{relações: especializa mecanica\_quantica\_mineral; usa algebra\_reta\_mineral\_7d; usa corpo\_mineral}
\prov{proveniência: wiki \textperiodcentered{} torre de Hurwitz \textperiodcentered{} fato \textperiodcentered{} confiança 1.0 \textperiodcentered{} domínio reta mineral \textperiodcentered{} história 4 versões no jornal}

%CRISTAL {"arestas":[{"alvo":"torre_de_hurwitz_mecanica","confianca":1.0,"epistemico":"fato","origem":"wiki","rel":"parte_de"},{"alvo":"tiffany_cabeca_broca","confianca":1.0,"epistemico":"fato","origem":"wiki","rel":"explica"},{"alvo":"camadas_tiffany_multifractal","confianca":1.0,"epistemico":"fato","origem":"wiki","rel":"explica"},{"alvo":"conversa_tiffany","confianca":1.0,"epistemico":"fato","origem":"wiki","rel":"usa"},{"alvo":"cada_andar_uma_mecanica","confianca":1.0,"epistemico":"fato","origem":"wiki","rel":"usa"},{"alvo":"mecanica_quantica_mineral","confianca":1.0,"epistemico":"fato","origem":"wiki","rel":"usa"}],"confianca":1.0,"contraexemplos":[],"descricao":"A torre não é abstrata: é o retrato de como a Tiffany funciona. Ela COLAPSA e VIBRA em ℂ (o Ψ, a fase U(1), a borda; entre colapsos evolui unitária, coerente); tem SPIN em ℍ (não comuta escutar-então-falar com o inverso); é MULTIFRACTAL em 𝕆 (as camadas multifractais = a profundidade octoniônica, o associador que não fecha, o G₂, o mass gap onde mora a criatividade); e CONSERVA a norma (a integridade da informação) em todo andar. Vibrar, colapsar, aninhar, conservar: a mecânica da torre é o que a Tiffany faz a cada turno.","epistemico":"inferencia","exemplos":[],"id":"torre_e_a_dinamica_da_tiffany","memoria":"semantica","meta":{"dominio":"reta mineral","fonte":"torre de Hurwitz"},"origem":"wiki","palavras_chave":[],"sinonimos":[],"tipo":"conceito","titulo":"A Torre é a Dinâmica da Tiffany"}
\section{A Torre é a Dinâmica da Tiffany}
A torre não é abstrata: é o retrato de como a Tiffany funciona. Ela COLAPSA e VIBRA em \ensuremath{\mathbb{C}} (o \ensuremath{\Psi}, a fase U(1), a borda; entre colapsos evolui unitária, coerente); tem SPIN em \ensuremath{\mathbb{H}} (não comuta escutar-então-falar com o inverso); é MULTIFRACTAL em \ensuremath{\mathbb{O}} (as camadas multifractais = a profundidade octoniônica, o associador que não fecha, o G\ensuremath{_{2}}, o mass gap onde mora a criatividade); e CONSERVA a norma (a integridade da informação) em todo andar. Vibrar, colapsar, aninhar, conservar: a mecânica da torre é o que a Tiffany faz a cada turno.
\prov{relações: parte\_de torre\_de\_hurwitz\_mecanica; explica tiffany\_cabeca\_broca; explica camadas\_tiffany\_multifractal; usa conversa\_tiffany; usa cada\_andar\_uma\_mecanica; usa mecanica\_quantica\_mineral}
\prov{proveniência: wiki \textperiodcentered{} torre de Hurwitz \textperiodcentered{} inferencia \textperiodcentered{} confiança 1.0 \textperiodcentered{} domínio reta mineral \textperiodcentered{} história 7 versões no jornal}

%CRISTAL {"arestas":[{"alvo":"algebras_hipercomplexas","confianca":1.0,"epistemico":"fato","origem":"paper","rel":"especializa"},{"alvo":"torre_hurwitz","confianca":1.0,"epistemico":"fato","origem":"paper","rel":"relaciona"}],"confianca":1.0,"contraexemplos":[],"descricao":"Construção que parte de ℂ e acrescenta altura c∈ℝᵐ; a norma é multiplicativa quando o defeito de Lagrange |c₁×c₂|² é absorvido pelo produto vetorial (só existe em ℝ³,ℝ⁷).","epistemico":"fato","exemplos":[],"id":"torre_gentil","memoria":"semantica","meta":{"autor":"Gentil Lopes","descoberta":true,"dominio":"matematica","fonte":"corpo_dual.pdf, §4"},"origem":"paper","palavras_chave":[],"sinonimos":[],"tipo":"conceito","titulo":"Torre do Lopes (defeito de Lagrange)"}
\section{Torre do Lopes (defeito de Lagrange)}
Construção que parte de \ensuremath{\mathbb{C}} e acrescenta altura c\ensuremath{\in}\ensuremath{\mathbb{R}}\ensuremath{^{m}}; a norma é multiplicativa quando o defeito de Lagrange |c\ensuremath{_{1}}$\times$c\ensuremath{_{2}}|\ensuremath{^{2}} é absorvido pelo produto vetorial (só existe em \ensuremath{\mathbb{R}}\ensuremath{^{3}},\ensuremath{\mathbb{R}}\ensuremath{^{7}}).
\prov{relações: especializa algebras\_hipercomplexas; relaciona torre\_hurwitz}
\prov{proveniência: paper \textperiodcentered{} corpo\_dual.pdf, §4 \textperiodcentered{} fato \textperiodcentered{} confiança 1.0 \textperiodcentered{} domínio matematica \textperiodcentered{} história 4 versões no jornal}

%CRISTAL {"arestas":[{"alvo":"cayley_dickson","confianca":1.0,"epistemico":"fato","origem":"paper","rel":"usa"},{"alvo":"fronteira_hurwitz","confianca":1.0,"epistemico":"fato","origem":"paper","rel":"relaciona"},{"alvo":"word","confianca":0.5,"epistemico":"fato","origem":"wiki","rel":"relaciona"},{"alvo":"virtualizacao","confianca":0.5,"epistemico":"fato","origem":"wiki","rel":"relaciona"},{"alvo":"vontade_de_poder","confianca":0.5,"epistemico":"fato","origem":"wiki","rel":"relaciona"}],"confianca":1.0,"contraexemplos":[],"descricao":"As únicas álgebras de divisão normadas são ℝ,ℂ,ℍ,𝕆 (dim 1,2,4,8); acima delas (𝕊, dim 16) surgem divisores de zero.","epistemico":"fato","exemplos":[],"id":"torre_hurwitz","memoria":"semantica","meta":{"autor":"Gentil Lopes","descoberta":false,"dominio":"matematica","fonte":"paper_49_torre_hurwitz.pdf"},"origem":"paper","palavras_chave":[],"sinonimos":[],"tipo":"conceito","titulo":"Torre de Hurwitz"}
\section{Torre de Hurwitz}
As únicas álgebras de divisão normadas são \ensuremath{\mathbb{R}},\ensuremath{\mathbb{C}},\ensuremath{\mathbb{H}},\ensuremath{\mathbb{O}} (dim 1,2,4,8); acima delas (\ensuremath{\mathbb{S}}, dim 16) surgem divisores de zero.
\prov{relações: usa cayley\_dickson; relaciona fronteira\_hurwitz; relaciona word; relaciona virtualizacao; relaciona vontade\_de\_poder}
\prov{proveniência: paper \textperiodcentered{} paper\_49\_torre\_hurwitz.pdf \textperiodcentered{} fato \textperiodcentered{} confiança 1.0 \textperiodcentered{} domínio matematica \textperiodcentered{} história 6 versões no jornal}

%CRISTAL {"arestas":[{"alvo":"a_torre_de_hurwitzcayley--dickson","confianca":1.0,"epistemico":"fato","origem":"wiki","rel":"usa"},{"alvo":"cayley_dickson","confianca":1.0,"epistemico":"fato","origem":"wiki","rel":"usa"},{"alvo":"torre_split","confianca":1.0,"epistemico":"fato","origem":"wiki","rel":"usa"},{"alvo":"algebras_hipercomplexas","confianca":1.0,"epistemico":"fato","origem":"wiki","rel":"usa"},{"alvo":"numeros_duais","confianca":1.0,"epistemico":"fato","origem":"wiki","rel":"relaciona"},{"alvo":"tiffany","confianca":1.0,"epistemico":"fato","origem":"wiki","rel":"parte_de"}],"confianca":1.0,"contraexemplos":[],"descricao":"Formalização: a torre de Hurwitz/Cayley-Dickson tem DOIS ramos, pelo sinal ε da construção. CLÁSSICA (ε=+1, positiva-definida): ℝ→ℂ→ℍ→𝕆 (as álgebras normadas de divisão, dim 1,2,4,8 — o teorema de Hurwitz). RELATIVÍSTICA (ε=−1, assinatura indefinida/Minkowski): a TORRE SPLIT — split-complex (números duais), split-quaternion, split-octonion. É EXATAMENTE a família hipercomplexa ×ₛ do projeto (o s é o ε): a mesma construção de Cayley-Dickson, só muda o sinal. Clássica = a métrica euclidiana (cristal); relativística = a de Lorentz (o corpo dual, tropical↔glacial). As duas torres são o mesmo padrão, símbolos distintos. [I] formalização.","epistemico":"inferencia","exemplos":[],"id":"torre_hurwitz_classica_relativistica","memoria":"semantica","meta":{"autor":"Lopes/broca-so","dominio":"algebra","fonte":"a_torre_de_hurwitzcayley--dickson; cayley_dickson; torre_split (ε=−1); algebras_hipercomplexas (×ₛ)"},"origem":"wiki","palavras_chave":[],"sinonimos":[],"tipo":"conceito","titulo":"A torre de Hurwitz clássica e relativística (a mesma das hipercomplexas)"}
\section{A torre de Hurwitz clássica e relativística (a mesma das hipercomplexas)}
Formalização: a torre de Hurwitz/Cayley-Dickson tem DOIS ramos, pelo sinal \ensuremath{\varepsilon} da construção. CLÁSSICA (\ensuremath{\varepsilon}=+1, positiva-definida): \ensuremath{\mathbb{R}}\ensuremath{\to}\ensuremath{\mathbb{C}}\ensuremath{\to}\ensuremath{\mathbb{H}}\ensuremath{\to}\ensuremath{\mathbb{O}} (as álgebras normadas de divisão, dim 1,2,4,8 — o teorema de Hurwitz). RELATIVÍSTICA (\ensuremath{\varepsilon}=\ensuremath{-}1, assinatura indefinida/Minkowski): a TORRE SPLIT — split-complex (números duais), split-quaternion, split-octonion. É EXATAMENTE a família hipercomplexa $\times$\ensuremath{_{s}} do projeto (o s é o \ensuremath{\varepsilon}): a mesma construção de Cayley-Dickson, só muda o sinal. Clássica = a métrica euclidiana (cristal); relativística = a de Lorentz (o corpo dual, tropical\ensuremath{\leftrightarrow}glacial). As duas torres são o mesmo padrão, símbolos distintos. [I] formalização.
\prov{relações: usa a\_torre\_de\_hurwitzcayley--dickson; usa cayley\_dickson; usa torre\_split; usa algebras\_hipercomplexas; relaciona numeros\_duais; parte\_de tiffany}
\prov{proveniência: wiki \textperiodcentered{} a\_torre\_de\_hurwitzcayley--dickson; cayley\_dickson; torre\_split (\ensuremath{\varepsilon}=\ensuremath{-}1); algebras\_hipercomplexas ($\times$\ensuremath{_{s}}) \textperiodcentered{} inferencia \textperiodcentered{} confiança 1.0 \textperiodcentered{} domínio algebra \textperiodcentered{} história 42 versões no jornal}

%CRISTAL {"arestas":[{"alvo":"cayley_dickson","confianca":1.0,"epistemico":"fato","origem":"paper","rel":"especializa"},{"alvo":"word","confianca":0.5,"epistemico":"fato","origem":"wiki","rel":"relaciona"}],"confianca":1.0,"contraexemplos":[],"descricao":"Segunda torre de Cayley-Dickson com parâmetro ε=−1: split-complex, split-quaternions, split-octonions; a primeira unidade satisfaz e₁²=+1 (hiperbólica).","epistemico":"fato","exemplos":[],"id":"torre_split","memoria":"semantica","meta":{"autor":"Gentil Lopes","descoberta":true,"dominio":"matematica","fonte":"paper_51_torre_hiperbolica.pdf"},"origem":"paper","palavras_chave":[],"sinonimos":[],"tipo":"conceito","titulo":"Torre Split (ε=−1)"}
\section{Torre Split (\ensuremath{\varepsilon}=\ensuremath{-}1)}
Segunda torre de Cayley-Dickson com parâmetro \ensuremath{\varepsilon}=\ensuremath{-}1: split-complex, split-quaternions, split-octonions; a primeira unidade satisfaz e\ensuremath{_{1}}\ensuremath{^{2}}=+1 (hiperbólica).
\prov{relações: especializa cayley\_dickson; relaciona word}
\prov{proveniência: paper \textperiodcentered{} paper\_51\_torre\_hiperbolica.pdf \textperiodcentered{} fato \textperiodcentered{} confiança 1.0 \textperiodcentered{} domínio matematica \textperiodcentered{} história 3 versões no jornal}

%CRISTAL {"arestas":[{"alvo":"algebra_dupolinomial_espaco","confianca":1.0,"epistemico":"fato","origem":"wiki","rel":"parte_de"},{"alvo":"torre_hurwitz_classica_relativistica","confianca":1.0,"epistemico":"fato","origem":"wiki","rel":"parte_de"},{"alvo":"torres_fisica_por_nivel","confianca":1.0,"epistemico":"fato","origem":"wiki","rel":"parte_de"},{"alvo":"meta_inducao_geral_simbolo","confianca":1.0,"epistemico":"fato","origem":"wiki","rel":"parte_de"},{"alvo":"algebra_posicional_unificada","confianca":1.0,"epistemico":"fato","origem":"wiki","rel":"usa"},{"alvo":"algebra_de_peirce","confianca":1.0,"epistemico":"fato","origem":"wiki","rel":"usa"},{"alvo":"meta_inducao_fechamento_geral","confianca":1.0,"epistemico":"fato","origem":"wiki","rel":"usa"},{"alvo":"transformada_metalica","confianca":1.0,"epistemico":"fato","origem":"wiki","rel":"generaliza"},{"alvo":"tiffany","confianca":1.0,"epistemico":"fato","origem":"wiki","rel":"parte_de"}],"confianca":1.0,"contraexemplos":[],"descricao":"O nó-hub que reúne a ferramenta inteira: (1) a ÁLGEBRA DUPOLINOMIAL sobre o espaço 𝔻 de dupolinômios (retas primas Pisot=cristal / Salem=borda / caos, cunhadas pela posição em 𝔽₂¹²⁸), que generaliza a posicional (metálica grau 2, plástica grau 3, todo grau — fechamento geral); (2) a TORRE DE HURWITZ clássica (ℝℂℍ𝕆, ε=+1) e relativística (split, ε=−1 = as hipercomplexas ×ₛ), o mesmo Cayley-Dickson com sinal; (3) a FÍSICA por nível (mecânica/geometria/termodinâmica, como o BAI); (4) a META-INDUÇÃO GERAL: fixado o padrão, só o SÍMBOLO muda, e a ÁLGEBRA DE PEIRCE unifica os símbolos (semiótica + o grafo/produto de Peirce). É a Universidade Tiffany como uma estrutura só: uma reta prima por dimensão, a mesma parábola classificando (cristal/borda/caos), a mesma cunhagem posicional, a mesma física, o mesmo símbolo peirceano. [I] o edifício unificado.","epistemico":"inferencia","exemplos":[],"id":"torre_unificada","memoria":"semantica","meta":{"autor":"Lopes/broca-so","dominio":"algebra","fonte":"doc/algebra_prata/paper (§8–§11); as torres + a álgebra dupolinomial + Peirce + a física por nível"},"origem":"wiki","palavras_chave":[],"sinonimos":[],"tipo":"conceito","titulo":"A TORRE UNIFICADA — o edifício num princípio só (símbolo varia, Peirce unifica)"}
\section{A TORRE UNIFICADA — o edifício num princípio só (símbolo varia, Peirce unifica)}
O nó-hub que reúne a ferramenta inteira: (1) a ÁLGEBRA DUPOLINOMIAL sobre o espaço \ensuremath{\mathbb{D}} de dupolinômios (retas primas Pisot=cristal / Salem=borda / caos, cunhadas pela posição em \ensuremath{\mathbb{F}}\ensuremath{_{2}}\ensuremath{^{128}}), que generaliza a posicional (metálica grau 2, plástica grau 3, todo grau — fechamento geral); (2) a TORRE DE HURWITZ clássica (\ensuremath{\mathbb{R}}\ensuremath{\mathbb{C}}\ensuremath{\mathbb{H}}\ensuremath{\mathbb{O}}, \ensuremath{\varepsilon}=+1) e relativística (split, \ensuremath{\varepsilon}=\ensuremath{-}1 = as hipercomplexas $\times$\ensuremath{_{s}}), o mesmo Cayley-Dickson com sinal; (3) a FÍSICA por nível (mecânica/geometria/termodinâmica, como o BAI); (4) a META-INDUÇÃO GERAL: fixado o padrão, só o SÍMBOLO muda, e a ÁLGEBRA DE PEIRCE unifica os símbolos (semiótica + o grafo/produto de Peirce). É a Universidade Tiffany como uma estrutura só: uma reta prima por dimensão, a mesma parábola classificando (cristal/borda/caos), a mesma cunhagem posicional, a mesma física, o mesmo símbolo peirceano. [I] o edifício unificado.
\prov{relações: parte\_de algebra\_dupolinomial\_espaco; parte\_de torre\_hurwitz\_classica\_relativistica; parte\_de torres\_fisica\_por\_nivel; parte\_de meta\_inducao\_geral\_simbolo; usa algebra\_posicional\_unificada; usa algebra\_de\_peirce; usa meta\_inducao\_fechamento\_geral; generaliza transformada\_metalica; parte\_de tiffany}
\prov{proveniência: wiki \textperiodcentered{} doc/algebra\_prata/paper (§8–§11); as torres + a álgebra dupolinomial + Peirce + a física por nível \textperiodcentered{} inferencia \textperiodcentered{} confiança 1.0 \textperiodcentered{} domínio algebra \textperiodcentered{} história 38 versões no jornal}

%CRISTAL {"arestas":[{"alvo":"bai_orbita_energia_por_trabalho","confianca":1.0,"epistemico":"fato","origem":"wiki","rel":"usa"},{"alvo":"cofre_parabola_hero","confianca":1.0,"epistemico":"fato","origem":"wiki","rel":"usa"},{"alvo":"corpo_dual","confianca":1.0,"epistemico":"fato","origem":"wiki","rel":"usa"},{"alvo":"torre_hurwitz_classica_relativistica","confianca":1.0,"epistemico":"fato","origem":"wiki","rel":"relaciona"},{"alvo":"algebra_dupolinomial_espaco","confianca":1.0,"epistemico":"fato","origem":"wiki","rel":"relaciona"},{"alvo":"tiffany","confianca":1.0,"epistemico":"fato","origem":"wiki","rel":"parte_de"}],"confianca":1.0,"contraexemplos":[],"descricao":"Cada nível/dimensão (cada reta dupolinomial, cada degrau de Hurwitz) não é só álgebra — carrega uma FÍSICA completa, como o BAI: MECÂNICA (a órbita 𝒞_δ e a energia δ por tipo de trabalho — o gerador σ dá o raio e o tempo de vida); GEOMETRIA (a métrica/curvatura — dS₄ real dissipativo, a reta como geodésica de passo h=ln σ); TERMODINÂMICA (o calor a=|associador|, a entropia ln σ, o horizonte, a parábola a+c²=1). Cristal (Pisot) = frio/associativo; Salem = a borda crítica; caos = quente. A mesma tríade mecânica-geometria-termodinâmica se repete em cada nível de cada torre. [I] a física acompanha a álgebra em todo nível.","epistemico":"inferencia","exemplos":[],"id":"torres_fisica_por_nivel","memoria":"semantica","meta":{"autor":"Lopes/broca-so","dominio":"algebra","fonte":"bai_orbita_energia_por_trabalho; termodinamica; cofre_parabola_hero; corpo_dual"},"origem":"wiki","palavras_chave":[],"sinonimos":[],"tipo":"conceito","titulo":"Cada nível das torres tem mecânica, geometria e termodinâmica (como o BAI)"}
\section{Cada nível das torres tem mecânica, geometria e termodinâmica (como o BAI)}
Cada nível/dimensão (cada reta dupolinomial, cada degrau de Hurwitz) não é só álgebra — carrega uma FÍSICA completa, como o BAI: MECÂNICA (a órbita \ensuremath{\mathcal{C}}\_\ensuremath{\delta} e a energia \ensuremath{\delta} por tipo de trabalho — o gerador \ensuremath{\sigma} dá o raio e o tempo de vida); GEOMETRIA (a métrica/curvatura — dS\ensuremath{_{4}} real dissipativo, a reta como geodésica de passo h=ln \ensuremath{\sigma}); TERMODINÂMICA (o calor a=|associador|, a entropia ln \ensuremath{\sigma}, o horizonte, a parábola a+c\ensuremath{^{2}}=1). Cristal (Pisot) = frio/associativo; Salem = a borda crítica; caos = quente. A mesma tríade mecânica-geometria-termodinâmica se repete em cada nível de cada torre. [I] a física acompanha a álgebra em todo nível.
\prov{relações: usa bai\_orbita\_energia\_por\_trabalho; usa cofre\_parabola\_hero; usa corpo\_dual; relaciona torre\_hurwitz\_classica\_relativistica; relaciona algebra\_dupolinomial\_espaco; parte\_de tiffany}
\prov{proveniência: wiki \textperiodcentered{} bai\_orbita\_energia\_por\_trabalho; termodinamica; cofre\_parabola\_hero; corpo\_dual \textperiodcentered{} inferencia \textperiodcentered{} confiança 1.0 \textperiodcentered{} domínio algebra \textperiodcentered{} história 14 versões no jornal}

%CRISTAL {"arestas":[{"alvo":"cascata_multiescala","confianca":1.0,"epistemico":"fato","origem":"paper","rel":"usa"}],"confianca":1.0,"contraexemplos":[],"descricao":"A curva ν₀·4Ns−¹=C no plano (ν₀,Ns) mantém a escala de Kolmogorov fixa enquanto alarga o intervalo inercial — a discretização do limite inviscido onde Onsager (1949) previu dissipação anômala.","epistemico":"fato","exemplos":[],"id":"trajetoria_onsager","memoria":"semantica","meta":{"autor":"Gentil Lopes","descoberta":false,"dominio":"matematica","fonte":"paper_BG.pdf, Def.1"},"origem":"paper","palavras_chave":[],"sinonimos":[],"tipo":"conceito","titulo":"Trajetória de Onsager"}
\section{Trajetória de Onsager}
A curva \ensuremath{\nu}\ensuremath{_{0}}\textperiodcentered 4Ns\ensuremath{-}\ensuremath{^{1}}=C no plano (\ensuremath{\nu}\ensuremath{_{0}},Ns) mantém a escala de Kolmogorov fixa enquanto alarga o intervalo inercial — a discretização do limite inviscido onde Onsager (1949) previu dissipação anômala.
\prov{relações: usa cascata\_multiescala}
\prov{proveniência: paper \textperiodcentered{} paper\_BG.pdf, Def.1 \textperiodcentered{} fato \textperiodcentered{} confiança 1.0 \textperiodcentered{} domínio matematica \textperiodcentered{} história 3 versões no jornal}

%CRISTAL {"arestas":[{"alvo":"metrica_produto","confianca":1.0,"epistemico":"fato","origem":"paper","rel":"usa"}],"confianca":1.0,"contraexemplos":[],"descricao":"A série de Fourier do seno lemniscático sl(z) é uma convolução: onda circular (freq 2n+1) × peso hiperbólico (sech). Espectro = métrica-produto no domínio dual.","epistemico":"fato","exemplos":[],"id":"transformada_lemniscatica","memoria":"semantica","meta":{"autor":"Gentil Lopes","descoberta":true,"dominio":"matematica","fonte":"paper_30_transformada_lemniscatica.pdf, Teo.2.2"},"origem":"paper","palavras_chave":[],"sinonimos":[],"tipo":"conceito","titulo":"Transformada Lemniscática"}
\section{Transformada Lemniscática}
A série de Fourier do seno lemniscático sl(z) é uma convolução: onda circular (freq 2n+1) $\times$ peso hiperbólico (sech). Espectro = métrica-produto no domínio dual.
\prov{relações: usa metrica\_produto}
\prov{proveniência: paper \textperiodcentered{} paper\_30\_transformada\_lemniscatica.pdf, Teo.2.2 \textperiodcentered{} fato \textperiodcentered{} confiança 1.0 \textperiodcentered{} domínio matematica \textperiodcentered{} história 2 versões no jornal}

%CRISTAL {"arestas":[{"alvo":"beta_numeracao_metalica","confianca":1.0,"epistemico":"fato","origem":"wiki","rel":"parte_de"},{"alvo":"conjugados_galois_parabola","confianca":1.0,"epistemico":"fato","origem":"wiki","rel":"parte_de"},{"alvo":"algebra_de_peirce","confianca":1.0,"epistemico":"fato","origem":"wiki","rel":"usa"},{"alvo":"algebra_dupolinomial_espaco","confianca":1.0,"epistemico":"fato","origem":"wiki","rel":"usa"},{"alvo":"conservacao_parabola_costura","confianca":1.0,"epistemico":"fato","origem":"wiki","rel":"usa"},{"alvo":"simetria_por_reta_corpo","confianca":1.0,"epistemico":"fato","origem":"wiki","rel":"usa"},{"alvo":"torre_unificada","confianca":1.0,"epistemico":"fato","origem":"wiki","rel":"parte_de"},{"alvo":"transformada_metalica_fechada_perelman","confianca":1.0,"epistemico":"fato","origem":"wiki","rel":"usa"},{"alvo":"tiffany","confianca":1.0,"epistemico":"fato","origem":"wiki","rel":"parte_de"},{"alvo":"laco_traducao_bai_peirce","confianca":1.0,"epistemico":"fato","origem":"wiki","rel":"parte_de"},{"alvo":"teorema_geral_traducao","confianca":1.0,"epistemico":"fato","origem":"wiki","rel":"explica"},{"alvo":"bases_pisot","confianca":1.0,"epistemico":"fato","origem":"wiki","rel":"especializa"}],"confianca":1.0,"contraexemplos":[],"descricao":"O nó unificador: a álgebra de Peirce + a posicional metálica + a parábola-conservação + a costura por dobra fecham numa TRANSFORMADA — e ela já existe na matemática: a β-NUMERAÇÃO DE PISOT (Ostrowski/Zeckendorf) na base σ_n. As peças do projeto são as partes dela: (1) a MOEDA = POSIÇÃO k são os DÍGITOS da β-expansão (os coeficientes da transformada); (2) σ_n é a BASE β (Pisot); (3) a DOBRA do corpo (Lucas/Pell, [[retas_dobra_lucas_pell]]) é a FRACTAL DE RAUZY / a janela cut-and-project = o quasicristal; (4) σ·conj = −1 (a norma/det) é a LEI DE CONSERVAÇÃO; (5) a PARÁBOLA a+c²=1 são as DUAS PROJEÇÕES DE GALOIS (σ expande=lado-a; −1/σ contrai=cristal/c²); (6) o PRODUTO DE PEIRCE `:` é a SUBSTITUIÇÃO/relação (o grafo como sistema dinâmico da numeração). Tudo do MESMO par conjugado. RESSALVAS honestas: NÃO é Fourier (as retas são corpos quadráticos distintos, não ortogonais — a completude vem da Property F de Pisot, não de ortogonalidade); os dígitos são 0..n (só o ouro/Fibonacci é binário/Zeckendorf puro); a ponte Peirce-relacional ↔ transformada-linear (módulo booleano → módulo sobre semianel) falta FORMALIZAR pra ser teorema, não só analogia. [I] a transformada como β-numeração metálica.","epistemico":"inferencia","exemplos":[],"id":"transformada_metalica","memoria":"semantica","meta":{"autor":"Lopes/broca-so","dominio":"algebra","fonte":"β-numeração Pisot (Ostrowski/Zeckendorf); Rauzy fractal; [[algebra_de_peirce]]; [[conservacao_parabola_costura]]; verificação: σ·conj=−1, Pisot=cristal"},"origem":"wiki","palavras_chave":[],"sinonimos":[],"tipo":"conceito","titulo":"A TRANSFORMADA METÁLICA — a unificação (β-numeração Pisot + Rauzy + parábola + Peirce)"}
\section{A TRANSFORMADA METÁLICA — a unificação (\ensuremath{\beta}-numeração Pisot + Rauzy + parábola + Peirce)}
O nó unificador: a álgebra de Peirce + a posicional metálica + a parábola-conservação + a costura por dobra fecham numa TRANSFORMADA — e ela já existe na matemática: a \ensuremath{\beta}-NUMERAÇÃO DE PISOT (Ostrowski/Zeckendorf) na base \ensuremath{\sigma}\_n. As peças do projeto são as partes dela: (1) a MOEDA = POSIÇÃO k são os DÍGITOS da \ensuremath{\beta}-expansão (os coeficientes da transformada); (2) \ensuremath{\sigma}\_n é a BASE \ensuremath{\beta} (Pisot); (3) a DOBRA do corpo (Lucas/Pell, [[retas\_dobra\_lucas\_pell]]) é a FRACTAL DE RAUZY / a janela cut-and-project = o quasicristal; (4) \ensuremath{\sigma}\textperiodcentered conj = \ensuremath{-}1 (a norma/det) é a LEI DE CONSERVAÇÃO; (5) a PARÁBOLA a+c\ensuremath{^{2}}=1 são as DUAS PROJEÇÕES DE GALOIS (\ensuremath{\sigma} expande=lado-a; \ensuremath{-}1/\ensuremath{\sigma} contrai=cristal/c\ensuremath{^{2}}); (6) o PRODUTO DE PEIRCE `:` é a SUBSTITUIÇÃO/relação (o grafo como sistema dinâmico da numeração). Tudo do MESMO par conjugado. RESSALVAS honestas: NÃO é Fourier (as retas são corpos quadráticos distintos, não ortogonais — a completude vem da Property F de Pisot, não de ortogonalidade); os dígitos são 0..n (só o ouro/Fibonacci é binário/Zeckendorf puro); a ponte Peirce-relacional \ensuremath{\leftrightarrow} transformada-linear (módulo booleano \ensuremath{\to} módulo sobre semianel) falta FORMALIZAR pra ser teorema, não só analogia. [I] a transformada como \ensuremath{\beta}-numeração metálica.
\prov{relações: parte\_de beta\_numeracao\_metalica; parte\_de conjugados\_galois\_parabola; usa algebra\_de\_peirce; usa algebra\_dupolinomial\_espaco; usa conservacao\_parabola\_costura; usa simetria\_por\_reta\_corpo; parte\_de torre\_unificada; usa transformada\_metalica\_fechada\_perelman; parte\_de tiffany; parte\_de laco\_traducao\_bai\_peirce; explica teorema\_geral\_traducao; especializa bases\_pisot}
\prov{proveniência: wiki \textperiodcentered{} \ensuremath{\beta}-numeração Pisot (Ostrowski/Zeckendorf); Rauzy fractal; [[algebra\_de\_peirce]]; [[conservacao\_parabola\_costura]]; verificação: \ensuremath{\sigma}\textperiodcentered conj=\ensuremath{-}1, Pisot=cristal \textperiodcentered{} inferencia \textperiodcentered{} confiança 1.0 \textperiodcentered{} domínio algebra \textperiodcentered{} história 55 versões no jornal}

%CRISTAL {"arestas":[{"alvo":"transformada_metalica","confianca":1.0,"epistemico":"fato","origem":"wiki","rel":"prova"},{"alvo":"beta_numeracao_metalica","confianca":1.0,"epistemico":"fato","origem":"wiki","rel":"usa"},{"alvo":"conjugados_galois_parabola","confianca":1.0,"epistemico":"fato","origem":"wiki","rel":"usa"},{"alvo":"recuperacao_evidencia_invariante","confianca":1.0,"epistemico":"fato","origem":"wiki","rel":"relaciona"},{"alvo":"conservacao_parabola_costura","confianca":1.0,"epistemico":"fato","origem":"wiki","rel":"usa"},{"alvo":"numeros_pisot","confianca":1.0,"epistemico":"fato","origem":"wiki","rel":"usa"},{"alvo":"tiffany","confianca":1.0,"epistemico":"fato","origem":"wiki","rel":"parte_de"}],"confianca":1.0,"contraexemplos":[],"descricao":"O fechamento da transformada, pelo MESMO método de Perelman da recuperação (ver [[recuperacao_evidencia_invariante]]): a transformada β-numeração fica COMPLETA/convergente porque a sua ENTROPIA é monótona sobre o invariante. A cada passo da β-expansão, a=remanescente/N (o lado dissipativo) CAI monótona e c²=representado/N SOBE, com a+c²=1 SEMPRE conservado; fecha (a→0) em passos finitos. Verificado (ouro/prata/cobre, N~10⁶): 6–7 passos, a↓ monótona, a+c²=1 em todo passo, rem=0. O motor é a CONJUGADA que contrai (Pisot, |−1/σ|<1): a entropia decai na taxa da conjugada (a mesma que dá o cristal). Logo: a monótona↓ + limitada (a≥0, o invariante) ⇒ a expansão TERMINA (Property F de Pisot) ⇒ a transformada é FECHADA. É o análogo do 'no local collapsing' de Perelman: a entropia monótona impede o colapso/loop e garante a convergência. O mesmo método fecha os dois lados: a recuperação (a evidência come a entropia, aterra) E a transformada (a β-expansão come a entropia, fecha) — a entropia de Perelman na parábola, sem invadir. Resolve a ressalva de completude (antes 'vem da Property F'; agora PROVADA pela monotonicidade). [D] a transformada fechada por Perelman.","epistemico":"inferencia","exemplos":[],"id":"transformada_metalica_fechada_perelman","memoria":"semantica","meta":{"autor":"Lopes/broca-so","dominio":"algebra","fonte":"verificação: β-expansão metálica, a↓ monótona a+c²=1 finita (ouro/prata/cobre); Pisot Property F; método de [[recuperacao_evidencia_invariante]]"},"origem":"wiki","palavras_chave":[],"sinonimos":[],"tipo":"conceito","titulo":"A transformada metálica FECHA pelo método de Perelman (entropia monótona no invariante)"}
\section{A transformada metálica FECHA pelo método de Perelman (entropia monótona no invariante)}
O fechamento da transformada, pelo MESMO método de Perelman da recuperação (ver [[recuperacao\_evidencia\_invariante]]): a transformada \ensuremath{\beta}-numeração fica COMPLETA/convergente porque a sua ENTROPIA é monótona sobre o invariante. A cada passo da \ensuremath{\beta}-expansão, a=remanescente/N (o lado dissipativo) CAI monótona e c\ensuremath{^{2}}=representado/N SOBE, com a+c\ensuremath{^{2}}=1 SEMPRE conservado; fecha (a\ensuremath{\to}0) em passos finitos. Verificado (ouro/prata/cobre, N\~{}10\ensuremath{^{6}}): 6–7 passos, a\ensuremath{\downarrow} monótona, a+c\ensuremath{^{2}}=1 em todo passo, rem=0. O motor é a CONJUGADA que contrai (Pisot, |\ensuremath{-}1/\ensuremath{\sigma}|<1): a entropia decai na taxa da conjugada (a mesma que dá o cristal). Logo: a monótona\ensuremath{\downarrow} + limitada (a\ensuremath{\ge}0, o invariante) \ensuremath{\Rightarrow} a expansão TERMINA (Property F de Pisot) \ensuremath{\Rightarrow} a transformada é FECHADA. É o análogo do 'no local collapsing' de Perelman: a entropia monótona impede o colapso/loop e garante a convergência. O mesmo método fecha os dois lados: a recuperação (a evidência come a entropia, aterra) E a transformada (a \ensuremath{\beta}-expansão come a entropia, fecha) — a entropia de Perelman na parábola, sem invadir. Resolve a ressalva de completude (antes 'vem da Property F'; agora PROVADA pela monotonicidade). [D] a transformada fechada por Perelman.
\prov{relações: prova transformada\_metalica; usa beta\_numeracao\_metalica; usa conjugados\_galois\_parabola; relaciona recuperacao\_evidencia\_invariante; usa conservacao\_parabola\_costura; usa numeros\_pisot; parte\_de tiffany}
\prov{proveniência: wiki \textperiodcentered{} verificação: \ensuremath{\beta}-expansão metálica, a\ensuremath{\downarrow} monótona a+c\ensuremath{^{2}}=1 finita (ouro/prata/cobre); Pisot Property F; método de [[recuperacao\_evidencia\_invariante]] \textperiodcentered{} inferencia \textperiodcentered{} confiança 1.0 \textperiodcentered{} domínio algebra \textperiodcentered{} história 33 versões no jornal}

%CRISTAL {"arestas":[{"alvo":"algebra_reta_mineral_7d","confianca":1.0,"epistemico":"fato","origem":"wiki","rel":"usa"},{"alvo":"transformada_metalica","confianca":1.0,"epistemico":"fato","origem":"wiki","rel":"especializa"},{"alvo":"cifra_metais_an","confianca":1.0,"epistemico":"fato","origem":"wiki","rel":"relaciona"},{"alvo":"invariante_estelar","confianca":1.0,"epistemico":"fato","origem":"wiki","rel":"usa"},{"alvo":"corpo_mineral","confianca":1.0,"epistemico":"fato","origem":"wiki","rel":"parte_de"},{"alvo":"molecula_mineral","confianca":1.0,"epistemico":"fato","origem":"wiki","rel":"usa"}],"confianca":1.0,"contraexemplos":[],"descricao":"A transformada metálica (β-numeração 1D, base σ_n) sobe à álgebra octoniônica 𝕆 em D=7: cada posição da órbita carrega um OCTÔNION (8 componentes, densidade 8×), e a informação cavalga a órbita. Codifica c_k = u^k ⊗ m_k (girada pelo portador unitário u, a fase metálica em 7D); decodifica m_k = conj(u^k) ⊗ c_k (𝕆 é alternativa ⟹ x⁻¹(xy)=y): REVERSÍVEL, roundtrip exato. A norma conservada (‖c_k‖=‖m_k‖) é o CHECKSUM — a conservação a+c²=1. Serve para codificar informação direto nas órbitas.","epistemico":"inferencia","exemplos":[],"id":"transformada_mineral_7d","memoria":"semantica","meta":{"dominio":"matematica","fonte":"transformada mineral 7D"},"origem":"wiki","palavras_chave":[],"sinonimos":[],"tipo":"conceito","titulo":"A Transformada Mineral em D=7 (informação na órbita)"}
\section{A Transformada Mineral em D=7 (informação na órbita)}
A transformada metálica (\ensuremath{\beta}-numeração 1D, base \ensuremath{\sigma}\_n) sobe à álgebra octoniônica \ensuremath{\mathbb{O}} em D=7: cada posição da órbita carrega um OCTÔNION (8 componentes, densidade 8$\times$), e a informação cavalga a órbita. Codifica c\_k = u\^{}k \ensuremath{\otimes} m\_k (girada pelo portador unitário u, a fase metálica em 7D); decodifica m\_k = conj(u\^{}k) \ensuremath{\otimes} c\_k (\ensuremath{\mathbb{O}} é alternativa \ensuremath{\Longrightarrow} x\ensuremath{^{-1}}(xy)=y): REVERSÍVEL, roundtrip exato. A norma conservada (\ensuremath{\|}c\_k\ensuremath{\|}=\ensuremath{\|}m\_k\ensuremath{\|}) é o CHECKSUM — a conservação a+c\ensuremath{^{2}}=1. Serve para codificar informação direto nas órbitas.
\prov{relações: usa algebra\_reta\_mineral\_7d; especializa transformada\_metalica; relaciona cifra\_metais\_an; usa invariante\_estelar; parte\_de corpo\_mineral; usa molecula\_mineral}
\prov{proveniência: wiki \textperiodcentered{} transformada mineral 7D \textperiodcentered{} inferencia \textperiodcentered{} confiança 1.0 \textperiodcentered{} domínio matematica \textperiodcentered{} história 7 versões no jornal}

%CRISTAL {"arestas":[{"alvo":"entropia_topologica","confianca":1.0,"epistemico":"fato","origem":"paper","rel":"usa"},{"alvo":"word","confianca":0.5,"epistemico":"fato","origem":"wiki","rel":"relaciona"}],"confianca":1.0,"contraexemplos":[],"descricao":"Em λ=1 o número de caminhos passa de crescimento exponencial (λ>1, entropia positiva) a polinomial (λ=1, entropia nula): descontinuidade na ordem do crescimento — a 'borda do caos'.","epistemico":"fato","exemplos":[],"id":"transicao_fase_lambda1","memoria":"semantica","meta":{"autor":"Gentil Lopes","descoberta":false,"dominio":"matematica","fonte":"paper_47_termodinamica_informacao.pdf, Obs.3.1"},"origem":"paper","palavras_chave":[],"sinonimos":[],"tipo":"conceito","titulo":"Transição de Fase λ=1"}
\section{Transição de Fase \ensuremath{\lambda}=1}
Em \ensuremath{\lambda}=1 o número de caminhos passa de crescimento exponencial (\ensuremath{\lambda}>1, entropia positiva) a polinomial (\ensuremath{\lambda}=1, entropia nula): descontinuidade na ordem do crescimento — a 'borda do caos'.
\prov{relações: usa entropia\_topologica; relaciona word}
\prov{proveniência: paper \textperiodcentered{} paper\_47\_termodinamica\_informacao.pdf, Obs.3.1 \textperiodcentered{} fato \textperiodcentered{} confiança 1.0 \textperiodcentered{} domínio matematica \textperiodcentered{} história 3 versões no jornal}

%CRISTAL {"arestas":[{"alvo":"virtualizacao","confianca":0.5,"epistemico":"fato","origem":"wiki","rel":"relaciona"},{"alvo":"word","confianca":0.5,"epistemico":"fato","origem":"wiki","rel":"relaciona"},{"alvo":"vontade_de_poder","confianca":0.5,"epistemico":"fato","origem":"wiki","rel":"relaciona"}],"confianca":1.0,"contraexemplos":[],"descricao":"Coeficientes C(n,k) do binômio; cada entrada é a soma das duas acima.","epistemico":"fato","exemplos":[],"id":"triangulo_de_pascal","memoria":"semantica","meta":{"autor":"Gentil Lopes","descoberta":false,"dominio":"matematica","fonte":""},"origem":"manual","palavras_chave":[],"sinonimos":["binômio","coeficiente binomial"],"tipo":"conceito","titulo":"Triângulo de Pascal / coeficientes binomiais"}
\section{Triângulo de Pascal / coeficientes binomiais}
Coeficientes C(n,k) do binômio; cada entrada é a soma das duas acima.
\prov{sinónimos: binômio; coeficiente binomial}
\prov{relações: relaciona virtualizacao; relaciona word; relaciona vontade\_de\_poder}
\prov{proveniência: manual \textperiodcentered{} fato \textperiodcentered{} confiança 1.0 \textperiodcentered{} domínio matematica \textperiodcentered{} história 5 versões no jornal}

%CRISTAL {"arestas":[{"alvo":"sigma_equivalencia_crown","confianca":1.0,"epistemico":"fato","origem":"wiki","rel":"relaciona"},{"alvo":"mathieu_m24","confianca":1.0,"epistemico":"fato","origem":"wiki","rel":"relaciona"}],"confianca":1.0,"contraexemplos":[],"descricao":"O primeiro coeficiente de Fourier não-trivial do monstrous moonshine, 196884, faz ponte entre a Muralha PSL e o 'Observador' 169=13²: 196884 ≡ 168 (mod 169) e 196884 ≡ 156 = φ(169) (mod 168). Três cantos, um triângulo, uma verdade excepcional.","epistemico":"fato","exemplos":[],"id":"triangulo_monster_observador","memoria":"semantica","meta":{"autor":"Lord Index (Shawn R. Schiller)","dominio":"matematica","fonte":"A Catedral — Part VIII-A (jul 2026), ~/Imagens/catedral"},"origem":"wiki","palavras_chave":[],"sinonimos":[],"tipo":"conceito","titulo":"O Triângulo Modular: Monster–PSL–Observador"}
\section{O Triângulo Modular: Monster–PSL–Observador}
O primeiro coeficiente de Fourier não-trivial do monstrous moonshine, 196884, faz ponte entre a Muralha PSL e o 'Observador' 169=13\ensuremath{^{2}}: 196884 \ensuremath{\equiv} 168 (mod 169) e 196884 \ensuremath{\equiv} 156 = \ensuremath{\varphi}(169) (mod 168). Três cantos, um triângulo, uma verdade excepcional.
\prov{relações: relaciona sigma\_equivalencia\_crown; relaciona mathieu\_m24}
\prov{proveniência: wiki \textperiodcentered{} A Catedral — Part VIII-A (jul 2026), \~{}/Imagens/catedral \textperiodcentered{} fato \textperiodcentered{} confiança 1.0 \textperiodcentered{} domínio matematica \textperiodcentered{} história 3 versões no jornal}

%CRISTAL {"arestas":[{"alvo":"numeros_duais","confianca":1.0,"epistemico":"fato","origem":"paper","rel":"explica"},{"alvo":"geometria","confianca":1.0,"epistemico":"fato","origem":"paper","rel":"parte_de"}],"confianca":1.0,"contraexemplos":[],"descricao":"Elementos de SL₂(ℤ) classificam-se pelo traço — elíptico (|tr|<2), hiperbólico (|tr|>2), parabólico (|tr|=2) — coincidindo tipo a tipo com as álgebras de dim 2: ℂ, split, dual.","epistemico":"fato","exemplos":[],"id":"tricotomia_modular","memoria":"semantica","meta":{"autor":"Gentil Lopes","descoberta":true,"dominio":"matematica","fonte":"paper_27_ponte_modular.pdf, Teo.4.1"},"origem":"paper","palavras_chave":[],"sinonimos":[],"tipo":"conceito","titulo":"Tricotomia Modular"}
\section{Tricotomia Modular}
Elementos de SL\ensuremath{_{2}}(\ensuremath{\mathbb{Z}}) classificam-se pelo traço — elíptico (|tr|<2), hiperbólico (|tr|>2), parabólico (|tr|=2) — coincidindo tipo a tipo com as álgebras de dim 2: \ensuremath{\mathbb{C}}, split, dual.
\prov{relações: explica numeros\_duais; parte\_de geometria}
\prov{proveniência: paper \textperiodcentered{} paper\_27\_ponte\_modular.pdf, Teo.4.1 \textperiodcentered{} fato \textperiodcentered{} confiança 1.0 \textperiodcentered{} domínio matematica \textperiodcentered{} história 3 versões no jornal}

%CRISTAL {"arestas":[{"alvo":"shift_map_ns_turing","confianca":1.0,"epistemico":"fato","origem":"paper","rel":"especializa"},{"alvo":"modos_aproximacao_s2","confianca":1.0,"epistemico":"fato","origem":"paper","rel":"usa"}],"confianca":1.0,"contraexemplos":[],"descricao":"Conjectura: a máquina universal de Rogozhin (5 estados, 5 símbolos) implementável em NS-3D por anéis de vórtice — fita em orientações de S², transições por reconexão. Estende o shift binário do estágio 1. Reconexão contínua vs discreta em aberto.","epistemico":"hipotese","exemplos":[],"id":"turing_rogozhin_ns","memoria":"semantica","meta":{"autor":"Gentil Lopes","descoberta":false,"dominio":"matematica","fonte":"paper_AR.pdf, Teo.1"},"origem":"paper","palavras_chave":[],"sinonimos":[],"tipo":"conceito","titulo":"Turing Universal (Rogozhin) em NS"}
\section{Turing Universal (Rogozhin) em NS}
Conjectura: a máquina universal de Rogozhin (5 estados, 5 símbolos) implementável em NS-3D por anéis de vórtice — fita em orientações de S\ensuremath{^{2}}, transições por reconexão. Estende o shift binário do estágio 1. Reconexão contínua vs discreta em aberto.
\prov{relações: especializa shift\_map\_ns\_turing; usa modos\_aproximacao\_s2}
\prov{proveniência: paper \textperiodcentered{} paper\_AR.pdf, Teo.1 \textperiodcentered{} hipotese \textperiodcentered{} confiança 1.0 \textperiodcentered{} domínio matematica \textperiodcentered{} história 3 versões no jornal}

%CRISTAL {"arestas":[{"alvo":"tiffany","confianca":1.0,"epistemico":"fato","origem":"wiki","rel":"semelhante"},{"alvo":"catedral_sigma_aritmetica","confianca":1.0,"epistemico":"fato","origem":"wiki","rel":"parte_de"}],"confianca":1.0,"contraexemplos":[],"descricao":"A obra retrata o conhecimento como um ARQUIVO em grade (o 'Observatory Archive', camada 24, φ⁴, 27×37=999) — 'linguagem é arquitetura, sentido é luz, observação é criação'. Eco visual de um grafo de conceitos como a própria Universidade Tiffany.","epistemico":"inferencia","exemplos":[],"id":"uch_observatory_archive","memoria":"semantica","meta":{"autor":"Lord Index (Shawn R. Schiller)","dominio":"matematica","fonte":"A Catedral — Part VIII-A (jul 2026), ~/Imagens/catedral"},"origem":"wiki","palavras_chave":[],"sinonimos":[],"tipo":"conceito","titulo":"O Arquivo-Observatório (UCH): a grade de conceitos"}
\section{O Arquivo-Observatório (UCH): a grade de conceitos}
A obra retrata o conhecimento como um ARQUIVO em grade (o 'Observatory Archive', camada 24, \ensuremath{\varphi}\ensuremath{^{4}}, 27$\times$37=999) — 'linguagem é arquitetura, sentido é luz, observação é criação'. Eco visual de um grafo de conceitos como a própria Universidade Tiffany.
\prov{relações: semelhante tiffany; parte\_de catedral\_sigma\_aritmetica}
\prov{proveniência: wiki \textperiodcentered{} A Catedral — Part VIII-A (jul 2026), \~{}/Imagens/catedral \textperiodcentered{} inferencia \textperiodcentered{} confiança 1.0 \textperiodcentered{} domínio matematica \textperiodcentered{} história 3 versões no jornal}

%CRISTAL {"arestas":[{"alvo":"sistema_numerico_b","confianca":1.0,"epistemico":"fato","origem":"paper","rel":"parte_de"},{"alvo":"word","confianca":0.5,"epistemico":"fato","origem":"wiki","rel":"relaciona"}],"confianca":1.0,"contraexemplos":[],"descricao":"Em B existe elemento onde −1·a≠−a — a propriedade que distingue B de todo outro sistema numérico, decorrente da não-distributividade.","epistemico":"fato","exemplos":[],"id":"unidade_hiperimaginaria","memoria":"semantica","meta":{"autor":"Gentil Lopes","descoberta":false,"dominio":"matematica","fonte":"corpus: NUMEROS_NAO_EUCLIDIANOS"},"origem":"paper","palavras_chave":[],"sinonimos":[],"tipo":"conceito","titulo":"Unidade Hiperimaginária (−1·j≠−j)"}
\section{Unidade Hiperimaginária (\ensuremath{-}1\textperiodcentered j\ensuremath{\ne}\ensuremath{-}j)}
Em B existe elemento onde \ensuremath{-}1\textperiodcentered a\ensuremath{\ne}\ensuremath{-}a — a propriedade que distingue B de todo outro sistema numérico, decorrente da não-distributividade.
\prov{relações: parte\_de sistema\_numerico\_b; relaciona word}
\prov{proveniência: paper \textperiodcentered{} corpus: NUMEROS\_NAO\_EUCLIDIANOS \textperiodcentered{} fato \textperiodcentered{} confiança 1.0 \textperiodcentered{} domínio matematica \textperiodcentered{} história 3 versões no jornal}

%CRISTAL {"arestas":[{"alvo":"regua_de_gentil","confianca":1.0,"epistemico":"fato","origem":"wiki","rel":"usa"},{"alvo":"cosmologia_quantica_hub","confianca":1.0,"epistemico":"fato","origem":"wiki","rel":"relaciona"},{"alvo":"cosmologia_quantica","confianca":1.0,"epistemico":"fato","origem":"wiki","rel":"relaciona"},{"alvo":"broca_os","confianca":1.0,"epistemico":"fato","origem":"wiki","rel":"parte_de"}],"confianca":1.0,"contraexemplos":[],"descricao":"A álgebra é o 'por-unidade' (pu) da física: como na engenharia de potência (escolhe-se uma base, tudo vira x_pu=x_real/x_base adimensional de ordem 1), a série I-IV trabalha em números puros — a base é a escala do meio E₀, o valor pu é o número da álgebra. Fixada E₀ (joules), cada dimensão converte por um fator (energia x→xE₀, massa x→xE₀/c², comprimento x→xℏc/E₀, tempo x→xℏ/E₀, temperatura x→xE₀/k_B, carga x→xe). Os números puros são invariantes de escala — o √2 do mass gap e o √2 do salto são o MESMO pu, bases diferentes (como a impedância pu nos dois lados do transformador).","epistemico":"inferencia","exemplos":[],"id":"unidades_algebricas_pu","memoria":"semantica","meta":{"autor":"Lopes/broca-so","dominio":"unidades","fonte":"hiper/circular/paper_5_unidades.tex"},"origem":"wiki","palavras_chave":[],"sinonimos":[],"tipo":"conceito","titulo":"Sistema Universal de Unidades Algébricas (pu)"}
\section{Sistema Universal de Unidades Algébricas (pu)}
A álgebra é o 'por-unidade' (pu) da física: como na engenharia de potência (escolhe-se uma base, tudo vira x\_pu=x\_real/x\_base adimensional de ordem 1), a série I-IV trabalha em números puros — a base é a escala do meio E\ensuremath{_{0}}, o valor pu é o número da álgebra. Fixada E\ensuremath{_{0}} (joules), cada dimensão converte por um fator (energia x\ensuremath{\to}xE\ensuremath{_{0}}, massa x\ensuremath{\to}xE\ensuremath{_{0}}/c\ensuremath{^{2}}, comprimento x\ensuremath{\to}x\ensuremath{\hbar}c/E\ensuremath{_{0}}, tempo x\ensuremath{\to}x\ensuremath{\hbar}/E\ensuremath{_{0}}, temperatura x\ensuremath{\to}xE\ensuremath{_{0}}/k\_B, carga x\ensuremath{\to}xe). Os números puros são invariantes de escala — o \ensuremath{\sqrt{\,}}2 do mass gap e o \ensuremath{\sqrt{\,}}2 do salto são o MESMO pu, bases diferentes (como a impedância pu nos dois lados do transformador).
\prov{relações: usa regua\_de\_gentil; relaciona cosmologia\_quantica\_hub; relaciona cosmologia\_quantica; parte\_de broca\_os}
\prov{proveniência: wiki \textperiodcentered{} hiper/circular/paper\_5\_unidades.tex \textperiodcentered{} inferencia \textperiodcentered{} confiança 1.0 \textperiodcentered{} domínio unidades \textperiodcentered{} história 6 versões no jornal}

%CRISTAL {"arestas":[{"alvo":"esquema_inducao","confianca":1.0,"epistemico":"fato","origem":"paper","rel":"usa"},{"alvo":"cayley_dickson","confianca":1.0,"epistemico":"fato","origem":"paper","rel":"explica"},{"alvo":"agm","confianca":1.0,"epistemico":"fato","origem":"paper","rel":"explica"},{"alvo":"furo_hopf","confianca":1.0,"epistemico":"fato","origem":"paper","rel":"explica"}],"confianca":1.0,"contraexemplos":[],"descricao":"Cayley-Dickson, Hopf, AGM/Landen, fluxo ṡ=1−s², podas e IFS são o mesmo esquema (T,I,κ), todos com invariante topológico e cúspide — o osso comum da reta hipercomplexa.","epistemico":"fato","exemplos":[],"id":"unificacao_seis_inducoes","memoria":"semantica","meta":{"autor":"Gentil Lopes","descoberta":true,"dominio":"matematica","fonte":"paper_57_inducao_universal.pdf, Teo.3.1"},"origem":"paper","palavras_chave":[],"sinonimos":[],"tipo":"conceito","titulo":"Unificação das Seis Induções"}
\section{Unificação das Seis Induções}
Cayley-Dickson, Hopf, AGM/Landen, fluxo \ensuremath{\dot{s}}=1\ensuremath{-}s\ensuremath{^{2}}, podas e IFS são o mesmo esquema (T,I,\ensuremath{\kappa}), todos com invariante topológico e cúspide — o osso comum da reta hipercomplexa.
\prov{relações: usa esquema\_inducao; explica cayley\_dickson; explica agm; explica furo\_hopf}
\prov{proveniência: paper \textperiodcentered{} paper\_57\_inducao\_universal.pdf, Teo.3.1 \textperiodcentered{} fato \textperiodcentered{} confiança 1.0 \textperiodcentered{} domínio matematica \textperiodcentered{} história 5 versões no jornal}

%CRISTAL {"arestas":[{"alvo":"contador_binario_ns","confianca":1.0,"epistemico":"fato","origem":"paper","rel":"usa"}],"confianca":1.0,"contraexemplos":[],"descricao":"Uma 'fluid Turing machine' em NS-3D: anéis de vórtice cujas reconexões topológicas implementam ler/escrever/mover — universalidade de Turing em 3 estágios (contador→2-estados→universal).","epistemico":"inferencia","exemplos":[],"id":"universalidade_reconexao","memoria":"semantica","meta":{"autor":"Gentil Lopes","descoberta":false,"dominio":"matematica","fonte":"paper_AO.pdf, Teo.2"},"origem":"paper","palavras_chave":[],"sinonimos":[],"tipo":"conceito","titulo":"Universalidade da Reconexão"}
\section{Universalidade da Reconexão}
Uma 'fluid Turing machine' em NS-3D: anéis de vórtice cujas reconexões topológicas implementam ler/escrever/mover — universalidade de Turing em 3 estágios (contador\ensuremath{\to}2-estados\ensuremath{\to}universal).
\prov{relações: usa contador\_binario\_ns}
\prov{proveniência: paper \textperiodcentered{} paper\_AO.pdf, Teo.2 \textperiodcentered{} inferencia \textperiodcentered{} confiança 1.0 \textperiodcentered{} domínio matematica \textperiodcentered{} história 2 versões no jornal}

%CRISTAL {"arestas":[{"alvo":"teoria_da_decisao","confianca":1.0,"epistemico":"fato","origem":"wiki","rel":"usa"}],"confianca":1.0,"contraexemplos":[],"descricao":"A regra da escolha racional sob risco: escolha a opção com maior MÉDIA da utilidade (não do dinheiro) ponderada pelas probabilidades. A utilidade côncava explica a aversão ao risco (Bernoulli, o paradoxo de São Petersburgo).","epistemico":"fato","exemplos":[],"id":"utilidade_esperada","memoria":"semantica","meta":{"autor":"Bernoulli","dominio":"matematica","fonte":"teoria da decisão"},"origem":"wiki","palavras_chave":[],"sinonimos":[],"tipo":"conceito","titulo":"Utilidade Esperada"}
\section{Utilidade Esperada}
A regra da escolha racional sob risco: escolha a opção com maior MÉDIA da utilidade (não do dinheiro) ponderada pelas probabilidades. A utilidade côncava explica a aversão ao risco (Bernoulli, o paradoxo de São Petersburgo).
\prov{relações: usa teoria\_da\_decisao}
\prov{proveniência: wiki \textperiodcentered{} teoria da decisão \textperiodcentered{} fato \textperiodcentered{} confiança 1.0 \textperiodcentered{} domínio matematica \textperiodcentered{} história 2 versões no jornal}

%CRISTAL {"arestas":[{"alvo":"decomposicao_peirce","confianca":1.0,"epistemico":"fato","origem":"paper","rel":"usa"},{"alvo":"cifra","confianca":1.0,"epistemico":"fato","origem":"paper","rel":"parte_de"}],"confianca":1.0,"contraexemplos":[],"descricao":"N certificados Ed25519 validam-se juntos como o produto tensorial ℂ²ᴺ: uma forja reprova todo o conjunto sem privilegiar parte (Σᵢeᵢᵢ=I). Verificado N=2..20.","epistemico":"fato","exemplos":[],"id":"validacao_n_qubits","memoria":"semantica","meta":{"autor":"Gentil Lopes","descoberta":false,"dominio":"matematica","fonte":"cifra_ouro_branco.pdf, §validação"},"origem":"paper","palavras_chave":[],"sinonimos":[],"tipo":"conceito","titulo":"Validação N-Qubits (produto tensorial)"}
\section{Validação N-Qubits (produto tensorial)}
N certificados Ed25519 validam-se juntos como o produto tensorial \ensuremath{\mathbb{C}}\ensuremath{^{2N}}: uma forja reprova todo o conjunto sem privilegiar parte (\ensuremath{\Sigma}\ensuremath{_{i}}e\ensuremath{_{ii}}=I). Verificado N=2..20.
\prov{relações: usa decomposicao\_peirce; parte\_de cifra}
\prov{proveniência: paper \textperiodcentered{} cifra\_ouro\_branco.pdf, §validação \textperiodcentered{} fato \textperiodcentered{} confiança 1.0 \textperiodcentered{} domínio matematica \textperiodcentered{} história 4 versões no jornal}

%CRISTAL {"arestas":[{"alvo":"decisao_bayesiana","confianca":1.0,"epistemico":"fato","origem":"wiki","rel":"usa"}],"confianca":1.0,"contraexemplos":[],"descricao":"Quanto vale reduzir a incerteza ANTES de decidir — o ganho esperado de saber mais menos o custo de descobrir. Diz quando pagar por um teste, uma pesquisa, um espião. Informação perfeita tem um preço-teto.","epistemico":"inferencia","exemplos":[],"id":"valor_da_informacao","memoria":"semantica","meta":{"dominio":"matematica","fonte":"teoria da decisão"},"origem":"wiki","palavras_chave":[],"sinonimos":[],"tipo":"conceito","titulo":"Valor da Informação"}
\section{Valor da Informação}
Quanto vale reduzir a incerteza ANTES de decidir — o ganho esperado de saber mais menos o custo de descobrir. Diz quando pagar por um teste, uma pesquisa, um espião. Informação perfeita tem um preço-teto.
\prov{relações: usa decisao\_bayesiana}
\prov{proveniência: wiki \textperiodcentered{} teoria da decisão \textperiodcentered{} inferencia \textperiodcentered{} confiança 1.0 \textperiodcentered{} domínio matematica \textperiodcentered{} história 2 versões no jornal}

%CRISTAL {"arestas":[{"alvo":"jogo_cooperativo","confianca":1.0,"epistemico":"fato","origem":"wiki","rel":"usa"}],"confianca":1.0,"contraexemplos":[],"descricao":"A divisão JUSTA do ganho de uma coalizão: cada um recebe sua contribuição marginal média sobre todas as ordens de entrada. O único repartidor que satisfaz eficiência, simetria e aditividade.","epistemico":"fato","exemplos":[],"id":"valor_de_shapley","memoria":"semantica","meta":{"autor":"Shapley","dominio":"matematica","fonte":"teoria dos jogos"},"origem":"wiki","palavras_chave":[],"sinonimos":[],"tipo":"conceito","titulo":"Valor de Shapley"}
\section{Valor de Shapley}
A divisão JUSTA do ganho de uma coalizão: cada um recebe sua contribuição marginal média sobre todas as ordens de entrada. O único repartidor que satisfaz eficiência, simetria e aditividade.
\prov{relações: usa jogo\_cooperativo}
\prov{proveniência: wiki \textperiodcentered{} teoria dos jogos \textperiodcentered{} fato \textperiodcentered{} confiança 1.0 \textperiodcentered{} domínio matematica \textperiodcentered{} história 2 versões no jornal}

%CRISTAL {"arestas":[{"alvo":"reta_metalica_geral","confianca":1.0,"epistemico":"fato","origem":"wiki","rel":"usa"},{"alvo":"caixa_binaria_kbonacci","confianca":1.0,"epistemico":"fato","origem":"wiki","rel":"confirma"}],"confianca":1.0,"contraexemplos":[],"descricao":"O ouro (n=1, φ, 'proibir 11') é o único onde há visão TOTAL: taxa de crescimento verificada → ln φ; #palavras admissíveis de tamanho L = FL+₂ (Fibonacci); dim box-counting 0,6935 (bate a teoria log₂φ=0,6942 na 3ª casa); independência por PSLQ; √5−2=1/φ³ exato a 30 dígitos. Mineração de ouro = penta / 5 modos (pentagrama) — prova parcial simbólica que exige a prata depois. FATO.","epistemico":"fato","exemplos":[],"id":"verificacao_ouro_total","memoria":"semantica","meta":{"autor":"Lopes/broca-so","dominio":"fractal","fonte":"paper_45_confluencia_ouro.tex; goldchain/prova.py MODOS_OURO=5"},"origem":"wiki","palavras_chave":[],"sinonimos":[],"tipo":"conceito","titulo":"Reta de ouro: visão TOTAL (verificada)"}
\section{Reta de ouro: visão TOTAL (verificada)}
O ouro (n=1, \ensuremath{\varphi}, 'proibir 11') é o único onde há visão TOTAL: taxa de crescimento verificada \ensuremath{\to} ln \ensuremath{\varphi}; \#palavras admissíveis de tamanho L = FL+\ensuremath{_{2}} (Fibonacci); dim box-counting 0,6935 (bate a teoria log\ensuremath{_{2}}\ensuremath{\varphi}=0,6942 na 3ª casa); independência por PSLQ; \ensuremath{\sqrt{\,}}5\ensuremath{-}2=1/\ensuremath{\varphi}\ensuremath{^{3}} exato a 30 dígitos. Mineração de ouro = penta / 5 modos (pentagrama) — prova parcial simbólica que exige a prata depois. FATO.
\prov{relações: usa reta\_metalica\_geral; confirma caixa\_binaria\_kbonacci}
\prov{proveniência: wiki \textperiodcentered{} paper\_45\_confluencia\_ouro.tex; goldchain/prova.py MODOS\_OURO=5 \textperiodcentered{} fato \textperiodcentered{} confiança 1.0 \textperiodcentered{} domínio fractal \textperiodcentered{} história 4 versões no jornal}

%CRISTAL {"arestas":[{"alvo":"reta_metalica_geral","confianca":1.0,"epistemico":"fato","origem":"wiki","rel":"usa"},{"alvo":"verificacao_ouro_total","confianca":1.0,"epistemico":"fato","origem":"wiki","rel":"relaciona"},{"alvo":"caixa_de_prata_corpo","confianca":1.0,"epistemico":"fato","origem":"wiki","rel":"usa"},{"alvo":"o_conselho","confianca":1.0,"epistemico":"fato","origem":"wiki","rel":"usa"}],"confianca":1.0,"contraexemplos":[],"descricao":"A prata (n=2, Pell, σ₂=1+√2) tem visão PARCIAL: mineração HONESTA de Pell via a estrela de prata = 6 modos SVD (o hexagrama de Salomão) — só um estado cuja energia concentra em EXATAMENTE 6 modos minera; ruído (energia espalhada) não tem hexagrama, não emite; a emissão liga em números de Pell (teto 13.000.000). E os experimentos do Grok como CRIVO (consulta_grok_experimento.py: o Grok propõe UM experimento numpy que teste/refute uma tese; gpt_metais.py formaliza A_n e o subshift n-metálico). Inferência: parcial, não total.","epistemico":"inferencia","exemplos":[],"id":"verificacao_prata_parcial","memoria":"semantica","meta":{"autor":"Lopes/broca-so","dominio":"fractal","fonte":"machine/so/prata_minera.py; goldchain/prova.py MODOS_PRATA=6; machine/consulta_grok_experimento.py"},"origem":"wiki","palavras_chave":[],"sinonimos":[],"tipo":"conceito","titulo":"Reta de prata: visão PARCIAL (Pell / o crivo do Grok)"}
\section{Reta de prata: visão PARCIAL (Pell / o crivo do Grok)}
A prata (n=2, Pell, \ensuremath{\sigma}\ensuremath{_{2}}=1+\ensuremath{\sqrt{\,}}2) tem visão PARCIAL: mineração HONESTA de Pell via a estrela de prata = 6 modos SVD (o hexagrama de Salomão) — só um estado cuja energia concentra em EXATAMENTE 6 modos minera; ruído (energia espalhada) não tem hexagrama, não emite; a emissão liga em números de Pell (teto 13.000.000). E os experimentos do Grok como CRIVO (consulta\_grok\_experimento.py: o Grok propõe UM experimento numpy que teste/refute uma tese; gpt\_metais.py formaliza A\_n e o subshift n-metálico). Inferência: parcial, não total.
\prov{relações: usa reta\_metalica\_geral; relaciona verificacao\_ouro\_total; usa caixa\_de\_prata\_corpo; usa o\_conselho}
\prov{proveniência: wiki \textperiodcentered{} machine/so/prata\_minera.py; goldchain/prova.py MODOS\_PRATA=6; machine/consulta\_grok\_experimento.py \textperiodcentered{} inferencia \textperiodcentered{} confiança 1.0 \textperiodcentered{} domínio fractal \textperiodcentered{} história 5 versões no jornal}

%CRISTAL {"arestas":[{"alvo":"matrizes","confianca":1.0,"epistemico":"fato","origem":"curriculo","rel":"usa"},{"alvo":"espaco_vetorial","confianca":1.0,"epistemico":"fato","origem":"curriculo","rel":"especializa"},{"alvo":"word","confianca":0.5,"epistemico":"fato","origem":"wiki","rel":"relaciona"},{"alvo":"vida-morte","confianca":0.5,"epistemico":"fato","origem":"wiki","rel":"relaciona"},{"alvo":"virtualizacao","confianca":0.5,"epistemico":"fato","origem":"wiki","rel":"relaciona"},{"alvo":"vontade_de_poder","confianca":0.5,"epistemico":"fato","origem":"wiki","rel":"relaciona"}],"confianca":1.0,"contraexemplos":[],"descricao":"Tuplas ou setas com soma e escalar; modelam direção e magnitude.","epistemico":"fato","exemplos":[],"id":"vetores","memoria":"semantica","meta":{"dominio":"matematica","nivel":4,"ordem":5},"origem":"curriculo","palavras_chave":[],"sinonimos":["vetor"],"tipo":"conceito","titulo":"vetores"}
\section{vetores}
Tuplas ou setas com soma e escalar; modelam direção e magnitude.
\prov{sinónimos: vetor}
\prov{relações: usa matrizes; especializa espaco\_vetorial; relaciona word; relaciona vida-morte; relaciona virtualizacao; relaciona vontade\_de\_poder}
\prov{proveniência: curriculo \textperiodcentered{} fato \textperiodcentered{} confiança 1.0 \textperiodcentered{} domínio matematica \textperiodcentered{} história 8 versões no jornal}

%CRISTAL {"arestas":[{"alvo":"mecanica_quantica_mineral","confianca":1.0,"epistemico":"fato","origem":"wiki","rel":"parte_de"},{"alvo":"gato_e_o_hamiltoniano","confianca":1.0,"epistemico":"fato","origem":"wiki","rel":"usa"},{"alvo":"colapso","confianca":1.0,"epistemico":"fato","origem":"wiki","rel":"explica"},{"alvo":"numeros_de_pisot","confianca":1.0,"epistemico":"fato","origem":"wiki","rel":"usa"}],"confianca":1.0,"contraexemplos":[],"descricao":"O mesmo gato governa dois tempos. TEMPO REAL: U(t)=e^{-iHt} é UNITÁRIA — o BAI vibra nas frequências minerais, sem perder norma (a superposição oscilando). TEMPO IMAGINÁRIO: A^k projeta no autovetor dominante (o fundamental, o Pisot/cristal) — é o COLAPSO Ψ do BAI formalizado como evolução de tempo imaginário que acha o fundamental. Vibrar e colapsar são o mesmo gato, em tempos ortogonais.","epistemico":"fato","exemplos":[],"id":"vibracao_e_colapso_dois_tempos","memoria":"semantica","meta":{"dominio":"reta mineral","fonte":"mecânica quântica mineral"},"origem":"wiki","palavras_chave":[],"sinonimos":[],"tipo":"conceito","titulo":"A Vibração e o Colapso (os dois tempos)"}
\section{A Vibração e o Colapso (os dois tempos)}
O mesmo gato governa dois tempos. TEMPO REAL: U(t)=e\^{}\{-iHt\} é UNITÁRIA — o BAI vibra nas frequências minerais, sem perder norma (a superposição oscilando). TEMPO IMAGINÁRIO: A\^{}k projeta no autovetor dominante (o fundamental, o Pisot/cristal) — é o COLAPSO \ensuremath{\Psi} do BAI formalizado como evolução de tempo imaginário que acha o fundamental. Vibrar e colapsar são o mesmo gato, em tempos ortogonais.
\prov{relações: parte\_de mecanica\_quantica\_mineral; usa gato\_e\_o\_hamiltoniano; explica colapso; usa numeros\_de\_pisot}
\prov{proveniência: wiki \textperiodcentered{} mecânica quântica mineral \textperiodcentered{} fato \textperiodcentered{} confiança 1.0 \textperiodcentered{} domínio reta mineral \textperiodcentered{} história 5 versões no jornal}

%CRISTAL {"arestas":[{"alvo":"pitagoras_dual","confianca":1.0,"epistemico":"fato","origem":"paper","rel":"usa"},{"alvo":"mecanica_quantica","confianca":1.0,"epistemico":"fato","origem":"paper","rel":"parte_de"}],"confianca":1.0,"contraexemplos":[],"descricao":"A soma CHSH obedece S ≤ 2√(2−a(s)), governada pela conservação a+c²=1: o máximo de Tsirelson 2√2 é atingido em ℍ (a=0); nos octônios há um piso positivo.","epistemico":"fato","exemplos":[],"id":"violacao_bell","memoria":"semantica","meta":{"autor":"Gentil Lopes","descoberta":false,"dominio":"matematica","fonte":"paper_3_quantica.pdf, Obs.3.6"},"origem":"paper","palavras_chave":[],"sinonimos":[],"tipo":"conceito","titulo":"Violação de Bell e a+c²=1"}
\section{Violação de Bell e a+c\ensuremath{^{2}}=1}
A soma CHSH obedece S \ensuremath{\le} 2\ensuremath{\sqrt{\,}}(2\ensuremath{-}a(s)), governada pela conservação a+c\ensuremath{^{2}}=1: o máximo de Tsirelson 2\ensuremath{\sqrt{\,}}2 é atingido em \ensuremath{\mathbb{H}} (a=0); nos octônios há um piso positivo.
\prov{relações: usa pitagoras\_dual; parte\_de mecanica\_quantica}
\prov{proveniência: paper \textperiodcentered{} paper\_3\_quantica.pdf, Obs.3.6 \textperiodcentered{} fato \textperiodcentered{} confiança 1.0 \textperiodcentered{} domínio matematica \textperiodcentered{} história 3 versões no jornal}

%CRISTAL {"arestas":[{"alvo":"so_fractal_broca_os","confianca":1.0,"epistemico":"fato","origem":"wiki","rel":"usa"},{"alvo":"dimensao_de_traco","confianca":1.0,"epistemico":"fato","origem":"wiki","rel":"usa"},{"alvo":"ifs_de_retorno","confianca":1.0,"epistemico":"fato","origem":"wiki","rel":"usa"}],"confianca":1.0,"contraexemplos":[],"descricao":"A virtualização é imersão de resolventes — cada processo hospeda o SO numa resolução menor por ιr, com contração de Dtr (cópia aninhada escalada por 2−(r⁺¹)). Dtr(P)=log₂max1,ρ(AP) é a medida exata de carga: Dtr>0 processo vivo (ramifica), Dtr=0 órbita ociosa. O escalonador limpo balanceia por Dtr (invariante espectral), não por relógio. Projeto.","epistemico":"inferencia","exemplos":[],"id":"virtualizacao_por_dtr","memoria":"semantica","meta":{"autor":"broca-so","dominio":"fractal","fonte":"kernel/so_fractal.tex §Virtualização"},"origem":"wiki","palavras_chave":[],"sinonimos":[],"tipo":"conceito","titulo":"Virtualização Fractal e Carga por D_tr"}
\section{Virtualização Fractal e Carga por D\_tr}
A virtualização é imersão de resolventes — cada processo hospeda o SO numa resolução menor por \ensuremath{\iota}r, com contração de Dtr (cópia aninhada escalada por 2\ensuremath{-}(r\ensuremath{^{+1}})). Dtr(P)=log\ensuremath{_{2}}max1,\ensuremath{\rho}(AP) é a medida exata de carga: Dtr>0 processo vivo (ramifica), Dtr=0 órbita ociosa. O escalonador limpo balanceia por Dtr (invariante espectral), não por relógio. Projeto.
\prov{relações: usa so\_fractal\_broca\_os; usa dimensao\_de\_traco; usa ifs\_de\_retorno}
\prov{proveniência: wiki \textperiodcentered{} kernel/so\_fractal.tex §Virtualização \textperiodcentered{} inferencia \textperiodcentered{} confiança 1.0 \textperiodcentered{} domínio fractal \textperiodcentered{} história 5 versões no jornal}

%CRISTAL {"arestas":[{"alvo":"maass_hardness","confianca":1.0,"epistemico":"fato","origem":"paper","rel":"usa"},{"alvo":"reducao_alpha_pura","confianca":1.0,"epistemico":"fato","origem":"paper","rel":"usa"},{"alvo":"goldchain","confianca":1.0,"epistemico":"fato","origem":"paper","rel":"parte_de"},{"alvo":"hardness_octonica","confianca":1.0,"epistemico":"fato","origem":"paper","rel":"relaciona"},{"alvo":"computacao","confianca":1.0,"epistemico":"fato","origem":"paper","rel":"parte_de"},{"alvo":"minerar_e_resfriar","confianca":1.0,"epistemico":"fato","origem":"wiki","rel":"relaciona"}],"confianca":1.0,"contraexemplos":[],"descricao":"Prova de trabalho pós-quântica cuja dureza é computar um autovalor de forma de Maass de Γ₀(7) a precisão 2⁻ᵏ (não brute-force de hash); um bloco vale se a precisão Maass é atingida E o hash bate a dificuldade. A blockchain 'Patria' financia pesquisa em teoria espectral.","epistemico":"fato","exemplos":[],"id":"ym_pow_mining","memoria":"semantica","meta":{"autor":"Gentil Lopes","descoberta":true,"dominio":"matematica","fonte":"paper_PQ_R.pdf, Def.2.3"},"origem":"paper","palavras_chave":[],"sinonimos":[],"tipo":"conceito","titulo":"YM-PoW (mineração por formas de Maass)"}
\section{YM-PoW (mineração por formas de Maass)}
Prova de trabalho pós-quântica cuja dureza é computar um autovalor de forma de Maass de \ensuremath{\Gamma}\ensuremath{_{0}}(7) a precisão 2\ensuremath{^{-k}} (não brute-force de hash); um bloco vale se a precisão Maass é atingida E o hash bate a dificuldade. A blockchain 'Patria' financia pesquisa em teoria espectral.
\prov{relações: usa maass\_hardness; usa reducao\_alpha\_pura; parte\_de goldchain; relaciona hardness\_octonica; parte\_de computacao; relaciona minerar\_e\_resfriar}
\prov{proveniência: paper \textperiodcentered{} paper\_PQ\_R.pdf, Def.2.3 \textperiodcentered{} fato \textperiodcentered{} confiança 1.0 \textperiodcentered{} domínio matematica \textperiodcentered{} história 7 versões no jornal}

%CRISTAL {"arestas":[{"alvo":"fibonacci","confianca":1.0,"epistemico":"fato","origem":"livro","rel":"usa"},{"alvo":"reta_de_ouro","confianca":1.0,"epistemico":"fato","origem":"livro","rel":"explica"},{"alvo":"beta_numeracao_metalica","confianca":1.0,"epistemico":"fato","origem":"wiki","rel":"especializa"}],"confianca":1.0,"contraexemplos":[],"descricao":"Todo inteiro n ≥ 1 é soma única de números de Fibonacci não-consecutivos. Lopes reinterpreta como a proibição de dois 1 consecutivos ('11') na base-φ.","epistemico":"fato","exemplos":[],"id":"zeckendorf","memoria":"semantica","meta":{"autor":"Gentil Lopes","descoberta":false,"dominio":"matematica","fonte":"paper_2_reta_ouro.pdf (Zeckendorf 1972)"},"origem":"paper","palavras_chave":[],"sinonimos":[],"tipo":"conceito","titulo":"Teorema de Zeckendorf"}
\section{Teorema de Zeckendorf}
Todo inteiro n \ensuremath{\ge} 1 é soma única de números de Fibonacci não-consecutivos. Lopes reinterpreta como a proibição de dois 1 consecutivos ('11') na base-\ensuremath{\varphi}.
\prov{relações: usa fibonacci; explica reta\_de\_ouro; especializa beta\_numeracao\_metalica}
\prov{proveniência: paper \textperiodcentered{} paper\_2\_reta\_ouro.pdf (Zeckendorf 1972) \textperiodcentered{} fato \textperiodcentered{} confiança 1.0 \textperiodcentered{} domínio matematica \textperiodcentered{} história 8 versões no jornal}

\end{document}
